\documentclass{missingbook}
\usepackage{wild}

\begin{document}

\frontmatter
\hypersetup{pdftitle={Optimizers in the Wild: Steepest Descent and the Dynamics of Training},pdfauthor={Written by Claude; conceived and edited by Xiaoyu Li}}
\mbtitlepage{Optimizers in the Wild}{Steepest Descent and the Dynamics of Training}{Written by Claude\\[2.5mm]Conceived and edited by Xiaoyu Li}{First edition \textperiodcentered{} October 2026}
\addchap{Preface}

This book began with a list of textbooks that should exist but do not. One entry read, in effect: textbooks on optimization end with convex analysis and with the theorem that stochastic gradient descent finds a stationary point, while what decides whether a large network trains well is something else entirely: Adam, warmup, weight decay, the learning-rate schedule, Muon, and the edge of stability. The theory of these practices exists, but it is scattered across several hundred papers, blog posts and theses, written in a dozen notations, and it has never been assembled into a book that a graduate student can learn from.

The book that resulted follows two lines. The first is a unification: almost every optimizer used to train neural networks is steepest descent under some norm, combined with moving averages and preconditioning. SGD, sign descent, Adam, Lion, Shampoo and Muon are all instances, and the norm of a whole network, assembled from the norms of its layers, tells how the learning rate must scale with width. This line was developed by Jeremy Bernstein, Laker Newhouse and their coauthors, and Newhouse's thesis at the Massachusetts Institute of Technology gives a connected account of it; a thesis, however, is not a textbook, and it does not treat the second line. The second line is about dynamics: what happens when these optimizers run with the step sizes that practice uses. Sharpness rises until the step size is barely stable and then stabilizes itself; warmup is needed because the first steps meet too much curvature; gradient noise sets a floor that a schedule removes, which explains the shapes of cosine and warmup--stable--decay schedules; weight decay drives normalized layers to an equilibrium; and the reasons that Adam beats gradient descent on transformers turn out to involve both lines at once.

Each of the fifteen chapters starts from a practice, a rule or default as it is used ``in the wild'', builds the smallest idealization in which the practice can be analysed, proves the main theorem of that idealization completely, compares its predictions with a simulation, and then says what the idealization misses and what the theory does and does not tell a practitioner to do. All chapters share one notation, summarized on the following pages, and each ends by translating the notation of its main sources into that of the book. An appendix collects the mathematics the chapters share.

The intended readers are graduate students and researchers in machine learning, optimization, statistics and applied mathematics, and practitioners who train large models and want to know which of their rules are understood, in what sense, and how far the understanding reaches. The prerequisites are undergraduate analysis, linear algebra and probability, and some acquaintance with neural networks. The prologue explains the plan of the book and works its first example.

The book is the third volume of \emph{The Missing Textbooks}. Its companion, \emph{The Hydrogen Atoms of Deep Learning}, collects solvable models of trained networks; one of them, the two-parameter product model of the edge of stability, is solved there and cited here rather than repeated.

A theory of practice has to be honest about how much of practice it explains. Some practices in this book are explained well, some partly, and some not at all, and the chapters say which. If the book succeeds, a reader will know not only why the recipes work where they do, but also exactly where the explanations stop.

\paragraph{A note on authorship.} This book was written by Claude, an AI model made by Anthropic, and conceived and edited by Xiaoyu Li, who proposed the subject, its two lines, and its position relative to existing work. The chapters were drafted from a common plan, their proofs re-derived and their numerical claims checked against simulations, and their references checked against sources wherever these could be reached; the colophon at the end of the book describes how it was made. Mistakes surely remain. Readers who find one are invited to treat it as the book treats its idealizations: as a precise statement of what is missed, and a pointer to what comes next.

\mbtableofcontents
\addchap{Notation}

One notation is used throughout the book. Where a chapter needs a local symbol that clashes with this list, the chapter says so; where a source paper uses a different notation, the chapter's notes translate it.

\newcommand{\notationrow}[2]{#1 & #2\\}
\newcommand{\notationhead}[1]{\multicolumn{2}{@{}l}{\textsc{\color{sea}#1}}\\[0.3ex]}
\noindent{\small
\renewcommand{\arraystretch}{1.1}%
\begin{tabular}{@{}>{\raggedright\arraybackslash}p{32mm}>{\raggedright\arraybackslash}p{84mm}@{}}
\notationhead{sizes}
\notationrow{$d$}{number of parameters, the dimension of the optimization variable}
\notationrow{$W\in\R^{d_{\mathrm{out}}\times d_{\mathrm{in}}}$}{a weight matrix; $G\in\R^{m\times n}$ a generic matrix or matrix gradient}
\notationrow{$N$, $l$}{number of layers and the index of a layer; widths $d_0,\dots,d_N$}
\notationrow{$n$, $B$, $K$}{number of samples, batch size, number of classes}
\notationrow{$T$, $t$}{number of steps and the step, counted from $t=0$; $t$ is continuous for flows; the online-learning arguments of \cref{ch:adaptive,ch:kronecker,ch:schedules} count rounds from $t=1$, as their sources do}
\addlinespace
\notationhead{objective}
\notationrow{$\Loss(\theta)$}{the objective: the training or the population loss, as each chapter says}
\notationrow{$\ell(\theta;\xi)$}{the loss on one example or one batch $\xi$}
\notationrow{$\Loss^\star$, $\theta^\star$}{the infimum of $\Loss$ and a minimizer; a superscript star marks an optimal value}
\notationrow{$\Delta_t$}{the suboptimality $\Loss(\theta_t)-\Loss^\star$}
\notationrow{$\theta_\star$}{a subscript star marks a target or a teacher}
\addlinespace
\notationhead{gradients and curvature}
\notationrow{$g_t$, $\zeta_t$}{a stochastic gradient at $\theta_t$, and its noise $g_t-\nabla\Loss(\theta_t)$}
\notationrow{$\Sigma(\theta)$}{the per-example covariance of the gradient, so that $\Cov(g_t)=\Sigma/B$}
\notationrow{$H$, $S$}{the Hessian $\nabla^2\Loss$ and the sharpness $S=\lmax(H)$}
\notationrow{$S_P$}{the preconditioned sharpness $\lmax(P^{-1}H)$}
\addlinespace
\notationhead{optimizers}
\notationrow{$\eta$, $\eta_t$}{step size (learning rate) and its schedule; $\eta_{\max}$ the peak, $T_{\mathrm w}$ the warmup length}
\notationrow{$\beta$; $\beta_1$, $\beta_2$, $\epsilon$}{momentum; Adam's averaging coefficients and its $\epsilon$}
\notationrow{$m_t$, $v_t$; $\hat m_t$, $\hat v_t$}{Adam's moving averages of $g$ and $g^2$, and their bias-corrected versions}
\notationrow{$\lambda$}{decoupled weight decay, $\theta\leftarrow(1-\eta\lambda)\theta-\eta u$; eigenvalues always carry a subscript, as in $\lambda_i$, $\lmax$, $\lmin$}
\end{tabular}

\medskip\noindent
\begin{tabular}{@{}>{\raggedright\arraybackslash}p{32mm}>{\raggedright\arraybackslash}p{84mm}@{}}
\notationhead{norms}
\notationrow{$\norm{\cdot}$, $\norm{\cdot}_*$}{a generic norm and its dual norm; an unadorned norm is never silently Euclidean}
\notationrow{$\norm{\cdot}_2$, $\norm{\cdot}_1$, $\norm{\cdot}_\infty$}{the $\ell_2$, $\ell_1$ and $\ell_\infty$ norms of a vector}
\notationrow{$\norm{W}_{\op}$, $\norm{W}_F$, $\norm{W}_{\nuc}$}{spectral, Frobenius and nuclear norms of a matrix; $\norm{W}_{p\to q}$ an induced norm}
\notationrow{$\norm{x}_{\rms}$}{the root-mean-square norm $\norm{x}_2/\sqrt{\dim x}$; $\norm{W}_{\rms\to\rms}=\sqrt{d_{\mathrm{in}}/d_{\mathrm{out}}}\,\norm{W}_{\op}$}
\notationrow{$\norm{x}_P$}{$\sqrt{x^\top Px}$ for a positive definite $P$}
\addlinespace
\notationhead{steepest descent}
\notationrow{$\lmo(g)$}{a minimizer of $\ip{g,u}$ over $\norm{u}\le1$, so that $\ip{g,\lmo(g)}=-\norm{g}_*$}
\notationrow{$g^\sharp$}{the sharp operator $\argmax_u\{\ip{g,u}-\tfrac12\norm{u}^2\}=-\norm{g}_*\lmo(g)$}
\notationrow{normalized step}{$\theta_{t+1}=\theta_t+\eta\lmo(g_t)$: a step of length $\eta$ in the chosen norm}
\notationrow{unnormalized step}{$\theta_{t+1}=\theta_t-\eta g_t^\sharp$}
\notationrow{$\sign$, $\polar$}{coordinatewise sign with $\sign0=0$; $\polar(G)=UV^\top$ for a compact singular value decomposition $G=U\diag(\sigma)V^\top$}
\addlinespace
\notationhead{smoothness}
\notationrow{$L$-smooth}{$\norm{\nabla\Loss(\theta)-\nabla\Loss(\theta')}_*\le L\norm{\theta-\theta'}$, with respect to a stated norm}
\notationrow{$(L_0,L_1)$}{$\norm{\nabla^2\Loss(\theta)}\le L_0+L_1\norm{\nabla\Loss(\theta)}$ (\cref{ch:spikes})}
\notationrow{$(K_0,K_1)$}{$\norm{\nabla^2\Loss(\theta)}\le K_0+K_1(\Loss(\theta)-\Loss^\star)$ (\cref{ch:warmup})}
\addlinespace
\notationhead{matrices, probability, asymptotics}
\notationrow{$\otimes$, $\vect$, $\preceq$}{Kronecker product, column-stacking vectorization, Löwner order}
\notationrow{$\cN$, $\E$, $\Var$, $\Cov$}{Gaussian law, expectation, variance, covariance}
\notationrow{$a\lesssim b$, $a\asymp b$}{$a\le Cb$, and $a\lesssim b\lesssim a$, for a constant $C$ depending only on parameters declared fixed}
\addlinespace
\notationhead{status of an argument}
\notationrow{$\blacksquare$}{end of a complete proof}
\notationrow{proof sketch}{an argument that omits steps, with a pointer to where they are}
\notationrow{$\Diamond$}{end of a heuristic derivation: an argument that is not a proof}
\notationrow{observation}{an empirical regularity, with its evidence}
\notationrow{folklore}{a rule used in practice for which no evidence can be cited}
\end{tabular}\par}


\mainmatter
% Prologue of Optimizers in the Wild.
\specialchapter{Prologue}{What the textbook proves and what the practitioner decides}{4}\label{ch:prologue}
\begin{refsection}[chapters/00-prologue.bib]

A course on optimization ends with a theorem like this one: stochastic gradient descent on a smooth function, with a small enough step size, finds a point where the gradient is small. The theorem is true, its proof fits on half a page, and it is silent about almost every decision that determines whether a large neural network trains well. Which optimizer? Adam, or the newer Muon? What learning rate (the practitioner's name for the step size), and should it change with the width of the network? How long a warmup, and why one at all? Cosine decay or a constant rate with a short cooldown? How much weight decay? How large a batch? Should gradients be clipped? These questions are settled every day by rules of thumb, by inherited recipes, and by expensive sweeps. This book is about the theory behind them.

\addsec{What the textbook proves}

The standard guarantee for nonconvex stochastic optimization goes back to \textcite{GhadimiLan2013}; its ancestor for root finding is the stochastic approximation of \textcite{RobbinsMonro1951}. We state it in the form that the rest of the book will generalize. Let $\Loss\colon\R^d\to\R$ be bounded below, with infimum $\Loss^\star$, and $L$-smooth in the Euclidean norm, which means that $\norm{\nabla\Loss(\theta)-\nabla\Loss(\theta')}_2\le L\norm{\theta-\theta'}_2$ for all $\theta,\theta'$. Stochastic gradient descent is the iteration $\theta_{t+1}=\theta_t-\eta g_t$, where $g_t$ is an unbiased estimate of $\nabla\Loss(\theta_t)$ whose noise $\zeta_t=g_t-\nabla\Loss(\theta_t)$ satisfies $\E[\norm{\zeta_t}_2^2\mid\theta_t]\le\sigma^2$.

\begin{theorem*}[Stochastic gradient descent finds approximate stationary points]
Let $\Delta_0=\Loss(\theta_0)-\Loss^\star$. If $0<\eta\le1/L$, then for every $T\ge1$,
\[
  \frac1T\sum_{t=0}^{T-1}\E\norm{\nabla\Loss(\theta_t)}_2^2\le\frac{2\Delta_0}{\eta T}+L\eta\sigma^2.
\]
If $\Delta_0>0$, then the step size $\eta=\min\{1/L,\sqrt{2\Delta_0/(L\sigma^2T)}\}$, which is $1/L$ when $\sigma=0$, makes the right side at most $4L\Delta_0/T+2\sigma\sqrt{2L\Delta_0/T}$.
\end{theorem*}

\begin{proof}
The plan is to show that each step decreases the loss in expectation by $\frac\eta2\norm{\nabla\Loss}_2^2$, up to a noise term, and to add up the decreases. Smoothness gives the descent inequality $\Loss(\theta')\le\Loss(\theta)+\ip{\nabla\Loss(\theta),\theta'-\theta}+\frac L2\norm{\theta'-\theta}_2^2$ (\cref{lem:toolkit:descent}). Apply it with $\theta=\theta_t$ and $\theta'=\theta_{t+1}=\theta_t-\eta g_t$ and take the conditional expectation given $\theta_t$. Because $\E[g_t\mid\theta_t]=\nabla\Loss(\theta_t)$ and $\E[\norm{g_t}_2^2\mid\theta_t]\le\norm{\nabla\Loss(\theta_t)}_2^2+\sigma^2$, and because $L\eta\le1$ in the second step, this gives
\[
\begin{aligned}
  \E[\Loss(\theta_{t+1})\mid\theta_t]&\le\Loss(\theta_t)-\eta\Bigl(1-\frac{L\eta}2\Bigr)\norm{\nabla\Loss(\theta_t)}_2^2+\frac{L\eta^2\sigma^2}2\\
  &\le\Loss(\theta_t)-\frac\eta2\norm{\nabla\Loss(\theta_t)}_2^2+\frac{L\eta^2\sigma^2}2.
\end{aligned}
\]
Taking expectations, summing over $t<T$, and using $\E\Loss(\theta_T)\ge\Loss^\star$ gives $\frac\eta2\sum_{t<T}\E\norm{\nabla\Loss(\theta_t)}_2^2\le\Delta_0+\frac{TL\eta^2\sigma^2}2$, which is the first claim after division by $\eta T/2$. For the second, there are two cases. If $\sigma>0$ and $\eta=\sqrt{2\Delta_0/(L\sigma^2T)}\le1/L$, then both terms equal $\sigma\sqrt{2L\Delta_0/T}$. Otherwise $\eta=1/L$ and $\sigma^2<2L\Delta_0/T$ (when $\sigma=0$, because $\Delta_0>0$), so the bound is at most $2L\Delta_0/T+\sigma^2<4L\Delta_0/T$. Each case gives at most the stated sum.
\end{proof}

Read as a description of practice, the theorem has four gaps. It measures progress in the Euclidean norm, and its only knob is a constant step size below $1/L$, while the optimizers that train large models take steps of a fixed length in other norms, average their gradients, and rescale each coordinate or each matrix by a running estimate of its scale. It assumes a global smoothness constant, while the curvature of a network changes during training and, as we shall see, settles at whatever value the step size allows. It treats the noise as a nuisance to be averaged away, while the noise sets a floor under the loss that the learning-rate schedule is designed to remove. And it asks for a stationary point, while practitioners ask for the lowest loss after a fixed budget of steps, and care which of many minimizers they reach.

\addsec{What the practitioner decides}

Here is a training recipe of the kind that was standard for large language models between 2020 and 2025. Its details come from the published descriptions of GPT-3 \parencite{BrownEtAl2020} and LLaMA \parencite{TouvronEtAl2023}, and the box says which paper gives which.

\begin{practice}[GPT-3, \textcite{BrownEtAl2020}; LLaMA, \textcite{TouvronEtAl2023}]
Train with Adam, $\beta_1=0.9$, $\beta_2=0.95$, and weight decay $0.1$ (both papers; LLaMA calls the optimizer AdamW, and GPT-3 cites \textcite{LoshchilovHutter2019} for the decay), with $\epsilon=10^{-8}$ (GPT-3; LLaMA does not state it). Warm the learning rate up, linearly over the first $375$ million tokens in GPT-3, and over $2000$ steps of unstated shape in LLaMA; then decay it along a cosine to $10\%$ of its peak (both). Clip the gradient at $1.0$ (both; GPT-3 clips the global norm). Use batches of millions of tokens: $0.5$ to $3.2$ million in GPT-3, by model size, after a ramp from $32$ thousand, and $4$ million in LLaMA. Lower the peak learning rate as the model grows: from $6\times10^{-4}$ to $0.6\times10^{-4}$ across the GPT-3 models.
\end{practice}

Each clause of this recipe is a decision with an alternative, and each alternative has its advocates. Since 2024 a growing number of runs replace AdamW on the hidden weight matrices by Muon, which orthogonalizes each matrix's momentum before stepping \parencite{JordanEtAl2024Muon,LiuEtAl2025Muon}; many replace the cosine by a warmup--stable--decay schedule that holds the learning rate constant and cools down only at the end (\cref{ch:schedules}); and $\mu$P and its successors prescribe how the learning rate of each layer must scale with width so that a sweep on a small model transfers to a large one (\cref{ch:modular}). The \emph{Deep Learning Tuning Playbook} \parencite{GodboleEtAl2023} is a systematic guide to making these decisions empirically. This book asks what can be said about them mathematically.

The questions, and the chapters that answer them, are these.
\begin{itemize}
\item \emph{Why do the update rules look the way they do?} AdamW and Muon, like the other optimizers of this book, are steepest descent under a norm, combined with moving averages and, in AdamW's case, a preconditioner (\cref{ch:steepest,ch:sign,ch:spectral,ch:momentum,ch:adaptive,ch:kronecker}).
\item \emph{How should the learning rate scale with width?} The norm of a whole network decides it (\cref{ch:modular}).
\item \emph{Why can the learning rate be so large?} Training runs at the edge of stability, where the curvature adapts to the step size (\cref{ch:eos}).
\item \emph{Why warm up?} The first steps meet curvature too large for the target learning rate (\cref{ch:warmup}).
\item \emph{Why clip, and where do loss spikes come from?} Curvature that grows with the gradient, an optimizer state that lags behind the gradient, and directions of growth that the loss cannot see (\cref{ch:spikes}).
\item \emph{How large a batch?} Up to a critical size set by the noise and the curvature, which preconditioning raises (\cref{ch:noise}).
\item \emph{Why decay the learning rate, and when?} The noise floor and the geometry of the loss along a ``river valley'' (\cref{ch:schedules}).
\item \emph{What does weight decay do?} It turns a step of fixed length in a norm into the Frank--Wolfe method on a ball of that norm, and it sets an equilibrium speed of rotation for layers followed by a normalization (\cref{ch:decay}).
\item \emph{Why Adam rather than SGD for transformers, and where does each optimizer end up?} Heavy-tailed classes, heterogeneous curvature, and large batches; and the maximum margin in the optimizer's norm (\cref{ch:transformers,ch:bias}).
\end{itemize}

\addsec{Two lines}

The book follows two lines of argument, which meet in its last part.

\paragraph{The unification line.} An optimizer turns a gradient into a step, and the step is rarely the gradient itself. The gradient of a loss is a linear functional on parameter space: it says how fast the loss changes in each direction. To turn it into a direction one needs a way to compare directions, that is, a norm. Given a norm $\norm{\cdot}$, a steepest direction is a unit vector $u$ that makes $\ip{g,u}$ most negative, written $\lmo(g)$, and a \emph{normalized} step of length $\eta$ is $\theta\leftarrow\theta+\eta\lmo(g)$. In the Euclidean norm this is gradient descent with a normalized gradient; in the max norm $\norm{\cdot}_\infty$ it is $\theta\leftarrow\theta-\eta\sign(g)$, sign descent, which is Adam with its moving averages and its $\epsilon$ switched off; for a weight matrix under the spectral norm it is $W\leftarrow W-\eta UV^\top$, where $G=U\diag(\sigma)V^\top$ is a compact singular value decomposition of the gradient, and this is Shampoo without accumulation \parencite{BernsteinNewhouse2024} and Muon without momentum \parencite{JordanEtAl2024Muon}. Two further ingredients complete the picture: moving averages, of the gradient (momentum, Adam's $\beta_1$) or of the iterates; and preconditioners, norms that adapt to the gradients seen so far (Adam's $\beta_2$, AdaGrad, Shampoo, SOAP). \Cref{tab:prologue:ingredients} lists the main optimizers of this book by their three ingredients.

\begin{table}[t]
\centering\small
\renewcommand{\arraystretch}{1.15}
\begin{tabular}{@{}>{\raggedright\arraybackslash}p{19mm}>{\raggedright\arraybackslash}p{29mm}>{\raggedright\arraybackslash}p{30mm}>{\raggedright\arraybackslash}p{29mm}@{}}
\toprule
\textsc{optimizer} & \textsc{norm} & \textsc{preconditioner} & \textsc{moving averages}\\
\midrule
SGD & $\ell_2$, unnormalized & none & momentum, optional\\
signSGD & $\ell_\infty$ & none & none\\
Lion & $\ell_\infty$ & none & one average of $g$, two coefficients\\
Adam & $\ell_\infty$ when memoryless & diagonal, from $g^2$ & of $g$ ($\beta_1$) and of $g^2$ ($\beta_2$)\\
AdaGrad & $\norm{\cdot}_{P_t}$, $P_t$ diagonal & diagonal, from the running sum of $g^2$ & none\\
Shampoo & spectral when memoryless & Kronecker factors, power $-\frac14$ & sums or averages of $GG^\top$ and $G^\top G$\\
K-FAC & Fisher norm, layer by layer & Kronecker factors, power $-1$ & of both factors\\
SOAP & $\ell_\infty$ in a rotated basis when memoryless & diagonal, in Shampoo's eigenbasis & as in Adam\\
Muon & spectral (RMS$\to$RMS) & none & Nesterov momentum\\
\bottomrule
\end{tabular}
\caption{Nine optimizers of Parts I and II, described by their three ingredients. ``Memoryless'' means that only the current gradient is used and that the small constant added for numerical stability is zero; the entries are made precise in the field guides of \cref{ch:steepest,ch:sign,ch:spectral,ch:momentum,ch:adaptive,ch:kronecker}.}
\label{tab:prologue:ingredients}
\end{table}

This line of thought has a precise history. Steepest descent in non-Euclidean norms is classical, and it was used to speed up combinatorial algorithms such as maximum flow \parencite{KelnerEtAl2014}. Spectral descent was proposed for neural networks by \textcite{CarlsonEtAl2015}; the observation that the optimizers of deep learning, with their moving averages switched off, are instances of steepest descent, and that the right norm for a network is assembled from norms on its layers, was developed by Jeremy Bernstein, Laker Newhouse, and coauthors \parencite{BernsteinNewhouse2024,LargeEtAl2024,BernsteinNewhouse2025} and is the subject of Newhouse's thesis \parencite{Newhouse2025}. Parts I and II of this book present it as a textbook theory, with complete proofs of the facts it rests on.

\paragraph{The dynamics line.} The second line studies what happens when these optimizers run with the step sizes that practice uses, which are far outside the regime of the theorem above. Full-batch gradient descent on a neural network drives the largest eigenvalue of the Hessian up until the step size is barely stable, and then holds it there \parencite{CohenEtAl2021}: the step size controls the curvature, not the other way around. The first steps of training meet curvature that would make the target step size unstable, which is one reason for warmup. Gradient noise adds to the loss a term proportional to the learning rate; holding the rate high makes fast progress along the flat directions of the landscape, and decaying it removes the noise term, which is what the cooldown of a schedule does. Weight decay, combined with normalization layers, drives the weights to an equilibrium norm, and the equilibrium sets the speed at which the weights rotate, an effective learning rate. Parts III and IV develop these mechanisms, each in the smallest model that exhibits it.

\paragraph{Where the lines meet.} Part V asks two questions that need both lines. Why does Adam train transformers so much faster than gradient descent? In the idealizations of \cref{ch:transformers} the answer involves the coordinatewise geometry of Adam (the first line), the heavy-tailed classes and heterogeneous curvature of the problem, and the regime of curvature and noise that the batch size sets (the second). Where does each optimizer end up? Each norm has its own implicit bias, a maximum-margin solution measured in that norm.

\addsec{A first example: one quadratic, two steps}

A single example already contains both lines. Let $h_1\ge h_2\ge\dots\ge h_d>0$ and consider the quadratic
\[
  \Loss(\theta)=\tfrac12\sum_{i=1}^dh_i\theta_i^2,
\]
whose Hessian $H=\diag(h_1,\dots,h_d)$ has sharpness $S=h_1$. We compare two steepest-descent steps: gradient descent, $\theta\leftarrow\theta-\eta\nabla\Loss(\theta)$, which is steepest descent in the Euclidean norm with the unnormalized step; and sign descent, $\theta\leftarrow\theta-\eta\sign(\nabla\Loss(\theta))$, which is normalized steepest descent in the max norm. Because $h_i>0$, the sign of the gradient is the sign of $\theta$.

\begin{theorem*}[Two steps on one quadratic]
Let $\theta_0\in\R^d$.
\begin{enumerate}
\item Gradient descent with step size $\eta>0$ gives $\theta_{t,i}=(1-\eta h_i)^t\theta_{0,i}$. It converges from every initialization if and only if $\eta<2/h_1$, and then each coordinate converges at the rate $\abs{1-\eta h_i}$, which is at least $1-2h_i/h_1$.
\item Sign descent with a constant step size $\eta>0$ moves each coordinate toward zero by exactly $\eta$ per step until it lies in $[0,\eta)$ in absolute value; from then on the coordinate alternates between two values $a_i$ and $a_i-\eta\sign(a_i)$ with $\abs{a_i}<\eta$, unless it is exactly zero. It never diverges and, unless every coordinate of $\theta_0$ is an integer multiple of $\eta$, never converges. For every $t\ge\max_i\abs{\theta_{0,i}}/\eta$, the loss $\Loss(\theta_t)$ is less than $\frac12\eta^2\tr H$.
\item Sign descent with positive step sizes $\eta_t$ that tend to zero and satisfy $\sum_t\eta_t=\infty$ converges to zero from every initialization.
\end{enumerate}
\end{theorem*}

\begin{proof}
The plan is to reduce to one coordinate and solve each recursion exactly. The coordinates evolve independently in both methods, so it suffices to follow one coordinate, $x_t=\theta_{t,i}$ with curvature $h=h_i$.

For (i), gradient descent gives $x_{t+1}=(1-\eta h)x_t$, a geometric sequence that tends to zero for every $x_0$ exactly when $\abs{1-\eta h}<1$, that is, when $0<\eta h<2$. All coordinates converge exactly when this holds for $h=h_1$, the largest curvature. If $\eta<2/h_1$, then $\eta h_i<2h_i/h_1$, so the rate $\abs{1-\eta h_i}\ge1-\eta h_i>1-2h_i/h_1$.

For (ii), sign descent gives $x_{t+1}=x_t-\eta\sign(x_t)$. If $\abs{x_t}\ge\eta$, then the step does not cross zero and $\abs{x_{t+1}}=\abs{x_t}-\eta$; so after $k=\lfloor\abs{x_0}/\eta\rfloor$ steps $\abs{x_k}=\abs{x_0}-k\eta<\eta$. If $x_k=0$, then the coordinate stays at zero, because $\sign0=0$. Otherwise the next step crosses zero: $x_{k+1}=x_k-\eta\sign(x_k)$ has the opposite sign and $\abs{x_{k+1}}=\eta-\abs{x_k}<\eta$, and the step after that returns to $x_{k+2}=x_{k+1}+\eta\sign(x_k)=x_k$. So the coordinate alternates between $a=x_k$ and $a-\eta\sign(a)$, both of absolute value less than $\eta$. It is eventually constant only when $x_k=0$, that is, when $x_0$ is a multiple of $\eta$. For $t\ge\max_i\abs{\theta_{0,i}}/\eta$ we have $t\ge k$ for every coordinate, so every coordinate is in its cycle or at zero, each term $\frac12h_i\theta_{t,i}^2$ is less than $\frac12h_i\eta^2$, and the loss is less than $\frac12\eta^2\sum_ih_i$.

For (iii), the step $x_{t+1}=x_t-\eta_t\sign(x_t)$ satisfies $\abs{x_{t+1}}=\abs{x_t}-\eta_t$ when $\abs{x_t}\ge\eta_t$, and $\abs{x_{t+1}}<\eta_t$ when $\abs{x_t}<\eta_t$: then $\abs{x_{t+1}}=\eta_t-\abs{x_t}$ if $x_t\ne0$, and $x_{t+1}=0$ if $x_t=0$. Let $\eps>0$ and choose $t_0$ with $\eta_t<\eps$ for all $t\ge t_0$. While $\abs{x_t}\ge\eps$ and $t\ge t_0$, the absolute value falls by $\eta_t$ at each step; since $\sum_t\eta_t=\infty$, it falls below $\eps$ at some step $t_1\ge t_0$. From then on it stays below $\eps$: if $\abs{x_t}<\eps$ and $t\ge t_0$, then either $\abs{x_{t+1}}=\abs{x_t}-\eta_t<\eps$ or $\abs{x_{t+1}}<\eta_t<\eps$. Since $\eps$ was arbitrary, $x_t\to0$.
\end{proof}

\begin{figure}[t]
\centering
\begin{tikzpicture}
\begin{axis}[wplot,xlabel={step $t$},ylabel={loss $\Loss(\theta_t)$},ymode=log,xmin=0,xmax=400,ymin=1e-9,ymax=2,legend pos=south west,legend columns=1]
\addplot[theory] table[x=t,y=gd_exact] {figures/data/00-prologue-quadratic.dat};
\addplot[theory3,line width=0.6pt] table[x=t,y=sign_exact] {figures/data/00-prologue-quadratic.dat};
\addplot[theory2] table[x=t,y=signlin_sim] {figures/data/00-prologue-quadratic.dat};
\addplot[simulation,each nth point=9] table[x=t,y=gd_sim] {figures/data/00-prologue-quadratic.dat};
\addplot[simulation,each nth point=9] table[x=t,y=sign_sim] {figures/data/00-prologue-quadratic.dat};
\addplot[guide,domain=0:400] {1.635*0.0004/6};
\addplot[guide] coordinates {(74.86,1e-9) (74.86,2)};
\legend{gradient descent, sign descent, sign descent with linear decay}
\end{axis}
\end{tikzpicture}
\caption{The two steps on $\Loss(\theta)=\frac12\sum_ih_i\theta_i^2$ with $h_i=1/i^2$, $d=100$, and $\theta_{0,i}$ uniform on $[\frac12,\frac32]$ (seed $0$). Gradient descent at $\eta=1.9$, just below its threshold $2/h_1=2$: the normal-mode formula (curve) and the iterates (points) agree to machine precision; the loss is soon dominated by the flat coordinates, whose rates $1-1.9\,h_i$ are close to one. Sign descent at $\eta=0.02$, the closed form of part (ii) of the theorem (curve) and the iterates (points): every coordinate travels at the same speed and reaches its two-cycle by the dashed vertical line $\max_i\abs{\theta_{0,i}}/\eta\approx75$; the loss then alternates between two values (the thin band, $4.1\times10^{-5}$ and $1.9\times10^{-4}$), around the dashed level $\eta^2\tr H/6\approx1.1\times10^{-4}$ expected for remainders $\abs{\theta_{0,i}}\bmod\eta$ spread uniformly over $[0,\eta)$. With the step size decayed linearly from $0.02$ to zero over the $400$ steps (dashed curve, iterates), sign descent goes far below that level, to $9\times10^{-10}$ at the last step.}
\label{fig:prologue:quadratic}
\end{figure}

Small as it is, the example contains the first lesson of each part of the book (\cref{fig:prologue:quadratic}). Gradient descent is limited by stability: its step size must stay below $2/h_1$, so the flat directions, with $h_i\ll h_1$, are learned no faster than at the rate $1-2h_i/h_1$, close to one; the condition number $h_1/h_d$ sets the speed. Sign descent is not limited by stability at all, and it treats every coordinate alike, whatever its curvature: on this problem it reaches its floor, a loss of about $10^{-4}$, in about $75$ steps, and gradient descent needs about $6500$ steps to get as low. This is the $\ell_\infty$ geometry of \cref{ch:sign}, and its advantage on badly conditioned, axis-aligned problems is one ingredient of the answer in \cref{ch:transformers}. But at a fixed step size a normalized step does not converge, except from special starting points; it oscillates at an amplitude set by the step, and only a schedule that takes the step to zero removes the oscillation. With noisy gradients a fixed step size leaves every optimizer on a floor set by the step, which is one reason that learning-rate schedules exist (\cref{ch:schedules}); and when the noise is skewed, a normalized step without memory can stall at a gradient of the size of the noise, however small the step, which is why normalized methods need momentum (\cref{ch:momentum}).

\addsec{What the quadratic gets wrong}

The quadratic is to optimization what the harmonic oscillator is to mechanics, and it is wrong in ways that organize the rest of the book. Its Hessian is constant, so the step size cannot change the curvature; in a network it does: gradient descent runs at the edge of stability (\cref{ch:eos}), and the curvature of the first steps, too high for the target step size, calls for a warmup (\cref{ch:warmup}). Its gradients are exact; with noise, every method has a floor that depends on the step size, and the batch size and the schedule become the main controls (\cref{ch:noise,ch:schedules}). It is axis-aligned; rotate it, and sign descent loses its advantage (\cref{ch:sign,ch:transformers}). It has a single minimizer; a model that can fit its training data can fit it in many ways, and the norm of the optimizer decides which one training approaches (\cref{ch:bias}). And its parameters are a single vector; a network's parameters are matrices arranged in layers, which calls for norms on matrices and for a norm on the whole network (\cref{ch:spectral,ch:modular}).

\addsec{How each chapter is built}

Each chapter starts in the wild, with a practice: a default, a rule of thumb, or a regularity seen in training runs, stated in a box with its typical values and its sources, together with the evidence for it. The second section builds the smallest idealization in which the practice can be analyzed, and a dictionary translates between the two. The theory follows, with at least one main theorem proved completely. A section on reading the theory compares its predictions with a simulation; in every chapter at least one figure plots an exact prediction as a curve against simulated points. The chapter then says what the theory misses, in three parts: what was idealized away, where the theory breaks, and which refinements repair it and what remains open. The last section returns to the wild, with field notes that summarize the practice, the idealization, what was proved, what is predicted, what is missed, and what a practitioner might take away.

In Parts I and II each optimizer gets a field-guide entry: its update rule, its norm, its preconditioner, its moving averages, its memory cost, and its typical settings. Status is marked throughout. A theorem with a proof ending in $\blacksquare$ is proved; a proof sketch says what it omits and where to find it; a heuristic derivation, ending in $\Diamond$, is an argument that is not a proof and is marked as such; an observation is an empirical regularity, stated with its evidence; and a rule used without evidence that can be cited is called folklore. Many practices in this book are folklore, and saying so is part of the theory.

Each chapter closes with exercises of four kinds (consolidation, variation, simulation, and research) and with notes that end by translating the notation of the main source papers into the notation of this book.

\addsec{A map of the book}

\Cref{tab:prologue:map} lists the chapters. Part I develops steepest descent: the geometry of a step in a general norm, the $\ell_\infty$ geometry of sign descent and Adam, the spectral geometry of Muon, and the norm of a whole network. Part II adds memory: moving averages, adaptive diagonal preconditioners, and Kronecker-factored ones. Part III is about stability: the edge of stability, warmup, and clipping and loss spikes. Part IV is about noise and time: batch size, schedules, and weight decay. Part V brings the two lines together. The epilogue keeps a ledger of the practices and the verdicts of the chapters on them, and the appendix collects the mathematics the chapters share.

\begin{table}[p]
\begin{fullwidth}
\centering\footnotesize
\renewcommand{\arraystretch}{1.14}
\begin{tabular}{@{}r>{\raggedright\arraybackslash}p{28mm}>{\raggedright\arraybackslash}p{32mm}>{\raggedright\arraybackslash}p{50mm}>{\raggedright\arraybackslash}p{38mm}@{}}
\toprule
& \textsc{in the wild} & \textsc{idealization} & \textsc{main result} & \textsc{misses}\\
\midrule
1 & Normalized updates & Steepest descent in a normed space & Duality maps of the norms in use; rates that depend on the norm only through $L$ and $\norm\cdot_*$ & A fixed norm; global smoothness\\
2 & Adam, signSGD, Lion & $\ell_\infty$ steepest descent & Sign descent wins when the gradient is spread over more coordinates than the curvature has stiff directions; Adam's step is at most $C(\beta_1,\beta_2)$ learning rates & Memory; why the coordinates of a network are good ones\\
3 & Muon, Shampoo & Spectral steepest descent & The polar factor is the steepest step; Muon's quintic keeps the singular values in a band; spectral beats Frobenius when the gradient's nuclear rank exceeds the inputs' stable rank & Momentum and noise; one matrix, one step\\
4 & Hyperparameter transfer & The norm of a network & Per-layer learning rates for width-independent feature updates & Alignment after the first step; horizon and batch size\\
\midrule
5 & Momentum $0.9$ & Quadratics; noisy normalized steps & Acceleration and stability regions; with noise, momentum is a larger step; without memory, normalized steps keep a bias & A reason for $0.9$\\
6 & Adam's $\beta_2$; AdaGrad & Online convex optimization & AdaGrad against the best diagonal norm in hindsight; when Adam fails; the memory that $\beta_2$ sets & Nonconvexity; why $0.95$ works\\
7 & Shampoo, K-FAC, SOAP & Kronecker-structured curvature & The Kronecker bound and Shampoo's regret; which power of the factors conditions the problem & Drifting, non-Kronecker statistics; gains at scale\\
\midrule
8 & Training at the stability limit & A cubic model of self-stabilization & Thresholds for GD, momentum and Adam; time averages at the edge & The cause of progressive sharpening\\
9 & Linear warmup & Curvature bounded by the loss & The safe step grows as the loss falls; a rising step is optimal among stable schedules & The curvature bound is measured, not derived\\
10 & Clipping, z-loss, QK-norm & $(L_0,L_1)$-smoothness & Clipped descent avoids the large-gradient penalty that every fixed step pays; mechanisms of spikes & Interacting causes of real spikes\\
\midrule
11 & Batch-size rules & The noisy quadratic model & Exact loss curves; linear scaling up to a critical batch size; preconditioning raises it & Curvature and noise that change during training\\
12 & Cosine and WSD & Convex last iterates; the river valley & Bounds for every schedule; the noise floor and the cooldown & A landscape that moves with the schedule\\
13 & AdamW's weight decay & Frank--Wolfe; scale invariance & Steepest descent with decay is Frank--Wolfe; rotational equilibrium at the angle $\arccos(1-\eta\lambda)$ per step & Layers that are not scale invariant\\
\midrule
14 & Adam for transformers & Zipf classes; block curvature; batch regimes & GD slows to logarithmic decay on heavy-tailed classes; the gap grows with batch size & How the mechanisms interact\\
15 & Adaptive methods and generalization & Separable data & Steepest descent reaches the maximum margin of its norm & Depth; finite training\\
\bottomrule
\end{tabular}
\caption{The chapters, each with the practice it starts from, the idealization it analyzes, its main result, and its most important omission.}
\label{tab:prologue:map}
\end{fullwidth}
\end{table}

\addsec{How to read this book}

The prerequisites are undergraduate analysis, linear algebra, and probability, together with familiarity with gradient descent and with neural networks at the level of an introductory course. Each chapter builds what it uses, and the shared tools (dual norms, the polar decomposition, Kronecker products and the matrix geometric mean, the stability of linear recurrences, the descent lemma, online convex optimization, and martingale and tail bounds) are collected in \cref{ch:toolkit}. \Cref{ch:steepest} fixes the vocabulary of the whole book and should be read first. After it, Parts I and II can be read in order as a course on the unification line, and Parts III and IV as a course on the dynamics line; \cref{ch:momentum} is useful for both. Part V uses both lines.

Three cautions. First, an idealization is not an explanation of training. Each chapter explains one practice in one idealized setting, and whether the same mechanism operates in a large network is an empirical question that the chapter examines separately. Second, practice changes quickly: the recipes quoted in this book are dated and sourced, and some will be out of date when it is read. The mechanisms should last longer than the recipes. Third, the theory is uneven. Some practices, such as the form of the update rules, are explained well; most, such as warmup and the shape of the learning-rate schedule, are explained in part; and some, such as the value $\beta_1=0.9$ or the rise of curvature during training, are not explained at all. The epilogue keeps the account.

\notesandrefs

The convergence theorem for stochastic gradient descent in the nonconvex smooth setting is due to \textcite{GhadimiLan2013}; \textcite{BottouCurtisNocedal2018} survey the theory of stochastic optimization for machine learning as it stood in 2018, before many of the practices in this book became standard, and \textcite{Nesterov2018} and \textcite{Polyak1987} are standard texts on gradient methods and their convergence theory. The training recipes quoted in the prologue are from the papers describing GPT-3 \parencite{BrownEtAl2020} and LLaMA \parencite{TouvronEtAl2023}; the \emph{Deep Learning Tuning Playbook} \parencite{GodboleEtAl2023} is a systematic practical guide to tuning such recipes. Adam is due to \textcite{KingmaBa2015} and its decoupled weight decay to \textcite{LoshchilovHutter2019}; Muon was introduced in a blog post by \textcite{JordanEtAl2024Muon} and scaled to large language models by \textcite{LiuEtAl2025Muon}.

The unification of optimizers as steepest descent under norms is set out in \textcite{BernsteinNewhouse2024}, a workshop paper that reads Adam, Shampoo, and Prodigy with their moving averages switched off as steepest descent in particular norms, and, for whole networks, in \textcite{LargeEtAl2024} and \textcite{BernsteinNewhouse2025}; \textcite{Newhouse2025} gives a connected account in a thesis, and \cref{ch:steepest} gives more of the history. The edge of stability was described by \textcite{CohenEtAl2021}. Benchmarks that compare optimizers under careful tuning are AlgoPerf \parencite{DahlEtAl2023} and the study of \textcite{WenEtAl2026Fantastic}, which found that the speedups of newer optimizers over a well-tuned AdamW shrink as models grow; their results are a useful corrective to every claim in this book that one optimizer is better than another.

\printchapterbib
\end{refsection}


\casepart{One Rule: Steepest Descent}{An optimizer turns a gradient into a step, and the step depends on a norm. In the Euclidean norm the step is the gradient; in the max norm it is a vector of signs; for a weight matrix under the spectral norm it is the orthogonalized gradient; for a whole network the norm is assembled from the norms of its layers, and that norm fixes how the learning rate must scale with width.}{18}
% Chapter 1 of Optimizers in the Wild: the geometry of a step.

\casechapter{The Geometry of a Step}{Steepest descent in a normed space}{Normalized updates}{ch:steepest}
\begin{refsection}[chapters/01-steepest.bib]

\begin{chaptersummary}
The optimizers that train large networks do not step along the gradient: Adam, signSGD, Lion and Muon let the gradient choose a direction and the learning rate a length. The idealization that explains them is steepest descent in a normed space, in which the gradient is a dual vector and a norm turns it into a step. We compute the steepest step for the norms that optimizers use, and prove that the guaranteed decrease of a step and the convergence rates depend on the norm only through the smoothness constant in that norm and the dual norm of the gradient. For quadratics both are explicit: the max-norm step wins when the gradient is denser than the curvature's effective rank, and for sign descent on a diagonal quadratic the bound is attained exactly, including a floor which says that a normalized step of fixed length never converges. The theory fixes the norm in advance, assumes a global smoothness constant, and cannot say which norm to choose.
\end{chaptersummary}

\section{In the wild: steps of a fixed size}\label{sec:steepest:wild}

Plain gradient descent takes a step proportional to the gradient: if the gradient doubles, so does the step. Few of the optimizers that train large networks work this way. Most let the gradient choose a direction and the learning rate a length; Kingma and Ba read Adam's learning rate as the radius of a trust region around the current parameters \parencite{KingmaBa2015}.\sidenote{The reading is in Section~2.1 of their paper, with two approximate bounds on the size of the step; \cref{ch:sign} derives the exact bound.} Bernstein and coauthors proposed signSGD to compress the gradients that workers exchange to one bit per coordinate \parencite{BernsteinEtAl2018}, Lion was found by a program search over update rules \parencite{ChenEtAl2023Lion}, and gradient clipping, part of the recipe quoted in the prologue, caps the step whenever the gradient is large (\cref{ch:spikes}).

\begin{practice}[Adam, \textcite{KingmaBa2015}; Lion, \textcite{ChenEtAl2023Lion}; LARS and LAMB, \textcite{YouGitmanGinsburg2017,YouEtAl2020}; Muon, \textcite{JordanEtAl2024Muon,LiuEtAl2025Muon}]
Let the gradient choose the direction and the learning rate the size of the step. Adam divides the momentum of each coordinate by the root-mean-square of its recent gradients, so that no coordinate moves much more than $\eta$ per step. In signSGD and Lion every coordinate moves by exactly $\eta$, in the direction of the sign of the gradient or of an interpolation between the gradient and its momentum; Lion's authors recommend a learning rate $3$ to $10$ times smaller than AdamW's, because a sign step is longer. LARS and LAMB rescale the update of each layer to a fixed fraction of the norm of the layer's weights. Muon replaces the momentum of each hidden weight matrix by an approximation of its orthogonal polar factor, whose singular values all equal one; to reuse AdamW's learning rate and weight decay, a scaled variant multiplies the polar factor of an $m\times n$ matrix by $0.2\sqrt{\max(m,n)}$, which gives the update a root-mean-square of $0.2$ times the learning rate.
\end{practice}

Because these rules fix the size of the step, the size of an update is measured, and matched across optimizers.

\begin{observation}[Update size; \textcite{LiuEtAl2025Muon}]\label[observation]{obs:steepest:rms}
In the language-model training runs reported by Liu and coauthors, the root-mean-square of AdamW's update of a weight matrix, in units of the learning rate, typically lay between $0.2$ and $0.4$. Rescaled to root-mean-square $0.2$, the Muon update could reuse the learning rate and weight decay tuned for AdamW.
\end{observation}

This is one report, from one family of models, in 2025; the older observation behind it is that Adam's update behaves like a sign when the gradients are consistent from step to step \parencite{BallesHennig2018}. The logic of the rule is that a learning rate is a length, and lengths in different norms can be converted. A second observation shows that normalization matters.

\begin{observation}[\textcite{KunstnerEtAl2023}]\label[observation]{obs:steepest:sign}
Transformers for language and translation, trained with full-batch gradients so that noise plays no role, still train much faster with Adam than with gradient descent, and sign descent with momentum closes most of the gap.
\end{observation}

On these tasks (\cref{ch:sign,ch:transformers}), much of what separates Adam from gradient descent lies in replacing the gradient by its sign. The rules of the box are one rule: take the direction in which the loss decreases fastest per unit of a chosen length, and move a prescribed distance in that length. Which length is right, and what does the choice buy?

\section{The idealization: a gradient is a dual vector}\label{sec:steepest:idealization}

Let $\Loss\colon\R^d\to\R$ be differentiable. Its derivative at $\theta$ is the linear functional $u\mapsto\ip{\nabla\Loss(\theta),u}=\sum_i\partial_i\Loss(\theta)\,u_i$; the pairing identifies a functional with its coefficients and measures no lengths (for matrices, $\ip{G,D}=\tr(G^\top D)$). The gradient lives in the dual space: it says how fast the loss changes along each direction, but it is not a direction in which to move, and its coordinates are in units of loss per unit of $\theta_i$. Gradient descent converts the gradient into a step with the Euclidean inner product of whatever coordinates the parameters are written in, and a linear change of coordinates changes the method (\cref{ex:steepest:quadnorm}). A norm supplies the map from the dual space back to the parameter space.

\begin{definition}[Dual norm]\label[definition]{def:steepest:dualnorm}
Let $\norm{\cdot}$ be a norm on $\R^d$. Its \term{dual norm}\index{dual norm} is $\norm{g}_*=\max_{\norm{u}\le1}\ip{g,u}$.
\end{definition}

The maximum exists because the unit ball is compact. By \cref{lem:toolkit:dualnorm}, $\norm{\cdot}_*$ is a norm whose dual is $\norm{\cdot}$, Hölder's inequality $\abs{\ip{g,u}}\le\norm{g}_*\norm{u}$ holds, and the dual of $\ell_p$ is $\ell_q$ with $1/p+1/q=1$. The dual norm of the gradient is the largest rate at which the loss decreases along a direction of unit length: the slope of the loss, measured with the chosen ruler.

\begin{definition}[Linear minimization oracle and sharp operator]\label[definition]{def:steepest:lmo}
For $g\in\R^d$, the \term{linear minimization oracle}\index{linear minimization oracle} $\lmo(g)$ is an element of the set $\argmin_{\norm{u}\le1}\ip{g,u}$ of \term{steepest unit directions}; when the set has several elements we fix a choice, and we take $\lmo(0)=0$. The \term{sharp operator}\index{sharp operator}\label[definition]{def:steepest:sharp} maps $g$ to $g^\sharp=-\norm{g}_*\lmo(g)$.
\end{definition}

The oracle is named after the subproblem of the Frank--Wolfe method \parencite{FrankWolfe1956}, and \cref{ch:decay} shows that steepest descent with decoupled weight decay is that method.\sidenote{The sharp notation follows \textcite{KelnerEtAl2014}, who used steepest descent in the $\ell_\infty$ norm to compute approximate maximum flows.} The steepest unit direction minimizes the directional derivative over directions of unit length: steepness is a ratio of a change of loss to a length, and it depends on how length is measured.

\begin{definition}[Steepest descent]\label[definition]{def:steepest:steepest}
Let $\eta>0$ and $g_t=\nabla\Loss(\theta_t)$. \term{Normalized steepest descent}\index{steepest descent!normalized} in the norm $\norm{\cdot}$ is the iteration $\theta_{t+1}=\theta_t+\eta\lmo(g_t)$, and \term{unnormalized steepest descent}\index{steepest descent!unnormalized} is $\theta_{t+1}=\theta_t-\eta g_t^\sharp$.
\end{definition}

By \cref{lem:steepest:basic} below, a normalized step has length exactly $\eta$ in the norm unless the gradient vanishes, and an unnormalized step has length $\eta\norm{g_t}_*$; this is what ``normalized'' means throughout the book.\sidenote{It has nothing to do with normalization layers, which normalize activations; \cref{ch:decay} studies how the two interact.} The step size $\eta$ is the learning rate of practice.\sidenote{For a normalized method the learning rate has the units of the parameters. For gradient descent it has the units of a squared parameter per unit of loss, so its good values change when the loss is rescaled, while those of a normalized method do not (\cref{lem:steepest:basic}(iv)).} In the Euclidean norm $g^\sharp=g$, and unnormalized steepest descent is gradient descent.

\begin{remark}[One formula]\label[remark]{rem:steepest:trustregion}
By the substitutions $u=\eta v$ and $u=-\eta w$ and \cref{lem:steepest:basic}(iii), both steps minimize the linear model of the loss under a control on the length of the step:
\[
  \eta\lmo(g_t)\in\argmin_{\norm{u}\le\eta}\ip{g_t,u},\qquad
  -\eta g_t^\sharp\in\argmin_{u}\Bigl\{\ip{g_t,u}+\frac1{2\eta}\norm{u}^2\Bigr\}.
\]
The first is a \term{trust-region}\index{trust region} step: the model is trusted on a ball of radius $\eta$ and minimized there \parencite{ConnGouldToint2000}. The second replaces the constraint by a penalty, and the two coincide when the radius is $\eta\norm{g_t}_*$. Every optimizer of Parts~I and~II takes a step of the first kind, with a norm that may change from step to step and a gradient that may be replaced by an average: this is the one formula of the book.
\end{remark}

\begin{assumption}[Smoothness in a norm]\label[assumption]{ass:steepest:smooth}
The loss $\Loss\colon\R^d\to\R$ is differentiable and bounded below, $\Loss^\star=\inf\Loss>-\infty$, and for some $L>0$ it is \term{$L$-smooth with respect to $\norm{\cdot}$}\index{smoothness!with respect to a norm}: $\norm{\nabla\Loss(\theta)-\nabla\Loss(\theta')}_*\le L\norm{\theta-\theta'}$ for all $\theta,\theta'$.
\end{assumption}

Displacements are measured in the norm and changes of gradient in the dual norm. All norms on $\R^d$ are equivalent, but the smoothness constants differ by factors as large as the dimension (\cref{prop:steepest:whichnorm}), and the content of a choice of norm lies in these factors. For a twice differentiable loss, the descent inequality below needs only $u^\top\nabla^2\Loss(\theta)\,u\le L\norm{u}^2$.

\begin{dictionary}
the norm $\norm{\cdot}$ & what the optimizer treats as a step of unit size: $\ell_2$ for SGD, $\ell_\infty$ for sign descent, the spectral norm of each matrix for Muon, a maximum over layers for LARS\\
the radius $\eta$ & the learning rate of a normalized optimizer, the length of every step in its norm\\
$\lmo(g)$ & the update direction before the learning rate: $-g/\norm{g}_2$, $-\sign(g)$, $-\polar(G)$\\
$\norm{g}_*$ & the decrease of the loss per unit step: $\norm{g}_1$ for sign descent, $\norm{G}_{\nuc}$ for Muon\\
the smoothness constant $L$ & the curvature that limits the step, in the norm: the sharpness $S$ for $\ell_2$\\
the floor $L\eta/2$ & the plateau of a run at a constant learning rate, removed only by a decaying schedule
\end{dictionary}

Two optimizers are derived in full here; greedy coordinate descent is steepest descent in $\ell_1$ (\cref{prop:steepest:dualitymaps}(iii)), and sign descent and Muon have their entries in \cref{ch:sign,ch:spectral}.\sidenote{Cauchy's note of 1847 proposed to solve a system of equations by moving, step after step, in the direction that decreases a sum of squares fastest.}

\begin{fieldguide}{GD and SGD}{\textcite{Cauchy1847}; \textcite{RobbinsMonro1951}}
\fgrow{update}{$\theta_{t+1}=\theta_t-\eta g_t$, with the gradient (GD) or a minibatch gradient (SGD)}
\fgrow{norm}{$\ell_2$, unnormalized ($g^\sharp=g$): the step has length $\eta\norm{g_t}_2$}
\fgrow{preconditioner}{none; a fixed $P\succ0$ gives the norm $\norm{\cdot}_P$ (\cref{ex:steepest:quadnorm})}
\fgrow{averages}{none; momentum in \cref{ch:momentum}}
\fgrow{state}{none (one number per parameter with momentum)}
\fgrow{in the wild}{with momentum $0.9$, common for convolutional networks, rare for transformers}
\end{fieldguide}

\begin{fieldguide}{LARS and LAMB}{\textcite{YouGitmanGinsburg2017,YouEtAl2020}}
\fgrow{update}{per layer, $W_l\leftarrow W_l-\eta\norm{W_l}_F\,U_l/\norm{U_l}_F$, with $U_l$ the gradient (LARS) or the Adam update (LAMB)}
\fgrow{norm}{$\max_l\norm{\Delta W_l}_F/\norm{W_l}_F$, rescaled by the weights at each step; without momentum, LARS is normalized steepest descent in it (\cref{prop:steepest:dualitymaps}(iv))}
\fgrow{preconditioner}{none (LARS); Adam's diagonal one (LAMB)}
\fgrow{averages}{momentum (LARS); Adam's two averages (LAMB)}
\fgrow{state}{one or two numbers per parameter}
\fgrow{in the wild}{large-batch training of image classifiers and language models; $\norm{W_l}_F/\norm{U_l}_F$ is the ``trust ratio''}
\end{fieldguide}

\paragraph{What was thrown away.} Noise (the gradients are exact), memory (each step depends only on the current gradient), change (the norm is fixed and the smoothness constant global) and structure (one vector, or a fixed list of blocks). Each has a chapter of its own: noise and memory change what is fed to the duality map, not the map (\cref{ch:momentum}); a norm that adapts to the gradients is a preconditioner (\cref{ch:adaptive,ch:kronecker}); changing curvature is the subject of Part~III; the norm of a network is built in \cref{ch:modular}. What is left comes first: given one exact gradient and one norm, what is the best step, and what does it guarantee?

\paragraph{What one had to notice.} The rules of the practice box are not normalizations bolted onto gradient descent; they are gradient descent in other norms. A sign vector is the steepest unit direction of the max norm, a polar factor that of the spectral norm. Steepest descent in a general norm is classical \parencite{BoydVandenberghe2004}, spectral descent was proposed for networks by \textcite{CarlsonEtAl2015}, and \textcite{BernsteinNewhouse2024} observed that Adam, Shampoo and Prodigy with their accumulators switched off are steepest descent in particular norms, a view that \textcite{LargeEtAl2024,BernsteinNewhouse2025} extended to whole networks. Seen this way, the choice of an optimizer is the choice of a norm, and it can be analyzed.

\section{The theory: the steepest step in each norm}\label{sec:steepest:maps}

We collect first the elementary facts that every later chapter uses.

\begin{lemma}[Basic properties]\label[lemma]{lem:steepest:basic}
Let $\norm{\cdot}$ be a norm on $\R^d$ and $g\in\R^d$.
\begin{enumerate}
\item $\ip{g,\lmo(g)}=-\norm{g}_*$, and $\norm{\lmo(g)}=1$ if $g\ne0$.
\item $\ip{g,g^\sharp}=\norm{g}_*^2$ and $\norm{g^\sharp}=\norm{g}_*$.
\item For every $u$, $\ip{g,u}-\frac12\norm{u}^2\le\frac12\norm{g}_*^2$, with equality at $u=g^\sharp$. If $g\ne0$, then the maximizers are the vectors $-\norm{g}_*v$ with $v$ a steepest unit direction; if $g=0$, then the only maximizer is $0$.
\item For $c>0$ the steepest unit directions of $cg$ and of $g$ coincide; so $\lmo$ can be chosen with $\lmo(cg)=\lmo(g)$, and then $(cg)^\sharp=c\,g^\sharp$.
\end{enumerate}
\end{lemma}

\begin{proof}
The plan is to combine Hölder's inequality with the symmetry of the unit ball. (i) By symmetry, $\min_{\norm{u}\le1}\ip{g,u}$ equals $-\max_{\norm{u}\le1}\ip{g,-u}$, which is $-\norm{g}_*$. If $g\ne0$, then $\norm{g}_*>0$, and Hölder's inequality gives $\norm{g}_*=-\ip{g,\lmo(g)}\le\norm{g}_*\norm{\lmo(g)}$, so $\norm{\lmo(g)}\ge1$; the reverse holds because $\lmo(g)$ lies in the ball. (ii) Multiply both identities of (i) by $\norm{g}_*$. (iii) By Hölder's inequality and $ab-\frac12b^2\le\frac12a^2$,
\[
  \ip{g,u}-\tfrac12\norm{u}^2\le\norm{g}_*\norm{u}-\tfrac12\norm{u}^2\le\tfrac12\norm{g}_*^2 .
\]
The second step is an equality only if $\norm{u}=\norm{g}_*$, and the first only if $\ip{g,u}=\norm{g}_*\norm{u}$. If $g=0$, then this forces $u=0$. If $g\ne0$, then it says that $u=\norm{g}_*w$ with $\norm{w}=1$ and $\ip{g,w}=\norm{g}_*$, that is, $u=-\norm{g}_*v$ with $v=-w$ a steepest unit direction; by (ii), $g^\sharp$ is one of these. (iv) The objective $c\ip{g,u}$ has the same minimizers, and $\norm{cg}_*=c\norm{g}_*$.
\end{proof}

Part (iii) says that $\frac12\norm{\cdot}_*^2$ is the convex conjugate of $\frac12\norm{\cdot}^2$ (\cref{lem:toolkit:conjugate}); part (iv), that a normalized step is insensitive to the scale of the loss. For $W\in\R^{m\times n}$ we write $We_j$ for the columns, $W^\top e_i$ for the rows, and $\norm{W}_{p\to q}=\max_{x\ne0}\norm{Wx}_q/\norm{x}_p$.

\begin{proposition}[Duality maps of the norms that optimizers use]\label[proposition]{prop:steepest:dualitymaps}
Let $g\in\R^d$ and $G\in\R^{m\times n}$ be nonzero, and let $M(g)=\argmin_{\norm{u}\le1}\ip{g,u}$ be the set of steepest unit directions of the norm at hand.
\begin{enumerate}
\item ($\ell_2$ and Frobenius.) The dual of $\norm{\cdot}_2$ is $\norm{\cdot}_2$, $M(g)=\{-g/\norm{g}_2\}$ and $g^\sharp=g$. Likewise $\norm{\cdot}_F$ is self-dual, $\lmo(G)=-G/\norm{G}_F$ and $G^\sharp=G$.
\item ($\ell_\infty$.) The dual norm is $\norm{g}_1$, and $M(g)=\{u\in[-1,1]^d:u_i=-\sign(g_i)\text{ whenever }g_i\ne0\}$. We choose $\lmo(g)=-\sign(g)$, so that $g^\sharp=\norm{g}_1\sign(g)$, and call normalized steepest descent in $\ell_\infty$ \term{sign descent}\index{sign descent}.
\item ($\ell_1$.) The dual norm is $\norm{g}_\infty$, and $M(g)$ is the convex hull of the vectors $-\sign(g_i)e_i$ with $\abs{g_i}=\norm{g}_\infty$. For one such $i$, $\lmo(g)=-\sign(g_i)e_i$ and $g^\sharp=g_ie_i$: unnormalized steepest descent in $\ell_1$ is greedy coordinate descent, the \term{Gauss--Southwell rule}\index{Gauss--Southwell rule}.
\item (Maxima over blocks.) Let $\R^d=\R^{d_1}\times\dots\times\R^{d_K}$ with $\theta=(\theta_1,\dots,\theta_K)$, give the $k$th factor a norm $\norm{\cdot}_{(k)}$ with dual $\norm{\cdot}_{(k),*}$ and steepest unit directions $M_k$, and let $\norm{\theta}=\max_k\norm{\theta_k}_{(k)}$. Then $\norm{g}_*=\sum_k\norm{g_k}_{(k),*}$, $M(g)=M_1(g_1)\times\dots\times M_K(g_K)$, $\lmo(g)=(\lmo_{(k)}(g_k))_k$, and $g^\sharp=-\bigl(\sum_j\norm{g_j}_{(j),*}\bigr)(\lmo_{(k)}(g_k))_k$.
\item (Columns and rows.) $\norm{W}_{1\to2}=\max_j\norm{We_j}_2$ and $\norm{W}_{2\to\infty}=\max_i\norm{W^\top e_i}_2$. Their duals are the sum of the column norms and the sum of the row norms, and $\lmo$ normalizes each column, respectively row, of $G$ to unit length and reverses its sign, mapping a zero column or row to zero.
\item (Spectral.) The dual of $\norm{\cdot}_{\op}$ is the nuclear norm $\norm{G}_{\nuc}=\sum_i\sigma_i(G)$, and $-\polar(G)\in M(G)$, so $G^\sharp=\norm{G}_{\nuc}\polar(G)$.
\end{enumerate}
\end{proposition}

\begin{proof}
The plan is to bound $\ip{g,u}$ from below over the unit ball by a chain of elementary inequalities, to exhibit a point where every inequality is an equality, and to read off the set of such points; (ii) and (v) are instances of (iv). If $\min_{\norm{u}\le1}\ip{g,u}=-c$, then $\norm{g}_*=c$ by symmetry of the ball.

(i) By the Cauchy--Schwarz inequality, $\ip{g,u}\ge-\norm{g}_2\norm{u}_2\ge-\norm{g}_2$ when $\norm{u}_2\le1$. Equality in the first step requires $u=-cg$ with $c\ge0$, and in the second $\norm{u}_2=1$, so $u=-g/\norm{g}_2$, which attains the bound; then $g^\sharp=g$. The Frobenius norm is the $\ell_2$ norm of the entries, and the trace pairing is the coordinate pairing of the entries.

(iv) If $\norm{u}\le1$, then $\norm{u_k}_{(k)}\le1$ for every $k$, and Hölder's inequality in each factor gives $\ip{g,u}=\sum_k\ip{g_k,u_k}\ge-\sum_k\norm{g_k}_{(k),*}$. Equality holds if and only if $\ip{g_k,u_k}=-\norm{g_k}_{(k),*}$ for every $k$, that is, $u_k\in M_k(g_k)$ for every $k$; for a block with $g_k=0$ this is the whole unit ball of the block. Such $u$ exist, for instance $u_k=\lmo_{(k)}(g_k)$, so the minimum is $-\sum_k\norm{g_k}_{(k),*}$ and the minimizers form the product of the sets $M_k(g_k)$.

(ii) The max norm is (iv) with $d$ blocks of size one, each with the absolute value, which is its own dual. For one coordinate, $\argmin_{\abs{u_i}\le1}g_iu_i$ is $\{-\sign(g_i)\}$ if $g_i\ne0$ and $[-1,1]$ if $g_i=0$.

(iii) For $\norm{u}_1\le1$,
\[
  \ip{g,u}\ge-\sum_i\abs{g_i}\abs{u_i}\ge-\norm{g}_\infty\sum_i\abs{u_i}\ge-\norm{g}_\infty .
\]
The first inequality is an equality if and only if $g_iu_i\le0$ for all $i$, the second if and only if $u_i=0$ whenever $\abs{g_i}<\norm{g}_\infty$, and the third, since $g\ne0$, if and only if $\norm{u}_1=1$. The vectors with these properties are $u=-\sum_{i\in I}c_i\sign(g_i)e_i$ with $c_i\ge0$, $\sum_ic_i=1$ and $I=\{i:\abs{g_i}=\norm{g}_\infty\}$: the convex hull of the vertices $-\sign(g_i)e_i$, $i\in I$. At a vertex, $g^\sharp=\norm{g}_\infty\sign(g_i)e_i=g_ie_i$.

(v) By the triangle inequality, $\norm{Wx}_2=\norm{\sum_jx_jWe_j}_2\le\norm{x}_1\max_j\norm{We_j}_2$, with equality at $x=e_j$ for a column of largest norm. By the Cauchy--Schwarz inequality, $\norm{Wx}_\infty=\max_i\abs{\ip{W^\top e_i,x}}\le\norm{x}_2\max_i\norm{W^\top e_i}_2$, with equality at $x=W^\top e_i/\norm{W^\top e_i}_2$ for a row of largest norm. So both norms are maxima, over the columns and over the rows, of the Euclidean norm, and (iv) with (i) gives their duals and oracles; for a zero block we choose zero.

(vi) By von Neumann's trace inequality (\cref{thm:toolkit:vonneumann}), every $D$ with $\norm{D}_{\op}=\sigma_1(D)\le1$ satisfies
\[
  \ip{G,D}=\tr(G^\top D)\ge-\sum_i\sigma_i(G)\sigma_i(D)\ge-\sigma_1(D)\sum_i\sigma_i(G)\ge-\norm{G}_{\nuc}.
\]
Let $G=U\diag(\sigma)V^\top$ be a compact singular value decomposition, with $r=\rank G\ge1$, $U\in\R^{m\times r}$, $V\in\R^{n\times r}$ and $U^\top U=V^\top V=I_r$. Then $(UV^\top)^\top UV^\top=VV^\top$ is an orthogonal projection of rank $r$, so the nonzero singular values of $UV^\top$ equal $1$ and $\norm{UV^\top}_{\op}=1$; and $\ip{G,-UV^\top}=-\tr(\diag(\sigma)V^\top V)=-\sum_i\sigma_i$. So $-UV^\top$ attains the bound.
\end{proof}

\begin{remark}[Ties, jumps and step sizes]\label[remark]{rem:steepest:ties}
The steepest unit direction is unique for the Euclidean and Frobenius norms, whose balls are strictly convex, and otherwise only away from ties: for $\ell_\infty$ when no $g_i$ vanishes, for $\ell_1$ when the largest $\abs{g_i}$ is attained once, for the spectral norm when $G$ has full rank (\cref{ch:spectral}). At a tie the direction jumps, as a coordinate of $-\sign(g)$ does when $g_i$ crosses zero, and noisy gradients straddle such ties all the time (\cref{ch:sign,ch:momentum}). The norms also give steps of very different sizes: a normalized step of length $\eta$ has root-mean-square $\eta$ in $\ell_\infty$ but $\eta/\sqrt d$ in $\ell_2$ (\cref{ex:steepest:rms}), and matching update sizes (\cref{obs:steepest:rms}) converts between the norms.
\end{remark}

\begin{remark}[Steepest descent is not mirror descent]\label[remark]{rem:steepest:mirror}
\term{Mirror descent}\index{mirror descent} \parencite{NemirovskyYudin1983,BeckTeboulle2003} carries the iterate, not the gradient, to the dual space: with a differentiable, strongly convex mirror map $\psi$ it sets $\nabla\psi(\theta_{t+1})=\nabla\psi(\theta_t)-\eta g_t$. With $\psi=\frac12\norm{\cdot}^2$, whose conjugate has the gradient $g\mapsto g^\sharp$ where it is differentiable, this would read $\theta_{t+1}=(\nabla\psi(\theta_t)-\eta g_t)^\sharp$, while steepest descent reads $\theta_{t+1}=\theta_t-\eta g_t^\sharp$. The two agree for all $\theta_t$ and $g_t$ exactly when the sharp operator is linear, that is, when the norm comes from an inner product, $\norm{u}=\norm{u}_P$, and both are then the preconditioned step $\theta_t-\eta P^{-1}g_t$. For the max, $\ell_1$ and spectral norms, $\frac12\norm{\cdot}^2$ is neither differentiable nor strictly convex, so mirror descent in these geometries uses other maps, such as the negative entropy, and it is analyzed through Bregman divergences rather than the descent inequality.
\end{remark}

\section{The theory: descent in a general norm}\label{sec:steepest:descent}

Smoothness gives a quadratic upper bound on the loss around every point, and steepest descent minimizes it; the main theorem of the chapter follows.

\begin{theorem}[Descent in a general norm]\label{thm:steepest:descent}
Let $\Loss$ satisfy \cref{ass:steepest:smooth}, let $\theta_0\in\R^d$, $\Delta_0=\Loss(\theta_0)-\Loss^\star$ and $g_t=\nabla\Loss(\theta_t)$.
\begin{enumerate}
\item (One step.) For all $\theta,u\in\R^d$,
\begin{equation}\label{eq:steepest:descentlemma}
  \Loss(\theta+u)\le\Loss(\theta)+\ip{\nabla\Loss(\theta),u}+\frac L2\norm{u}^2 .
\end{equation}
The right side is smallest at $u=-\nabla\Loss(\theta)^\sharp/L$, where it equals $\Loss(\theta)-\norm{\nabla\Loss(\theta)}_*^2/(2L)$. At $u=\eta\lmo(\nabla\Loss(\theta))$ it is at most $\Loss(\theta)-\eta\norm{\nabla\Loss(\theta)}_*+L\eta^2/2$.
\item (Unnormalized, nonconvex.) Unnormalized steepest descent with step size $1/L$ satisfies, for all $T\ge1$,
\[
  \min_{t<T}\norm{g_t}_*^2\le\frac1T\sum_{t<T}\norm{g_t}_*^2\le\frac{2L\Delta_0}{T}.
\]
\item (Normalized, nonconvex.) Normalized steepest descent with step size $\eta>0$ satisfies, for all $T\ge1$,
\begin{equation}\label{eq:steepest:normalized}
  \frac1T\sum_{t<T}\norm{g_t}_*\le\frac{\Delta_0}{\eta T}+\frac{L\eta}2 ,
\end{equation}
and if $\Delta_0>0$, the right side is smallest, equal to $\sqrt{2L\Delta_0/T}$, at $\eta=\sqrt{2\Delta_0/(LT)}$.
\item (Unnormalized, convex.) If moreover $\Loss$ is convex with a minimizer $\theta^\star$, and $R=\sup\{\norm{\theta-\theta^\star}:\Loss(\theta)\le\Loss(\theta_0)\}$ is finite, then unnormalized steepest descent with step size $1/L$ satisfies $\Delta_T\le2LR^2/(T+4)$ for all $T\ge0$.
\end{enumerate}
\end{theorem}

\begin{proof}
The plan is to prove the quadratic upper bound \eqref{eq:steepest:descentlemma} once, to minimize it over the two kinds of steps, and to add up the decreases; the convex rate needs one more inequality, which turns a gap into a lower bound on the gradient.

(i) Fix $\theta$ and $u$, and let $\varphi(s)=\Loss(\theta+su)$ for $s\in[0,1]$. The gradient is Lipschitz, hence continuous, so $\varphi$ is continuously differentiable with $\varphi'(s)=\ip{\nabla\Loss(\theta+su),u}$. By the fundamental theorem of calculus, then Hölder's inequality and smoothness,
\[
\begin{aligned}
  \Loss(\theta+u)-\Loss(\theta)-\ip{\nabla\Loss(\theta),u}&=\int_0^1\ip{\nabla\Loss(\theta+su)-\nabla\Loss(\theta),u}\dd s\\
  &\le\int_0^1Ls\norm{u}^2\dd s=\frac L2\norm{u}^2,
\end{aligned}
\]
which is \eqref{eq:steepest:descentlemma} (\cref{lem:toolkit:descent}). Write $g=\nabla\Loss(\theta)$. By Hölder's inequality the right side of \eqref{eq:steepest:descentlemma} is at least $\Loss(\theta)-\norm{g}_*\norm{u}+\frac L2\norm{u}^2$, a quadratic in $\norm{u}$ whose minimum over $\norm{u}\ge0$ is $\Loss(\theta)-\norm{g}_*^2/(2L)$. At $u=-g^\sharp/L$, \cref{lem:steepest:basic}(ii) gives $\ip{g,u}=-\norm{g}_*^2/L$ and $\norm{u}=\norm{g}_*/L$, so the right side equals this minimum. At $u=\eta\lmo(g)$, \cref{lem:steepest:basic}(i) gives $\ip{g,u}=-\eta\norm{g}_*$ and $\norm{u}\le\eta$.

(ii) By (i), $\Loss(\theta_{t+1})\le\Loss(\theta_t)-\norm{g_t}_*^2/(2L)$. Summing over $t<T$ and using $\Loss(\theta_T)\ge\Loss^\star$ gives $\sum_{t<T}\norm{g_t}_*^2\le2L\Delta_0$; a minimum is at most an average.

(iii) By (i), $\Loss(\theta_{t+1})\le\Loss(\theta_t)-\eta\norm{g_t}_*+L\eta^2/2$; if $g_t=0$, then the step is zero and the inequality holds too. Adding for $t<T$ gives $\eta\sum_{t<T}\norm{g_t}_*\le\Delta_0+TL\eta^2/2$, which is \eqref{eq:steepest:normalized} after division by $\eta T$. By the inequality between the arithmetic and geometric means, $\frac{\Delta_0}{\eta T}+\frac{L\eta}2\ge2\sqrt{L\Delta_0/(2T)}=\sqrt{2L\Delta_0/T}$, with equality when the two terms are equal, which for $\Delta_0>0$ happens at $\eta=\sqrt{2\Delta_0/(LT)}$.

(iv) The losses do not increase, by (i), so every $\theta_t$ lies in the sublevel set of $\theta_0$ and $\norm{\theta_t-\theta^\star}\le R$. By convexity and Hölder's inequality, $\Delta_t\le\ip{g_t,\theta_t-\theta^\star}\le R\norm{g_t}_*$, and with (i) this gives $\Delta_{t+1}\le\Delta_t-\Delta_t^2/(2LR^2)$. If some $\Delta_t$ is zero, then so are all later ones, and there is nothing to prove; so suppose that every $\Delta_t$ is positive, which forces $R>0$. Dividing the recursion by $\Delta_t\Delta_{t+1}$ and using $\Delta_{t+1}\le\Delta_t$,
\[
  \frac1{\Delta_{t+1}}-\frac1{\Delta_t}\ge\frac{\Delta_t}{2LR^2\,\Delta_{t+1}}\ge\frac1{2LR^2},
\]
so $1/\Delta_T\ge1/\Delta_0+T/(2LR^2)$. Finally, $\nabla\Loss(\theta^\star)=0$, so \eqref{eq:steepest:descentlemma} at $\theta^\star$ with $u=\theta_0-\theta^\star$ gives $\Delta_0\le LR^2/2$, that is, $1/\Delta_0\ge4/(2LR^2)$. Hence $1/\Delta_T\ge(T+4)/(2LR^2)$.
\end{proof}

\paragraph{Where the norm enters.} The theorem holds for every norm at once, and the norm enters only through the smoothness constant, the dual norm of the gradient and, for the convex rate, the radius $R$. Choosing a norm is choosing which smoothness constant and which gradient norm one pays; the guarantees come in different currencies ($\norm{g}_1$ for sign descent, $\norm{g}_2$ for gradient descent), and \cref{prop:steepest:whichnorm} converts between them for quadratics. Part (iv) is Nesterov's argument for gradient descent \parencite{Nesterov2018}, with the dual norm in place of the Euclidean one.\sidenote{In the Euclidean case $R$ can be replaced by $\norm{\theta_0-\theta^\star}_2$, because gradient descent with step size $1/L$ never moves away from a minimizer of a convex smooth loss (\cref{lem:toolkit:descent}). The argument expands the square of a Euclidean distance, and in a general norm the radius of the sublevel set takes its place.}

\paragraph{Normalized and unnormalized.} At their best step sizes the two methods have the same guarantee: by (ii) and the Cauchy--Schwarz inequality, the average of $\norm{g_t}_*$ under unnormalized descent is at most $\sqrt{2L\Delta_0/T}$, the value of (iii) at its best step. They differ in what they risk. An unnormalized step of size $\eta$ has the guaranteed decrease $\eta(1-L\eta/2)\norm{g_t}_*^2$, which vanishes at $\eta=2/L$, the stability threshold of the prologue. A normalized step never moves farther than $\eta$, so \eqref{eq:steepest:normalized} holds for every $\eta>0$: there is no threshold, a step size that is too large only raises the floor, and the method is clipped by construction (\cref{ch:spikes}).

\paragraph{The floor.} At a fixed step size the right side of \eqref{eq:steepest:normalized} tends to $L\eta/2$, not to zero\index{floor of a normalized method}: a normalized step has length $\eta$ even where the gradient is tiny, so the iterates cannot settle, as the prologue showed for sign descent. A decaying learning rate removes the floor, since with step sizes $\eta_t$ the same argument gives $\sum_t\eta_t\norm{g_t}_*\le\Delta_0+\frac L2\sum_t\eta_t^2$ (\cref{ex:steepest:schedule}); that a schedule removes the floor of a constant learning rate is a theme of the book (\cref{ch:schedules,ch:decay}). The floor takes over at $T=2\Delta_0/(L\eta^2)$, where the two terms of \eqref{eq:steepest:normalized} are equal. To compare norms we need their smoothness constants.

\begin{proposition}[Which norm? Smoothness constants of a quadratic]\label[proposition]{prop:steepest:whichnorm}
Let $\Loss(\theta)=\frac12\theta^\top H\theta-\ip{b,\theta}$ with $H$ symmetric, positive semidefinite and nonzero. For a norm $\norm{\cdot}$ let $L_{\norm\cdot}$ be the smallest $L$ for which $\Loss$ is $L$-smooth with respect to $\norm{\cdot}$.
\begin{enumerate}
\item For every norm, $L_{\norm\cdot}=\max_{\norm{u}\le1}u^\top Hu$.
\item $L_{\ell_2}=\lmax(H)$, $L_{\ell_1}=\max_iH_{ii}$, and $L_{\ell_\infty}=\norm{H}_{\infty\to1}=\max_{s\in\{-1,1\}^d}s^\top Hs$, which satisfies $\tr H\le L_{\ell_\infty}\le d\,\lmax(H)$, with $L_{\ell_\infty}=\tr H$ when $H$ is diagonal.
\item At a point with gradient $g\ne0$, the guaranteed decrease $\norm{g}_*^2/(2L_{\norm\cdot})$ of \cref{thm:steepest:descent}(i) is larger for $\ell_\infty$ than for $\ell_2$ if and only if
\begin{equation}\label{eq:steepest:whichnorm}
  \frac{\norm{g}_1^2}{\norm{g}_2^2}>\frac{\norm{H}_{\infty\to1}}{\lmax(H)},
\end{equation}
and the right side is at least $\tr H/\lmax(H)$, with equality when $H$ is diagonal.
\item For every norm, the step $-g^\sharp/L_{\norm\cdot}$ decreases $\Loss$ by at least $\norm{g}_*^2/(2L_{\norm\cdot})$ and at most twice as much, and by exactly the lower value if and only if $\lmo(g)$ maximizes $u^\top Hu$ over the unit ball; for $\ell_\infty$ and diagonal $H$, this holds whenever no coordinate of $g$ vanishes.
\end{enumerate}
\end{proposition}

\begin{proof}
The plan is to reduce every constant to the maximum of the quadratic form $u^\top Hu$ over the unit ball, and to evaluate that maximum at the vertices of the ball. (i) Since $\nabla\Loss(\theta)-\nabla\Loss(\theta')=H(\theta-\theta')$, the best constant is $\sup_{u\ne0}\norm{Hu}_*/\norm{u}=\sup\{v^\top Hu:\norm{u}\le1,\ \norm{v}\le1\}$. Taking $v=u$ shows that this is at least $\max_{\norm{u}\le1}u^\top Hu$. Conversely, $(u,v)\mapsto v^\top Hu$ is a positive semidefinite bilinear form, so the Cauchy--Schwarz inequality gives $v^\top Hu\le(u^\top Hu)^{1/2}(v^\top Hv)^{1/2}\le\max\{u^\top Hu,v^\top Hv\}$.

(ii) For $\ell_2$, (i) is the Rayleigh quotient. For $\ell_1$ and $\ell_\infty$, the function $u\mapsto u^\top Hu$ is convex, and a convex function on a polytope attains its maximum at a vertex, because every point of the polytope is a convex combination of vertices. The vertices of the $\ell_1$ ball are $\pm e_i$, where the function equals $H_{ii}$; those of the $\ell_\infty$ ball are the sign vectors. That $\max_ss^\top Hs=\max_{\norm{u}_\infty\le1}\norm{Hu}_1=\norm{H}_{\infty\to1}$ is (i) with the dual norm $\ell_1$. For $s$ uniform on $\{-1,1\}^d$, $\E[s_is_j]=\ind[i=j]$, so $\E[s^\top Hs]=\tr H$, and the maximum is at least the mean. For every $s$, $s^\top Hs\le\lmax(H)\norm{s}_2^2=d\,\lmax(H)$; and if $H$ is diagonal, then $s^\top Hs=\tr H$ for every $s$.

(iii) Compare the guarantees $\norm{g}_1^2/(2L_{\ell_\infty})$ and $\norm{g}_2^2/(2L_{\ell_2})$, and use (ii).

(iv) Let $L=L_{\norm\cdot}$, $u=\lmo(g)$ and $q=u^\top Hu$, so that $0\le q\le L$ by (i) and \cref{lem:steepest:basic}(i). For a quadratic, \eqref{eq:steepest:descentlemma} holds with equality once $\frac L2\norm{v}^2$ is replaced by $\frac12v^\top Hv$; with $v=-g^\sharp/L=(\norm{g}_*/L)\,u$,
\[
  \Loss(\theta)-\Loss(\theta+v)=\frac{\norm{g}_*^2}{L}-\frac{\norm{g}_*^2}{2L^2}\,q=\frac{\norm{g}_*^2}{2L}\Bigl(2-\frac qL\Bigr),
\]
which lies between $\norm{g}_*^2/(2L)$ and $\norm{g}_*^2/L$ and equals the first exactly when $q=L$. For $\ell_\infty$ and diagonal $H$, $q=\sum_iH_{ii}\sign(g_i)^2=\tr H=L$ when every $g_i\ne0$.
\end{proof}

The left side of \eqref{eq:steepest:whichnorm} is the \term{density}\index{density of a gradient} of the gradient, between $1$ for a gradient with one nonzero coordinate and $d$ for one with coordinates of equal size; the ratio $\tr H/\lmax(H)$, between $1$ and $\rank H$, is the \term{effective rank}\index{effective rank} of the curvature. So the max norm wins a step, on the guarantees, when the gradient is spread over more coordinates than the curvature over directions: a dense gradient against a curvature of low effective rank. A gradient with independent Gaussian coordinates has density close to $2d/\pi$, so for diagonal $H$ the max norm then wins unless the effective rank exceeds $2d/\pi$, that is, unless the curvature is nearly isotropic. By (iv) the guarantees decide the comparison of the true decreases whenever they differ by more than a factor of two; \cref{ch:sign,ch:transformers} measure both sides in networks. But $\norm{H}_{\infty\to1}$ is NP-hard to compute \parencite{Rohn2000}, since for a graph Laplacian it is four times the maximum cut (\cref{ex:steepest:maxcut}); for diagonal $H$ the difficulty disappears, and there the theory can be tested exactly.

\section{Reading the theory}\label{sec:steepest:reading}

The theory predicts the best step length $\norm{g}_*/L$ and a guaranteed decrease within a factor of two of the truth on convex quadratics, the comparison \eqref{eq:steepest:whichnorm} between norms, a $1/T$ transient above a floor $L\eta/2$, and the convex rate $2LR^2/(T+4)$. \Cref{fig:steepest:geometry} shows the geometry of a step.\sidenote{The emblem on the title page of this book is the left panel of \cref{fig:steepest:geometry}, drawn for $g$ instead of $-g$.}

\begin{figure}[t]
\begin{fullwidth}
\centering
\begin{tikzpicture}
\begin{groupplot}[group style={group size=2 by 1,horizontal sep=15mm},wplot,width=57mm,height=57mm,axis equal image]
\nextgroupplot[xmin=-1.7,xmax=1.7,ymin=-1.7,ymax=1.7,xtick={-1,0,1},ytick={-1,0,1},xlabel={$u_1$},ylabel={$u_2$}]
\addplot[sea,line width=0.9pt] table[x=x,y=y] {figures/data/01-steepest-geometry-circle.dat};
\addplot[seadeep,line width=0.9pt] coordinates {(1,1) (-1,1) (-1,-1) (1,-1) (1,1)};
\addplot[sandline!55!black,line width=0.9pt] coordinates {(1,0) (0,1) (-1,0) (0,-1) (1,0)};
\addplot[sandline!55!black,densely dashed,line width=0.6pt,quiver={u=\thisrow{du},v=\thisrow{dv}},skip coords between index={1}{3}] table[x=ax,y=ay] {figures/data/01-steepest-geometry-lmo.dat};
\addplot[sea,densely dashed,line width=0.6pt,quiver={u=\thisrow{du},v=\thisrow{dv}},skip coords between index={0}{1},skip coords between index={2}{3}] table[x=ax,y=ay] {figures/data/01-steepest-geometry-lmo.dat};
\addplot[seadeep,densely dashed,line width=0.6pt,quiver={u=\thisrow{du},v=\thisrow{dv}},skip coords between index={0}{2}] table[x=ax,y=ay] {figures/data/01-steepest-geometry-lmo.dat};
\addplot[muted,-stealth,line width=0.7pt,quiver={u=\thisrow{x}*1.62,v=\thisrow{y}*1.62},skip coords between index={0}{1},skip coords between index={2}{3}] table[x expr=0,y expr=0] {figures/data/01-steepest-geometry-lmo.dat};
\addplot[sandline!55!black,-stealth,line width=1.1pt,quiver={u=\thisrow{x},v=\thisrow{y}},skip coords between index={1}{3}] table[x expr=0,y expr=0] {figures/data/01-steepest-geometry-lmo.dat};
\addplot[sea,-stealth,line width=1.1pt,quiver={u=\thisrow{x},v=\thisrow{y}},skip coords between index={0}{1},skip coords between index={2}{3}] table[x expr=0,y expr=0] {figures/data/01-steepest-geometry-lmo.dat};
\addplot[seadeep,-stealth,line width=1.1pt,quiver={u=\thisrow{x},v=\thisrow{y}},skip coords between index={0}{2}] table[x expr=0,y expr=0] {figures/data/01-steepest-geometry-lmo.dat};
\fill[ink] (axis cs:0,0) circle (1.1pt);
\node[font=\small,text=muted,anchor=south,inner sep=1.5pt] at (axis cs:1.40,0.80) {$-g$};
\node[font=\small,text=seadeep,anchor=south west,inner sep=1pt] at (axis cs:-1.0,1.0) {$\ell_\infty$};
\node[font=\small,text=sea,anchor=north east,inner sep=1pt] at (axis cs:-0.72,-0.72) {$\ell_2$};
\node[font=\small,text=sandline!55!black,anchor=north,inner sep=1.5pt] at (axis cs:0.62,-0.42) {$\ell_1$};
\nextgroupplot[xmin=-2.75,xmax=1.15,ymin=-1.2,ymax=2.7,xtick={-2,-1,0,1},ytick={-1,0,1,2},xlabel={$\theta_1$},ylabel={$\theta_2$},legend style={at={(0.97,0.97)},anchor=north east,font=\scriptsize}]
\addplot[muted!55,line width=0.4pt,forget plot] table[x=xa,y=ya] {figures/data/01-steepest-geometry-levels.dat};
\addplot[muted!55,line width=0.4pt,forget plot] table[x=xb,y=yb] {figures/data/01-steepest-geometry-levels.dat};
\addplot[muted!55,line width=0.4pt,forget plot] table[x=xc,y=yc] {figures/data/01-steepest-geometry-levels.dat};
\addplot[muted!55,line width=0.4pt,forget plot] table[x=xd,y=yd] {figures/data/01-steepest-geometry-levels.dat};
\addplot[sea,line width=0.8pt,mark=*,mark size=0.9pt] table[x=x2,y=y2] {figures/data/01-steepest-geometry-traj.dat};
\addplot[seadeep,line width=0.8pt,mark=square*,mark size=0.9pt] table[x=xinf,y=yinf] {figures/data/01-steepest-geometry-traj.dat};
\addplot[sandline!55!black,line width=0.8pt,mark=triangle*,mark size=1.1pt] table[x=x1,y=y1] {figures/data/01-steepest-geometry-traj.dat};
\fill[ink] (axis cs:-2.4,1.3) circle (1.2pt);
\legend{$\ell_2$,$\ell_\infty$,$\ell_1$}
\end{groupplot}
\end{tikzpicture}
\end{fullwidth}
\caption{The geometry of a step. Left: the unit balls of $\ell_1$ (diamond), $\ell_2$ (circle) and $\ell_\infty$ (square), the negative gradient $-g$ at $28^\circ$ with $\norm{g}_2=1$, and the steepest unit direction $\lmo(g)$ of each norm (thick arrows), where the ball touches its supporting line $\{u:\ip{g,u}=-\norm{g}_*\}$ (dashed). The lines lie at distances $\norm{g}_\infty=0.88$, $\norm{g}_2=1$ and $\norm{g}_1=1.35$ from the origin: the dual norm measures how far the ball reaches toward $-g$. Right: $40$ normalized steps in each norm with $\eta=0.2$ from one start (dot), on the quadratic with Hessian $R\diag(1,0.12)R^\top$, $R$ the rotation by $30^\circ$; all three runs end on cycles of period two. Script: \texttt{01-steepest-geometry.py}.}
\label{fig:steepest:geometry}
\end{figure}

\paragraph{One step.} The left panel of \cref{fig:steepest:descent} tests the first two predictions on a diagonal quadratic in $d=1000$ dimensions with curvatures $h_i=1/i$, so that $\lmax(H)=1$ and $\tr H=7.49$, at a random point. Its gradient has density $30.1$ against an effective rank of $7.49$, so \eqref{eq:steepest:whichnorm} predicts that the sign step wins, and the guarantees differ by a factor of $4.02$ ($2.43$ against $0.60$). The sign step's true decrease equals its guarantee at every step length, as \cref{prop:steepest:whichnorm}(iv) requires for a diagonal Hessian. The normalized gradient step beats its guarantee, because the curvature along $g/\norm{g}_2$ is $q=0.65$, not $\lmax=1$: at its recommended length $\norm{g}_2=1.10$ it gains $0.82$, which is $2-q=1.35$ times its guarantee. The sign step still wins, by a factor of $2.97$ at the recommended lengths and $2.60$ against the best Euclidean step. The exact agreement of the sign step must fail for a non-diagonal Hessian, where the curvature along $-\sign(g)$ is below $\norm{H}_{\infty\to1}$, and a rotation of the coordinates can reverse the comparison (\cref{ch:sign}).

\begin{figure}[t]
\begin{fullwidth}
\centering
\begin{tikzpicture}
\begin{groupplot}[group style={group size=2 by 1,horizontal sep=17mm},wplot,width=82mm,height=49mm]
\nextgroupplot[xlabel={step length $\eta$},ylabel={decrease $\Loss(\theta_0)-\Loss(\theta_1)$},xmin=0,xmax=2.4,ymin=-0.4,ymax=2.75,xtick={0,0.5,1,1.5,2},legend style={at={(0.98,0.98)},anchor=north east,font=\scriptsize},legend columns=1]
\addplot[guide,forget plot] coordinates {(0,0) (2.4,0)};
\addplot[theory] table[x=eta,y=bdinf] {figures/data/01-steepest-descent-onestep.dat};
\addplot[theory2] table[x=eta,y=bdtwo] {figures/data/01-steepest-descent-onestep.dat};
\addplot[simulation,each nth point=3] table[x=eta,y=acinf] {figures/data/01-steepest-descent-onestep.dat};
\addplot[simulation2,each nth point=3] table[x=eta,y=actwo] {figures/data/01-steepest-descent-onestep.dat};
\legend{guarantee $\eta\norm{g}_1-\eta^2\tr H/2$,guarantee $\eta\norm{g}_2-\eta^2\lmax/2$,sign step,normalized gradient step}
\nextgroupplot[xmode=log,ymode=log,xlabel={number of steps $T$},ylabel={$\frac1T\sum_{t<T}\norm{\nabla\Loss(\theta_t)}_1$},xmin=1,xmax=1e5,ymin=1e-2,ymax=40,xtick={1,10,100,1000,10000,100000},ytick={0.01,0.1,1,10},legend style={at={(0.03,0.03)},anchor=south west,font=\scriptsize}]
\addplot[theory] table[x=T,y=bdA] {figures/data/01-steepest-descent-floor.dat};
\addplot[simulation] table[x=T,y=avgA] {figures/data/01-steepest-descent-floor.dat};
\addplot[theory,forget plot] table[x=T,y=bdB] {figures/data/01-steepest-descent-floor.dat};
\addplot[simulation,forget plot] table[x=T,y=avgB] {figures/data/01-steepest-descent-floor.dat};
\addplot[theory3,densely dotted] table[x=T,y=env] {figures/data/01-steepest-descent-floor.dat};
\addplot[guide,forget plot,domain=1:1e5] {0.374274};
\addplot[guide,forget plot,domain=1:1e5] {0.0374274};
\node[font=\scriptsize,text=ink,anchor=south] at (axis cs:3e4,0.40) {$\eta=0.1$};
\node[font=\scriptsize,text=ink,anchor=north] at (axis cs:3e4,0.034) {$\eta=0.01$};
\legend{bound \eqref{eq:steepest:normalized},sign descent,best bound $\sqrt{2\Delta_0\tr H/T}$}
\end{groupplot}
\end{tikzpicture}
\end{fullwidth}
\caption{The descent theorem on $\Loss(\theta)=\frac12\sum_ih_i\theta_i^2$ with $d=1000$ and $h_i=1/i$, from $\theta_0$ with independent standard Gaussian coordinates (numpy seed $9$, whose ratio of guarantees, $4.02$, is near the median $3.91$ over seeds $0$ to $199$; $\Delta_0=3.64$). Left: the decrease achieved by one sign step $-\eta\sign(g)$ and one normalized gradient step $-\eta g/\norm{g}_2$ (points), against the guarantees of \cref{thm:steepest:descent}(i) with $L=\tr H$ and $L=\lmax(H)$ (curves). Right: the running average of $\norm{\nabla\Loss(\theta_t)}_1$ along sign descent at two step sizes (points), against the bound \eqref{eq:steepest:normalized} computed in advance from $\Delta_0$, $\tr H$ and $\eta$ (curves); dashed, the floors $\eta\tr H/2$; dotted, the best bound at each horizon. Script: \texttt{01-steepest-descent.py}.}
\label{fig:steepest:descent}
\end{figure}

\paragraph{The floor.} The right panel follows sign descent on the same problem for $10^5$ steps at two step sizes. Early in each run the points lie below the curve: the bound charges the whole initial gap $\Delta_0$ to the first $T$ steps, while a run that has moved each coordinate by only $T\eta$ has spent part of it. Once every coordinate has reached its cycle, after $\max_i\abs{\theta_{0,i}}/\eta$ steps (about $31$ for $\eta=0.1$ and $310$ for $\eta=0.01$), the points lie on the curve; at $T=10^5$ the two agree to within three parts in a million. Each run then follows its curve down to its floor $\eta\tr H/2$, reached near $T=10^2$ for $\eta=0.1$ and near $T=10^4$ for $\eta=0.01$, and touches the best bound (dotted) only at its crossover $2\Delta_0/(\eta^2\tr H)$, which is $97$ and $9.7\times10^3$.

\begin{proposition}[When the bound is attained]\label[proposition]{prop:steepest:attained}
Let $\Loss(\theta)=\frac12\theta^\top H\theta$ with $H$ positive semidefinite.
\begin{enumerate}
\item If normalized steepest descent in a norm $\norm{\cdot}$ with step size $\eta$ has a cycle of period two, $\theta\mapsto\theta+\eta u\mapsto\theta$, then the average of $\norm{\nabla\Loss}_*$ over the cycle is $\frac\eta2u^\top Hu\le L_{\norm\cdot}\eta/2$.
\item Let $H=\diag(h_1,\dots,h_d)$ with all $h_i>0$, and let sign descent start at a point $\theta_0$ none of whose coordinates is an integer multiple of $\eta$. Then for every $T\ge1$,
\begin{equation}\label{eq:steepest:identity}
  \frac1T\sum_{t<T}\norm{g_t}_1=\frac{\Delta_0-\Delta_T}{\eta T}+\frac{\eta\tr H}2 ,
\end{equation}
and from step $\max_i\abs{\theta_{0,i}}/\eta$ on the iterates lie on a cycle of period two, with $\Delta_T<\frac12\eta^2\tr H$. So the two sides of \eqref{eq:steepest:normalized}, with $L=\tr H=L_{\ell_\infty}$, differ by exactly $\Delta_T/(\eta T)$, and the running average tends to the floor $\eta\tr H/2$.
\end{enumerate}
\end{proposition}

\begin{proof}
The plan is to evaluate the exact expansion of the quadratic along the steps. (i) Write $\theta'=\theta+\eta u$, $g=\nabla\Loss(\theta)$ and $g'=\nabla\Loss(\theta')=g+\eta Hu$. The cycle means $u=\lmo(g)$ and $-u=\lmo(g')$, so by \cref{lem:steepest:basic}(i), $\norm{g}_*=-\ip{g,u}$ and $\norm{g'}_*=\ip{g',u}=\ip{g,u}+\eta u^\top Hu$. Adding, $\norm{g}_*+\norm{g'}_*=\eta u^\top Hu$, and $u^\top Hu\le L_{\norm\cdot}$ by \cref{prop:steepest:whichnorm}(i).

(ii) Each coordinate of $\theta_t$ changes by $\pm\eta$ per step, so under the hypothesis no coordinate of $\theta_t$, and hence none of $g_t=H\theta_t$, ever vanishes. Then $\sign(g_t)^\top H\sign(g_t)=\tr H$, and the exact expansion gives $\Loss(\theta_{t+1})=\Loss(\theta_t)-\eta\norm{g_t}_1+\frac{\eta^2}2\tr H$ at every step; adding over $t<T$ gives \eqref{eq:steepest:identity}. By the prologue's theorem on sign descent, from step $\max_i\abs{\theta_{0,i}}/\eta$ on each coordinate alternates between two values of absolute value less than $\eta$, so the iterates cycle and $\Delta_T<\frac12\eta^2\tr H$. Hence $\Delta_T/(\eta T)\to0$.
\end{proof}

For normalized gradient descent the same computation gives \eqref{eq:steepest:identity} with $\frac\eta2\cdot\frac1T\sum_tq_t$ in place of $\frac{\eta\tr H}2$, where $q_t$ is the curvature along $g_t/\norm{g_t}_2$, so its bound is attained only by runs that step along a top eigenvector. Its cycles lie along eigenvectors, since on a cycle $g$ and $g+\eta Hu$ are both parallel to $u$, and by part (i) a cycle along an eigenvector with eigenvalue $\lambda_k$ has average gradient norm $\eta\lambda_k/2$. In the plane of \cref{fig:steepest:geometry} all three methods end on cycles of maximal curvature and attain their floors exactly (\cref{ex:steepest:plane}); in high dimension a Euclidean run can linger on a cycle along an eigenvector with $\lambda_k<\lmax$.

\paragraph{The best step depends on the criterion.} By \cref{prop:steepest:attained}(ii) the best fixed step for the criterion of the theorem is very nearly $\sqrt{2\Delta_0/(T\tr H)}$, which shrinks like $T^{-1/2}$. For the final loss, the criterion of a practitioner with a budget, the best step on the same problem shrinks like $T^{-1}$: it must cover the distance $\norm{\theta_0-\theta^\star}_\infty$ in the time available, and a longer step only raises the loss on the final cycle (\cref{ex:steepest:beststep}). How the best learning rate scales with the horizon depends on what is measured.

\paragraph{Worst-case rates do not rank the norms.} The convex rate compares norms through $LR^2$. For a diagonal quadratic the sublevel set is an ellipsoid with axes along the coordinates, whose $\ell_\infty$ and $\ell_2$ radii both equal its longest semi-axis $R$. So $L_{\ell_\infty}R^2=\tr H\cdot R^2\ge L_{\ell_2}R^2$: the convex rate never prefers the max norm on such a problem, while the step of \cref{fig:steepest:descent} prefers it by a factor of three. The worst point of a sublevel set has its gradient on one flat coordinate, the worst case for the max norm, and the gradient at hand need not look like that (\cref{ex:steepest:pl}).

\section{What the theory misses}\label{sec:steepest:misses}

\paragraph{What was idealized away.} A fixed norm: practical optimizers change their norm at every step, through Adam's second moments or Shampoo's Kronecker factors, and a theory of one norm describes them only with their accumulators switched off (\cref{ch:adaptive,ch:kronecker}). A global smoothness constant: the curvature of a network changes during training, and at practical step sizes it settles at the largest value the step size tolerates (\cref{ch:eos}). Exact gradients: with noise, $\lmo(g_t)$ is a biased estimate of $\lmo(\nabla\Loss(\theta_t))$, because the oracle is not linear (\cref{ch:momentum}). No memory, and one vector instead of matrices in layers (\cref{ch:spectral,ch:modular}).

\paragraph{Where it breaks.} At the step sizes of practice the one-step guarantee is void. Full-batch gradient descent on neural networks drives the sharpness up to $2/\eta$ and trains there \parencite{CohenEtAl2021}, where the Euclidean guarantee $\eta(1-L\eta/2)\norm{g}_2^2$, with $L$ the sharpness, is zero; the loss still falls, over long times and not monotonically (\cref{ch:eos}). Smoothness constants are not constant: along the training of language models and image classifiers, \textcite{ZhangEtAl2020Clipping} measured a local smoothness that grows with the gradient norm, so a global $L$ is absent or uselessly large (\cref{ch:spikes}). The rates do not rank optimizers: on transformers, the gap between Adam and gradient descent is attributed to mechanisms such as the sign-like geometry of the step \parencite{KunstnerEtAl2023}, not to a worst-case rate (\cref{ch:transformers}). The right norm is a modeling choice that the theorem cannot identify, since it holds for all norms at once and its inputs are not computable for networks. And the jumps of the oracle at ties break the deterministic picture once gradients are noisy: \textcite{KarimireddyEtAl2019} give simple convex problems on which sign descent with stochastic gradients does not converge to the optimum (\cref{ch:sign}).

\paragraph{Repairs and open problems.} Smoothness that is local, anisotropic or tied to the gradient repairs the second failure: \textcite{XieEtAl2025} analyze Adam under smoothness in the max norm and argue that this constant is favorable for language models in their natural coordinates, and $(L_0,L_1)$-smoothness \parencite{ZhangEtAl2020Clipping} and curvature bounded by the loss are the subjects of \cref{ch:spikes,ch:warmup}. Stochastic steepest descent with momentum in general norms is analyzed by \textcite{PethickEtAl2025}, with a norm-constrained oracle in each layer, and by \textcite{Kovalev2025}, as a trust-region method (\cref{ch:momentum}); norms built from the architecture are the modular norms of \textcite{LargeEtAl2024,BernsteinNewhouse2025} (\cref{ch:modular}). For normalized gradient descent at a fixed step, \textcite{AroraLiPanigrahi2022} show that near a manifold of minimizers the oscillation at the floor drives the iterates toward lower sharpness, so on a loss that is not quadratic the floor has dynamics of its own. Open: a choice of norm from data (\cref{ex:steepest:choose}); the right norm for embeddings, attention and normalization layers; and a theory for the step sizes of practice.

\section{Back in the wild}\label{sec:steepest:back}

\fieldnotes{Fixed-length steps: per coordinate (Lion, roughly Adam), per matrix (Muon), per layer (LARS, LAMB).}
{Steepest descent in one fixed norm; exact gradients; a loss that is $L$-smooth in that norm.}
{Duality maps; the decrease $\norm{g}_*^2/(2L)$, the rates and the floor $L\eta/2$; smoothness constants of quadratics.}
{No stability threshold for normalized steps; the $\ell_\infty$ guarantee beats $\ell_2$'s if $\norm{g}_1^2/\norm{g}_2^2>\norm{H}_{\infty\to1}/\lmax(H)$.}
{Global smoothness (edge of stability), noise (the oracle jumps), memory, layers, the choice of norm.}
{Match update sizes when moving a learning rate between optimizers (suggested); decay a normalized step to go below its floor (proved in the idealization).}

The theory tells a practitioner what a learning rate is: in every optimizer of the practice box it is the radius of a ball in some norm, and the norm decides which directions a step can take. It explains why matching update sizes is a sensible way to move a learning rate between optimizers, and why normalized optimizers tolerate step sizes at which gradient descent diverges, at the price of a floor that only a schedule removes. It does not say which norm to use, because the comparison depends on the gradients that training will produce. Part~I goes on to the max norm (\cref{ch:sign}), the spectral norm (\cref{ch:spectral}) and the norm of a network (\cref{ch:modular}); the floor returns as the floor that momentum and schedules remove (\cref{ch:momentum,ch:schedules}), and the oracle as the Frank--Wolfe step of decoupled weight decay (\cref{ch:decay}).

\begin{exercises}
\item \consolidation (Duality maps of $\ell_p$.) Let $1<p<\infty$, $1/p+1/q=1$ and $g\ne0$. (a) Show that the dual of $\norm{\cdot}_p$ is $\norm{\cdot}_q$ and that the unique steepest unit direction is $\lmo(g)_i=-\sign(g_i)\abs{g_i}^{q-1}/\norm{g}_q^{q-1}$, so that $g^\sharp_i=\norm{g}_q^{2-q}\sign(g_i)\abs{g_i}^{q-1}$. (b) Find the limits of $\lmo(g)$ as $p\to\infty$ and $p\to1$, and compare with \cref{prop:steepest:dualitymaps}(ii) and (iii); if the largest $\abs{g_i}$ is attained $k$ times, which point of the set in (iii) is the limit? \hint{equality holds in $\sum_ig_iu_i\ge-\norm{g}_q\norm{u}_p$ exactly when $\abs{u_i}^p$ is proportional to $\abs{g_i}^q$ and the signs are opposite.}\label{ex:steepest:lp}

\item \consolidation (Sums and maxima of norms.) With blocks as in \cref{prop:steepest:dualitymaps}(iv) and scales $s_k>0$: (a) show that the dual of $\max_ks_k\norm{\theta_k}_{(k)}$ is $\sum_k\norm{g_k}_{(k),*}/s_k$ with oracle $(\lmo_{(k)}(g_k)/s_k)_k$, so normalized descent in a weighted maximum norm is blockwise normalized descent with learning rates $\eta/s_k$; (b) show that the dual of $\sum_ks_k\norm{\theta_k}_{(k)}$ is $\max_k\norm{g_k}_{(k),*}/s_k$, and that its steepest unit directions put the whole step on the blocks that attain the maximum; (c) show that unnormalized descent in the norm of (a) couples all blocks through the single factor $\sum_j\norm{g_j}_{(j),*}/s_j$ (\cref{ch:modular}).\label{ex:steepest:blocks}

\item \consolidation (The constants of $\ell_\infty$ and $\ell_1$.) (a) For a graph Laplacian $H=D-A$, show that $s^\top Hs=4\abs{E(S,S^c)}$ for $s\in\{-1,1\}^d$ and $S=\{i:s_i=1\}$, where $E(S,S^c)$ is the set of edges leaving $S$; conclude that $L_{\ell_\infty}$ is four times the maximum cut, so that computing it is NP-hard. (b) For $H=\diag(h)+c\mathbf1\mathbf1^\top$ with $h_i,c\ge0$, show that $L_{\ell_\infty}=\sum_ih_i+cd^2$; for $h=0$ this attains the upper bound $d\,\lmax(H)$ of \cref{prop:steepest:whichnorm}(ii). (c) Show that the $\ell_1$ guarantee beats the $\ell_2$ guarantee if and only if $\norm{g}_\infty^2/\norm{g}_2^2>\max_iH_{ii}/\lmax(H)$, and evaluate both sides for $H=I+\mathbf1\mathbf1^\top$.\label{ex:steepest:maxcut}

\item \consolidation (Step length and update size.) (a) Show that a normalized step of length $\eta$ has root-mean-square $\eta$ in $\ell_\infty$ (if no $g_i$ vanishes), $\eta/\sqrt d$ in $\ell_2$ and $\ell_1$, and $\eta/\sqrt{\max(m,n)}$ in the spectral norm of a full-rank $m\times n$ matrix, and that the rescaled Muon update $0.2\sqrt{\max(m,n)}\,\eta\polar(G)$ has root-mean-square $0.2\eta$. (b) For $W\in\R^{d_{\mathrm{out}}\times d_{\mathrm{in}}}$ and $\norm{W}_{\rms\to\rms}=\sqrt{d_{\mathrm{in}}/d_{\mathrm{out}}}\,\norm{W}_{\op}$, show that the steepest unit direction is $-\sqrt{d_{\mathrm{out}}/d_{\mathrm{in}}}\,\polar(G)$, and compare how the root-mean-squares of the two rules scale with the width of a square matrix (\cref{ch:modular}).\label{ex:steepest:rms}

\item \variation (A quadratic norm is a preconditioner.) Let $P\succ0$ and $\norm{u}_P=(u^\top Pu)^{1/2}$. (a) Show that the dual norm is $\norm{g}_{P^{-1}}$, $\lmo(g)=-P^{-1}g/\norm{g}_{P^{-1}}$ and $g^\sharp=P^{-1}g$: unnormalized steepest descent in $\norm{\cdot}_P$ is preconditioned gradient descent. (b) For $\Loss(\theta)=\frac12\theta^\top H\theta$ show that $L_{\norm{\cdot}_P}=\lmax(P^{-1}H)$, the preconditioned sharpness $S_P$. (c) For $H\succ0$, show that among quadratic norms the guarantee $g^\top P^{-1}g/(2\lmax(P^{-1}H))$ is largest for $P=H$, where it is the whole gap $\frac12g^\top H^{-1}g$: Newton's method. (d) For $\theta=A\varphi$ with $A$ invertible, show that steepest descent on $\varphi\mapsto\Loss(A\varphi)$ in the norm $\varphi\mapsto\norm{A\varphi}$ gives $\varphi_t=A^{-1}\theta_t$, while gradient descent in $\varphi$ is steepest descent in $\theta$ for the norm $\norm{A^{-1}\theta}_2$. \hint{for (c), with $\mu=\lmax(P^{-1}H)=\lmax(H^{1/2}P^{-1}H^{1/2})$, show that $P^{-1}\preceq\mu H^{-1}$.}\label{ex:steepest:quadnorm}

\item \variation (Polyak--Łojasiewicz in a dual norm; see \cref{lem:toolkit:pl}.) Suppose that, besides \cref{ass:steepest:smooth}, $\norm{\nabla\Loss(\theta)}_*^2\ge2\mu(\Loss(\theta)-\Loss^\star)$ for all $\theta$. (a) Show that unnormalized steepest descent with step size $1/L$ has $\Delta_t\le(1-\mu/L)^t\Delta_0$. (b) For $\Loss(\theta)=\frac12\sum_ih_i\theta_i^2$ with $h_1\ge\dots\ge h_d>0$, show that the best constant is $\mu=h_d$ both for $\ell_2$ and for $\ell_\infty$ (whose dual is $\ell_1$), so the guaranteed rates are $1-h_d/h_1$ and $1-h_d/\tr H$. (c) Conclude that the worst-case linear rate always prefers $\ell_2$ on a diagonal quadratic, and reconcile this with \cref{fig:steepest:descent}.\label{ex:steepest:pl}

\item \variation (A schedule removes the floor.) Run normalized steepest descent with step sizes $\eta_t>0$. (a) Show that $\sum_{t<T}\eta_t\norm{g_t}_*\le\Delta_0+\frac L2\sum_{t<T}\eta_t^2$. (b) Deduce that $\liminf_t\norm{g_t}_*=0$ if $\sum_t\eta_t=\infty$ and $\sum_t\eta_t^2<\infty$. (c) With $\eta_t=\eta_0/\sqrt{t+1}$, show that $\min_{t<T}\norm{g_t}_*\le\bigl(\Delta_0+\frac L2\eta_0^2(1+\ln T)\bigr)/\bigl(2\eta_0(\sqrt{T+1}-1)\bigr)$, and compare with the prologue's sign descent with a decaying step.\label{ex:steepest:schedule}

\item \simulation (Three norms in the plane.) For the quadratic of \cref{fig:steepest:geometry}, right (script \texttt{01-steepest-geometry.py}), (a) check that $L_{\ell_2}=1$, $L_{\ell_\infty}=H_{11}+H_{22}+2\abs{H_{12}}=1.882$ and $L_{\ell_1}=\max_iH_{ii}=0.780$. (b) For $\eta\in\{0.2,0.05,0.01\}$, run each method for $400/\eta$ steps and average the dual norm of the gradient over the last $100$ steps; you should find $L_{\norm\cdot}\eta/2$ in each norm, to five digits. (c) Explain this with \cref{prop:steepest:attained}(i), by finding the direction of each final cycle and checking that it maximizes $u^\top Hu$ over the unit ball.\label{ex:steepest:plane}

\item \simulation (The best fixed step depends on what is measured.) Run sign descent with a fixed step $\eta$ on the quadratic of \cref{fig:steepest:descent}, from its starting point (script \texttt{01-steepest-descent.py}). For each of $T=10^2$, $10^3$ and $10^4$, find by a grid search over $\eta\in[10^{-4},1]$ (a) the $\eta$ that minimizes $\frac1T\sum_{t<T}\norm{g_t}_1$, and (b) the $\eta$ that minimizes $\Loss(\theta_T)$. You should find (a) within the grid spacing of $\sqrt{2\Delta_0/(T\tr H)}$, with minimum $\sqrt{2\Delta_0\tr H/T}$ to three digits, and (b) within about ten percent of $\max_i\abs{\theta_{0,i}}/T$. Explain (a) with \cref{prop:steepest:attained}(ii), and (b) by the time a step needs to cover the distance $\norm{\theta_0}_\infty$.\label{ex:steepest:beststep}

\item \research (Choosing a norm from data.) \Cref{ex:steepest:quadnorm}(c) shows that among quadratic norms the one-step guarantee singles out the Hessian. (a) For weighted max norms $\max_ia_i\abs{u_i}$ on a diagonal quadratic, find the weights that maximize the guarantee $\norm{g}_*^2/(2L_{\norm\cdot})$ at a given gradient, and show that the resulting step is Newton's. (b) For $h_i=i^{-a}$ and Gaussian $\theta_0$, map numerically the exponents and target losses for which normalized $\ell_\infty$ descent at its best fixed step beats normalized $\ell_2$ descent, and test whether \eqref{eq:steepest:whichnorm} at $\theta_0$ predicts the winner. Open: a rule that chooses a norm online, from gradients alone, with a nonconvex guarantee that improves on the best fixed norm of a family in hindsight; AdaGrad's guarantee (\cref{ch:adaptive}) is of this kind for diagonal quadratic norms and convex losses.\label{ex:steepest:choose}
\end{exercises}

\notesandrefs

Steepest descent goes back to \textcite{Cauchy1847}, and its stochastic form to \textcite{RobbinsMonro1951}. \Textcite[Section~9.4]{BoydVandenberghe2004} treat it in a general norm, define normalized and unnormalized steepest descent directions, and work out the quadratic norms and the $\ell_1$ norm, for which it is coordinate descent; \textcite{NutiniEtAl2015} showed that the speed of the Gauss--Southwell rule is governed by strong convexity in the $\ell_1$ norm. The descent inequality and the convex rate are standard \parencite{Nesterov2018}, and \textcite{KelnerEtAl2014} used steepest descent in the $\ell_\infty$ norm for maximum flow. Mirror descent is due to \textcite{NemirovskyYudin1983}, in its modern form to \textcite{BeckTeboulle2003}; the linear minimization oracle is the subproblem of the method of \textcite{FrankWolfe1956}.

In deep learning, spectral descent was proposed by \textcite{CarlsonEtAl2015}, sign descent entered through gradient compression \parencite{BernsteinEtAl2018} and the analysis of Adam as a sign method \parencite{BallesHennig2018}, and layerwise normalization through large-batch training \parencite{YouGitmanGinsburg2017,YouEtAl2020}. \Textcite{BernsteinNewhouse2024} read Adam, Shampoo and Prodigy without accumulators as steepest descent in particular norms, \textcite{LargeEtAl2024} built a norm for whole networks from the norms of their modules, \textcite{BernsteinNewhouse2025} made the duality map of that norm the central object of a theory of optimizers, and \textcite{Newhouse2025} gives a connected account. The bound \eqref{eq:steepest:normalized} for exact gradients is elementary, and we do not know who stated it first, nor have we seen the exactness of \cref{prop:steepest:attained} stated. The hardness of $\norm{H}_{\infty\to1}$ is due to \textcite{Rohn2000}, and \textcite{BallesPedregosaLeRoux2020} relate the smoothness of a quadratic in the max norm to the concentration of its Hessian on the diagonal.

\paragraph{Translating the literature.} \Textcite{BoydVandenberghe2004} call our $\lmo(g)$ a normalized steepest descent direction $\Delta x_{\mathrm{nsd}}$, and our $-g^\sharp$ the unnormalized steepest descent direction $\Delta x_{\mathrm{sd}}=\norm{\nabla f(x)}_*\Delta x_{\mathrm{nsd}}$. \Textcite{KelnerEtAl2014} write $x^\#$ for our $x^\sharp$. \Textcite{BernsteinNewhouse2024} write $\norm{\cdot}^\dagger$ for our dual norm and $\lambda$ for the coefficient of the quadratic penalty, our $L$ (or $1/\eta$), so their steepest descent step $-(\norm{g}^\dagger/\lambda)\argmax_{\norm{t}=1}g^\top t$ is our $-g^\sharp/L$; in \textcite{BernsteinNewhouse2025} the map $g\mapsto\argmax_{\norm{t}=1}\ip{g,t}$, our $-\lmo(g)$, is the duality map that produces the dualized gradient. \Textcite{PethickEtAl2025} write $\lmo(d)\in\argmin_{x\in\cD}\ip{d,x}$ for a norm ball $\cD$ of radius $\rho$, which is $\rho$ times our $\lmo(d)$, so their update $x^{k+1}=x^k+\gamma_k\lmo(d^k)$ is our normalized step with $\eta=\gamma_k\rho$. Our $L_{\ell_\infty}=\norm{H}_{\infty\to1}$ is also written $\norm{H}_{\infty,1}$.

\printchapterbib
\end{refsection}

% Chapter 2 of Optimizers in the Wild: sign descent and the l-infinity geometry of Adam, signSGD and Lion.

\casechapter{Sign Descent}{The $\ell_\infty$ geometry of Adam, signSGD and Lion}{signSGD, Lion, Adam without memory}{ch:sign}
\begin{refsection}[chapters/02-sign.bib]

\begin{chaptersummary}
Adam without its moving averages, signSGD and Lion move every coordinate by the same amount: they take the step $-\eta\,\sign(\cdot)$, a step of fixed length in the max norm. We prove that this step gains $\eta\norm{\nabla\Loss}_1$ against a curvature cost $\frac12\eta^2\norm{\vec L}_1$, that the trade is exact on separable quadratics, and that a rotation of the coordinates can raise the cost by the factor $d\lmax/\tr H$, which gradient descent does not notice. With minibatches the sign follows the median of the gradient rather than its mean; the resulting bias $\norm{\vec\sigma}_1/\sqrt B$ is sharp, a rare class realizes it, and symmetric noise avoids it. Adam's step is at most $C(\beta_1,\beta_2)$ times the learning rate in each coordinate, $7.27$ at the defaults; the idealization drops the memory that makes Adam work, and it cannot say why the coordinates of real networks are good ones.
\end{chaptersummary}

\section{In the wild: one step length for every coordinate}\label{sec:sign:wild}

Adam divides a moving average of gradients by the square root of a moving average of squared gradients \parencite{KingmaBa2015}. The division removes the scale of each coordinate: multiply the gradient of one parameter by a thousand, and Adam's step in that parameter does not change. Its learning rate is therefore not a multiplier of the gradient but a length, roughly the distance each parameter moves in one step. Two other optimizers make the same choice exactly. \term{SignSGD}\index{signSGD} updates with the sign of a stochastic gradient; in data-parallel training each worker then sends one bit per coordinate, and a server that takes the majority vote of the signs sends one bit back \parencite{BernsteinEtAl2018,BernsteinEtAl2019}. The idea goes back to Rprop, which used only the signs of a full-batch gradient,\sidenote{Rprop keeps a step size for each weight, grows it while the weight's partial derivative keeps its sign and shrinks it when the sign flips; the size of the gradient never enters the update \parencite{RiedmillerBraun1993}. RMSProp was introduced in a lecture course as a minibatch version of Rprop \parencite{TielemanHinton2012}.} and to 1-bit SGD, which quantized each gradient coordinate to one bit for distributed training and carried the quantization error over to the next step \parencite{SeideEtAl2014}. Lion, found by a symbolic search over update rules, updates with the sign of an interpolation between the gradient and a momentum, and keeps one number of state per parameter where Adam keeps two \parencite{ChenEtAl2023}.

\begin{practice}[Adam, \textcite{KingmaBa2015}; GPT-3, \textcite{BrownEtAl2020}; Lion, \textcite{ChenEtAl2023}; signSGD, \textcite{BernsteinEtAl2019}]
Adam with $\beta_1=0.9$, $\beta_2=0.999$, $\epsilon=10^{-8}$, the published defaults, and $\beta_2=0.95$ in large language models since GPT-3. Lion with $\beta_1=0.9$, $\beta_2=0.99$, a learning rate $3$ to $10$ times smaller than AdamW's and a weight decay $3$ to $10$ times larger, so that their product, the strength of the decay, stays about the same. In distributed training, send $\sign(g)$, one bit per coordinate, and aggregate by majority vote.
\end{practice}

Three observations tie these optimizers together.

\begin{observation}[Sign descent with momentum behaves like Adam; \textcite{KunstnerEtAl2023}, \textcite{ZhaoEtAl2025}]\label[observation]{obs:sign:signum}
On transformer language models the gap between Adam and stochastic gradient descent does not shrink as the batch grows to the full training set: Adam improves with the batch size and gradient descent much less, and with large batches Adam behaves like sign descent with momentum. In a comparison of optimizers for autoregressive language models across sizes and hyperparameters, signed momentum (Signum) recovered both the performance of Adam and its stability across hyperparameters.
\end{observation}

\Textcite{KunstnerEtAl2023} concluded that heavy-tailed noise, proposed earlier as the reason for the gap, is not its main cause, and \textcite{ZhaoEtAl2025} found Adam, Adafactor, Lion and Sophia comparable, with only stochastic gradient descent behind. Both studies are empirical and at moderate scale, and neither says why the sign helps.

\begin{observation}[Adam is a sign with a variance-dependent length; \textcite{BallesHennig2018}]\label[observation]{obs:sign:dissect}
Adam's update of each coordinate can be read as the sign of the momentum times a factor that decreases with the relative variance of that coordinate's gradient. Studied separately, the sign is the aspect responsible for adverse effects on generalization that had been reported for adaptive methods.
\end{observation}

\begin{observation}[Rotating the parameters slows Adam; \textcite{XieMohamadiLi2025}]\label[observation]{obs:sign:rotation}
Training GPT-2 with Adam after a rotation of the parameter space is markedly slower than in the original coordinates, while stochastic gradient descent, whose trajectory rotates with the coordinates, is unaffected.
\end{observation}

The three point at one object, the step $-\eta\,\sign(g)$, and one question: what does it buy, compared with $-\eta g$, and what does it cost? The answer depends on the norm in which steps are measured, on the coordinates, on the shape of the noise, and on memory.

\section{The idealization: steepest descent in the max norm}\label{sec:sign:ideal}

\Cref{ch:steepest} turned a gradient into a step with a norm: the normalized steepest-descent step is $\eta\lmo(g)$ with $\lmo(g)\in\argmin_{\norm u\le1}\ip{g,u}$. For the max norm the dual norm is $\norm{g}_1$ (\cref{lem:toolkit:dualnorm}), and $\ip{g,u}$ is minimized over the cube $\{\norm{u}_\infty\le1\}$ by $u_i=-\sign(g_i)$ where $g_i\ne0$ and by every $u_i\in[-1,1]$ where $g_i=0$; the convention $\sign0=0$ picks the shortest minimizer. Normalized steepest descent in $\ell_\infty$ is sign descent.

\begin{definition}[Sign descent and signSGD]\label[definition]{def:sign:signdescent}
Given step sizes $\eta_t>0$, \term{sign descent}\index{sign descent} is $\theta_{t+1}=\theta_t-\eta_t\sign(\nabla\Loss(\theta_t))$, and signSGD is $\theta_{t+1}=\theta_t-\eta_t\sign(g_t)$ for a stochastic gradient $g_t$, with the sign taken coordinatewise. When no coordinate of the gradient vanishes, the step has length $\eta_t$ in $\ell_\infty$ and in the RMS norm, and $\eta_t\sqrt d$ in $\ell_2$.
\end{definition}

The learning rate of the practice sections and the step size of the theory are the same number $\eta$; for a sign method it is the distance that every coordinate moves. The curvature bound that suits the max norm is the one with which \textcite{BernsteinEtAl2018} posed the problem.

\begin{assumption}[Coordinatewise smoothness]\label[assumption]{ass:sign:smooth}
The loss $\Loss\colon\R^d\to\R$ is differentiable and bounded below by $\Loss^\star$, and there are $L_1,\dots,L_d\ge0$ such that for all $\theta,u\in\R^d$
\[
  \Loss(\theta+u)\le\Loss(\theta)+\ip{\nabla\Loss(\theta),u}+\frac12\sum_{i=1}^dL_iu_i^2 .
\]
We write $\vec L=(L_1,\dots,L_d)$, $\norm{\vec L}_1=\sum_iL_i$ and $L_{\max}=\max_iL_i$.
\end{assumption}

For a twice-differentiable loss, \term{coordinatewise smoothness}\index{coordinatewise smoothness} says that $\nabla^2\Loss(\theta)\preceq\diag(\vec L)$ at every $\theta$. Since $\frac12\sum_iL_iu_i^2$ is at most $\frac12\norm{\vec L}_1\norm{u}_\infty^2$ and at most $\frac12L_{\max}\norm{u}_2^2$, it gives the quadratic upper bounds of \cref{ch:steepest} (\cref{lem:toolkit:descent}) in two geometries at once: the constant $\norm{\vec L}_1$ in $\ell_\infty$ and $L_{\max}$ in $\ell_2$.

\begin{remark}[Coordinatewise and $\ell_\infty$ smoothness]\label[remark]{rem:sign:sdp}
For a quadratic with Hessian $H\succeq0$ the admissible $\vec L$ are those with $\diag(\vec L)\succeq H$, and each has $\norm{\vec L}_1=s^\top\diag(\vec L)s\ge s^\top Hs$ for every sign vector $s$. Hence $\norm{\vec L}_1\ge\norm{H}_{\infty\to1}=\max_{s\in\{\pm1\}^d}s^\top Hs$, the exact $\ell_\infty$ constant of \cref{ch:steepest}, with equality for diagonal $H$ and $\vec L=\diag(H)$; conversely the best $\norm{\vec L}_1$ is at most $\frac\pi2\norm{H}_{\infty\to1}$ (\cref{ex:sign:sdp}, after \textcite{Nesterov1998}).
\end{remark}

\begin{assumption}[Minibatch noise]\label[assumption]{ass:sign:noise}
The stochastic gradient is $g_t=\nabla\Loss(\theta_t)+\zeta_t$ where, conditionally on $\theta_0,g_0,\dots,g_{t-1}$, each coordinate $\zeta_{t,i}$ has mean zero and variance at most $\sigma_i^2/B$. We write $\vec\sigma=(\sigma_1,\dots,\sigma_d)$.
\end{assumption}

A batch of $B$ independent examples whose per-example gradients have coordinate variances at most $\sigma_i^2$ satisfies the assumption. Adam enters with the conventions of the book: $m_{-1}=v_{-1}=0$, $m_t=\beta_1m_{t-1}+(1-\beta_1)g_t$, $v_t=\beta_2v_{t-1}+(1-\beta_2)g_t^2$ with squares taken coordinatewise, $\hat m_t=m_t/(1-\beta_1^{t+1})$, $\hat v_t=v_t/(1-\beta_2^{t+1})$, and $\theta_{t+1}=\theta_t-\eta\,\hat m_t/(\sqrt{\hat v_t}+\epsilon)$.

\begin{dictionary}
the step $-\eta\sign(g)$, of length $\eta$ in $\ell_\infty$ & an update that moves every parameter by the learning rate\\
the dual norm $\norm{\nabla\Loss}_1$ & the progress per unit of learning rate, summed over parameters\\
$\vec L$ and $\norm{\vec L}_1$ & the curvature that each parameter sees along its own axis, and its sum\\
$\norm{g}_1^2/\norm{g}_2^2$, between $1$ and $d$ & the number of parameters that receive a gradient of comparable size\\
$\tr H/\lmax$ & the effective number of stiff directions\\
an orthogonal matrix $Q$ & a change of basis of parameter space, such as mixing the weights of a layer\\
$\sigma_i/\sqrt B$; the median of $g_{t,i}$ & the minibatch noise of a parameter; what a sign update follows\\
$C(\beta_1,\beta_2)$ & Adam's largest step of a parameter, in units of the learning rate
\end{dictionary}

\begin{fieldguide}{signSGD and Signum}{\textcite{BernsteinEtAl2018}}
\fgrow{update}{$\theta_{t+1}=\theta_t-\eta\,\sign(g_t)$; Signum: $\theta_{t+1}=\theta_t-\eta\,\sign(m_t)$, $m_t=\beta m_{t-1}+(1-\beta)g_t$}
\fgrow{norm}{$\ell_\infty$; $\lmo(g)=-\sign(g)$}
\fgrow{preconditioner}{none (equivalently $\diag\abs{g_t}^{-1}$, rebuilt at every step)}
\fgrow{averages}{none (signSGD); EMA of $g$ (Signum)}
\fgrow{state}{none (signSGD); one number per parameter (Signum)}
\fgrow{variants}{majority vote over $M$ workers, $\theta_{t+1}=\theta_t-\eta\,\sign\bigl(\sum_m\sign(g^{(m)}_t)\bigr)$ \parencite{BernsteinEtAl2019}; error feedback \parencite{SeideEtAl2014,KarimireddyEtAl2019}}
\fgrow{caution}{without momentum the step follows the median of the minibatch gradient, not its mean (\cref{prop:sign:rareclass})}
\end{fieldguide}

\begin{fieldguide}{Lion}{\textcite{ChenEtAl2023}}
\fgrow{update}{$c_t=\beta_1m_{t-1}+(1-\beta_1)g_t$, \ $m_t=\beta_2m_{t-1}+(1-\beta_2)g_t$,\\ $\theta_{t+1}=\theta_t-\eta\,\bigl(\sign(c_t)+\lambda\theta_t\bigr)$}
\fgrow{norm}{$\ell_\infty$}
\fgrow{preconditioner}{none}
\fgrow{averages}{one EMA of $g$ (coefficient $\beta_2$), read out with weight $1-\beta_1$ on the current gradient}
\fgrow{state}{one number per parameter}
\fgrow{in the wild}{$\beta_1=0.9$, $\beta_2=0.99$, a learning rate several times smaller than AdamW's (\cref{ex:sign:lion})}
\fgrow{variants}{$\beta_1=\beta_2$ is Signum; with weight decay it solves a box-constrained problem \parencite{ChenLiuLiangLiu2024}}
\end{fieldguide}

\begin{fieldguide}{Adam}{\textcite{KingmaBa2015}}
\fgrow{update}{$\theta_{t+1}=\theta_t-\eta\,\hat m_t/(\sqrt{\hat v_t}+\epsilon)$; full entry in \cref{ch:adaptive}}
\fgrow{norm}{$\ell_\infty$ if $\beta_1=\beta_2=\epsilon=0$ (\cref{prop:sign:adam0}); otherwise steps of at most $C(\beta_1,\beta_2)\,\eta$ per coordinate (\cref{prop:sign:adambound})}
\fgrow{preconditioner}{diagonal, $\diag(\sqrt{\hat v_t})+\epsilon$}
\fgrow{averages}{EMA of $g$ ($\beta_1$) and of $g^2$ ($\beta_2$)}
\fgrow{state}{two numbers per parameter}
\fgrow{in the wild}{$\beta_1=0.9$, $\beta_2\in\{0.95,0.999\}$, $\epsilon=10^{-8}$}
\end{fieldguide}

\paragraph{What was thrown away.} The idealization keeps the step and discards the rest: the moving averages of Adam, Lion and Signum (they return only through the size of Adam's step and a heuristic for its mean, and in full in \cref{ch:momentum,ch:adaptive}), Adam's $\epsilon$, the variation of curvature along a trajectory, which a global bound in the coordinates replaces, and every feature of the noise except its variance. These are the right first simplifications because the sign is the one ingredient that all of these optimizers share and that gradient descent lacks.

\paragraph{What one had to notice.} Without memory, dividing by $\sqrt v$ is dividing by $\abs g$, so Adam's update is a vector of signs, a corner of the unit ball of the max norm. Once this is seen, Adam, Rprop, signSGD and Lion are one method in one geometry, and the questions are those of \cref{ch:steepest}: which dual norm measures the gradient, and which smoothness constant pays for the step. \Textcite{BernsteinNewhouse2024} put this observation at the center of a classification of optimizers by norms.

\section{The theory: the max-norm step}\label{sec:sign:theory}

\begin{proposition}[Adam without memory is sign descent]\label[proposition]{prop:sign:adam0}
Let $\epsilon\ge0$ and write Adam's update as $\theta_{t+1}=\theta_t-\eta u_t$, with $u_{t,i}=\hat m_{t,i}/(\sqrt{\hat v_{t,i}}+\epsilon)$ and $0/0=0$. (i) If $\beta_1=\beta_2=0$, then $u_{t,i}=\frac{\abs{g_{t,i}}}{\abs{g_{t,i}}+\epsilon}\sign(g_{t,i})$ for every $t$, with or without bias correction, and $u_t=\sign(g_t)$ if $\epsilon=0$; the same holds for RMSProp, $\theta_{t+1}=\theta_t-\eta g_t/(\sqrt{v_t}+\epsilon)$, with coefficient $0$. (ii) With bias correction, $u_0$ is given by the same formula for all $\beta_1,\beta_2\in[0,1)$: if $\epsilon=0$, Adam's first step is a sign step whatever the size of the first gradient.
\end{proposition}

\begin{proof}
(i) With $\beta_1=\beta_2=0$, $m_t=g_t$ and $v_t=g_t^2$, and the correction factors $1-0^{t+1}$ equal $1$, so $u_{t,i}=g_{t,i}/(\abs{g_{t,i}}+\epsilon)$; RMSProp with coefficient $0$ has the same update. (ii) Here $m_0=(1-\beta_1)g_0$ and $v_0=(1-\beta_2)g_0^2$, so $\hat m_0=g_0$ and $\hat v_0=g_0^2$.
\end{proof}

The factor $\abs g/(\abs g+\epsilon)$ makes memoryless Adam sign descent where the gradient is much larger than $\epsilon$ and gradient descent with step size $\eta/\epsilon$ where it is much smaller (\cref{ex:sign:epsilon}). With $\epsilon=0$, memoryless Adam, memoryless RMSProp and signSGD are one algorithm, and its guarantee is the descent inequality of \cref{ch:steepest} in the max norm.

\begin{theorem}[Sign descent in the max norm]\label{thm:sign:descent}
Let \cref{ass:sign:smooth} hold and let $\theta_{t+1}=\theta_t-\eta_t\sign(\nabla\Loss(\theta_t))$ with $\eta_t>0$. Then for every $t$
\begin{equation}\label{eq:sign:onestep}
  \Loss(\theta_{t+1})\le\Loss(\theta_t)-\eta_t\norm{\nabla\Loss(\theta_t)}_1+\tfrac12\eta_t^2\norm{\vec L}_1,
\end{equation}
where $\norm{\vec L}_1$ may be replaced by the sum of the $L_i$ with $\partial_i\Loss(\theta_t)\ne0$; and for every $T\ge1$
\begin{equation}\label{eq:sign:rate}
  \frac{\sum_{t<T}\eta_t\norm{\nabla\Loss(\theta_t)}_1}{\sum_{t<T}\eta_t}\le\frac{\Delta_0+\frac12\norm{\vec L}_1\sum_{t<T}\eta_t^2}{\sum_{t<T}\eta_t}.
\end{equation}
For a constant step size $\eta$ the right side is $\Delta_0/(\eta T)+\frac12\eta\norm{\vec L}_1$; if $\norm{\vec L}_1>0$ and $\eta=\sqrt{2\Delta_0/(\norm{\vec L}_1T)}$, it is $\sqrt{2\Delta_0\norm{\vec L}_1/T}$, which therefore bounds $\min_{t<T}\norm{\nabla\Loss(\theta_t)}_1$.
\end{theorem}

\begin{proof}
The plan is to apply \cref{ass:sign:smooth} to the step $u=-\eta_t\sign(\nabla\Loss(\theta_t))$ and to sum. Fix $t$, and write $g=\nabla\Loss(\theta_t)$ and $s=\sign(g)$. The linear term of the assumption is $\ip{g,-\eta_ts}=-\eta_t\sum_i\abs{g_i}=-\eta_t\norm{g}_1$. The quadratic term is $\frac12\eta_t^2\sum_iL_is_i^2$, and $s_i^2$ is $1$ where $g_i\ne0$ and $0$ elsewhere. This proves \eqref{eq:sign:onestep}. Rearranged, \eqref{eq:sign:onestep} bounds $\eta_t\norm{\nabla\Loss(\theta_t)}_1$ by $\Loss(\theta_t)-\Loss(\theta_{t+1})+\frac12\eta_t^2\norm{\vec L}_1$; summed over $t<T$, the loss differences telescope to $\Loss(\theta_0)-\Loss(\theta_T)\le\Delta_0$, and dividing by $\sum_t\eta_t$ gives \eqref{eq:sign:rate}. The stated constant $\eta$ minimizes $\Delta_0/(\eta T)+\frac12\eta\norm{\vec L}_1$, and a minimum is at most an average. If $\norm{\vec L}_1=0$, the assumption puts $\Loss$ below each of its tangent planes, so $\Loss$ is concave; being bounded below on $\R^d$, it is constant, and there is nothing to prove.
\end{proof}

The comparison with gradient descent comes from reading the same assumption in three geometries.

\begin{corollary}[Three guaranteed decreases]\label[corollary]{cor:sign:compare}
Let \cref{ass:sign:smooth} hold with every $L_i>0$, and let $g=\nabla\Loss(\theta)\ne0$. The gradient step $-g/L_{\max}$, the sign step $-(\norm{g}_1/\norm{\vec L}_1)\sign(g)$ and the diagonally preconditioned step $-\diag(\vec L)^{-1}g$ decrease $\Loss$ by at least $\norm{g}_2^2/(2L_{\max})$, $\norm{g}_1^2/(2\norm{\vec L}_1)$ and $\sum_ig_i^2/(2L_i)$ respectively. The third is at least the second, with equality if and only if $\abs{g_i}/L_i$ does not depend on $i$; the second is at least the first if and only if
\begin{equation}\label{eq:sign:density}
  \frac{\norm{g}_1^2}{\norm{g}_2^2}\ \ge\ \frac{\norm{\vec L}_1}{L_{\max}} .
\end{equation}
\end{corollary}

\begin{proof}
The plan is to evaluate \cref{ass:sign:smooth} at the three steps and to compare the results. The first and third bounds are the assumption at $u=-g/L_{\max}$, with $\sum_iL_iu_i^2\le L_{\max}\norm{u}_2^2$, and at $u_i=-g_i/L_i$. The second is \eqref{eq:sign:onestep} at $\eta=\norm{g}_1/\norm{\vec L}_1$, which maximizes $\eta\norm{g}_1-\frac12\eta^2\norm{\vec L}_1$. Writing $\abs{g_i}=(\abs{g_i}/\sqrt{L_i})\sqrt{L_i}$, the Cauchy--Schwarz inequality gives $\norm{g}_1^2\le\norm{\vec L}_1\sum_ig_i^2/L_i$, with equality exactly when $\abs{g_i}/\sqrt{L_i}$ is proportional to $\sqrt{L_i}$. The last claim is a rearrangement.
\end{proof}

Both sides of \eqref{eq:sign:density} lie between $1$ and $d$. The left side counts the coordinates over which the gradient is spread; the right side is an \term{effective rank}\index{effective rank} of the curvature in the coordinates of the parameters, which for a quadratic with diagonal Hessian and $\vec L=\diag(H)$ is $\tr H/\lmax$. This is the reading of the comparison of normalized steps in \cref{prop:steepest:whichnorm}, which has the exact constant $\norm{H}_{\infty\to1}$ in place of $\norm{\vec L}_1$: sign descent's guarantee is the better one when the gradient is spread over more coordinates than the curvature has stiff directions. A network has millions of parameters, most of which receive gradient, and the Hessians of trained networks have a few large outlier eigenvalues above a bulk of small ones \parencite{GhorbaniKrishnanXiao2019}; were such a Hessian diagonal in the parameters' coordinates, the advantage could approach $d$ over the effective rank. Whether it is close to diagonal is the subject of \cref{prop:sign:rotation}. The equality case says when the sign is as good as the best diagonal preconditioner: for a quadratic $\abs{g_i}/L_i$ is the length of the Newton step in coordinate $i$, so this happens when all coordinates are equally far from their targets.

Is the cost $\frac12\eta^2\norm{\vec L}_1$ an artifact of the proof? On a quadratic it is exact.

\begin{proposition}[The floor of sign descent on a quadratic]\label[proposition]{prop:sign:floor}
Let $\Loss(\theta)=\frac12(\theta-\theta^\star)^\top H(\theta-\theta^\star)$ with $H\succeq0$, run sign descent with a constant step size $\eta$, and write $g_t=\nabla\Loss(\theta_t)$ and $s_t=\sign(g_t)$.
\begin{enumerate}
\item For every $T\ge1$,
\begin{equation}\label{eq:sign:identity}
  \frac1T\sum_{t<T}\Bigl(\norm{g_t}_1-\frac\eta2\,s_t^\top Hs_t\Bigr)=\frac{\Delta_0-\Delta_T}{\eta T},
\end{equation}
and $\Delta_t$ is bounded; so $\frac1T\sum_{t<T}\norm{g_t}_1$ and $\frac1T\sum_{t<T}\frac12\eta s_t^\top Hs_t$ have the same limit points, all at most $\frac12\eta\norm{H}_{\infty\to1}$.
\item If $H=\diag(h_1,\dots,h_d)$ with all $h_i>0$, let $r_i\in[0,\eta)$ be the remainder of $\abs{\theta_{0,i}-\theta^\star_i}$ modulo $\eta$. From step $\lfloor\abs{\theta_{0,i}-\theta^\star_i}/\eta\rfloor$ on, the distance of coordinate $i$ from $\theta^\star_i$ alternates between $r_i$ and $\eta-r_i$ if $r_i>0$, and is $0$ if $r_i=0$. Hence $\frac1T\sum_{t<T}\norm{g_t}_1\to\frac12\eta\sum_{i:\,r_i>0}h_i$, which is $\frac12\eta\tr H$ unless some $\abs{\theta_{0,i}-\theta^\star_i}$ is a multiple of $\eta$.
\end{enumerate}
\end{proposition}

\begin{proof}
The plan is to use that the descent inequality is an identity for a quadratic, and to solve the separable case one coordinate at a time. (i) The identity is $\Loss(\theta_t-\eta s_t)=\Loss(\theta_t)-\eta\norm{g_t}_1+\frac12\eta^2s_t^\top Hs_t$; summing over $t<T$ and dividing by $\eta T$ gives \eqref{eq:sign:identity}. For boundedness, let $H\ne0$, let $\lambda_+$ be its smallest positive eigenvalue and $L_\infty=\norm{H}_{\infty\to1}$; then $s^\top Hs\le L_\infty$ for every $s\in\{-1,0,1\}^d$, since the convex function $x\mapsto x^\top Hx$ is maximized over the cube at a vertex. As $g_t$ lies in the range of $H$, $\Delta_t=\frac12g_t^\top H^+g_t\le\norm{g_t}_1^2/(2\lambda_+)$. If $\norm{g_t}_1\ge\frac12\eta L_\infty$ the identity gives $\Delta_{t+1}\le\Delta_t$, and otherwise $\Delta_{t+1}\le\Delta_t+\frac12\eta^2L_\infty<\eta^2L_\infty^2/(8\lambda_+)+\frac12\eta^2L_\infty$; by induction $\Delta_t$ never exceeds the larger of $\Delta_0$ and this number.
(ii) Coordinate $i$ evolves on its own: $x_t=\theta_{t,i}-\theta^\star_i$ satisfies $x_{t+1}=x_t-\eta\sign(x_t)$, so $\abs{x_{t+1}}=\bigl\lvert\abs{x_t}-\eta\bigr\rvert$ if $x_t\ne0$, and $x_{t+1}=0$ if $x_t=0$. The distance falls by $\eta$ per step for $\lfloor\abs{x_0}/\eta\rfloor$ steps and reaches $r_i$; if $r_i=0$ it stays there, and if $r_i>0$ the next distances are $\eta-r_i$ and $\abs{\eta-r_i-\eta}=r_i$. Over each period $h_i\abs{x_t}$ averages $\frac12\eta h_i$.
\end{proof}

For a separable quadratic the best constants are $\vec L=\diag(H)$, so the long-run average of $\norm{\nabla\Loss}_1$ equals the floor $\frac12\eta\norm{\vec L}_1$ of \cref{thm:sign:descent}, which cannot be improved. At a constant step size sign descent ends on a \term{two-cycle}\index{two-cycle}, the many-dimensional version of the prologue's example (\cref{ch:prologue}); only a decreasing schedule removes the floor (\cref{ch:schedules}).

Gradient descent does not depend on the orthonormal basis in which the parameters are written; sign descent does.

\begin{proposition}[Sign descent is not rotation invariant]\label[proposition]{prop:sign:rotation}
Let $H\succeq0$ have largest eigenvalue $\lmax$ with unit eigenvector $v_1$, and for an orthogonal $Q$ let $\Loss_Q(\phi)=\Loss(Q\phi)$, whose Hessian $Q^\top HQ$ has largest eigenvalue $\lmax$. Then (i) $\tr H\le\norm{Q^\top HQ}_{\infty\to1}\le d\lmax$, with the lower bound attained when the columns of $Q$ are eigenvectors of $H$ and the upper bound when $Q^\top v_1=\mathbf1/\sqrt d$; (ii) $\norm{Q^\top HQ}_{\infty\to1}\ge\lmax\norm{Q^\top v_1}_1^2$; (iii) if $Q$ is uniformly distributed on the orthogonal group, $\E\norm{Q^\top v_1}_1^2=1+\frac2\pi(d-1)$. So rotations change the $\ell_\infty$ smoothness constant by factors up to $d\lmax/\tr H$, and a random rotation multiplies it on average by at least $(1+\frac2\pi(d-1))\lmax/\tr H$.
\end{proposition}

\begin{proof}
The plan is to bound $s^\top As$, for $A=Q^\top HQ$, by averaging over random signs, by the top eigenvalue, and at the signs of the top eigenvector; recall $\norm{A}_{\infty\to1}=\max_{s\in\{\pm1\}^d}s^\top As$. (i) For independent uniform signs, $\E[s^\top As]=\sum_{i,j}A_{ij}\E[s_is_j]=\tr A=\tr H$, so some sign vector reaches $\tr H$; and $s^\top As\le\lmax\norm{s}_2^2=d\lmax$. If the columns of $Q$ are eigenvectors, $A$ is diagonal and nonnegative and $s^\top As=\tr A$ for all $s$. Some orthogonal $Q$ maps $\mathbf1/\sqrt d$ to $v_1$, and for it (ii) gives $\norm{A}_{\infty\to1}\ge d\lmax$. (ii) Let $s=\sign(Q^\top v_1)$, zeros replaced by $1$, and let $(\lambda_j,v_j)$ be the eigenpairs of $H$, all $\lambda_j\ge0$. Then $s^\top As=\sum_j\lambda_j\ip{v_j,Qs}^2\ge\lmax\ip{Q^\top v_1,s}^2=\lmax\norm{Q^\top v_1}_1^2$. (iii) The vector $u=Q^\top v_1$ is uniform on the sphere, so it has the law of $z/\norm z_2$ for $z\sim\cN(0,I_d)$ and is independent of $\norm z_2$. For $i\ne j$, $\frac2\pi=\E\abs{z_iz_j}=\E\norm z_2^2\,\E\abs{u_iu_j}=d\,\E\abs{u_iu_j}$; hence $\E\norm u_1^2=\sum_i\E u_i^2+\sum_{i\ne j}\E\abs{u_iu_j}=1+d(d-1)\frac{2}{\pi d}$.
\end{proof}

A random rotation also spreads the gradient, $\norm{Q^\top g}_1^2\approx\frac2\pi d\norm{g}_2^2$, so after it both sides of the criterion \eqref{eq:sign:density}, in its exact form, are about $\frac2\pi d$ or more on the right: the advantage is gone, not reversed. The guarantee of sign descent is a property of the coordinate system, and \cref{fig:sign:rotation} shows the dynamics following it.

\section{The theory: noise and memory}\label{sec:sign:noise}

A stochastic gradient has the wrong sign in some coordinates, and a sign method moves those coordinates the wrong way by a full step. The cost is a question about one random variable.

\begin{lemma}[Sign errors]\label[lemma]{lem:sign:signerror}
Let $X$ be a random variable with mean $\mu\ne0$ and variance at most $s^2$, where $s>0$. (i) $\Pr[\sign X\ne\sign\mu]\le s^2/(s^2+\mu^2)$. (ii) $\mu\,\E[\sign X]\ge\abs{\mu}-2\abs{\mu}s^2/(s^2+\mu^2)\ge\abs{\mu}-s$. (iii) If $X-\mu$ is symmetric about $0$ and unimodal, then $\mu\,\E[\sign X]\ge\min\bigl\{\mu^2/(\sqrt3\,s),\ \frac23\abs{\mu}\bigr\}>0$.
\end{lemma}

\begin{proof}
The plan is Cantelli's inequality for (i) and (ii) and Gauss's inequality for (iii). Replacing $X$ by $-X$ if necessary, let $\mu>0$ and write $\zeta=X-\mu$. (i) The event is $\{\zeta\le-\mu\}$, which for every $c\ge0$ lies in $\{(c-\zeta)^2\ge(c+\mu)^2\}$; by Markov's inequality its probability is at most $(s^2+c^2)/(c+\mu)^2$, which is $s^2/(s^2+\mu^2)$ at $c=s^2/\mu$ (Cantelli's inequality\index{Cantelli's inequality}). (ii) Pointwise $\sign X\ge1-2\,\ind[X\le0]$; take expectations and use (i). The second inequality is $2\mu s\le s^2+\mu^2$. (iii) A unimodal law with mode $0$ has no atoms away from $0$, and by symmetry $\Pr[\zeta<-\mu]=\Pr[\zeta>\mu]$, so $\E[\sign X]=\Pr[\zeta>-\mu]-\Pr[\zeta>\mu]=\Pr[\abs\zeta<\mu]$. Gauss's inequality\index{Gauss's inequality} (\cref{lem:toolkit:gauss}), for a unimodal law with mode $0$ and second moment $\tau^2\le s^2$, bounds $\Pr[\abs\zeta\ge k]$ by $1-k/(\sqrt3\tau)$ if $k\le2\tau/\sqrt3$ and by $4\tau^2/(9k^2)$ otherwise. Both bounds increase with $\tau$, so they hold with $s$ for $\tau$, and at $k=\mu$ they give $\Pr[\abs\zeta<\mu]\ge\mu/(\sqrt3s)$ if $\mu\le2s/\sqrt3$ and $\Pr[\abs\zeta<\mu]\ge1-\frac{4s^2}{9\mu^2}\ge\frac23$ otherwise.
\end{proof}

By (ii), the expected progress $\mu\E\sign X$ falls short of the noiseless $\abs\mu$ by at most the noise level $s$. Summed over coordinates, this is the stochastic version of \cref{thm:sign:descent}.

\begin{theorem}[signSGD]\label{thm:sign:signsgd}
Let \cref{ass:sign:smooth,ass:sign:noise} hold, and run signSGD, $\theta_{t+1}=\theta_t-\eta_t\sign(g_t)$, with deterministic $\eta_t>0$. Then for every $T\ge1$
\[
  \frac{\sum_{t<T}\eta_t\,\E\norm{\nabla\Loss(\theta_t)}_1}{\sum_{t<T}\eta_t}\le\frac{\Delta_0+\frac12\norm{\vec L}_1\sum_{t<T}\eta_t^2}{\sum_{t<T}\eta_t}+\frac{\norm{\vec\sigma}_1}{\sqrt B}.
\]
For a constant step size the right side is $\Delta_0/(\eta T)+\frac12\eta\norm{\vec L}_1+\norm{\vec\sigma}_1/\sqrt B$. With $\eta=\sqrt{2\Delta_0/(\norm{\vec L}_1T)}$ and $B=T$, after $N=BT$ stochastic gradients, it is $\bigl(\sqrt{2\Delta_0\norm{\vec L}_1}+\norm{\vec\sigma}_1\bigr)N^{-1/4}$.
\end{theorem}

\begin{proof}
The plan is the proof of \cref{thm:sign:descent} in conditional expectation, with \cref{lem:sign:signerror}(ii) for the linear term. Let $\cF_t$ be generated by $\theta_0,g_0,\dots,g_{t-1}$, and write $\mu=\nabla\Loss(\theta_t)$. \Cref{ass:sign:smooth} at $u=-\eta_t\sign(g_t)$ gives
\[
  \E[\Loss(\theta_{t+1})\mid\cF_t]\le\Loss(\theta_t)-\eta_t\sum_{i=1}^d\mu_i\,\E[\sign g_{t,i}\mid\cF_t]+\tfrac12\eta_t^2\norm{\vec L}_1 .
\]
Conditionally on $\cF_t$, the coordinate $g_{t,i}$ has mean $\mu_i$ and variance at most $\sigma_i^2/B$, so \cref{lem:sign:signerror}(ii) gives $\mu_i\E[\sign g_{t,i}\mid\cF_t]\ge\abs{\mu_i}-\sigma_i/\sqrt B$ if $\mu_i\ne0$ and $\sigma_i>0$; the bound also holds if $\mu_i=0$, when the left side vanishes, and if $\sigma_i=0$, when $g_{t,i}=\mu_i$. Sum over $i$, take expectations, sum over $t<T$, use $\E\Loss(\theta_T)\ge\Loss^\star$ and divide by $\sum_t\eta_t$. With the stated $\eta$ and $B=T$ the first two terms equal $\sqrt{2\Delta_0\norm{\vec L}_1/T}$, the third is $\norm{\vec\sigma}_1/\sqrt T$, and $\sqrt T=N^{1/4}$.
\end{proof}

With a batch that grows with the horizon, signSGD reaches the rate $N^{-1/4}$ of stochastic gradient descent, with $\ell_1$ norms of the gradient and of the noise in place of $\ell_2$ norms; this is the comparison of \textcite{BernsteinEtAl2018}, whose noise term has the factor $2$ of Markov's inequality where Cantelli's gives $1$. At a fixed batch size, however, the term $\norm{\vec\sigma}_1/\sqrt B$ survives every step size. This is not a weakness of the proof: the sign of a minibatch gradient is right in most batches when the median of the gradient has the sign of its mean, and a rare event can separate the two.

\begin{proposition}[A rare class defeats signSGD]\label[proposition]{prop:sign:rareclass}
Let labels $y\in\{0,1\}$ be independent with $\Pr[y=1]=\pi\in(0,\frac12)$, let a model predict $\varsigma(b)=1/(1+e^{-b})$ with one bias $b\in\R$, and let $\Loss(b)=-\pi\log\varsigma(b)-(1-\pi)\log(1-\varsigma(b))$. SignSGD on batches of $B$ labels is $b_{t+1}=b_t-\eta\,\sign(\varsigma(b_t)-\bar y_t)$, where $\bar y_t$ is the fraction of ones in batch $t$. (i) $\Loss$ satisfies \cref{ass:sign:smooth,ass:sign:noise} with $L_1=\frac14$ and $\sigma_1^2=\pi(1-\pi)$, and is minimized where $\varsigma(b^\star)=\pi$. (ii) If $q=(1-\pi)^B>\frac12$, then for every $\eta>0$ and $b_0$, almost surely $\limsup_tb_t/t\le-\eta(2q-1)<0$; hence $\Loss(b_t)\to\infty$ and $\nabla\Loss(b_t)\to-\pi$.
\end{proposition}

\begin{proof}
(i) Here $\Loss'(b)=\varsigma(b)-\pi$ and $0<\Loss''(b)=\varsigma(b)(1-\varsigma(b))\le\frac14$, and the minibatch gradient $\varsigma(b)-\bar y_t$ averages $B$ independent terms $\varsigma(b)-y$ of mean $\varsigma(b)-\pi$ and variance $\pi(1-\pi)$.
(ii) The plan is to bound the drift of $b_t$ and to control its fluctuations by a martingale argument. With probability $q$ batch $t$ has no label $1$; then $\bar y_t=0<\varsigma(b_t)$ and the step is $-\eta$. A step $+\eta$ needs $\bar y_t>\varsigma(b_t)$, which has probability at most $1-q$. So $D_t=\E[b_{t+1}-b_t\mid b_0,\dots,b_t]\le-\eta(2q-1)$ whatever $b_t$ is. The differences $M_t=b_{t+1}-b_t-D_t$ are martingale differences with $\abs{M_t}\le2\eta$, and by the Azuma--Hoeffding inequality (\cref{thm:toolkit:azuma}) $\Pr\bigl[\abs{\sum_{s<t}M_s}\ge\eps t\bigr]\le2e^{-\eps^2t/(8\eta^2)}$ for each $\eps>0$. These are summable in $t$, so by the Borel--Cantelli lemma $\frac1t\sum_{s<t}M_s\to0$ almost surely, and $b_t/t=b_0/t+\frac1t\sum_{s<t}(D_s+M_s)$ has $\limsup$ at most $-\eta(2q-1)$. Then $\varsigma(b_t)\to0$, $\Loss(b_t)\ge-\pi\log\varsigma(b_t)\to\infty$ and $\Loss'(b_t)\to-\pi$.
\end{proof}

The batches without the rare label are the majority, and in them the gradient is positive, so the sign follows the majority. The mean gradient $\varsigma(b)-\pi$ is negative wherever the predicted probability is too small, but the median of the minibatch gradient is $\varsigma(b)>0$ for every $b$: signSGD descends the median, which here has no zero.\sidenote{In one dimension this is classical: with $g=\theta-\xi$, signSGD is stochastic subgradient descent on $\E\abs{\theta-\xi}$, whose minimizers are the medians of $\xi$.} Once $\varsigma(b_t)<1/B$, every batch with a label $1$ pushes $b$ up, and the drift is exactly $-\eta(2q-1)$ (\cref{fig:sign:noise}). The condition $(1-\pi)^B>\frac12$ holds for every class rarer than $1-2^{-1/B}\approx\ln2/B$, and a softmax layer over heavy-tailed classes, such as tokens, has many (\cref{ch:transformers}). \Textcite{ReddiKaleKumar2018} built their counterexample to Adam with a small $\beta_2$ from the same ingredient, a large gradient that comes rarely against small ones of the opposite sign that come often (\cref{ch:adaptive}). The proposition also shows that the constant in \cref{thm:sign:signsgd} cannot be lowered: as $\eta\to0$ and $T\to\infty$ the average of $\abs{\nabla\Loss(b_t)}$ is $\pi$ and the bound is $\sqrt{\pi(1-\pi)/B}$, and their ratio tends to $1$ for $B=1$ and $\pi\uparrow\frac12$, and to $\sqrt{\ln2}\approx0.83$ for large $B$ and $\pi$ just below $1-2^{-1/B}$. For symmetric noise the median is the mean, and the bias disappears.

\begin{proposition}[Symmetric unimodal noise: no bias]\label[proposition]{prop:sign:symmetric}
Let \cref{ass:sign:smooth,ass:sign:noise} hold with every $\sigma_i>0$, and let each $\zeta_{t,i}$ be, conditionally on the past, symmetric about $0$ and unimodal. With $\Phi_B(\mu)=\sum_i\min\{\sqrt B\mu_i^2/(\sqrt3\sigma_i),\frac23\abs{\mu_i}\}$, signSGD with a constant step size $\eta$ satisfies
\begin{equation}\label{eq:sign:mixed}
  \frac1T\sum_{t<T}\E\,\Phi_B\bigl(\nabla\Loss(\theta_t)\bigr)\le\frac{\Delta_0}{\eta T}+\frac\eta2\norm{\vec L}_1 .
\end{equation}
\end{proposition}

\begin{proof}
Repeat the proof of \cref{thm:sign:signsgd} with \cref{lem:sign:signerror}(iii), for $s=\sigma_i/\sqrt B$, in place of (ii): the one-step inequality becomes $\E[\Loss(\theta_{t+1})\mid\cF_t]\le\Loss(\theta_t)-\eta\,\Phi_B(\nabla\Loss(\theta_t))+\frac12\eta^2\norm{\vec L}_1$, both sides of the per-coordinate bound being $0$ when $\mu_i=0$, and the sum telescopes.
\end{proof}

There is no bias term: with $\eta\propto T^{-1/2}$ the right side of \eqref{eq:sign:mixed} is $O(T^{-1/2})$ at every batch size. The batch size enters only the measure of progress, which is $\frac23\abs{\mu_i}$ in coordinates whose gradient exceeds the noise $\sigma_i/\sqrt B$ by a constant factor, and $\sqrt B\mu_i^2/(\sqrt3\sigma_i)$ in coordinates dominated by noise. In the second regime signSGD is a preconditioned gradient method.

\begin{heuristic}[Heuristic derivation of the effective step]
Let the noise be Gaussian and $\abs{\mu_i}\ll\sigma_i/\sqrt B$. Then $\Var(\sign g_{t,i})\approx1$ and
\[
  \E[\sign g_{t,i}]=\erf\Bigl(\frac{\sqrt B\,\mu_i}{\sqrt2\,\sigma_i}\Bigr)\approx\sqrt{\frac{2B}\pi}\,\frac{\mu_i}{\sigma_i},
\]
so a step of signSGD has the mean and twice the variance, up to the factor $\pi/2$ in place of $2$, of a step of stochastic gradient descent with the step size $\alpha_i=\eta\sqrt{2B/\pi}/\sigma_i$: signSGD is gradient descent preconditioned by $\diag(\sigma_i)^{-1}$, with noise inflated by $\sqrt{\pi/2}$. On $\Loss(\theta)=\frac12h\theta^2$, the stationary variance $\alpha\frac\pi2\frac{\sigma^2}{B}/(2h)$ of that method gives the floor
\begin{equation}\label{eq:sign:gaussfloor}
  \E\abs{\nabla\Loss}\approx(2\pi)^{-1/4}\sqrt{\eta h\sigma}\;B^{-1/4},
\end{equation}
which vanishes as $\eta\to0$. Keeping $\alpha_i\propto B$, the linear scaling rule of gradient descent (\cref{ch:noise}), requires $\eta\propto\sqrt B$: the \term{square-root scaling rule}\index{square-root scaling rule}, which \textcite{MalladiEtAl2022} derived for adaptive methods from stochastic differential equations.
\end{heuristic}

Memory changes the picture in two ways. Momentum averages gradients over time before the sign is taken, which moves the median towards the mean; \cref{ch:momentum} proves that normalized steepest descent with momentum converges in every norm, and \cref{fig:sign:noise} shows Signum learning the rare class that defeats signSGD. Adam's second moment changes the length of the step, and the first question is how long the step can get.

\begin{proposition}[The size of Adam's step]\label[proposition]{prop:sign:adambound}
Let $0\le\beta_1<1$, $0<\beta_2<1$, $\rho=\beta_1^2/\beta_2$ and $\epsilon=0$, and consider one coordinate of Adam without bias correction. For every $t\ge0$ and every gradient sequence with $v_t>0$,
\begin{equation}\label{eq:sign:adambound}
  \frac{\abs{m_t}}{\sqrt{v_t}}\le\frac{1-\beta_1}{\sqrt{1-\beta_2}}\Bigl(\sum_{k=0}^t\rho^k\Bigr)^{1/2},
\end{equation}
with equality if and only if $g_{t-k}=c\,(\beta_1/\beta_2)^k$ for $k=0,\dots,t$ and some $c\ne0$ (with $0^0=1$). Consequently, if $\beta_1^2<\beta_2$, the supremum over all $t$ and all gradient sequences is
\begin{equation}\label{eq:sign:C}
  C(\beta_1,\beta_2)=\frac{1-\beta_1}{\sqrt{(1-\beta_2)(1-\beta_1^2/\beta_2)}},
\end{equation}
and if $\beta_1^2\ge\beta_2$ the ratio is unbounded. With bias correction the bound \eqref{eq:sign:adambound} is multiplied by $\sqrt{1-\beta_2^{t+1}}/(1-\beta_1^{t+1})$; it is then $1$ at $t=0$ and below $C(\beta_1,\beta_2)$ for all $t$ (\cref{ex:sign:biascorr}).
\end{proposition}

\begin{proof}
The plan is to write $m_t$ as a weighted sum of gradients and to apply the Cauchy--Schwarz inequality with the weights of $v_t$. From $m_{-1}=v_{-1}=0$, $m_t=(1-\beta_1)\sum_{k=0}^t\beta_1^kg_{t-k}$ and $v_t=(1-\beta_2)\sum_{k=0}^t\beta_2^kg_{t-k}^2$. Splitting $\beta_1^kg_{t-k}=(\beta_1^k\beta_2^{-k/2})(\beta_2^{k/2}g_{t-k})$,
\[
  \Bigl\lvert\sum_{k=0}^t\beta_1^kg_{t-k}\Bigr\rvert\le\Bigl(\sum_{k=0}^t\rho^k\Bigr)^{1/2}\Bigl(\sum_{k=0}^t\beta_2^kg_{t-k}^2\Bigr)^{1/2}=\Bigl(\sum_{k=0}^t\rho^k\Bigr)^{1/2}\Bigl(\frac{v_t}{1-\beta_2}\Bigr)^{1/2},
\]
and multiplying by $1-\beta_1$ gives \eqref{eq:sign:adambound}. Equality holds exactly when $(\beta_2^{k/2}g_{t-k})_k$ is parallel to $(\beta_1^k\beta_2^{-k/2})_k$, that is, $g_{t-k}=c(\beta_1/\beta_2)^k$, with $c\ne0$ because $v_t>0$. If $\rho<1$ the sums $\sum_{k\le t}\rho^k$ increase to $1/(1-\rho)$, and the equality sequences make \eqref{eq:sign:C} the supremum; if $\rho\ge1$ the sums are unbounded. At $t=0$ the bias correction factor is $\sqrt{1-\beta_2}/(1-\beta_1)$, the reciprocal of the bound.
\end{proof}

The bound is $C=7.27$ at the defaults $(0.9,0.999)$, $2.35$ at $(0.9,0.99)$, $1.16$ at $(0.9,0.95)$, and $1$ whenever $\beta_1=\beta_2$. A single large gradient after a run of zeros gives only $(1-\beta_1)/\sqrt{1-\beta_2}$, which is $3.16$ at the defaults;\sidenote{\Textcite[Section~2.1]{KingmaBa2015} describe the effective step as approximately bounded by this single-spike value or by the learning rate. \Cref{prop:sign:adambound} gives the exact supremum.} the worst case is a gradient that grows by the factor $\beta_2/\beta_1$ at every step, so that the slow second moment lags behind the fast first moment. Adam is thus sign descent of its momentum with a bounded multiplier: each coordinate moves by $\eta$ times a number in $[0,C]$ in the direction of $-\sign(m_t)$, and $\norm{\theta_{t+1}-\theta_t}_\infty\le C\eta$, a trust region in the max norm. A smaller $\beta_2$ keeps the multiplier near $1$, and the condition $\beta_1<\sqrt{\beta_2}$ for a bounded step reappears in convergence analyses of Adam and AMSGrad \parencite{ReddiKaleKumar2018,ZhangEtAl2022}. The typical multiplier, as opposed to the worst one, is set by the signal-to-noise ratio.

\begin{heuristic}[Heuristic derivation of Adam's mean step]
Let the gradients of coordinate $i$ be independent over steps with mean $\mu_i$ and variance $s_i^2$, and let $\beta_2$ be so close to $1$ that $\hat v_{t,i}$, an average of about $1/(1-\beta_2)$ squared gradients, concentrates near $\mu_i^2+s_i^2$; $\hat m_{t,i}$ has mean $\mu_i$ and variance $(1-\beta_1)^2\sum_k\beta_1^{2k}s_i^2=\frac{1-\beta_1}{1+\beta_1}s_i^2$. Treating $\hat v_{t,i}$ as a constant,
\begin{equation}\label{eq:sign:snr}
  \E\Bigl[\frac{\hat m_{t,i}}{\sqrt{\hat v_{t,i}}}\Bigr]\approx\frac{\sign(\mu_i)}{\sqrt{1+s_i^2/\mu_i^2}},
\end{equation}
and the mean square of the update is about $\bigl(\mu_i^2+\frac{1-\beta_1}{1+\beta_1}s_i^2\bigr)/(\mu_i^2+s_i^2)$. The mean update is the sign shrunk by the factor $(1+s_i^2/\mu_i^2)^{-1/2}$ of \textcite{BallesHennig2018}: near $\sign(\mu_i)$ for a consistent gradient and near $\mu_i/s_i$ for noise, where Adam, like signSGD, is gradient descent preconditioned by the noise. The root mean square of the update lies between $\sqrt{(1-\beta_1)/(1+\beta_1)}=0.23$, for pure noise at $\beta_1=0.9$, and $1$.
\end{heuristic}

For $(\beta_1,\beta_2)=(0.9,0.999)$ and Gaussian gradients, simulation confirms \eqref{eq:sign:snr} and the mean square to within $0.002$ for signal-to-noise ratios from $0$ to $3$, and in $2\times10^6$ steps of pure noise the multiplier never exceeded $1.5$, far from the worst case $7.27$. SignSGD's mean update under Gaussian noise is exactly $\erf\bigl(\mu_i/(\sqrt2s_i)\bigr)$, within $0.07$ of \eqref{eq:sign:snr}: in this regime Adam is very nearly a sign method with a different nonlinearity.

\section{Reading the theory}\label{sec:sign:reading}

\paragraph{The guarantee is exact, and it belongs to the basis.} \Cref{fig:sign:rotation} runs gradient descent and sign descent on $\Loss(\theta)=\frac12\sum_ih_i\theta_i^2$ with $d=256$ and $h_i=i^{-2}$, so that $\tr H=1.64$ and the effective rank is below $2$, and on the same quadratic in randomly rotated coordinates. In panel (b) the long-run average of $\norm{\nabla\Loss}_1$ on the diagonal problem lies on the line $\frac12\eta\tr H$ of \cref{prop:sign:floor} to all printed digits. On the rotated problem it lies on $\frac12\eta\hat L_\infty$, where $\hat L_\infty=161.2$ is the largest $s^\top Q^\top HQs$ that a local search over sign vectors finds; the certified lower bound of \cref{prop:sign:rotation}(ii) is $160.9$, and its mean over rotations $163.3$. In every run the final two-cycle used the best sign vector of the search, so the floor $\frac12\eta\norm{\cdot}_{\infty\to1}$ of \cref{thm:sign:descent} is attained in both bases, and the rotation multiplies it by $98$, against the $\frac2\pi d\lmax/\tr H=99$ of \cref{prop:sign:rotation}.

\begin{figure}[t]
\begin{fullwidth}
\centering
\begin{tikzpicture}
\begin{groupplot}[group style={group size=2 by 1,horizontal sep=17mm},wplot,width=82mm,height=54mm]
\nextgroupplot[xmode=log,ymode=log,xlabel={step $t$},ylabel={loss $\Loss(\theta_t)-\Loss^\star$},xmin=1,xmax=1e4,ymin=1e-5,ymax=1,
  legend style={at={(0.02,0.02)},anchor=south west,font=\scriptsize},legend cell align=left]
\addplot[theory] table[x=t,y=gdexact] {figures/data/02-sign-rotation-loss.dat};
\addplot[theory3] table[x=t,y=signexact] {figures/data/02-sign-rotation-loss.dat};
\addplot[simulation2] table[x=t,y=gdrot] {figures/data/02-sign-rotation-loss.dat};
\addplot[simulation] table[x=t,y=signdiag] {figures/data/02-sign-rotation-loss.dat};
\addplot[only marks,mark=triangle*,mark size=2pt,color=sea,mark options={fill=sea,draw=white,line width=0.2pt}] table[x=t,y=signrot] {figures/data/02-sign-rotation-loss.dat};
\legend{gradient descent (exact; both bases),sign descent (exact; diagonal),gradient descent (rotated),sign descent (diagonal),sign descent (rotated)}
\nextgroupplot[xmode=log,ymode=log,xlabel={step size $\eta$},ylabel={long-run average of $\norm{\nabla\Loss(\theta_t)}_1$},xmin=7e-4,xmax=0.14,ymin=3e-4,ymax=20,
  legend style={at={(0.02,0.98)},anchor=north west,font=\scriptsize},legend cell align=left]
\addplot[theory] table[x=eta,y=diagtheory] {figures/data/02-sign-rotation-floor.dat};
\addplot[theory3] table[x=eta,y=rottheory] {figures/data/02-sign-rotation-floor.dat};
\addplot[simulation] table[x=eta,y=diagsim] {figures/data/02-sign-rotation-floor.dat};
\addplot[simulation2] table[x=eta,y=rotsim] {figures/data/02-sign-rotation-floor.dat};
\legend{$\frac12\eta\tr H$ (diagonal),$\frac12\eta\hat L_\infty$ (rotated),diagonal (run),rotated (run)}
\end{groupplot}
\end{tikzpicture}
\end{fullwidth}
\caption{Gradient descent and sign descent on a quadratic with curvatures $h_i=i^{-2}$, $d=256$, $\theta^\star=0$ and $\theta_0\sim\cN(0,I)$, in the eigenbasis and after one uniformly random rotation (seed in \texttt{02-sign-rotation.py}). Left: the loss of gradient descent with step size $1/\lmax$ and of sign descent with $\eta=0.01$; curves are exact closed forms, points are runs. Gradient descent is the same in both bases; sign descent is far ahead on the diagonal problem and is caught on a two-cycle with loss $2.1\times10^{-2}$ after about $300$ steps on the rotated one. Right: the average of $\norm{\nabla\Loss}_1$ over the final cycle against $\eta$, with the floor of \cref{prop:sign:floor} and the floor $\frac12\eta\hat L_\infty$ of \cref{thm:sign:descent} for the rotated Hessian, $\hat L_\infty$ found by local search.}
\label{fig:sign:rotation}
\end{figure}

\paragraph{A time scale without curvature.} On the diagonal problem coordinate $i$ reaches its two-cycle after $\lfloor\abs{\theta_{0,i}}/\eta\rfloor$ steps whatever its curvature, so sign descent with $\eta=0.01$ is at its floor, near $2\times10^{-5}$, after about $300$ steps. Gradient descent with step size $1/\lmax$ needs about $\lmax/h_i=i^2$ steps for coordinate $i$, and its expected loss $\frac12\sum_ii^{-2}(1-i^{-2})^{2t}$ decays like $\frac{\sqrt\pi}{4\sqrt{2t}}\approx0.31\,t^{-1/2}$ while $t\ll d^2$: it is $1.5\times10^{-2}$ at $t=300$ and $1.3\times10^{-3}$ at $t=10^4$. This is the regime in which \cref{cor:sign:compare} favors the sign: curvature of effective rank below $2$, a spread gradient, and distances of the same order in all coordinates.

\paragraph{The trap after a rotation.} On the rotated problem gradient descent is unchanged to rounding error. Sign descent makes fast progress for a hundred steps and is then caught on a two-cycle with loss $2.1\times10^{-2}$, a thousand times its diagonal floor, after which gradient descent is ahead. The centers of the two-cycles form a box in gradient space, which the inverse Hessian stretches along the directions of small curvature; when these are not coordinate axes, the box contains points of high loss (\cref{ex:sign:cycles}). A decreasing step size breaks the cycle; Adam, whose effective step changes over time, escapes it but still loses orders of magnitude to the rotation (\cref{ex:sign:adamrot}).

\paragraph{Median against mean.} \Cref{fig:sign:noise}(a) shows one coordinate of signSGD with $B=16$ under two noises of the same variance. The iterates move on a lattice of spacing $\eta$ and form a birth--death chain, whose stationary law is computed exactly. With Gaussian noise the floor follows \eqref{eq:sign:gaussfloor} and vanishes like $\sqrt\eta$, as \cref{prop:sign:symmetric} requires. With rare-event noise, in which a fraction $p=0.04$ of the examples push hard one way and the rest gently the other way, $(1-p)^B=0.52>\frac12$, and as $\eta\to0$ the floor rises to $c=\sigma\sqrt{p/(1-p)}=0.82\,\sigma/\sqrt B$, the gradient at the point where the median of the minibatch gradient vanishes, $82\%$ of the bias term of \cref{thm:sign:signsgd}. A smaller step makes it worse; the lowest floor is near $\eta=0.03$. Panel (b) is \cref{prop:sign:rareclass} for a class of frequency $\pi=0.01$: the bias drifts at exactly the predicted speed for $B=16,32,64$, the speed vanishes at the threshold $B=\ln2/(-\ln(1-\pi))=69$, and Signum with $\beta=0.9$ holds the bias at its optimum.

\begin{figure}[t]
\begin{fullwidth}
\centering
\begin{tikzpicture}
\begin{groupplot}[group style={group size=2 by 1,horizontal sep=17mm},wplot,width=82mm,height=54mm]
\nextgroupplot[xmode=log,ymode=log,xlabel={step size $\eta$},ylabel={stationary $\E\abs{\nabla\Loss(\theta_t)}$},xmin=8e-4,xmax=0.4,ymin=7e-3,ymax=0.6,
  legend style={at={(0.98,0.02)},anchor=south east,font=\scriptsize},legend cell align=left]
\addplot[guide] table[x=eta,y=bound] {figures/data/02-sign-noise-eta.dat};
\addplot[guide,sandline!55!black] table[x=eta,y=bias] {figures/data/02-sign-noise-eta.dat};
\addplot[theory] table[x=eta,y=gaussexact] {figures/data/02-sign-noise-eta.dat};
\addplot[theory2] table[x=eta,y=heuristic] {figures/data/02-sign-noise-eta.dat};
\addplot[theory3] table[x=eta,y=rareexact] {figures/data/02-sign-noise-eta.dat};
\addplot[simulation2] table[x=eta,y=gausssim] {figures/data/02-sign-noise-eta.dat};
\addplot[simulation] table[x=eta,y=raresim] {figures/data/02-sign-noise-eta.dat};
\legend{bound of \cref{thm:sign:signsgd},median gap $c$,Gaussian (exact),heuristic \eqref{eq:sign:gaussfloor},rare event (exact),Gaussian (runs),rare event (runs)}
\nextgroupplot[xlabel={step $t$},ylabel={bias $b_t$},xmin=0,xmax=20000,ymin=-150,ymax=8,xtick={0,5000,10000,15000,20000},
  xticklabels={$0$,$5000$,$10000$,$15000$,$20000$},legend style={at={(0.02,0.02)},anchor=south west,font=\scriptsize},legend cell align=left]
\addplot[theory] table[x=t,y=theory16] {figures/data/02-sign-noise-rare.dat};
\addplot[theory2] table[x=t,y=theory32] {figures/data/02-sign-noise-rare.dat};
\addplot[theory3] table[x=t,y=theory64] {figures/data/02-sign-noise-rare.dat};
\addplot[simulation,mark size=1.2pt] table[x=t,y=sim16] {figures/data/02-sign-noise-rare.dat};
\addplot[simulation,mark size=1.2pt,forget plot] table[x=t,y=sim32] {figures/data/02-sign-noise-rare.dat};
\addplot[simulation,mark size=1.2pt,forget plot] table[x=t,y=sim64] {figures/data/02-sign-noise-rare.dat};
\addplot[simulation2,mark size=1.3pt] table[x=t,y=signum16] {figures/data/02-sign-noise-rare.dat};
\legend{drift for $B=16$,drift for $B=32$,drift for $B=64$,signSGD (runs),Signum{,} $B=16$ (runs)}
\end{groupplot}
\end{tikzpicture}
\end{fullwidth}
\caption{SignSGD and the shape of the noise. Left: one coordinate, $\Loss(\theta)=\frac12h\theta^2$ with $h=\sigma=1$ and $B=16$, under Gaussian per-example noise and under the rare-event noise $\sigma(X-p)/\sqrt{p(1-p)}$ with $X$ Bernoulli($p$), $p=0.04$. Curves: exact stationary values of the lattice chain; points: $1000$ runs averaged over at least $4000$ steps after a burn-in; dashed: the bound $\frac12\eta h+\sigma/\sqrt B$ and the median gap $c$. Right: the bias of a class of frequency $\pi=0.01$ under signSGD with $\eta=0.01$ from $b^\star$, averaged over $200$ runs; lines are the exact drift $-\eta(2(1-\pi)^B-1)t$; squares are Signum ($\beta=0.9$). Script: \texttt{02-sign-noise.py}.}
\label{fig:sign:noise}
\end{figure}

\section{What the theory misses}\label{sec:sign:misses}

\paragraph{What was idealized away.} Memory first: the theorems are about the memoryless step, Adam's averages enter only through a bound on its step and a heuristic for its mean, and nothing here proves that Adam converges or explains why $\beta_1=0.9$ and $\beta_2$ near $1$ work (\cref{ch:momentum,ch:adaptive}). Adam's $\epsilon$ is zero except in \cref{ex:sign:epsilon}. The curvature is bounded by one vector $\vec L$ everywhere, a strong assumption for networks, whose sharpness adapts to the step size (\cref{ch:eos}); the proofs use $\vec L$ only along the trajectory, but nothing tells us its value in advance. The noise is described by its variance, yet the outcome depends on its shape: the same variance gives no bias when the noise is symmetric and divergence when it is skewed. And most statements use a constant step size, while practice decays it (\cref{ch:schedules}).

\paragraph{Where it breaks.} (1) SignSGD without momentum fails at small batch sizes: the rare class of \cref{prop:sign:rareclass} is convex and smooth, and its loss diverges. \Textcite{KarimireddyEtAl2019} give further examples in which signSGD does not converge, and \textcite{SafaryanRichtarik2021} prove convergence under a condition on the probability that each sign is correct, which fails exactly when the median and the mean of a coordinate's gradient have different signs. (2) Adam with long memories is not sign descent: at $\beta_2=0.999$ its multiplier can reach $7.27$, and typically it is a smooth function of the signal-to-noise ratio \eqref{eq:sign:snr}, so the theorems for sign descent do not transfer to it. (3) \Cref{cor:sign:compare} compares upper bounds, and worst-case rates do not explain the size of the gap between Adam and gradient descent on transformers, which \textcite{KunstnerEtAl2023} measure in full batch (\cref{ch:transformers}). (4) The advantage of the max norm depends on the basis (\cref{prop:sign:rotation}), and the theory does not say why the coordinates of trained networks are favorable. \Textcite{XieMohamadiLi2025} relate Adam's sensitivity to rotations to a smoothness constant of the $\ell_\infty$ type, and \textcite{ZhangEtAl2024} find transformer Hessians close to block diagonal, with blocks that correspond to groups of parameters; both are measurements, not derivations. (5) The floor of sign descent can be far worse than \cref{prop:sign:floor} suggests: after a rotation, the trapping two-cycle has a loss a thousand times higher (\cref{fig:sign:rotation}).

\paragraph{Repairs and open problems.} Momentum is the main repair: averaging before taking the sign reduces the noise and moves the median towards the mean, normalized steepest descent with momentum converges in every norm (\cref{ch:momentum}), and Signum and Lion train language models about as well as Adam \parencite{ZhaoEtAl2025}. Error feedback, which adds the part of the gradient lost to compression to the next gradient, makes sign compression converge \parencite{SeideEtAl2014,KarimireddyEtAl2019}; majority vote reduces the noise of the sign when it is symmetric and unimodal \parencite{BernsteinEtAl2019} (\cref{ex:sign:vote}). Analyses of Adam in the $\ell_\infty$ geometry \parencite{XieMohamadiLi2025}, its implicit bias towards $\ell_\infty$ max-margin solutions \parencite{ZhangZouCao2024} (\cref{ch:bias}), and the $\ell_\infty$-constrained problems that AdamW and Lion solve with weight decay \parencite{XieLi2024,ChenLiuLiangLiu2024} (\cref{ch:decay}) carry the max-norm picture to the full algorithms; Shampoo and SOAP remove the dependence on the basis by rotating into an estimated eigenbasis (\cref{ch:kronecker}). Open: why the coordinates of trained networks are good ones for sign methods (\cref{ex:sign:basis}); a convergence theory of Adam that uses both its geometry and its memory and explains its speed; and the batch sizes and momenta at which sign methods stop ignoring rare events in heavy-tailed data (\cref{ex:sign:rare}).

\section{Back in the wild}\label{sec:sign:back}

\fieldnotes{Adam, Lion and signSGD move every parameter by about the learning rate per step: Lion and signSGD exactly, Adam within the factor $C(\beta_1,\beta_2)$.}
{Sign descent as normalized steepest descent in $\ell_\infty$, with $\nabla^2\Loss\preceq\diag(\vec L)$ and minibatch noise of variance $\sigma_i^2/B$.}
{Decrease $\eta\norm{\nabla\Loss}_1-\frac12\eta^2\norm{\vec L}_1$, exact on separable quadratics; rotations move $\norm{H}_{\infty\to1}$ between $\tr H$ and $d\lmax$; a sharp bias $\norm{\vec\sigma}_1/\sqrt B$, absent for symmetric unimodal noise; $\abs{m_t}/\sqrt{v_t}\le C(\beta_1,\beta_2)$, $7.27$ at the defaults, $1$ if $\beta_1=\beta_2$.}
{Sign descent wins when the gradient is spread over more coordinates than the curvature has stiff directions; a random rotation costs about $\frac2\pi d\lmax/\tr H$; signSGD ignores classes rarer than about $\ln2/B$; sign methods follow the square-root scaling rule (derived).}
{The memory that makes Adam and Lion work; why the coordinates of networks are good ones; the size of the Adam--SGD gap.}
{Use sign updates with momentum, not without, at small batch sizes (proved in the model); expect the advantage of Adam-like methods to depend on the basis (proved for quadratics, observed for GPT-2); a smaller $\beta_2$ keeps Adam's step near the learning rate (proved bound).}

What the theory tells a practitioner is modest and specific. A sign method with momentum suits problems whose gradients are spread over many parameters and whose curvature is concentrated in a few directions that the coordinates respect; whether a given network is such a problem the theory does not say. A sign method without momentum should not be trusted at small batch sizes, because it follows the median of the gradient, and the median ignores rare events. Adam's learning rate is the length of a step to within the factor $C(\beta_1,\beta_2)$, independent of the scale of the gradients when $\epsilon$ is negligible, which is why Lion, whose step always has RMS norm $1$, needs a smaller learning rate than Adam, whose update RMS is usually well below $1$. Claims beyond these, such as that Lion generalizes better than AdamW, are folklore as far as this chapter goes.

The spectral norm plays the role of the max norm for matrices (\cref{ch:spectral}); for a layer, sign descent moves the output by an amount proportional to the fan-in, which is why $\mu$P scales Adam's learning rate by $1/\mathrm{fan\_in}$ (\cref{ch:modular}). Momentum repairs the bias of the sign (\cref{ch:momentum}), the second moment turns the sign into a preconditioner (\cref{ch:adaptive}), the dependence on the basis returns with Shampoo and SOAP and with rotated transformers and rare tokens (\cref{ch:kronecker,ch:transformers}), and $C(\beta_1,\beta_2)$ bounds Adam's largest steps when an instability begins (\cref{ch:spikes}).

\begin{exercises}
\item \consolidation (Bias correction and equal coefficients.)\label{ex:sign:biascorr} Let $\rho=\beta_1^2/\beta_2<1$ and $\epsilon=0$. (a) Show that with bias correction $\abs{\hat m_t}/\sqrt{\hat v_t}\le B_t$, where
\[
  B_t=\frac{(1-\beta_1)\sqrt{1-\beta_2^{t+1}}}{(1-\beta_1^{t+1})\sqrt{1-\beta_2}}\Bigl(\frac{1-\rho^{t+1}}{1-\rho}\Bigr)^{1/2},\qquad B_0=1 .
\]
(b) Show that $B_t\le C(\beta_1,\beta_2)$ for every $t$, with equality for every $t$ if and only if $\beta_1=\beta_2$. (c) For $(0.9,0.999)$, compute $B_{10}$ and $B_{100}$ (answers: $1.05$ and $2.25$). (d) Let $\beta_1=\beta_2\in(0,1)$. Show that $\hat s_t^2=\hat v_t-\hat m_t^2$ is a weighted variance of the past gradients and that $\hat m_t/\sqrt{\hat v_t}=\sign(\hat m_t)(1+\hat s_t^2/\hat m_t^2)^{-1/2}$, so that \eqref{eq:sign:snr} is an identity in empirical moments; deduce $\abs{\hat m_t}\le\sqrt{\hat v_t}$, with equality if and only if $g_0=\dots=g_t$. \hint{with $a=\beta_1^{t+1}$ and $b=\beta_2^{t+1}$, the claim $B_t\le C$ is $(1-b)(1-a^2/b)\le(1-a)^2$.}

\item \consolidation\hard (Two-cycles of sign descent.)\label{ex:sign:cycles} Let $\Loss(\theta)=\frac12(\theta-\theta^\star)^\top H(\theta-\theta^\star)$ with $H\succ0$, and let $s\in\{\pm1\}^d$. (a) Show that $\{\theta,\theta-\eta s\}$ is an orbit of period two of sign descent if and only if $\sign(Hs)=s$ and the center $\bar\theta=\theta-\frac12\eta s$ satisfies $\abs{(H(\bar\theta-\theta^\star))_i}<\frac12\eta\abs{(Hs)_i}$ for every $i$. (b) Show that over the cycle $\norm{\nabla\Loss}_1$ averages $\frac12\eta s^\top Hs$ and the loss averages $\frac12(\bar\theta-\theta^\star)^\top H(\bar\theta-\theta^\star)+\frac18\eta^2s^\top Hs$. (c) For diagonal $H$, show that every two-cycle has average loss below $\frac14\eta^2\tr H$; in general, show that the centers form $\theta^\star+H^{-1}W$ for a box $W$, and explain the trap of \cref{fig:sign:rotation}. (d) Show that $s\mapsto\sign(Hs)$ (with $\sign0$ replaced by $1$) does not decrease $s^\top Hs$, and that a maximizer of $s^\top Hs$ satisfies $\sign(Hs)=s$. \hint{in (d), $(s'-s)^\top H(s'-s)\ge0$ for $s'=\sign(Hs)$, and flipping $s_i$ changes $s^\top Hs$ by $4H_{ii}-4s_i(Hs)_i$.}

\item \consolidation (Cantelli is sharp.) For $\mu,s,\delta>0$ let $X=\mu+\zeta$, where $\zeta=-\mu-\delta$ with probability $p=s^2/(s^2+(\mu+\delta)^2)$ and $\zeta=p(\mu+\delta)/(1-p)$ otherwise. Check that $\zeta$ has mean $0$ and variance $s^2$, compute $\mu-\mu\E\sign X$, and show that its supremum over $\mu$ and $\delta$ is $s$: the constant of \cref{lem:sign:signerror}(ii) cannot be lowered.

\item \consolidation\harder (Coordinatewise against $\ell_\infty$ smoothness.)\label{ex:sign:sdp} Let $H\succeq0$. (a) Show that $\min\{\norm{\vec L}_1:\diag(\vec L)\succeq H\}=\max\{\tr(HX):X\succeq0,\ X_{ii}=1\}$. (b) For feasible $X=VV^\top$, $\gamma\sim\cN(0,I)$ and $s=\sign(V\gamma)$, show that $\E[ss^\top]=\frac2\pi\arcsin(X)$ entrywise and that $\arcsin(X)\succeq X$. (c) Deduce $\norm{H}_{\infty\to1}\ge\frac2\pi\tr(HX)$ and the factor $\pi/2$ of \cref{rem:sign:sdp} \parencite{Nesterov1998}. \hint{(a) is semidefinite duality with strictly feasible points on both sides; in (b), $\E[\sign(a^\top\gamma)\sign(b^\top\gamma)]=1-2\angle(a,b)/\pi$, and the series of $\arcsin$ has nonnegative coefficients, so the Schur product theorem applies.}

\item \variation (Majority vote.)\label{ex:sign:vote} Let $M$ be odd, and let $M$ workers compute independent minibatch gradients whose coordinate noises are symmetric and unimodal with variances at most $\sigma_i^2/B$; the server steps by $-\eta\sign\bigl(\sum_m\sign(g^{(m)}_{t})\bigr)$. (a) With \cref{lem:sign:signerror}(iii) and Hoeffding's inequality, show that the vote has the wrong sign in coordinate $i$ with probability at most $\exp\bigl(-\frac M2\min\{S_i^2/3,\frac49\}\bigr)$, where $S_i=\sqrt B\abs{\partial_i\Loss}/\sigma_i$. (b) Prove the analogue of \cref{thm:sign:signsgd}. (c) Compare with one worker and batch $MB$ under Gaussian noise, whose error probability is $\Phi(-\sqrt MS_i)\le\frac12e^{-MS_i^2/2}$.

\item \variation (Lion.)\label{ex:sign:lion} (a) Show that Lion with $\beta_1=\beta_2$ is Signum. (b) Show that $c_t=\sum_{k=0}^tw_kg_{t-k}$ with $w_0=1-\beta_1$ and $w_k=\beta_1(1-\beta_2)\beta_2^{k-1}$ for $k\ge1$, that the weights sum to $1-\beta_1\beta_2^t$, and that their mean lag $\sum_kkw_k$ tends to $\beta_1/(1-\beta_2)$, which is $90$ steps for the defaults; compare with Signum at $\beta=\beta_2$. (c) Lion's update has RMS norm $1$ when no coordinate of $c_t$ vanishes. Using \eqref{eq:sign:snr}, explain why matching AdamW's movement per step calls for a Lion learning rate several times smaller.

\item \variation (Memoryless Adam with $\epsilon$.)\label{ex:sign:epsilon} Let $\theta_{t+1}=\theta_t-\eta\,\psi(\nabla\Loss(\theta_t))$ with $\psi(x)=x/(\abs x+\epsilon)$ coordinatewise. (a) Under \cref{ass:sign:smooth}, show that $\Loss(\theta_{t+1})\le\Loss(\theta_t)-\eta\sum_ig_i^2/(\abs{g_i}+\epsilon)+\frac12\eta^2\norm{\vec L}_1$ for $g=\nabla\Loss(\theta_t)$. (b) Show that $x^2/(\abs x+\epsilon)\ge\frac12\min\{\abs x,x^2/\epsilon\}$, and deduce a bound of the form \eqref{eq:sign:mixed} with $\epsilon$ in the role of $\sigma_i/\sqrt B$. (c) When does $\epsilon=10^{-8}$ matter?

\item \variation (Block-diagonal curvature.) Let $H=\diag(H_1,\dots,H_K)$ with blocks of sizes $d_1,\dots,d_K$. (a) Show that $\norm{H}_{\infty\to1}=\sum_k\norm{H_k}_{\infty\to1}$. (b) If each block is rotated by its own orthogonal matrix, show that the $\ell_\infty$ constant is at most $\sum_kd_k\lmax(H_k)$, and compare with $d\lmax(H)$ for a rotation that mixes the blocks. (c) Relate this to the block structure that \textcite{ZhangEtAl2024} measure in transformers.

\item \simulation (Shapes of noise.) With \texttt{02-sign-noise.py}: (a) vary $p$ and $B$ and check that the rare-event floor at small $\eta$ collapses as soon as $(1-p)^B<\frac12$. (b) Replace the rare-event noise by $\sigma(E-1)$ with $E$ exponential of mean $1$. Show that as $\eta\to0$ the floor at $B=1$ is the median gap $\sigma(1-\ln2)=0.31\sigma$, and that for $B\ge2$ it is within $3\%$ of $\sigma/(3B)$, the median gap of a gamma law. Conclude that for a fixed skewed law the bias decays like $1/B$, and that the worst case $1/\sqrt B$ of \cref{thm:sign:signsgd} needs events of probability comparable to $1/B$.

\item \simulation (Adam after a rotation.)\label{ex:sign:adamrot} With \texttt{02-sign-rotation.py}, run Adam with $(\beta_1,\beta_2)=(0.9,0.999)$, $\epsilon=0$ and $\eta=0.01$ on both problems. You should find a loss near $10^{-9}$ after $1000$ steps on the diagonal problem and near $3\times10^{-3}$ on the rotated one, still ahead of gradient descent. Show that along the rotated run the preconditioned sharpness $\eta\lmax(P^{1/2}Q^\top HQP^{1/2})$, with $P=\diag(\hat v_t)^{-1/2}$, grows towards the stability threshold $(2+2\beta_1)/(1-\beta_1)=38$ of momentum (\cref{ch:momentum,ch:eos}), and explain why Adam's advantage there has nothing to do with the coordinates.

\item \research (When are the coordinates good?)\label{ex:sign:basis} (a) Show that $\norm{H}_{\infty\to1}\le\sum_j\lambda_j\norm{v_j}_1^2$ for $H=\sum_j\lambda_jv_jv_j^\top$, and find conditions on the localization of the eigenvectors under which $\norm{H}_{\infty\to1}/\tr H=O(1)$. (b) Estimate both sides for the Hessian of a small network with Hessian--vector products and the iteration of \cref{ex:sign:cycles}(d). What is open is a mechanism that explains why the Hessians of trained networks are closer to diagonal, or block diagonal, in the parameters' coordinates than a random rotation would make them.

\item \research (Rare events and memory.)\label{ex:sign:rare} (a) For Signum with momentum $\beta$ in the model of \cref{prop:sign:rareclass}, argue that the update follows the sign of an average of about $B(1+\beta)/(1-\beta)$ labels, derive the class frequency below which Signum fails, and check it by simulation. (b) For a softmax output layer with Zipf frequencies $\pi_k\propto1/k$ over $K$ classes, find the fraction of classes that signSGD and Signum learn as a function of $B$ and $\beta$. What is open is whether this mechanism affects the rare tokens of language models trained with Adam, whose $\beta_2<1$ restores the mean (\cref{ch:adaptive,ch:transformers}).
\end{exercises}

\notesandrefs

Sign methods are older than their analysis (\cref{sec:sign:wild}). \Textcite{BernsteinEtAl2018} gave the first nonconvex analysis of signSGD and Signum under coordinatewise smoothness, with the comparison of the densities of gradients and curvatures that \cref{cor:sign:compare} restates, and \textcite{BernsteinEtAl2019} analysed majority vote under symmetric unimodal noise. \Textcite{KarimireddyEtAl2019} showed that signSGD can fail to converge and that error feedback repairs it; \textcite{SafaryanRichtarik2021} replaced the noise assumptions by a condition on the probability that each sign is correct. We know of no source for \cref{prop:sign:floor,prop:sign:rareclass} in this form; both are elementary, and the second is the median-against-mean mechanism of the counterexample of \textcite{ReddiKaleKumar2018}.

Steepest descent in the max norm is a tool outside machine learning as well; \textcite{KelnerEtAl2014} used it in an almost-linear-time algorithm for approximate maximum flow. \Textcite{BallesPedregosaLeRoux2020} compared coordinatewise and $\ell_\infty$ smoothness for sign descent and related the latter to the concentration of the Hessian on its diagonal; the factor $\pi/2$ of \cref{rem:sign:sdp} is Nesterov's theorem on the semidefinite relaxation of quadratic maximization over the cube \parencite{Nesterov1998}. \Textcite{BernsteinNewhouse2024} organized optimizers by norms and identified Adam without its averages as steepest descent in $\ell_\infty$, and \textcite{XieMohamadiLi2025} developed the $\ell_\infty$ theory of Adam and tested it by rotating GPT-2's parameters.

\Textcite{KingmaBa2015} introduced Adam and discussed the size of its effective step; the supremum $C(\beta_1,\beta_2)$ takes one application of the Cauchy--Schwarz inequality. \Textcite{BallesHennig2018} separated Adam's sign from its variance adaptation, \textcite{KunstnerEtAl2023} found that Adam behaves like sign descent with momentum as the batch grows, and \textcite{ZhaoEtAl2025} that Signum matches Adam on language models. \Textcite{MalladiEtAl2022} derived the square-root scaling rule for adaptive methods. Lion was found by \textcite{ChenEtAl2023} through a symbolic program search; \textcite{ChenLiuLiangLiu2024} showed with a Lyapunov function that Lion with weight decay solves a box-constrained problem, \textcite{XieLi2024} proved the corresponding statement for AdamW, and \textcite{ZhangZouCao2024} that Adam on separable data converges in direction to an $\ell_\infty$ max-margin classifier.

\paragraph{Translating the literature.} \Textcite{BernsteinEtAl2018} write $f$ for our $\Loss$, $x_k$ for $\theta_t$, $g_k$ for the full gradient $\nabla\Loss(\theta_t)$ and $\tilde g_k$ for the stochastic gradient $g_t$, $\delta_k$ for the step size $\eta_t$, $n_k$ for the batch size $B$ and $K$ for the number of steps $T$; their $\vec L$ and $\vec\sigma$ are ours. \Textcite{KingmaBa2015} write $\alpha$ for the step size and start at $t=1$ with $m_0=v_0=0$, so their $\hat m_t=m_t/(1-\beta_1^t)$ is our $\hat m_{t-1}$. \Textcite{ChenEtAl2023} write Lion with $c_t$ for the signed interpolation, $m_t$ for the momentum and $\lambda$ for the decoupled weight decay, as we do. Some implementations of the sign map $0$ to $1$; in this chapter that changes events of probability zero only.

\printchapterbib
\end{refsection}

% Chapter 3 of Optimizers in the Wild: orthogonalized updates (Muon, Shampoo, Newton--Schulz).
\newcommand{\spectralrnuc}{r_{\mathrm{nuc}}}
\newcommand{\spectralrst}{r_{\mathrm{st}}}

\casechapter{Orthogonalized Updates}{The spectral norm, Muon and Newton--Schulz}{Muon and Shampoo}{ch:spectral}
\begin{refsection}[chapters/03-spectral.bib]

\begin{chaptersummary}
Muon replaces the momentum of each hidden weight matrix by an approximately orthogonal matrix, computed by five steps of a matrix polynomial; memoryless Shampoo computes the same matrix. For one weight matrix under the spectral norm we prove that the steepest direction is the polar factor $UV^\top$ of the gradient (a whole face of them when the gradient is rank deficient), that the cubic Newton--Schulz iteration converges to it while Muon's quintic does not converge and holds all but the smallest singular values in $[0.681,1.203]$, and that on a linear layer with $d_{\mathrm{out}}\ge d_{\mathrm{in}}$ a spectral step guarantees more decrease than a gradient step exactly when the gradient's nuclear rank exceeds the stable rank of the inputs. The theory concerns one matrix and one step: it nearly locates the crossover in a solvable case, misjudges it for generic targets, and cannot say why the gains of matrix optimizers shrink at scale.
\end{chaptersummary}

\section{In the wild: an optimizer that orthogonalizes}\label{sec:spectral:wild}

Adam updates each entry of a weight matrix on its own, as if the matrix were a list of numbers. Since 2024 another kind of optimizer has spread through the training of language models: it treats each weight matrix as a linear map and replaces its update by an orthogonal matrix. The best known is \term{Muon}\index{Muon}, described by Keller Jordan and coauthors in a blog post of December 2024 after it had set speed records on two community benchmarks \parencite{JordanEtAl2024Muon}. The name stands for MomentUm Orthogonalized by Newton--Schulz.\sidenote{The reference code calls its orthogonalization routine \texttt{zeropower\_via\_newtonschulz5}: for $G=U\Sigma V^\top$ the target $UV^\top$ is $U\Sigma^0V^\top$, the ``zeroth power'' of $G$ \parencite{Jordan2024}.}

\begin{practice}[Muon, \textcite{JordanEtAl2024Muon}, code \textcite{Jordan2024}; Moonlight, \textcite{LiuEtAl2025Muon}]
For each hidden weight matrix $W\in\R^{d_{\mathrm{out}}\times d_{\mathrm{in}}}$ (convolution filters flattened to matrices) with gradient $G_t$: set $M_t=\beta M_{t-1}+(1-\beta)G_t$ with $\beta=0.95$ and $N_t=(1-\beta)G_t+\beta M_t$ (Nesterov form); cast $N_t$ to bfloat16, divide it by its Frobenius norm, and apply five times $X\leftarrow aX+b(XX^\top)X+c(XX^\top)^2X$ with $(a,b,c)=(3.4445,-4.7750,2.0315)$; step $W\leftarrow W-\eta\sqrt{\max(1,d_{\mathrm{out}}/d_{\mathrm{in}})}\,X$ with $\eta=0.02$ by default. Train embeddings, output layer, gains, and biases with AdamW. At scale (Moonlight, 2025): add decoupled weight decay ($0.1$ in the released code) and use the factor $0.2\sqrt{\max(d_{\mathrm{out}},d_{\mathrm{in}})}$ instead, so that the update has the same root-mean-square size for every shape and shares AdamW's learning rate and weight decay.
\end{practice}

According to the documentation of the reference code, the coefficients were chosen to make the slope of the polynomial at zero large, increasing it even past the point where the iteration stops converging to one, so that the output is roughly $US'V^\top$, with the diagonal of $S'$ spread over $[0.5,1.5]$, instead of $UV^\top$; in the author's experience this does not hurt \parencite{Jordan2024}.

\begin{observation}[Muon at small and at large scale]\label{obs:spectral:records}
(i) On the NanoGPT speedrun, which trains a transformer of the size of GPT-2 small to a fixed validation loss as fast as possible, Muon improved the record by a factor of $1.35$; on a CIFAR-10 speedrun it cut the time to $94\%$ accuracy from $3.3$ to $2.6$ A100-seconds \parencite{JordanEtAl2024Muon,Jordan2024}. (ii) In Moonshot AI's scaling-law experiments Muon matched AdamW with about $52\%$ of the training FLOPs, and it trained Moonlight, a mixture of experts with 3B active and 16B total parameters, on 5.7T tokens \parencite{LiuEtAl2025Muon}; Kimi K2, with 1T total and 32B active parameters, was pre-trained on 15.5T tokens with a variant, MuonClip, and its report states that the run had no loss spike \parencite{KimiTeam2025}. (iii) Shampoo, which computes the same update when its memory is switched off (\cref{prop:spectral:shampoo}), won the external-tuning ruleset of the AlgoPerf competition \parencite{KasimbegEtAl2025}.
\end{observation}

These are records, company reports, and one competition, not controlled comparisons, and (ii) comes from the group that uses the method.\sidenote{The README of the reference code also lists a transformer trained to the performance of GPT-2 XL for \$175 of compute \parencite{Jordan2024}.}

\begin{observation}[Spread-out gradients, flat updates]\label{obs:spectral:shape}
The gradients of hidden weight matrices of transformers have a few dominant singular values and a long tail. Low-rank methods such as GaLore, which project each gradient onto a subspace of small dimension, rest on the first fact \parencite{ZhaoEtAl2024GaLore}. In NanoGPT-scale training, \textcite{DavisDrusvyatskiy2025}, in a preprint, find that the activations entering a layer keep a low stable rank throughout training, while the gradients keep a large ratio $\norm{G}_{\nuc}/\norm{G}_F$. Muon's update has all but its smallest singular values near one.
\end{observation}

The facts are compatible: singular values $(1,\eps,\dots,\eps)$ with $k$ copies of $\eps=k^{-1/2}$ give $\norm{G}_F^2/\norm{G}_{\op}^2=2$ but $\norm{G}_{\nuc}^2/\norm{G}_F^2=(1+\sqrt k)^2/2$, and setting every singular value to one gives the tail the weight of the head.

\begin{observation}[The advantage shrinks with scale; \textcite{WenEtAl2026Fantastic}]\label{obs:spectral:shrink}
In a comparison of ten optimizers on language models of $0.1$B to $1.2$B parameters, trained on one to eight times the Chinchilla-optimal amount of data with each optimizer carefully tuned, the matrix-based optimizers, Muon and SOAP among them, were the fastest; but their speedup over well-tuned AdamW fell from about $1.4\times$ at $0.1$B parameters to about $1.1\times$ at $1.2$B.
\end{observation}

This is one careful study on one data distribution, and no theorem we know explains it. Why orthogonalize an update, and why does a preconditioner produce the same matrix? How do five multiplications by a polynomial come close to an orthogonal matrix, and why do Muon's coefficients not converge? When does a step that discards the gradient's singular values beat the gradient itself?

\section{The idealization: steepest descent on one matrix}\label{sec:spectral:model}

Fix one weight matrix and freeze the rest of the network. The parameter is $W\in\R^{m\times n}$, with $m=d_{\mathrm{out}}$ and $n=d_{\mathrm{in}}$ for a layer $x\mapsto Wx$, and $\R^{m\times n}$ carries the inner product $\ip{A,B}=\tr(A^\top B)=\sum_{ij}A_{ij}B_{ij}$, so that the gradient $G=\nabla\Loss(W)$ is again an $m\times n$ matrix. The singular values $\sigma_1(W)\ge\sigma_2(W)\ge\dots\ge0$ define the \term{spectral norm}\index{spectral norm} $\norm{W}_{\op}=\sigma_1(W)$, the operator norm from $\ell_2$ to $\ell_2$; the Frobenius norm $\norm{W}_F=(\sum_i\sigma_i(W)^2)^{1/2}$, the Euclidean norm of the entries; and the \term{nuclear norm}\index{nuclear norm} $\norm{W}_{\nuc}=\sum_i\sigma_i(W)$. Measuring vectors by their root-mean-square size $\norm{x}_{\rms}=\norm{x}_2/\sqrt{\dim x}$ gives the \term{RMS-to-RMS norm}\index{RMS-to-RMS norm}
\[
  \norm{W}_{\rms\to\rms}=\sup_{x\ne0}\frac{\norm{Wx}_{\rms}}{\norm{x}_{\rms}}=\sqrt{\frac nm}\,\norm{W}_{\op},
\]
the factor by which the layer can enlarge the typical size of an input coordinate. As in \cref{ch:steepest}, normalized steepest descent is $W_{t+1}=W_t+\eta\lmo(G_t)$ with $\lmo(G)\in\argmin_{\norm{D}\le1}\ip{G,D}$, a step of length $\eta$ in the chosen norm, and unnormalized steepest descent is $W_{t+1}=W_t-\eta G_t^\sharp$ with $G^\sharp=-\norm{G}_*\lmo(G)$; the step size $\eta$ is the learning rate of the practice box.

\begin{assumption}[One matrix, spectral geometry]\label{ass:spectral:smooth}
The loss $\Loss\colon\R^{m\times n}\to\R$ is bounded below and differentiable, and it is $L_{\op}$-smooth with respect to the spectral norm: $\norm{\nabla\Loss(W)-\nabla\Loss(W')}_{\nuc}\le L_{\op}\norm{W-W'}_{\op}$ for all $W,W'$. When the Frobenius geometry is compared, $L_F$ is the constant with the Frobenius norm on both sides.
\end{assumption}

The nuclear norm on the left is the dual of the spectral norm (\cref{thm:spectral:lmo}).

\begin{dictionary}
a matrix $W\in\R^{m\times n}$ with the spectral norm & one hidden weight matrix (an attention projection, an MLP matrix, a flattened convolution)\\
the gradient $G$ & the Nesterov combination of momentum and gradient that Muon orthogonalizes\\
the polar factor $\polar(G)=UV^\top$ & the output of Newton--Schulz, exact only after many steps\\
a normalized step of length $\eta$ in $\norm\cdot_{\op}$ & Muon's learning rate, ``in units of spectral norm per update''\\
the RMS-to-RMS norm & how far a layer's outputs move per unit input; Muon's factor $\sqrt{\max(1,d_{\mathrm{out}}/d_{\mathrm{in}})}$\\
nuclear rank of $G$; stable rank of the inputs $X$ & how many directions carry the gradient; the effective dimension of the activations entering the layer
\end{dictionary}

\begin{fieldguide}{Muon}{\textcite{JordanEtAl2024Muon}}
\fgrow{update}{$M_t=\beta M_{t-1}+(1-\beta)G_t$, \ $N_t=(1-\beta)G_t+\beta M_t$, \ $W_{t+1}=W_t-\eta\,c_{\mathrm{shape}}\,\mathrm{NS}_5(N_t)$ with $c_{\mathrm{shape}}=\sqrt{\max(1,d_{\mathrm{out}}/d_{\mathrm{in}})}$}
\fgrow{norm}{spectral; RMS-to-RMS when $d_{\mathrm{out}}\ge d_{\mathrm{in}}$ (\cref{prop:spectral:rms})}
\fgrow{preconditioner}{none: orthogonalization is the duality map of the norm (\cref{thm:spectral:lmo})}
\fgrow{averages}{moving average of $G$, $\beta=0.95$, Nesterov form (\cref{ch:momentum})}
\fgrow{state}{one matrix per weight matrix}
\fgrow{in the wild}{$\mathrm{NS}_5$: five quintic Newton--Schulz steps in bfloat16; $\eta=0.02$; AdamW for embeddings, output layer, and gains; weight decay and update-size matching at scale}
\fgrow{caution}{orthogonalization is approximate: all but the smallest singular values in $[0.681,1.203]$ (\cref{prop:spectral:quintic})}
\end{fieldguide}

\begin{fieldguide}{Shampoo}{\textcite{GuptaKorenSinger2018}}
\fgrow{update}{$L_t=L_{t-1}+G_tG_t^\top$, \ $R_t=R_{t-1}+G_t^\top G_t$ (from $\eps I$), \ $W_{t+1}=W_t-\eta L_t^{-1/4}G_tR_t^{-1/4}$}
\fgrow{norm}{spectral when memoryless (\cref{prop:spectral:shampoo})}
\fgrow{preconditioner}{Kronecker-factored, from $L_t$ and $R_t$ (\cref{ch:kronecker})}
\fgrow{averages}{sums, or moving averages, of $GG^\top$ and $G^\top G$}
\fgrow{state}{an $m\times m$ and an $n\times n$ matrix per weight matrix}
\fgrow{in the wild}{the AlgoPerf external-tuning winner}
\end{fieldguide}

\begin{fieldguide}{Spectral descent}{\textcite{CarlsonEtAl2015a}}
\fgrow{update}{$W_{t+1}=W_t-\eta\norm{G_t}_{\nuc}\polar(G_t)$, unnormalized; a preconditioned variant (2015b)}
\fgrow{norm}{spectral (Schatten-$\infty$) on weight matrices}
\fgrow{in the wild}{restricted Boltzmann machines (2015a), feedforward and convolutional networks (2015b)}
\end{fieldguide}

\paragraph{What was thrown away.} The other layers are \cref{ch:modular}'s subject, minibatch noise and momentum \cref{ch:momentum}'s, and weight decay \cref{ch:decay}'s; here gradients are exact and memory is off. We keep two features often idealized away, because they are where Muon departs from exact steepest descent: the orthogonalization is a polynomial iteration with few steps, run in low precision. For the comparison with gradient descent we freeze the layer's inputs and use the square loss, the smallest model in which the two geometries can disagree about which step is better and both smoothness constants are explicit.

\paragraph{What one had to notice.} A weight matrix acts on vectors, and a training step must control how much the layer's outputs change: an operator norm of the update, not the Euclidean norm of its entries.\sidenote{The RMS-to-RMS version of this requirement, with updates of spectral norm proportional to $\sqrt{d_{\mathrm{out}}/d_{\mathrm{in}}}$, is the spectral condition for feature learning of \textcite{YangSimonBernstein2023}; the modular norm of \textcite{LargeEtAl2024} extends it to whole networks (\cref{ch:modular}).} Once the norm is chosen, duality forces the update: the steepest direction is the polar factor \parencite{CarlsonEtAl2015a,BernsteinNewhouse2024}. And the polar factor is an odd function of the matrix, so odd polynomials, which need only matrix products, can approximate it.

\section{The theory: the polar factor and how to compute it}\label{sec:spectral:theory}

Throughout, $G\in\R^{m\times n}$ is nonzero of rank $r$ with compact singular value decomposition $G=U_r\Sigma V_r^\top$: the columns $u_1,\dots,u_r$ of $U_r$ and $v_1,\dots,v_r$ of $V_r$ are orthonormal, and $\Sigma=\diag(\sigma_1,\dots,\sigma_r)$ with $\sigma_1\ge\dots\ge\sigma_r>0$. We complete $U_r$ and $V_r$ to orthogonal matrices $[U_r\ U_\perp]$ and $[V_r\ V_\perp]$. The \term{polar factor}\index{polar factor} is $\polar(G)=U_rV_r^\top$, and $\polar(0)=0$; it does not depend on the choice of singular vectors, because $\polar(G)=G(G^\top G)^{-1/2}$ with the inverse square root taken on the range (\cref{prop:spectral:shampoo}), and $G=\polar(G)(G^\top G)^{1/2}$ is the polar decomposition of \cref{thm:toolkit:polar}.

\subsection{The steepest direction}

\begin{theorem}[Spectral steepest descent]\label{thm:spectral:lmo}
Let $G\in\R^{m\times n}$ be nonzero, of rank $r$, with the notation above.
\begin{enumerate}
\item $\max\{\ip{G,D}:\norm{D}_{\op}\le1\}=\norm{G}_{\nuc}$. Hence the dual of the spectral norm is the nuclear norm, and the dual of the nuclear norm is the spectral norm.
\item The maximizers are exactly the matrices
\begin{equation}\label{eq:spectral:face}
  D=U_rV_r^\top+U_\perp ZV_\perp^\top,\qquad Z\in\R^{(m-r)\times(n-r)},\quad\norm{Z}_{\op}\le1 .
\end{equation}
The maximizer is unique, and then equal to $\polar(G)$, if and only if $r=\min(m,n)$.
\item Among the maximizers, $\polar(G)$ is the unique one of least Frobenius norm and the unique one of least rank. So $\lmo_{\op}(G)=-\polar(G)$ is a valid oracle, the only one when $G$ has full rank; with this choice $G^\sharp=\norm{G}_{\nuc}\polar(G)$, and normalized and unnormalized steepest descent in the spectral norm are $W_{t+1}=W_t-\eta\polar(G_t)$ and $W_{t+1}=W_t-\eta\norm{G_t}_{\nuc}\polar(G_t)$.
\end{enumerate}
\end{theorem}

\begin{proof}
The plan is to expand $\ip{G,D}$ along the singular vectors of $G$, bound each term by $\norm{D}_{\op}$, and read off when every bound is attained. Since $G=\sum_{i\le r}\sigma_iu_iv_i^\top$ and $\ip{uv^\top,D}=u^\top Dv$,
\begin{equation}\label{eq:spectral:expand}
  \ip{G,D}=\sum_{i=1}^r\sigma_i\,u_i^\top Dv_i ,\qquad u_i^\top Dv_i\le\norm{Dv_i}_2\le\norm{D}_{\op},
\end{equation}
by the Cauchy--Schwarz inequality.

(i) If $\norm{D}_{\op}\le1$, then every term of \eqref{eq:spectral:expand} is at most $\sigma_i$, so $\ip{G,D}\le\norm{G}_{\nuc}$; for $D=U_rV_r^\top$, which has $\norm{D}_{\op}=1$ and $Dv_i=u_i$, every term equals $\sigma_i$. The maximum over the unit ball is the dual norm, and the second claim holds because every norm is the dual of its dual (\cref{lem:toolkit:dualnorm}).

(ii) Let $\norm{D}_{\op}\le1$ attain the maximum. The terms are at most $\sigma_i>0$ and add up to $\sum_i\sigma_i$, so $u_i^\top Dv_i=1$ for each $i\le r$. Then $1=u_i^\top(Dv_i)\le\norm{Dv_i}_2\le1$, and equality in the Cauchy--Schwarz inequality forces $Dv_i=u_i$; the same argument applied to $D^\top$ gives $D^\top u_i=v_i$. In the bases $[U_r\ U_\perp]$ and $[V_r\ V_\perp]$ the blocks of $D$ are therefore $U_r^\top DV_r=I_r$, $U_\perp^\top DV_r=U_\perp^\top U_r=0$, $U_r^\top DV_\perp=(D^\top U_r)^\top V_\perp=V_r^\top V_\perp=0$, and $Z\defeq U_\perp^\top DV_\perp$, with $\norm{Z}_{\op}\le\norm{D}_{\op}\le1$: $D$ has the form \eqref{eq:spectral:face}. Conversely, a matrix of the form \eqref{eq:spectral:face} is $\diag(I_r,Z)$ in these bases, so its spectral norm is $\max\{1,\norm{Z}_{\op}\}=1$, and it satisfies $Dv_i=u_i$, so $\ip{G,D}=\norm{G}_{\nuc}$. The set is a single point exactly when $Z$ is an empty matrix, that is, when $m=r$ or $n=r$.

(iii) In the same bases $\norm{D}_F^2=r+\norm{Z}_F^2$ and $\rank D=r+\rank Z$, both smallest exactly when $Z=0$. The oracle is the negative of a maximizer. For the sharp operator, every $D$ satisfies $\ip{G,D}-\frac12\norm{D}_{\op}^2\le\norm{G}_{\nuc}\norm{D}_{\op}-\frac12\norm{D}_{\op}^2\le\frac12\norm{G}_{\nuc}^2$, with equality if and only if $\norm{D}_{\op}=\norm{G}_{\nuc}$ and $D/\norm{D}_{\op}$ is a maximizer; $D=\norm{G}_{\nuc}\polar(G)$ is one (\cref{lem:toolkit:conjugate}).
\end{proof}

An orthogonalized update keeps the singular vectors of the gradient and discards its singular values: every direction $u_iv_i^\top$ in which the loss decreases receives the same step, however small its $\sigma_i$. For a rank-deficient gradient, the rule for a linear layer trained on a batch of $B<\min(m,n)$ examples (each contributes a rank-one term $\delta x^\top$), the steepest directions form the face \eqref{eq:spectral:face}, and the polar factor is its most conservative point: it moves nothing outside the row and column spaces of $G$.\sidenote{The face \eqref{eq:spectral:face} is the subdifferential of the nuclear norm at $G$, as \cref{lem:toolkit:subdiff} predicts for every norm: the maximizers of $\ip{G,\cdot}$ over the unit ball of a norm are the subgradients of the dual norm at $G$.} And the guarantees of \cref{ch:steepest} measure the gradient in the nuclear norm.

\begin{proposition}[Memoryless Shampoo is orthogonalized descent]\label{prop:spectral:shampoo}
Take negative powers of positive semidefinite matrices on their range (positive eigenvalues are raised to the power, zero eigenvalues stay zero). For every nonzero $G$,
\[
  (GG^\top)^{-1/4}G(G^\top G)^{-1/4}=(GG^\top)^{-1/2}G=G(G^\top G)^{-1/2}=\polar(G),
\]
and for $\eps>0$, $(\eps I+GG^\top)^{-1/4}G(\eps I+G^\top G)^{-1/4}=U_r\diag\bigl(\sigma_i(\sigma_i^2+\eps)^{-1/2}\bigr)V_r^\top$.
\end{proposition}

\begin{proof}
Since $GG^\top=U_r\Sigma^2U_r^\top$ and $G^\top G=V_r\Sigma^2V_r^\top$, the powers are taken on $\Sigma$, and
\[
  (GG^\top)^{-1/4}G(G^\top G)^{-1/4}=U_r\Sigma^{-1/2}\,\Sigma\,\Sigma^{-1/2}V_r^\top=U_rV_r^\top ;
\]
the other two products are $U_r\Sigma^{-1}\Sigma V_r^\top$ and $U_r\Sigma\Sigma^{-1}V_r^\top$. With damping, $(\eps I+GG^\top)^{-1/4}$ is $U_r(\Sigma^2+\eps I)^{-1/4}U_r^\top+\eps^{-1/4}U_\perp U_\perp^\top$, whose second term vanishes against $G$ because $U_\perp^\top U_r=0$; the right factor acts in the same way.
\end{proof}

\term{Shampoo}\index{Shampoo} preconditions the gradient by fourth roots of the accumulated statistics $L_t=\eps I+\sum_{s\le t}G_sG_s^\top$ and $R_t=\eps I+\sum_{s\le t}G_s^\top G_s$ \parencite{GuptaKorenSinger2018}. Keeping only the current gradient and letting $\eps\to0$ turns its update into $-\eta\polar(G_t)$: memoryless Shampoo is Muon without momentum \parencite{BernsteinNewhouse2024}. With $\eps>0$ it is a soft polar factor, $\sigma\mapsto\sigma(\sigma^2+\eps)^{-1/2}$, which like Newton--Schulz with few steps does not lift the smallest singular values. Accumulation is \cref{ch:kronecker}'s subject.

Under \cref{ass:spectral:smooth}, normalized spectral descent with $G_t=\nabla\Loss(W_t)$ satisfies, by the descent lemma in the spectral norm (\cref{lem:toolkit:descent}), the identity $\ip{G,\polar(G)}=\norm{G}_{\nuc}$, and the bound $\norm{\polar(G)}_{\op}\le1$,
\begin{equation}\label{eq:spectral:rate}
  \Loss(W_{t+1})\le\Loss(W_t)-\eta\norm{G_t}_{\nuc}+\tfrac12L_{\op}\eta^2 ,
\end{equation}
and summing over $t<T$ gives $\frac1T\sum_{t<T}\norm{G_t}_{\nuc}\le(\Loss(W_0)-\Loss^\star)/(\eta T)+\frac12L_{\op}\eta$, the spectral case of \cref{thm:steepest:descent}; with stochastic gradients and momentum, \cref{ch:momentum}'s analysis applies in every norm, and the published analyses of Muon follow that pattern \parencite{LiHong2025,Kovalev2025,ShenEtAl2025,PethickEtAl2025}. Why the spectral norm, rather than the Frobenius norm or the $\ell_\infty$ norm of \cref{ch:sign}? Because it controls what a step does to the outputs.

\begin{proposition}[The RMS-to-RMS norm]\label{prop:spectral:rms}
Let $W\in\R^{m\times n}$ act on inputs $x\in\R^n$, and let $\Delta W\ne0$.
\begin{enumerate}
\item $\norm{\Delta W\,x}_{\rms}\le\norm{\Delta W}_{\rms\to\rms}\norm{x}_{\rms}$ with $\norm{\Delta W}_{\rms\to\rms}=\sqrt{n/m}\,\norm{\Delta W}_{\op}$, with equality if and only if $x$ is a right singular vector of $\Delta W$ for its largest singular value or $x=0$; for $\Delta W=\sigma uv^\top$, if and only if $x$ is parallel to $v$.
\item One oracle is $\lmo_{\rms\to\rms}(G)=-\sqrt{m/n}\,\polar(G)$, the only one when $G$ has full rank, so the normalized step $\Delta W=-\eta\sqrt{m/n}\,\polar(G)$ changes every output by at most $\eta\norm{x}_{\rms}$ in RMS, whatever the widths.
\item If $m\ge n$ and $\rank G=n$, then this step changes every output by exactly $\eta\norm{x}_{\rms}$.
\end{enumerate}
\end{proposition}

\begin{proof}
(i) Since $\norm{\Delta Wx}_2^2=x^\top\Delta W^\top\Delta W\,x\le\sigma_1(\Delta W)^2\norm{x}_2^2$, with equality exactly on the top eigenspace of $\Delta W^\top\Delta W$,
\[
  \norm{\Delta Wx}_{\rms}=\frac{\norm{\Delta Wx}_2}{\sqrt m}\le\frac{\norm{\Delta W}_{\op}\norm{x}_2}{\sqrt m}=\sqrt{\frac nm}\,\norm{\Delta W}_{\op}\norm{x}_{\rms}.
\]
(ii) The unit ball of $\norm\cdot_{\rms\to\rms}$ is the spectral ball of radius $\sqrt{m/n}$, so its oracle is $\sqrt{m/n}$ times that of \cref{thm:spectral:lmo}; the bound follows from (i). (iii) If $\rank G=n\le m$, then $V_r$ is square and orthogonal, so $\norm{\polar(G)x}_2=\norm{V_r^\top x}_2=\norm{x}_2$ and $\norm{\Delta Wx}_{\rms}=\eta\sqrt{m/n}\,\norm{x}_2/\sqrt m=\eta\norm{x}_{\rms}$.
\end{proof}

So Muon's shape factor converts the spectral norm into the RMS-to-RMS norm: for $d_{\mathrm{out}}\ge d_{\mathrm{in}}$ its step is the normalized RMS-to-RMS step of size $\eta$, and a full-rank update moves every output by exactly $\eta\norm{x}_{\rms}$ in RMS; for $d_{\mathrm{out}}<d_{\mathrm{in}}$ the code uses the factor one, a step of RMS-to-RMS norm $\eta\sqrt{d_{\mathrm{in}}/d_{\mathrm{out}}}$. Moonlight's factor normalizes another quantity: a full-rank polar factor has entrywise RMS $\max(m,n)^{-1/2}$, so $0.2\sqrt{\max(m,n)}\,\polar(G)$ has entrywise RMS $0.2$ for every shape, as Moonlight intends, and an RMS-to-RMS norm that grows like $\sqrt{d_{\mathrm{in}}}$ when $d_{\mathrm{out}}\ge d_{\mathrm{in}}$. Which rule keeps the best learning rate fixed as the width grows is \cref{ch:modular}'s question.\sidenote{The README of the reference code states that the learning rate ``should have built-in muP scaling'', so that it need not be retuned as the model grows \parencite{Jordan2024}.}

\subsection{Computing the polar factor}

A singular value decomposition is slow on accelerators built for matrix products. But the polar factor maps singular values to one and keeps singular vectors, and odd matrix polynomials act on singular values in the same way.

\begin{lemma}[Odd polynomials act on singular values]\label{lem:spectral:oddpoly}
For an odd polynomial $p(s)=\sum_{j=0}^Jc_js^{2j+1}$ and $X\in\R^{m\times n}$ put $p(X)=\sum_jc_j(XX^\top)^jX$. If $X=UDV^\top$ with $U^\top U=I$, $V^\top V=I$, and $D$ diagonal, then $p(X)=U\,p(D)\,V^\top$. Hence the iteration $X_{k+1}=p_{k+1}(X_k)$ from $X_0=G/\gamma$, $\gamma>0$, satisfies $X_k=U_r\diag\bigl(p_k\circ\dots\circ p_1(\sigma_i/\gamma)\bigr)V_r^\top$.
\end{lemma}

\begin{proof}
Because $V^\top V=I$, we have $XX^\top=UD^2U^\top$, and induction on $j$ with $U^\top U=I$ gives $(XX^\top)^jX=(UD^2U^\top)^jUDV^\top=UD^{2j+1}V^\top$; summing gives $p(X)=Up(D)V^\top$. The diagonal $p(D)$ may have negative or zero entries, but the identity needs only $U^\top U=V^\top V=I$, so induction on $k$ applies.
\end{proof}

The \term{Newton--Schulz iteration}\index{Newton--Schulz iteration} for the polar factor uses the cubic $p(s)=\frac32s-\frac12s^3$, the first-order correction toward $X^\top X=I$; Muon's quintic has a steeper start.\sidenote{\textcite{Schulz1933} gave the iteration $X\leftarrow X(2I-AX)$ for the inverse of $A$, which named the family; the cubic iteration for the orthogonal factor appears in \textcite{BjorckBowie1971}, and in machine learning it is also called Bj\"orck orthogonalization. \Textcite[Chapter 8]{Higham2008} treats iterations for the polar decomposition.} A quintic step costs three matrix products, $XX^\top$, its square, and the product with $X$, and Muon transposes tall matrices so that $XX^\top$ is the smaller Gram matrix.

\begin{theorem}[Newton--Schulz]\label{thm:spectral:ns}
Let $G\ne0$ have rank $r$, condition number $\kappa=\sigma_1/\sigma_r$, and stable rank $\spectralrst(G)=\norm{G}_F^2/\norm{G}_{\op}^2$. Let $p(s)=\frac32s-\frac12s^3$, $X_0=G/\gamma$, and $X_{k+1}=\frac32X_k-\frac12X_kX_k^\top X_k$.
\begin{enumerate}
\item If $\gamma>\sigma_1/\sqrt3$, then $X_k\to\polar(G)$.
\item If $\gamma\ge\sigma_1$, for instance $\gamma=\norm{G}_F$, then for each $s_i=\sigma_i/\gamma\in(0,1]$ the sequence $p^{\circ k}(s_i)$ is nondecreasing with limit $1$, and
\[
  \norm{X_k-\polar(G)}_{\op}=1-p^{\circ k}(s_r),\qquad 1-p(s)=\tfrac12(1-s)^2(2+s).
\]
\item Let $\gamma=\norm{G}_F$, $0<\eps\le\frac12$, and write $\lceil x\rceil_+=\max\{0,\lceil x\rceil\}$. Then $\norm{X_k-\polar(G)}_{\op}\le\eps$ for every
\[
  k\ \ge\ \Bigl\lceil\frac{\log\bigl(\kappa\sqrt{\spectralrst(G)}/2\bigr)}{\log(11/8)}\Bigr\rceil_+ +\Bigl\lceil\log_2\frac{\log(2/(3\eps))}{\log(4/3)}\Bigr\rceil .
\]
\end{enumerate}
\end{theorem}

\begin{proof}
The plan is to reduce to the scalar map with \cref{lem:spectral:oddpoly} and then follow $s\mapsto p(s)$: a monotone climb toward one, then quadratic convergence. By the lemma, $X_k-\polar(G)=U_r\diag(p^{\circ k}(s_i)-1)V_r^\top$, so $\norm{X_k-\polar(G)}_{\op}=\max_i\abs{p^{\circ k}(s_i)-1}$.

Since $p(s)=\frac12s(3-s^2)$ and $p'(s)=\frac32(1-s^2)$, the map increases on $[0,1]$ from $p(0)=0$ to $p(1)=1$ and decreases on $[1,\sqrt3]$ to $p(\sqrt3)=0$; so $p$ maps $(0,\sqrt3)$ into $(0,1]$. On $(0,1)$, $p(s)/s=\frac12(3-s^2)>1$, so $s<p(s)<1$; a sequence started in $(0,1]$ is nondecreasing and bounded by one, and its limit is a fixed point of $p$ in $(0,1]$, which must be $1$ because $p(s)=s$ forces $s^2=1$. A start in $(1,\sqrt3)$ is mapped into $(0,1)$ by one step. This proves (i). In (ii) all $s_i$ lie in $(0,1]$, and because $p$ is increasing there the sequence from $s_r$ stays the smallest; the factorization of $1-p$ is checked by expanding.

For (iii), $s_r=\sigma_r/\norm{G}_F=1/(\kappa\sqrt{\spectralrst(G)})$. While $s\le\frac12$, $p(s)/s\ge\frac{11}8$, and once $s\ge\frac12$, $p(s)\ge p(\frac12)=\frac{11}{16}$, so the sequence never returns below $\frac12$; hence $p^{\circ k_1}(s_r)\ge\frac12$ after the $k_1$ steps of the first term in (iii). From then on $e_j=1-p^{\circ(k_1+j)}(s_r)\in[0,\frac12]$ satisfies $e_{j+1}=\frac12e_j^2(3-e_j)\le\frac32e_j^2$ by (ii), so $f_j=\frac32e_j$ has $f_{j+1}\le f_j^2$ and $f_0\le\frac34$. Thus $e_j\le\frac23(\frac34)^{2^j}$, which is at most $\eps$ once $2^j\ge\log(2/(3\eps))/\log(4/3)$, a ratio of at least one as $\eps\le\frac12$.
\end{proof}

The Frobenius normalization is safe, since $\norm{G}_F\ge\norm{G}_{\op}$, and it costs only the factor $\sqrt{\spectralrst(G)}$, small for the low-stable-rank gradients of \cref{obs:spectral:shape}. The climb is what is slow: near zero the cubic multiplies a singular value by only $\frac32$ per step, so lifting $s=10^{-2}$ above $\frac12$ takes ten steps and reaching an error of $10^{-3}$ fourteen (the bound gives eighteen). Zero singular values stay zero, so in exact arithmetic the iteration returns the polar factor even for rank-deficient $G$.

\begin{proposition}[Muon's quintic does not converge]\label{prop:spectral:quintic}
Let $q(s)=as+bs^3+cs^5$ with $(a,b,c)=(3.4445,-4.7750,2.0315)$, with positive critical points $s_1=0.5545\ldots$ and $s_2=1.0501\ldots$.
\begin{enumerate}
\item $q>0$ on $(0,\infty)$; $q$ increases on $[0,s_1]$, decreases on $[s_1,s_2]$, and increases on $[s_2,\infty)$.
\item $J=[q(s_2),q(s_1)]=[0.6818\ldots,1.2023\ldots]$ satisfies $q(J)=J'\subseteq J$ with $J'=[0.6818\ldots,1.1344\ldots]$, and $q(J')=J'$.
\item If $s\in(0,1]$ and $q^{\circ k}(s)\ge q(s_2)$, then $q^{\circ j}(s)\in J$ for all $j\ge k$. After five steps every $s\ge1.6\cdot10^{-3}$ lies in $J$, and every $s\ge5.2\cdot10^{-3}$ lies in $J'$.
\item The fixed points of $q$ in $(0,\infty)$ are $s_a=0.8680\ldots$ and $s_b=1.2637\ldots$, with $q'(s_a)=-1.58\ldots$ and $q'(s_b)=6.47\ldots$. For no $s\in(0,1]$ outside a countable set do the iterates $q^{\circ k}(s)$ converge.
\end{enumerate}
\end{proposition}

\begin{proof}
The plan is to locate the critical and fixed points through quadratics in $u=s^2$ and to follow a lower envelope of the iterates by monotonicity; decimals are evaluations of explicit expressions. (i) $q(s)/s=a+bu+cu^2$ has discriminant $b^2-4ac=-5.19\ldots<0$, so $q>0$; $q'(s)=a+3bu+5cu^2$ is positive at $u=0$ and vanishes at $u=s_1^2=0.3075\ldots$ and $u=s_2^2=1.1028\ldots$.

(ii) and (iii). By (i), for $0<\ell\le s_1$ the values of $q$ on $[\ell,q(s_1)]$ lie in $[\min\{q(\ell),q(s_2)\},q(s_1)]$, since on its three monotone pieces they lie between $q(\ell)$ and $q(s_1)$, between $q(s_2)$ and $q(s_1)$, and between $q(s_2)$ and $q(q(s_1))=0.9465\ldots$; for $s_1<\ell\le q(s_1)$ only the last two pieces occur. As $q(s_2)>s_1$, this gives $q(J)\subseteq J$; the maximum of $q$ on $J$ is at its left end, where $q$ decreases, and equals $q(q(s_2))=1.1344\ldots$, so $q(J)=J'$, and the same computation gives $q(J')=J'$. Starting from $[s,1]\subseteq[s,q(s_1)]$, after $k$ steps every point of $[s,1]$ is in $[q^{\circ k}(s),q(s_1)]$ while $q^{\circ k}(s)\le s_1$, and in $J$ once $q^{\circ k}(s)\ge q(s_2)$. As $q$ increases on $[0,s_1]$, it suffices to check $q^{\circ5}(1.6\cdot10^{-3})=0.718\ldots$ and $q^{\circ4}(5.2\cdot10^{-3})=0.6836\ldots$, both above $q(s_2)$.

(iv) The fixed points solve $a-1+bu+cu^2=0$, with roots $u=0.7535\ldots$ and $u=1.5970\ldots$. Every orbit from $(0,1]$ enters $J$, because $q(s)/s$ decreases in $u$ on $(0,s_1]$, so that $q(s)\ge(q(s_1)/s_1)s>2s$ there, and an iterate in $(s_1,q(s_1)]$ is mapped into $J$. If $q^{\circ k}(s)\to s^*$, then $s^*$ is a fixed point by continuity and lies in the closed interval $J$, which excludes $s_b$; so $s^*=s_a$. As $\abs{q'(s_a)}>1$, there are $\delta>0$ and $\rho>1$ with $\abs{q(x)-s_a}\ge\rho\abs{x-s_a}$ for $\abs{x-s_a}<\delta$; iterates that stay within $\delta$ of $s_a$ without reaching it move away geometrically, so some iterate equals $s_a$. Then $s$ lies in $\bigcup_kq^{-k}(\{s_a\})$, which is countable because a point has at most five preimages under $q$.
\end{proof}

The quintic makes a deliberate trade. Its slope at zero is $3.4445$ instead of $\frac32$, so $s=10^{-2}$ passes $\frac12$ after four steps instead of ten, and five steps bring every normalized singular value above $1.6\cdot10^{-3}$ into $[0.681,1.203]$. The price is that $q(1)=0.701$ and the fixed points repel: the singular values never settle but move about inside $J$, of which the documentation's $[0.5,1.5]$ is a loose description. The step along different singular directions then varies by a factor of up to $q(s_1)/q(s_2)=1.76$, which costs at most a factor $(q(s_2)/q(s_1))^2=0.32$ in the guaranteed decrease \eqref{eq:spectral:rate} (\cref{ex:spectral:approx}).

\section{The theory: when a spectral step beats a gradient step}\label{sec:spectral:whenhelps}

Consider one linear layer $W\in\R^{d_{\mathrm{out}}\times d_{\mathrm{in}}}$ trained on $n$ examples, with inputs $X\in\R^{d_{\mathrm{in}}\times n}$ and targets $Y\in\R^{d_{\mathrm{out}}\times n}$, by
\[
  \Loss(W)=\frac1{2n}\norm{WX-Y}_F^2,\qquad A=\frac1nXX^\top\ne0 ,
\]
and let $\lambda_1\ge\dots\ge\lambda_{d_{\mathrm{in}}}\ge0$ be the eigenvalues of $A$, with orthonormal eigenvectors $q_1,\dots,q_{d_{\mathrm{in}}}$.

\begin{definition}[Stable rank and nuclear rank]\label{def:spectral:ranks}
For a nonzero matrix $M$, the \term{stable rank}\index{stable rank} is $\spectralrst(M)=\norm{M}_F^2/\norm{M}_{\op}^2$ and the \term{nuclear rank}\index{nuclear rank} is $\spectralrnuc(M)=\norm{M}_{\nuc}^2/\norm{M}_F^2$. Both lie in $[1,\rank M]$, and each equals $\rank M$ exactly when the nonzero singular values of $M$ are equal.
\end{definition}

The bounds follow from $\sigma_1^2\le\sum_i\sigma_i^2\le(\rank M)\sigma_1^2$ and, by the Cauchy--Schwarz inequality, $\sum_i\sigma_i^2\le(\sum_i\sigma_i)^2\le(\rank M)\sum_i\sigma_i^2$. For the inputs, $\spectralrst(X)=\tr A/\lambda_1$.

\begin{lemma}[The two smoothness constants of a linear layer]\label{lem:spectral:layer}
\begin{enumerate}
\item $\Loss(W+\Delta)=\Loss(W)+\ip{\nabla\Loss(W),\Delta}+\frac12\tr(\Delta A\Delta^\top)$, with $\nabla\Loss(W)=\frac1n(WX-Y)X^\top$.
\item In every norm, the smallest smoothness constant of $\Loss$ is $\sup\{\tr(\Delta A\Delta^\top):\norm{\Delta}\le1\}$.
\item $L_F=\lambda_1$ and $L_{\op}=\sum_{i\le k}\lambda_i$ with $k=\min(d_{\mathrm{out}},d_{\mathrm{in}})$; so $L_{\op}=\tr A$ exactly when $d_{\mathrm{out}}\ge\rank A$, for instance when $d_{\mathrm{out}}\ge d_{\mathrm{in}}$.
\end{enumerate}
\end{lemma}

\begin{proof}
The plan is to read everything off the exact quadratic (i), which follows by expanding $\norm{(W+\Delta)X-Y}_F^2$ with $\norm{\Delta X}_F^2=\tr(\Delta XX^\top\Delta^\top)$. (ii) By (i), $\nabla\Loss(W+\Delta)-\nabla\Loss(W)=\Delta A$, and $\ip{\Delta,\Delta'}_A\defeq\ip{\Delta A,\Delta'}=\tr(A\Delta^\top\Delta')$ is a symmetric positive semidefinite bilinear form. With $L=\sup_{\norm{E}\le1}\ip{E,E}_A$, the Cauchy--Schwarz inequality for this form gives $\norm{\Delta A}_*=\sup_{\norm{\Delta'}\le1}\ip{\Delta,\Delta'}_A\le L\norm{\Delta}$, and $\norm{\Delta A}_*\ge\ip{\Delta,\Delta}_A/\norm{\Delta}$ for $\Delta\ne0$ shows that no smaller constant works. (iii) With rows $\delta_j$ of $\Delta$, $\tr(\Delta A\Delta^\top)=\sum_j\delta_j^\top A\delta_j\le\lambda_1\norm{\Delta}_F^2$, with equality at $\Delta=e_1q_1^\top$. For the spectral norm put $P=\Delta^\top\Delta$ and $w_i=q_i^\top Pq_i$, so that $\tr(\Delta A\Delta^\top)=\sum_i\lambda_iw_i$. If $\norm{\Delta}_{\op}\le1$, then $0\preceq P\preceq I$ and $\rank P\le k$, so $w_i\in[0,1]$ and $\sum_iw_i=\tr P\le k$, and the sum is at most $\sum_{i\le k}\lambda_i$; $\Delta=\sum_{i\le k}e_iq_i^\top$ attains it and has spectral norm one, as $k\le d_{\mathrm{out}}$. It equals $\tr A$ exactly when $\rank A\le k$, that is, when $\rank A\le d_{\mathrm{out}}$.
\end{proof}

\begin{theorem}[When spectral steps beat Frobenius steps]\label{thm:spectral:whenhelps}
Let $G=\nabla\Loss(W)\ne0$, and compare the spectral step $-\polar(G)\norm{G}_{\nuc}/L_{\op}$ with the gradient step $-G/L_F$.
\begin{enumerate}
\item The spectral step decreases $\Loss$ by at least $\norm{G}_{\nuc}^2/(2L_{\op})$ and the gradient step by at least $\norm{G}_F^2/(2L_F)$. Each step minimizes the upper bound of the descent lemma in its norm, and so does the normalized step of the best length.
\item The spectral guarantee is the larger if and only if $\spectralrnuc(G)>L_{\op}/L_F$. When $d_{\mathrm{out}}\ge d_{\mathrm{in}}$ this reads
\[
  \frac{\norm{G}_{\nuc}^2}{\norm{G}_F^2}>\frac{\norm{X}_F^2}{\norm{X}_{\op}^2},\qquad\text{that is,}\qquad\spectralrnuc(G)>\spectralrst(X):
\]
the nuclear rank of the gradient must exceed the stable rank of the inputs.
\item The actual decreases are $(2-\rho_{\op})$ and $(2-\rho_F)$ times the guarantees, with $\rho_{\op}=\tr(V_r^\top AV_r)/L_{\op}$ and $\rho_F=\tr(GAG^\top)/(L_F\norm{G}_F^2)$ in $(0,1]$. If $d_{\mathrm{out}}\ge d_{\mathrm{in}}=\rank G$, then $\rho_{\op}=1$: the spectral guarantee is attained, and the ratio of the actual decreases lies between $\frac12$ and $1$ times the ratio of the guarantees. The gradient guarantee is attained if and only if every row of $G$ lies in the top eigenspace of $A$.
\end{enumerate}
\end{theorem}

\begin{proof}
The plan is to evaluate the exact quadratic of \cref{lem:spectral:layer}(i) along each step. (i) For a norm with constant $L$, \cref{lem:spectral:layer}(ii) gives $\Loss(W+\Delta)\le\Loss(W)+\ip{G,\Delta}+\frac L2\norm{\Delta}^2$. At fixed length $\norm{\Delta}=\tau$ the smallest $\ip{G,\Delta}$ is $-\tau\norm{G}_*$, at $\Delta=\tau\lmo(G)$, and $-\tau\norm{G}_*+\frac L2\tau^2$ is smallest at $\tau=\norm{G}_*/L$, with value $-\norm{G}_*^2/(2L)$; by \cref{thm:spectral:lmo} this step is $-\polar(G)\norm{G}_{\nuc}/L_{\op}$ in the spectral norm and $-G/L_F$ in the Frobenius norm. (ii) The comparison is $\spectralrnuc(G)>L_{\op}/L_F$, and \cref{lem:spectral:layer}(iii) gives $L_{\op}/L_F=\tr A/\lambda_1=\spectralrst(X)$ when $d_{\mathrm{out}}\ge d_{\mathrm{in}}$.

(iii) Along $-\tau\polar(G)$, $\Loss$ changes by $-\tau\norm{G}_{\nuc}+\frac{\tau^2}2\tr(U_rV_r^\top AV_rU_r^\top)$, and the trace equals $\tr(V_r^\top AV_r)=\rho_{\op}L_{\op}$; at $\tau=\norm{G}_{\nuc}/L_{\op}$ the decrease is $(2-\rho_{\op})\norm{G}_{\nuc}^2/(2L_{\op})$. The same computation along $-G/L_F$ gives the factor $2-\rho_F$. The rows of $G=\frac1n(WX-Y)X^\top$, and hence the columns of $V_r$, lie in the range of $A$, where $A$ is positive definite, so $\rho_F>0$ and $\rho_{\op}>0$. Next, $\rho_F\le1$ because $\tr(GAG^\top)\le\lambda_1\norm{G}_F^2$, with equality exactly when every row of $G$ is an eigenvector for $\lambda_1$; and $\rho_{\op}\le1$ by the argument of \cref{lem:spectral:layer}(iii) applied to $P=V_rV_r^\top$, which has rank $r\le k$. If $\rank G=d_{\mathrm{in}}\le d_{\mathrm{out}}$, then $V_r$ is square and orthogonal and $\tr(V_r^\top AV_r)=\tr A=L_{\op}$; the ratio of actual decreases is then the ratio of guarantees divided by $2-\rho_F\in[1,2)$.
\end{proof}

A gradient step is limited by the stiffest input direction, through $\lambda_1$, and paid in the Frobenius norm of the gradient; a spectral step is limited by the total input energy $\tr A$ and paid in the nuclear norm. The spectral step wins when the gradient is spread over more directions than the inputs effectively occupy, which needs anisotropic inputs: $\spectralrnuc(G)\le\rank A$, so with $A=\lambda_1I$ the stable rank $d_{\mathrm{in}}$ is never beaten, and gradient descent with step $1/\lambda_1$ solves the problem in one step. This is the comparison of \textcite{DavisDrusvyatskiy2025}.\sidenote{They also prove that post-activation matrices have low stable rank at Gaussian initialization, in random-feature regression, feedforward networks, and transformer blocks \parencite{DavisDrusvyatskiy2025}.} One family makes every quantity explicit.

\begin{proposition}[The aligned case is sign descent on the spectrum]\label{prop:spectral:aligned}
Let $d_{\mathrm{out}}\ge d_{\mathrm{in}}$, $Y=W_\star X$, and $A=Q\diag(a)Q^\top$ with $Q$ orthogonal and $a_1\ge\dots\ge a_{d_{\mathrm{in}}}>0$, so that $\Loss(W)=\frac12\tr\bigl((W-W_\star)A(W-W_\star)^\top\bigr)$, and let $W_0-W_\star=U\diag(e)Q^\top$ with $U^\top U=I$. Then every iterate of gradient descent and of spectral descent has the form $W_t-W_\star=U\diag(e^{(t)})Q^\top$, and
\begin{enumerate}
\item gradient descent with step $\eta$ is $e^{(t+1)}_i=(1-\eta a_i)e^{(t)}_i$, gradient descent on $\frac12\sum_ia_ie_i^2$;
\item spectral descent with steps $\eta_t$ is $e^{(t+1)}=e^{(t)}-\eta_t\sign(e^{(t)})$, sign descent on the same quadratic (\cref{ch:sign});
\item with $g_i=a_ie_i$, $\norm{G}_{\nuc}=\norm{g}_1$, $\norm{G}_F=\norm{g}_2$, $L_{\op}=\sum_ia_i$, and $L_F=a_1$, and the criterion of \cref{thm:spectral:whenhelps} becomes $\norm{g}_1^2/\norm{g}_2^2>\sum_ia_i/a_1$, the criterion of \cref{ch:steepest} for preferring $\ell_\infty$ steps to $\ell_2$ steps on a diagonal quadratic.
\end{enumerate}
\end{proposition}

\begin{proof}
If $W-W_\star=U\diag(e)Q^\top$, then $G=(W-W_\star)A=U\diag(a_ie_i)Q^\top$, and $(W-W_\star)(I-\eta A)=U\diag\bigl(e_i(1-\eta a_i)\bigr)Q^\top$. As $U$ and $Q$ have orthonormal columns, the singular values of $G$ are the $\abs{a_ie_i}$ and $\polar(G)=U\diag(\sign(e_i))Q^\top$, with zeros where $e_i=0$. Part (iii) reads the norms off the singular values and uses \cref{lem:spectral:layer}(iii).
\end{proof}

So the spectral comparison contains the $\ell_\infty$ comparison of \cref{ch:steepest,ch:sign} as its diagonal case, and there Muon is sign descent on the singular values of the error.\index{aligned case} Outside the aligned case the singular vectors of the gradient move from step to step and the reduction fails (\cref{sec:spectral:misses}).

\section{Reading the theory}\label{sec:spectral:reading}

\paragraph{Newton--Schulz, step by step.} \Cref{fig:spectral:ns} draws the maps $p_k\circ\dots\circ p_1$ of \cref{lem:spectral:oddpoly} for $k\le5$, with singular values of matrices iterated in bfloat16 as in Muon's code. The cubic's curves rise like $(\frac32)^ks$ and flatten at one, where the error squares at every step; after five steps only singular values above $0.27$ are within $10^{-2}$ of one. Muon's curves rise like $3.4445^ks$, overshoot to $1.2$, and fold back, filling the band $J$ for every $s$ above $1.6\cdot10^{-3}$; below that the output stays small, so Muon leaves the weakest singular directions underweighted rather than equalizing them. The third iteration, \term{Polar Express}\index{Polar Express}, uses a different quintic at each step: \textcite{AmselEtAl2026} choose at step $k$ the odd quintic closest to one in the worst case over the interval $[\ell_k,u_k]$ that contains the singular values, prove this greedy choice optimal among compositions, and in their code start from $\ell_0=10^{-3}$ with a safety factor against rounding. Its curves swing between nearly zero and nearly two and then contract: after five steps the singular values with $s\ge10^{-3}$ lie in $[0.857,1.131]$, after six in $[0.9967,1.0016]$. For $k\ge3$ the compositions oscillate faster than the plot resolves and fill their bands.

\begin{figure}[t]
\begin{fullwidth}
\centering
\begin{tikzpicture}
\begin{groupplot}[group style={group size=3 by 1,horizontal sep=8mm},wplot,width=62mm,height=47mm,
  xmode=log,xmin=1e-4,xmax=1,xtick={1e-4,1e-3,1e-2,1e-1,1},
  xticklabels={$10^{-4}$,$10^{-3}$,$10^{-2}$,$10^{-1}$,$1$},xlabel={$s=\sigma/\norm{G}_F$},
  title style={yshift=-1ex}]
\nextgroupplot[title={cubic Newton--Schulz},ylabel={after $k$ steps},ymin=0,ymax=1.32,ytick={0,0.5,1}]
\addplot[guide] coordinates {(1e-4,1) (1,1)};
\addplot[sea!22,line width=0.9pt] table[x=s,y=cub1] {figures/data/03-spectral-ns-curves.dat};
\addplot[sea!40,line width=0.9pt] table[x=s,y=cub2] {figures/data/03-spectral-ns-curves.dat};
\addplot[sea!60,line width=0.9pt] table[x=s,y=cub3] {figures/data/03-spectral-ns-curves.dat};
\addplot[sea!80,line width=0.9pt] table[x=s,y=cub4] {figures/data/03-spectral-ns-curves.dat};
\addplot[seadeep,line width=1.1pt] table[x=s,y=cub5] {figures/data/03-spectral-ns-curves.dat};
\addplot[simulation2,mark size=1.2pt] table[x=s,y=cub1] {figures/data/03-spectral-ns-points.dat};
\addplot[simulation,mark size=1.4pt] table[x=s,y=cub5] {figures/data/03-spectral-ns-points.dat};
\node[font=\scriptsize,text=seadeep,anchor=south east] at (axis cs:0.06,0.62) {$k=5$};
\node[font=\scriptsize,text=sea!70,anchor=west] at (axis cs:0.12,0.12) {$k=1$};
\nextgroupplot[title={Muon's quintic},ymin=0,ymax=1.32,ytick={0,0.5,1}]
\fill[seamist] (axis cs:1e-4,0.6818) rectangle (axis cs:1,1.2024);
\addplot[guide] coordinates {(1e-4,1) (1,1)};
\addplot[sea!22,line width=0.9pt] table[x=s,y=muon1] {figures/data/03-spectral-ns-curves.dat};
\addplot[sea!40,line width=0.9pt] table[x=s,y=muon2] {figures/data/03-spectral-ns-curves.dat};
\addplot[sea!60,line width=0.9pt] table[x=s,y=muon3] {figures/data/03-spectral-ns-curves.dat};
\addplot[sea!80,line width=0.9pt] table[x=s,y=muon4] {figures/data/03-spectral-ns-curves.dat};
\addplot[seadeep,line width=1.1pt] table[x=s,y=muon5] {figures/data/03-spectral-ns-curves.dat};
\addplot[simulation2,mark size=1.2pt] table[x=s,y=muon1] {figures/data/03-spectral-ns-points.dat};
\addplot[simulation,mark size=1.4pt] table[x=s,y=muon5] {figures/data/03-spectral-ns-points.dat};
\node[font=\scriptsize,text=ink,anchor=north west] at (axis cs:1.2e-4,1.3) {band $J$};
\nextgroupplot[title={Polar Express},ymin=0,ymax=2.05,ytick={0,0.5,1,1.5,2}]
\addplot[guide] coordinates {(1e-4,1) (1,1)};
\addplot[sea!22,line width=0.9pt] table[x=s,y=pe1] {figures/data/03-spectral-ns-curves.dat};
\addplot[sea!40,line width=0.9pt] table[x=s,y=pe2] {figures/data/03-spectral-ns-curves.dat};
\addplot[sea!60,line width=0.9pt] table[x=s,y=pe3] {figures/data/03-spectral-ns-curves.dat};
\addplot[sea!80,line width=0.9pt] table[x=s,y=pe4] {figures/data/03-spectral-ns-curves.dat};
\addplot[seadeep,line width=1.1pt] table[x=s,y=pe5] {figures/data/03-spectral-ns-curves.dat};
\addplot[simulation2,mark size=1.2pt] table[x=s,y=pe1] {figures/data/03-spectral-ns-points.dat};
\addplot[simulation,mark size=1.4pt] table[x=s,y=pe5] {figures/data/03-spectral-ns-points.dat};
\end{groupplot}
\end{tikzpicture}
\end{fullwidth}
\caption{What $k$ steps of three odd-polynomial iterations do to a singular value $s$ of the Frobenius-normalized input. Curves: the compositions $p_k\circ\dots\circ p_1(s)$ for $k=1,\dots,5$, light to dark (for Polar Express, of $s/1.01$, as in its code; coefficients by the recipe of its code). Points: $u_i^\top X_kv_i$ for the singular pairs $(u_i,v_i)$ of six $64\times128$ test matrices with prescribed singular values and random singular vectors (seed 3), iterated in bfloat16 exactly as in Muon's reference code; squares $k=1$, circles $k=5$. In float64 the same iterations agree with the curves to $3\cdot10^{-5}$. Shaded: the band $J=[0.681,1.203]$ of \cref{prop:spectral:quintic}. Script: \texttt{03-spectral-ns.py}.}
\label{fig:spectral:ns}
\end{figure}

\paragraph{Where the lemma stops being exact.} In float64 the points lie on the curves, up to the constant $10^{-7}$ added to the Frobenius norm. In bfloat16, which keeps eight significant bits,\sidenote{bfloat16 keeps the eight exponent bits of float32 and only seven explicit mantissa bits, so its relative rounding error is at most $2^{-8}$; matrix products typically accumulate in float32 and round the result.} the deviations grow with the slopes of the composed maps: after five steps they stay below $1.6\cdot10^{-3}$ for the cubic and $3.1\cdot10^{-2}$ for Muon's quintic, but reach about $0.25$ for Polar Express, whose early polynomials are steep near the top of the spectrum and multiply the rounding error of about $2^{-8}$ by their slopes. Below the rounding level, rounding seeds every singular direction, and five quintic steps multiply small singular values by up to $3.4445^5\approx485$: an exactly rank-one $256\times512$ input leaves Muon's iteration in bfloat16 with about $55$ singular values above $0.1$, the largest near $0.4$ (\cref{ex:spectral:bf16}). They lie almost entirely in the block $U_\perp^\top X_5V_\perp$, with norm below one, so rounding moves the output along the face \eqref{eq:spectral:face}: in bfloat16 Newton--Schulz picks another steepest direction.

\paragraph{The crossover.} \Cref{fig:spectral:crossover} tests \cref{thm:spectral:whenhelps} on the aligned problems of \cref{prop:spectral:aligned}, with $d_{\mathrm{in}}=d_{\mathrm{out}}=128$, input spectrum $a_i=i^{-1/2}$ of stable rank $21.2$, and initial errors $e_i=i^{-\alpha}$. As $\alpha$ grows the error concentrates on the stiff input directions, the gradient $g_i=i^{-1/2-\alpha}$ loses nuclear rank, and gradient descent, fast on stiff directions, gains. At $\alpha=0$ one spectral step removes the error, while gradient descent needs $17$ steps to reach $1\%$ of the initial loss. The theorem predicts the crossover where the nuclear rank of the first gradient equals $21.2$, at $\alpha=0.450$; the one-step decreases cross at $\alpha=0.386$, and the step counts at $\alpha=0.388$. The prediction errs in the direction \cref{thm:spectral:whenhelps}(iii) predicts: the spectral guarantee is attained, while gradient descent beats its own by the factor $2-\rho_F$, which falls from $1.55$ at $\alpha=0$ to $1.03$ at $\alpha=1.5$. The theorem's scope ends at one step, and this is a case in which the singular vectors of the gradient never move; \cref{sec:spectral:misses} shows what happens when they do.

\begin{figure}[t]
\begin{fullwidth}
\centering
\begin{tikzpicture}
\begin{groupplot}[group style={group size=2 by 1,horizontal sep=17mm},wplot,width=80mm,height=47mm,xmin=0,xmax=1.5,
  xlabel={decay $\alpha$ of the initial error, $e_i=i^{-\alpha}$}]
\nextgroupplot[ylabel={one-step decrease $/\,\Loss(W_0)$},ymin=0,ymax=1.05]
\addplot[guide] coordinates {(0.450,0) (0.450,1.05)};
\addplot[theory] table[x=alpha,y=onestep_S] {figures/data/03-spectral-crossover-theory.dat};
\addplot[theory3] table[x=alpha,y=onestep_F] {figures/data/03-spectral-crossover-theory.dat};
\addplot[theory2] table[x=alpha,y=guar_F] {figures/data/03-spectral-crossover-theory.dat};
\addplot[simulation2] table[x=alpha,y=sim_onestep_S] {figures/data/03-spectral-crossover-sim.dat};
\addplot[simulation] table[x=alpha,y=sim_onestep_F] {figures/data/03-spectral-crossover-sim.dat};
\node[font=\scriptsize,text=sea,align=right,anchor=south east] at (axis cs:1.48,0.27) {spectral\\(exact $=$ guarantee)};
\node[font=\scriptsize,text=sandline!55!black,anchor=south east] at (axis cs:0.3,0.68) {gradient};
\node[font=\scriptsize,text=seadeep,align=left,anchor=north west] at (axis cs:0.02,0.2) {gradient\\guarantee};
\nextgroupplot[ylabel={steps to reach $10^{-2}\Loss(W_0)$},ymin=0,ymax=25]
\addplot[guide] coordinates {(0.450,0) (0.450,25)};
\addplot[theory,const plot mark mid] table[x=alpha,y=iters_S] {figures/data/03-spectral-crossover-theory.dat};
\addplot[theory3,const plot mark mid] table[x=alpha,y=iters_F] {figures/data/03-spectral-crossover-theory.dat};
\addplot[simulation2] table[x=alpha,y=sim_iters_S] {figures/data/03-spectral-crossover-sim.dat};
\addplot[simulation] table[x=alpha,y=sim_iters_F] {figures/data/03-spectral-crossover-sim.dat};
\node[font=\scriptsize,text=sea,anchor=south east] at (axis cs:0.8,20.6) {spectral};
\node[font=\scriptsize,text=sandline!55!black,anchor=west] at (axis cs:0.95,6.5) {gradient};
\end{groupplot}
\end{tikzpicture}
\end{fullwidth}
\caption{Spectral versus gradient steps on matrix least squares in the aligned case: $128\times128$ matrices, inputs with spectrum $a_i=i^{-1/2}$ (stable rank $21.2$), initial error $U\diag(i^{-\alpha})Q^\top$ with Haar-random $U$ and eigenbasis $Q$ (seed 11), both methods with the step sizes of \cref{thm:spectral:whenhelps}. Left: decrease of the loss in the first step, relative to the initial loss. Right: steps to reach $1\%$ of the initial loss. Curves: the scalar recursions of \cref{prop:spectral:aligned}; points (squares spectral, circles gradient): the matrix iterations, with polar factors from a singular value decomposition, which agree with the curves to $6\cdot10^{-16}$ and in every step count. Dashed line: the crossover $\alpha=0.450$ predicted by the theorem; the observed crossovers are at $\alpha=0.386$ (left) and $0.388$ (right). Script: \texttt{03-spectral-crossover.py}.}
\label{fig:spectral:crossover}
\end{figure}

\section{What the theory misses}\label{sec:spectral:misses}

\paragraph{What was idealized away.} A network has many weight matrices that share a step size and feed into one another (\cref{ch:modular}). Gradients here are exact and memory is off, while Muon averages minibatch gradients with momentum $0.95$, without which a normalized step stalls at a noise floor (\cref{ch:momentum}). Embeddings, the output layer, and gains are trained by AdamW for reasons the spectral norm does not supply (\cref{ch:modular}). The theorems about the oracle assume exact orthogonalization; \cref{prop:spectral:quintic} and the bfloat16 measurements quantify the gap for one step only. Weight decay is absent (\cref{ch:decay}). The comparison with gradient descent freezes the layer's inputs, uses the square loss, and compares worst-case guarantees for one step.

\paragraph{Where it breaks.} (1) \emph{One step does not predict many.} In the aligned case of \cref{fig:spectral:crossover} the theorem nearly does, placing the crossover at $\alpha=0.45$ instead of $0.39$; for generic targets it does not. With the same inputs and target singular values $i^{-\alpha}$, but right singular vectors drawn independently of the eigenvectors of the inputs, the ratio of guarantees stays above $2$ for every $\alpha\le0.5$, so the first spectral step always beats the first gradient step; yet in our runs (five seeds) spectral descent needed more steps than gradient descent to reach $1\%$ of the initial loss once $\alpha\ge0.4$: $36$ against $17$ at $\alpha=0.4$ (\cref{ex:spectral:generic}). A spectral step moves every singular direction by the same amount, overshoots the small ones, and must correct them later, as sign descent does (\cref{ch:sign}), and a first-step guarantee cannot see this. In the analysis of \textcite{ShenEtAl2025}, Muon matches the rate of gradient descent in a star-convex setting and can beat it when the Hessians are of low rank and close to block diagonal. (2) \emph{The advantage shrinks with scale} (\cref{obs:spectral:shrink}), and a theory of one matrix and one step has no scale in it. (3) \emph{Muon at scale needs more than the norm.} Moonshot AI found weight decay crucial for scaling Muon \parencite{LiuEtAl2025Muon}, and for Kimi K2 Moonshot AI report that attention logits exploded under Muon, which MuonClip prevents by rescaling the query and key weights after an update whenever the largest attention logit exceeds a threshold \parencite{KimiTeam2025}. A normalized step bounds the change of a matrix per step, not its accumulated size: after $t$ steps its spectral norm may have grown by as much as $\eta t$, and the bound on a logit bilinear in the query and key weights like $(\eta t)^2$ (a heuristic reading, not an analysis of those runs). (4) \emph{How exact must the orthogonalization be?} The reference code states that the approximation does not hurt \parencite{Jordan2024}, while \textcite{AmselEtAl2026} report consistent gains in validation loss for GPT-2-sized models from more accurate polynomials; one-step theory cannot settle this.

\paragraph{Repairs and open problems.} Better polynomials address the accuracy of the approximation \parencite{AmselEtAl2026}, and Gram Newton--Schulz, which runs the same iteration on the Gram matrix $XX^\top$, its cost \parencite{ZhangEtAl2026Gram}. Norm control addresses the growth of the weights: with decoupled weight decay Muon solves a problem constrained to a spectral-norm ball \parencite{ChenLiLiu2026,PethickEtAl2025} (\cref{ch:decay}), and Hyperball fixes the Frobenius norms of the weights and the updates, with token-equivalent speedups of $20$--$30\%$ over weight-decay baselines for Muon up to $1.2$B parameters, in a preprint \parencite{WenEtAl2026Hyperball}. Preconditioners with memory such as SOAP \parencite{VyasEtAl2025} are \cref{ch:kronecker}'s subject, and convergence with momentum and noise is \cref{ch:momentum}'s \parencite{LiHong2025,Kovalev2025}. Open: a multistep comparison of spectral and Euclidean descent outside the aligned case; why the gradients of trained networks keep a high nuclear rank while their activations keep a low stable rank, which \textcite{DavisDrusvyatskiy2025} prove for spiked random-feature models after a burn-in; why the advantage of matrix optimizers shrinks with scale; and how accurate, and in what precision, the orthogonalization must be.

\section{Back in the wild}\label{sec:spectral:back}

\fieldnotes{Muon orthogonalizes the momentum of each hidden matrix with five quintic Newton--Schulz steps in bfloat16; memoryless Shampoo computes the same matrix.}
{Steepest descent on one matrix under the spectral (RMS-to-RMS) norm; a linear layer with frozen inputs for the comparison with gradient descent.}
{The steepest direction is $\polar(G)$, with a whole face of them for rank-deficient $G$; the cubic converges, quadratically after a climb of $O(\log(\kappa\sqrt{\spectralrst(G)}))$ steps; Muon's quintic does not converge and keeps all but the smallest singular values in $[0.681,1.203]$; spectral beats gradient guarantees iff $\spectralrnuc(G)>\spectralrst(X)$ (when $d_{\mathrm{out}}\ge d_{\mathrm{in}}$).}
{Newton--Schulz step counts; the band of Muon's update; the crossover in the aligned case ($\alpha=0.45$ predicted, $0.39$ observed); full-rank RMS-to-RMS steps move every output by exactly $\eta\norm{x}_{\rms}$.}
{Momentum and noise, the other layers, many steps outside the aligned case, low precision beyond one step, why gradients keep a high nuclear rank, and why the advantage shrinks with scale.}
{Orthogonalized steps pay off on layers whose inputs have low stable rank and whose gradients are spread over many directions (proved for one step in the idealization); a band of singular values around one is good enough (folklore, supported by the speed records); at scale, add weight decay and an update-size rule (observed).}

What the theory gives a practitioner is a reading of Muon's recipe. The orthogonalization is the duality map of the spectral norm, not an ad hoc normalization; the shape factor converts it to the norm that measures how far a layer's outputs move; and the quintic trades accuracy for speed in a way one plot shows before a single run. The rank criterion says which layers should benefit, at the cost of two singular value decompositions per layer. What the theory does not give is the size of the benefit over a long run, or a reason for its decline with scale.

The rest of the unification line builds on this chapter: \cref{ch:modular} assembles the layers' norms into a norm for the network, starting from \cref{prop:spectral:rms}; \cref{ch:momentum} supplies the momentum that makes normalized steps converge; \cref{ch:kronecker} returns to Shampoo with its memory; and \cref{ch:decay} shows that Muon with weight decay is Frank--Wolfe on a spectral-norm ball.

\begin{exercises}
\item \consolidation (The face of steepest directions.) (a) For $G=e_1e_1^\top\in\R^{2\times2}$, show that the maximizers of $\ip{G,D}$ over $\norm{D}_{\op}\le1$ are the matrices $\diag(1,w)$ with $\abs w\le1$, and that Newton--Schulz returns $w=0$. (b) Show that all of them give the same guarantee \eqref{eq:spectral:rate}, but that for $\Loss(W)=\frac12\norm{W+e_1e_1^\top}_F^2$, whose gradient at $W=0$ is $e_1e_1^\top$, the loss after the step $-\eta\diag(1,w)$ depends on $w$. (c) For $m\ge n$, use $\norm{G-Q}_F^2=\norm{G}_F^2+n-2\ip{G,Q}$ to show that the semi-orthogonal $Q$ nearest to $G$ are the points of the face with $Z^\top Z=I$ (\cref{thm:toolkit:polar}).\label{ex:spectral:face}

\item \consolidation (The basin of the cubic.) Let $p(s)=\frac32s-\frac12s^3$. (a) Show that $p^{\circ k}(s)\to1$ for $s\in(0,\sqrt3)$, that $p(\sqrt3)=0$, and that $p^{\circ k}(s)\to-1$ for $s\in(\sqrt3,s_*)$, where $s_*=2.1478\ldots$ solves $p(s_*)=-\sqrt3$; so $(0,\sqrt3)$ is the largest interval $(0,c)$ on which the iteration converges to one. (b) Show that $\abs{p(s)}>s$ for $s>\sqrt5$, and that the iterates then diverge.\label{ex:spectral:basin}

\item \consolidation (One example.) For one example the gradient of a linear layer is $G=\delta x^\top$. (a) Show that $\polar(G)=\delta x^\top/(\norm{\delta}_2\norm{x}_2)$, so that memoryless Muon on one example is normalized gradient descent. (b) For $\Delta W=-\eta\sqrt{m/n}\,\polar(G)$ show that $\norm{\Delta Wx'}_{\rms}=\eta\abs{\cos\angle(x,x')}\norm{x'}_{\rms}$.\label{ex:spectral:rankone}

\item \consolidation (Designing a quintic.) For $a\in\R$ let $p_a(s)=as+\frac{5-4a}2s^3+\frac{2a-3}2s^5$. (a) Show that $p_a(1)=1$, $p_a'(1)=0$, $p_a''(1)=8a-15$, and that $a=\frac32$ gives the cubic. (b) Show that $1-p_{15/8}(s)=\frac18(1-s)^3(8+9s+3s^2)$: cubic convergence, the polynomial Polar Express uses once $[\ell_k,u_k]$ shrinks to a point. (c) Show that no $p_a$ with Muon's slope $3.4445$ at zero has a local maximum at one, and that Muon's quintic is not a $p_a$.\label{ex:spectral:quintic}

\item \variation (Approximate orthogonalization.) Let the update be $X=U_r\diag(\tau_i)V_r^\top$ with $0<\tau_{\min}\le\tau_i\le\tau_{\max}$ instead of $\polar(G)$. (a) Under \cref{ass:spectral:smooth}, show that $\Loss(W-\eta X)\le\Loss(W)-\eta\tau_{\min}\norm{G}_{\nuc}+\frac12L_{\op}\eta^2\tau_{\max}^2$, and that the best guaranteed decrease is $(\tau_{\min}/\tau_{\max})^2$ times that of the polar factor. (b) Evaluate it for Muon's band $J$ ($0.32$), for Polar Express after five steps on $s\ge10^{-3}$ ($0.57$), and for five cubic steps on $s\ge0.3$ ($0.99$).\label{ex:spectral:approx}

\item \variation (Momentum and scale.) Show that $\polar(cM)=\polar(M)$ for $c>0$, that in exact arithmetic Muon's Newton--Schulz output is unchanged when its input is multiplied by $c$ (up to the $10^{-7}$ of the normalization), and deduce that the moving-average form of the practice box and the summed form of Moonlight's code, $\tilde M_t=\beta\tilde M_{t-1}+G_t$ with input $G_t+\beta\tilde M_t$, give the same updates from a zero start.\label{ex:spectral:momentum}

\item \variation (Weight decay bounds the spectral norm.) Let $W_{t+1}=(1-\eta\lambda)W_t-\eta\polar(G_t)$ with $0<\eta\lambda\le1$. Show that $\norm{W_t}_{\op}\le1/\lambda+(1-\eta\lambda)^t\max\{0,\norm{W_0}_{\op}-1/\lambda\}$, and write the update as a convex combination of $W_t$ and a point of the ball $\norm{W}_{\op}\le1/\lambda$ (\cref{ch:decay}).\label{ex:spectral:decay}

\item \simulation (Rounding fills the null space.) Run five steps of Muon's quintic, as in \texttt{03-spectral-ns.py}, in float64 and in bfloat16 on the product $G$ of Gaussian matrices of sizes $256\times r$ and $r\times512$, $r\in\{1,4,16,64\}$, and count the output's singular values above $0.1$: exactly $r$ in float64, but $52$--$56$, $99$--$101$, $132$--$134$, and $170$--$173$ in bfloat16 in our four runs, the largest spurious one between $0.19$ and $0.43$. Explain this with the relative rounding error $2^{-8}$ and \cref{fig:spectral:ns}.\label{ex:spectral:bf16}

\item \simulation (Generic targets.) Modify \texttt{03-spectral-crossover.py} so that the initial error is $U\diag(i^{-\alpha})V^\top$ with $V$ Haar-random and independent of $Q$, and average over five seeds. Check that $\spectralrnuc(G_0)/\spectralrst(X)$ stays above $2$ for $\alpha\le0.5$, that gradient descent needs $17$ steps for every $\alpha$, and that spectral descent needs about $1$, $5$, $15$, $36$, and $63$ at $\alpha=0$, $0.2$, $0.3$, $0.4$, and $0.5$. Why does $\alpha$ not matter for gradient descent here?\label{ex:spectral:generic}

\item \research (Nuclear rank versus stable rank in trained networks.) Measure in a small transformer the nuclear rank of each hidden layer's minibatch gradient and the stable rank of its inputs along training, and test whether \cref{thm:spectral:whenhelps} predicts the layers on which Muon helps. \Textcite{DavisDrusvyatskiy2025} prove that the ratio grows with dimension in spiked random-feature models; as far as we know, no proof explains why trained deep networks keep a high nuclear rank.\label{ex:spectral:research1}

\item \research (Polynomials for low precision.) \Textcite{AmselEtAl2026} find the optimal composition of odd polynomials in exact arithmetic. Add a relative rounding error $u$ at every multiplication and find the five quintics that minimize the worst deviation of the output singular values from one on $[\ell,1]$. As far as we know this is open; \cref{fig:spectral:ns} suggests that the answer is less steep than Polar Express at first.\label{ex:spectral:research2}
\end{exercises}

\notesandrefs
Steepest descent in the spectral norm entered machine learning with \textcite{CarlsonEtAl2015a}, for restricted Boltzmann machines, and \textcite{CarlsonEtAl2015b}, with preconditioning for feedforward and convolutional networks. Shampoo is due to \textcite{GuptaKorenSinger2018}. \Textcite{BernsteinNewhouse2024} observed that Adam, Shampoo, and Prodigy are steepest descent under particular norms once their moving averages are switched off, and \textcite{BernsteinNewhouse2025} and \textcite{LargeEtAl2024} built norms and duality maps for whole networks.

Muon was introduced by \textcite{JordanEtAl2024Muon} and scaled by \textcite{LiuEtAl2025Muon} and the Kimi team \parencite{KimiTeam2025}. Analyses include \textcite{LiHong2025}, \textcite{Kovalev2025}, who reads orthogonalized methods as non-Euclidean trust-region methods, \textcite{ShenEtAl2025}, and \textcite{PethickEtAl2025}, whose methods built on linear minimization oracles over norm balls unify several optimizers and give the Scion optimizer; \textcite{ChenLiLiu2026} show that Muon with decoupled weight decay solves a spectral-norm-constrained problem. \Textcite{WenEtAl2026Fantastic} benchmark ten optimizers with careful tuning, and \textcite{WenEtAl2026Hyperball} replace weight decay by fixed norms.

The iteration for the inverse is due to \textcite{Schulz1933} and the iterative computation of the best orthogonal approximation to \textcite{BjorckBowie1971}; Chapter~8 of \textcite{Higham2008} is devoted to the polar decomposition. \Textcite{AmselEtAl2026} derive optimal polynomial compositions for Muon, and \textcite{ZhangEtAl2026Gram} reorganize the iteration around the Gram matrix. The rank comparison of \cref{thm:spectral:whenhelps}(ii) is due to \textcite{DavisDrusvyatskiy2025}; the exact value of $L_{\op}$ when $d_{\mathrm{out}}<\rank A$ (\cref{lem:spectral:layer}), the tightness statements of \cref{thm:spectral:whenhelps}(iii), the aligned reduction of \cref{prop:spectral:aligned}, and the band of \cref{prop:spectral:quintic} are worked out here, and we know of no earlier reference for them.

\paragraph{Translating the literature.} In Muon's reference code \parencite{Jordan2024} \texttt{G} is the Nesterov combination of momentum and gradient, \texttt{X} our iterate $X_k$, and \texttt{steps} our $k$; \texttt{zeropower\_via\_newtonschulz5} approximates our $\polar(G)$; \texttt{lr} is our $\eta$, ``in units of spectral norm per update''; and the factor \texttt{max(1, update.size(-2) / update.size(-1))**0.5} is $\sqrt{\max(1,d_{\mathrm{out}}/d_{\mathrm{in}})}$ for PyTorch's (out, in) layout of a weight. \Textcite{GuptaKorenSinger2018} write $L_t$ and $R_t$ for the left and right statistics, as in our field guide. Moonlight's code writes $A,B$ for the two dimensions of a weight and scales the learning rate by $0.2\sqrt{\max(A,B)}$. \Textcite{DavisDrusvyatskiy2025} call our nuclear rank the squared nuclear-to-Frobenius ratio of the gradient and compare it with the stable rank of the incoming activations, our $\spectralrst(X)$. The Polar Express code \parencite{AmselEtAl2026} optimizes each polynomial on an interval $[\ell,u]$ of normalized singular values, our $[\ell_k,u_k]$.

\printchapterbib
\end{refsection}

% Chapter 4 of Optimizers in the Wild: norms for networks (the modular norm, the spectral condition, muP).
\newcommand{\modularh}{d_{\mathrm h}}
\newcommand{\modularin}{d_{\mathrm{in}}}
\newcommand{\modularout}{d_{\mathrm{out}}}

\casechapter{Norms for Networks}{The modular norm, the spectral condition and $\mu$P}{Hyperparameter transfer across width}{ch:modular}
\begin{refsection}[chapters/04-modular.bib]

\begin{chaptersummary}
Hyperparameters tuned on a narrow network are reused on a wide one, and this works only if each layer's step is scaled with the widths in the right way: in $\mu$P, Adam's learning rate for hidden weights is divided by the width, Muon's orthogonalized update is multiplied by $\sqrt{\modularout/\modularin}$, and SGD needs per-layer rates proportional to fan-out over fan-in. We measure features by their root mean square, each layer by its RMS-to-RMS operator norm, and the network by the maximum of its layers' norms. We prove that steepest descent in this norm is per-layer normalized descent; compute exactly how far one step of SGD, of sign descent and of spectral descent moves a layer's features on one example, which yields all three width rules; and prove that weights and updates of width-independent operator norm give features and feature changes of width-independent size, with a lower bound that rules out cancellation between layers. The theory says how learning rates must scale with width, not what they should be, and its picture of a rank-one gradient aligned with the features holds only early in training.
\end{chaptersummary}

\section{In the wild: learning rates that survive widening}\label{sec:modular:wild}

No one tunes the learning rate of a network with tens of billions of parameters by trying twenty values, because each trial costs as much as the network. Practitioners tune a network a hundred times smaller and carry the result over. Whether that works depends on the \term{parametrization}\index{parametrization}: the three choices, made for every layer as functions of its widths, of the scale of the initial weights, of a constant multiplier on the layer's output, and of the layer's learning rate. In the \term{standard parametrization}\index{parametrization!standard} (SP) of the deep-learning libraries, initial weights have variance proportional to $1/\text{fan-in}$, there are no multipliers, and one learning rate serves every layer. \Textcite{YangHu2021} showed that the standard and the neural-tangent parametrizations admit no infinite-width limit that learns features, and they identified the parametrization under which every layer's features move by an amount of order one in every step, at every width: the \term{maximal update parametrization}\index{maximal update parametrization ($\mu$P)}, $\mu$P. \Textcite{YangEtAl2021} then found that in $\mu$P many optimal hyperparameters stay put as the network grows, and turned the observation into a procedure.\sidenote{The NeurIPS 2021 paper of \textcite{YangEtAl2021} appeared on arXiv in 2022 as \emph{Tensor Programs V}; the series' fourth paper, \textcite{YangHu2021}, contains the classification of parametrizations into feature-learning and kernel limits.}

\begin{practice}[$\mu$Transfer, \textcite{YangEtAl2021}; Muon, \textcite{JordanEtAl2024}; Moonlight, \textcite{LiuEtAl2025}]
Tune the learning rate and the other hyperparameters on a narrow proxy, widen the network, and keep them. With Adam in $\mu$P: initialize input and hidden weights with variance $1/\text{fan-in}$ and the output layer much smaller (variance $1/\text{fan-in}^2$, or zero); keep the learning rate of the input layer; divide the learning rates of hidden and output weights by the width multiplier, the ratio of the width to the proxy's width. Yang et al.\ transferred hyperparameters from a 13M-parameter proxy to BERT-large, and from a 40M-parameter proxy to a 6.7B-parameter GPT-3 at about 7\% of its pretraining cost. With Muon on the hidden matrices: multiply the orthogonalized update of a $\modularout\times\modularin$ matrix by $\sqrt{\max(1,\modularout/\modularin)}$ (Jordan et al.), or by $0.2\sqrt{\max(\modularout,\modularin)}$, which gives its entries the root mean square of an AdamW update so that AdamW's learning rate and weight decay can be reused (Liu et al.); embeddings, the output layer and the gains stay on AdamW.
\end{practice}

The effect is visible at the scale of a laptop. \Cref{fig:modular:curves} trains a two-hidden-layer network at five widths on a task that requires feature learning, sweeping the learning rate. In the standard parametrization the curves slide to the left as the width grows, by about a factor of two per doubling; in $\mu$P and under normalized spectral steps they stay where they are, and the wider networks are better at every learning rate near the optimum.

\begin{figure}[t]
\begin{fullwidth}
\centering
\begin{tikzpicture}
\begin{groupplot}[group style={group size=3 by 1,horizontal sep=7mm,y descriptions at=edge left},wplot,width=62mm,height=50mm,
  ymin=0.08,ymax=0.62,unbounded coords=jump,xlabel={$\log_2$ base learning rate},ylabel={loss after training},
  every axis plot/.append style={line width=0.9pt,mark=*,mark size=1.0pt}]
\nextgroupplot[title={Adam, standard parametrization},xmin=-12.5,xmax=-4.5,xtick={-12,-10,-8,-6}]
\addplot[sea!30] table[x=log2lr,y=adamsp64] {figures/data/04-modular-transfer-curves.dat};
\addplot[sea!50] table[x=log2lr,y=adamsp128] {figures/data/04-modular-transfer-curves.dat};
\addplot[sea!75] table[x=log2lr,y=adamsp256] {figures/data/04-modular-transfer-curves.dat};
\addplot[sea] table[x=log2lr,y=adamsp512] {figures/data/04-modular-transfer-curves.dat};
\addplot[seadeep] table[x=log2lr,y=adamsp1024] {figures/data/04-modular-transfer-curves.dat};
\nextgroupplot[title={Adam, $\mu$P},xmin=-10.5,xmax=-3.5,xtick={-10,-8,-6,-4}]
\addplot[sea!30] table[x=log2lr,y=adammup64] {figures/data/04-modular-transfer-curves.dat};
\addplot[sea!50] table[x=log2lr,y=adammup128] {figures/data/04-modular-transfer-curves.dat};
\addplot[sea!75] table[x=log2lr,y=adammup256] {figures/data/04-modular-transfer-curves.dat};
\addplot[sea] table[x=log2lr,y=adammup512] {figures/data/04-modular-transfer-curves.dat};
\addplot[seadeep] table[x=log2lr,y=adammup1024] {figures/data/04-modular-transfer-curves.dat};
\nextgroupplot[title={normalized spectral steps},xmin=-7.5,xmax=-0.5,xtick={-7,-5,-3,-1},
  legend style={at={(0.97,0.03)},anchor=south east,font=\scriptsize},legend cell align=left]
\addplot[sea!30] table[x=log2lr,y=spec64] {figures/data/04-modular-transfer-curves.dat};
\addplot[sea!50] table[x=log2lr,y=spec128] {figures/data/04-modular-transfer-curves.dat};
\addplot[sea!75] table[x=log2lr,y=spec256] {figures/data/04-modular-transfer-curves.dat};
\addplot[sea] table[x=log2lr,y=spec512] {figures/data/04-modular-transfer-curves.dat};
\addplot[seadeep] table[x=log2lr,y=spec1024] {figures/data/04-modular-transfer-curves.dat};
\legend{$64$,$128$,$256$,$512$,$1024$}
\end{groupplot}
\end{tikzpicture}
\end{fullwidth}
\caption{Learning-rate transfer in a small network. A ReLU network with input dimension $32$, two hidden layers of width $64$ (lightest) to $1024$ (darkest) and one output learns $y=\sum_{j\le4}((q_j^\top x)^2-1)/\sqrt8$ for orthonormal $q_j$ from fresh Gaussian minibatches of $32$ inputs, $150$ steps at a constant learning rate; the loss is measured afterwards on $4096$ fresh inputs. Left: Adam, standard parametrization (one learning rate). Middle: Adam in $\mu$P (learning rates $\eta$, $\eta\,64/\modularh$, $\eta\,64/\modularh$ for the input, hidden and output layers, output layer initialized $\sqrt{\modularh/64}$ times smaller). Right: every layer takes the normalized step $\eta\lmo_{\rms\to\rms}$ of its minibatch gradient, Muon without momentum. One seed, the same initial weights and data for every learning rate; script \texttt{04-modular-transfer.py}.}
\label{fig:modular:curves}
\end{figure}

\begin{observation}[Transfer under $\mu$P; \textcite{YangEtAl2021}, \textcite{VyasEtAl2023}]\label{obs:modular:transfer}
Under the standard parametrization the learning rate that minimizes the loss after a fixed amount of training decreases as the width grows. Under $\mu$P the minimizer is nearly independent of the width, and on simple tasks the loss curves and even the predictions of wide networks of different widths nearly coincide throughout training.
\end{observation}

The observation is robust for width and for the learning rate, in transformers, ResNets and multilayer perceptrons. Its limits are documented too. \Textcite{VyasEtAl2023} found that on harder tasks and late in training the networks of different widths drift apart. \Textcite{Lingle2024}, in an unrefereed study of transformers with up to 10 billion parameters, found that transfer holds in most configurations but fails in some, among them trainable gains in the normalization layers and the Lion optimizer. And $\mu$P is not the only rule that transfers.

\begin{observation}[Every parametrization can transfer; \textcite{EverettEtAl2024}]\label{obs:modular:everett}
Training tens of thousands of models with three optimizers and four parametrizations, up to $26.8$ billion parameters, Everett et al.\ found that every parametrization transfers once each layer receives its own learning-rate exponent. For Adam in the standard parametrization they recommend an embedding learning rate independent of the width $\modularh$ and hidden and readout learning rates proportional to $1/\modularh$, a prescription that outperformed $\mu$P in their comparison. Adam's $\epsilon$ must be scaled with the width as well, or removed (their Adam-atan2).
\end{observation}

A third observation, recent and not yet tested as widely, qualifies all of this. \Textcite{KossonEtAl2026} report that in language-model training the alignment assumptions behind $\mu$P, which concern how a layer's inputs line up with its weights and its updates, hold only at the start of training; afterwards it is weight decay, not $\mu$P, that keeps the size of the updates consistent across widths. We return to it in \cref{sec:modular:misses}.

The width rules in use look unrelated. One divides Adam's hidden learning rate by the width, one multiplies Muon's update by the square root of an aspect ratio, one leaves SGD's hidden learning rate alone and scales its input layer up. What do they have in common, why does the default fail, and which quantity is the same at every width?

\section{The idealization: layers, norms and one step}\label{sec:modular:ideal}

\begin{assumption}[A network of layers]\label[assumption]{ass:modular:network}
The parameters are $\theta=(W_1,\dots,W_N)$ with $W_l\in\R^{d_l\times d_{l-1}}$. An input $x=h_0\in\R^{d_0}$ produces the pre-activations $u_l=W_lh_{l-1}$, the features $h_l=\phi(u_l)$ for $l<N$, and the output $f=u_N\in\R^{d_N}$. The activation $\phi$ acts entrywise, $\phi(0)=0$, and $\abs{\phi(a)-\phi(b)}\le\abs{a-b}$ for all real $a,b$, as for ReLU, $\tanh$ and the identity.
\end{assumption}

The loss of an example $\xi=(x,y)$ is $\ell(\theta;\xi)=\psi(f)$ for a differentiable function $\psi$ of the output, such as $\psi(f)=\frac12\norm{f-y}_2^2$. Backpropagation computes the \term{backpropagated signals}\index{backpropagated signal} $\delta_l=\nabla_{u_l}\ell$ from $\delta_N=\nabla\psi(f)$ by $\delta_l=\phi'(u_l)\odot(W_{l+1}^\top\delta_{l+1})$, where $\odot$ is the entrywise product, and the gradient of layer $l$ on one example is the outer product
\begin{equation}\label{eq:modular:rankone}
  G_l=\nabla_{W_l}\ell=\delta_l\,h_{l-1}^\top .
\end{equation}
For a minibatch, $G_l$ is the average of such outer products and has rank at most the batch size. We write $\modularin=d_{l-1}$ and $\modularout=d_l$ (fan-in and fan-out) when one layer is in view, and $\modularh$ for a common hidden width.

\paragraph{How to measure.} A vector of $m$ entries of order one has Euclidean norm of order $\sqrt m$ but \term{RMS norm}\index{RMS norm} $\norm{v}_\rms=\norm{v}_2/\sqrt m$ of order one, so ``every activation has a typical size of one'' reads ``$\norm{h_l}_\rms=\Theta(1)$''. Features are measured in the RMS norm, layers in the induced operator norm $\norm{W}_{\rms\to\rms}=\max_{v\ne0}\norm{Wv}_\rms/\norm{v}_\rms$ (\cref{ch:spectral}), and backpropagated signals, which pair with features, in the dual of the RMS norm, $\norm{v}_{\rms,*}=\max_{\norm{z}_\rms\le1}\ip{v,z}=\sqrt m\,\norm{v}_2$. A quantity is \term{width-independent} if it is bounded above and below by constants that do not depend on $d_1,\dots,d_{N-1}$; the depth $N$, the input and output dimensions $d_0$ and $d_N$, the data and the base step size are fixed.

\begin{assumption}[One step from one example]\label[assumption]{ass:modular:onestep}
The step $\Delta\theta=(\Delta W_1,\dots,\Delta W_N)$ is computed from the loss of one example $\xi=(x,y)$, so that every layer gradient is the rank-one matrix \eqref{eq:modular:rankone}, and its effect is measured on the same input $x$: $\Delta u_l$ and $\Delta h_l$ denote the changes of $u_l$ and $h_l$ when $\theta$ becomes $\theta+\Delta\theta$.
\end{assumption}

\begin{definition}[The feature-learning desideratum; \textcite{YangSimonBernstein2023}]\label[definition]{def:modular:desideratum}
A rule that chooses initial weights and steps as functions of the widths \term{learns features}\index{feature learning} if $\norm{h_l}_\rms$ and $\norm{\Delta h_l}_\rms$ are width-independent for every hidden layer $l<N$, and the output change $\norm{\Delta f}_\rms$ is bounded independently of the widths.
\end{definition}

A rule under which $\norm{\Delta h_l}_\rms\to0$ as the widths grow is \term{lazy}\index{lazy training}: the network then trains like its linearization at initialization, a kernel method \parencite{JacotGabrielHongler2018,ChizatOyallonBach2019}. A rule under which some $\norm{\Delta h_l}_\rms\to\infty$, or $\norm{\Delta f}_\rms\to\infty$, is unstable at large width unless its learning rate shrinks.

\begin{dictionary}
the RMS norm $\norm{h}_\rms$ & the typical size of one activation, which unit scaling and u-$\mu$P keep near one\\
the RMS$\to$RMS norm $\norm{W}_{\rms\to\rms}=\sqrt{\modularin/\modularout}\,\norm{W}_{\op}$ & the largest gain of a layer on its activations\\
the dual norm $\norm{\delta}_{\rms,*}=\sqrt m\,\norm{\delta}_2$ & the size of the gradient with respect to a layer's pre-activations, which shrinks like $1/\text{width}$ per entry\\
the scales $s_l$ of the max norm (\cref{def:modular:maxnorm}) & per-layer learning-rate multipliers: $\mu$P's rates, Modula's masses\\
normalized steepest descent in the max norm & per-layer normalized optimizers: LARS, Muon with its shape factor, Modula\\
one step from one example & the first steps of training, when gradients are low-rank and aligned with the activations that produced them\\
$\norm{\Delta h_l}_\rms=\Theta(1)$ at every width & feature learning, as opposed to lazy (kernel) training\\
a width rule & a column of a $\mu$P table
\end{dictionary}

\paragraph{What was thrown away.} Almost everything that makes a modern network: attention, normalization layers, embeddings, residual connections and depth that grows, biases, weight decay, the optimizer's memory, and all examples but one. What is kept is the question of scale. Each simplification is the right one to make first. A single example makes every layer's gradient rank one, so that its geometry is that of a pair of vectors; one step isolates the effect of a learning rate from the history of training; and the question ``how must the learning rate depend on the width?'' is answered by exponents, which a calculation that is exact up to constants can find. Early in training, gradients are in fact of low rank and aligned with the features, which is where the rules matter most: a learning rate that is too large at a large width diverges in the first steps.

\paragraph{What one had to notice.} Two things. First, that the size of a weight update must be measured by its effect on the features, which is an operator norm, and not by the size of its entries or its Frobenius norm. A rank-one matrix with entries $\pm\eta$ has operator norm $\eta\sqrt{\modularout\modularin}$, a random matrix with the same entries one of order $\eta(\sqrt{\modularout}+\sqrt{\modularin})$, and an entrywise rule cannot tell them apart; this is the point of \textcite{YangSimonBernstein2023}. Second, that once every layer is measured in such a norm, the network can be given the maximum of its layers' norms, and then steepest descent in that one norm hands each layer its own normalized step: per-layer learning rates are not a hack added to an optimizer but the duality map of a norm \parencite{LargeEtAl2024,BernsteinNewhouse2025}.

\begin{fieldguide}{$\mu$P Adam}{\textcite{YangHu2021}, \textcite{YangEtAl2021}}
\fgrow{update}{Adam on every layer; learning rate $\eta$ for the input layer and $\eta\,d^{\circ}/\modularin$ for hidden and output weights, $d^{\circ}$ the proxy's width}
\fgrow{norm}{with $\beta_1=\beta_2=0$ (sign descent, \cref{ch:sign}), normalized steepest descent in $\max_ls_l\max_{ij}\abs{(W_l)_{ij}}$ with $s_l=\modularin/d^{\circ}$ for hidden and output layers (\cref{thm:modular:dual}), which controls the RMS$\to$RMS norms (\cref{ex:modular:signnorm})}
\fgrow{preconditioner}{diagonal, Adam's (\cref{ch:adaptive})}
\fgrow{averages}{Adam's $\beta_1$, $\beta_2$}
\fgrow{state}{two numbers per parameter}
\fgrow{in the wild}{initialization variance $1/\text{fan-in}$, output layer near zero; the \texttt{mup} package divides the hidden learning rate by the width multiplier}
\end{fieldguide}

\begin{fieldguide}{Modular normalized descent}{\textcite{LargeEtAl2024}}
\fgrow{update}{$W_l\leftarrow W_l+(\eta/s_l)\lmo_{(l)}(M_l)$, $M_l$ the gradient or its momentum; for a linear layer $\lmo_{\rms\to\rms}(M)=-\sqrt{\modularout/\modularin}\,\polar(M)$}
\fgrow{norm}{the modular norm, a scaled maximum of the modules' norms (\cref{def:modular:maxnorm}), RMS$\to$RMS for linear layers}
\fgrow{preconditioner}{none}
\fgrow{averages}{optional momentum}
\fgrow{state}{one number per parameter with momentum}
\fgrow{in the wild}{the Modula library, with scales computed from user-assigned masses; Muon (\cref{ch:spectral}) on one matrix, with Newton--Schulz in place of the polar factor}
\end{fieldguide}

\begin{fieldguide}{LARS and LAMB}{\textcite{YouGitmanGinsburg2017}, \textcite{YouEtAl2020}}
\fgrow{update}{$W_l\leftarrow W_l-\eta\tau\,(\norm{W_l}_F/\norm{G_l}_F)\,G_l$, with momentum and the trust coefficient $\tau$; LAMB applies the same ratio to Adam's update}
\fgrow{norm}{a maximum of Frobenius norms with the state-dependent scales $s_l=1/\norm{W_l}_F$: each layer moves by a fixed fraction of its own norm}
\fgrow{preconditioner}{none (LARS); diagonal (LAMB)}
\fgrow{averages}{heavy-ball momentum (LARS); Adam's (LAMB)}
\fgrow{state}{one (LARS) or two (LAMB) numbers per parameter}
\fgrow{in the wild}{ResNet-50 at batch size 32K; BERT at batch size 32K in 76 minutes}
\end{fieldguide}

\section{The theory: the norm of a network}\label{sec:modular:norm}

We first collect the facts about the RMS geometry of one layer, most of which \cref{ch:spectral} proved in the language of the spectral norm.

\begin{lemma}[RMS geometry of a layer]\label[lemma]{lem:modular:rms}
Let $W,G\in\R^{\modularout\times\modularin}$ and $v\in\R^{\modularout}$.
\begin{enumerate}
\item $\norm{W}_{\rms\to\rms}=\sqrt{\modularin/\modularout}\,\norm{W}_{\op}$.
\item The dual of $\norm{\cdot}_{\rms\to\rms}$ under $\ip{G,W}=\tr(G^\top W)$ is $\norm{G}_{\rms\to\rms,*}=\sqrt{\modularout/\modularin}\,\norm{G}_{\nuc}$, and $-\sqrt{\modularout/\modularin}\,\polar(G)\in\lmo_{\rms\to\rms}(G)$ when $G\ne0$.
\item $\norm{W^\top v}_{\rms,*}\le\norm{W}_{\rms\to\rms}\,\norm{v}_{\rms,*}$.
\item If $G=ab^\top$ with $a\ne0$ and $b\ne0$, then $\polar(G)=\frac{a}{\norm a_2}\frac{b^\top}{\norm b_2}$, $\norm{G}_{\rms\to\rms}=\norm{a}_\rms\norm{b}_{\rms,*}$, $\norm{G}_{\rms\to\rms,*}=\norm{a}_{\rms,*}\norm{b}_\rms$, and $\norm{Gb}_\rms=\norm{G}_{\rms\to\rms}\norm{b}_\rms$.
\end{enumerate}
\end{lemma}

\begin{proof}
(i) For $v\ne0$, $\norm{Wv}_\rms/\norm{v}_\rms=\sqrt{\modularin/\modularout}\,\norm{Wv}_2/\norm{v}_2$; take the maximum over $v$. (ii) For a norm $\norm{\cdot}$ and $c>0$, the dual of $c\norm{\cdot}$ is $c^{-1}\norm{\cdot}_*$, because $\max_{c\norm{U}\le1}\ip{G,U}=c^{-1}\max_{\norm{U'}\le1}\ip{G,U'}$, and the substitution $U=U'/c$ shows that $\lmo_{c\norm\cdot}(G)=c^{-1}\lmo_{\norm\cdot}(G)$. The dual of the spectral norm is the nuclear norm and $-\polar(G)\in\lmo_{\op}(G)$ (\cref{thm:toolkit:vonneumann}; \cref{ch:spectral}). (iii) $\sqrt{\modularin}\norm{W^\top v}_2\le\sqrt{\modularin}\norm{W}_{\op}\norm{v}_2=\sqrt{\modularin/\modularout}\,\norm{W}_{\op}\cdot\sqrt{\modularout}\norm{v}_2$. (iv) The compact singular value decomposition of $ab^\top$ has the single singular value $\norm a_2\norm b_2$ with singular vectors $a/\norm a_2$ and $b/\norm b_2$, so $\norm{G}_{\op}=\norm{G}_{\nuc}=\norm a_2\norm b_2$, and (i) and (ii) give the two norms. Finally $Gb=\norm{b}_2^2\,a$, whose RMS norm is $\norm{a}_\rms\norm{b}_2^2=\norm{a}_\rms\norm{b}_{\rms,*}\norm{b}_\rms$.
\end{proof}

Part (iii) says that the backward pass contracts dual norms by the same factor as the forward pass contracts RMS norms; part (iv), that a rank-one matrix acts on its own right singular vector with its full operator norm. Both are used repeatedly below.

\subsection{Steepest descent in a maximum of layer norms}

\begin{definition}[The max norm of a network]\label[definition]{def:modular:maxnorm}
Give each layer a norm $\norm{\cdot}_{(l)}$ on $\R^{d_l\times d_{l-1}}$, with dual norm $\norm{\cdot}_{(l),*}$, and a scale $s_l>0$. The \term{max norm}\index{modular norm} of the network is $\norm{\theta}=\max_ls_l\norm{W_l}_{(l)}$, on the parameter space with the inner product $\ip{g,\theta}=\sum_l\tr(G_l^\top W_l)$.
\end{definition}

For a sequential network, the \term{modular norm} of \textcite{LargeEtAl2024} is of this form, as \textcite{BernsteinNewhouse2024} write it: the layer norms are operator norms chosen for each kind of module (RMS$\to$RMS for a linear layer) and the scales are computed from two numbers attached to every module, its mass, which sets its share of the feature learning, and its sensitivity to its input. The positivity, homogeneity and triangle inequality of $\norm{\cdot}$ follow from those of the layer norms, since a maximum of sums is at most the sum of the maxima. Steepest descent in the max norm is computed layer by layer.

\begin{theorem}[Steepest descent in the norm of a network]\label{thm:modular:dual}
For the max norm of \cref{def:modular:maxnorm} and every $g=(G_1,\dots,G_N)$:
\begin{enumerate}
\item the dual norm is $\norm{g}_*=\sum_{l=1}^N\norm{G_l}_{(l),*}/s_l$;
\item $u=(U_1,\dots,U_N)$ belongs to $\lmo(g)$ if and only if $s_lU_l\in\lmo_{(l)}(G_l)$ for every $l$;
\item normalized steepest descent $\theta_{t+1}=\theta_t+\eta\lmo(g_t)$ is the layerwise update $W_{l,t+1}=W_{l,t}+(\eta/s_l)\lmo_{(l)}(G_{l,t})$: normalized descent in every layer, with the per-layer learning rate $\eta/s_l$;
\item unnormalized steepest descent $\theta_{t+1}=\theta_t-\eta g_t^\sharp$ is
\[
  W_{l,t+1}=W_{l,t}+\frac{\eta}{s_l}\Bigl(\sum_{k=1}^N\frac{\norm{G_{k,t}}_{(k),*}}{s_k}\Bigr)\lmo_{(l)}(G_{l,t}):
\]
the same directions, every layer's step multiplied by one global factor, the dual norm of the whole gradient.
\end{enumerate}
\end{theorem}

\begin{proof}
The plan is to observe that the unit ball of the max norm is a product of layer balls and that the linear function $\ip{g,\cdot}$ is a sum of layer terms, so that minimizing it over the ball splits into one problem per layer. The unit ball is $\{\theta:\norm{W_l}_{(l)}\le1/s_l\text{ for every }l\}$, the product of the balls $\cB_l=\{U:\norm{U}_{(l)}\le1/s_l\}$.

(i) The maximum of $\sum_l\ip{G_l,U_l}$ over a product of sets is the sum of the maxima of the terms. By the definition of the dual norm and its homogeneity, $\max_{U\in\cB_l}\ip{G_l,U}=\norm{G_l}_{(l),*}/s_l$.

(ii) If every $U_l$ minimizes $\ip{G_l,\cdot}$ over $\cB_l$, then $u$ minimizes the sum over the product. Conversely, if some $U_k$ does not, replace it by a minimizer, which exists because $\cB_k$ is compact: the sum strictly decreases, so $u$ was not a minimizer. The minimizers of $\ip{G_l,\cdot}$ over $\cB_l$ are $1/s_l$ times those over the unit ball of $\norm{\cdot}_{(l)}$, by the substitution $U=U'/s_l$; that is, $U_l$ is a minimizer if and only if $s_lU_l\in\lmo_{(l)}(G_l)$.

(iii) This is (ii) with $u=\lmo(g_t)$ multiplied by $\eta$.

(iv) By \cref{lem:toolkit:conjugate}, $g^\sharp=-\norm{g}_*\lmo(g)$ for every norm; insert (i) and (ii).
\end{proof}

Two degenerate cases deserve a word. If $G_l=0$, every point of the layer's ball is a minimizer, and implementations take $U_l=0$, which keeps the layer still. If $g=0$, both updates are zero.

\begin{remark}[Which combination of layer norms?]\label[remark]{rem:modular:sum}
The maximum is a modelling choice. The sum $\sum_ls_l\norm{W_l}_{(l)}$ has the dual norm $\max_l\norm{G_l}_{(l),*}/s_l$, and its $\lmo$ puts the whole step on the layers that attain this maximum (\cref{ex:modular:sum}): steepest descent in a sum of layer norms is greedy block-coordinate descent, close in spirit to the duality-structure gradient descent of \textcite{Flynn2017}, which updates one greedily chosen layer per step. The maximum is the combination in which a step of size $\eta$ moves every layer by its full share at once, which is what the desideratum of \cref{def:modular:desideratum} asks for.
\end{remark}

\Cref{thm:modular:dual} turns the question ``how should the learning rate of each layer depend on the widths?'' into the question ``which norm should each layer carry?''. Once every layer's norm measures the effect of its weights on its features in width-independent units, the scales $s_l$ can be constants, and one learning rate $\eta$ serves every width. With smoothness with respect to the max norm, the descent bound of \cref{ch:steepest} reads $\Loss(\theta+\eta\lmo(g))\le\Loss(\theta)-\eta\sum_l\norm{G_l}_{(l),*}/s_l+\frac L2\eta^2$: the guaranteed decrease adds up over the layers.

\subsection{One step of three optimizers on one layer}

Which layer norms have this property? The answer comes from computing exactly what one step of each optimizer does to the layer's output, on the example that produced the step.

\begin{theorem}[One step of three optimizers on one layer]\label{thm:modular:widthrules}
Let $W\in\R^{\modularout\times\modularin}$ be a layer whose input on an example is $x\ne0$ and whose backpropagated signal is $\delta\ne0$, so that its gradient is $G=\delta x^\top$. Write $\rho(x)=\norm{x}_1/(\sqrt{\modularin}\,\norm{x}_2)\in(0,1]$ for the density of $x$, and $\sigma(v)=\norm{\sign v}_\rms$ for the square root of the fraction of nonzero entries of a vector $v$. Let $\eta>0$.
\begin{enumerate}
\item (SGD) The step $\Delta W=-\eta G$ gives $\Delta Wx=-\eta\norm{x}_2^2\,\delta$ and
\[
  \norm{\Delta Wx}_\rms=\eta\,\modularin\norm{\delta}_\rms\norm{x}_\rms^2,\qquad \norm{\Delta W}_{\rms\to\rms}=\eta\,\modularin\norm{\delta}_\rms\norm{x}_\rms .
\]
\item (Sign descent) The step $\Delta W=-\eta\sign(G)$ gives $\Delta Wx=-\eta\norm{x}_1\sign(\delta)$ and
\[
  \norm{\Delta Wx}_\rms=\eta\,\modularin\,\rho(x)\,\sigma(\delta)\norm{x}_\rms,\qquad \norm{\Delta W}_{\rms\to\rms}=\eta\,\modularin\,\sigma(\delta)\,\sigma(x).
\]
\item (Normalized spectral step) The step $\Delta W=\eta\lmo_{\rms\to\rms}(G)=-\eta\sqrt{\modularout/\modularin}\,\polar(G)$ gives $\Delta Wx=-\eta\sqrt{\modularout}\,\norm{x}_\rms\,\delta/\norm{\delta}_2$ and
\[
  \norm{\Delta Wx}_\rms=\eta\norm{x}_\rms,\qquad \norm{\Delta W}_{\rms\to\rms}=\eta .
\]
\end{enumerate}
Consequently the step sizes that move the layer's output on $x$ by $\norm{\Delta Wx}_\rms=c\norm{x}_\rms$ are
\begin{equation}\label{eq:modular:rules}
  \eta_{\mathrm{SGD}}=\frac{c}{\modularin\norm{\delta}_\rms\norm{x}_\rms},\qquad
  \eta_{\mathrm{sign}}=\frac{c}{\modularin\,\rho(x)\,\sigma(\delta)},\qquad
  \eta_{\mathrm{spec}}=c .
\end{equation}
\end{theorem}

\begin{proof}
(i) $\Delta Wx=-\eta\,\delta\,(x^\top x)$. Its RMS norm is $\eta\norm{x}_2^2\norm{\delta}_\rms=\eta\,\modularin\norm{x}_\rms^2\norm{\delta}_\rms$, and by \cref{lem:modular:rms}(iv), $\norm{\Delta W}_{\rms\to\rms}=\eta\norm{\delta}_\rms\norm{x}_{\rms,*}=\eta\norm{\delta}_\rms\,\modularin\norm{x}_\rms$.

(ii) Since $\sign(\delta_ix_j)=\sign(\delta_i)\sign(x_j)$, also when one factor is zero, $\sign(G)=\sign(\delta)\sign(x)^\top$, and $\sign(x)^\top x=\norm{x}_1$. Hence $\Delta Wx=-\eta\norm{x}_1\sign(\delta)$, whose RMS norm is $\eta\norm{x}_1\sigma(\delta)$, and $\norm{x}_1=\rho(x)\sqrt{\modularin}\norm{x}_2=\rho(x)\,\modularin\norm{x}_\rms$. By \cref{lem:modular:rms}(iv), $\norm{\Delta W}_{\rms\to\rms}=\eta\norm{\sign\delta}_\rms\norm{\sign x}_{\rms,*}=\eta\,\sigma(\delta)\,\modularin\,\sigma(x)$.

(iii) By \cref{lem:modular:rms}(ii) and (iv), $\Delta W=-\eta\sqrt{\modularout/\modularin}\,\frac{\delta}{\norm\delta_2}\frac{x^\top}{\norm x_2}$, so $\Delta Wx=-\eta\sqrt{\modularout/\modularin}\,\norm{x}_2\,\delta/\norm{\delta}_2=-\eta\sqrt{\modularout}\,\norm{x}_\rms\,\delta/\norm{\delta}_2$, whose RMS norm is $\eta\norm{x}_\rms$; and $\norm{\Delta W}_{\rms\to\rms}=\eta$ because $\polar(G)$ has operator norm one. Solving each identity for $\eta$ gives \eqref{eq:modular:rules}.
\end{proof}

The three rules depend on different things, and that is the whole story of this chapter in one display. SGD's step depends on the size of the backpropagated signal $\norm{\delta}_\rms$, which is set by the layers above, and on the size of the input, squared. Sign descent is blind to the size of $\delta$ and of $x$, but its step grows with the fan-in, in proportion to the density of the input: a dense input ($\rho(x)$ of order one, as for Gaussian or ReLU features) receives a change of order $\eta\,\modularin$, a one-hot input a change of order $\eta\sqrt{\modularin}\,\norm{x}_\rms$ (\cref{ex:modular:embedding}). The normalized spectral step depends on nothing: it moves the output by exactly $\eta$ in RMS. In each case the output moves along $-\delta$ or $-\sign\delta$, the direction that decreases the loss, and in each case $\norm{\Delta Wx}_\rms$ equals $\norm{\Delta W}_{\rms\to\rms}\norm{x}_\rms$, up to the factor $\rho(x)/\sigma(x)$ for sign descent: a single-example update is aligned with the input it came from. That is why the RMS$\to$RMS norm of the update, and not the size of its entries, is the quantity a width rule must control.

\section{The theory: the spectral condition}\label{sec:modular:spectral}

\Cref{thm:modular:widthrules} concerns one layer. A network adds two questions: whether the steps of different layers can cancel each other's effect on the features, and how large the backpropagated signal is. The first has a clean answer.

\begin{lemma}[No cancellation between layers]\label[lemma]{lem:modular:nocancel}
Under \cref{ass:modular:network}, fix one example and a direction $\Delta\theta=(\Delta W_1,\dots,\Delta W_N)$, and suppose that $\phi$ is differentiable at every entry of $u_1,\dots,u_{N-1}$. The first-order changes $\dot u_l=\frac{\dd}{\dd\epsilon}u_l(\theta+\epsilon\Delta\theta)\big|_{\epsilon=0}$ satisfy
\begin{equation}\label{eq:modular:nocancel}
  \ip{\delta_l,\dot u_l}=\sum_{k=1}^l\ip{G_k,\Delta W_k}\qquad(l=1,\dots,N).
\end{equation}
Consequently, if every layer descends, $\ip{G_k,\Delta W_k}\le0$ for $k\le l$, and $\delta_l\ne0$, then $\norm{\dot u_l}_2\ge\sum_{k\le l}\abs{\ip{G_k,\Delta W_k}}/\norm{\delta_l}_2$.
\end{lemma}

\begin{proof}
The plan is an induction on $l$ along the forward pass. Differentiating $u_l=W_lh_{l-1}$ and $h_{l-1}=\phi(u_{l-1})$ gives $\dot u_l=\Delta W_l\,h_{l-1}+W_l\dot h_{l-1}$ and $\dot h_{l-1}=\phi'(u_{l-1})\odot\dot u_{l-1}$, with $\dot h_0=0$. The first term pairs with $\delta_l$ to the gradient term, $\ip{\delta_l,\Delta W_lh_{l-1}}=\tr(\Delta W_l^\top\delta_lh_{l-1}^\top)=\ip{G_l,\Delta W_l}$. The second pairs with $\delta_l$ to the previous layer's left side, because the backward recursion moves the transpose and the derivative across:
\[
  \ip{\delta_l,W_l\dot h_{l-1}}=\ip{W_l^\top\delta_l,\;\phi'(u_{l-1})\odot\dot u_{l-1}}=\ip{\phi'(u_{l-1})\odot W_l^\top\delta_l,\;\dot u_{l-1}}=\ip{\delta_{l-1},\dot u_{l-1}} .
\]
For $l=1$ the second term vanishes, which starts the induction, and adding the two terms proves \eqref{eq:modular:nocancel}. If all terms on the right of \eqref{eq:modular:nocancel} are nonpositive, then $\abs{\ip{\delta_l,\dot u_l}}=\sum_{k\le l}\abs{\ip{G_k,\Delta W_k}}$, and the Cauchy--Schwarz inequality $\abs{\ip{\delta_l,\dot u_l}}\le\norm{\delta_l}_2\norm{\dot u_l}_2$ gives the bound.
\end{proof}

The identity says that the first-order change of a layer's pre-activations, projected on its backpropagated signal, is the first-order change of the loss caused by the layers up to it. When every layer takes a descent step, all these contributions have the same sign, and the steps of the lower layers cannot undo the effect of a layer's own step in the direction $-\delta_l$. Taking only the term $k=l$, and writing $\hat\delta_l=\delta_l/\norm{\delta_l}_2$,
\begin{equation}\label{eq:modular:direct}
  \norm{\dot u_l}_2\ \ge\ \ip{-\hat\delta_l,\;\Delta W_l\,h_{l-1}} :
\end{equation}
the feature change is at least the projection of the layer's own direct change on $-\hat\delta_l$. For SGD and for normalized spectral steps the direct change is parallel to $-\delta_l$ (\cref{thm:modular:widthrules}), and \eqref{eq:modular:direct} is its full length; for sign descent it is its length times $\cos\angle(\delta_l,\sign\delta_l)=\rho(\delta_l)/\sigma(\delta_l)$. We can now state the condition that makes a network learn features at every width.

\begin{theorem}[The spectral condition]\label{thm:modular:spectralcondition}
Under \cref{ass:modular:network}, suppose that $\norm{W_l}_{\rms\to\rms}\le a_l$ and $\norm{\Delta W_l}_{\rms\to\rms}\le b_l$ for every $l$, and put $A_l=\prod_{k\le l}a_k$ and $B_l=\prod_{k\le l}(a_k+b_k)$, with $A_0=B_0=1$. For every input $x$ and every $l$:
\begin{enumerate}
\item (forward) $\norm{h_l}_\rms\le\norm{u_l}_\rms\le A_l\norm{x}_\rms$;
\item (update) $\norm{\Delta h_l}_\rms\le\norm{\Delta u_l}_\rms\le(B_l-A_l)\norm{x}_\rms$;
\item (rank-one lower bound) if $\phi$ is differentiable at the entries of $u_1,\dots,u_{N-1}$, and every layer takes the normalized step $\Delta W_k=\eta_k\lmo_{\rms\to\rms}(G_k)$ of the loss of one example with input $x$, then $\norm{\dot u_l}_\rms\ge\eta_l\norm{h_{l-1}}_\rms$ whenever $\delta_l\ne0$;
\item (deep linear networks) if $\phi$ is the identity, then
\[
  \norm{\Delta u_l-\dot u_l}_\rms\le\Bigl(B_l-A_l-\sum_{k\le l}b_k\prod_{j\le l,\,j\ne k}a_j\Bigr)\norm{x}_\rms ,
\]
and, if moreover $a_k\le1$ and the steps of (iii) all have $\eta_k=\eta$,
\begin{equation}\label{eq:modular:sandwich}
  \eta\norm{h_{l-1}}_\rms-\bigl((1+\eta)^l-1-l\eta\bigr)\norm{x}_\rms\ \le\ \norm{\Delta h_l}_\rms\ \le\ \bigl((1+\eta)^l-1\bigr)\norm{x}_\rms .
\end{equation}
\end{enumerate}
\end{theorem}

\begin{proof}
The plan is to propagate norm bounds through the layers, using that $\phi$ does not increase distances, and to obtain the lower bound from \cref{lem:modular:nocancel}.

(i) Since $\phi(0)=0$ and $\phi$ is $1$-Lipschitz entrywise, $\norm{\phi(u)}_\rms\le\norm{u}_\rms$ and $\norm{\phi(u)-\phi(u')}_\rms\le\norm{u-u'}_\rms$. So $\norm{h_{l-1}}_\rms\le\norm{u_{l-1}}_\rms$ and $\norm{u_l}_\rms\le a_l\norm{h_{l-1}}_\rms$, and induction on $l$ gives (i).

(ii) The changed network has layers $W_k+\Delta W_k$ of RMS$\to$RMS norm at most $a_k+b_k$, so (i) applied to it bounds its features: $\norm{h_{l-1}+\Delta h_{l-1}}_\rms\le B_{l-1}\norm{x}_\rms$. Write the change of the pre-activations as
\[
  \Delta u_l=(W_l+\Delta W_l)(h_{l-1}+\Delta h_{l-1})-W_lh_{l-1}=W_l\,\Delta h_{l-1}+\Delta W_l\,(h_{l-1}+\Delta h_{l-1}) .
\]
With $D_l=B_l-A_l$ and the induction hypothesis $\norm{\Delta h_{l-1}}_\rms\le\norm{\Delta u_{l-1}}_\rms\le D_{l-1}\norm{x}_\rms$, which holds for $l=1$ because $\Delta h_0=0$ and $D_0=0$,
\[
  \norm{\Delta u_l}_\rms\le\bigl(a_lD_{l-1}+b_lB_{l-1}\bigr)\norm{x}_\rms=\bigl((a_l+b_l)B_{l-1}-a_lA_{l-1}\bigr)\norm{x}_\rms=D_l\norm{x}_\rms ,
\]
and $\norm{\Delta h_l}_\rms\le\norm{\Delta u_l}_\rms$ because $\phi$ is $1$-Lipschitz.

(iii) By \cref{lem:modular:rms}(ii), every layer descends: $\ip{G_k,\Delta W_k}=-\eta_k\norm{G_k}_{\rms\to\rms,*}\le0$. By \cref{lem:modular:rms}(iv) applied to $G_l=\delta_lh_{l-1}^\top$, the term $k=l$ has absolute value $\eta_l\norm{\delta_l}_{\rms,*}\norm{h_{l-1}}_\rms=\eta_l\sqrt{d_l}\,\norm{\delta_l}_2\norm{h_{l-1}}_\rms$. \Cref{lem:modular:nocancel}, keeping only this term, gives $\norm{\dot u_l}_2\ge\eta_l\sqrt{d_l}\,\norm{h_{l-1}}_\rms$; divide by $\sqrt{d_l}$.

(iv) For a linear network, $u_l+\Delta u_l=(W_l+\Delta W_l)\cdots(W_1+\Delta W_1)x$. Expanding the product gives one term for each set $S\subseteq\{1,\dots,l\}$ of layers that contribute their $\Delta W_k$; the empty set gives $u_l$, the singletons give $\dot u_l$, and the term of $S$ has RMS norm at most $\prod_{k\in S}b_k\prod_{k\notin S}a_k\norm{x}_\rms$. Summing over $\abs S\ge2$ gives the first claim, since $\sum_{S}\prod_{S}b_k\prod_{S^c}a_k=B_l$. The bound is a polynomial with nonnegative coefficients in the $a_k$, so with $a_k\le1$ and $b_k=\eta$ it is at most $\sum_{\abs S\ge2}\eta^{\abs S}=(1+\eta)^l-1-l\eta$. The left inequality of \eqref{eq:modular:sandwich} follows from (iii) and the triangle inequality, the right one from (ii) and the same monotonicity, because $\Delta h_l=\Delta u_l$.
\end{proof}

In operator norms, the hypothesis $a_l,b_l=\Theta(1)$ reads $\norm{W_l}_{\op}=\Theta(\sqrt{d_l/d_{l-1}})$ and $\norm{\Delta W_l}_{\op}=\Theta(\sqrt{d_l/d_{l-1}})$, the \term{spectral condition}\index{spectral condition} of \textcite{YangSimonBernstein2023}. Under it, the forward and update bounds of the theorem do not involve the widths; under the normalized steps of \cref{thm:modular:dual} with RMS$\to$RMS layer norms, the lower bound does not either.

\begin{corollary}[Feature learning at every width]\label[corollary]{cor:modular:featurelearning}
In a deep linear network with $\norm{W_k}_{\rms\to\rms}\le1$ for all $k$, let every layer take the normalized RMS$\to$RMS step of size $\eta$ of one example's loss, and suppose that the forward pass does not collapse: $\norm{h_{l-1}}_\rms\ge c\,\norm{x}_\rms$ for some $c>0$. Then
\[
  \Bigl(c\,\eta-\tbinom l2\eta^2(1+\eta)^{l-2}\Bigr)\norm{x}_\rms\ \le\ \norm{\Delta h_l}_\rms\ \le\ \bigl((1+\eta)^l-1\bigr)\norm{x}_\rms ,
\]
bounds in which no width appears.
\end{corollary}

\begin{proof}
By Taylor's theorem with the Lagrange remainder, $(1+\eta)^l-1-l\eta=\binom l2(1+\xi)^{l-2}\eta^2$ for some $\xi\in(0,\eta)$, which is at most $\binom l2(1+\eta)^{l-2}\eta^2$ for $l\ge2$ (for $l=1$ both sides are zero). Insert this and the hypothesis into \eqref{eq:modular:sandwich}.
\end{proof}

The upper bound is attained to first order. If every layer is square and orthogonal and the steps are the normalized spectral steps of one example, then all $l$ contributions to $\dot u_l$ point along $-\hat\delta_l$ and have RMS norm $\eta\norm{x}_\rms$, so $\norm{\dot u_l}_\rms=l\eta\norm{x}_\rms$, which is the first-order term of $(1+\eta)^l-1$ (\cref{ex:modular:tight}). The lower bound is attained at $l=1$, where the direct term is the only one.

\subsection{The backward pass and the size of the signal}

SGD's rule \eqref{eq:modular:rules} depends on $\norm{\delta_l}_\rms$, which the layers above set. Its size is controlled by the same operator norms.

\begin{proposition}[The backward pass in dual norms]\label[proposition]{prop:modular:backward}
Under \cref{ass:modular:network}, with $\phi$ differentiable at the entries of $u_1,\dots,u_{N-1}$, suppose $\norm{W_k}_{\rms\to\rms}\le a_k$ for every $k$. For one example:
\begin{enumerate}
\item $\norm{\delta_l}_{\rms,*}\le\bigl(\prod_{k=l+1}^Na_k\bigr)\norm{\nabla\psi(f)}_{\rms,*}$, that is, $\norm{\delta_l}_\rms\le\frac1{d_l}\bigl(\prod_{k>l}a_k\bigr)\norm{\nabla\psi(f)}_{\rms,*}$;
\item $\norm{G_l}_{\rms\to\rms,*}=\norm{\delta_l}_{\rms,*}\norm{h_{l-1}}_\rms\le\bigl(\prod_{k\ne l}a_k\bigr)\norm{\nabla\psi(f)}_{\rms,*}\,\norm{x}_\rms$;
\item in the max norm with RMS$\to$RMS layer norms and scales $s_l$, $\norm{\nabla_\theta\ell}_*\le\bigl(\sum_ls_l^{-1}\prod_{k\ne l}a_k\bigr)\norm{\nabla\psi(f)}_{\rms,*}\,\norm{x}_\rms$.
\end{enumerate}
For the average loss over several examples, (ii) and (iii) hold with the right sides averaged over the examples.
\end{proposition}

\begin{proof}
(i) Because $\phi$ is $1$-Lipschitz, $\abs{\phi'}\le1$ where it exists, so $\norm{\delta_l}_{\rms,*}=\norm{\phi'(u_l)\odot W_{l+1}^\top\delta_{l+1}}_{\rms,*}\le\norm{W_{l+1}^\top\delta_{l+1}}_{\rms,*}\le a_{l+1}\norm{\delta_{l+1}}_{\rms,*}$ by \cref{lem:modular:rms}(iii); iterate down from $\delta_N=\nabla\psi(f)$, and use $\norm{v}_\rms=\norm{v}_{\rms,*}/m$ for $v\in\R^m$. (ii) The identity is \cref{lem:modular:rms}(iv); bound $\norm{\delta_l}_{\rms,*}$ by (i) and $\norm{h_{l-1}}_\rms$ by \cref{thm:modular:spectralcondition}(i). (iii) Insert (ii) into \cref{thm:modular:dual}(i). For several examples, the gradient is the average of the single-example gradients, and the triangle inequality for the dual norms gives the averaged bounds.
\end{proof}

For the square loss in RMS units, $\psi(f)=\frac12\norm{f-y}_\rms^2$, the quantity $\norm{\nabla\psi(f)}_{\rms,*}$ is the RMS residual $\norm{f-y}_\rms$. Part (i) says that under the spectral condition every backpropagated signal is at most of order $1/d_l$ per entry, and (iii) that the loss is Lipschitz in the max norm with a constant that does not grow with width: the dual norm of the gradient, which a normalized steepest-descent step converts into a guaranteed decrease (\cref{ch:steepest}), stays bounded. That is the Lipschitz half of the analysis of \textcite{LargeEtAl2024}, for this network. Upper bounds are all that operator norms give; that the bounds are attained at a random initialization is a heuristic about random matrices, and it is what the width rules need.

\begin{heuristic}[Heuristic derivation of the width exponents]
Fix the depth, $d_0$, $d_N$, an example with $\norm{x}_\rms=1$ and $\norm{\nabla\psi(f)}_{\rms,*}$ of order one, let all hidden widths equal $\modularh\to\infty$, and let the initial weights of layer $l$ be independent $\cN(0,v_l^2)$. For such a matrix $W\in\R^{\modularout\times\modularin}$ and vectors $z\in\R^{\modularin}$, $z'\in\R^{\modularout}$ that do not depend on it, the law of large numbers gives $\norm{Wz}_2\approx v\sqrt{\modularout}\norm{z}_2$ and $\norm{W^\top z'}_2\approx v\sqrt{\modularin}\norm{z'}_2$, while $\norm{W}_{\op}\approx v(\sqrt{\modularin}+\sqrt{\modularout})$. In the units of this chapter the forward gain $\norm{Wz}_\rms/\norm{z}_\rms\approx v\sqrt{\modularin}$ and the backward gain $\norm{W^\top z'}_{\rms,*}/\norm{z'}_{\rms,*}\approx v\,\modularin/\sqrt{\modularout}$ are within a factor $1+\sqrt{\modularin/\modularout}$, respectively $1+\sqrt{\modularout/\modularin}$, of $\norm{W}_{\rms\to\rms}\approx v\sqrt{\modularin}\,(1+\sqrt{\modularin/\modularout})$. The forward pass traverses the input and hidden layers, which have $\modularin\le\modularout$, and the backward pass the hidden and output layers, which have $\modularin\ge\modularout$; so both bounds of \cref{thm:modular:spectralcondition}(i) and \cref{prop:modular:backward}(i) are attained up to constants. ReLU halves the squared norms, and initializations with $v_l^2=2/d_{l-1}$ compensate.

So if every layer has $\norm{W_l}_{\rms\to\rms}=\Theta(1)$, which for the output layer means entries of size $1/\modularh$ (variance $1/\text{fan-in}^2$, $\mu$P's choice), then $\norm{h_l}_\rms=\Theta(1)$ and $\norm{\delta_l}_\rms=\Theta(1/\modularh)$. In the standard parametrization the output layer has entries of size $\modularh^{-1/2}$, so $\norm{W_N}_{\rms\to\rms}=\Theta(\modularh^{1/2})$, and the backward gain through it raises the signals of all lower layers to $\norm{\delta_l}_\rms=\Theta(\modularh^{-1/2})$. Now insert these sizes, and the densities $\rho,\sigma=\Theta(1)$ of Gaussian and ReLU features, into \cref{thm:modular:widthrules}. The step sizes that move every layer's output by $\Theta(1)$ are those of \cref{tab:modular:rules}: $\eta_l\propto d_l/d_{l-1}$ for SGD, $\eta_l\propto1/d_{l-1}$ for sign descent, and $\eta_l\propto1$ for normalized spectral steps. In the standard parametrization one learning rate $\eta$ moves the input, hidden and output layers by $\Theta(\eta d_0\modularh^{-1/2})$, $\Theta(\eta\modularh^{1/2})$ and $\Theta(\eta\modularh)$ under SGD, and by $\Theta(\eta)$, $\Theta(\eta\modularh)$ and $\Theta(\eta\modularh)$ under sign descent. Keeping the output stable requires $\eta=O(1/\modularh)$ in both cases, and then the hidden layers of SGD and the input layer of sign descent learn nothing in the limit.
\end{heuristic}

\begin{table}[t]
\centering\small
\caption{Width exponents of the per-layer step size that moves every layer's output by $\Theta(1)$, for one input layer of width $\modularh\times d_0$, hidden layers $\modularh\times\modularh$ and an output layer $d_N\times\modularh$, with $d_0$ and $d_N$ fixed. Derived from \cref{thm:modular:widthrules} and the heuristic sizes of $\delta_l$; the sign-descent column is the $\mu$P rule for Adam, the SGD column the $\mu$P rule for SGD in a parametrization without multipliers. The last row is the standard parametrization with one learning rate $\eta$.}\label{tab:modular:rules}
\renewcommand{\arraystretch}{1.15}
\begin{tabular}{@{}lccc@{}}
\toprule
& SGD & sign descent (Adam) & normalized spectral\\
\midrule
input layer & $\eta_1\propto\modularh$ & $\eta_1\propto1$ & $\eta_1\propto1$\\
hidden layers & $\eta_l\propto1$ & $\eta_l\propto1/\modularh$ & $\eta_l\propto1$\\
output layer & $\eta_N\propto1/\modularh$ & $\eta_N\propto1/\modularh$ & $\eta_N\propto1$\\
general rule & $d_l/d_{l-1}$ (fan-out over fan-in) & $1/d_{l-1}$ (one over fan-in) & $1$, with $\sqrt{d_l/d_{l-1}}$ inside $\lmo$\\
\midrule
SP, per-layer change & $\eta d_0\modularh^{-1/2},\ \eta\modularh^{1/2},\ \eta\modularh$ & $\eta,\ \eta\modularh,\ \eta\modularh$ & ---\\
\bottomrule
\end{tabular}
\end{table}

The sign-descent column is the Adam column of $\mu$P \parencite{YangHu2021,YangEtAl2021}, which \textcite{YangLittwin2023} justify for adaptive optimizers in the infinite-width limit, and the shape factor in the spectral column is the one that \textcite{BernsteinNewhouse2025} derive by dualizing the RMS$\to$RMS norm, from which they recover $\mu$P. The derivation also explains why the three rules are one rule. Each multiplies the optimizer's natural step by the constant that converts its norm into the RMS$\to$RMS norm on aligned updates: for SGD the constant depends on the signal, for sign descent it is the fan-in, which is the norm-equivalence constant between the largest entry of a matrix and its RMS$\to$RMS norm (\cref{ex:modular:signnorm}), and for spectral steps it is one.

\subsection{Smoothness and depth}

The Lipschitz bound of \cref{prop:modular:backward} has a second-order companion. For a deep linear network and the square loss, the second derivative of $\ell$ along a direction $\Delta\theta$ with $\max_l\norm{\Delta W_l}_{\rms\to\rms}\le1$ is $\norm{\dot f}^2_\rms+\ip{f-y,\ddot f}_\rms$, where $\ip{a,b}_\rms=\ip{a,b}/d_N$, and the expansion in the proof of \cref{thm:modular:spectralcondition}(iv) bounds $\norm{\dot f}_\rms$ by $\sum_k\prod_{j\ne k}a_j\norm{x}_\rms$ and $\norm{\ddot f}_\rms$ by $2\sum_{j<k}\prod_{i\ne j,k}a_i\norm{x}_\rms$: on the set $\{\norm{W_l}_{\rms\to\rms}\le a_l\}$ the loss is smooth in the max norm with a constant that does not depend on the widths. \Textcite{LargeEtAl2024} prove width-independent Lipschitz and smoothness bounds of this kind for networks assembled from modules that satisfy their conditions of being well normed and sharp, with constants given by recursive formulas over the architecture. For nonlinear layers the conditions carry real content. In RMS units the second-order term of an entrywise nonlinearity, $\phi''(u)\odot\dot u\odot\dot u$, is bounded by $\norm{\phi''}_\infty\norm{\dot u}_\infty\norm{\dot u}_\rms$ but not by a multiple of $\norm{\dot u}_\rms^2$ that is independent of the width: a change concentrated on few coordinates is penalized more than its RMS norm shows (\cref{ex:modular:curvature}).

Depth is a separate scaling. In a residual network $h_l=h_{l-1}+N^{-\alpha}F_l(h_{l-1})$ whose branches satisfy $\norm{F_l(h)}_\rms\le a\norm{h}_\rms$, the argument of \cref{thm:modular:spectralcondition}(i) gives $\norm{h_N}_\rms\le(1+aN^{-\alpha})^N\norm{x}_\rms$, which is bounded uniformly in $N$ if and only if $\alpha\ge1$, whereas branch outputs orthogonal to the stream give $\norm{h_N}_\rms^2\le(1+a^2N^{-2\alpha})^N\norm{x}^2_\rms$, bounded for $\alpha\ge1/2$ (\cref{ex:modular:depth}). Both exponents are in use. \Textcite{BordelonEtAl2024} and \textcite{YangEtAl2024} scale the residual branches and their updates by $1/\sqrt N$ (Depth-$\mu$P) and observe hyperparameter transfer across depth, while \textcite{DeyEtAl2025} argue for $1/N$ (CompleteP), which they find necessary for every layer to learn features beyond its linearization; it requires rescaling the learning rates of normalization layers and biases, AdamW's weight decay and $\epsilon$ as well. We state these results and do not prove them.

\section{Reading the theory}\label{sec:modular:reading}

The theory makes three kinds of prediction. (1) \emph{Per-layer exponents}: the learning rates of \cref{tab:modular:rules}, under which the optimal base learning rate should not depend on the width. (2) \emph{The failure of one learning rate}: in the standard parametrization the output layer, and for Adam also the hidden layers, move by $\Theta(\eta\modularh)$ per step, so the largest stable and, in short training, the best learning rate scale like $1/\modularh$, a slope of $-1$ on a $\log$--$\log$ plot, for SGD and for Adam alike. (3) \emph{The size of one step}: at a fixed base learning rate the change of the hidden features grows like $\modularh^{1/2}$ under SGD in the standard parametrization and like $\modularh$ under sign descent, the first step of Adam, and is width-independent under the three width-aware rules; for normalized spectral steps it lies between the two bounds of \cref{thm:modular:spectralcondition}. Leading constants follow from the Gaussian computations of the heuristic: for the hidden layer of the network of \cref{fig:modular:curves}, with $f-y=1$ and $\norm{x}_\rms=1$, the direct change per unit of step is $\sqrt{\modularh/2}$ for SGD and $\modularh/\sqrt{2\pi}$ for sign descent in the standard parametrization, and $\sqrt{32}$ and $64/\sqrt{2\pi}$ in $\mu$P with base width $64$ (\cref{ex:modular:constants}).

\begin{figure}[t]
\begin{fullwidth}
\centering
\begin{tikzpicture}
\begin{groupplot}[group style={group size=2 by 2,horizontal sep=17mm,vertical sep=17mm},wplot,width=80mm,height=52mm,
  xmode=log,log basis x=2,xtick={64,128,256,512,1024,2048,4096},xticklabels={$2^6$,$2^7$,$2^8$,$2^9$,$2^{10}$,$2^{11}$,$2^{12}$},
  xlabel={hidden width $\modularh$}]
\nextgroupplot[title={(a) optimal learning rate},xmin=54,xmax=1210,ymin=-5.6,ymax=1.3,
  ylabel={$\log_2\bigl(\eta^\star(\modularh)/\eta^\star(64)\bigr)$},
  legend style={at={(0.02,0.03)},anchor=south west,font=\scriptsize},legend columns=2]
\addplot[guide,domain=64:1024] {0};
\addplot[guide,domain=64:1024] {-ln(x/64)/ln(2)};
\addplot[simulation2] table[x=n,y=adamsprel] {figures/data/04-modular-transfer-optima.dat};
\addplot[only marks,mark=square,mark size=1.6pt,color=seadeep] table[x=n,y=adammuprel] {figures/data/04-modular-transfer-optima.dat};
\addplot[only marks,mark=diamond*,mark size=2.0pt,color=sea] table[x=n,y=specrel] {figures/data/04-modular-transfer-optima.dat};
\addplot[only marks,mark=o,mark size=1.7pt,color=sandline!55!black] table[x=n,y=sgdmuprel] {figures/data/04-modular-transfer-optima.dat};
\addplot[simulation] table[x=n,y=sgdspstablerel] {figures/data/04-modular-transfer-optima.dat};
\legend{,,Adam SP,Adam $\mu$P,spectral,SGD $\mu$P,SGD SP${}^{\dagger}$}
\nextgroupplot[title={(b) one SGD step},xmin=54,xmax=4800,ymode=log,ymin=3,ymax=80,
  ylabel={$\norm{\Delta u_2}_\rms/\eta$},legend style={at={(0.02,0.97)},anchor=north west,font=\scriptsize}]
\addplot[theory,domain=64:4096,samples=40] {sqrt(x/2)};
\addplot[theory2,domain=64:4096] {sqrt(32)};
\addplot[simulation] table[x=n,y=sgdsp] {figures/data/04-modular-onestep.dat};
\addplot[only marks,mark=o,mark size=1.7pt,color=sandline!55!black] table[x=n,y=sgdmup] {figures/data/04-modular-onestep.dat};
\legend{$\sqrt{\modularh/2}$,$\sqrt{32}$,SP,$\mu$P}
\nextgroupplot[title={(c) one sign-descent (first Adam) step},xmin=54,xmax=4800,ymode=log,ymin=15,ymax=3000,
  ylabel={$\norm{\Delta u_2}_\rms/\eta$},legend style={at={(0.02,0.97)},anchor=north west,font=\scriptsize}]
\addplot[theory,domain=64:4096] {x/sqrt(2*pi)};
\addplot[theory2,domain=64:4096] {64/sqrt(2*pi)};
\addplot[simulation2] table[x=n,y=adamsp] {figures/data/04-modular-onestep.dat};
\addplot[only marks,mark=square,mark size=1.6pt,color=seadeep] table[x=n,y=adammup] {figures/data/04-modular-onestep.dat};
\legend{$\modularh/\sqrt{2\pi}$,$64/\sqrt{2\pi}$,SP,$\mu$P}
\nextgroupplot[title={(d) one normalized spectral step},xmin=54,xmax=4800,ymin=0,ymax=5.6,
  ylabel={$\norm{\Delta u_2}_\rms/\eta$},legend style={at={(0.98,0.5)},anchor=east,font=\scriptsize}]
\addplot[theory3] table[x=n,y=spechi] {figures/data/04-modular-onestep.dat};
\addplot[only marks,mark=diamond*,mark size=2.0pt,color=sea] table[x=n,y=spec] {figures/data/04-modular-onestep.dat};
\addplot[theory] table[x=n,y=speclo] {figures/data/04-modular-onestep.dat};
\legend{upper bound $a_1+a_2$,measured,lower bound $\norm{h_1}_\rms$}
\end{groupplot}
\end{tikzpicture}
\end{fullwidth}
\caption{Width rules, measured. (a) Optimal base learning rate relative to width $64$, from the sweeps of \cref{fig:modular:curves} (vertex of a parabola through the best grid point and its neighbours in $(\log_2\eta,\log\text{loss})$); dashed lines: the predicted slopes $0$ and $-1$. ${}^{\dagger}$Under SGD in the standard parametrization the loss hardly depends on the learning rate (the network barely learns this task in $150$ steps), so the points show its largest stable grid learning rate instead. (b)--(d) One step from one example, $f-y=1$, $\norm{x}_\rms=1$, base step $10^{-5}$, of the same network at initialization; points: the change of the second hidden layer's pre-activations divided by the step, all three layers updated, mean of $8$ seeds. Curves in (b), (c): the hidden layer's own contribution predicted by \cref{thm:modular:widthrules} with the heuristic sizes of $\delta$; the gap above them is the input layer's contribution, which does not grow with width. (d): the bounds of \cref{thm:modular:spectralcondition} at first order, computed from the initial weights. Scripts \texttt{04-modular-transfer.py} and \texttt{04-modular-onestep.py}.}
\label{fig:modular:widthrules}
\end{figure}

\Cref{fig:modular:widthrules} compares the predictions with measurements on the network of \cref{fig:modular:curves}. Panel (a) shows prediction (2). The optimum of Adam in the standard parametrization moves by $-0.96$ octaves per doubling of the width (least-squares slope over the five widths, $-1.03$ from width $128$ on), against the predicted $-1$, and the largest stable learning rate of SGD in the standard parametrization halves exactly with every doubling from width $128$ on. The three width-aware rules keep the optimum within about one octave over a sixteenfold range of widths: slopes of $-0.21$ for Adam in $\mu$P, $-0.19$ for spectral steps and $+0.16$ for SGD in $\mu$P, the last noisy because its optimum sits at the edge of stability. The small negative drift of Adam in $\mu$P and of spectral steps is not predicted by the leading-order theory; it is a finite-width effect of the kind \textcite{VyasEtAl2023} and \textcite{EverettEtAl2024} report.

Panels (b)--(d) show prediction (3) with constants. When only the hidden layer steps, its measured change agrees with the predicted contribution to within $8\%$ at every width, and to within $1\%$ for widths of $2048$ and $4096$. When all layers step, the change is larger by the input layer's contribution, which is width-independent under every rule, so that under the standard parametrization the points approach the curves from above, with local slopes rising to $0.49$ (SGD) and $0.99$ (sign descent) at the largest widths, while under $\mu$P they stay a constant factor above them. For normalized spectral steps the change is $2.4$ to $2.5$ at every width, between the lower bound $\norm{h_1}_\rms\approx1$, which is the hidden layer's exact contribution, and the upper bound $a_1+a_2\approx4.4$.

\paragraph{Where the agreement must fail.} Prediction (3) is a statement about one step at initialization, and it is accurate there. Prediction (2) relies on more: that the best learning rate after $150$ steps is limited by the same layers that limit the first step, and that later gradients stay aligned with the features. Minibatch gradients have rank up to $32$ here, Adam's moments mix steps, and the alignment of updates with features decays as training proceeds \parencite{KossonEtAl2026}; at a fixed learning rate the change of features then grows more slowly with width than the rank-one calculation says, which would flatten the slope of (a). In this experiment the slope survives because training is short. At small widths the exponents are blurred by lower-order terms, such as the factor $1+\sqrt{d_0/\modularh}$ in the norm of the input layer and the input layer's contribution in (b)--(c).

\paragraph{Muon's shape factors.} The factor $\sqrt{\modularout/\modularin}$ is the RMS$\to$RMS normalization of the polar factor exactly, and with it a Muon step has $\norm{\Delta W}_{\rms\to\rms}=\eta$ for every shape. Jordan et al.'s clamp, $\sqrt{\max(1,\modularout/\modularin)}$, agrees for $\modularout\ge\modularin$ and gives a wide-to-narrow matrix the larger norm $\sqrt{\modularin/\modularout}$, a constant factor for a fixed architecture shape. Liu et al.'s factor $0.2\sqrt{\max(\modularout,\modularin)}$ gives $\norm{\Delta W}_{\rms\to\rms}=0.2\sqrt{\modularin}$ for $\modularout\ge\modularin$, which grows with width: under it, the theory predicts, the base learning rate must shrink like $\modularh^{-1/2}$ to transfer (\cref{ex:modular:muonscale}). RMS matching serves a different purpose, reusing AdamW's tuned learning rate at a given width, and there it does what it was designed to do. For a square matrix and a full-rank polar factor, every input receives the gain $\norm{\Delta W}_{\rms\to\rms}$ exactly, aligned or not: the normalized spectral step is the one step whose effect does not depend on the alignment assumptions that the other two rules need.

\section{What the theory misses}\label{sec:modular:misses}

\paragraph{What was idealized away.} \emph{One example and one step.} Minibatch gradients are sums of rank-one terms, with stable rank well below the batch size but above one, and Adam's moments and Muon's momentum mix many steps; once updates are not aligned with the inputs they are applied to, $\norm{\Delta Wh}_\rms$ falls below $\norm{\Delta W}_{\rms\to\rms}\norm{h}_\rms$, by a factor that depends on the update's spectrum and on the alignment, and the exponents of \cref{tab:modular:rules} change (\cref{ex:modular:batch}). \emph{Orders of magnitude in place of constants.} The theory fixes exponents; the factor of $2$ between the operator norm and the typical gain of a square Gaussian matrix, the densities $\rho$ and $\sigma$, and the base learning rate are outside it. \emph{The RMS norm itself.} It treats all coordinates alike, which is right for delocalized features and wrong for the few large activations that trained transformers develop, and it controls entrywise nonlinearities only to first order. \emph{The architecture.} Embeddings receive one-hot inputs, for which the right input norm is $\ell_1$ and not RMS (\cref{ex:modular:embedding}); attention and normalization layers have no norm in this chapter. \emph{Everything that is not width.} Weight decay, the training horizon, the batch size and depth are fixed.

\paragraph{Where it breaks.} (1) The theory does not single out $\mu$P. \Textcite{EverettEtAl2024} found that all four parametrizations they tested transfer once each layer has the right exponent, and that a per-layer prescription for the standard parametrization beat $\mu$P at scale; they also found that the exponents depend on how aligned the updates are, which is a property of training and not of the architecture. (2) The alignment that the rank-one calculation assumes is transient. \Textcite{KossonEtAl2026} measured it in language-model training and found that it holds only at the start, after which the relative size of the updates is set by weight decay, and they argue that $\mu$P acts mainly as an implicit warmup. (3) Optimal learning rates also depend on the length of training: \textcite{BjorckEtAl2025} find that the best learning rate of a language model shrinks as the token horizon grows, following a scaling law of its own, about which the norm of the network is silent. (4) Weight decay interacts with the width rules: with AdamW, $\mu$P's scaling of the learning rate breaks down during training unless the weight decay is scaled too \parencite{WangAitchison2025}. (5) In the ablations of \textcite{Lingle2024}, transfer failed with trainable gains in the normalization layers, with Lion, with decoupled weight decay, and with the standard $1/\sqrt{d}$ scaling of attention logits, all components outside this chapter's model. (6) Depth needs its own rule, and the two candidate exponents disagree about whether hyperparameters transfer across depth at all \parencite{YangEtAl2024,BordelonEtAl2024,DeyEtAl2025}. (7) Finally, the theory predicts how the optimum scales, not where it is, and not why it transfers: the optimum's location is a property of the loss landscape that the norm does not see.

\paragraph{Repairs and open problems.} The landscape explanation has been taken up directly. \Textcite{NociEtAl2024} found that under $\mu$P the sharpness along the training trajectory is nearly independent of width and depth, which they call super consistency, while under the neural-tangent parametrization it changes with width, and they link the presence or absence of super consistency to the presence or absence of learning-rate transfer; in the language of \cref{ch:eos}, the optimal learning rate sits near the stability threshold, which then does not move. For the components that break transfer, u-$\mu$P \parencite{BlakeEtAl2025} combines $\mu$P with unit scaling, gives near-optimal defaults, and trains in 8-bit floating point; Depth-$\mu$P and CompleteP extend the rules to depth; and the scaling of weight decay with width is the subject of \cref{ch:decay} \parencite{WangAitchison2025,KossonEtAl2026}. A more radical repair fixes the norms instead of scaling the steps: Hyperball \parencite{WenEtAl2026} constrains the Frobenius norm of every matrix and normalizes every update, so that the learning rate sets an angular step, and reports better transfer across widths and depths in an unrefereed study of models up to $1.2$ billion parameters. The spectral-condition derivation has been extended to LAMB, Sophia, Shampoo and Muon \parencite{GuptaEtAl2026}, and per-layer rules for matrix-preconditioned optimizers have been worked out by \textcite{QiuEtAl2025}. Open are a theory of where the optimum lies, not only how it scales; norms for attention, embeddings and normalization layers derived rather than chosen; the dynamics of alignment over training, which decides the exponents in practice; and transfer across the training horizon and the batch size, which no norm on the weights can see.

\section{Back in the wild}\label{sec:modular:back}

\fieldnotes{Tune on a narrow proxy and widen: in $\mu$P, Adam's hidden and output learning rates are divided by the width multiplier; Muon's update carries $\sqrt{\modularout/\modularin}$ or an RMS-matching factor; Modula normalizes every module.}
{A network as a chain of layers, each with an operator norm (RMS$\to$RMS for linear layers), the network normed by the maximum; one step from one example, so rank-one gradients; features of RMS $\Theta(1)$.}
{Steepest descent in the max norm is per-layer normalized descent; exact one-step feature changes for SGD, sign descent and spectral steps; width-free forward, update and backward bounds under the spectral condition; no cancellation between layers at first order.}
{Per-layer rates $\propto$ fan-out/fan-in (SGD), $1/\text{fan-in}$ (Adam), $1$ (spectral); in SP the optimum scales like $1/\text{width}$; one-step feature changes $\propto\modularh^{1/2}$, $\modularh$ and $1$ (measured: slopes $-0.96$ and about $0$ for the optimum).}
{Alignment after the first steps; constants and the optimum's location; embeddings, attention, normalization; weight decay, training horizon, batch size; depth.}
{Measure a step by its RMS$\to$RMS norm and scale each layer's learning rate so that this norm is width-independent: proved in the idealization to give width-independent feature changes. Expect the optimum's location to transfer across width, not across horizon or batch size (observed, not proved).}

What the chapter tells a practitioner is a units check rather than a recipe. A learning rate is a number only after a norm has been chosen; the RMS$\to$RMS norm is the one in which ``the same learning rate'' means ``the same change of the features'' at every width, and every working width rule, $\mu$P's $1/\text{fan-in}$ for Adam, Muon's $\sqrt{\modularout/\modularin}$, the per-layer exponents of Everett et al., is a way of expressing that norm in the optimizer's own currency. The theory does not tell anyone which base learning rate to use, nor that a proxy tuned for one training horizon is tuned for another; the belief that a $\mu$P proxy can stand in for every hyperparameter of a large run goes beyond the evidence in \cref{sec:modular:misses}.

The chapter closes Part~I. \Cref{ch:steepest} gave the dual of a maximum of norms in the abstract; \cref{ch:sign,ch:spectral} computed the duality maps of the entrywise and spectral norms of one tensor; here the layers are assembled into one norm, and the per-layer learning rate turns out to be part of the norm. The memory that practical optimizers add to these steps is the subject of \cref{ch:momentum,ch:adaptive,ch:kronecker}; the per-layer rate returns as the effective learning rate of a normalized layer under weight decay in \cref{ch:decay}, and the stability threshold that super consistency keeps fixed is the subject of \cref{ch:eos}.

\begin{exercises}
\item \consolidation (Sums and other combinations.) (a) Show that the dual of $\norm{\theta}=\sum_ls_l\norm{W_l}_{(l)}$ is $\max_l\norm{G_l}_{(l),*}/s_l$, and that $\lmo(g)$ is the set of $u$ that vanish outside the layers attaining the maximum and are convex combinations of their scaled $\lmo$'s. (b) For $1<p<\infty$, show that the dual of $\bigl(\sum_l(s_l\norm{W_l}_{(l)})^p\bigr)^{1/p}$ is $\bigl(\sum_l(\norm{G_l}_{(l),*}/s_l)^q\bigr)^{1/q}$ with $1/p+1/q=1$, and find its $\lmo$. (c) For two layers with $\norm{G_1}_{(1),*}/s_1=2\norm{G_2}_{(2),*}/s_2$, compare the steps of the max, sum and $\ell_2$ combinations.\label{ex:modular:sum}

\item \consolidation (Sign descent and the largest entry.) (a) Show that the $\ell_1\to\ell_\infty$ operator norm of $W$ is $\norm{W}_{\max}=\max_{ij}\abs{W_{ij}}$, so that sign descent is normalized steepest descent in it. (b) Show that $\norm{W}_{\rms\to\rms}\le\modularin\norm{W}_{\max}$, with equality if and only if $W=c\,ab^\top$ for sign vectors $a\in\{\pm1\}^{\modularout}$, $b\in\{\pm1\}^{\modularin}$. (c) Deduce the $\mu$P rule $\eta\propto1/\text{fan-in}$ for sign descent as the statement that $\eta\norm{\cdot}_{\max}$-steps have RMS$\to$RMS norm at most $\eta\,\modularin$, with equality for single-example steps on dense inputs.\label{ex:modular:signnorm}

\item \consolidation (Embeddings.) An embedding layer $W\in\R^{\modularout\times\modularin}$ receives one-hot inputs $x=e_j$, so $\norm{x}_1=1$ while $\norm{x}_\rms=\modularin^{-1/2}$. (a) Compute $\norm{\Delta We_j}_\rms$ for one step of SGD, of sign descent and of the normalized RMS$\to$RMS step, and show that the last is $\eta\,\modularin^{-1/2}$: the spectral norm is the wrong norm for a vocabulary. (b) Show that the $\ell_1\to\rms$ operator norm is the largest column RMS, $\max_j\norm{W_{:j}}_\rms$, that its dual is $\sqrt{\modularout}\sum_j\norm{G_{:j}}_2$, and that its $\lmo$ normalizes the nonzero columns of $-G$ to RMS one. (c) Show that its normalized step moves the embedding of token $j$ by exactly $\eta$ in RMS, and that the width exponents of the input-layer row of \cref{tab:modular:rules} follow for one-hot inputs.\label{ex:modular:embedding}

\item \consolidation (Tightness.) Let all layers of a deep linear network be square and orthogonal, $\norm{x}_\rms=1$, and let each layer take the normalized RMS$\to$RMS step of size $\eta$ of one example's loss. (a) Show that $\delta_k=W_{k+1}^\top\cdots W_l^\top\delta_l$ and that the contribution of layer $k$ to $\dot u_l$ is $-\eta\sqrt{d_l}\,\hat\delta_l$. (b) Deduce $\norm{\dot u_l}_\rms=l\eta$, the first-order term of the upper bound of \cref{thm:modular:spectralcondition}(ii). (c) Check numerically that $\norm{\Delta u_l}_\rms/\eta\to l$ as $\eta\to0$.\label{ex:modular:tight}

\item \consolidation (Constants.) For the network of \cref{fig:modular:curves} with He initialization, $f-y=1$ and $\norm{x}_\rms=1$, show heuristically that $\norm{h_1}_\rms\approx1$, $\rho(h_1)\approx1/\sqrt\pi$, $\sigma(h_1)\approx\sigma(\delta_2)\approx1/\sqrt2$, and $\norm{\delta_2}_\rms\approx(2\modularh)^{-1/2}$ in the standard parametrization and $\sqrt{32}/\modularh$ in $\mu$P with base width $64$. Deduce the four curves of \cref{fig:modular:widthrules}(b)--(c).\label{ex:modular:constants}

\item \variation (Minibatches.) Let $G=\frac1B\sum_{i=1}^B\delta_ix_i^\top$ with orthogonal inputs of equal norm and orthogonal signals. (a) Show that SGD moves the layer's output on $x_j$ by $\frac\eta B\norm{x_j}_2^2\,\delta_j$, a factor $B$ less than the single-example step of the same learning rate. (b) Show that $\polar(G)=\sum_i\hat\delta_i\hat x_i^\top$ and that the normalized spectral step moves the output on every $x_j$ by exactly $\eta\norm{x_j}_\rms$, as for one example. (c) For sign descent with $B=2$ and Gaussian vectors, estimate $\norm{\Delta Wx_1}_\rms$ and compare.\label{ex:modular:batch}

\item \variation (Depth.) Let $h_l=h_{l-1}+N^{-\alpha}F_l(h_{l-1})$, $l=1,\dots,N$, with $\norm{F_l(h)}_\rms\le a\norm{h}_\rms$. (a) Show that $\norm{h_N}_\rms\le(1+aN^{-\alpha})^N\norm{x}_\rms\le e^{aN^{1-\alpha}}\norm{x}_\rms$, uniformly bounded in $N$ for $\alpha\ge1$. (b) If every branch output is orthogonal to its input, show that $\norm{h_N}^2_\rms\le(1+a^2N^{-2\alpha})^N\norm{x}_\rms^2$, bounded for $\alpha\ge1/2$. (c) Relate the two exponents to Depth-$\mu$P and CompleteP.\label{ex:modular:depth}

\item \variation (Muon's shape factors.) For an update $\gamma\polar(M)$ of a $\modularout\times\modularin$ matrix of full rank, compute $\norm{\gamma\polar(M)}_{\rms\to\rms}$ and its entrywise RMS for $\gamma=\sqrt{\modularout/\modularin}$, $\gamma=\sqrt{\max(1,\modularout/\modularin)}$ and $\gamma=0.2\sqrt{\max(\modularout,\modularin)}$. Show that with the third factor all updates have entrywise RMS $0.2$, and that a hidden $\modularh\times\modularh$ matrix needs a base learning rate $\propto\modularh^{-1/2}$ for width-independent feature changes. What does the second factor do to the down-projection of a transformer's MLP, of shape $d\times4d$?\label{ex:modular:muonscale}

\item \variation (Weight decay under $\mu$P.) With decoupled weight decay, $W\leftarrow(1-\eta_l\lambda_l)W-\eta_lU$, the weights are an exponential moving average of the updates with time scale $1/(\eta_l\lambda_l)$ steps (\cref{ch:decay}). Show that keeping this time scale fixed while $\eta_l\propto1/\modularh$ requires $\lambda_l\propto\modularh$, and compare with the advice of \textcite{WangAitchison2025} that the optimal weight decay grows with model size under $\mu$P learning rates.\label{ex:modular:decay}

\item \variation (Curvature of a nonlinearity.) Let $\phi=\tanh$ and $u=0$. (a) For $\dot u=\sqrt m\,e_1\in\R^m$, which has $\norm{\dot u}_\rms=1$, compute $\norm{\phi''(u+t\dot u)\odot\dot u\odot\dot u}_\rms$ at a small $t\ne0$ and show that it grows like $\sqrt m$. (b) For $\dot u$ with entries $\pm1$, show that it stays bounded. (c) Conclude that a width-independent smoothness constant in RMS units needs a condition on $\norm{\dot u}_\infty$, and discuss which updates of \cref{thm:modular:widthrules} satisfy it.\label{ex:modular:curvature}

\item \simulation (Alignment over training.) Modify \texttt{04-modular-transfer.py} to record, for the hidden matrix at widths $64$ to $1024$, the alignment ratio $\norm{\Delta W_2h_1}_\rms/(\norm{\Delta W_2}_{\rms\to\rms}\norm{h_1}_\rms)$ averaged over the batch, and the stable rank $\norm{\Delta W_2}_F^2/\norm{\Delta W_2}_{\op}^2$ of the update, at steps $1,10,50,150$, for Adam in $\mu$P and for spectral steps. Observe that the ratio of Adam's update starts near the rank-one value and decays, and explain why the spectral update's ratio is bounded below by the fraction of the input that lies in its row space.\label{ex:modular:align}

\item \simulation (Breaking transfer.) In \texttt{04-modular-transfer.py}: (a) replace the factor $\sqrt{\modularout/\modularin}$ of the spectral steps by $0.2\sqrt{\max(\modularout,\modularin)}$ and measure the slope of the optimum (prediction: $-1/2$); (b) give the output layer of the $\mu$P configurations the standard initialization and observe the effect on the optimum and on the loss; (c) double the number of steps and check whether the slopes of \cref{fig:modular:widthrules}(a) change.\label{ex:modular:break}

\item \research (A norm for attention.) Self-attention is not a linear layer: the logits are quadratic in the query and key weights and the softmax is Lipschitz only in a range of logits. Propose a norm on $(W_Q,W_K,W_V,W_O)$ under which one normalized step changes the attention output by a width-independent amount, derive the scaling of the logits it implies (the $1/d$ versus $1/\sqrt d$ question of \textcite{Lingle2024}), and test it against the alignment measurements of \cref{ex:modular:align}. As far as we know, no derivation of the attention norm from the feature-learning desideratum alone is accepted.\label{ex:modular:attention}

\item \research (Exponents beyond alignment.) Suppose that after $t$ steps the update of a hidden layer has stable rank $r(t,\modularh)$ and alignment ratio $\kappa(t,\modularh)$ with the features. Derive the width exponents of \cref{tab:modular:rules} as functions of the growth rates of $r$ and $\kappa$, and compare with the exponents measured by \textcite{EverettEtAl2024}. A model in which $r$ and $\kappa$ are derived rather than measured is open.\label{ex:modular:research}
\end{exercises}

\notesandrefs
The scale of initial weights was tied to the fan-in long before width became a variable: \textcite{GlorotBengio2010} chose variances that keep activations and gradients of constant size through depth. Wide networks became a theory with the neural tangent kernel \parencite{JacotGabrielHongler2018}, in which features do not move, and with the mean-field limit of two-layer networks \parencite{MeiMontanariNguyen2018}, in which they do; \textcite{ChizatOyallonBach2019} showed that the difference is a matter of scaling. \Textcite{YangHu2021} classified the parametrizations of networks of every depth into those with a kernel limit and those that learn features, singled out $\mu$P, and extended the program to adaptive optimizers with \textcite{YangLittwin2023}. $\mu$Transfer \parencite{YangEtAl2021} made the classification a tuning procedure, used, for example, for the Cerebras-GPT models \parencite{DeyEtAl2023}.

Per-layer normalized steps came first from large-batch training: LARS \parencite{YouGitmanGinsburg2017} and LAMB \parencite{YouEtAl2020} rescale each layer's step to a fixed fraction of its weights. \Textcite{BernsteinEtAl2020} derived such relative steps from a bound on the change of a deep network's function, deep relative trust, and \textcite{Flynn2017} analysed descent with layerwise norms. \Textcite{YangSimonBernstein2023} replaced the tensor-program derivation of $\mu$P by the spectral condition, a few lines of linear algebra, which \cref{thm:modular:spectralcondition} follows; \textcite{LargeEtAl2024} built the modular norm for general architectures and the Modula library, and \textcite{BernsteinNewhouse2024,BernsteinNewhouse2025} wrote the optimizers in use as steepest descent in such norms and computed the duality maps of linear, embedding and convolutional layers. \Textcite{Newhouse2025} presents this line in a thesis. Muon's shape factor is in \textcite{JordanEtAl2024}, and the RMS-matching variant in \textcite{LiuEtAl2025}. \Textcite{PethickEtAl2025} develop optimizers built on linear minimization oracles over norm balls, with a norm for deep architectures chosen layer by layer for transfer across model sizes.

\Cref{lem:modular:nocancel} and its use in \cref{thm:modular:spectralcondition}(iii) are, as far as we know, not stated in this form in the literature, although the alignment of the steps of different layers is implicit in the arguments of \textcite{YangSimonBernstein2023}. The constants of \cref{fig:modular:widthrules}(b)--(c) are our computation. The empirical record on transfer is large and moving fast; the studies of \textcite{EverettEtAl2024}, \textcite{Lingle2024}, \textcite{KossonEtAl2026} and \textcite{BjorckEtAl2025} are the place to start, and depth is treated in \textcite{BordelonEtAl2024}, \textcite{YangEtAl2024} and \textcite{DeyEtAl2025}.

\paragraph{Translating the literature.} \Textcite{YangSimonBernstein2023} write $n_\ell$ for the width of layer $\ell$ (our $d_l$) and $\norm{\cdot}_*$ for the \emph{spectral} norm, which is our $\norm{\cdot}_{\op}$; in this book a star always marks a dual norm. Their Desideratum~1, $\norm{h_\ell}_2=\Theta(\sqrt{n_\ell})$ and $\norm{\Delta h_\ell}_2=\Theta(\sqrt{n_\ell})$, is our $\norm{h_l}_\rms,\norm{\Delta h_l}_\rms=\Theta(1)$, and their Condition~1 is our $\norm{W_l}_{\rms\to\rms},\norm{\Delta W_l}_{\rms\to\rms}=\Theta(1)$. \Textcite{YangHu2021} write each weight as a width-dependent multiplier times a trainable matrix whose initialization and learning rate also scale with the width $n$ (an ``abc-parametrization''); multiplier, initialization and learning rate can be traded for one another without changing the dynamics, which is why the same dynamics appears in different tables, and our \cref{tab:modular:rules} puts all width dependence into the initialization and the per-layer learning rates. ``Fan-in'' and ``fan-out'' are our $d_{l-1}$ and $d_l$, the ``width multiplier'' is $\modularh$ divided by the proxy's width. \Textcite{LargeEtAl2024} and \textcite{BernsteinNewhouse2025} call the map $g\mapsto\argmax_{\norm t=1}\ip{g,t}$ the duality map, which is $-\lmo(g)$ here, and the scales of our max norm are computed in their framework from the masses and sensitivities of the modules. \Textcite{PethickEtAl2025} take the linear minimization oracle over a norm ball of radius $\rho$, which is $\rho$ times our $\lmo$. Muon's documentation writes the shape of a weight matrix as $A\times B$, our $\modularout\times\modularin$.

\printchapterbib
\end{refsection}


\casepart{Memory: Averages and Preconditioners}{Practical optimizers remember. A moving average of gradients filters noise, and it is what lets normalized steps converge; a moving average of squared gradients, or of products of gradients, builds a norm adapted to the problem. Momentum, Adam, Shampoo and SOAP are such memories added to steepest descent.}{36}
% Chapter 5 of Optimizers in the Wild: moving averages of gradients and of iterates.
\newcommand{\momentummed}{\operatorname{med}}
\newcommand{\momentumzbar}{\bar\zeta}

\casechapter{Moving Averages}{Momentum, noise, and why normalized steps need memory}{Momentum $0.9$}{ch:momentum}
\begin{refsection}[chapters/05-momentum.bib]

\begin{chaptersummary}
Almost every optimizer in use averages its gradients with a coefficient near $0.9$, and every normalized optimizer, from sign descent with momentum to Lion and Muon, has the average built in. We study the average on a quadratic with exact gradients, where it accelerates and has stability regions of its own (among them the threshold $38$ behind Adam's edge of stability); on a quadratic with additive noise, where heavy-ball momentum is, to first order, a larger step size; and for normalized steepest descent in a general norm with stochastic gradients. The main theorem bounds the gradient of normalized steepest descent with momentum in every norm at the rate $T^{-1/4}$, and a matching example shows that without memory a normalized step can stall at a gradient of the size of the noise, whatever its step size. Simulations separate the bias of a normalized step, which memory removes, from its fluctuations, which memory enlarges at a fixed step size; none of this explains why $0.9$ in particular works so widely.
\end{chaptersummary}

\section{In the wild: momentum \texorpdfstring{$0.9$}{0.9}, almost everywhere}\label{sec:momentum:wild}

The oldest ingredient of a modern optimizer is also the most uniform. Gradient descent with a memory of its past steps was introduced by \textcite{Polyak1964} as the \term{heavy-ball method}\index{heavy-ball method}, appeared in the training of neural networks together with back-propagation \parencite{RumelhartHintonWilliams1986}, and has been part of almost every widely used optimizer since. Its coefficient has barely moved.

\begin{practice}[ResNet, \textcite{HeEtAl2016}; Adam, \textcite{KingmaBa2015}; Lion, \textcite{ChenEtAl2023}; Muon, \textcite{JordanEtAl2024Muon}]
Average the gradients before stepping, with a coefficient of about $0.9$. Convolutional networks: SGD with momentum $0.9$ in the form $b_t=0.9\,b_{t-1}+g_t$, $\theta_{t+1}=\theta_t-\eta\,b_t$, with $\eta=0.1$ and step decay for ResNets on ImageNet. Transformers: Adam or AdamW with $\beta_1=0.9$, the moving average $m_t=0.9\,m_{t-1}+0.1\,g_t$ (GPT-3 and LLaMA, \cite{BrownEtAl2020,TouvronEtAl2023}). Lion: step along the sign of $0.9\,m_{t-1}+0.1\,g_t$ and update the state $m_t$ with coefficient $0.99$. Muon: Nesterov momentum with coefficient $0.95$, then orthogonalize. Separately, many runs keep a slow moving average of the weights for evaluation, or average the iterates in place of a decaying schedule.
\end{practice}

Two conventions hide in this list. PyTorch's SGD accumulates $b_t=\beta b_{t-1}+g_t$, whose weights sum to $1/(1-\beta)$, while Adam's first moment $m_t=\beta m_{t-1}+(1-\beta)g_t$ has weights that sum to $1-\beta^{t+1}$, and to one after Adam's bias correction. The two averages differ by the factor $1-\beta$: at $\beta=0.9$ a learning rate of $0.1$ in the first convention moves the weights as far as a learning rate of $1$ in the second. Which of the two is ``the'' learning rate matters as soon as $\beta$ changes, and it is the first thing the theory will sort out.\sidenote{PyTorch's \texttt{torch.optim.SGD} also has a ``dampening'' argument that multiplies $g_t$ by one minus its value; with dampening equal to $\beta$ it computes Adam's average. Most recipes leave it at zero.} Behind the practice are three regularities of different strength.

\begin{observation}[Momentum as a larger step; \textcite{KidambiEtAl2018}, \textcite{ShallueEtAl2019}, \textcite{ZhangEtAl2019}]\label[observation]{obs:momentum:rescaled}
With small batches, SGD with momentum $\beta$ and learning rate $\eta$ in PyTorch's convention trains about as well as plain SGD with learning rate $\eta/(1-\beta)$, once both are tuned. Momentum pays off at large batches, where it extends the range of batch sizes over which doubling the batch halves the number of steps to a target.
\end{observation}

\textcite{ShallueEtAl2019} measured the number of steps to a target as a function of the batch size for image and language workloads, and found that momentum extends the regime of perfect scaling to larger batches; \textcite{ZhangEtAl2019} reproduced the pattern in a noisy quadratic model, and \textcite{KidambiEtAl2018} showed that the classical momentum schemes do not improve on SGD for some stochastic least-squares problems. For small learning rates, \textcite{WangEtAl2024} explain the observation with stochastic differential equations. It concerns SGD; the dependence on the batch size is the subject of \cref{ch:noise}.

\begin{observation}[Normalized steps carry memory; \textcite{BernsteinEtAl2018}, \textcite{KarimireddyEtAl2019}, \textcite{KunstnerEtAl2023}]\label[observation]{obs:momentum:normalized}
Every widely used normalized optimizer takes its step along an average of past gradients. Without the average and without large batches, sign descent can stall or fail to converge; with momentum and large batches, it comes close to Adam on transformers.
\end{observation}

The optimizers in question are Signum, the sign of an exponential moving average \parencite{BernsteinEtAl2018}, Lion and Muon. The evidence that they need the average is of two kinds. Experimentally, \textcite{KunstnerEtAl2023} found that the advantage of sign descent with momentum over SGD grows with the batch size. Theoretically, sign descent with stochastic gradients can fail to converge \parencite{KarimireddyEtAl2019}, and normalized SGD was first analysed with minibatches whose size grows as the target accuracy shrinks \parencite{HazanLevyShalevShwartz2015}, then with momentum and without large minibatches \parencite{CutkoskyMehta2020}. How much of the need is due to small batches rather than to other features of the training runs is less well documented than the practice itself.

\begin{observation}[Momentum moves the edge of stability; \textcite{CohenEtAl2022}]\label[observation]{obs:momentum:edge}
In full-batch training with Adam and $\beta_1=0.9$, the largest eigenvalue of the preconditioned Hessian rises to about $38/\eta$ and then hovers there, as the sharpness of gradient descent hovers at $2/\eta$ \parencite{CohenEtAl2021}.
\end{observation}

The number $38$ is $2(1+\beta_1)/(1-\beta_1)$ at $\beta_1=0.9$, a stability threshold that we derive in \cref{prop:momentum:regions}; why training climbs to the threshold is the subject of \cref{ch:eos}. So momentum is used everywhere, helps SGD only in some regimes, is indispensable for normalized methods, and changes where training sits. What does an average of past gradients do that a different step size cannot do, and why do normalized steps need it when SGD can do without?

\section{The idealization: one filter, three settings}\label{sec:momentum:model}

A moving average is a linear filter. Unrolled, $m_t=(1-\beta)\sum_{s\le t}\beta^{t-s}g_s$ is a weighted sum of past gradients with a memory of about $1/(1-\beta)$ steps; half of its weight lies on the last $\ln2/\ln(1/\beta)\approx6.6$ steps when $\beta=0.9$. The methods of practice differ in where the filter sits and how it is normalized.

\begin{definition}[Momentum methods]\label[definition]{def:momentum:forms}
Let $g_t$ be the gradient or stochastic gradient at $\theta_t$, $\eta>0$ and $0\le\beta<1$.
\begin{enumerate}
\item \emph{Heavy ball}: $\theta_{t+1}=\theta_t-\eta g_t+\beta(\theta_t-\theta_{t-1})$, with $\theta_{-1}=\theta_0$. Its \emph{PyTorch form} is $b_t=\beta b_{t-1}+g_t$, $\theta_{t+1}=\theta_t-\eta b_t$, with $b_{-1}=0$.
\item \emph{EMA momentum}: $m_t=\beta m_{t-1}+(1-\beta)g_t$, $\theta_{t+1}=\theta_t-\eta m_t$, with $m_{-1}=0$; this is Adam's first moment.
\item \emph{Quasi-hyperbolic momentum} \parencite{MaYarats2019}: $m_t$ as in (ii), $u_t=(1-\nu)g_t+\nu m_t$ with $0\le\nu\le1$, and $\theta_{t+1}=\theta_t-\eta u_t$. It is EMA momentum for $\nu=1$ and gradient descent for $\nu=0$; for $\nu=\beta$ we call it \emph{Nesterov momentum in EMA form}, the form used by Muon.
\item \emph{Nesterov momentum}: $\theta_{t+1}=y_t-\eta\nabla\Loss(y_t)$ with $y_t=\theta_t+\beta(\theta_t-\theta_{t-1})$. Its PyTorch form is $b_t=\beta b_{t-1}+g_t$, $\theta_{t+1}=\theta_t-\eta(g_t+\beta b_t)$.
\item \emph{Normalized steepest descent with momentum}: $m_t$ as in (ii) and $\theta_{t+1}=\theta_t+\eta u_t$ with $u_t\in\lmo(m_t)$.
\end{enumerate}
\end{definition}

The forms are related by changes of the step size, which is why the two conventions of practice can be compared at all.

\begin{lemma}[One filter, several step sizes]\label[lemma]{lem:momentum:forms}
(i) The PyTorch form of heavy ball with step size $\eta$ produces the iterates of heavy ball with step size $\eta$. (ii) EMA momentum with step size $\eta$ produces the iterates of heavy ball with step size $\eta(1-\beta)$. (iii) The PyTorch form of Nesterov momentum with step size $\eta$ produces the iterates of Nesterov momentum in EMA form with step size $\eta/(1-\beta)$. (iv) On a quadratic, along an eigenvector of the Hessian with eigenvalue $h$, both forms of Nesterov momentum satisfy $x_{t+1}=(1-\eta h)\bigl((1+\beta)x_t-\beta x_{t-1}\bigr)$.
\end{lemma}

\begin{proof}
For (i), $\theta_{t+1}-\theta_t=-\eta b_t=-\eta\beta b_{t-1}-\eta g_t=\beta(\theta_t-\theta_{t-1})-\eta g_t$ for $t\ge1$, and $\theta_1=\theta_0-\eta g_0$. For (ii), the same computation with $\eta m_t$ in place of $\eta b_t$ gives $\theta_{t+1}-\theta_t=\beta(\theta_t-\theta_{t-1})-\eta(1-\beta)g_t$. For (iii), both averages start at zero, so $b_t=m_t/(1-\beta)$ and $g_t+\beta b_t=\bigl((1-\beta)g_t+\beta m_t\bigr)/(1-\beta)$. For (iv), the first form gives $x_{t+1}=(1-\eta h)y_t$ with $y_t=(1+\beta)x_t-\beta x_{t-1}$. In the second, $x_t-x_{t+1}=\eta(hx_t+\beta b_t)$ and $b_t=\beta b_{t-1}+hx_t$; eliminating $\beta b_t$ and $\beta^2b_{t-1}$ with the first identity at times $t$ and $t-1$ gives the same recursion.
\end{proof}

For a normalized step the conventions do not matter at all: $\lmo(cm)=\lmo(m)$ for $c>0$, so the PyTorch average, the EMA and Adam's bias-corrected EMA give the same iterates. We study the filter in three settings.

\begin{assumption}[Quadratic loss, additive noise]\label[assumption]{ass:momentum:quadratic}
$\Loss(\theta)=\frac12(\theta-\theta^\star)^\top H(\theta-\theta^\star)$ on $\R^d$, where $H$ is symmetric with eigenvalues $h_1\ge\dots\ge h_d>0$ and orthonormal eigenvectors $v_1,\dots,v_d$; we write $L=h_1$, $\mu=h_d$ and $\kappa=L/\mu$. The stochastic gradient is $g_t=H(\theta_t-\theta^\star)+\zeta_t$, where $\zeta_0,\zeta_1,\dots$ are independent and identically distributed, independent of the initialization, with mean zero and covariance $C$. Exact gradients are the case $C=0$.
\end{assumption}

\begin{assumption}[Smooth loss, unbiased noise]\label[assumption]{ass:momentum:smooth}
$\Loss\colon\R^d\to\R$ is bounded below by $\Loss^\star$ and $L$-smooth with respect to a norm $\norm{\cdot}$. The stochastic gradient is $g_t=\nabla\Loss(\theta_t)+\zeta_t$, where $\E[\zeta_t\mid\cF_t]=0$ and $\E[\norm{\zeta_t}_2^2\mid\cF_t]\le\sigma^2$ for the $\sigma$-algebra $\cF_t$ generated by $\theta_0,\zeta_0,\dots,\zeta_{t-1}$. The constant $\chi=\max_{z\ne0}\norm{z}_*/\norm{z}_2$ compares the dual norm with the Euclidean norm.
\end{assumption}

Setting (a) is \cref{ass:momentum:quadratic} with $C=0$: acceleration and stability. Setting (b) is the same with $C\ne0$, the \term{noisy quadratic model}\index{noisy quadratic model} of \textcite{ZhangEtAl2019} that \cref{ch:noise} develops: averaging. Setting (c) is \cref{ass:momentum:smooth} with normalized steps. Its constant $\chi$ is $1$ for the Euclidean norm, $\sqrt d$ for $\ell_\infty$, whose dual is $\ell_1$, and $\sqrt{\min(m,n)}$ for the spectral norm on $m\times n$ matrices, whose dual is the nuclear norm.

\begin{dictionary}
the momentum coefficient $\beta$ & $0.9$ for SGD, Adam's $\beta_1$ and Lion's step; $0.95$ for Muon; $0.99$ for Lion's state\\
the memory $1/(1-\beta)$ & the number of recent minibatches that the optimizer state remembers\\
the effective step $\eta/(1-\beta)$ of heavy ball & the learning rate that matters for SGD with momentum in PyTorch's convention\\
a coordinate with curvature $h$ & an eigenvector of the Hessian; the top one has $h=S$, the sharpness\\
the threshold $\eta h<2(1+\beta)/(1-\beta)$ & the edge of stability of Adam-style momentum, $38/\eta$ at $\beta_1=0.9$\\
additive noise of covariance $C$ & minibatch noise $\Sigma/B$\\
the momentum error $m_t-\nabla\Loss(\theta_t)$ & the gap between the optimizer state and the full-batch gradient at the current weights\\
its lag, of order $L\eta/(1-\beta)$ & stale gradients, computed at weights up to $1/(1-\beta)$ steps old\\
skewed noise & rare large per-example gradients
\end{dictionary}

\begin{fieldguide}{SGD with momentum}{\textcite{Polyak1964}}
\fgrow{update}{heavy ball $\theta_{t+1}=\theta_t-\eta g_t+\beta(\theta_t-\theta_{t-1})$, or in PyTorch's form $b_t=\beta b_{t-1}+g_t$, $\theta_{t+1}=\theta_t-\eta b_t$; EMA form $m_t=\beta m_{t-1}+(1-\beta)g_t$, $\theta_{t+1}=\theta_t-\eta m_t$}
\fgrow{norm}{$\ell_2$, unnormalized}
\fgrow{preconditioner}{none}
\fgrow{averages}{exponential moving average of $g$; PyTorch's form multiplies the step by $1/(1-\beta)$}
\fgrow{state}{one number per parameter}
\fgrow{in the wild}{$\beta=0.9$, $\eta=0.1$, weight decay $10^{-4}$ for ResNets \parencite{HeEtAl2016}}
\fgrow{caution}{in PyTorch's form a change of $\beta$ changes the effective step $\eta/(1-\beta)$}
\end{fieldguide}

\begin{fieldguide}{Nesterov momentum}{\textcite{Nesterov1983}; \textcite{SutskeverEtAl2013}}
\fgrow{update}{$b_t=\beta b_{t-1}+g_t$, $\theta_{t+1}=\theta_t-\eta(g_t+\beta b_t)$; in EMA form $\theta_{t+1}=\theta_t-\eta\bigl(\beta m_t+(1-\beta)g_t\bigr)$}
\fgrow{norm}{$\ell_2$, unnormalized; spectral and normalized inside Muon (\cref{ch:spectral})}
\fgrow{preconditioner}{none}
\fgrow{averages}{the EMA of $g$ with the fresh gradient counted twice: $\beta m_t+(1-\beta)g_t=\beta^2m_{t-1}+(1-\beta^2)g_t$}
\fgrow{state}{one number per parameter}
\fgrow{in the wild}{an option of PyTorch's SGD; the default in Muon, with $\beta=0.95$ \parencite{JordanEtAl2024Muon}}
\end{fieldguide}

\begin{fieldguide}{Lion}{\textcite{ChenEtAl2023}}
\fgrow{update}{in Lion's notation, $c_t=\beta_1m_{t-1}+(1-\beta_1)g_t$, $\theta_{t+1}=\theta_t-\eta\bigl(\sign(c_t)+\lambda\theta_t\bigr)$, $m_t=\beta_2m_{t-1}+(1-\beta_2)g_t$}
\fgrow{norm}{$\ell_\infty$, normalized: sign descent (\cref{ch:sign})}
\fgrow{preconditioner}{none}
\fgrow{averages}{one EMA of $g$ with two coefficients: $\beta_2$ for the state and $\beta_1$ for the step; it is the sign of quasi-hyperbolic momentum with $\beta=\beta_2$ and $\nu=\beta_1/\beta_2$, and Signum when $\beta_1=\beta_2$}
\fgrow{state}{one number per parameter, against Adam's two}
\fgrow{in the wild}{$\beta_1=0.9$, $\beta_2=0.99$; a learning rate several times smaller and a weight decay several times larger than AdamW's}
\fgrow{variants}{with weight decay, a method constrained to an $\ell_\infty$ ball (\cref{ch:decay})}
\end{fieldguide}

\begin{fieldguide}{Schedule-free SGD and AdamW}{\textcite{DefazioEtAl2024}}
\fgrow{update}{$y_t=(1-\beta)z_t+\beta x_t$, $z_{t+1}=z_t-\gamma g(y_t)$, $x_{t+1}=(1-c_{t+1})x_t+c_{t+1}z_{t+1}$ with $c_{t+1}=1/(t+1)$: gradients at $y_t$, evaluation at $x_t$}
\fgrow{norm}{$\ell_2$, unnormalized; Adam's diagonal metric in the AdamW version}
\fgrow{preconditioner}{none, or Adam's $\diag(\sqrt{\hat v_t})+\epsilon$}
\fgrow{averages}{a uniform average $x_t$ of the iterates $z_t$, and an interpolation $y_t$ that acts as momentum}
\fgrow{state}{one vector $z$ besides the weights, which hold $y$ in training and $x$ in evaluation}
\fgrow{in the wild}{$\beta=0.9$ and no decaying schedule; Schedule-Free AdamW won the self-tuning track of the 2024 AlgoPerf competition \parencite{KasimbegEtAl2025}}
\end{fieldguide}

\paragraph{What was thrown away.} In settings (a) and (b), everything that makes a network more than a quadratic: the Hessian is constant, so the step size cannot change the curvature (\cref{ch:eos}), and the noise is additive, independent of the position, and has a finite variance. In setting (c) the loss is general but the guarantee is a worst-case bound on gradients, not a description of the dynamics. Throughout, $\beta$ and $\eta$ are constant. The quadratic is the right first model because a moving average is a linear filter, and the filter's interaction with curvature and noise is exactly solvable when the dynamics is linear too; the worst-case setting is the right second model because the normalized step is not linear, and the question there is what an average can guarantee at all.

\paragraph{What one had to notice.} A linear filter commutes with a linear update. Averaging the gradients and then stepping is the same as stepping and then averaging the steps, and with small steps the iterates average the steps anyway: for SGD, momentum is a change of step size in disguise, at least when the noise dominates. A normalized step is not linear in the gradient, and there the order matters. A small step size averages normalized noisy gradients, $\frac1k\sum_s\lmo(g_s)$, which keeps whatever bias the normalization introduces; momentum normalizes an average, $\lmo\bigl(\frac1k\sum_sg_s\bigr)$, which removes it. The rest of the chapter makes these two sentences precise.

\section{The theory: momentum on a quadratic}\label{sec:momentum:quadratic}

Under \cref{ass:momentum:quadratic}, every method of \cref{def:momentum:forms} except the normalized one is linear, and the coordinates $x_t=\ip{\theta_t-\theta^\star,v_i}$ along the eigenvectors evolve separately. For quasi-hyperbolic momentum with exact gradients, the pair formed by $x_t$ and the coordinate $n_{t-1}$ of $m_{t-1}$ evolves by
\[
  \begin{pmatrix}x_{t+1}\\ n_t\end{pmatrix}=\begin{pmatrix}1-\eta h(1-\nu\beta)&-\eta\nu\beta\\(1-\beta)h&\beta\end{pmatrix}\begin{pmatrix}x_t\\ n_{t-1}\end{pmatrix},
\]
with $h=h_i$, as one checks from $u_t=\nu\beta m_{t-1}+(1-\nu\beta)g_t$. The trace of the matrix is $1+\beta-\eta h(1-\nu\beta)$ and its determinant is $\beta-\eta h\beta(1-\nu)$. A linear recursion tends to zero from every initial state if and only if both roots of its characteristic polynomial lie in the open unit disk, and the Jury test decides this for a quadratic polynomial (\cref{lem:toolkit:recurrence}).

\begin{proposition}[Stability regions]\label[proposition]{prop:momentum:regions}
Let $h>0$, $0\le\beta<1$ and $0\le\nu\le1$. On a coordinate with curvature $h$, quasi-hyperbolic momentum with exact gradients is stable if and only if
\begin{equation}\label{eq:momentum:qhm}
  \eta h<\frac{2(1+\beta)}{1+\beta-2\nu\beta}.
\end{equation}
Consequently the stability thresholds for $\eta h$ are: $2(1+\beta)/(1-\beta)$ for EMA momentum, which is $38$ at $\beta=0.9$; $2(1+\beta)/\bigl((1-\beta)(1+2\beta)\bigr)$ for Nesterov momentum in EMA form, about $13.6$ at $\beta=0.9$; $2(1+\beta)$ for heavy ball in both of its forms, $3.8$ at $\beta=0.9$; and $2(1+\beta)/(1+2\beta)$ for Nesterov momentum in both of its forms, about $1.36$ at $\beta=0.9$.
\end{proposition}

\begin{proof}
The plan is to apply the Jury test to the characteristic polynomial $z^2+a_1z+a_0$ of the matrix above and then to translate the special cases with \cref{lem:momentum:forms}. Write $u=\eta h>0$, so that $a_1=-\bigl(1+\beta-u(1-\nu\beta)\bigr)$ and $a_0=\beta\bigl(1-u(1-\nu)\bigr)$. The test asks for $\abs{a_0}<1$ and $\abs{a_1}<1+a_0$. The inequality $-a_1<1+a_0$ reads $u(1-\nu\beta)>u\beta(1-\nu)$, that is, $u(1-\beta)>0$, which holds. The inequality $a_1<1+a_0$ reads
\[
  u(1-\nu\beta)+u\beta(1-\nu)<2(1+\beta),\qquad\text{that is,}\qquad u(1+\beta-2\nu\beta)<2(1+\beta).
\]
Next, $a_0\le\beta<1$. Finally, $a_0>-1$ means $u\beta(1-\nu)<1+\beta$, which follows from the previous inequality because $1+\beta-2\nu\beta-2\beta(1-\nu)=1-\beta\ge0$. So the recursion is stable exactly when $u(1+\beta-2\nu\beta)<2(1+\beta)$, which is \eqref{eq:momentum:qhm} because $1+\beta-2\nu\beta\ge1-\beta>0$.

The case $\nu=1$ is EMA momentum, and $\nu=\beta$ gives $1+\beta-2\beta^2=(1-\beta)(1+2\beta)$. By \cref{lem:momentum:forms}(i) and (ii), heavy ball with step $\eta$ in either form is EMA momentum with step $\eta/(1-\beta)$, which is stable when $\eta h/(1-\beta)<2(1+\beta)/(1-\beta)$. By \cref{lem:momentum:forms}(iii) and (iv), both forms of Nesterov momentum with step $\eta$ follow the recursion of Nesterov momentum in EMA form with step $\eta/(1-\beta)$, which is stable when $\eta h/(1-\beta)<2(1+\beta)/\bigl((1-\beta)(1+2\beta)\bigr)$.
\end{proof}

Adam's convention widens the stable range of $\eta h$ by the factor $(1+\beta)/(1-\beta)$, which is $19$ at $\beta=0.9$; PyTorch's convention widens it by $1+\beta<2$, because it also multiplies the effective step by $1/(1-\beta)$. Nesterov's look-ahead shrinks the range below gradient descent's $2$: evaluating the gradient at the extrapolated point amplifies the coordinate that flips sign at every step. These are the thresholds that \cref{ch:eos} meets at the edge of stability, and the first prediction of the chapter: with Adam's $\beta_1=0.9$ the edge sits at $38/\eta$ (\cref{obs:momentum:edge}). For the stiff coordinates, which decide stability, momentum is thus a way to take larger steps; for the flat coordinates it does something else.

\begin{theorem}[Heavy ball on quadratics; \textcite{Polyak1964}]\label{thm:momentum:heavyball}
Under \cref{ass:momentum:quadratic} with exact gradients and $\mu<L$, let $r=(\sqrt\kappa-1)/(\sqrt\kappa+1)$ and run heavy ball with
\[
  \eta=\frac{4}{(\sqrt L+\sqrt\mu)^2},\qquad\beta=r^2,\qquad\theta_{-1}=\theta_0.
\]
\begin{enumerate}
\item For every $h\in[\mu,L]$, both roots of $z^2-(1+\beta-\eta h)z+\beta$ have modulus exactly $r$.
\item For every $t\ge0$, $\norm{\theta_t-\theta^\star}_2\le\bigl(1+(1+r)t\bigr)\,r^t\,\norm{\theta_0-\theta^\star}_2$.
\item For every $\eta'>0$ and every $\beta'\in\R$, heavy ball with step $\eta'$ and coefficient $\beta'$ has a characteristic root of modulus at least $r$ at $h=\mu$ or at $h=L$.
\end{enumerate}
\end{theorem}

Gradient descent contracts every coordinate by $\abs{1-\eta h}$, which is at least $(\kappa-1)/(\kappa+1)$ at the best step $2/(L+\mu)$; to reduce the error by a factor $e$ it needs about $\kappa/2$ steps, heavy ball about $\sqrt\kappa/2$.

\begin{proof}
The plan is to show that the chosen parameters put every eigenvalue into the band where the characteristic roots are complex, to solve the recursion there in closed form, and to prove optimality with a rescaled Jury test.

(i) The roots of $z^2-az+\beta$ with $a=1+\beta-\eta h$ are complex conjugate or double when $a^2\le4\beta$, and then their product is $\beta$, so both have modulus $\sqrt\beta=r$. The condition $\abs{1+\beta-\eta h}\le2\sqrt\beta$ is $(1-r)^2\le\eta h\le(1+r)^2$. Here $1-r=2/(\sqrt\kappa+1)$ and $1+r=2\sqrt\kappa/(\sqrt\kappa+1)$, so $(1-r)^2=4\mu/(\sqrt L+\sqrt\mu)^2=\eta\mu$ and $(1+r)^2=4L/(\sqrt L+\sqrt\mu)^2=\eta L$. Hence $\eta h\in[(1-r)^2,(1+r)^2]$ for every $h\in[\mu,L]$.

(ii) Fix an eigenvalue $h$ and its coordinate $x_t$, which satisfies $x_{t+1}=ax_t-r^2x_{t-1}$ with $x_1=(1-\eta h)x_0$, because $\theta_{-1}=\theta_0$. By (i), $a=2r\cos\varphi$ for some $\varphi\in[0,\pi]$; note $r>0$ because $\mu<L$. The rescaled sequence $y_t=x_t/r^t$ satisfies $y_{t+1}=2\cos\varphi\,y_t-y_{t-1}$, with $y_0=x_0$ and $y_1=(1-\eta h)x_0/r=(2\cos\varphi-r)x_0$, since $1-\eta h=a-\beta$. If $0<\varphi<\pi$, the solutions of this recursion are $A\cos t\varphi+B\sin t\varphi$, and matching $y_0$ and $y_1$ gives
\[
  x_t=r^tx_0\Bigl(\cos t\varphi+(\cos\varphi-r)\frac{\sin t\varphi}{\sin\varphi}\Bigr).
\]
The inequality $\abs{\sin t\varphi}\le t\abs{\sin\varphi}$ holds by induction on $t$, since $\abs{\sin(t+1)\varphi}\le\abs{\sin t\varphi}+\abs{\sin\varphi}$; with $\abs{\cos\varphi-r}\le1+r$ it gives $\abs{x_t}\le(1+(1+r)t)r^t\abs{x_0}$. If $\varphi\in\{0,\pi\}$, which happens only at $h=\mu$ or $h=L$, the same bound follows by continuity, since $x_t$ is a polynomial in $h$. Summing the squares of the coordinates gives (ii).

(iii) Suppose that both roots have modulus at most $\rho<1$ at $h=\mu$ and at $h=L$; we show $\rho\ge r$. If $\rho=0$ there is nothing to prove unless $r>0$, and then the roots at $h=\mu$ and $h=L$ would both be $0$, forcing $\beta'=0$ and $\eta'\mu=1=\eta'L$, which is impossible. So let $\rho>0$. Write $a_h=1+\beta'-\eta'h$. The product of the roots is $\beta'$, so $\abs{\beta'}\le\rho^2$. For every $s>\rho$, the polynomial $w^2-(a_h/s)w+\beta'/s^2$, whose roots are those of the original divided by $s$, has both roots in the open unit disk, and the Jury test gives $\abs{a_h}<s+\beta'/s$; letting $s\downarrow\rho$, $\abs{a_h}\le\rho+\beta'/\rho$. At $h=\mu$ this gives $\eta'\mu\ge1+\beta'-\rho-\beta'/\rho=(1-\rho)(1-\beta'/\rho)$, and at $h=L$ it gives $\eta'L\le(1+\rho)(1+\beta'/\rho)$. Since $\abs{\beta'/\rho}\le\rho<1$, the lower bound is positive, and dividing,
\[
  \kappa=\frac{\eta'L}{\eta'\mu}\le\frac{1+\rho}{1-\rho}\cdot\frac{1+\beta'/\rho}{1-\beta'/\rho}\le\Bigl(\frac{1+\rho}{1-\rho}\Bigr)^2,
\]
because $x\mapsto(1+x)/(1-x)$ is increasing on $(-1,1)$ and $\beta'/\rho\le\rho$. Hence $\sqrt\kappa\le(1+\rho)/(1-\rho)$, which is $\rho\ge r$.
\end{proof}

The bound in (ii) is attained at $h=L$, where the double root $-r$ makes the coordinate $(-r)^t(1+(1+r)t)x_0$; the linear factor is the price of a Jordan block, and of tuning the parameters so that the band of complex roots fits $[\mu,L]$ exactly. The theorem is a theorem about quadratics. For general strongly convex functions with the same $\mu$ and $L$, heavy ball with these parameters need not converge \parencite{LessardRechtPackard2016} (\cref{ex:momentum:lessard}); Nesterov's method keeps the accelerated rates for every smooth convex function, with $\Loss(\theta_t)-\Loss^\star=O(L\norm{\theta_0-\theta^\star}_2^2/t^2)$ in the convex case and a contraction of order $1-1/\sqrt\kappa$ per step in the strongly convex case \parencite{Nesterov1983,Nesterov2018}, and in high dimension no method that queries only gradients can improve these rates by more than constant factors.\sidenote{Polyak's 1964 paper proves the local rate of heavy ball for twice differentiable functions, which is the quadratic case near a minimizer. \textcite{Polyak1987} gives a textbook account, and \textcite[Chapter 2]{Nesterov2018} the lower bounds and the accelerated method.}

Acceleration is a statement about exact gradients. With noise, the same filter has a second effect, and it is the one that \cref{obs:momentum:rescaled} reports.

\begin{theorem}[In the noise-dominated regime momentum is a larger step]\label{thm:momentum:noise}
Let $h>0$, $\sigma>0$, $0\le\beta<1$ and $0<\eta h<2(1+\beta)$. Let $\zeta_0,\zeta_1,\dots$ be independent and identically distributed with mean zero and variance $\sigma^2$, independent of $(x_0,x_{-1})$, which have finite second moments, and let
\[
  x_{t+1}=x_t-\eta(hx_t+\zeta_t)+\beta(x_t-x_{t-1}).
\]
\begin{enumerate}
\item As $t\to\infty$, $\E x_t\to0$ and $\E x_t^2\to V$, where
\begin{equation}\label{eq:momentum:variance}
  V=\frac{\eta\sigma^2(1+\beta)}{h(1-\beta)\bigl(2(1+\beta)-\eta h\bigr)}=\frac{\eta\sigma^2}{2h(1-\beta)}\Bigl(1-\frac{\eta h}{2(1+\beta)}\Bigr)^{-1}.
\end{equation}
\item If $\eta h<2(1-\beta)$, gradient descent with step $\eta/(1-\beta)$ and the same noise has stationary variance $\frac{\eta\sigma^2}{2h(1-\beta)}\bigl(1-\frac{\eta h}{2(1-\beta)}\bigr)^{-1}$, which is at least $V$. The two agree to first order in $\eta h/(1-\beta)$.
\item If $\eta h\le(1-\beta)^2/4$, the characteristic roots of the recursion are real and distinct, $\E x_t$ is a combination of their powers, and the larger root $z_+$ satisfies
\[
  1-\frac{\eta h}{1-\beta}\Bigl(1-\frac{2\eta h}{(1-\beta)^2}\Bigr)^{-1}\le z_+\le1-\frac{\eta h}{1-\beta},
\]
while gradient descent with step $\eta/(1-\beta)$ contracts its mean by exactly $1-\eta h/(1-\beta)$.
\item EMA momentum $m_t=\beta m_{t-1}+(1-\beta)(hx_t+\zeta_t)$, $x_{t+1}=x_t-\eta m_t$, with $\eta h<2(1+\beta)/(1-\beta)$, has stationary variance $\frac{\eta\sigma^2}{2h}\bigl(1-\frac{\eta(1-\beta)h}{2(1+\beta)}\bigr)^{-1}$, which agrees to first order in $\eta h$ with that of gradient descent with step $\eta$.
\end{enumerate}
\end{theorem}

\begin{proof}
The plan is to solve the Lyapunov equation of the two-dimensional recursion exactly and then compare. The state $X_t=(x_t,x_{t-1})^\top$ satisfies $X_{t+1}=AX_t+w_t$ with
\[
  A=\begin{pmatrix}a&-\beta\\1&0\end{pmatrix},\qquad a=1+\beta-\eta h,\qquad w_t=\begin{pmatrix}-\eta\zeta_t\\0\end{pmatrix},
\]
and $w_t$ is independent of $X_t$. By \cref{prop:momentum:regions} the spectral radius of $A$ is less than $1$, so $\E X_t=A^t\E X_0\to0$, and by \cref{lem:toolkit:lyapunov} the second-moment matrix $\E X_tX_t^\top$ converges to the unique solution $\Sigma$ of $\Sigma=A\Sigma A^\top+Q$, with $Q=\diag(\eta^2\sigma^2,0)$. Write $\Sigma=\bigl(\begin{smallmatrix}V&c\\c&V'\end{smallmatrix}\bigr)$. The entry $(2,2)$ of the equation gives $V'=V$, the entry $(1,2)$ gives $c=aV-\beta c$, so $c=aV/(1+\beta)$, and the entry $(1,1)$ gives $V=a^2V-2a\beta c+\beta^2V+\eta^2\sigma^2$. Substituting $c$,
\[
  V\Bigl(1-\beta^2-\frac{1-\beta}{1+\beta}\,a^2\Bigr)=\eta^2\sigma^2,\qquad\text{and}\qquad1-\beta^2-\frac{1-\beta}{1+\beta}\,a^2=\frac{1-\beta}{1+\beta}(1+\beta-a)(1+\beta+a).
\]
With $1+\beta-a=\eta h$ and $1+\beta+a=2(1+\beta)-\eta h$ this is \eqref{eq:momentum:variance}, and the second form follows by factoring out $2(1+\beta)$.

For (ii), gradient descent is the case $\beta=0$ of (i), which needs $\eta h/(1-\beta)<2$; with step $\eta/(1-\beta)$ the second form of \eqref{eq:momentum:variance} becomes the stated expression. It is at least $V$ because $1-\beta\le1+\beta$, and both expressions equal $\frac{\eta\sigma^2}{2h(1-\beta)}\bigl(1+O(\eta h/(1-\beta))\bigr)$.

For (iii), write $b=1-\beta$, $u=\eta h$ and $z=1-\delta$. The characteristic polynomial becomes $\delta^2-(b+u)\delta+u$, whose discriminant $(b+u)^2-4u\ge b^2-4u$ is positive when $u\le b^2/4$. Both $\delta$-roots are positive, because their sum $b+u$ and product $u$ are; let $\delta$ be the smaller one, so that $z_+=1-\delta$, and note $\delta\le(b+u)/2<1$. From $\delta(b-\delta)=u(1-\delta)>0$ we get $\delta<b$, and then $\delta=u(1-\delta)/(b-\delta)\ge u/b$, because $b(1-\delta)\ge b-\delta$. For the other side, $\delta=2u/\bigl(b+u+\sqrt{(b+u)^2-4u}\bigr)\le2u/b$, so $\delta\le u/(b-\delta)\le u/(b-2u/b)$, which is the stated lower bound on $z_+$. The mean $\E x_t$ follows the noiseless recursion, so it is a combination of $z_+^t$ and $z_-^t$.

For (iv), EMA momentum with step $\eta$ is heavy ball with step $\eta(1-\beta)$ and the same noise (\cref{lem:momentum:forms}(ii)); substitute into (i) and compare with (i) at $\beta=0$.
\end{proof}

For the noisy quadratic model of \cref{ass:momentum:quadratic}, each coordinate is such a recursion, driven by the noise coordinate $\ip{\zeta_t,v_i}$ of variance $c_i=v_i^\top Cv_i$; correlations between noise coordinates do not change the variance of each coordinate. Hence the stationary excess loss of heavy ball is $\sum_i\frac{\eta c_i}{4(1-\beta)}\bigl(1-\frac{\eta h_i}{2(1+\beta)}\bigr)^{-1}$, the formula that \cref{ch:noise} uses.

The theorem says two things. When $\eta h\ll(1-\beta)^2$, which is where the noise dominates and the step is small, heavy ball is gradient descent with step $\eta/(1-\beta)$: the mean decays at the same rate and the noise floor is the same, up to relative corrections of order $\eta h/(1-\beta)^2$ and $\eta h/(1-\beta)$. In Adam's convention momentum then changes nothing at first order. And where the step is limited by stability rather than by noise, the two differ: heavy ball stays stable up to an effective step $2(1+\beta)/((1-\beta)h)$, while gradient descent diverges at $2/h$, and its floor at a given effective step is lower. Both effects are in \cref{fig:momentum:linear}.

\section{The theory: normalized steps with memory}\label{sec:momentum:normalized}

We now drop the quadratic and the linearity. The algorithm is \cref{def:momentum:forms}(v): $m_t=\beta m_{t-1}+(1-\beta)g_t$ from a given $m_{-1}$, and $\theta_{t+1}=\theta_t+\eta u_t$ with $u_t\in\lmo(m_t)$, a step of length at most $\eta$ in the norm $\norm{\cdot}$ along a steepest direction of the average (\cref{ch:steepest}). The selection $u_t$ is any measurable function of $m_t$; when $m_t=0$, every point of the unit ball is allowed. Two lemmas carry the argument. The first says that a steepest direction of $m$ is a good direction for $g$ when $m$ is close to $g$.

\begin{lemma}[A steepest direction of a nearby vector]\label[lemma]{lem:momentum:direction}
For all $m,g\in\R^d$ and every $u$ with $\norm u\le1$ and $\ip{m,u}=-\norm m_*$,
\[
  \ip{g,u}\le-\norm m_*+\norm{m-g}_*\le-\norm g_*+2\norm{m-g}_*.
\]
\end{lemma}

\begin{proof}
Split $\ip{g,u}=\ip{m,u}+\ip{g-m,u}$ and bound the second term by $\norm{g-m}_*\norm u\le\norm{g-m}_*$ (Hölder's inequality, \cref{lem:toolkit:dualnorm}); this gives the first inequality. The second is the triangle inequality $\norm m_*\ge\norm g_*-\norm{m-g}_*$. If $m=0$, every $u$ in the unit ball qualifies and both sides equal $\ip{g,u}\le\norm g_*$.
\end{proof}

The second lemma bounds the \term{momentum error}\index{momentum error} $\delta_t=m_t-\nabla\Loss(\theta_t)$, the gap between the average and the gradient at the current point.

\begin{lemma}[The error of the average]\label[lemma]{lem:momentum:error}
Under \cref{ass:momentum:smooth}, let $m_t=\beta m_{t-1}+(1-\beta)g_t$ for $t\ge0$, where the iterates satisfy $\norm{\theta_{t+1}-\theta_t}\le\eta$, and let $\delta_{-1}=m_{-1}-\nabla\Loss(\theta_0)$. Then for every $t\ge0$,
\[
  \E\norm{\delta_t}_*\le\beta^{t+1}\norm{\delta_{-1}}_*+\frac{\beta L\eta}{1-\beta}+\chi\sigma\sqrt{\frac{1-\beta}{1+\beta}}.
\]
\end{lemma}

\begin{proof}
The plan is to unroll the recursion of $\delta_t$ into a decaying initial error, a lag and a noise average, and to bound each. With the convention $\theta_{-1}=\theta_0$,
\[
  \delta_t=\beta\delta_{t-1}+\beta\bigl(\nabla\Loss(\theta_{t-1})-\nabla\Loss(\theta_t)\bigr)+(1-\beta)\zeta_t\qquad(t\ge0),
\]
which one checks by inserting $g_t=\nabla\Loss(\theta_t)+\zeta_t$ into the definition of $m_t$. Unrolled,
\[
  \delta_t=\beta^{t+1}\delta_{-1}+\sum_{s=1}^t\beta^{t-s+1}\bigl(\nabla\Loss(\theta_{s-1})-\nabla\Loss(\theta_s)\bigr)+(1-\beta)N_t,\qquad N_t=\sum_{s=0}^t\beta^{t-s}\zeta_s.
\]
Smoothness gives $\norm{\nabla\Loss(\theta_{s-1})-\nabla\Loss(\theta_s)}_*\le L\norm{\theta_{s-1}-\theta_s}\le L\eta$, and $\sum_{s=1}^t\beta^{t-s+1}\le\beta/(1-\beta)$, so the middle sum has dual norm at most $\beta L\eta/(1-\beta)$. For the noise, $\E\norm{N_t}_*\le\chi\E\norm{N_t}_2\le\chi(\E\norm{N_t}_2^2)^{1/2}$. In $\E\norm{N_t}_2^2$ the cross terms vanish: for $r<s$, $\zeta_r$ is $\cF_s$-measurable, so $\E\ip{\zeta_r,\zeta_s}=\E\ip{\zeta_r,\E[\zeta_s\mid\cF_s]}=0$. Hence $\E\norm{N_t}_2^2=\sum_s\beta^{2(t-s)}\E\norm{\zeta_s}_2^2\le\sigma^2/(1-\beta^2)$, and $(1-\beta)\chi\sigma/\sqrt{1-\beta^2}=\chi\sigma\sqrt{(1-\beta)/(1+\beta)}$.
\end{proof}

The average trades one error for another. Its noise falls like $\sqrt{1-\beta}$, because it averages about $1/(1-\beta)$ independent gradients; its lag grows like $1/(1-\beta)$, because those gradients were computed at points up to $1/(1-\beta)$ steps away, each step of length $\eta$. The theorem converts this trade into a rate.

\begin{theorem}[Normalized steepest descent with momentum]\label{thm:momentum:normalized}
Under \cref{ass:momentum:smooth}, let $0\le\beta<1$, $\eta>0$ and $m_{-1}\in\R^d$, and run normalized steepest descent with momentum (\cref{def:momentum:forms}(v)). Then for every $T\ge1$, with $\Delta_0=\Loss(\theta_0)-\Loss^\star$,
\begin{equation}\label{eq:momentum:normalized}
  \frac1T\sum_{t=0}^{T-1}\E\norm{\nabla\Loss(\theta_t)}_*\le\frac{\Delta_0}{\eta T}+\frac{L\eta}2+\frac{2\beta L\eta}{1-\beta}+2\chi\sigma\sqrt{\frac{1-\beta}{1+\beta}}+\frac{2\beta\norm{m_{-1}-\nabla\Loss(\theta_0)}_*}{(1-\beta)T}.
\end{equation}
\end{theorem}

\begin{proof}
The plan is to bound the decrease of the loss at each step with \cref{lem:momentum:direction}, to add up, and to control the accumulated momentum error with \cref{lem:momentum:error}. By $L$-smoothness and the descent lemma (\cref{lem:toolkit:descent}), and because $\norm{u_t}\le1$,
\[
  \Loss(\theta_{t+1})\le\Loss(\theta_t)+\eta\ip{\nabla\Loss(\theta_t),u_t}+\frac{L\eta^2}2.
\]
\Cref{lem:momentum:direction} with $m=m_t$ and $g=\nabla\Loss(\theta_t)$ gives $\ip{\nabla\Loss(\theta_t),u_t}\le-\norm{\nabla\Loss(\theta_t)}_*+2\norm{\delta_t}_*$. Sum over $t<T$, use $\Loss(\theta_T)\ge\Loss^\star$, and take expectations:
\[
  \eta\sum_{t<T}\E\norm{\nabla\Loss(\theta_t)}_*\le\Delta_0+2\eta\sum_{t<T}\E\norm{\delta_t}_*+\frac{TL\eta^2}2.
\]
The steps have length $\norm{\theta_{t+1}-\theta_t}=\eta\norm{u_t}\le\eta$, so \cref{lem:momentum:error} applies, and $\sum_{t<T}\beta^{t+1}\le\beta/(1-\beta)$ gives
\[
  \sum_{t<T}\E\norm{\delta_t}_*\le\frac{\beta\norm{\delta_{-1}}_*}{1-\beta}+T\Bigl(\frac{\beta L\eta}{1-\beta}+\chi\sigma\sqrt{\frac{1-\beta}{1+\beta}}\Bigr).
\]
Insert this and divide by $\eta T$.
\end{proof}

The terms of \eqref{eq:momentum:normalized} have distinct origins: $\Delta_0/(\eta T)$ is the distance to travel at speed $\eta$; $L\eta/2$ is the curvature that a step of length $\eta$ meets; $2\beta L\eta/(1-\beta)$ is the lag; $2\chi\sigma\sqrt{(1-\beta)/(1+\beta)}$ is the noise left in the average; and the last term is the initial error, which decays in about $1/(1-\beta)$ steps. Two consequences follow by choosing parameters.

\begin{corollary}[Rates]\label[corollary]{cor:momentum:rate}
(i) With $1-\beta=T^{-1/2}$ and $\eta=(\Delta_0/L)^{1/2}T^{-3/4}$, the right side of \eqref{eq:momentum:normalized} is at most
\[
  \bigl(3\sqrt{L\Delta_0}+2\chi\sigma\bigr)T^{-1/4}+2\norm{m_{-1}-\nabla\Loss(\theta_0)}_*T^{-1/2}+\tfrac12\sqrt{L\Delta_0}\,T^{-3/4}.
\]
(ii) Without momentum, $\beta=0$, the right side is $\Delta_0/(\eta T)+L\eta/2+2\chi\sigma$, whose noise term does not decrease with $\eta$ or $T$.
\end{corollary}

\begin{proof}
For (i), $\Delta_0/(\eta T)=\sqrt{L\Delta_0}\,T^{-1/4}$, $L\eta/2=\frac12\sqrt{L\Delta_0}\,T^{-3/4}$, $2\beta L\eta/(1-\beta)\le2L\eta\,T^{1/2}=2\sqrt{L\Delta_0}\,T^{-1/4}$, $2\chi\sigma\sqrt{(1-\beta)/(1+\beta)}\le2\chi\sigma T^{-1/4}$, and $2\beta/((1-\beta)T)\le2T^{-1/2}$. Part (ii) is \eqref{eq:momentum:normalized} at $\beta=0$.
\end{proof}

So a point with $\E\norm{\nabla\Loss}_*\le\eps$ is found after $T=O(\eps^{-4})$ steps with constant memory and no large batches; in the Euclidean norm this is the order of the best possible guarantee for stochastic first-order methods under these assumptions \parencite{ArjevaniEtAl2023}. Whether the floor $2\chi\sigma$ of (ii) is real or an artifact of the proof is the next question, and the answer is that it is real.

\begin{proposition}[Without memory, a floor of the size of the noise]\label[proposition]{prop:momentum:floor}
Let $\Loss(\theta)=\theta^2/2$ on $\R$, let $b>0$, and let $\zeta_0,\zeta_1,\dots$ be independent with $\Pr[\zeta_t=3b]=\frac14$ and $\Pr[\zeta_t=-b]=\frac34$, so that $\E\zeta_t=0$ and $\sigma^2=\E\zeta_t^2=3b^2$. Run normalized steepest descent without momentum, $\theta_{t+1}=\theta_t+\eta u_t$ with $u_t=-\sign(\theta_t+\zeta_t)$ when $\theta_t+\zeta_t\ne0$ and $u_t\in[-1,1]$ otherwise. Then for every $\eta>0$, every $\theta_0$ and every $t\ge0$,
\[
  \E\bigl[(b-\theta_t)_+\bigr]\le\eta\log_2\Bigl(\bigl(\tfrac78\bigr)^t2^{(b-\theta_0)_+/\eta}+32\Bigr).
\]
Consequently $\liminf_{t\to\infty}\E\abs{\nabla\Loss(\theta_t)}\ge b-5\eta=\sigma/\sqrt3-5\eta$.
\end{proposition}

\begin{proof}
The plan is to show that below $b$ the iterate performs a random walk with drift toward $b$, and to control its distance below $b$ with an exponential supermartingale. Let $V_t=(b-\theta_t)_+$ and $W_t=2^{V_t/\eta}$. We show $\E[W_{t+1}\mid\theta_t]\le\frac78W_t+4$ in four cases.
\emph{If $-3b<\theta_t\le b-\eta$:} with probability $\frac14$, $\theta_t+\zeta_t=\theta_t+3b>0$ and $\theta_{t+1}=\theta_t-\eta$; with probability $\frac34$, $\theta_t+\zeta_t=\theta_t-b<0$ and $\theta_{t+1}=\theta_t+\eta\le b$. So $W_{t+1}$ is $2W_t$ or $W_t/2$, and $\E[W_{t+1}\mid\theta_t]=(\frac12+\frac38)W_t=\frac78W_t$.
\emph{If $\theta_t\le-3b$:} with probability $\frac34$ the step is $+\eta$, which gives $V_{t+1}\le\max(V_t-\eta,0)$ and $W_{t+1}\le\max(W_t/2,1)$; otherwise $V_{t+1}\le V_t+\eta$ and $W_{t+1}\le2W_t$. So $\E[W_{t+1}\mid\theta_t]\le\frac78W_t+\frac34$.
\emph{If $b-\eta<\theta_t<b$:} $V_t<\eta$ and $V_{t+1}\le V_t+\eta<2\eta$, so $W_{t+1}<4$.
\emph{If $\theta_t\ge b$:} $\theta_t+\zeta_t\ge\theta_t-b\ge0$, so the step is $-\eta$ or, when the gradient is zero, at least $-\eta$; hence $V_{t+1}\le\eta$ and $W_{t+1}\le2$.
Taking expectations, $\E W_{t+1}\le\frac78\E W_t+4$, and by induction $\E W_t\le(\frac78)^tW_0+32$. Jensen's inequality for the convex function $v\mapsto2^{v/\eta}$ gives $2^{\E V_t/\eta}\le\E W_t$, which is the first claim. As $t\to\infty$ its right side tends to $\eta\log_232=5\eta$, and $\abs{\nabla\Loss(\theta_t)}=\abs{\theta_t}\ge\theta_t\ge b-V_t$.
\end{proof}

In one dimension every normalized step is a sign step, so the proposition covers all norms at once. The mechanism is a \term{bias of the normalized step}\index{bias of a normalized step}: when $\abs{\theta}<b$, the stochastic gradient has the wrong sign three times in four, because the noise is usually slightly negative and occasionally very positive, and the expected step $-\eta\,\E\sign(\theta+\zeta)=+\eta/2$ pushes uphill, until the iterate settles near $\theta=b$. No step size removes this: a small step averages the signs, and the average of the signs is what is wrong. Momentum removes it, because the sign is taken of an average: by \cref{thm:momentum:normalized} the floor falls to $\eta/2+2\beta\eta/(1-\beta)+2\sigma\sqrt{(1-\beta)/(1+\beta)}$. The bias needs asymmetric noise. If the law of $\zeta_t$ given $\cF_t$ is symmetric, then $\E[\ip{\nabla\Loss(\theta_t),u_t}\mid\cF_t]\le0$ in every norm (\cref{ex:momentum:symmetric}): a normalized step never points uphill on average, although its useful part can be small.

\subsection{Averaging the iterates}

Momentum averages gradients. Model averages and schedule-free methods average iterates, which is the other way to use a filter. The classical result is that the average of the iterates of stochastic approximation is asymptotically optimal \parencite{Ruppert1988,PolyakJuditsky1992}. A short argument shows that, for convex losses, averaging can replace a decaying schedule: evaluate the gradient at the running average itself.

\begin{proposition}[Averaging instead of a schedule; \textcite{Cutkosky2019}]\label[proposition]{prop:momentum:averaging}
Let $F$ be convex on $\R^d$, and let a stochastic oracle return at each query point $x$ a vector $\tilde g$ whose conditional mean, given the past, is a subgradient of $F$ at $x$ and whose conditional second moment $\E\norm{\tilde g}_2^2$ is at most $G^2$. Let $z_1=x_1$ and, for $t\ge1$, query the oracle at $x_t$ to get $\tilde g_t$, and set $z_{t+1}=z_t-\gamma\tilde g_t$ and $x_{t+1}=\frac{t}{t+1}x_t+\frac1{t+1}z_{t+1}$. Then for every $x^\star$ and $T\ge1$,
\[
  \E F(x_T)-F(x^\star)\le\frac{\norm{z_1-x^\star}_2^2}{2\gamma T}+\frac{\gamma G^2}2,
\]
which is $DG/\sqrt T$ for $\gamma=D/(G\sqrt T)$ and $D=\norm{z_1-x^\star}_2$.
\end{proposition}

\begin{proof}
The plan is to show that $T(F(x_T)-F(x^\star))$ is at most the linear regret of the sequence $z_t$, and to bound the regret as for online gradient descent. Let $g_t$ be the conditional mean of $\tilde g_t$, a subgradient at $x_t$. Since $tx_t=(t-1)x_{t-1}+z_t$, we have $z_t-x^\star=t(x_t-x^\star)-(t-1)(x_{t-1}-x^\star)$, and therefore
\[
  \ip{g_t,z_t-x^\star}=\ip{g_t,x_t-x^\star}+(t-1)\ip{g_t,x_t-x_{t-1}}\ge F(x_t)-F(x^\star)+(t-1)\bigl(F(x_t)-F(x_{t-1})\bigr),
\]
by convexity. The right side is $t(F(x_t)-F(x^\star))-(t-1)(F(x_{t-1})-F(x^\star))$, which telescopes: $\sum_{t\le T}\ip{g_t,z_t-x^\star}\ge T(F(x_T)-F(x^\star))$. On the other side, $\norm{z_{t+1}-x^\star}_2^2=\norm{z_t-x^\star}_2^2-2\gamma\ip{\tilde g_t,z_t-x^\star}+\gamma^2\norm{\tilde g_t}_2^2$, so $\sum_{t\le T}\ip{\tilde g_t,z_t-x^\star}\le\norm{z_1-x^\star}_2^2/(2\gamma)+\frac\gamma2\sum_{t\le T}\norm{\tilde g_t}_2^2$ (\cref{thm:toolkit:ogd}). Since $z_t$ is determined by the past, $\E\ip{\tilde g_t,z_t-x^\star}=\E\ip{g_t,z_t-x^\star}$, and taking expectations of both bounds gives the claim.
\end{proof}

This is schedule-free SGD with $\beta=1$, gradients at the average. With $\beta=0$, gradients at $z_t$, it is SGD with a constant step and Polyak--Ruppert averaging, and the same bound follows from Jensen's inequality. \textcite{DefazioEtAl2024} prove the bound $DG/\sqrt T$ for every interpolation $y_t=(1-\beta)z_t+\beta x_t$ with $\beta\in[0,1]$, and use $\beta=0.9$ in practice. The constant step $\gamma$ with a uniform average matches the guarantee of a carefully decayed step for the last iterate (\cref{ch:schedules}).

A moving average of the weights, $\bar\theta_t=\beta_e\bar\theta_{t-1}+(1-\beta_e)\theta_t$, does something different on a noisy quadratic, which \cref{ex:momentum:modelema} computes exactly. A coordinate whose relaxation time $1/(\eta h)$ is short compared with the window $1/(1-\beta_e)$ is averaged over many independent fluctuations, and its variance becomes that of gradient descent with the step $(1-\beta_e)/h$; a coordinate with a long relaxation time is left alone, and the average lags behind it by $\beta_e/(1-\beta_e)$ steps.

\begin{heuristic}[Heuristic derivation: a model average is a cooldown of the stiff coordinates]
In the language of \cref{ch:schedules}, lowering the step from $\eta$ to $\eta'$ divides the noise floor of each coordinate by about $\eta/\eta'$ once the coordinate has relaxed, which takes about $1/(\eta'h)$ steps. The model average achieves the same reduction, with $\eta'=(1-\beta_e)/h$, for every coordinate with $\eta h\gg1-\beta_e$, and achieves it without slowing the iterates themselves; it does nothing for the flat coordinates, which are the ones that a cooldown also leaves nearly unchanged. So a weight average with window $1/(1-\beta_e)$ acts like a cooldown of length $1/(1-\beta_e)$ restricted to the stiff coordinates, at the price of a lag of $1/(1-\beta_e)$ steps on the flat ones.
\end{heuristic}

\section{Reading the theory}\label{sec:momentum:reading}

\paragraph{Thresholds.} \Cref{prop:momentum:regions} gives exact limits on $\eta S$, where $S$ is the largest curvature: $2$ for gradient descent; $2(1+\beta)=3.8$ for heavy ball and PyTorch's SGD at $\beta=0.9$; $2(1+\beta)/(1-\beta)=38$ for Adam-style momentum; $13.6$ for Nesterov momentum in EMA form at $\beta=0.9$ and $26.9$ at Muon's $\beta=0.95$; and $1.36$ for Nesterov momentum in PyTorch's form. \Cref{fig:momentum:linear} (left) checks them against the measured growth of the iterates. The EMA threshold grows without bound as $\beta\to1$, which is why the edge of stability of Adam with $\beta_1=0.9$ lies nineteen times higher than that of gradient descent at the same $\eta$, while the PyTorch thresholds stay below $4$; the difference is the factor $1/(1-\beta)$ that the PyTorch convention puts into the effective step.

\paragraph{Acceleration.} Polyak's parameters for a condition number $\kappa$ are $\beta=\bigl((\sqrt\kappa-1)/(\sqrt\kappa+1)\bigr)^2$; inverting, $\beta=0.9$ is the optimal choice for $\kappa\approx1.4\cdot10^3$, and $\beta=0.99$ for $\kappa\approx1.6\cdot10^5$. On a quadratic with $\kappa=10^3$ and exact gradients, heavy ball reaches a loss $10^{-8}$ times the initial one in about $220$ steps against about $3900$ for gradient descent (\cref{ex:momentum:noisyaccel}). With noise the order reverses: Polyak's parameters put the stiffest coordinate close to the stability limit, where \eqref{eq:momentum:variance} is large, and the noise floor of heavy ball is then sixteen times that of gradient descent in the same example. Acceleration and noise reduction pull the parameters in opposite directions.

\paragraph{Noise.} \Cref{fig:momentum:linear} (right) plots the stationary loss of heavy ball divided by its first-order value $\eta\sigma^2/(4(1-\beta))$, against the effective step $\eta h/(1-\beta)$. All curves start at $1$: in the noise-dominated regime, momentum $\beta$ with step $\eta$ is plain SGD with step $\eta/(1-\beta)$ (\cref{obs:momentum:rescaled}). They separate only near the stability limit of SGD, $\eta h/(1-\beta)=2$, which momentum pushes out to $2(1+\beta)/(1-\beta)$. In the language of \cref{ch:noise}: at small batches the best step is small compared with the stability limit, and momentum offers nothing a larger learning rate does not; at large batches the best step is the stability limit, and momentum lets the effective step grow by the factor $(1+\beta)/(1-\beta)$, with a lower floor than SGD would have at that step.

\begin{figure}[t]
\begin{fullwidth}
\centering
\begin{tikzpicture}
\begin{groupplot}[group style={group size=2 by 1,horizontal sep=17mm},wplot,width=82mm,height=56mm]
\nextgroupplot[xlabel={momentum $\beta$},ylabel={largest stable $\eta h$},ymode=log,xmin=0,xmax=1,ymin=1,ymax=600,
  legend style={at={(0.02,0.98)},anchor=north west,font=\scriptsize},legend cell align=left]
\addplot[theory] table[x=beta,y=ema] {figures/data/05-momentum-thresholds.dat};
\addplot[theory2] table[x=beta,y=nema] {figures/data/05-momentum-thresholds.dat};
\addplot[theory3] table[x=beta,y=hb] {figures/data/05-momentum-thresholds.dat};
\addplot[muted,line width=0.9pt,dashdotted] table[x=beta,y=nest] {figures/data/05-momentum-thresholds.dat};
\addplot[simulation] table[x=beta,y=ema] {figures/data/05-momentum-thresholds-sim.dat};
\addplot[simulation,forget plot] table[x=beta,y=nema] {figures/data/05-momentum-thresholds-sim.dat};
\addplot[simulation,forget plot] table[x=beta,y=hb] {figures/data/05-momentum-thresholds-sim.dat};
\addplot[simulation,forget plot] table[x=beta,y=nest] {figures/data/05-momentum-thresholds-sim.dat};
\addplot[guide,domain=0:1] {2};
\legend{EMA (Adam),EMA Nesterov (Muon),heavy ball (PyTorch),Nesterov (PyTorch),measured}
\nextgroupplot[xlabel={effective step $\eta h/(1-\beta)$},ylabel={stationary loss / first-order value},xmode=log,ymode=log,xmin=0.01,xmax=45,ymin=0.9,ymax=30,
  legend style={at={(0.02,0.98)},anchor=north west,font=\scriptsize},legend cell align=left]
\addplot[theory3] table[x=u,y=r0] {figures/data/05-momentum-variance.dat};
\addplot[theory2] table[x=u,y=r5] {figures/data/05-momentum-variance.dat};
\addplot[theory] table[x=u,y=r9] {figures/data/05-momentum-variance.dat};
\addplot[simulation] table[x=u,y=ratio] {figures/data/05-momentum-variance-sim.dat};
\addplot[guide] coordinates {(2,0.9) (2,30)};
\addplot[guide] coordinates {(6,0.9) (6,30)};
\addplot[guide] coordinates {(38,0.9) (38,30)};
\legend{$\beta=0$ (SGD),$\beta=0.5$,$\beta=0.9$,Monte Carlo}
\end{groupplot}
\end{tikzpicture}
\end{fullwidth}
\caption{The linear theory of momentum. Left: the stability thresholds of \cref{prop:momentum:regions} (curves), and the thresholds measured by running each method as code on $\Loss(x)=x^2/2$ with $401$ step sizes within $2\%$ of the prediction and recording where the growth rate of the state over $8000$ steps turns nonnegative (points; all within $10^{-4}$ of the curves). The dashed line is gradient descent's threshold $2$. Right: the stationary loss of heavy ball divided by its first-order value $\eta\sigma^2/(4(1-\beta))$, the exact formula \eqref{eq:momentum:variance} (curves) against Monte Carlo with $2000$ chains, $2000$ burn-in steps and $6000$ measured steps, $h=\sigma=1$, seed $0$ (points; all within $0.3\%$). Dashed lines: the stability limits $2$, $6$ and $38$. Scripts \texttt{05-momentum-thresholds.py} and \texttt{05-momentum-variance.py}.}
\label{fig:momentum:linear}
\end{figure}

\paragraph{Normalized steps: bias and fluctuations.} \Cref{thm:momentum:normalized} is a worst-case bound, and it is far from the behaviour of typical problems; two exactly analysable examples show in what way. The first is \cref{prop:momentum:floor} with momentum: the quadratic $h\theta^2/2$ on $\R$ with the skewed two-point noise, $\sigma=1$, and sign steps along $m_t$. The second is the opposite extreme, an isotropic quadratic $h\norm\theta_2^2/2$ in dimension $d=1000$ with Gaussian noise and $\ell_2$-normalized steps. In both, the stationary gradient has two parts, which the following derivations compute.

\begin{heuristic}[Heuristic derivation of the floors of a normalized step with memory]
\emph{One dimension.} Write $m_t=h\tilde\theta_t+\momentumzbar_t$, where $\tilde\theta_t$ is the moving average of the iterates, within $\beta\eta/(1-\beta)$ of $\theta_t$, and $\momentumzbar_t$ is the moving average of the noise, a stationary sequence with correlation time about $1/(1-\beta)$ whose law is that of $\momentumzbar=(1-\beta)\sum_{k\ge0}\beta^k\zeta_k$. For small $\eta$ the iterate hardly moves during a correlation time, and over many of them the steps average to $-\eta(1-2F(-h\theta))$, with $F$ the distribution function of $\momentumzbar$. This drift vanishes where $F(-h\theta)=\frac12$: the stationary gradient concentrates, as $\eta\to0$, at the \emph{bias}
\begin{equation}\label{eq:momentum:bias}
  \E\nabla\Loss(\theta)\to-\momentummed(\momentumzbar),
\end{equation}
independent of $\eta$ and $h$. Without memory, $-\momentummed(\zeta)=b=\sigma/\sqrt3$. With memory, $\momentumzbar$ is nearly Gaussian, with variance $\kappa_2=\sigma^2(1-\beta)/(1+\beta)$ and third cumulant $\kappa_3=\E[\zeta^3](1-\beta)^3/(1-\beta^3)$, and the Cornish--Fisher expansion of the median, $\momentummed(\momentumzbar)\approx-\kappa_3/(6\kappa_2)$, gives
\[
  -\momentummed(\momentumzbar)\approx\frac{\E[\zeta^3]}{6\sigma^2}\cdot\frac{(1-\beta)(1+\beta)}{1+\beta+\beta^2}\approx\frac{\E[\zeta^3]}{9\sigma^2}(1-\beta)\qquad(\beta\to1).
\]
Around the rest point the drift is $-2\eta hf\cdot(\theta-\theta_c)$, with $f$ the density of $\momentumzbar$ at its median, and the signs fluctuate with integrated autocovariance $D$. A diffusion approximation gives $\Var\theta\approx\eta D/(4hf)$. For Gaussian $\momentumzbar$, $f=(2\pi\kappa_2)^{-1/2}$ and the signs of a Gaussian sequence with correlations $\beta^k$ have $D=1+\frac4\pi\sum_{k\ge1}\arcsin\beta^k\approx2\ln2/(1-\beta)$, so the \emph{fluctuations} are
\begin{equation}\label{eq:momentum:fluct}
  \Var\nabla\Loss(\theta)\approx\frac{\sqrt{2\pi\kappa_2}}4\,h\eta D\approx\frac{\sqrt{2\pi}\ln2}2\cdot\frac{h\eta\sigma}{\sqrt{1-\beta^2}}.
\end{equation}
\emph{High dimension.} For isotropic Gaussian noise with $\E\norm{\zeta}_2^2=\sigma^2$, the norm $\norm{m_t}_2$ concentrates as $d\to\infty$ at a constant $c$. Replacing it by $c$ makes the iteration linear: EMA momentum with step $\eta/c$ on every coordinate. \Cref{thm:momentum:noise}(iv) and the covariance computed in its proof give the stationary second moments of the coordinates of $\theta_t$ and of $m_t=c(\theta_t-\theta_{t+1})/\eta$, and the condition $\E\norm{m_t}_2^2=c^2$ is the quadratic equation $2(1+\beta)c^2-\eta(1-\beta)hc-2(1-\beta)\sigma^2=0$. Its positive root gives
\begin{equation}\label{eq:momentum:iso}
  \E\norm{\nabla\Loss(\theta)}_2^2=\frac{h\eta c(1+\beta)}{2(1-\beta)}\approx\frac{h\eta\sigma}2\sqrt{\frac{1+\beta}{1-\beta}}\qquad\text{when }\eta h\ll\sigma\sqrt{(1+\beta)/(1-\beta)}.\qedhere
\end{equation}
\end{heuristic}

\Cref{fig:momentum:normalized} compares both derivations with simulations. In one dimension (left), the mean of the stationary gradient follows the exact limit \eqref{eq:momentum:bias}, computed from the distribution of $\momentumzbar$, from $0.577$ without momentum down to $6\cdot10^{-4}$ at $\beta=0.995$: memory reduces the bias linearly in $1-\beta$, much faster than the factor $\sqrt{1-\beta}$ of \cref{thm:momentum:normalized}, which bounds the error of the average rather than the bias of its sign. The standard deviation follows \eqref{eq:momentum:fluct} for $\beta\ge0.9$, and it grows like $\eta^{1/2}(1-\beta)^{-1/4}$: a longer memory makes the sign errors persist longer. At a fixed step size the two parts cross, here near $1-\beta\approx0.05$, and the best $\beta$ balances them, which suggests $1-\beta\propto(h\eta/\sigma)^{2/5}$ in this example (\cref{ex:momentum:bestbeta}); the bound \eqref{eq:momentum:normalized} balances its own lag and noise terms at $1-\beta\approx2(L\eta/(\chi\sigma))^{2/3}$. In high dimension (right), the noise is symmetric and there is no bias to remove; the floor \eqref{eq:momentum:iso}, which matches the simulations to $0.2\%$, rises with $\beta$ at fixed $\eta$. The reason is the one that \cref{thm:momentum:noise} gave for SGD: when $\norm{m_t}$ concentrates, a normalized step is a linear step of size $\eta/\norm{m_t}$, averaging shrinks $\norm{m_t}$ from $\sigma$ to about $\sigma\sqrt{(1-\beta)/(1+\beta)}$, and the effective step grows by the same factor. The two examples are the two halves of the answer to the question of \cref{sec:momentum:wild}. Memory removes the bias of a nonlinear step, which no step size can remove; where the step is effectively linear, memory is a rescaled step size again, and must be paid for with a smaller $\eta$. This is the schedule that \cref{cor:momentum:rate} prescribes: $1-\beta\propto T^{-1/2}$ together with $\eta\propto T^{-3/4}$.

\begin{figure}[t]
\begin{fullwidth}
\centering
\begin{tikzpicture}
\begin{groupplot}[group style={group size=2 by 1,horizontal sep=17mm},wplot,width=82mm,height=58mm,xmode=log,x dir=reverse]
\nextgroupplot[xlabel={$1-\beta$},ylabel={stationary gradient},ymode=log,xmin=0.003,xmax=1.4,ymin=1e-4,ymax=4,
  legend style={at={(0.98,0.02)},anchor=south east,font=\scriptsize},legend cell align=left]
\addplot[theory,unbounded coords=jump] table[x=x,y=bias] {figures/data/05-momentum-skew.dat};
\addplot[theory2] table[x=x,y=sd] {figures/data/05-momentum-skew.dat};
\addplot[guide] table[x=x,y=bound] {figures/data/05-momentum-skew.dat};
\addplot[simulation,error bars/.cd,y dir=both,y explicit] table[x=x,y=mean,y error expr=2*\thisrow{se}] {figures/data/05-momentum-skew-sim.dat};
\addplot[simulation2] table[x=x,y=sd] {figures/data/05-momentum-skew-sim.dat};
\addplot[sea,only marks,mark=o,mark size=2.6pt,line width=0.8pt,forget plot] coordinates {(1,0.57735)};
\legend{exact bias as $\eta\rightarrow0$,fluctuation (heuristic),bound of the theorem,simulated mean,simulated s.d.}
\nextgroupplot[xlabel={$1-\beta$},ylabel={root mean square of $\norm{\nabla\Loss}_2$},ymode=log,xmin=0.004,xmax=1.4,ymin=0.04,ymax=4,
  legend style={at={(0.02,0.98)},anchor=north west,font=\scriptsize},legend cell align=left]
\addplot[theory] table[x=x,y=rmsA] {figures/data/05-momentum-isotropic.dat};
\addplot[theory3] table[x=x,y=rmsB] {figures/data/05-momentum-isotropic.dat};
\addplot[guide] table[x=x,y=boundA] {figures/data/05-momentum-isotropic.dat};
\addplot[simulation] table[x=x,y=rmsA] {figures/data/05-momentum-isotropic-sim.dat};
\addplot[simulation2] table[x=x,y=rmsB] {figures/data/05-momentum-isotropic-sim.dat};
\legend{theory; $\eta=0.01$,theory; $\eta=0.1$,bound of the theorem; $\eta=0.01$,simulated; $\eta=0.01$,simulated; $\eta=0.1$}
\end{groupplot}
\end{tikzpicture}
\end{fullwidth}
\caption{Normalized steps with memory; $1-\beta$ decreases to the right. Left: sign steps with momentum on $\theta^2/2$ with the skewed noise of \cref{prop:momentum:floor} ($\sigma=1$), $\eta=10^{-4}$; $2000$ chains ($4000$ for $\beta\ge0.95$) from $\theta_0=0$, $25000$ burn-in and $40000$ measured steps, seed $0$. Points: mean of the stationary gradient with two standard errors, and its standard deviation. Curves: the exact limit \eqref{eq:momentum:bias} for $\eta\to0$, computed from the fixed-point equation of the distribution of $\momentumzbar$ (drawn for $\beta\ge\frac12$, and at $\beta=0$ as the open circle; for $0<\beta<\frac12$ the law of $\momentumzbar$ is a Cantor-type measure), the Gaussian heuristic \eqref{eq:momentum:fluct}, and the stationary part of the bound \eqref{eq:momentum:normalized}. Right: $\ell_2$-normalized steps with momentum on $\norm\theta_2^2/2$ in $\R^{1000}$ with isotropic Gaussian noise, $\sigma=1$; four runs, $3000$ burn-in and $8000$ measured steps, seed $0$ (points), against \eqref{eq:momentum:iso} (curves), with the stationary part of \eqref{eq:momentum:normalized} for $\eta=0.01$ (dashed). Scripts \texttt{05-momentum-skew.py} and \texttt{05-momentum-isotropic.py}.}
\label{fig:momentum:normalized}
\end{figure}

Where must the agreement fail? The heavy-ball curves of \cref{fig:momentum:linear} are exact for the model, so only Monte Carlo error separates them from the points. The bias \eqref{eq:momentum:bias} is a limit as $\eta\to0$; at $\eta=10^{-3}$ the simulated mean falls below it by a quarter at $\beta=0.99$, where the lag $\beta\eta/(1-\beta)=0.1$ is no longer small against the spread $0.07$ of $\momentumzbar$. The Gaussian fluctuation formula fails for small $\beta$, where $\momentumzbar$ is far from Gaussian; at $\beta=0$ the two-point noise confines the iterate to a lattice next to $b$ and the spread is of order $\eta$. The high-dimensional formula has corrections of relative order $1/d$.

\section{What the theory misses}\label{sec:momentum:misses}

\paragraph{What was idealized away.} The quadratic settings keep the Hessian fixed, so the curvature cannot respond to the step size; in a network it does, and momentum's wider stability region is then not a margin of safety but the place where training settles (\cref{obs:momentum:edge}, \cref{ch:eos}). The noise is additive and independent of the position; near a minimum, minibatch noise scales with the loss and roughly with the curvature, which the noisy quadratic model of \cref{ch:noise} reintroduces as $C\propto H$. The noise has finite variance, which gradient noise in deep networks may not have \parencite{SimsekliSagunGurbuzbalaban2019}. The theorem on normalized steps is a bound on the average gradient over a run, with constants that are worst-case over all smooth losses and noise laws; the examples of \cref{fig:momentum:normalized} show that it can overstate the floor by orders of magnitude and can get the dependence on $\beta$ wrong in sign. And $\beta$ and $\eta$ are constant, while practice decays $\eta$, sometimes warms $\beta$ up \parencite{SutskeverEtAl2013}, and couples the two through bias correction.

\paragraph{Where it breaks.} (1) \Cref{thm:momentum:heavyball} is about quadratics: on a strongly convex function that is a quadratic in each of three pieces, heavy ball with Polyak's parameters settles on a cycle of period three \parencite{LessardRechtPackard2016} (\cref{ex:momentum:lessard}), and with noise the accelerated parameters raise the floor (\cref{ex:momentum:noisyaccel}). Acceleration is not why momentum is used in deep learning: \textcite{KidambiEtAl2018} show that the classical schemes do not accelerate SGD on stochastic least squares, and \cref{thm:momentum:noise} shows why. (2) The rates of \cref{cor:momentum:rate} prescribe $1-\beta\propto T^{-1/2}$, that is, $\beta$ between $0.999$ and $0.9999$ for runs of $10^6$ to $10^8$ steps, while practice keeps $\beta_1=0.9$ for runs of every length. Either the worst-case trade-off between noise and lag is not the one that sets $\beta$ in practice, or the losses of practice are far from worst case; the theory does not say which. (3) With heavy-tailed noise, an average of $1/(1-\beta)$ gradients concentrates more slowly than $\sqrt{1-\beta}$ and the variance assumption fails (\cref{ex:momentum:heavytail}). (4) \Cref{prop:momentum:floor} needs skewed noise. With symmetric noise a normalized step without memory is never biased uphill, its floor vanishes as $\eta\to0$, and the right panel of \cref{fig:momentum:normalized} shows memory raising the floor at a fixed step; so \cref{obs:momentum:normalized} is explained by the theory only to the extent that real gradient noise is asymmetric in the coordinates that a normalized step treats separately, or that its effective dimension is low. How much of the observed need for memory comes from bias, and how much from the effective step size, has not been measured as far as we know. (5) The linear theory says that in Adam's convention momentum changes neither the noise floor nor the rate of the flat coordinates to first order; in practice, $\beta_1$ is tuned jointly with the learning rate, and the interaction with the edge of stability, where the factor $(1+\beta)/(1-\beta)$ matters, may be the more important effect (\cref{ch:eos}).

\paragraph{Repairs and open problems.} Variance reduction replaces the plain average by one that corrects for the lag: STORM \parencite{CutkoskyOrabona2019} adds the gradient difference at consecutive points to the average, and under a mean-squared smoothness assumption reaches $\E\norm{\nabla\Loss}\le\eps$ in $\tilde O(\eps^{-3})$ steps instead of $O(\eps^{-4})$; \textcite{CutkoskyMehta2020} already improve the rate of normalized momentum under second-order smoothness. AdEMAMix \parencite{PagliardiniAblinGrangier2025} mixes a fast average with a very slow one, $\beta=0.9$ and about $0.9999$, and reports that gradients thousands of steps old still help language-model training, which the worst-case lag term of \eqref{eq:momentum:normalized} does not anticipate. Averaging iterates instead of gradients, in schedule-free methods \parencite{DefazioEtAl2024} and model averages \parencite{IzmailovEtAl2018}, attacks the noise floor without the lag of the step. General norms are treated by \textcite{PethickEtAl2025}, whose analysis covers Muon-type steps. Open: a theory that predicts the value $0.9$, or explains why it transfers across scales; a theory of normalized momentum under heavy-tailed, state-dependent noise that is sharp for typical rather than worst-case problems; and a measurement of how biased the normalized steps of real optimizers are.

\section{Back in the wild}\label{sec:momentum:back}

\fieldnotes{Average the gradients with $\beta\approx0.9$ (SGD, Adam's $\beta_1$, Lion), $0.95$ in Muon; every normalized optimizer carries momentum.}
{Quadratics with exact or additively noisy gradients; normalized steepest descent in a general norm with a smooth loss and bounded-variance noise.}
{Polyak's acceleration with explicit constants and its optimality; the stability regions of heavy ball, Nesterov and EMA momentum; the exact stationary variance; a $T^{-1/4}$ rate for normalized steepest descent with momentum in every norm; a floor of order $\sigma$ without memory; averaging in place of a schedule.}
{Thresholds $2(1+\beta)$, $38$ and $1.36$ at $\beta=0.9$; momentum as the step $\eta/(1-\beta)$ when noise dominates; for normalized steps, a bias that falls like $1-\beta$ and fluctuations that grow like $\eta^{1/2}(1-\beta)^{-1/4}$.}
{Curvature that responds to the step, state-dependent and heavy-tailed noise, schedules of $\beta$, and a reason for the value $0.9$.}
{Proved in the idealization: with PyTorch's convention, retune $\eta$ by $1-\beta$ when $\beta$ changes; keep memory in normalized optimizers, since it removes a bias that no step size removes. Suggested: increase $\beta$ as the step size falls. Folklore: $\beta=0.9$ for everything.}

What the theory tells a practitioner is mostly a matter of bookkeeping, and it is proved. With PyTorch's convention, the effective learning rate of SGD with momentum is $\eta/(1-\beta)$, so a sweep over $\beta$ at a fixed $\eta$ is mostly a sweep over the learning rate; with Adam's convention it is not, and the visible effect of $\beta_1$ is on stability, where the threshold grows like $2(1+\beta)/(1-\beta)$. For normalized optimizers, memory is not optional, and its value is not a lower floor at a fixed step size but the removal of a bias; when $\beta$ is raised, $\eta$ should come down, in the ratio $\eta\propto(1-\beta)^{3/2}$ of \cref{cor:momentum:rate} if the worst case is a guide. The theory does not say that $0.9$ is right, and the success of a single value across scales and architectures remains folklore with a good track record.

The thresholds of this chapter return in \cref{ch:eos} as the places where training settles, and the noisy quadratic model returns in \cref{ch:noise}, where momentum, like preconditioning, moves the critical batch size. The bias of normalized steps connects to the small-batch failure of sign descent in \cref{ch:sign} and to Muon in \cref{ch:spectral}; Adam's $\beta_1$ is the same average, with a preconditioner, in \cref{ch:adaptive}. Averaging the iterates is the other half of \cref{ch:schedules}, where a cooldown and a model average compete to remove the same noise floor, and Lion's and Muon's weight decay make the methods of this chapter constrained ones in \cref{ch:decay}.

\begin{exercises}
\item \consolidation (The band of complex roots.) Show that the roots of $z^2-(1+\beta-\eta h)z+\beta$ are complex or double exactly when $(1-\sqrt\beta)^2\le\eta h\le(1+\sqrt\beta)^2$, and that they then have modulus $\sqrt\beta$. For $\beta=0.9$, deduce that every coordinate with $0.0026\le\eta h\le3.797$ contracts at the same rate $\sqrt{0.9}\approx0.949$, and explain why the best heavy-ball parameters for the spectrum $[\mu,L]$ are those for which this band is $[\eta\mu,\eta L]$.\label{ex:momentum:band}

\item \consolidation (More regions.) (a) Verify \cref{lem:momentum:forms}(iv) for the PyTorch form of Nesterov momentum. (b) For the two-coefficient update $u_t=\beta_am_{t-1}+(1-\beta_a)g_t$ with $m_t=\beta_bm_{t-1}+(1-\beta_b)g_t$ and $0\le\beta_b<\beta_a<1$, show that the stability threshold for $\eta h$ is $\min\{(1-\beta_b)/(\beta_a-\beta_b),\,2(1+\beta_b)/(1+\beta_b-2\beta_a)\}$, where the second term is omitted when $1+\beta_b\le2\beta_a$. (For $\beta_a\le\beta_b$ this update is quasi-hyperbolic momentum with $\beta=\beta_b$ and $\nu=\beta_a/\beta_b$.)\label{ex:momentum:regions}

\item \consolidation (Bias correction.) (a) Adam divides $m_t$ by $1-\beta^{t+1}$. Show that, with this correction, the first step of EMA momentum is $-\eta g_0$, and that the correction changes the steps of a linear method by a factor that tends to one with a delay of about $1/(1-\beta)$ steps. (b) Show that it does not change the iterates of normalized steepest descent with momentum at all.\label{ex:momentum:bias}

\item \consolidation (The noisy quadratic model with momentum.) Under \cref{ass:momentum:quadratic} with $C=\diag(c_1,\dots,c_d)/B$ in the eigenbasis, show that EMA momentum is stable if and only if $\eta h_1<2(1+\beta)/(1-\beta)$, and that its stationary excess loss is $\sum_i\frac{\eta c_i}{4B}\bigl(1-\frac{\eta(1-\beta)h_i}{2(1+\beta)}\bigr)^{-1}$. Compare with gradient descent at the same $\eta$ (\cref{ch:noise}).\label{ex:momentum:nqm}

\item \consolidation (Tuning the bound.) (a) In \eqref{eq:momentum:normalized}, take $1-\beta=\sqrt{2L\Delta_0}/(\chi\sigma\sqrt T)$, assumed at most $1$, and $\eta=\sqrt{\Delta_0(1-\beta)/(2LT)}$, and show that the first four terms add up to at most $4(2L\Delta_0)^{1/4}(\chi\sigma)^{1/2}T^{-1/4}+L\eta/2$. (b) At a fixed $\eta$ and for $\beta$ close to $1$, approximate the $\beta$-dependent part of \eqref{eq:momentum:normalized} by $2L\eta/(1-\beta)+\chi\sigma\sqrt{2(1-\beta)}$, and show that this is smallest at $1-\beta=2(L\eta/(\chi\sigma))^{2/3}$.\label{ex:momentum:tuning}

\item \variation (Lion and Nesterov in the theorem.) For quasi-hyperbolic directions $\theta_{t+1}=\theta_t+\eta\lmo(u_t)$, $u_t=(1-\nu)g_t+\nu m_t$, show that \cref{thm:momentum:normalized} holds with the lag term $2\nu\beta L\eta/(1-\beta)$, the noise factor $\sqrt{(1-\nu\beta)^2+\nu^2\beta^2(1-\beta)/(1+\beta)}$ in place of $\sqrt{(1-\beta)/(1+\beta)}$, and $\nu$ times the last term. Evaluate the noise factor and the lag coefficient $2\nu\beta/(1-\beta)$ for Signum with $\beta=0.9$ ($0.229$; $18$), Lion with $(\beta_1,\beta_2)=(0.9,0.99)$, that is $\nu\beta=0.9$ and $\beta=0.99$ ($0.119$; $180$), and Muon's Nesterov momentum with $\beta=0.95$ ($0.174$; $36.1$).\label{ex:momentum:lion}

\item \variation (Symmetric noise.) (a) Show that if $\zeta$ and $-\zeta$ have the same law, then $\E\ip{a,u(a+\zeta)}\le0$ for every $a$, every norm and every measurable selection $u(g)\in\lmo(g)$. \hint{pair $\zeta$ with $-\zeta$ and use $\ip{a\pm\zeta,u(a\pm\zeta)}=-\norm{a\pm\zeta}_*$.} (b) For sign descent with independent Gaussian noise of standard deviations $s_i$, show that $\E\ip{a,\sign(a+\zeta)}=\sum_i\abs{a_i}\erf\bigl(\abs{a_i}/(\sqrt2s_i)\bigr)$, so that for $\abs{a_i}\ll s_i$ the expected step is that of gradient descent preconditioned by $\sqrt{2/\pi}\,\diag(1/s_i)$. With momentum, and treating the noise of the average as independent of the current point, show that the preconditioner is multiplied by $\sqrt{(1+\beta)/(1-\beta)}$.\label{ex:momentum:symmetric}

\item \variation\hard (The noise term cannot be improved.) Let $0<\beta<1$, $p=(1-\beta)/10$, and let $\zeta$ equal $\sigma\sqrt{(1-p)/p}$ with probability $p$ and $-\sigma\sqrt{p/(1-p)}$ otherwise. Show that $\E\zeta=0$, $\E\zeta^2=\sigma^2$, and that the stationary average $\momentumzbar=(1-\beta)\sum_{k\ge0}\beta^k\zeta_k$ has $\momentummed(\momentumzbar)\le-\sigma\sqrt{1-\beta}/6$. Explain, with the heuristic \eqref{eq:momentum:bias}, why this shows that the factor $\sqrt{1-\beta}$ in \cref{thm:momentum:normalized} cannot be improved for all noise laws with variance $\sigma^2$. \hint{condition on the event that none of the last $K=\lceil3/(1-\beta)\rceil$ samples is large, whose probability is at least $0.6$ by Bernoulli's inequality, and bound the contribution of the older samples, which carry the weight $\beta^K\le e^{-3}$, with Chebyshev's inequality.}\label{ex:momentum:worst}

\item \variation (A model average.) Let $x_t$ follow gradient descent with step $\eta$ and noise of variance $\sigma^2$ on a coordinate with curvature $h$, in its stationary state, and let $\bar x_t=\beta_e\bar x_{t-1}+(1-\beta_e)x_t$. (a) Show that the stationary variance of $\bar x$ is $V\frac{1-\beta_e}{1+\beta_e}\cdot\frac{1+a\beta_e}{1-a\beta_e}$, with $a=1-\eta h$ and $V=\eta\sigma^2/(h(2-\eta h))$. (b) Show that it is close to $\sigma^2(1-\beta_e)/(2h^2)$, the variance of gradient descent with step $(1-\beta_e)/h$, when $1-\beta_e\ll\eta h\ll1$, and close to $V$ when $\eta h\ll1-\beta_e$; show that a drifting coordinate is tracked with a lag of $\beta_e/(1-\beta_e)$ steps.\label{ex:momentum:modelema}

\item \simulation (Acceleration drowns in noise.) On a quadratic in $\R^{100}$ with eigenvalues $h_i$ spaced geometrically from $10^{-3}$ to $1$ and $\theta_0-\theta^\star=(1,\dots,1)$, compare gradient descent with step $2/(L+\mu)$ and heavy ball with Polyak's parameters. Without noise, observe that heavy ball reaches $10^{-8}$ times the initial loss in about $220$ steps and gradient descent in about $3900$. With independent Gaussian noise of standard deviation $10^{-2}$ in each coordinate, observe that both stall, heavy ball at a loss about sixteen times higher (about $0.91$ against $0.057$), and derive both floors coordinate by coordinate from \cref{thm:momentum:noise}.\label{ex:momentum:noisyaccel}

\item \simulation (The best memory at a fixed step.) With \texttt{05-momentum-skew.py}, measure $\E\abs{\nabla\Loss}$ for $\eta\in\{10^{-3},10^{-4},10^{-5}\}$ and $\beta$ on a fine grid. Check that the bias does not depend on $\eta$, that the standard deviation scales like $\eta^{1/2}$, and that the best $1-\beta$ shrinks roughly like $\eta^{2/5}$, as balancing \eqref{eq:momentum:bias} against \eqref{eq:momentum:fluct} suggests.\label{ex:momentum:bestbeta}

\item \simulation (Heavy ball beyond quadratics.) Let $F(x)=\frac{25}2x^2$ for $x<1$, $F(x)=\frac12x^2+24x-12$ for $1\le x<2$, and $F(x)=\frac{25}2x^2-24x+36$ for $x\ge2$. Check that $F$ is $1$-strongly convex and $25$-smooth. Run heavy ball with Polyak's parameters for $\mu=1$ and $L=25$, that is $\eta=\frac19$ and $\beta=\frac49$, from $x_0=x_{-1}=3.3$, and observe that the iterates settle on a cycle through about $0.6465$, $-1.8024$ and $2.1159$. Run Nesterov's method with $\eta=1/L$ and $\beta=(\sqrt\kappa-1)/(\sqrt\kappa+1)$ from the same point and observe that it converges.\label{ex:momentum:lessard}

\item \research (Why $0.9$?) The theory of this chapter sends $\beta$ to one as the run gets longer (\cref{cor:momentum:rate}) or the step smaller (\cref{fig:momentum:normalized}), and practice keeps $\beta_1=0.9$ from small models to the largest. Build a model in which a fixed $\beta$ is optimal across scales. Candidates are the edge of stability, where a larger $\beta$ widens the stable range but slows the response of the average to a change of curvature (\cref{ch:eos}); learning-rate schedules, which change the balance of lag and noise during a run (\cref{ch:schedules}); and the batch-size regime (\cref{ch:noise}). As far as we know, no model predicts the value.\label{ex:momentum:why}

\item \research (Heavy-tailed noise.) Suppose that the noise is $\alpha$-stable with $1<\alpha<2$, mean zero and skewness parameter not zero, so that its variance is infinite. (a) Show that the average $\momentumzbar$ is again $\alpha$-stable, with the same skewness and a scale multiplied by $(1-\beta)(1-\beta^\alpha)^{-1/\alpha}\approx\alpha^{-1/\alpha}(1-\beta)^{1-1/\alpha}$, and deduce from \eqref{eq:momentum:bias} that the bias of a sign step falls only like $(1-\beta)^{1-1/\alpha}$. (b) Open: a convergence theory of normalized steepest descent with momentum under noise with only a $p$-th moment, $1<p<2$, that is sharp for the noise of real networks.\label{ex:momentum:heavytail}
\end{exercises}

\notesandrefs
Momentum is as old as accelerated iterative methods. \Textcite{Polyak1964} introduced the heavy-ball method and its local rate; \textcite{Nesterov1983} found the accelerated method that bears his name and its global rate for smooth convex functions; both are in the textbooks of \textcite{Polyak1987} and \textcite{Nesterov2018}. \Textcite{RumelhartHintonWilliams1986} used a momentum term to train networks by back-propagation, and \textcite{SutskeverEtAl2013} showed that momentum with a well-chosen schedule and initialization trains deep and recurrent networks that were thought to need second-order methods, writing Nesterov's method as a momentum update with a look-ahead gradient.

The equivalence of momentum and a larger step in the presence of noise was observed and analysed several times. \Textcite{FlammarionBach2015} treat averaging and acceleration as second-order linear recursions that differ by the step size; \textcite{KidambiEtAl2018} prove that the classical schemes do not accelerate SGD on stochastic least squares and propose a scheme that does; \textcite{ZhangEtAl2019} make the noisy quadratic model a tool for predicting which algorithmic choices matter at which batch sizes; \textcite{ShallueEtAl2019} provide the measurements; and \textcite{WangEtAl2024} prove an equivalence for small learning rates in the sense of stochastic differential equations. \Cref{thm:momentum:noise} is the exact two-dimensional computation behind these results.

Normalized SGD was analysed by \textcite{HazanLevyShalevShwartz2015} with growing minibatches, and signSGD and Signum by \textcite{BernsteinEtAl2018}, under assumptions that include large batches or unimodal symmetric noise; \textcite{KarimireddyEtAl2019} gave examples where sign compression fails and repaired it with error feedback. \Textcite{CutkoskyMehta2020} proved the Euclidean case of \cref{thm:momentum:normalized}, with momentum and without large batches, and an improved rate under second-order smoothness; general norms, with Muon and sign descent as instances, are treated by \textcite{PethickEtAl2025}. Our proof is a short version of this argument, organized around \cref{lem:momentum:direction,lem:momentum:error}. The lower bound $\Omega(\eps^{-4})$ for stochastic first-order methods is due to \textcite{ArjevaniEtAl2023}. We know of no reference for \cref{prop:momentum:floor} in this form, for the median formula \eqref{eq:momentum:bias}, or for the self-consistent floor \eqref{eq:momentum:iso}; the last two are derived, not proved.

Averaging the iterates of stochastic approximation is due to \textcite{Ruppert1988} and \textcite{PolyakJuditsky1992}; stochastic weight averaging brought it to deep learning \parencite{IzmailovEtAl2018}. \Cref{prop:momentum:averaging} is the anytime online-to-batch conversion of \textcite{Cutkosky2019}, which \textcite{DefazioEtAl2024} combine with an interpolation to obtain schedule-free methods. The stability thresholds of \cref{prop:momentum:regions} are classical for heavy ball; the threshold $38$ for Adam's first moment was pointed out by \textcite{CohenEtAl2022}, and quasi-hyperbolic momentum, which contains all the forms, is due to \textcite{MaYarats2019}.

\paragraph{Translating the literature.} \Textcite{Polyak1987} writes the heavy-ball method as $x_{k+1}=x_k-\alpha\nabla f(x_k)+\beta(x_k-x_{k-1})$: his $\alpha$ is our $\eta$. \Textcite{SutskeverEtAl2013} write $v_{t+1}=\mu v_t-\varepsilon\nabla f(\theta_t+\mu v_t)$, $\theta_{t+1}=\theta_t+v_{t+1}$ for Nesterov's method: their $\mu$ is our $\beta$, their $\varepsilon$ our $\eta$, and $\theta_t+\mu v_t$ is our $y_t$. PyTorch's \texttt{momentum} is our $\beta$ in the PyTorch forms of \cref{def:momentum:forms}, with \texttt{dampening} zero. \Textcite{ChenEtAl2023} write Lion with coefficients $\beta_1$ and $\beta_2$, which are not Adam's: their $\beta_2$ is our $\beta$ and their $\beta_1$ is our $\nu\beta$. \Textcite{MaYarats2019} write quasi-hyperbolic momentum with our $\nu$ and $\beta$ and a step size $\alpha$. \Textcite{DefazioEtAl2024} use $\beta$ for the interpolation of schedule-free methods, which is not a momentum coefficient on gradients, and $\gamma$ for the step size; we keep their letters in the field guide. \Textcite{CutkoskyMehta2020} state their result for the Euclidean norm, where our $\chi$ is $1$.

\printchapterbib
\end{refsection}

% Chapter 6 of Optimizers in the Wild: adaptive preconditioning (AdaGrad, RMSProp, Adam).
\newcommand{\adaptiveclip}{\operatorname{clip}}
\newcommand{\adaptiveKL}{\operatorname{KL}}

\casechapter{Adaptive Preconditioning}{AdaGrad, RMSProp and Adam}{Adam's $\beta_2$ and $\epsilon$}{ch:adaptive}
\begin{refsection}[chapters/06-adaptive.bib]

\begin{chaptersummary}
Adam divides each coordinate of its step by a running root mean square of that coordinate's gradients; $\beta_2$ sets how fast the running average forgets, and $\epsilon$ the gradient size below which it stops dividing. A step size per coordinate is a norm, so these methods are steepest descent in a diagonal norm that the gradients choose. We prove that AdaGrad, whose norm never forgets, has regret within $\sqrt2$ of the best diagonal norm chosen in hindsight, and that its scalar version converges at every step size without knowing the smoothness constant. We then prove what forgetting costs: Adam converges to the worst point of a one-dimensional convex problem when its memory $1/(1-\beta_2)$ is short compared with the period of a rare large gradient, and after a quiet stretch a long memory gives steps of up to $6.6$ step sizes and an overshoot of about $2\ln2/(1-\beta_2)$ steps. Adafactor's factored second moment is the best rank-one approximation in I-divergence. The complete theory is worst case and convex; it does not say why $\beta_2=0.95$ works for language models.
\end{chaptersummary}

\section{In the wild: Adam's two dials}\label{sec:adaptive:wild}

Adam divides the momentum of each coordinate by a running root mean square of that coordinate's gradients. Two hyperparameters set the denominator: $\beta_2$ fixes how fast the running average forgets, and $\epsilon$ fixes a floor below which a gradient is no longer divided by its own size. Both are rarely tuned, and the values in use have moved.

\begin{practice}[Adam, \textcite{KingmaBa2015}; GPT-3, \textcite{BrownEtAl2020}; LLaMA, \textcite{TouvronEtAl2023}]
Train with Adam or AdamW (\cref{ch:decay}), $\beta_1=0.9$, bias correction on. Keep $\beta_2=0.999$, the default of Kingma and Ba, except for large language models, which use $\beta_2=0.95$. Keep $\epsilon=10^{-8}$, the default of the original paper and of common libraries. For sparse convex models, such as click prediction, use AdaGrad's per-coordinate learning rate $\eta/(\sum_{s\le t}g_{s,i}^2)^{1/2}$. To save memory, Adafactor stores row and column sums of squared gradients instead of $v$.
\end{practice}

For language models the value of $\beta_2$ drifted down: RoBERTa lowered it from $0.999$ to $0.98$, reporting better stability with large batches \parencite{LiuEtAl2019}, and GPT-3 and LLaMA used $0.95$. PaLM moved the other way: it used Adafactor without factorization and $\beta_2=1-k^{-0.8}$ at step $k$, which its authors found more stable than $\beta_2=0.99$ because the second moments of rare embedding tokens are poorly estimated over short windows \parencite{ChowdheryEtAl2022}.\sidenote{The increasing schedule $1-k^{-0.8}$ is from \textcite{ShazeerStern2018}, who also argued that a slowly decaying second-moment estimate can be out of date, and proposed to clip the root mean square of each update.}

\begin{observation}[$\beta_2$ and stability; \textcite{LiuEtAl2019}, \textcite{ChowdheryEtAl2022}, \textcite{MolybogEtAl2023}]\label[observation]{obs:adaptive:beta2}
In the training of large language models (2019--2023), lowering $\beta_2$ from $0.999$ to $0.98$ or $0.95$ has been reported to improve stability, and so has raising it toward $1$ for rare tokens; \textcite{MolybogEtAl2023} relate loss spikes in large runs to Adam updates that become large and nearly uncorrelated with the gradient. These are reports from training runs or an unrefereed preprint; none isolates a mechanism by an ablation across scales.
\end{observation}

The reports point in opposite directions, the first hint that $\beta_2$ trades one failure for another.

\begin{observation}[\textcite{ReddiKaleKumar2018}; \textcite{ZhangEtAl2022}]\label[observation]{obs:adaptive:reddi}
For every $\beta_1<\sqrt{\beta_2}<1$ there is an online convex problem on an interval on which Adam's average regret does not tend to zero, and a stochastic convex problem on which Adam does not converge to the minimizer \parencite{ReddiKaleKumar2018}. Adam remained the default, and the proposed fix, AMSGrad, did not replace it. \Textcite{ZhangEtAl2022} pointed to the order of choices: the counterexamples are built after the hyperparameters, practitioners tune the hyperparameters after the problem, and with the problem fixed first and $\beta_2$ large enough Adam converges.
\end{observation}

The dial $\epsilon$ matters more than its size suggests: \textcite{WortsmanEtAl2024} report that in small transformers the default $10^{-8}$ becomes comparable to some gradients as the models are made wider, which shrinks the updates, and that a much smaller $\epsilon$ avoids this.

AdaGrad came first. \Textcite{McMahanStreeter2010} and \textcite{DuchiHazanSinger2011} scaled each coordinate's step by the inverse root of its accumulated squared gradients, and the method spread through large sparse problems such as click prediction \parencite{McMahanEtAl2013}. For deep networks the accumulated sum grows without bound and ends learning too early; RMSProp replaced it by an exponential moving average \parencite{TielemanHinton2012},\sidenote{RMSProp was never published as a paper; it appeared in a lecture of Hinton's 2012 online course, which is how it is still cited.} and Adam added momentum and a bias correction \parencite{KingmaBa2015}. The step that never forgets has a guarantee and was abandoned for deep networks; the step that forgets has a counterexample and is the default. Which property of the denominator does each guarantee use, and what do $\beta_2$ and $\epsilon$ control?

\section{The idealization: a norm chosen by the gradients}\label{sec:adaptive:idealization}

A step size per coordinate is a norm. For a symmetric positive definite matrix $P$ write $\norm{u}_P=(u^\top Pu)^{1/2}$.

\begin{lemma}[Preconditioning is steepest descent in a quadratic norm]\label[lemma]{lem:adaptive:quadnorm}
Let $P\succ0$. The dual norm of $\norm{\cdot}_P$ is $\norm{\cdot}_{P^{-1}}$; for $g\neq0$, $\lmo(g)=-P^{-1}g/\norm{g}_{P^{-1}}$ is unique; and $g^\sharp=P^{-1}g$. So $\theta-\eta P^{-1}g$ is unnormalized steepest descent with respect to $\norm{\cdot}_P$, and for $P=\diag(p)$ it is gradient descent with step size $\eta/p_i$ in coordinate $i$.
\end{lemma}

\begin{proof}
By the Cauchy--Schwarz inequality for $P^{-1/2}g$ and $P^{1/2}u$, $\ip{g,u}\ge-\norm{g}_{P^{-1}}\norm{u}_P$, with equality exactly when $u=-cP^{-1}g$ for some $c\ge0$. Over the unit ball the minimum of $\ip{g,u}$ is therefore $-\norm{g}_{P^{-1}}$, attained only at $-P^{-1}g/\norm{g}_{P^{-1}}$, a point of norm one since $\norm{P^{-1}g}_P=\norm{g}_{P^{-1}}$. Then $g^\sharp=-\norm{g}_{P^{-1}}\lmo(g)=P^{-1}g$, which also holds for $g=0$.
\end{proof}

We write squares, roots and quotients of vectors coordinatewise, count rounds from $t=1$, and call $\eta$ the step size; practitioners call it the learning rate.

\begin{definition}[Three adaptive methods]\label[definition]{def:adaptive:methods}
Let $g_t$ be the gradient or a stochastic gradient at $\theta_t$, $\eta>0$, $\epsilon\ge0$; accumulators start at $0$, and $0/0=0$.
\begin{enumerate}
\item \term{AdaGrad}\index{AdaGrad}: $s_t^2=s_{t-1}^2+g_t^2$ and $\theta_{t+1}=\theta_t-\eta\,g_t/(s_t+\epsilon)$. Thus $s_{t,i}=\norm{g_{1:t,i}}_2$, with $g_{1:t,i}=(g_{1,i},\dots,g_{t,i})$ the history of coordinate $i$.
\item \term{RMSProp}\index{RMSProp}: $v_t=\beta_2v_{t-1}+(1-\beta_2)g_t^2$ and $\theta_{t+1}=\theta_t-\eta\,g_t/(\sqrt{v_t}+\epsilon)$, $0\le\beta_2<1$.
\item \term{Adam}\index{Adam}: $m_t=\beta_1m_{t-1}+(1-\beta_1)g_t$, $v_t$ as in RMSProp, the \term{bias-corrected}\index{bias correction} moments $\hat m_t=m_t/(1-\beta_1^t)$, $\hat v_t=v_t/(1-\beta_2^t)$, and $\theta_{t+1}=\theta_t-\eta\,\hat m_t/(\sqrt{\hat v_t}+\epsilon)$, $0\le\beta_1<1$.
\end{enumerate}
On a constraint set each update is followed by a projection. On a box, the projection in every diagonal quadratic norm is coordinatewise clipping, because $\sum_ip_i(\theta_i-y_i)^2$ separates.
\end{definition}

By \cref{lem:adaptive:quadnorm}, AdaGrad is unnormalized steepest descent in $\norm{\cdot}_{P_t}$ with $P_t=\diag(s_t+\epsilon)$, RMSProp uses $P_t=\diag(\sqrt{v_t}+\epsilon)$, and Adam uses the bias-corrected matrix and applies the step to the moving average $\hat m_t$ of \cref{ch:momentum}. One property will decide everything: AdaGrad's $P_t$ is nondecreasing in $t$, and those of RMSProp and Adam are not. A moving average with coefficient $\beta_2$ gives the gradient of $k$ steps ago the weight $(1-\beta_2)\beta_2^k$; we call $1/(1-\beta_2)$, the mean lag plus one, its \term{memory}\index{memory of a moving average}: $20$ steps for $\beta_2=0.95$, $1000$ for $0.999$.\sidenote{The half-life of the weights, $\ln2/\ln(1/\beta_2)$, is about $0.69$ memories: $13.5$ steps for $\beta_2=0.95$ and $693$ for $0.999$.}

\begin{fieldguide}{AdaGrad}{\textcite{DuchiHazanSinger2011}; \textcite{McMahanStreeter2010}}
\fgrow{update}{$s_t^2=s_{t-1}^2+g_t^2$, \ $\theta_{t+1}=\theta_t-\eta\,g_t/(s_t+\epsilon)$}
\fgrow{norm}{$\norm\cdot_{P_t}$, $P_t=\diag(s_t+\epsilon)$ (\cref{lem:adaptive:quadnorm})}
\fgrow{preconditioner}{diagonal, nondecreasing; the best diagonal norm in hindsight up to $\sqrt2$ (\cref{prop:adaptive:hindsight})}
\fgrow{averages}{none: a running sum of $g^2$}
\fgrow{state}{one number per parameter}
\fgrow{in the wild}{sparse convex models; AdaGrad-Norm (\cref{thm:adaptive:adagradnorm}); full matrix in \cref{ch:kronecker}}
\end{fieldguide}

\begin{fieldguide}{RMSProp}{\textcite{TielemanHinton2012}}
\fgrow{update}{$v_t=\beta_2v_{t-1}+(1-\beta_2)g_t^2$, \ $\theta_{t+1}=\theta_t-\eta\,g_t/(\sqrt{v_t}+\epsilon)$}
\fgrow{norm}{$\norm\cdot_{P_t}$, $P_t=\diag(\sqrt{v_t}+\epsilon)$; $\ell_\infty$ if $\beta_2=\epsilon=0$ (sign descent, \cref{ch:sign})}
\fgrow{preconditioner}{diagonal; forgets with memory $1/(1-\beta_2)$}
\fgrow{averages}{EMA of $g^2$}
\fgrow{state}{one number per parameter}
\fgrow{in the wild}{$\beta_2=0.9$ in the lecture that introduced it}
\end{fieldguide}

\begin{fieldguide}{Adam}{\textcite{KingmaBa2015}}
\fgrow{update}{$m_t=\beta_1m_{t-1}+(1-\beta_1)g_t$, \ $v_t=\beta_2v_{t-1}+(1-\beta_2)g_t^2$, \ $\theta_{t+1}=\theta_t-\eta\,\hat m_t/(\sqrt{\hat v_t}+\epsilon)$, \ $\hat m_t=m_t/(1-\beta_1^t)$, $\hat v_t=v_t/(1-\beta_2^t)$}
\fgrow{norm}{$\ell_\infty$ if $\beta_1=\beta_2=\epsilon=0$ (\cref{ch:sign}); in general $\norm\cdot_{P_t}$, $P_t=\diag(\sqrt{\hat v_t}+\epsilon)$}
\fgrow{preconditioner}{diagonal; $\to\diag(\E g^2)^{1/2}$ for stationary gradients as $\beta_2\to1$ (\cref{rem:adaptive:limits})}
\fgrow{averages}{EMAs of $g$ ($\beta_1$, \cref{ch:momentum}) and of $g^2$ ($\beta_2$), bias-corrected}
\fgrow{state}{two numbers per parameter}
\fgrow{step}{bounded per coordinate if $\beta_1^2<\beta_2$ (\cref{ch:sign}); up to $6.6\eta$ after a quiet stretch at $(0.9,0.999)$ (\cref{prop:adaptive:quietloud})}
\fgrow{in the wild}{$\beta_1=0.9$; $\beta_2=0.999$, or $0.95$ for language models; $\epsilon=10^{-8}$}
\fgrow{caution}{fails if the memory is short compared with the period of rare large gradients (\cref{prop:adaptive:reddi})}
\end{fieldguide}

\begin{fieldguide}{AdamW}{\textcite{LoshchilovHutter2019}}
\fgrow{update}{$\theta_{t+1}=(1-\eta\lambda)\theta_t-\eta\,\hat m_t/(\sqrt{\hat v_t}+\epsilon)$}
\fgrow{preconditioner}{as Adam; the decay is not preconditioned, unlike an $L^2$ penalty added to $\Loss$}
\fgrow{state}{two numbers per parameter}
\fgrow{in the wild}{$\lambda=0.1$ in GPT-3 and LLaMA; full entry in \cref{ch:decay}}
\end{fieldguide}

\begin{fieldguide}{Adafactor}{\textcite{ShazeerStern2018}}
\fgrow{update}{for $W\in\R^{m\times n}$: $r_t=\beta_2r_{t-1}+(1-\beta_2)G_t^2\mathbf1$, \ $c_t=\beta_2c_{t-1}+(1-\beta_2)(G_t^2)^\top\mathbf1$, \ $\hat V_t=r_tc_t^\top/(\mathbf1^\top r_t)$, step $G_t/\hat V_t^{1/2}$ entrywise}
\fgrow{preconditioner}{diagonal; the best rank-one approximation of Adam's $v$ in I-divergence (\cref{prop:adaptive:adafactor})}
\fgrow{averages}{EMAs of row and column sums of $G^2$, $\beta_2=1-t^{-0.8}$; often no momentum}
\fgrow{state}{$m+n$ numbers per $m\times n$ matrix}
\fgrow{in the wild}{PaLM, without factorization; with update clipping and relative step sizes}
\end{fieldguide}

The theory uses two settings. The first is worst case: convex losses that may change at every round, chosen by an adversary who knows the method; the gradients of a training run, which change as the parameters move, are the motivation.

\begin{assumption}[Online convex optimization on a box]\label[assumption]{ass:adaptive:oco}
The feasible set is a box $\Theta=\prod_i[a_i,b_i]$ with sides $D_i=b_i-a_i$ and $\ell_\infty$-diameter $D_\infty=\max_iD_i$. In round $t=1,2,\dots$ the learner plays $\theta_t\in\Theta$, a convex $f_t\colon\R^d\to\R$ is revealed (it may depend on $\theta_1,\dots,\theta_t$), and the learner receives $g_t\in\partial f_t(\theta_t)$. The \term{regret}\index{regret} against $\theta^\star\in\Theta$ is $R_T(\theta^\star)=\sum_{t=1}^T(f_t(\theta_t)-f_t(\theta^\star))$.
\end{assumption}

If $f_t=\ell(\cdot;\xi_t)$ with independent samples and $\Loss=\E\ell$, then $\E f_t(\theta_t)=\E\Loss(\theta_t)$, and Jensen's inequality turns a regret bound into $\E\Loss(\bar\theta_T)-\Loss(\theta^\star)\le\E R_T(\theta^\star)/T$ for the average $\bar\theta_T$ of the iterates. The second setting is a fixed smooth objective.

\begin{assumption}[Smoothness of unknown size]\label[assumption]{ass:adaptive:smooth}
$\Loss\colon\R^d\to\R$ is bounded below by $\Loss^\star$ and $L$-smooth with respect to a norm $\norm{\cdot}$; the method is not told $L$.
\end{assumption}

\begin{dictionary}
the norm $\norm{\cdot}_{P_t}$, $P_t=\diag(\sqrt{\hat v_t})+\epsilon I$ & Adam's denominator; $\eta/(\sqrt{\hat v_{t,i}}+\epsilon)$, the effective learning rate of parameter $i$\\
the losses $f_t$ and the regret $R_T$ & minibatch losses; the summed training loss minus that of the best fixed parameters in hindsight\\
the box and its diameter $D_\infty$ & the range of a parameter, set by initialization and weight decay\\
the history norm $\norm{g_{1:T,i}}_2$ & $\sqrt T$ times the gradient RMS of parameter $i$, small for rare features\\
the memory $1/(1-\beta_2)$ & the number of steps that $v$ averages; times $B$, in examples\\
the period $n$ of a rare large gradient & a rare token whose gradient is large when it appears\\
a quiet-then-loud coordinate & the embedding of a rare token; a unit that revives\\
an unknown smoothness constant $L$ & the learning rate one would otherwise tune\\
$\epsilon$ & the gradient scale below which Adam stops normalizing
\end{dictionary}

\paragraph{What was thrown away.} In the first setting, convexity and a fixed objective: the bounds hold against an adversary and say nothing about a particular run. The box, which training does not have. Noise, in the theorem of the second setting. And the coordinate system: a diagonal preconditioner adapts to the coordinates in which the parameters are written (\cref{ch:sign,ch:transformers}). What is kept is the ingredient that all the methods share, a norm chosen from the gradients seen so far, and the property that separates them, whether that norm forgets.

\paragraph{What one had to notice.} \Textcite{McMahanStreeter2010} and \textcite{DuchiHazanSinger2011} saw that the regret bound of online gradient descent depends on the norm in which it is run, so that the norm can be chosen online to minimize the bound it will be charged; the minimizer is the inverse root of the accumulated squared gradients (\cref{prop:adaptive:hindsight}). RMSProp and Adam keep the choice of norm and give up its justification: a norm that forgets no longer minimizes a bound on the whole history.

\section{The theory: AdaGrad and the best norm in hindsight}\label{sec:adaptive:adagrad}

\begin{theorem}[Diagonal AdaGrad; \textcite{DuchiHazanSinger2011}, \textcite{McMahanStreeter2010}]\label{thm:adaptive:adagrad}
Under \cref{ass:adaptive:oco}, run AdaGrad with $\epsilon=0$ and step size $\eta>0$ from $\theta_1\in\Theta$, projected on $\Theta$. For every $\theta^\star\in\Theta$ and $T\ge1$,
\begin{equation}\label{eq:adaptive:adagrad}
  R_T(\theta^\star)\le\sum_{i=1}^d\Bigl(\frac{D_i^2}{2\eta}+\eta\Bigr)\norm{g_{1:T,i}}_2\le\Bigl(\frac{D_\infty^2}{2\eta}+\eta\Bigr)\sum_{i=1}^d\norm{g_{1:T,i}}_2 ,
\end{equation}
and with $\eta=D_\infty/\sqrt2$, $R_T(\theta^\star)\le\sqrt2\,D_\infty\sum_i\norm{g_{1:T,i}}_2$.
\end{theorem}

\begin{proof}
The plan is to linearize the losses, which splits the regret into $d$ one-dimensional regrets, and to bound each by the argument for online gradient descent with a step size that changes from round to round.

\emph{Splitting.} By the subgradient inequality $f_t(\theta_t)-f_t(\theta^\star)\le\sum_ig_{t,i}(\theta_{t,i}-\theta^\star_i)$. The projection is coordinatewise clipping, so coordinate $i$ evolves by $\theta_{t+1,i}=\adaptiveclip_{[a_i,b_i]}(\theta_{t,i}-\eta g_{t,i}/s_{t,i})$, through its own gradients only. Hence $R_T(\theta^\star)\le\sum_ir_i$ with $r_i=\sum_{t\le T}g_{t,i}(\theta_{t,i}-\theta^\star_i)$, and it suffices to show $r_i\le(D_i^2/(2\eta)+\eta)\norm{g_{1:T,i}}_2$.

\emph{One coordinate.} Fix $i$ and drop it from the notation. A round with $s_t=0$ has $g_t=0$: it contributes nothing and does not move the iterate. If all $g_t$ vanish, $r=0$; otherwise let $\tau$ be the first round with $g_\tau\neq0$, so that $\eta_t=\eta/s_t$ is finite and nonincreasing for $t\ge\tau$. Clipping to an interval containing $\theta^\star$ does not increase the distance to it, so with $e_t=(\theta_t-\theta^\star)^2$ we have $e_{t+1}\le(\theta_t-\eta_tg_t-\theta^\star)^2=e_t-2\eta_tg_t(\theta_t-\theta^\star)+\eta_t^2g_t^2$, that is,
\[
  g_t(\theta_t-\theta^\star)\le\frac{e_t-e_{t+1}}{2\eta_t}+\frac{\eta_t}2g_t^2 .
\]
We sum over $\tau\le t\le T$ and bound the two sums in turn. Regrouping the first by the $e_t$, dropping $-e_{T+1}/(2\eta_T)\le0$, and using $e_t\le D^2$ and $1/\eta_t-1/\eta_{t-1}\ge0$,
\[
  \sum_{t=\tau}^T\frac{e_t-e_{t+1}}{2\eta_t}\le\frac{e_\tau}{2\eta_\tau}+\sum_{t=\tau+1}^Te_t\Bigl(\frac1{2\eta_t}-\frac1{2\eta_{t-1}}\Bigr)\le\frac{D^2}{2\eta_T}=\frac{D^2s_T}{2\eta}.
\]
For the second we use the AdaGrad lemma (\cref{lem:toolkit:adagrad}): for $a_t\ge0$ and $S_t=\sum_{u\le t}a_u$, $a_t/\sqrt{S_t}\le2(\sqrt{S_t}-\sqrt{S_{t-1}})$, because $\sqrt{S_t}-\sqrt{S_{t-1}}=a_t/(\sqrt{S_t}+\sqrt{S_{t-1}})$; the sum telescopes. With $a_t=g_t^2$ it gives $\sum_t\frac{\eta_t}2g_t^2=\frac\eta2\sum_tg_t^2/s_t\le\eta s_T$. Adding, $r\le(D^2/(2\eta)+\eta)s_T$ with $s_T=\norm{g_{1:T}}_2$. Finally $D_i\le D_\infty$, and $\eta=D_\infty/\sqrt2$ minimizes $D_\infty^2/(2\eta)+\eta$.
\end{proof}

Besides its definition, the proof used one property of the step sizes: they never increase. Without it the coefficients $1/\eta_t-1/\eta_{t-1}$ can be negative and the first sum no longer telescopes; this is the property that RMSProp and Adam lack, and where \cref{prop:adaptive:reddi} breaks them. Run with one scalar $s_t=(\sum_{u\le t}\norm{g_u}_2^2)^{1/2}$, Euclidean projections and no splitting, the same proof bounds online gradient descent on a convex set of Euclidean diameter $D_2$:
\begin{equation}\label{eq:adaptive:ogd}
  R_T(\theta^\star)\le\Bigl(\frac{D_2^2}{2\eta}+\eta\Bigr)\Bigl(\sum_{t=1}^T\norm{g_t}_2^2\Bigr)^{1/2},
\end{equation}
within $\sqrt2$, at $\eta=D_2/\sqrt2$, of the bound $D_2(\sum_t\norm{g_t}_2^2)^{1/2}$ for the best constant step chosen in hindsight (\cref{thm:toolkit:ogd}).\sidenote{This scalar step $D_2/(\sqrt2\,s_t)$ is AdaGrad-Norm. Unlike the classical tuned step $D/(G\sqrt t)$, it needs no bound $G$ on the gradients.} On a cube of side $D_\infty$, $D_2=\sqrt d\,D_\infty$ and $\sum_i\norm{g_{1:T,i}}_2\le\sqrt d\,(\sum_t\norm{g_t}_2^2)^{1/2}$, so AdaGrad's bound is never larger than \eqref{eq:adaptive:ogd}, and can be smaller by a factor $\sqrt d$.

\begin{proposition}[The best diagonal norm in hindsight]\label[proposition]{prop:adaptive:hindsight}
Let $g_1,\dots,g_T\in\R^d$, $G_i=\norm{g_{1:T,i}}_2$ not all zero, and $0/0=0$, $x/0=\infty$ for $x>0$. Then
\[
  \min\Bigl\{\sum_{t=1}^T\norm{g_t}^2_{\diag(p)^{-1}}\colon p\in\R^d_{\ge0},\ \sum_ip_i\le1\Bigr\}=\Bigl(\sum_{i=1}^dG_i\Bigr)^2 ,
\]
attained at $p_i=G_i/\sum_jG_j$.
Consequently online gradient descent in the best fixed diagonal norm with the best constant step, both chosen in hindsight, has the regret bound $D_\infty\sum_iG_i$ on a box of $\ell_\infty$-diameter $D_\infty$, and AdaGrad's bound is $\sqrt2$ times this.
\end{proposition}

\begin{proof}
The objective is $\sum_iG_i^2/p_i$. If it is finite, $p_i>0$ wherever $G_i>0$, and by Cauchy--Schwarz $\sum_iG_i=\sum_{G_i>0}(G_i/\sqrt{p_i})\sqrt{p_i}\le(\sum_iG_i^2/p_i)^{1/2}(\sum_ip_i)^{1/2}$, which is at most $(\sum_iG_i^2/p_i)^{1/2}$. At $p_i=G_i/\sum_jG_j$ the objective equals $(\sum_iG_i)^2$. For the consequence, the proof of \cref{thm:adaptive:adagrad} in the norm $\norm{\cdot}_P$, $P=\diag(p)$, with a constant step gives the regret bound $D_P^2/(2\eta)+\frac\eta2\sum_t\norm{g_t}^2_{P^{-1}}$ with $D_P=\sup_{\Theta}\norm{\theta-\theta'}_P\le D_\infty(\sum_ip_i)^{1/2}$; the best $\eta$ gives $D_P(\sum_t\norm{g_t}^2_{P^{-1}})^{1/2}$, which is invariant under $p\mapsto cp$, and minimizing over $p$ gives $D_\infty\sum_iG_i$.
\end{proof}

AdaGrad pays a factor $\sqrt2$ for not knowing the future. The choice of norm matters when the coordinates are used at very different rates.

\begin{corollary}[Sparse features]\label[corollary]{cor:adaptive:sparse}
In \cref{thm:adaptive:adagrad} on a cube of side $D_\infty$, suppose $\abs{g_{t,i}}\le1$, and $g_{t,i}=0$ unless feature $i$ is active in round $t$, which happens with probability $p_i$. With $\eta=D_\infty/\sqrt2$ for AdaGrad and $\eta=D_2/\sqrt2$ in \eqref{eq:adaptive:ogd}, for every $\theta^\star\in\Theta$,
\begin{gather*}
  \E R_T^{\mathrm{AdaGrad}}(\theta^\star)\le\sqrt{2T}\,D_\infty\sum_{i=1}^d\sqrt{p_i},\\
  \E R_T^{\mathrm{OGD}}(\theta^\star)\le\sqrt{2T}\,D_\infty\Bigl(d\sum_{i=1}^dp_i\Bigr)^{1/2}.
\end{gather*}
For $p_i=i^{-\alpha}$ the ratio of the right sides tends to a constant if $\alpha<1$, grows like $(1-\frac\alpha2)\zeta(\alpha)^{1/2}d^{(\alpha-1)/2}$ if $1<\alpha<2$, and like $(\pi/\sqrt6)\sqrt d/\ln d$ if $\alpha=2$.
\end{corollary}

\begin{proof}
Pathwise $\sum_tg_{t,i}^2$ is at most the number $N_i$ of rounds with feature $i$ active, and $\E N_i=p_iT$. By Jensen's inequality $\E\norm{g_{1:T,i}}_2\le(p_iT)^{1/2}$ and $\E(\sum_t\norm{g_t}_2^2)^{1/2}\le(T\sum_ip_i)^{1/2}$; insert these in \eqref{eq:adaptive:adagrad} and \eqref{eq:adaptive:ogd} with $D_2=\sqrt d\,D_\infty$. The asymptotics follow from $\sum_{i\le d}i^{-s}\sim d^{1-s}/(1-s)$ for $s<1$, $\sum_{i\le d}i^{-1}\sim\ln d$, and $\sum_{i\le d}i^{-s}\to\zeta(s)$ for $s>1$, with $\zeta(2)=\pi^2/6$.
\end{proof}

The scalar method has a second property, unrelated to regret: it does not need the smoothness constant. \term{AdaGrad-Norm}\index{AdaGrad!AdaGrad-Norm} in a norm $\norm{\cdot}$, started at $\theta_0$ with $b_0>0$, iterates for $t=0,1,\dots$
\begin{equation}\label{eq:adaptive:adagradnorm}
  b_{t+1}^2=b_t^2+\norm{\nabla\Loss(\theta_t)}_*^2,\qquad \theta_{t+1}=\theta_t-\frac{\eta}{b_{t+1}}\nabla\Loss(\theta_t)^\sharp ;
\end{equation}
in the Euclidean norm this is the method of \textcite{WardWuBottou2019}. Since $b_{t+1}\ge\norm{\nabla\Loss(\theta_t)}_*$, every step has length at most $\eta$.

\begin{theorem}[AdaGrad-Norm adapts to the smoothness; after \textcite{WardWuBottou2019}]\label{thm:adaptive:adagradnorm}
Under \cref{ass:adaptive:smooth}, let $\eta,b_0>0$, $\Delta_0=\Loss(\theta_0)-\Loss^\star$, $\Lambda=\max\{0,\ln(\eta L/b_0)\}$ and $Q=2\Delta_0/\eta+2\eta L\Lambda$. Then \eqref{eq:adaptive:adagradnorm} satisfies (i) $\Loss(\theta_t)\le\Loss(\theta_0)+\eta^2L\Lambda$ for every $t$, and (ii) for every $T\ge1$
\[
  \sum_{t=0}^{T-1}\norm{\nabla\Loss(\theta_t)}_*^2\le\bigl(Q+\max\{b_0,\eta L\}\bigr)^2,
\]
so that $\min_{t<T}\norm{\nabla\Loss(\theta_t)}_*^2\le(Q+\max\{b_0,\eta L\})^2/T$.
\end{theorem}

\begin{proof}
The plan is to split the run at the first step at which $b_{t+1}$ exceeds $\eta L$: before it the loss can rise, but only as much as $b$ grows; after it every step descends; and together the two phases keep $b_t$ bounded. Write $a_t=\norm{\nabla\Loss(\theta_t)}_*^2=b_{t+1}^2-b_t^2$. The sharp operator has $\ip{\nabla\Loss(\theta_t),\nabla\Loss(\theta_t)^\sharp}=a_t$ and $\norm{\nabla\Loss(\theta_t)^\sharp}=a_t^{1/2}$, so the descent lemma (\cref{lem:toolkit:descent}) gives
\begin{equation}\label{eq:adaptive:onestep}
  \Loss(\theta_{t+1})\le\Loss(\theta_t)-\frac{\eta}{b_{t+1}}\Bigl(1-\frac{\eta L}{2b_{t+1}}\Bigr)a_t .
\end{equation}
(A zero gradient causes no difficulty: the iterate stays put.) Let $t_0$ be the first $t$ with $b_{t+1}\ge\eta L$, or $\infty$ if there is none.

\emph{Before $t_0$.} For $t<t_0$, \eqref{eq:adaptive:onestep} gives $\Loss(\theta_{t+1})-\Loss(\theta_t)\le\frac{\eta^2L}2a_t/b_{t+1}^2$, and $a_t/b_{t+1}^2=1-b_t^2/b_{t+1}^2\le\ln(b_{t+1}^2/b_t^2)$ because $1-1/x\le\ln x$. The sum over $t<s$ telescopes to $\ln(b_s^2/b_0^2)\le2\Lambda$ for $1\le s\le t_0$, since $b_s<\eta L$ there. Hence $\Loss(\theta_s)\le\Loss(\theta_0)+\eta^2L\Lambda$ and $\sum_{t<s}a_t=b_s^2-b_0^2\le\max\{0,\eta^2L^2-b_0^2\}$ for $s\le t_0$.

\emph{From $t_0$ on.} For $t\ge t_0$, $b_{t+1}\ge\eta L$, the bracket in \eqref{eq:adaptive:onestep} is at least $\frac12$, and the loss decreases; this proves (i). Summing \eqref{eq:adaptive:onestep} over $t_0\le t<T$ and using $\Loss(\theta_T)\ge\Loss^\star$ gives $\frac\eta2\sum_{t=t_0}^{T-1}a_t/b_{t+1}\le\Delta_0+\eta^2L\Lambda$, so $\sum_{t=t_0}^{T-1}a_t\le b_TQ$ because $b_{t+1}\le b_T$.

\emph{Both phases.} Adding, $b_T^2-b_0^2=\sum_{t<T}a_t\le\max\{0,\eta^2L^2-b_0^2\}+Qb_T$, so $b_T^2\le c+Qb_T$ with $c=\max\{b_0^2,\eta^2L^2\}$. The positive root of $x^2-Qx-c$ is at most $Q+\sqrt c$, so $\sum_{t<T}a_t\le b_T^2\le(Q+\sqrt c)^2$.
\end{proof}

Gradient descent with a fixed step needs $\eta<2/L$ to converge at all, and attains $\min_t\norm{\nabla\Loss(\theta_t)}_*^2\le2L\Delta_0/T$ at $\eta=1/L$ (\cref{ch:steepest}). AdaGrad-Norm converges at the rate $1/T$ for every $\eta$ and $b_0$; too large a step costs a rise of the loss by at most $\eta^2L\ln(\eta L/b_0)$ and a constant of order $(\eta L)^2\ln^2(\eta L/b_0)$. In one dimension $b_t$ grows until $\eta/b_t$ is about $1/L$ and stops:\sidenote{On $\Loss(\theta)=\frac12\theta^2$ from $\theta_0=1$ with $\eta=100$ and $b_0=10^{-3}$, $b_t$ settles at $99.0$, so that the final step size is $1.01/L$.} the method finds a stable step size, and loses this ability as soon as $b$ forgets (\cref{ex:adaptive:rmsnorm}). With stochastic gradients the step $\eta/b_{t+1}$ depends on the gradient it multiplies; \textcite{WardWuBottou2019} handle this and obtain a rate of order $\ln T/\sqrt T$ under bounded gradients and noise.

\section{The theory: what forgetting costs}\label{sec:adaptive:forgetting}

A moving average with $\beta_2<1$ forgets, and the forgetting has a price in every analysis. It shows first in the AdaGrad lemma.

\begin{lemma}[The AdaGrad lemma with forgetting]\label[lemma]{lem:adaptive:forgetting}
Let $a_t\ge0$, $0<\beta\le1$, $\epsilon>0$, $v_0\ge0$ and $v_t=\beta v_{t-1}+(1-\beta)a_t$. Then $(1-\beta)\sum_{t=1}^Ta_t/(\epsilon+v_t)\le\ln\frac{\epsilon+v_T}{\epsilon+v_0}+T\ln\frac1\beta$.
\end{lemma}

\begin{proof}
Since $\epsilon+v_t\ge\beta(\epsilon+v_{t-1})+(1-\beta)a_t$ and $1-x\le\ln(1/x)$ for $x>0$, we get $(1-\beta)a_t/(\epsilon+v_t)\le1-\beta(\epsilon+v_{t-1})/(\epsilon+v_t)\le\ln\frac{\epsilon+v_t}{\epsilon+v_{t-1}}+\ln\frac1\beta$, and the sum telescopes.
\end{proof}

For constant $a_t$ both sides grow like $T(1-\beta)$ (\cref{ex:adaptive:lemma}), so the term $T\ln(1/\beta)$ is genuine, and negligible only if $1-\beta_2\lesssim1/T$: an analysis that sums Adam's steps over a run pays for forgetting unless the memory is comparable to the run. That a short memory can fail outright is the next proposition, the counterexample of \textcite{ReddiKaleKumar2018} redone for an arbitrary period: a rare large gradient points one way, frequent small ones the other.\sidenote{Reddi, Kale and Kumar use period $3$ with $\beta_2=1/(1+C^2)$ for their first theorem, and a long period for the theorem that covers every $\beta_2$; a general period makes the role of the memory explicit.}

\begin{proposition}[Reddi's counterexample, period $n$]\label[proposition]{prop:adaptive:reddi}
Let $n\ge3$, $C>n-1$, $\Theta=[-1,1]$, and $f_t(\theta)=C\theta$ if $t\equiv1\pmod n$, $f_t(\theta)=-\theta$ otherwise; every partial sum $\sum_{t\le T}f_t$ is minimized at $-1$. Run Adam with $\beta_1=0$, $\epsilon=0$, $0\le\beta_2<1$, with or without bias correction, projected on $\Theta$, from $\theta_1\in\Theta$, with nonincreasing steps $\eta_t>0$ such that $\eta_t\to0$, $\eta_{t+1}/\eta_t\to1$ and $\sum_t\eta_t=\infty$ (for example $\eta/\sqrt t$). Let $q=(1-\beta_2)/(1-\beta_2^n)$ and
\begin{equation}\label{eq:adaptive:kappa}
  \kappa_n(\beta_2,C)=-\frac{C}{\sqrt{1+q(C^2-1)}}+\sum_{k=1}^{n-1}\frac1{\sqrt{1+q\beta_2^k(C^2-1)}} .
\end{equation}
\begin{enumerate}
\item If $\kappa_n(\beta_2,C)>0$, then $\theta_t\to1$ and $R_T/T\to2(C-n+1)/n$; if $\kappa_n(\beta_2,C)<0$, then $\theta_t\to-1$ and $R_T/T\to0$.
\item $\kappa_n(0,C)=n-2>0$; and as $\beta_2\to1$, $\kappa_n(\beta_2,C)$ tends to $(n-1-C)\sqrt{n/(C^2+n-1)}$, which is negative.
\item For every $\beta_2\in[0,1)$ there are $n$ and $C>n-1$ with $\kappa_n(\beta_2,C)>0$.
\end{enumerate}
\end{proposition}

\begin{proof}
The plan is to show that $v_t$ becomes periodic, so that each period moves the iterate by about $\eta_t\kappa_n$, and that a drift of fixed sign carries the iterate to the corresponding end despite the clipping.

\emph{Periodicity.} For the lag $k\in\{0,\dots,n-1\}$ since the last large gradient let $w_k=1+q\beta_2^k(C^2-1)$. These numbers solve the recursion of $v$: $\beta_2w_{k-1}+1-\beta_2=w_k$ for $k\ge1$, and $\beta_2w_{n-1}+(1-\beta_2)C^2=w_0$ because $q\beta_2^n+1-\beta_2=q$. The difference between $v_t$ and this periodic solution is multiplied by $\beta_2$ at each step, so $v_t$ and $\hat v_t$ approach $w_{k(t)}$, $k(t)$ the lag of round $t$. Since $\abs{g_t}\ge1$, $\hat v_t\ge(1-\beta_2)g_t^2$, so $\Delta_t=g_t/\sqrt{\hat v_t}$ satisfies $\abs{\Delta_t}\le(1-\beta_2)^{-1/2}$ and $\Delta_t-g_t/\sqrt{w_{k(t)}}\to0$. Over period $j$ (rounds $(j-1)n<t\le jn$) the unclipped displacement $U_j=-\sum_t\eta_t\Delta_t$ and the variation $M_j=\sum_t\eta_t\abs{\Delta_t}$ therefore satisfy $M_j\to0$ and $U_j/\eta_{jn}\to\kappa_n$, because $\eta_t/\eta_{jn}\to1$ uniformly over the period.

\emph{Clipping.} Let $\kappa_n>0$ and $j_0$ be such that $U_j\ge\kappa_n\eta_{jn}/2$ and $M_j\le\frac14$ for $j\ge j_0$. A period that starts at $\theta\le1-M_j$ never reaches the upper end, and clipping at $-1$ only raises the iterate, so it ends at least at $\theta+\kappa_n\eta_{jn}/2$. A period that starts above $1-M_j$ keeps the iterate above $1-2M_j$, since clipping at $1$ lowers nothing below $1$ minus the movement that follows. Since $\sum_j\eta_{jn}\ge\frac1n\sum_{t>n}\eta_t=\infty$ and $\theta_t\le1$, periods of the second kind occur infinitely often; after one of them, in period $j_1$, every period starts at least at $1-2\sup_{j\ge j_1}M_j$ (a period of the first kind does not lower the start of the next), and every iterate is at least $1-3\sup_{j\ge j_1}M_j$. Hence $\theta_t\to1$. For $\kappa_n<0$, apply the argument to $-\theta_t$.

\emph{Regret and limits.} Every partial sum of the $g_t$ is positive, so $\theta^\star=-1$, $R_T=\sum_{t\le T}g_t(\theta_t+1)$, and (i) follows from $\sum_{t\le T}g_t=(C-n+1)T/n+O(1)$. For (ii): if $\beta_2=0$, then $q=1$, $w_0=C^2$, $w_k=1$ for $k\ge1$, and $\kappa_n=-1+(n-1)$; as $\beta_2\to1$, $q\to1/n$ and $w_k\to1+(C^2-1)/n$. For (iii) let $\beta_2>0$ and $C=n$. Since $1-\beta_2\le q\le1$, the first term of \eqref{eq:adaptive:kappa} is at least $-1/\sqrt q\ge-(1-\beta_2)^{-1/2}$, and the terms with $\beta_2^kn^2\le1$, which are all $k\ge K$ for some $K$ of order $\ln n$, are at least $2^{-1/2}$. So $\kappa_n(\beta_2,n)\ge(n-K)/\sqrt2-(1-\beta_2)^{-1/2}>0$ for large $n$.
\end{proof}

With $\beta_2=0$ Adam is sign descent (\cref{ch:sign}), which moves by $\eta_t$ whatever the size of the gradient: it counts votes, and the small gradients outvote the large one.\sidenote{RMSProp is Adam with $\beta_1=0$ and no bias correction, so the proposition covers it too.} With a memory much longer than the period, $\hat v_t$ is nearly constant and Adam behaves like AdaGrad. In between, the large gradient inflates $\hat v$ right after it, so its own step is moderate, but by the end of the period $\hat v$ has forgotten it and the small gradients take nearly full steps; for the choice $n=3$, $\beta_2=1/(1+C^2)$ of \textcite{ReddiKaleKumar2018}, $\kappa_3\to2^{-1/2}$ (\cref{ex:adaptive:reddichoice}). Their fix, \term{AMSGrad}\index{AMSGrad}, replaces $v_t$ by its running maximum, which restores the monotonicity that \cref{thm:adaptive:adagrad} used; here it converges to $-1$.

The published convergence theory of Adam lives on the convergent side of this boundary, with a memory tied to the problem or the horizon. \Textcite{DefossezEtAl2022} assume $L$-smoothness and stochastic gradients bounded in $\ell_\infty$ almost surely; with $1-\beta_2=1/T$ and steps of order $1/\sqrt T$ they bound the average of $\E\norm{\nabla\Loss(\theta_t)}_2^2$ over the run by a constant times $d\ln T/\sqrt T$, with a factor of order $1/(1-\beta_1)$ for momentum; for fixed $\beta_2$ their bound does not vanish. \Textcite{ZhangEtAl2022} treat finite sums with random reshuffling and unbounded gradients: if $\beta_1<\sqrt{\beta_2}$ and $\beta_2$ exceeds a problem-dependent threshold, Adam reaches a neighborhood of stationary points of size proportional to the gradient noise, and for a range of small $\beta_2$ it diverges. \Textcite{LiRakhlinJadbabaie2023} reach $\eps$-stationary points with $O(\eps^{-4})$ stochastic gradients when the curvature may grow with the gradient (\cref{ch:spikes}), again with $\beta_2$ near one. \Cref{lem:adaptive:forgetting} shows where the requirement enters, and \cref{prop:adaptive:reddi} that it cannot be dropped.

Forgetting has a second, opposite cost. A long memory makes the preconditioner slow, and a coordinate whose gradient was small for a long time and then becomes large is divided by an out-of-date $v$. The next proposition is the exact form of this \term{quiet-then-loud}\index{quiet-then-loud} event.

\begin{proposition}[Quiet then loud]\label[proposition]{prop:adaptive:quietloud}
Consider one coordinate of Adam without bias correction, $0\le\beta_1<1$, $0<\beta_2<1$, $\epsilon\ge0$. Let $\abs{g_s}\le\delta$ for $s<t_0$ and $g_s=G$ for $t_0\le s\le t_0+k$, where $G>0$, $\delta\ge0$ and $\beta_1\delta\le(1-\beta_1)G$, and let $a=\beta_1^{k+1}$, $b=\beta_2^{k+1}$. The step $u_k=\eta\,m_{t_0+k}/(\sqrt{v_{t_0+k}}+\epsilon)$ satisfies
\begin{equation}\label{eq:adaptive:quiet}
  \frac{(1-a)G-a\delta}{\sqrt{b\delta^2+(1-b)G^2}+\epsilon}\le\frac{u_k}\eta\le\frac{(1-a)G+a\delta}{\sqrt{1-b}\,G+\epsilon} .
\end{equation}
\begin{enumerate}
\item As $\delta/G\to0$, $u_k/\eta\to\rho_k=(1-\beta_1^{k+1})/(\sqrt{1-\beta_2^{k+1}}+\epsilon/G)$. With $\epsilon=0$ the first loud step is $\rho_0=(1-\beta_1)/\sqrt{1-\beta_2}$, and $\rho_k\to1$.
\item With $\epsilon=0$ the excess of the loud steps over sign steps, $E(\beta_1,\beta_2)=\sum_{k\ge0}(\rho_k-1)$, is finite, and $E=2\ln2/(1-\beta_2)+O((1-\beta_2)^{-1/2})$ as $\beta_2\to1$, for fixed $\beta_1$.
\end{enumerate}
\end{proposition}

\begin{proof}
The plan is to unroll the two averages over the loud stretch and bound what they inherited. Since $m$ and $v$ start at $0$, they are averages with weights of total at most one, so $\abs{m_{t_0-1}}\le\delta$ and $0\le v_{t_0-1}\le\delta^2$. Over the loud rounds $m_{t_0+k}=a\,m_{t_0-1}+(1-a)G$ and $v_{t_0+k}=b\,v_{t_0-1}+(1-b)G^2$, so $(1-a)G-a\delta\le m_{t_0+k}\le(1-a)G+a\delta$, a nonnegative lower bound because $a\delta\le\beta_1^k(1-\beta_1)G\le(1-a)G$, and $(1-b)G^2\le v_{t_0+k}\le b\delta^2+(1-b)G^2$. Dividing gives \eqref{eq:adaptive:quiet}, and (i) follows. For (ii) let $\nu=\ln(1/\beta_2)$ and $f(x)=(1-e^{-x})^{-1/2}-1$, positive and decreasing, so that \mbox{$\rho_k-1=f((k+1)\nu)-\beta_1^{k+1}(1-\beta_2^{k+1})^{-1/2}$}. Comparing a sum of a decreasing function with its integral, $\int_\nu^\infty f\le\nu\sum_{j\ge1}f(j\nu)\le\int_0^\infty f$, and $\int_0^\infty f=\int_0^12\dd w/(1+w)=2\ln2$ by the substitution $1-e^{-x}=w^2$. Since $1-e^{-x}\ge\min\{x,1\}/2$, $\int_0^\nu f\le2\sqrt{2\nu}$, and the second part is at most $\sum_j\beta_1^j\max\{(2/(j\nu))^{1/2},\sqrt2\}=O(\nu^{-1/2})$. Finally $1/\nu=1/(1-\beta_2)+O(1)$.
\end{proof}

With $\beta_1=0.9$ the first loud step is $3.16\eta$ for $\beta_2=0.999$, $\eta$ for $0.99$ and $0.45\eta$ for $0.95$; the steps then grow, because the momentum fills up faster than $v$, and peak at $6.57\eta$, $2.13\eta$ and $1.10\eta$ eleven to nineteen rounds after the jump. Each $\rho_k$ is a limit of ratios $\abs{m_t}/\sqrt{v_t}$, so the peaks respect the bound $(1-\beta_1)/\sqrt{(1-\beta_2)(1-\beta_1^2/\beta_2)}$ of \cref{ch:sign} (\cref{prop:sign:adambound}): $7.27$, $2.35$ and $1.17$. The excess $E$ is the integrated lag of the preconditioner, $1213$ step sizes for $\beta_2=0.999$ and fewer than $2$ for $0.95$. Bias correction changes nothing once $t_0(1-\beta_2)$ is large (\cref{ex:adaptive:bias}).

\begin{remark}[What $\epsilon$ does]\label[remark]{rem:adaptive:eps}
The step $\eta\hat m_{t,i}/(\sqrt{\hat v_{t,i}}+\epsilon)$ is Adam's if $\sqrt{\hat v_{t,i}}\gg\epsilon$ and momentum SGD's with step size $\eta/\epsilon$ if $\sqrt{\hat v_{t,i}}\ll\epsilon$: $\epsilon$ switches from sign-like to gradient-like behavior at the gradient scale $\epsilon$, and caps every step at $\eta\abs{\hat m_{t,i}}/\epsilon$. In \cref{prop:adaptive:quietloud} it matters once it is comparable with $\sqrt{1-\beta_2}\,G$; on a quadratic it turns the two-cycle of sign descent into convergence once it exceeds half the step size times the curvature (\cref{ex:adaptive:eps}). It is set small so that every gradient that carries signal is normalized, and it matters when gradients shrink to its size \parencite{WortsmanEtAl2024}; \textcite{EverettEtAl2024} remove it by replacing $m/(\sqrt v+\epsilon)$ with a scaled $\arctan$ of $m/\sqrt v$.
\end{remark}

\begin{remark}[Adam's two limits]\label[remark]{rem:adaptive:limits}
With $\beta_1=\beta_2=\epsilon=0$ Adam is sign descent. At the other end, let the gradients at a fixed $\theta$ be independent with $\E g_{t,i}^2=\sigma_i^2>0$ and finite fourth moments. Then $\E\hat v_{t,i}=\sigma_i^2$ and $\Var v_{t,i}\le\frac{1-\beta_2}{1+\beta_2}\Var(g_{t,i}^2)$, so $\hat v_{t,i}\to\sigma_i^2$ in probability as $\beta_2\to1$ with $t(1-\beta_2)\to\infty$, and Adam becomes momentum SGD preconditioned by $\diag(\E g^2)^{-1/2}$. In the noisy quadratic of \cref{ch:noise} (diagonal $H$, gradient $H\theta+\zeta$, $\Cov\zeta=\sigma^2H/B$), $\E g_i^2=h_i^2\theta_i^2+\sigma^2h_i/B$: where noise dominates the preconditioner is proportional to $H^{1/2}$, where signal dominates it is $\abs{\nabla\Loss}$, the sign-descent limit. It is the root of the second moment of the minibatch gradient, neither the Hessian nor the Fisher information \parencite{KunstnerHennigBalles2019}.
\end{remark}

Adafactor stores, for each weight matrix, moving averages of the row and column sums of $G_t^2$ and multiplies them, and this solves a matrix approximation problem exactly. For entrywise nonnegative $V$ and $R$ of the same size the \term{I-divergence}\index{I-divergence} (generalized Kullback--Leibler divergence) is $D(V\Vert R)=\sum_{ij}(V_{ij}\ln(V_{ij}/R_{ij})-V_{ij}+R_{ij})$, with $0\ln0=0$, and $D=\infty$ if some $V_{ij}>0=R_{ij}$.

\begin{proposition}[Adafactor's factorization; \textcite{ShazeerStern2018}]\label[proposition]{prop:adaptive:adafactor}
Let $V\in\R^{m\times n}$ be entrywise nonnegative and nonzero, $S=\mathbf1^\top V\mathbf1$. Among nonnegative rank-one matrices $xy^\top$, $D(V\Vert xy^\top)$ is minimized exactly at $R^\star=(V\mathbf1)(\mathbf1^\top V)/S$, the row sums times the column sums over the total. The minimizer has the row and column sums of $V$, and equals $V$ if $V$ has rank one. Consequently Adafactor's estimate $r_tc_t^\top/(\mathbf1^\top r_t)$ is $R^\star$ for Adam's matrix $V_t$ of moving averages of $G_t^2$.
\end{proposition}

\begin{proof}
The plan is to write the excess divergence over $R^\star$ as a sum of three nonnegative terms. Let $r=V\mathbf1$, $c=V^\top\mathbf1$, $X=\sum_ix_i$, $Y=\sum_jy_j$, and assume $D(V\Vert xy^\top)<\infty$, so $x_iy_j>0$ wherever $V_{ij}>0$. Expanding, $D(V\Vert xy^\top)=\sum_{ij}V_{ij}\ln V_{ij}-\sum_ir_i\ln x_i-\sum_jc_j\ln y_j-S+XY$, and $D(V\Vert R^\star)$ is the case $x=r$, $y=c/S$. Subtracting, and writing $\sum_ir_i\ln(r_i/x_i)=S\adaptiveKL(r/S\Vert x/X)+S\ln(S/X)$ and likewise for $c$ and $y$, we find with $\phi(z)=z-1-\ln z\ge0$
\[
  D(V\Vert xy^\top)-D(V\Vert R^\star)=S\,\adaptiveKL\Bigl(\frac rS\Big\Vert\frac xX\Bigr)+S\,\adaptiveKL\Bigl(\frac cS\Big\Vert\frac yY\Bigr)+S\,\phi\Bigl(\frac{XY}S\Bigr).
\]
All three terms vanish exactly when $x/X=r/S$, $y/Y=c/S$ and $XY=S$, that is, when $xy^\top=rc^\top/S$. The row sums of $R^\star$ are $r(\mathbf1^\top c)/S=r$ and its column sums $c$; if $V=xy^\top$ then $D(V\Vert V)=0$ forces $V=R^\star$. Row and column sums commute with moving averages, so $r_t$, $c_t$ and $\mathbf1^\top r_t$ are those of $V_t$.
\end{proof}

The optimum has $D(V\Vert R^\star)=S$ times the mutual information of the row and column indices under $V/S$ (\cref{ex:adaptive:adafactor}): the factored second moment is exact when the squared gradient of an entry is a row effect times a column effect, and its error measures their dependence.

\section{Reading the theory}\label{sec:adaptive:reading}

\paragraph{Where the norm matters.} \Cref{fig:adaptive:regret} runs both learners of \cref{cor:adaptive:sparse} on online logistic regression with binary features of frequencies $p_i=i^{-\alpha}$. The corollary predicts an advantage for AdaGrad growing like $d^{1/4}$ for $\alpha=1.5$ and like $\sqrt d/\ln d$ for $\alpha=2$: the scalar step must be small enough for the frequent features and is then too small for the rare ones. The bounds are loose, by a factor of about $12$ for both learners at $d=1000$, so they need not predict the ratio of the regrets; yet the measured ratio follows the ratio of the bounds within $10\%$, over the whole range for $\alpha=1.5$ and up to $d=2000$ for $\alpha=2$. Beyond, for $\alpha=2$, it falls below ($10.0$ against $13.1$ at $d=10^4$): features rarer than $1/T$ appear once or never, and the bound charges each of them $\sqrt{p_iT}$, more than the regret they cause.

\begin{figure}[t]
\begin{fullwidth}
\centering
\begin{tikzpicture}
\begin{groupplot}[group style={group size=2 by 1,horizontal sep=17mm},wplot,width=82mm,height=46mm,xmode=log,xmin=8,xmax=12500]
\nextgroupplot[xlabel={dimension $d$},ylabel={regret, $T=10^4$},ymode=log,ymin=50,ymax=80000,legend style={at={(0.03,0.97)},anchor=north west,font=\scriptsize},legend columns=2]
\addplot[theory2] table[x=d,y=bndO20] {figures/data/06-adaptive-regret-bounds.dat};
\addplot[theory] table[x=d,y=bndA20] {figures/data/06-adaptive-regret-bounds.dat};
\addplot[simulation2] table[x=d,y=regO20] {figures/data/06-adaptive-regret-sim.dat};
\addplot[simulation] table[x=d,y=regA20] {figures/data/06-adaptive-regret-sim.dat};
\legend{OGD bound,AdaGrad bound,OGD,AdaGrad}
\nextgroupplot[xlabel={dimension $d$},ylabel={regret of OGD / AdaGrad},ymin=0,ymax=14,legend style={at={(0.03,0.97)},anchor=north west,font=\scriptsize}]
\addplot[theory] table[x=d,y=bratio20] {figures/data/06-adaptive-regret-bounds.dat};
\addplot[theory3] table[x=d,y=bratio15] {figures/data/06-adaptive-regret-bounds.dat};
\addplot[simulation] table[x=d,y=ratio20] {figures/data/06-adaptive-regret-sim.dat};
\addplot[simulation2] table[x=d,y=ratio15] {figures/data/06-adaptive-regret-sim.dat};
\legend{{bounds, $\alpha=2$},{bounds, $\alpha=1.5$},{measured, $\alpha=2$},{measured, $\alpha=1.5$}}
\end{groupplot}
\end{tikzpicture}
\end{fullwidth}
\caption{AdaGrad against online gradient descent (OGD) with the scalar adaptive step on online logistic regression over $[-1,1]^d$: feature $i$ is present with probability $i^{-\alpha}$, labels come from a logistic teacher with weights $\pm1$, and regret is against the best fixed weights in hindsight; $T=10^4$, means over $3$ runs (seeds in \texttt{06-adaptive-regret.py}). Left ($\alpha=2$): regrets (points) and the bounds of \cref{cor:adaptive:sparse} (curves). Right: ratio of the regrets against ratio of the bounds.}
\label{fig:adaptive:regret}
\end{figure}

\paragraph{Memory against period.} \Cref{fig:adaptive:reddi} draws the boundary $\kappa_n=0$ of \cref{prop:adaptive:reddi} for a rare gradient twice as large as needed to win, $C=2(n-1)$, with Adam runs on a grid of periods and memories; every run ends where the proposition says. The critical memory is $1.46$ at $n=3$ ($\beta_2=0.32$), $40.1$ at $n=100$ and $398$ at $n=1000$, approaching $0.398\,n$ (\cref{ex:adaptive:period}). Read with the practice box: a memory of $20$ steps ($\beta_2=0.95$) is outvoted by such a pattern from a period of about $50$ steps on, and a memory of $1000$ steps from about $2500$. Momentum moves the boundary little ($34$ instead of $40$ at $n=100$ for $\beta_1=0.9$, \cref{ex:adaptive:momentum}); random arrivals move it up, to a memory of about one period (\cref{ex:adaptive:random}). Near the curve the drift is so small that a finite run reaches neither end, and there the comparison must fail.

\begin{figure}[t]
\centering
\begin{tikzpicture}
\begin{axis}[wplot,height=0.5\linewidth,xmode=log,ymode=log,xmin=2.5,xmax=3500,ymin=0.8,ymax=13000,xlabel={period $n$ of the rare gradient},ylabel={memory $1/(1-\beta_2)$},legend style={at={(0.5,1.03)},anchor=south,font=\scriptsize},legend columns=2]
\addplot[theory] table[x=n,y=memcrit] {figures/data/06-adaptive-reddi-boundary.dat};
\addplot[guide] table[x=n,y=memasym] {figures/data/06-adaptive-reddi-boundary.dat};
\addplot[simulation] table[x=n,y=ok] {figures/data/06-adaptive-reddi-runs.dat};
\addplot[simulation2] table[x=n,y=fail] {figures/data/06-adaptive-reddi-runs.dat};
\addplot[muted,line width=0.4pt,dotted,domain=2.5:3500] {20};
\addplot[muted,line width=0.4pt,dotted,domain=2.5:3500] {1000};
\node[font=\scriptsize,text=muted,anchor=south west,inner sep=1pt] at (axis cs:1200,20.5) {$\beta_2=0.95$};
\node[font=\scriptsize,text=muted,anchor=south west,inner sep=1pt] at (axis cs:1200,1025) {$0.999$};
\legend{{$\kappa_n=0$},{$0.398\,n$},{runs ending at $-1$},{runs ending at $+1$}}
\end{axis}
\end{tikzpicture}
\caption{Where Adam fails on Reddi's problem with period $n$ and $C=2(n-1)$. Curve: the memory at which the drift \eqref{eq:adaptive:kappa} vanishes; dashed: its asymptote; dotted: the memories of $\beta_2=0.95$ and $0.999$. Points: Adam with $\beta_1=0$, bias correction, $\epsilon=10^{-8}$, $\eta_t=0.5/\sqrt t$, from $0$, for $\max(2\cdot10^5,300n)$ rounds; circles end near the minimizer $-1$, squares near the worst point $+1$ (\texttt{06-adaptive-reddi.py}).}
\label{fig:adaptive:reddi}
\end{figure}

\paragraph{The lag of the preconditioner.} \Cref{fig:adaptive:quiet} feeds Adam, as implemented in common libraries, $20000$ quiet gradients of size $10^{-3}$ and then the constant gradient $1$. The steps after the jump follow $\rho_k$ to three digits, and return to the sign step on the time scale $1/(1-\beta_2)$; $\epsilon=0.01$ lowers the first step from $3.16\eta$ to $2.40\eta$. The excess displacement grows like $2\ln2/(1-\beta_2)$. With $\eta=3\cdot10^{-4}$, a coordinate that wakes up with a constant gradient overshoots by $0.36$ at $\beta_2=0.999$, a large change for a weight, and by about $5\cdot10^{-4}$ at $0.95$. The agreement must fail when the quiet gradients are not small compared with the loud one, when $\epsilon$ is comparable with $\sqrt{1-\beta_2}\,G$, and above all when the loud gradient reacts to the steps, as it does in a network: the excess is then an envelope, not a prediction.

\begin{figure}[t]
\begin{fullwidth}
\centering
\begin{tikzpicture}
\begin{groupplot}[group style={group size=2 by 1,horizontal sep=17mm},wplot,width=82mm,height=46mm]
\nextgroupplot[xmode=log,xmin=0.8,xmax=25000,ymin=0,ymax=7.2,xlabel={rounds since the jump, $k+1$},ylabel={step length $\abs{u_k}/\eta$},legend style={at={(0.97,0.97)},anchor=north east,font=\scriptsize}]
\addplot[guide,domain=0.8:25000] {1};
\addplot[theory3] table[x=kp1,y=th95] {figures/data/06-adaptive-quiet-theory.dat};
\addplot[theory2] table[x=kp1,y=th99] {figures/data/06-adaptive-quiet-theory.dat};
\addplot[theory] table[x=kp1,y=th999] {figures/data/06-adaptive-quiet-theory.dat};
\addplot[muted,line width=0.9pt,dashdotted,no marks] table[x=kp1,y=th999e] {figures/data/06-adaptive-quiet-theory.dat};
\addplot[simulation] table[x=kp1,y=sim95] {figures/data/06-adaptive-quiet-sim.dat};
\addplot[simulation] table[x=kp1,y=sim99] {figures/data/06-adaptive-quiet-sim.dat};
\addplot[simulation] table[x=kp1,y=sim999] {figures/data/06-adaptive-quiet-sim.dat};
\addplot[simulation2] table[x=kp1,y=sim999e] {figures/data/06-adaptive-quiet-sim.dat};
\legend{,{$\beta_2=0.95$},{$\beta_2=0.99$},{$\beta_2=0.999$},{$0.999$, $\epsilon=0.01$}}
\nextgroupplot[xmode=log,ymode=log,xmin=15,xmax=13000,ymin=1,ymax=30000,xlabel={memory $1/(1-\beta_2)$},ylabel={excess displacement $E$},legend style={at={(0.03,0.97)},anchor=north west,font=\scriptsize}]
\addplot[guide] table[x=mem,y=lead] {figures/data/06-adaptive-quiet-excess.dat};
\addplot[theory] table[x=mem,y=exact] {figures/data/06-adaptive-quiet-excess.dat};
\addplot[simulation] table[x=mem,y=sim] {figures/data/06-adaptive-quiet-excess-sim.dat};
\legend{{$2\ln2/(1-\beta_2)$},{exact series},{Adam}}
\end{groupplot}
\end{tikzpicture}
\end{fullwidth}
\caption{Quiet then loud: $20000$ Gaussian gradients of standard deviation $10^{-3}$, then the constant gradient $1$; Adam with $\beta_1=0.9$, bias correction and $\epsilon=10^{-8}$ unless stated (seed $0$, \texttt{06-adaptive-quiet.py}). Left: step lengths (points) against $\rho_k$ of \cref{prop:adaptive:quietloud} (curves); dashed, the sign step. Right: the excess over sign steps, summed over $30/(1-\beta_2)$ rounds (points), against the series of \cref{prop:adaptive:quietloud}(ii) and its leading term.}
\label{fig:adaptive:quiet}
\end{figure}

\section{What the theory misses}\label{sec:adaptive:misses}

\paragraph{What was idealized away.} The regret theory holds for every sequence of convex losses, so it uses none of the structure of a fixed nonconvex objective, and it bounds the summed loss along the run, not the final loss; it needs a bounded domain and a projection, and its bounds grow with a diameter nobody knows. The theorem on AdaGrad-Norm is for exact gradients. The propositions about forgetting concern one coordinate with an imposed gradient sequence that does not react to the steps it causes. Throughout, $\epsilon$ is zero or a floor, and the coordinates are fixed.

\paragraph{Where it breaks.} (1) AdaGrad has the best guarantee of the chapter and does not train deep networks: its steps decay for good, which is why RMSProp replaced its sum by an average \parencite{TielemanHinton2012}; against an adversary a decaying step is what one needs, so the worst-case theory cannot see the cost. (2) Adam's convergence theory needs $\beta_2$ near one in a problem-dependent way \parencite{DefossezEtAl2022,ZhangEtAl2022}; language models use $\beta_2=0.95$, a memory of $20$ steps that fails on rare-gradient patterns of period $50$ (\cref{prop:adaptive:reddi}), and yet rare tokens are learned. The theory does not say which gradient patterns occur in a run. (3) The preconditioner is not curvature: it estimates the second moment of the minibatch gradient (\cref{rem:adaptive:limits}), and its close relative, the empirical Fisher, can be a poor model of curvature even in simple problems \parencite{KunstnerHennigBalles2019}. (4) Regret bounds do not rank Adam against SGD on transformers: the gap persists in full batch, and sign descent with momentum closes most of it \parencite{KunstnerEtAl2023}; nothing here predicts that (\cref{ch:transformers}). (5) AMSGrad removes the counterexample and restores the bound, yet no recipe of the practice box uses it: the counterexample is not what limits Adam in practice.

\paragraph{Repairs and open problems.} For convergence: AMSGrad \parencite{ReddiKaleKumar2018}; clipped Adam with a moving average of the iterates, with optimal nonconvex rates and a provable benefit of coordinatewise adaptivity when scales differ across coordinates \parencite{AhnCutkosky2024}; and analyses under smoothness that grows with the gradient \parencite{LiRakhlinJadbabaie2023}. For the preconditioner: Kronecker and full-matrix versions (\cref{ch:kronecker}), and one second moment per block of parameters in Adam-mini \parencite{ZhangEtAl2024}. For $\epsilon$: Adam-atan2 \parencite{EverettEtAl2024}. For $\beta_2$: \textcite{MarekEtAl2025} propose to keep its half-life fixed in tokens when the batch size changes, which makes the memory a property of the data. Open: a model of the gradient sequence of a run from which the right memory can be computed; how often quiet-then-loud events occur in large models, and whether they cause the spikes of \cref{obs:adaptive:beta2}; and why $\beta_2=0.95$ suffices for language models.

\section{Back in the wild}\label{sec:adaptive:back}

\fieldnotes{Adam with $\beta_2\in\{0.95,0.999\}$, $\epsilon=10^{-8}$; AdaGrad for sparse convex problems; Adafactor to save memory.}
{Steepest descent in a diagonal norm chosen from the gradients; online convex losses on a box; smooth losses with unknown $L$; one coordinate with a periodic or a quiet-then-loud gradient.}
{AdaGrad's regret is within $\sqrt2$ of the best diagonal norm in hindsight; AdaGrad-Norm converges at rate $1/T$ for every step size; Adam ends at the worst point exactly when $\kappa_n>0$; the steps after a quiet stretch; Adafactor is the I-divergence projection of $v$.}
{A $\sqrt d/\ln d$ gain on Zipf-sparse features; failure below a memory of $0.4$ periods when $C=2(n-1)$; a first loud step of $3.16\eta$ at $\beta_2=0.999$ and $0.45\eta$ at $0.95$; an overshoot of about $2\ln2/(1-\beta_2)$ steps.}
{Nonconvexity, the structure of real gradient sequences, $\epsilon$ at scale, rotations; why $0.95$ works.}
{$\beta_2$ sets a memory, which must be long compared with the gaps between rare informative gradients and short enough to follow changes of scale (proved in the idealizations, suggested for runs). If a run spikes after quiet stretches, a lower $\beta_2$ shrinks the overshoot in proportion to the memory (suggested). For sparse convex problems, AdaGrad (proved).}

The propositions turn two pieces of folklore into statements about time scales. The memory $1/(1-\beta_2)$ must be long compared with the gaps between the rare gradients that carry information, or they are outvoted: that is PaLM's choice for rare tokens and \cref{prop:adaptive:reddi}. And every quiet-then-loud event costs an overshoot proportional to the memory: that is the case for lowering $\beta_2$, and \cref{prop:adaptive:quietloud} says by how much, a factor of several hundred from $0.999$ to $0.95$. AdaGrad, with infinite memory, meets the first constraint and pays for the second in another currency: its steps shrink for good. Neither says which constraint binds in a given run, and the default $\epsilon$ is a convention, not a result. What the theory establishes is that a step size per coordinate is a choice of norm, made online almost as well as in hindsight, and that each way of forgetting it trades one failure for another.

Adam's other ingredients are elsewhere: sign descent in \cref{ch:sign}, $\beta_1$ in \cref{ch:momentum}, its edge of stability in \cref{ch:eos}, its first steps in \cref{ch:warmup}, quiet-then-loud spikes in \cref{ch:spikes}, the preconditioner $H^{1/2}$ and the critical batch size in \cref{ch:noise}, weight decay in \cref{ch:decay}, and full-matrix and Kronecker AdaGrad in \cref{ch:kronecker}.

\begin{exercises}
\item \consolidation (The AdaGrad lemma and its price.) (a) Prove $\sum_{t\le T}a_t/\sqrt{S_t}\le2\sqrt{S_T}$ by induction on $T$, for $a_t\ge0$ and $S_t=\sum_{u\le t}a_u$. (b) Show that the constant $2$ cannot be improved. (c) For constant $a_t=a>0$ and $v_0=0$, show that both sides of \cref{lem:adaptive:forgetting} grow like $T(1-\beta)$ up to a factor that tends to $1$ as $\epsilon/a\to0$, $\beta\to1$.\label{ex:adaptive:lemma}

\item \consolidation (Bias correction.) (a) For independent $g_t$ with mean $\mu$ and second moment $\sigma^2$, show $\E m_t=(1-\beta_1^t)\mu$ and $\E v_t=(1-\beta_2^t)\sigma^2$. (b) Show that with bias correction and $\epsilon=0$ the first step is $-\eta\sign(g_1)$ for all $\beta_1,\beta_2$. (c) Show that in \cref{prop:adaptive:quietloud} bias correction multiplies the step by between $(1-\beta_2^{t_0})^{1/2}$ and $(1-\beta_1^{t_0})^{-1}$.\label{ex:adaptive:bias}

\item \consolidation (Reddi's choice, and AMSGrad.) (a) For $n=3$ and $\beta_2=1/(1+C^2)$, show $\kappa_3\to2^{-1/2}$ as $C\to\infty$. (b) AMSGrad replaces $v_t$ by $\max_{s\le t}v_s$. In \cref{prop:adaptive:reddi} without bias correction, show that the maximum tends to $w_0$, that the drift per period tends to $(n-1-C)/\sqrt{w_0}<0$, and that AMSGrad converges to $-1$.\label{ex:adaptive:reddichoice}

\item \consolidation\hard (Memory against period.) Let $C=c\,(n-1)$ with $c>1$ fixed and let $n\to\infty$ with $(1-\beta_2)n\to x>0$. (a) Show $\kappa_n\sqrt q\to-1+(2/(cx))(e^{x/2}-1)$. (b) Deduce that the critical memory is asymptotic to $n/(2y_c)$, where $y_c>0$ solves $e^y-1=cy$; for $c=2$ it is $0.398\,n$. (c) Explain why it grows faster than $n$ as $c\downarrow1$ (for $C=n$ it is about $n^2/4$). \hint{for lags of order $n$, $q\beta_2^kC^2\to\infty$; the sum over $k$ becomes a Riemann sum of $e^{xs/2}$.}\label{ex:adaptive:period}

\item \consolidation (Adafactor.) (a) For $V=\bigl(\begin{smallmatrix}4&1\\1&1\end{smallmatrix}\bigr)$, compute $R^\star$ and the best rank-one approximation in Frobenius norm, and check that only $R^\star$ has the row sums of $V$. (b) Show that $D(V\Vert R^\star)=S\cdot I$, with $I$ the mutual information of the row and column indices under $V/S$. (c) Show that Adafactor is exact if every $G_t^2=x_ty^\top$ with one $y$.\label{ex:adaptive:adafactor}

\item \variation\hard (Full-matrix AdaGrad.) Let $P_t=(\sum_{s\le t}g_sg_s^\top)^{1/2}$, let $\Theta$ be convex with Euclidean diameter $D_2$, and project in $\norm\cdot_{P_t}$. (a) Show $R_T\le(D_2^2/(2\eta)+\eta)\tr P_T$. (b) Show $\min\{\sum_tg_t^\top S^{-1}g_t\colon S\succ0,\ \tr S\le1\}=(\tr P_T)^2$ if $P_T\succ0$. \hint{use $P_t+\delta I$, $\delta\to0$; $P_t\succeq P_{t-1}$ by \cref{thm:toolkit:loewner}; concavity of $\tr X^{1/2}$ gives $\tr A^{1/2}-\tr(A-gg^\top)^{1/2}\ge\frac12g^\top A^{-1/2}g$.}\label{ex:adaptive:fullmatrix}

\item \variation (Forgetting breaks AdaGrad-Norm.) Replace the sum in \eqref{eq:adaptive:adagradnorm} by $b_{t+1}^2=\beta_2b_t^2+(1-\beta_2)\norm{\nabla\Loss(\theta_t)}_*^2$ and run it on $\Loss(\theta)=\frac h2\theta^2$. (a) Show that $\pm\eta/2$, with $b=h\eta/2$, is an orbit of period two for every $\beta_2<1$. (b) Show that the iterates do not converge unless some iterate is exactly $0$. (c) Simulate $\beta_2=0.99$, $\eta=0.3$, $h=1$, $\theta_0=1$, $b_0=10^{-3}$: the loss falls below $10^{-40}$ while $b$ remembers the early gradients, then returns to $h\eta^2/8$. Why does $\beta_2=0.999$ seem to converge in double precision?\label{ex:adaptive:rmsnorm}

\item \variation (Adam without memory, with $\epsilon$.) On $\Loss=\frac h2\theta^2$, Adam with $\beta_1=\beta_2=0$ is $\theta\mapsto\theta-\eta h\theta/(h\abs\theta+\epsilon)$. (a) Show that $\theta_t\to0$ from every start if $\epsilon\ge\eta h/2$. (b) Show that for $\epsilon<\eta h/2$ there is an orbit of period two with amplitude $\eta/2-\epsilon/h$, below the amplitude $\eta/2$ of the symmetric two-cycle of sign descent. (c) Relate the threshold to the stability of SGD with step size $\eta/\epsilon$.\label{ex:adaptive:eps}

\item \variation (Momentum in Reddi's problem.) With $\beta_1>0$, show that the periodic first moment at lag $k$ is $-1+(C+1)\beta_1^k(1-\beta_1)/(1-\beta_1^n)$ and the drift per period is $-\sum_km_k/\sqrt{w_k}$. For $\beta_1=0.9$, $C=2(n-1)$ and memories above $1/(1-\beta_1^2)$, find the critical memory for $n=100$ and $1000$ (about $34$ and $392$, against $40$ and $398$ for $\beta_1=0$).\label{ex:adaptive:momentum}

\item \simulation (Sparse features.) Run \texttt{06-adaptive-regret.py} with $\alpha=1$ and $\alpha=3$. Predict the ratio of the bounds of \cref{cor:adaptive:sparse} at $d=10^4$ (about $1.6$ and $42$), compare with the measured ratio, and explain the deviation through the features with $p_iT<1$.\label{ex:adaptive:sparse}

\item \simulation (Random arrivals.) Let the large gradient of \cref{prop:adaptive:reddi} arrive at each round independently with probability $1/n$, as in the stochastic counterexample of \textcite{ReddiKaleKumar2018}. For $n=10,100,1000$ and $C=2(n-1)$, estimate the stationary drift $-\E[g_t/\sqrt{v_t}]$ for $\beta_1=0$ as a function of the memory, locate its zero (about one period, against $0.4$ for the periodic case), and explain why random gaps hurt.\label{ex:adaptive:random}

\item \research (Choosing the memory.) Build a model of one coordinate whose gradient has both rare informative spikes (\cref{prop:adaptive:reddi}) and changes of scale (\cref{prop:adaptive:quietloud}), and find the memory that minimizes a natural loss. How does it scale with the batch size, if a step consumes $B$ examples, and is the rule of \textcite{MarekEtAl2025} recovered? Then measure, in a small transformer, how often coordinates are quiet and then loud. As far as we know both questions are open.\label{ex:adaptive:research1}

\end{exercises}

\notesandrefs
AdaGrad appeared at COLT 2010, in \textcite{McMahanStreeter2010} and in the work of Duchi, Hazan and Singer whose journal version, \textcite{DuchiHazanSinger2011}, also gives the full-matrix version and sparse-data examples like \cref{cor:adaptive:sparse}. Online gradient descent and its regret bound are due to \textcite{Zinkevich2003}. \Textcite{ReddiKaleKumar2018} traced Adam's failure to the quantity that our proof of \cref{thm:adaptive:adagrad} needs to be nonnegative, which started the theory of \cref{sec:adaptive:forgetting}. \Cref{thm:adaptive:adagradnorm} is the deterministic part of \textcite{WardWuBottou2019}, in a general norm with our own constants.

\Textcite{BernsteinNewhouse2024} present Adam without its moving averages as steepest descent in $\ell_\infty$, the view of \cref{ch:sign}; this chapter adds the norm that the averages define. Adafactor's I-divergence factorization is from \textcite{ShazeerStern2018}. We know no reference for the drift $\kappa_n$, the critical memory $0.398\,n$ or the excess $2\ln2/(1-\beta_2)$; they are elementary consequences of the two models.

\paragraph{Translating the literature.} \Textcite{DuchiHazanSinger2011} write $x_t$ for our $\theta_t$ and use our $g_{1:t,i}$ and $s_{t,i}$; their $H_t=\delta I+\diag(s_t)$ is our $P_t$ with $\delta$ for $\epsilon$, and their $D_\infty$ is a diameter, as ours. \Textcite{KingmaBa2015} write $\alpha$ for the step size $\eta$, and otherwise our notation. \Textcite{ReddiKaleKumar2018} write $\alpha_t$ for $\eta_t$ and $V_t=\diag(v_t)$; their $\Gamma_{t+1}=\sqrt{V_{t+1}}/\alpha_{t+1}-\sqrt{V_t}/\alpha_t$ is the increment of our $P_t/\eta_t$. \Textcite{DefossezEtAl2022} write $F$, $N$, $\alpha$ for $\Loss$, $T$, $\eta$, and $R$ for the bound on the gradients. \Textcite{ShazeerStern2018} write $R_t$, $C_t$ for our $r_t$, $c_t$ and $\hat\beta_{2t}$ for the increasing coefficient.

\printchapterbib
\end{refsection}

% Chapter 7 of Optimizers in the Wild: Kronecker preconditioners (Shampoo, K-FAC, SOAP).

\casechapter{Kronecker Preconditioners}{Shampoo, K-FAC and SOAP}{Shampoo, K-FAC, SOAP}{ch:kronecker}
\begin{refsection}[chapters/07-kronecker.bib]

\begin{chaptersummary}
Shampoo, K-FAC, and SOAP precondition the gradient of each weight matrix by two small matrices, one for its rows and one for its columns, instead of the enormous matrix of full-matrix AdaGrad or the Fisher information; we study them on the vectorized gradient and on a layer loss whose Hessian is a Kronecker product. We prove that Shampoo's factors dominate full-matrix AdaGrad's statistics up to the rank of the gradients, which gives Shampoo its regret bound; that with stationary statistics a Kronecker preconditioner with power $s$ turns the condition number $\kappa$ into $\kappa^{\abs{1-qs}}$, where $q$ describes the statistics, so that Shampoo's power $1/4$, the power $1/2$, and K-FAC's inverse each win in a different regime; and that SOAP, Adam in Shampoo's eigenbasis, runs on Kronecker curvature as on an isotropic quadratic. The theory does not explain why these methods beat a well-tuned Adam when they do, and their measured advantage shrinks as models grow.
\end{chaptersummary}

\section{In the wild: two small matrices per layer}\label{sec:kronecker:wild}

Adam treats a weight matrix as a bag of numbers: each entry gets a learning rate from its own past (\cref{ch:adaptive}), and nothing in the update knows that rows belong to output units and columns to inputs. Full-matrix AdaGrad, which preconditions the vectorized gradient by the inverse square root of the summed outer products of past gradients, sees every correlation, but needs $(mn)^2$ numbers for an $m\times n$ matrix: $2.8\times10^{14}$ for a $4096\times4096$ layer.\sidenote{At four bytes each, about a petabyte for one layer. Even the diagonal that Adam keeps has $1.7\times10^7$ entries.} The optimizers of this chapter keep one small matrix for the rows and one for the columns, $m^2+n^2$ numbers, and multiply the gradient by them on both sides.

\begin{practice}[Shampoo, \textcite{GuptaKorenSinger2018,AnilEtAl2020}; K-FAC, \textcite{MartensGrosse2015}; SOAP, \textcite{VyasEtAl2025}]
For each weight matrix $W\in\R^{m\times n}$ with minibatch gradient $G_t$, keep $L_t\in\R^{m\times m}$, a sum or moving average of $G_tG_t^\top$, and $R_t\in\R^{n\times n}$, the same for $G_t^\top G_t$; every so many steps recompute inverse roots, and step $W\leftarrow W-\eta(L_t+\eps I)^{-1/4}G_t(R_t+\eps I)^{-1/4}$ (Shampoo). Large-scale implementations take the step length from Adam or AdaGrad (grafting), add momentum, and some use the power $-1/2$. K-FAC averages $aa^\top$ over a layer's inputs $a$ and $\delta\delta^\top$ over the gradients $\delta$ with respect to its outputs, and steps with $(\E[\delta\delta^\top]+\eps I)^{-1}G_t(\E[aa^\top]+\eps I)^{-1}$, with an adaptive damping $\eps$. SOAP runs Adam on $Q_L^\top G_tQ_R$, where $Q_L$ and $Q_R$ hold eigenvectors of Shampoo's $L_t$ and $R_t$, recomputed every $f$ steps.
\end{practice}

\Textcite{GuptaKorenSinger2018} introduced \term{Shampoo}\index{Shampoo} for tensors of every order, with one factor per mode and the power $-1/(2k)$ for order $k$, which is $-1/4$ for matrices, and proved its regret bound, the main theorem of this chapter.\sidenote{For a tensor of order $k$ they prove an analogue of \cref{thm:kronecker:bound} with $k$ factors, each raised to the power $1/k$; AdaGrad's square root then gives $1/(2k)$.} \Textcite{AnilEtAl2020} made it practical at scale, on Transformers for translation, BERT, and ResNet-50 among other models, computing the roots rarely and in parallel with training, and grafting its step length onto a first-order method \parencite{AgarwalEtAl2020Grafting}; \textcite{ShiEtAl2023} describe a PyTorch implementation. \term{K-FAC}\index{K-FAC} comes from natural gradient descent \parencite{Amari1998}\index{natural gradient}: \textcite{MartensGrosse2015} approximated each layer's block of the Fisher matrix by the Kronecker product of the second moments of the layer's inputs and of its back-propagated gradients, so that inverting it takes two small inversions.\sidenote{\Textcite{GeorgeEtAl2018} kept K-FAC's Kronecker eigenbasis but re-estimated the scale of each eigendirection from the gradients, an eigenvalue-corrected K-FAC that anticipates SOAP.} SOAP \parencite{VyasEtAl2025} is Shampoo with Adam inside.

\begin{observation}[AlgoPerf; \textcite{DahlEtAl2023,KasimbegEtAl2025}]\label[observation]{obs:kronecker:algoperf}
AlgoPerf times training algorithms to fixed validation targets on a set of workloads. In its external-tuning ruleset a submission supplies one hyperparameter search space for all workloads, from which five trials are drawn; in its self-tuning ruleset the submission must set its hyperparameters itself. In the first competition, with $18$ submissions from $10$ teams, a submission based on Distributed Shampoo \parencite{ShiEtAl2023} won the external-tuning ruleset, training about $28\%$ faster on average than the baseline and beating Adam-type methods also in wall-clock time; one based on Schedule-Free AdamW won the self-tuning ruleset.
\end{observation}

\begin{observation}[Language models; \textcite{VyasEtAl2025,WenEtAl2026Fantastic}]\label[observation]{obs:kronecker:scale}
On language models with $360$ million and $660$ million parameters trained with large batches, SOAP needed over $40\%$ fewer iterations and over $35\%$ less wall-clock time than AdamW to reach the same loss, and about $20\%$ less of each than Shampoo. A later study that tuned every optimizer separately at every size found that the speedup of matrix-based optimizers over AdamW, SOAP's included, shrinks as the models grow, from about $1.4\times$ at $0.1$ billion parameters to about $1.1\times$ at $1.2$ billion.
\end{observation}

\begin{observation}[The heuristics carry weight; \textcite{VyasEtAl2025,EschenhagenEtAl2025}]\label[observation]{obs:kronecker:heuristics}
Shampoo as used differs from Shampoo as analyzed: grafting, $\eps$, moving averages, the power, and the interval between root computations are set by experiment, and computing the eigendecomposition less often degrades Shampoo, increasingly with the interval. \Textcite{EschenhagenEtAl2025} find that grafting from Adam compensates for stale and mis-scaled eigenvalues of the factors, and that correcting the eigenvalues directly makes grafting unnecessary.
\end{observation}

The three methods come with three stories: an online-learning algorithm with a regret bound, an approximation of the natural gradient, and a repair of the first by way of Adam. Why can two small matrices stand in for a large one, and at what price? What do they buy on curvature with Kronecker structure, and which power of them should one take? In an idealization the three methods differ mainly in two choices: the power of the factors, and the basis of a diagonal correction.

\section{The idealization: AdaGrad on the vectorized gradient}\label{sec:kronecker:ideal}

Throughout, $W\in\R^{m\times n}$ is one weight matrix ($m=d_{\mathrm{out}}$ and $n=d_{\mathrm{in}}$ for a linear layer), $G_t$ is its gradient at step $t$, and $g_t=\vect(G_t)\in\R^{mn}$, where $\vect$ stacks the columns. From \cref{lem:toolkit:kron} we use $\vect(MXN)=(N^\top\otimes M)\vect(X)$, so that $\vect(uv^\top)=v\otimes u$; the mixed product rule $(M\otimes N)(M'\otimes N')=MM'\otimes NN'$; that the eigenvalues of $M\otimes N$ are the products of those of $M$ and $N$; that $(M\otimes N)^s=M^s\otimes N^s$ for positive semidefinite $M,N$ and $s>0$; and that $M\otimes N\preceq M'\otimes N'$ when $0\preceq M\preceq M'$ and $0\preceq N\preceq N'$, where $\preceq$ is the Löwner order.\sidenote{Some papers stack rows instead of columns, which reverses the order of every Kronecker product (Exercise~\ref{ex:kronecker:vec}); it is a common source of transposed formulas in this literature.}

\paragraph{Full-matrix AdaGrad.} With $\Gamma_t=\sum_{\tau\le t}g_\tau g_\tau^\top$, \term{full-matrix AdaGrad}\index{AdaGrad!full-matrix} \parencite{DuchiHazanSinger2011} steps $\vect(W_{t+1})=\vect(W_t)-\eta(\eps I+\Gamma_t)^{-1/2}g_t$: unnormalized steepest descent in the norm $\norm{\cdot}_{P_t}$ with $P_t=(\eps I+\Gamma_t)^{1/2}$ (\cref{ch:adaptive}). \Cref{thm:adaptive:adagrad} gives the regret of the diagonal version, and the proof of \cref{thm:kronecker:regret} below, with $r=1$ and $(\eps I+\Gamma_t)^{1/2}$ as preconditioner, gives at $\eta=D/\sqrt2$ the full-matrix bound $\sqrt2D\tr((\eps I+\Gamma_T)^{1/2})$. The trouble is the size of $\Gamma_t$. A preconditioner $P_R\otimes P_L$, with $P_L\in\R^{m\times m}$ and $P_R\in\R^{n\times n}$ positive definite, is cheap by comparison: $(P_R\otimes P_L)^{-1}\vect(G)=\vect(P_L^{-1}GP_R^{-1})$.

\begin{definition}[Kronecker preconditioners]\label[definition]{def:kronecker:shampoo}
Let $\eps>0$ and $s>0$, and let $L_t=\eps I_m+\sum_{\tau\le t}G_\tau G_\tau^\top$ and $R_t=\eps I_n+\sum_{\tau\le t}G_\tau^\top G_\tau$. \term{Shampoo with power $s$} updates $W_{t+1}=W_t-\eta\,L_t^{-s}G_tR_t^{-s}$, which on $\vect(W)$ is preconditioned gradient descent with the \term{Kronecker preconditioner}\index{Kronecker preconditioner} $R_t^s\otimes L_t^s$; Shampoo proper is the case $s=1/4$. A \term{frozen} Kronecker preconditioner uses fixed positive definite $P_L$ and $P_R$ in place of $L_t^s$ and $R_t^s$.
\end{definition}

Shampoo is to the spectral norm what Adam is to the max norm: with $\eps=0$ and a single gradient its step is $\polar(G_t)$, Muon's update without momentum (\cref{prop:spectral:shampoo}), and with sums over many gradients it is a slowly changing preconditioner, as Adam moves between sign descent at $\beta_1=\beta_2=0$ and a slowly changing diagonal preconditioner (\cref{ch:sign,ch:adaptive}; \cref{prop:kronecker:between} below). To ask what the factors buy, we need a loss whose curvature they can match.

\begin{assumption}[Kronecker curvature]\label[assumption]{ass:kronecker:quadratic}
Let $A\in\R^{m\times m}$ and $C\in\R^{n\times n}$ be symmetric positive definite, with eigenvalues $\alpha_1\ge\dots\ge\alpha_m>0$ and $\gamma_1\ge\dots\ge\gamma_n>0$ and condition numbers $\kappa_A=\alpha_1/\alpha_m$ and $\kappa_C=\gamma_1/\gamma_n$. The loss is the \term{Kronecker quadratic}\index{Kronecker quadratic}
\[
  \Loss(W)=\tfrac12\tr\bigl((W-W_\star)^\top A\,(W-W_\star)\,C\bigr)=\tfrac12\bigl\lVert A^{1/2}(W-W_\star)C^{1/2}\bigr\rVert_F^2 .
\]
\end{assumption}

This is the expected loss of a linear layer $a\mapsto Wa$ whose inputs have second moment $\E[aa^\top]=C$ and whose output enters a loss that is quadratic with Hessian $A$: if $\ell(W;a)=\frac12(Wa-W_\star a)^\top A(Wa-W_\star a)$, then $\E[\ell(W;a)]=\Loss(W)$. A linear layer with a Gaussian likelihood is an instance (\cref{prop:kronecker:kfac}). The input-side matrix is called $C$, not $B$, because $B$ is the batch size; the Hessian is then $C\otimes A$. For a matrix $M$ with positive eigenvalues we write $\kappa(M)=\lmax(M)/\lmin(M)$.

\begin{dictionary}
$W\in\R^{m\times n}$, its gradients $G_t$, and $\Gamma_T=\sum_tg_tg_t^\top$ & a weight matrix, its minibatch gradients, and full-matrix AdaGrad's statistics, never formed in practice\\
the factors $L_t$, $R_t$; the rank bound $r$ & Shampoo's two matrices; $r$ is at most the number of examples (or tokens) in a minibatch\\
the step size $\eta$ & the learning rate (two names for one number)\\
the power $s$ & the exponent of the inverse roots: $\frac14$ in the theory, $\frac12$ in some implementations, $1$ in K-FAC\\
the Kronecker quadratic, $H=C\otimes A$ & a layer with input second moment $C$ and output-side curvature $A$, as K-FAC models it\\
$\E[gg^\top]\propto H^q$ with $q=1$; with $q=2$ & gradient noise near a minimum when the Fisher matrix is the Hessian; gradients dominated by an isotropic error, as after a random initialization
\end{dictionary}

\begin{fieldguide}{Shampoo}{\textcite{GuptaKorenSinger2018}}
\fgrow{update}{$L_t=L_{t-1}+G_tG_t^\top$, \ $R_t=R_{t-1}+G_t^\top G_t$ \ ($L_0=R_0=\eps I$), \ $W_{t+1}=W_t-\eta\,L_t^{-1/4}G_tR_t^{-1/4}$}
\fgrow{norm}{spectral when memoryless: one gradient and $\eps=0$ give $\polar(G_t)$}
\fgrow{preconditioner}{$R_t^{1/4}\otimes L_t^{1/4}$; its square dominates $\Gamma_t/r$ (\cref{thm:kronecker:bound})}
\fgrow{averages}{sums or moving averages of $GG^\top$ and $G^\top G$; momentum in practice}
\fgrow{state}{$m^2+n^2$ numbers for an $m\times n$ matrix, plus inverse roots}
\fgrow{in the wild}{power $-\frac14$ or $-\frac12$; infrequent roots; step length grafted from Adam}
\end{fieldguide}

\begin{fieldguide}{K-FAC}{\textcite{MartensGrosse2015}}
\fgrow{update}{averages $\bar S_a$ of $aa^\top$ and $\bar S_\delta$ of $\delta\delta^\top$; \ $W\leftarrow W-\eta(\bar S_\delta+\eps I)^{-1}G(\bar S_a+\eps I)^{-1}$}
\fgrow{norm}{the Fisher (natural-gradient) norm, layer by layer}
\fgrow{preconditioner}{$\bar S_a\otimes\bar S_\delta$, damped; power $-1$}
\fgrow{averages}{moving averages of both factors; momentum}
\fgrow{state}{$m^2+n^2$ numbers per layer, plus inverses}
\fgrow{in the wild}{targets sampled from the model (Fisher) or the data (empirical Fisher)}
\end{fieldguide}

\begin{fieldguide}{SOAP}{\textcite{VyasEtAl2025}}
\fgrow{update}{every $f$ steps $Q_L,Q_R\leftarrow$ eigenvectors of averages of $GG^\top$, $G^\top G$; every step Adam on $Q_L^\top G_tQ_R$, rotated back}
\fgrow{norm}{$\ell_\infty$ in the rotated coordinates when memoryless}
\fgrow{preconditioner}{diagonal in Shampoo's eigenbasis}
\fgrow{averages}{Adam's two moving averages, kept in the current basis}
\fgrow{state}{$2mn+2(m^2+n^2)$ numbers}
\fgrow{in the wild}{Adam's $\beta_1,\beta_2$ and the preconditioning frequency $f$}
\end{fieldguide}

\paragraph{What was thrown away.} Everything but one matrix: the other layers, the nonlinearities, and the couplings between layers that K-FAC also drops. The curvature is fixed, and in \cref{sec:kronecker:curvature} so are the statistics behind the factors; momentum, grafting, $\eps$, and the infrequency of the roots are left out; the regret theory is online and convex. These are the right first simplifications because the Kronecker structure is exact for one linear layer with independent inputs and back-propagated gradients (\cref{prop:kronecker:kfac}), and stationary statistics are what the factors' averages estimate: whatever the factors can do, they do best here.

\paragraph{What one had to notice.} Two facts. The summed outer products of matrix gradients are dominated, up to the rank, by a Kronecker product of the Gram sums $\sum GG^\top$ and $\sum G^\top G$; seeing this takes the matrix geometric mean, whose square root, composed with AdaGrad's, gives Shampoo's power $1/4$. And a layer's per-example gradient $\delta a^\top$ has the second moment $\E[(aa^\top)\otimes(\delta\delta^\top)]$, which factorizes when $a$ and $\delta$ are independent; this makes Shampoo's and K-FAC's factors estimates of the same two matrices.

\section{The theory: the Kronecker bound and Shampoo's regret}\label{sec:kronecker:bound}

The comparison of Shampoo with full-matrix AdaGrad rests on one gradient at a time.

\begin{lemma}[One gradient]\label[lemma]{lem:kronecker:single}
Let $G\in\R^{m\times n}$ have rank $r\ge1$, and let $g=\vect(G)$. Then
\begin{equation}\label{eq:kronecker:single}
  gg^\top\ \preceq\ r\,(G^\top G)^{1/2}\otimes(GG^\top)^{1/2},
\end{equation}
and the factor $r$ cannot be replaced by a smaller number.
\end{lemma}

\begin{proof}
The plan is to compute both quadratic forms in the singular basis of $G$ and compare them by the Cauchy--Schwarz inequality. Write a compact singular value decomposition $G=U\Sigma V^\top$, with $U\in\R^{m\times r}$ and $V\in\R^{n\times r}$ having orthonormal columns and $\Sigma=\diag(\sigma_1,\dots,\sigma_r)$ positive, so that $(GG^\top)^{1/2}=U\Sigma U^\top$ and $(G^\top G)^{1/2}=V\Sigma V^\top$. For $X\in\R^{m\times n}$ and $x=\vect(X)$, put $Y=U^\top XV$ and $Z=\Sigma^{1/2}Y\Sigma^{1/2}$. By the vec identity, the right side of \eqref{eq:kronecker:single} without the factor $r$ has the quadratic form
\[
  x^\top\bigl((V\Sigma V^\top)\otimes(U\Sigma U^\top)\bigr)x=\ip{X,U\Sigma U^\top XV\Sigma V^\top}=\tr\bigl(Y^\top\Sigma Y\Sigma\bigr)=\norm{Z}_F^2 ,
\]
and the left side has $(g^\top x)^2=\ip{G,X}^2=(\tr(\Sigma U^\top XV))^2=(\tr Z)^2$. The Cauchy--Schwarz inequality $(\tr Z)^2=\ip{I_r,Z}^2\le r\norm{Z}_F^2$ proves \eqref{eq:kronecker:single}. Equality holds for $Z=I_r$, that is, for $X=U\Sigma^{-1}V^\top$, where the two forms equal $r^2$ and $r$; so no smaller factor works.
\end{proof}

For a sum of gradients the two sides of \eqref{eq:kronecker:single} cannot simply be added, because the square root of a sum is not the sum of the square roots; the \term{matrix geometric mean}\index{matrix geometric mean} repairs this. The geometric mean $P\#Q$ of positive definite $P$ and $Q$ is the largest symmetric $Y$ with $\bigl(\begin{smallmatrix}P&Y\\Y&Q\end{smallmatrix}\bigr)\succeq0$, and it equals $(PQ)^{1/2}$ when $PQ=QP$ (\cref{thm:toolkit:geomean}). Block matrices add, so the geometric mean is jointly concave, and for the commuting pairs $M\otimes I$ and $I\otimes N$ this says that $(M,N)\mapsto M^{1/2}\otimes N^{1/2}$ is superadditive: $\sum_tM_t^{1/2}\otimes N_t^{1/2}\preceq(\sum_tM_t)^{1/2}\otimes(\sum_tN_t)^{1/2}$ for positive semidefinite $M_t$ and $N_t$ (\cref{cor:toolkit:kronmean}).\sidenote{\Textcite[Chapter~4]{Bhatia2007} gives an account of the geometric mean, including its joint concavity.}

\begin{theorem}[The Kronecker bound; \textcite{GuptaKorenSinger2018}]\label{thm:kronecker:bound}
Let $G_1,\dots,G_T\in\R^{m\times n}$ have rank at most $r\ge1$, let $g_t=\vect(G_t)$ and $\eps>0$, and let $L_T=\eps I_m+\sum_{t=1}^TG_tG_t^\top$ and $R_T=\eps I_n+\sum_{t=1}^TG_t^\top G_t$. Then
\begin{equation}\label{eq:kronecker:bound}
  \eps I_{mn}+\frac1r\sum_{t=1}^Tg_tg_t^\top\ \preceq\ \eps I_{mn}+\sum_{t=1}^T(G_t^\top G_t)^{1/2}\otimes(G_tG_t^\top)^{1/2}\ \preceq\ R_T^{1/2}\otimes L_T^{1/2} .
\end{equation}
\end{theorem}

\begin{proof}
The first inequality follows from \cref{lem:kronecker:single}: a nonzero $G_t$ of rank $r_t$ has $g_tg_t^\top\preceq r_t(G_t^\top G_t)^{1/2}\otimes(G_tG_t^\top)^{1/2}$, and $r_t\le r$; gradients equal to zero contribute nothing to either side. The second is \cref{cor:toolkit:kronmean} applied to the $T+1$ pairs $(\eps I_n,\eps I_m)$ and $(G_t^\top G_t,G_tG_t^\top)$ for $t\le T$: the first pair contributes $(\eps I_n)^{1/2}\otimes(\eps I_m)^{1/2}=\eps I_{mn}$, so the left side of the corollary is the middle term of \eqref{eq:kronecker:bound}, and the sums of the pairs are $R_T$ and $L_T$.
\end{proof}

Letting $\eps\to0$, the outer inequality also holds without $\eps$. In words: the square of Shampoo's preconditioner dominates full-matrix AdaGrad's statistics divided by the rank. To see how large the constant is, and what Shampoo does when the statistics themselves are a Kronecker product, it helps to view the factors as functions of $\Gamma$.

\begin{proposition}[Shampoo's factors are partial traces]\label[proposition]{prop:kronecker:partial}
For positive semidefinite $\Gamma\in\R^{mn\times mn}$ define the \term{partial traces}\index{partial trace} $L(\Gamma)\in\R^{m\times m}$ and $R(\Gamma)\in\R^{n\times n}$ by $\tr(L(\Gamma)M)=\tr(\Gamma(I_n\otimes M))$ and $\tr(R(\Gamma)N)=\tr(\Gamma(N\otimes I_m))$ for all $M\in\R^{m\times m}$ and $N\in\R^{n\times n}$.
\begin{enumerate}
\item If $\Gamma=\sum_tg_tg_t^\top$, then $L(\Gamma)=\sum_tG_tG_t^\top$ and $R(\Gamma)=\sum_tG_t^\top G_t$; likewise $L(\E[gg^\top])=\E[GG^\top]$ and $R(\E[gg^\top])=\E[G^\top G]$.
\item If $\Gamma=\Phi\otimes\Psi$ with $\Phi\in\R^{n\times n}$ and $\Psi\in\R^{m\times m}$ positive semidefinite, then $L(\Gamma)=(\tr\Phi)\Psi$, $R(\Gamma)=(\tr\Psi)\Phi$, and for every $s>0$
\begin{equation}\label{eq:kronecker:power}
  R(\Gamma)^s\otimes L(\Gamma)^s=(\tr\Gamma)^s\,\Gamma^s .
\end{equation}
\end{enumerate}
\end{proposition}

\begin{proof}
The trace pairing determines $L(\Gamma)$ and $R(\Gamma)$, which are linear in $\Gamma$, so for (i) it suffices to take $\Gamma=gg^\top$: then $\tr(gg^\top(I_n\otimes M))=g^\top\vect(MG)=\tr(GG^\top M)$ and $\tr(gg^\top(N\otimes I_m))=g^\top\vect(GN^\top)=\tr(NG^\top G)$. For (ii), $\tr((\Phi\otimes\Psi)(I_n\otimes M))=\tr(\Phi)\tr(\Psi M)$ and $\tr((\Phi\otimes\Psi)(N\otimes I_m))=\tr(\Phi N)\tr(\Psi)$; then \eqref{eq:kronecker:power} follows from $(\Phi\otimes\Psi)^s=\Phi^s\otimes\Psi^s$ and $\tr(\Phi\otimes\Psi)=\tr\Phi\tr\Psi$.
\end{proof}

So for Kronecker statistics, and with $\eps=0$, Shampoo's preconditioner is $(\tr\Gamma)^{1/4}\Gamma^{1/4}$, a multiple of the square root of AdaGrad's $\Gamma^{1/2}$, and the power $1/2$ gives AdaGrad's $\Gamma^{1/2}$ up to a scalar. The Kronecker product nearest to $\Gamma$ in Frobenius norm is a rank-one approximation of a rearrangement of $\Gamma$ \parencite{VanLoanPitsianis1993}, and one step of the power method for it, from identity matrices, returns $R(\Gamma)$ and $L(\Gamma)$ (Exercise~\ref{ex:kronecker:constant}); \textcite{MorwaniEtAl2025} made this observation and found the one-step approximation close to optimal in neural networks, so that the square of Shampoo's preconditioner approximates full-matrix AdaGrad's up to a scalar.

\begin{proposition}[How large the constant is]\label[proposition]{prop:kronecker:constant}
For positive semidefinite $\Gamma\ne0$ let $c^\star(\Gamma)$ be the smallest $c$ with $\Gamma\preceq c\,R(\Gamma)^{1/2}\otimes L(\Gamma)^{1/2}$. For gradients of rank at most $r$, \cref{thm:kronecker:bound} gives $c^\star\le r$, with equality for a single gradient of rank $r$ (\cref{lem:kronecker:single}). Moreover:
\begin{enumerate}
\item If $G_t=\sum_{i=1}^ks_{ti}u_iv_i^\top$ for orthonormal $u_1,\dots,u_k$, orthonormal $v_1,\dots,v_k$ and a matrix $S=(s_{ti})\in\R^{T\times k}$ without zero columns, then $c^\star$ is the largest eigenvalue of the cosine matrix $\diag(d)^{-1/2}S^\top S\diag(d)^{-1/2}$, where $d_i=\sum_ts_{ti}^2$; it lies between $1$ and $k$.
\item If $\Gamma=\Phi\otimes\Psi$, then $c^\star=(\lmax(\Gamma)/\tr\Gamma)^{1/2}$, which can be as small as $(mn)^{-1/2}$.
\end{enumerate}
\end{proposition}

\begin{proof}
The plan in (i) is to find a common invariant subspace. The vectors $w_i=v_i\otimes u_i=\vect(u_iv_i^\top)$ are orthonormal, $g_t=\sum_is_{ti}w_i$, and $\Gamma=\sum_{i,j}(S^\top S)_{ij}w_iw_j^\top$ vanishes off their span. By orthonormality $L(\Gamma)=\sum_id_iu_iu_i^\top$ and $R(\Gamma)=\sum_id_iv_iv_i^\top$, so $K=R(\Gamma)^{1/2}\otimes L(\Gamma)^{1/2}$ has the eigenvectors $w_i$ with eigenvalues $d_i$, and leaves their span and its orthogonal complement invariant. On the span, $\Gamma\preceq cK$ reads $S^\top S\preceq c\diag(d)$, that is, $\diag(d)^{-1/2}S^\top S\diag(d)^{-1/2}\preceq cI$; on the complement, $\Gamma=0\preceq cK$. A unit diagonal puts the largest eigenvalue between $1$ and the trace $k$. For (ii), \eqref{eq:kronecker:power} with $s=\frac12$ gives $K=(\tr\Gamma)^{1/2}\Gamma^{1/2}$, and $\Gamma\preceq cK$ means $\lambda_i\le c(\tr\Gamma)^{1/2}\lambda_i^{1/2}$ for every eigenvalue $\lambda_i$ of $\Gamma$.
\end{proof}

The factor $r$ is the price of coherence, which the regret bound pays in full (\cref{fig:kronecker:constant}).

The losses are now convex functions $\ell_t$ on $\R^{m\times n}$, revealed one at a time, and the \term{regret}\index{regret} against a fixed $W^\star$ is $\sum_t(\ell_t(W_t)-\ell_t(W^\star))$.

\begin{theorem}[Shampoo's regret; \textcite{GuptaKorenSinger2018}]\label{thm:kronecker:regret}
Let $\ell_1,\dots,\ell_T$ be convex and differentiable on $\R^{m\times n}$, and run Shampoo ($s=1/4$, $\eps>0$) from $W_1$: $G_t=\nabla\ell_t(W_t)$, with $L_t$ and $R_t$ as in \cref{def:kronecker:shampoo}, and $W_{t+1}=W_t-\eta L_t^{-1/4}G_tR_t^{-1/4}$. Suppose that every $G_t$ has rank at most $r$. Then for every $W^\star$, with $D=\max_{t\le T}\norm{W_t-W^\star}_F$,
\begin{equation}\label{eq:kronecker:regret}
  \sum_{t=1}^T\bigl(\ell_t(W_t)-\ell_t(W^\star)\bigr)\ \le\ \Bigl(\frac{D^2}{2\eta}+r\eta\Bigr)\tr\bigl(L_T^{1/4}\bigr)\tr\bigl(R_T^{1/4}\bigr).
\end{equation}
\end{theorem}

\begin{proof}
The plan is to apply the regret bound of online gradient descent with changing preconditioners, and to compare Shampoo's preconditioner with full-matrix AdaGrad's by \cref{thm:kronecker:bound}. Write $x_t=\vect(W_t)$, $x^\star=\vect(W^\star)$ and $P_t=R_t^{1/4}\otimes L_t^{1/4}$, so that $x_{t+1}=x_t-\eta P_t^{-1}g_t$ and $\tr(P_t)=\tr(L_t^{1/4})\tr(R_t^{1/4})$. \Cref{thm:toolkit:ogd}(i), with the preconditioners $P_t/\eta$, the whole space as feasible set, and $P_0=0$, gives
\[
  \sum_{t=1}^T\bigl(\ell_t(W_t)-\ell_t(W^\star)\bigr)\le\frac1{2\eta}\sum_{t=1}^T(x_t-x^\star)^\top(P_t-P_{t-1})(x_t-x^\star)+\frac\eta2\sum_{t=1}^Tg_t^\top P_t^{-1}g_t .
\]

\emph{The distance terms.} The factors grow, so $L_t^{1/4}\succeq L_{t-1}^{1/4}$ and $R_t^{1/4}\succeq R_{t-1}^{1/4}$ by \cref{thm:toolkit:loewner}, and $P_t\succeq P_{t-1}$ by the monotonicity of Kronecker products. Since $v^\top Mv\le\norm{v}_2^2\tr M$ for $M\succeq0$ and $\norm{x_t-x^\star}_2\le D$, the first sum is at most $D^2\sum_t\tr(P_t-P_{t-1})=D^2\tr(P_T)$.

\emph{The gradient terms.} Put $\bar\Gamma_t=\eps I_{mn}+\frac1r\sum_{\tau\le t}g_\tau g_\tau^\top$. \Cref{thm:kronecker:bound} at time $t$ gives $\bar\Gamma_t\preceq R_t^{1/2}\otimes L_t^{1/2}=P_t^2$, hence $\bar\Gamma_t^{1/2}\preceq P_t$ and $P_t^{-1}\preceq\bar\Gamma_t^{-1/2}$ by \cref{thm:toolkit:loewner} with exponent $\frac12$. With $h_t=g_t/\sqrt r$ we have $\bar\Gamma_t=\bar\Gamma_{t-1}+h_th_t^\top$ and $\bar\Gamma_0=\eps I_{mn}\succ0$, so by the matrix form of \cref{lem:toolkit:adagrad}
\[
  \sum_{t=1}^Tg_t^\top P_t^{-1}g_t\le r\sum_{t=1}^Th_t^\top\bar\Gamma_t^{-1/2}h_t\le2r\tr\bigl(\bar\Gamma_T^{1/2}\bigr)\le2r\tr(P_T).
\]
The two bounds give \eqref{eq:kronecker:regret}.
\end{proof}

The bound holds for every step size, with $D$ measured along the run. If each step is followed by a projection in the norm $\norm{\cdot}_{P_t}$ onto a closed convex set that contains $W_1$ and $W^\star$ and has Frobenius diameter $D$, the same proof applies, because \cref{thm:toolkit:ogd}(i) allows the projection; then $D$ is known in advance, and $\eta=D/\sqrt{2r}$ gives $\sqrt{2r}\,D\tr(L_T^{1/4})\tr(R_T^{1/4})$, the form in which \textcite{GuptaKorenSinger2018} state the bound. Full-matrix AdaGrad's bound is $\sqrt2D\tr((\eps I+\Gamma_T)^{1/2})$; Shampoo's is no smaller, since $\sqrt{r}\tr(P_T)\ge\sqrt r\tr(\bar\Gamma_T^{1/2})=\tr((r\eps I+\Gamma_T)^{1/2})$, and that is the price of storing $m^2+n^2$ numbers instead of $(mn)^2$.

\section{The theory: the power and the basis}\label{sec:kronecker:curvature}

The regret bound is a guarantee, not a speed. To see what Kronecker preconditioners buy, we run them on the Kronecker quadratic, where everything can be computed.

\begin{proposition}[Kronecker preconditioners on Kronecker curvature]\label[proposition]{prop:kronecker:conditioning}
Under \cref{ass:kronecker:quadratic}:
\begin{enumerate}
\item $\nabla\Loss(W)=A(W-W_\star)C$, and the Hessian of $\Loss$ on $\vect(W)$ is $H=C\otimes A$, with eigenvalues $\alpha_i\gamma_j$ and condition number $\kappa(H)=\kappa_A\kappa_C$.
\item Let $P_L\in\R^{m\times m}$ and $P_R\in\R^{n\times n}$ be positive definite, and $W_{t+1}=W_t-\eta P_L^{-1}\nabla\Loss(W_t)P_R^{-1}$. The matrices $P_L^{-1}A$ and $P_R^{-1}C$ have positive eigenvalues; let $\mu_{\min}$ and $\mu_{\max}$ be the smallest and largest product of an eigenvalue of the first and one of the second, and $\kappa=\mu_{\max}/\mu_{\min}=\kappa(P_L^{-1}A)\kappa(P_R^{-1}C)$. With $\eta=2/(\mu_{\min}+\mu_{\max})$,
\begin{equation}\label{eq:kronecker:rate}
  \Loss(W_t)\le\Bigl(\frac{\kappa-1}{\kappa+1}\Bigr)^{2t}\Loss(W_0)\qquad\text{for every }t\ge0,
\end{equation}
with equality for some $W_0\ne W_\star$, and no constant step size gives a smaller factor for every $W_0$.
\item Gradient descent ($P_L=I$, $P_R=I$) has $\kappa=\kappa_A\kappa_C$; one-sided preconditioning ($P_L=A$, $P_R=I$) has $\kappa=\kappa_C$; K-FAC's two-sided step ($P_L=A$, $P_R=C$) has $\kappa=1$ and with $\eta=1$ reaches $W_\star$ in one step; and frozen Shampoo factors $P_L=L^s$, $P_R=R^s$ with $L=c_LA^q$, $R=c_RC^q$, for constants $c_L,c_R>0$ and $q\ge0$, have $\kappa=(\kappa_A\kappa_C)^{\abs{1-qs}}$.
\end{enumerate}
\end{proposition}

\begin{proof}
The plan is to read the Hessian off a second-order expansion and the rate off the eigenvalues of the symmetrized iteration. (i) With $E=W-W_\star$ and symmetric $A$ and $C$, $\Loss(W+\Delta)=\Loss(W)+\ip{AEC,\Delta}+\frac12\tr(\Delta^\top A\Delta C)$, and by the vec identity $\tr(\Delta^\top A\Delta C)=\vect(\Delta)^\top(C\otimes A)\vect(\Delta)$.

(ii) With $e_t=\vect(W_t-W_\star)$ and $P=P_R\otimes P_L$, $e_{t+1}=(I-\eta P^{-1}H)e_t$. Put $z_t=P^{1/2}e_t$; then $z_{t+1}=(I-\eta S)z_t$ with
\[
  S=P^{-1/2}HP^{-1/2}=\bigl(P_R^{-1/2}CP_R^{-1/2}\bigr)\otimes\bigl(P_L^{-1/2}AP_L^{-1/2}\bigr),
\]
symmetric positive definite, with eigenvalues the products of those of the two factors, which are those of the similar matrices $P_R^{-1}C$ and $P_L^{-1}A$. The loss is $\frac12e_t^\top He_t=\frac12z_t^\top Sz_t$. In an orthonormal eigenbasis $(b_k)$ of $S$, with eigenvalues $\mu_k$ and $z_0=\sum_k\zeta_kb_k$,
\[
  \Loss(W_t)=\tfrac12\sum_k\mu_k(1-\eta\mu_k)^{2t}\zeta_k^2\le\max_k(1-\eta\mu_k)^{2t}\,\Loss(W_0),
\]
with equality when $z_0$ is an eigenvector for an eigenvalue attaining the maximum. The function $\eta\mapsto\max\{\abs{1-\eta\mu_{\min}},\abs{1-\eta\mu_{\max}}\}$ is smallest where the two terms agree, at $\eta=2/(\mu_{\min}+\mu_{\max})$, where it is $(\kappa-1)/(\kappa+1)$.

(iii) For K-FAC, $P^{-1}H=I$. For frozen Shampoo factors, $P_L^{-1}A=c_L^{-s}A^{1-qs}$ has the eigenvalues $c_L^{-s}\alpha_i^{1-qs}$ and the condition number $\kappa_A^{\abs{1-qs}}$, and likewise $\kappa(P_R^{-1}C)=\kappa_C^{\abs{1-qs}}$.
\end{proof}

Shampoo's factors are powers of $A$ and $C$ when they average gradients whose second moment is a power of the Hessian. Two cases matter. Gradients taken at errors $E=W-W_\star$ that are isotropic random matrices, $\E[\vect(E)\vect(E)^\top]=\sigma^2I$, as after a random initialization, have $g=H\vect(E)$ and $\E[gg^\top]=\sigma^2H^2$. Gradients dominated by noise with covariance proportional to the Hessian, as near the minimum of a model whose Fisher matrix equals its Hessian (\cref{prop:kronecker:kfac}), have $\E[gg^\top]=\sigma^2H$.

\begin{corollary}[Which power]\label[corollary]{cor:kronecker:power}
Under \cref{ass:kronecker:quadratic}, let Shampoo's factors be the second moments $L=\E[GG^\top]$ and $R=\E[G^\top G]$ of random gradients with $\E[gg^\top]=\sigma^2H^q$, for some $\sigma>0$ and $q\ge0$. Then the frozen step $W\leftarrow W-\eta L^{-s}\nabla\Loss(W)R^{-s}$ has, at its best constant step size, the rate \eqref{eq:kronecker:rate} with
\begin{equation}\label{eq:kronecker:exponent}
  \kappa=\kappa(H)^{\abs{1-qs}} .
\end{equation}
\end{corollary}

\begin{proof}
Since $H^q=C^q\otimes A^q$, \cref{prop:kronecker:partial} gives $L=\sigma^2\tr(C^q)A^q$ and $R=\sigma^2\tr(A^q)C^q$, and \cref{prop:kronecker:conditioning}(iii) applies.
\end{proof}

The exponent $\abs{1-qs}$ is a V-shaped law\index{effective condition number} in the power:
\[
  \begin{array}{l|cccc}
    & s=0\ \text{(GD)} & s=\frac14\ \text{(Shampoo)} & s=\frac12 & s=1\ \text{(K-FAC)}\\\hline
    q=2 & 1 & \frac12 & 0 & 1\\
    q=1 & 1 & \frac34 & \frac12 & 0
  \end{array}
\]
In neither case does Shampoo's power remove the conditioning: it leaves $\kappa(H)^{1/2}$ for isotropic errors and $\kappa(H)^{3/4}$ for Hessian-shaped noise. The power $1/2$ is exact for $q=2$, where it is full-matrix AdaGrad (\cref{prop:kronecker:partial}). The inverse, K-FAC's power, is exact for $q=1$, where the factors estimate curvature, and for $q=2$ no better than gradient descent: then $L^{-1}A\propto A^{-1}$, which inverts the spectrum instead of flattening it. This is one precise form of the warning that the second moment of the gradients is not the curvature \parencite{KunstnerHennigBalles2019}: the inverse is right for the one, its square root for the other.\sidenote{The same arithmetic says why AdaGrad and Adam take square roots: on a diagonal quadratic with isotropic errors, the second moment of the $i$th gradient coordinate is proportional to $h_i^2$, and only its square root, not its inverse, flattens the curvatures $h_i$.}

\begin{proposition}[K-FAC's approximation]\label[proposition]{prop:kronecker:kfac}
Let a layer with weights $W\in\R^{m\times n}$ receive the input $a\in\R^n$ and the back-propagated output gradient $\delta\in\R^m$, so that its per-example gradient is $G=\delta a^\top$, with $(a,\delta)$ random with finite fourth moments.
\begin{enumerate}
\item $\vect(G)\vect(G)^\top=(aa^\top)\otimes(\delta\delta^\top)$, so the layer's block of the Fisher matrix (with targets sampled from the model) or of the empirical Fisher matrix\index{Fisher matrix!empirical} (with the data's targets) is $F=\E[(aa^\top)\otimes(\delta\delta^\top)]$.
\item If $a$ and $\delta$ are independent, then $F=\E[aa^\top]\otimes\E[\delta\delta^\top]$, so K-FAC is exact, and Shampoo's per-example factors $\E[GG^\top]=\E\norm{a}_2^2\,\E[\delta\delta^\top]$ and $\E[G^\top G]=\E\norm{\delta}_2^2\,\E[aa^\top]$ are multiples of K-FAC's.
\item For the mean $G_B$ of $B$ independent copies of $G$, $\E[G_BG_B^\top]=\frac1B\E[\norm{a}_2^2\,\delta\delta^\top]+(1-\frac1B)\bar G\bar G^\top$ with $\bar G=\E[\delta a^\top]$, and similarly for $G_B^\top G_B$.
\end{enumerate}
\end{proposition}

\begin{proof}
(i) $\vect(\delta a^\top)=a\otimes\delta$, and $(a\otimes\delta)(a\otimes\delta)^\top=(aa^\top)\otimes(\delta\delta^\top)$ by the mixed product rule. (ii) The entries $a_ja_l\delta_i\delta_k$ of $(aa^\top)\otimes(\delta\delta^\top)$ have factoring expectations, and $GG^\top=\norm{a}_2^2\delta\delta^\top$, $G^\top G=\norm{\delta}_2^2aa^\top$. (iii) Of the terms of $G_BG_B^\top=B^{-2}\sum_{k,k'}\delta_ka_k^\top a_{k'}\delta_{k'}^\top$, the $B$ with $k=k'$ have expectation $\E[\norm{a}_2^2\delta\delta^\top]$ and the others $\bar G\bar G^\top$.
\end{proof}

For example, let $a\sim\cN(0,C)$ and let the model predict $y\sim\cN(Wa,A^{-1})$, with negative log-likelihood $\frac12(y-Wa)^\top A(y-Wa)$; for data generated with weights $W_\star$ the expected loss is the Kronecker quadratic plus a constant. With targets sampled from the model, $\delta=-A\xi$ with $\xi\sim\cN(0,A^{-1})$ independent of $a$, so $\E[\delta\delta^\top]=A$ and $F=C\otimes A=H$ at every $W$: K-FAC is Newton's method and converges in one step. At the minimum the data's targets give the same $\delta$, so the gradients that Shampoo accumulates there have the per-example factors $\tr(C)A$ and $\tr(A)C$, the statistics with $q=1$, and a minibatch's factors are these divided by $B$. Away from the minimum the data's $\delta$ depends on $a$, and by \cref{prop:kronecker:kfac}(iii) the mean gradient dominates a minibatch's factors far from the minimum and the per-example noise near it; the right power drifts with them. Three statements remain.

\begin{proposition}[Shampoo between Muon and AdaGrad]\label[proposition]{prop:kronecker:between}
\begin{enumerate}
\item For a single gradient and $\eps=0$, with pseudo-inverse powers, $(GG^\top)^{-1/4}G(G^\top G)^{-1/4}=\polar(G)$.
\item If $G_\tau=U\diag(\sigma_\tau)V^\top$ for $\tau\le t$, where $U\in\R^{m\times k}$ and $V\in\R^{n\times k}$ have orthonormal columns and $\sigma_\tau\in\R^k$ has entries of either sign, then
\[
  L_t^{-1/4}G_tR_t^{-1/4}=U\diag\biggl(\frac{\sigma_{t,i}}{\bigl(\eps+\sum_{\tau\le t}\sigma_{\tau,i}^2\bigr)^{1/2}}\biggr)V^\top ,
\]
diagonal AdaGrad on the coefficients. With exponential moving averages (coefficient $\beta_2$) in place of the sums it is RMSProp on them, and for $\beta_2=0$ and $\eps\to0$ it tends to $U\diag(\sign\sigma_t)V^\top=\polar(G_t)$.
\end{enumerate}
\end{proposition}

\begin{proof}
(i) is \cref{prop:spectral:shampoo}: in a compact singular value decomposition $G=U\Sigma V^\top$ the powers act on $\Sigma$, and $\Sigma^{-1/2}\Sigma\Sigma^{-1/2}=I$. (ii) By orthonormality, $L_t=\eps I+U\diag(\varsigma_t)U^\top$ with $\varsigma_{t,i}=\sum_{\tau\le t}\sigma_{\tau,i}^2$, so $L_t^{-1/4}U=U\diag(\eps+\varsigma_t)^{-1/4}$, and likewise $R_t^{-1/4}V=V\diag(\eps+\varsigma_t)^{-1/4}$; multiplying out gives the formula. With moving averages $\varsigma_{t,i}$ becomes the average of $\sigma_{\tau,i}^2$, and for $\beta_2=0$ the $i$th coefficient is $\sigma_{t,i}/(\eps+\sigma_{t,i}^2)^{1/2}\to\sign\sigma_{t,i}$.
\end{proof}

Without memory Shampoo is spectral steepest descent; with memory, in a fixed singular basis, it is AdaGrad or RMSProp on the coefficients in that basis, as Adam is RMSProp with momentum on the coordinates (\cref{ch:sign,ch:spectral}; \textcite{BernsteinNewhouse2024}). Part (ii) also shows when the power $1/4$ is right: if moreover $\sum_{\tau\le t}\sigma_{\tau,i}\sigma_{\tau,j}=0$ for $i\ne j$, then $\Gamma_t=\sum_i\varsigma_{t,i}w_iw_i^\top$ with $w_i=v_i\otimes u_i$, and Shampoo's step is exactly full-matrix AdaGrad's. For Kronecker statistics the right power is $1/2$ (\cref{prop:kronecker:partial}); which structure a network's gradients have is an empirical question (Exercise~\ref{ex:kronecker:kron}).

\begin{proposition}[Shampoo with the power $1/2$ is Adafactor in its eigenbasis; \textcite{VyasEtAl2025}]\label[proposition]{prop:kronecker:adafactor}
Let $L_t=\sum_{\tau\le t}G_\tau G_\tau^\top=Q_L\diag(\lambda)Q_L^\top$ and $R_t=\sum_{\tau\le t}G_\tau^\top G_\tau=Q_R\diag(\mu)Q_R^\top$ be positive definite, put $G'_\tau=Q_L^\top G_\tau Q_R$, and let $V=\sum_{\tau\le t}G'_\tau\odot G'_\tau$ (entrywise squares). Adafactor's rank-one estimate $\hat V=(V\mathbf1)(\mathbf1^\top V)/(\mathbf1^\top V\mathbf1)$ of $V$ (\cref{ch:adaptive}) has entries $\hat V_{ij}=\lambda_i\mu_j/\tr(L_t)$, and Adafactor's step, rotated back, is (with $\oslash$ and $\hat V^{\odot1/2}$ the entrywise quotient and square root)
\[
  Q_L\bigl(G'_t\oslash\hat V^{\odot1/2}\bigr)Q_R^\top=\sqrt{\tr(L_t)}\;L_t^{-1/2}G_tR_t^{-1/2}.
\]
\end{proposition}

\begin{proof}
The row sums of $V$ are $\sum_\tau(G'_\tau G_\tau'^\top)_{ii}=(Q_L^\top L_tQ_L)_{ii}=\lambda_i$, the column sums are $(Q_R^\top R_tQ_R)_{jj}=\mu_j$, and the total is $\tr(L_t)$. So the rotated step is $\sqrt{\tr(L_t)}\diag(\lambda)^{-1/2}G'_t\diag(\mu)^{-1/2}$, and rotating back gives $\sqrt{\tr(L_t)}\,Q_L\diag(\lambda)^{-1/2}Q_L^\top G_tQ_R\diag(\mu)^{-1/2}Q_R^\top$.
\end{proof}

The same holds with moving averages. \term{SOAP}\index{SOAP} keeps the eigenbasis and replaces Adafactor's rank-one $\hat V$ by Adam's full matrix of second moments. On Kronecker curvature this removes the conditioning, because Adam with $\epsilon=0$ does not see a positive rescaling of a coordinate's gradients.\sidenote{\Textcite{KingmaBa2015} list invariance to diagonal rescaling of the gradients among Adam's properties. It is exact only for $\epsilon=0$; for coordinates with gradients much smaller than $\epsilon$, Adam is momentum gradient descent with step size $\eta/\epsilon$ (\cref{ch:adaptive}).}

\begin{proposition}[Adam in the Kronecker eigenbasis]\label[proposition]{prop:kronecker:soap}
Under \cref{ass:kronecker:quadratic}, let $A=U_A\diag(\alpha)U_A^\top$ and $C=U_C\diag(\gamma)U_C^\top$, and run Adam with $\epsilon=0$ (with or without bias correction, reading $0/0$ as $0$), or AdaGrad, on the rotated gradients $U_A^\top\nabla\Loss(W_t)U_C$, rotating each update back. Then the rotated errors $\tilde E_t=U_A^\top(W_t-W_\star)U_C$ do not depend on $\alpha$ and $\gamma$; in particular they are those of the same method on the isotropic loss $\frac12\norm{W-W_\star}_F^2$. If Shampoo's factors are $f(A)$ and $h(C)$ for injective $f,h$ on $(0,\infty)$, as in \cref{cor:kronecker:power} with $q>0$, then every orthonormal eigenbasis of the factors is one of $A$ and $C$, and SOAP with these frozen factors and $\epsilon=0$ is this method.
\end{proposition}

\begin{proof}
The plan is an induction on $t$, based on one invariance. In the rotated coordinates the gradient is $\diag(\alpha)\tilde E\diag(\gamma)$, with entries $\alpha_i\gamma_j\tilde E_{ij}$. If every gradient of one coordinate is multiplied by $c>0$, then Adam's moments become $cm$ and $c^2v$, and its update, a function of $m/\sqrt v$, does not change; nor does AdaGrad's $g/(\sum g^2)^{1/2}$. (A coordinate whose gradients have all been zero has $m=v=0$ in both problems and does not move.) So if the rotated errors of the problem with curvatures $\alpha_i\gamma_j$ and of the isotropic problem agree up to step $t$, then their gradients in coordinate $(i,j)$ differ by the factor $\alpha_i\gamma_j$ at every step so far, their updates agree, and so do their errors at step $t+1$. For the last claim, $f(A)=U_A\diag(f(\alpha))U_A^\top$ and, $f$ being injective, $f(A)$ and $A$ have the same eigenspaces; likewise $h(C)$ and $C$. In a fixed basis, SOAP's moving averages of rotated gradients are rotated moving averages, so SOAP is Adam in that basis.
\end{proof}

Adam is curvature-free on every diagonal quadratic; it is the rotation of the problem that hurts it (\cref{ch:sign,ch:transformers}), and on Kronecker curvature with stationary statistics SOAP undoes the rotation. Curvature-free does not mean fast: each mode moves as it would on the isotropic quadratic, at the method's pace there. And the basis must be right, which with estimated factors it is only approximately (\cref{fig:kronecker:ladder}).

\section{Reading the theory}\label{sec:kronecker:reading}

\paragraph{The ladder.} Take $\kappa_A=\kappa_C=30$, so $\kappa(H)=900$. Once the slowest mode dominates, the best constant step reduces the loss tenfold in about $\ln10/(2\ln\frac{\kappa+1}{\kappa-1})\approx0.58\kappa$ steps: $520$ for gradient descent, $95$ for Shampoo with Hessian-shaped statistics ($\kappa=900^{3/4}\approx164$), $17$ for Shampoo with isotropic-error statistics or the power $\frac12$ with Hessian-shaped ones ($\kappa=30$), and one for K-FAC with exact factors. In \cref{fig:kronecker:ladder} (left), factors estimated from $1000$ sampled gradients give these conditionings to within one percent ($164.7$ and $29.7$), and the losses follow the exact formulas over seventeen decades. AdaGrad in Shampoo's eigenbasis, the idealized SOAP of \cref{prop:kronecker:soap}, beats Shampoo with Hessian-shaped statistics, and with step size $0.5$ instead of $0.3$ also the other ($7\times10^{-22}$ against $2\times10^{-17}$ at $t=250$); in the original coordinates it gains only a decade on gradient descent. In the eigenbasis of the estimated factors, off by about a tenth of a radian, SOAP stays within a factor of three of its prediction over ten decades; how a basis error enters the rate, \cref{prop:kronecker:soap} does not say.

\begin{figure}[t]
\begin{fullwidth}
\centering
\begin{tikzpicture}
\begin{groupplot}[group style={group size=2 by 1,horizontal sep=17mm},wplot,width=82mm,height=50mm]
\nextgroupplot[xlabel={step $t$},ylabel={$\Loss(W_t)/\Loss(W_0)$},ymode=log,xmin=0,xmax=250,ymin=1e-18,ymax=4,
  ytick={1,1e-4,1e-8,1e-12,1e-16},
  legend style={at={(0.5,1.03)},anchor=south,font=\scriptsize},legend columns=3]
\addplot[muted,line width=1.0pt] table[x=t,y=gd_th] {figures/data/07-kronecker-ladder.dat};
\addplot[sandline!55!black,line width=1.1pt] table[x=t,y=shq1_th] {figures/data/07-kronecker-ladder.dat};
\addplot[sea,line width=1.1pt] table[x=t,y=shq2_th] {figures/data/07-kronecker-ladder.dat};
\addplot[seadeep,line width=1.0pt,dashed] table[x=t,y=soap_th] {figures/data/07-kronecker-ladder.dat};
\addplot[muted,line width=0.7pt,densely dotted] table[x=t,y=ada_sim] {figures/data/07-kronecker-ladder.dat};
\addplot[only marks,mark=*,mark size=1.3pt,color=muted,mark options={fill=muted,draw=white,line width=0.2pt},each nth point=10,filter discard warning=false,forget plot] table[x=t,y=gd_sim] {figures/data/07-kronecker-ladder.dat};
\addplot[only marks,mark=square*,mark size=1.2pt,color=sandline!55!black,mark options={fill=sandline!55!black,draw=white,line width=0.2pt},each nth point=10,filter discard warning=false,forget plot] table[x=t,y=shq1_sim] {figures/data/07-kronecker-ladder.dat};
\addplot[only marks,mark=*,mark size=1.4pt,color=sea,mark options={fill=sea,draw=white,line width=0.2pt},each nth point=10,filter discard warning=false,forget plot] table[x=t,y=shq2_sim] {figures/data/07-kronecker-ladder.dat};
\addplot[only marks,mark=triangle*,mark size=1.6pt,color=seadeep,mark options={fill=seadeep,draw=white,line width=0.2pt},each nth point=10,filter discard warning=false,forget plot] table[x=t,y=soap_sim] {figures/data/07-kronecker-ladder.dat};
\legend{gradient descent,{Shampoo, $q=1$},{Shampoo, $q=2$},{AdaGrad, eigenbasis},{AdaGrad, original basis}}
\nextgroupplot[xlabel={power $s$ of the factors},ylabel={exponent $\log\kappa/\log\kappa(H)$},xmin=0,xmax=1,ymin=0,ymax=1.08,
  xtick={0,0.25,0.5,0.75,1},xticklabels={$0$,$\frac14$,$\frac12$,$\frac34$,$1$},
  legend style={at={(0.5,1.03)},anchor=south,font=\scriptsize},legend columns=2]
\addplot[guide,forget plot] coordinates {(0.25,0) (0.25,1.08)};
\addplot[guide,forget plot] coordinates {(0.5,0) (0.5,1.08)};
\addplot[sandline!55!black,line width=1.1pt] table[x=s,y=q1_th] {figures/data/07-kronecker-power.dat};
\addplot[sea,line width=1.1pt] table[x=s,y=q2_th] {figures/data/07-kronecker-power.dat};
\addplot[only marks,mark=square*,mark size=1.6pt,color=sandline!55!black,mark options={fill=sandline!55!black,draw=white,line width=0.2pt},forget plot] table[x=s,y=q1_sim] {figures/data/07-kronecker-power.dat};
\addplot[only marks,mark=*,mark size=1.7pt,color=sea,mark options={fill=sea,draw=white,line width=0.2pt},forget plot] table[x=s,y=q2_sim] {figures/data/07-kronecker-power.dat};
\node[font=\scriptsize,text=muted,anchor=south west,rotate=90] at (axis cs:0.24,0.02) {Shampoo};
\node[font=\scriptsize,text=muted,anchor=south west,rotate=90] at (axis cs:0.99,0.24) {K-FAC};
\legend{{$q=1$: $\abs{1-s}$},{$q=2$: $\abs{1-2s}$}}
\end{groupplot}
\end{tikzpicture}
\end{fullwidth}
\caption{Kronecker preconditioners on the Kronecker quadratic with $m=16$, $n=24$, $\kappa_A=\kappa_C=30$ and Haar-random eigenvectors; Shampoo's factors are estimated from $1000$ sampled gradients with statistics $\E[gg^\top]\propto H^q$. Left: relative loss from one random start; curves are the exact formulas of \cref{prop:kronecker:conditioning,prop:kronecker:soap} for exact factors, marks (every tenth step) the matrix iterations with estimated factors, at the best constant step size. AdaGrad has step size $0.3$ in Shampoo's eigenbasis (exact basis for the curve, the estimated $q=1$ basis for the marks) and $1.0$ in the original coordinates (dotted, simulation only). Right: the exponent of \eqref{eq:kronecker:exponent}; lines $\abs{1-qs}$, marks from the asymptotic rate with estimated factors. Script: \texttt{07-kronecker-ladder.py}.}
\label{fig:kronecker:ladder}
\end{figure}

\paragraph{Which power.} The right panel measures \cref{cor:kronecker:power}: for each power $s$ the iteration runs with the estimated factors at its best step size, and its asymptotic rate $\rho$ gives $\kappa=(1+\rho)/(1-\rho)$. The marks follow $\abs{1-qs}$ to within about $0.01$, except at the bottoms of the V, where the sampling error of the factors, a few percent per eigenvalue, leaves the exponents $0.03$ and $0.04$ ($\kappa\approx1.2$ and $1.35$), a floor that shrinks like the inverse square root of the number of gradients (Exercise~\ref{ex:kronecker:simV}). Three lessons follow. The power matters as much as preconditioning: an inverse that over-corrects can be as bad as none. Shampoo's $1/4$ never hurts but never finishes the job. And the right power is set by what the factors estimate, which is why K-FAC inverts curvature estimates while AdaGrad-type methods take square roots of second moments of gradients.

\paragraph{How loose the constant is.} In \cref{fig:kronecker:constant} a single gradient of rank $4$ attains $c^\star=r$, and a coherent sequence, all multiples of one matrix, stays there. In a shared singular basis $c^\star$ falls toward the top eigenvalue of the limiting cosine matrix: $1+\frac34(r-1)=3.25$ for independent coefficients uniform on $(0,1)$, whose cosines tend to $\E[s]^2/\E[s^2]=\frac34$, and $1$ for sign-symmetric ones. Generic gradients $X_tY_t^\top$, with Gaussian factors, have statistics tending to $rI_n\otimes I_m$ and a constant falling toward $(mn)^{-1/2}$. The factor $\sqrt r$ in \cref{thm:kronecker:regret} is a worst case.

\begin{figure}[t]
\centering
\begin{tikzpicture}
\begin{axis}[wplot,width=\linewidth,height=0.33\linewidth,xmode=log,ymode=log,log basis x=2,xmin=0.8,xmax=5200,ymin=0.04,ymax=6,
  xlabel={number of gradients $T$},ylabel={tight constant $c^\star$},
  xtick={1,4,16,64,256,1024,4096},xticklabels={$1$,$4$,$16$,$64$,$256$,$1024$,$4096$},
  ytick={0.0625,0.25,1,4},yticklabels={$\frac1{16}$,$\frac14$,$1$,$4$},
  legend style={at={(0.5,1.03)},anchor=south,font=\scriptsize},legend columns=4]
\addplot[guide,forget plot] table[x=T,y=th_coherent] {figures/data/07-kronecker-constant.dat};
\addplot[guide,forget plot] table[x=T,y=th_positive] {figures/data/07-kronecker-constant.dat};
\addplot[guide,forget plot] table[x=T,y=th_signed] {figures/data/07-kronecker-constant.dat};
\addplot[guide,forget plot] table[x=T,y=th_generic] {figures/data/07-kronecker-constant.dat};
\addplot[only marks,mark=diamond*,mark size=2.0pt,color=muted,mark options={fill=muted,draw=white,line width=0.2pt}] table[x=T,y=coherent] {figures/data/07-kronecker-constant.dat};
\addplot[only marks,mark=square*,mark size=1.6pt,color=sandline!55!black,mark options={fill=sandline!55!black,draw=white,line width=0.2pt}] table[x=T,y=positive] {figures/data/07-kronecker-constant.dat};
\addplot[only marks,mark=*,mark size=1.8pt,color=sea,mark options={fill=sea,draw=white,line width=0.2pt}] table[x=T,y=signed] {figures/data/07-kronecker-constant.dat};
\addplot[only marks,mark=triangle*,mark size=2.0pt,color=seadeep,mark options={fill=seadeep,draw=white,line width=0.2pt}] table[x=T,y=generic] {figures/data/07-kronecker-constant.dat};
\legend{coherent,{shared basis, positive},{shared basis, signed},generic}
\end{axis}
\end{tikzpicture}
\caption{The tight constant $c^\star$ for $T$ gradients of rank $4$ in $\R^{16\times16}$ (marks: mean of five seeds, computed from the $256\times256$ matrices) and its limits from \cref{prop:kronecker:constant} (dashed): $4$, $3.25$, $1$ and $\frac1{16}$. Script: \texttt{07-kronecker-constant.py}.}
\label{fig:kronecker:constant}
\end{figure}

\paragraph{Costs and the regret bound.} Shampoo stores $m^2+n^2$ numbers for an $m\times n$ matrix, against $(mn)^2$; the roots cost of order $m^3+n^3$ operations, as much as a forward and backward pass of a $4096\times4096$ layer on some thousands of tokens, which is why practice computes them every so many steps. The regret bound cannot explain an advantage. Since $\tr(M^{1/4})\ge(\tr M)^{1/4}$ for $M\succeq0$, $\tr(L_T^{1/4})\tr(R_T^{1/4})\ge(\tr L_T\tr R_T)^{1/4}\ge(\sum_t\norm{G_t}_F^2)^{1/2}$, so the right side of \eqref{eq:kronecker:regret} is at least $\sqrt{2r}D(\sum_t\norm{G_t}_F^2)^{1/2}$, while online gradient descent with steps $\eta_t=D/(2\sum_{\tau\le t}\norm{G_\tau}_F^2)^{1/2}$ has regret at most $\sqrt2D(\sum_t\norm{G_t}_F^2)^{1/2}$ (\cref{thm:toolkit:ogd,lem:toolkit:adagrad}). With the same Frobenius distance, Shampoo's guarantee is never better than gradient descent's, though it can beat diagonal AdaGrad's by up to $\sqrt{mn}$ (Exercise~\ref{ex:kronecker:onesided}).

\section{What the theory misses}\label{sec:kronecker:misses}

\paragraph{What was idealized away.} The Kronecker structure itself, exact for one linear layer whose inputs and back-propagated gradients are independent (\cref{prop:kronecker:kfac}); in a network the gradient flowing back into a layer depends on its output and hence on its input, minibatches add the mean term of \cref{prop:kronecker:kfac}(iii), and layers are coupled (Exercise~\ref{ex:kronecker:twolayer}). Stationarity: the factors average over a trajectory along which the statistics drift from error-dominated to noise-dominated, so \cref{cor:kronecker:power} describes a frozen snapshot, and plain Shampoo on a deterministic quadratic never becomes stationary (Exercise~\ref{ex:kronecker:transient}). Convex online losses, in the regret theory. Exact roots at every step, where practice recomputes them infrequently and in finite precision, and the inverse fourth root of an ill-conditioned factor depends strongly on $\eps$. Grafting, momentum, weight decay, and the parameters that are not matrices.

\paragraph{Where it breaks.} (1) The regret bound does not rank optimizers: it is never better than that of gradient descent with an adaptive step (\cref{sec:kronecker:reading}), so it does not explain \cref{obs:kronecker:algoperf,obs:kronecker:scale}. (2) The theorem needs the power $1/4$, implementations also use $1/2$, and \cref{cor:kronecker:power} shows that no power is right for all statistics; \textcite{MorwaniEtAl2025} find the square of Shampoo's preconditioner the better approximation of full-matrix AdaGrad in networks. (3) Performance rests on heuristics the theory does not cover (\cref{obs:kronecker:heuristics}). (4) For K-FAC the curvature story is incomplete: the empirical Fisher can be a poor curvature estimate \parencite{KunstnerHennigBalles2019}, and \textcite{Benzing2022} found that a variant keeping only the input-side factor, a one-sided preconditioner in the sense of \cref{prop:kronecker:conditioning}, performs as well as K-FAC or better, which suggests that K-FAC's advantage is not mainly second-order. (5) The advantage of matrix-based optimizers over a tuned AdamW shrinks with scale \parencite{WenEtAl2026Fantastic}, which a one-layer idealization cannot address. (6) SOAP is curvature-free only in the exact eigenbasis; the theory does not say how a basis error enters the rate, nor how often to refresh the basis, which \cref{obs:kronecker:heuristics} says matters.

\paragraph{Repairs and open problems.} Correcting the eigenvalues in the Kronecker eigenbasis repairs the mismatch between the factors and the true scales: eigenvalue-corrected K-FAC \parencite{GeorgeEtAl2018}, SOAP \parencite{VyasEtAl2025}, and the decoupled updates of eigenvalues and eigenbasis studied by \textcite{EschenhagenEtAl2025}. Better Kronecker approximations iterate the power method of \cref{prop:kronecker:partial} \parencite{VanLoanPitsianis1993,MorwaniEtAl2025}. Open: how close trained networks' gradient statistics are to Kronecker products (Exercise~\ref{ex:kronecker:kron}); a theory of the power, or of a schedule for it, under drifting statistics (Exercise~\ref{ex:kronecker:schedule}); a nonconvex theory in which Shampoo or SOAP provably beats Adam; and why the advantage of matrix-based optimizers shrinks with scale.

\section{Back in the wild}\label{sec:kronecker:back}

\fieldnotes{Precondition each weight matrix by inverse roots of row and column second-moment matrices (Shampoo), by inverses of input and output-gradient covariances (K-FAC), or by Adam in the eigenbasis of those matrices (SOAP).}
{Full-matrix AdaGrad on $\vect(W)$; the Kronecker quadratic with Hessian $C\otimes A$; stationary statistics $\E[gg^\top]\propto H^q$.}
{The Kronecker bound, with constant $r$ attained by single gradients; Shampoo's regret; $\kappa\mapsto\kappa^{\abs{1-qs}}$ for frozen factors; K-FAC exact under independence; Shampoo with power $\frac12$ is Adafactor in its eigenbasis; Adam in that basis is curvature-free.}
{Shampoo's power $\frac14$ leaves $\kappa^{1/2}$ or $\kappa^{3/4}$; the power $\frac12$ is full-matrix AdaGrad for Kronecker statistics; inverses are Newton for curvature-like statistics and over-correct for gradient-like ones.}
{Why these methods beat a tuned Adam when they do; grafting, $\eps$ and stale roots; non-Kronecker and drifting statistics; shrinking gains at scale.}
{Match the power to what the factors estimate: inverses for curvature, inverse square roots for second moments of gradients (proved in the idealization); keep SOAP's eigenbasis fresh (observed); expect the gain over a tuned AdamW to shrink with scale (observed).}

The theory tells a practitioner three things: a Kronecker preconditioner genuinely approximates full-matrix AdaGrad; the power of the factors is the decisive choice and should match what they estimate, so the theory's $1/4$ is chosen for the proof, not for speed, consistent with the use of $1/2$ with grafting and of inverses on curvature estimates; and SOAP's advantage comes from the basis, in which Adam is curvature-free. It does not say when this beats a well-tuned AdamW; the benchmarks say that it sometimes does, by margins that shrink with scale, for reasons not understood.

The chapter closes the unification line of Parts~I and~II. Without memory Shampoo is the orthogonalized gradient of \cref{ch:spectral}, as Adam is the sign descent of \cref{ch:sign}; with memory both become preconditioners, diagonal (\cref{ch:adaptive}) or Kronecker-factored, and SOAP runs the first in the second's basis. The dynamics line takes over with preconditioned sharpness (\cref{ch:eos}), the critical batch size (\cref{ch:noise}), and the gap between Adam and gradient descent (\cref{ch:transformers}).

\begin{exercises}
\item \consolidation (Vectorization.) (a) Show that $\vect(uv^\top)=v\otimes u$ and $\ip{X,AYC}=\vect(X)^\top(C\otimes A)\vect(Y)$ for symmetric $C$, and that $(C\otimes A)(v_j\otimes u_i)=\gamma_j\alpha_i(v_j\otimes u_i)$ for eigenvectors $Au_i=\alpha_iu_i$, $Cv_j=\gamma_jv_j$. (b) With the row-stacking $\overline{\vect}(X)=\vect(X^\top)$, show that $\overline{\vect}(MXN)=(M\otimes N^\top)\overline{\vect}(X)$, so that every Kronecker product of this chapter reverses its order; write \eqref{eq:kronecker:bound} in that convention.\label{ex:kronecker:vec}

\item \consolidation (Geometric means.) (a) For commuting positive semidefinite $P,Q$, show that $\bigl(\begin{smallmatrix}P&(PQ)^{1/2}\\(PQ)^{1/2}&Q\end{smallmatrix}\bigr)\succeq0$. (b) For positive definite $P,Q$, show with a Schur complement that $\bigl(\begin{smallmatrix}P&Y\\Y&Q\end{smallmatrix}\bigr)\succeq0$ if and only if $YP^{-1}Y\preceq Q$, and deduce, using \cref{thm:toolkit:loewner} and $\abs{Z}\succeq Z$ for symmetric $Z$, that for commuting $P,Q$ every such $Y$ satisfies $Y\preceq(PQ)^{1/2}$.\label{ex:kronecker:geomean}

\item \consolidation (The constant.) Let $u_1,u_2$ and $v_1,v_2$ be orthonormal pairs. With \cref{prop:kronecker:constant}, compute $c^\star$ for (a) the single gradient $u_1v_1^\top+u_2v_2^\top$; (b) the gradients $u_1v_1^\top+u_2v_2^\top$ and $u_1v_1^\top-u_2v_2^\top$; (c) the gradients $u_1v_1^\top+u_2v_2^\top$ and $u_1v_1^\top$. (d) With the rearrangement $\cR$ defined by $\cR(\Phi\otimes\Psi)=\vect(\Phi)\vect(\Psi)^\top$ and extended linearly, show that $\cR(\Gamma)\vect(I_m)=\vect(R(\Gamma))$ and $\cR(\Gamma)^\top\vect(I_n)=\vect(L(\Gamma))$: one step of the power method for the nearest Kronecker product, from identities, returns the partial traces. \hint{(c) gives $1+1/\sqrt2$.}\label{ex:kronecker:constant}

\item \consolidation (K-FAC for a two-layer linear network.) Let $f(x)=W_2W_1x$ with $W_1\in\R^{d_1\times d_0}$, $W_2\in\R^{d_2\times d_1}$, $x\sim\cN(0,\Sigma)$, and the model $y\sim\cN(f(x),I)$ with sampled targets. (a) Show that the Fisher matrix has diagonal blocks $F_{11}=\Sigma\otimes W_2^\top W_2$ and $F_{22}=W_1\Sigma W_1^\top\otimes I_{d_2}$, on which K-FAC is exact. (b) Show that $F_{12}=\Sigma W_1^\top\otimes W_2^\top$, which K-FAC drops. (c) Show that the full Fisher matrix is singular, because $(W_1,W_2)\mapsto(MW_1,W_2M^{-1})$ leaves $f$ unchanged.\label{ex:kronecker:twolayer}

\item \variation (One-sided Shampoo.) Let $G$ have rank $r$ and $g=\vect(G)$. (a) Show that $gg^\top\preceq r(I_n\otimes GG^\top)$ and $gg^\top\preceq r(G^\top G\otimes I_m)$, with sharp constants. \hint{$\ip{G,X}=\ip{VV^\top,G^\top X}$ for $G=U\Sigma V^\top$.} (b) Deduce that one-sided Shampoo, $W_{t+1}=W_t-\eta L_t^{-1/2}G_t$, has regret at most $(\frac{D^2}{2\eta}+r\eta)\,n\tr(L_T^{1/2})$. (c) For $G_t=c_tuv^\top$ with fixed unit vectors $u,v$, show that as $\eps\to0$ the bound of \cref{thm:kronecker:regret} at its best step size tends to $\sqrt2D(\sum_tc_t^2)^{1/2}$, while diagonal AdaGrad's bound with $D$ as $\ell_\infty$ radius is $\sqrt2D\norm{u}_1\norm{v}_1(\sum_tc_t^2)^{1/2}$, up to $\sqrt{mn}$ times larger.\label{ex:kronecker:onesided}

\item \variation (The power $\frac12$ and grafting.) (a) With $\eps=0$, show that $L_t^{-1/2}G_tR_t^{-1/2}$ is homogeneous of degree $-1$ in the gradients, while Shampoo's step is homogeneous of degree $0$. (b) Grafting replaces a step $\Delta$ by $\Delta\norm{\Delta_{\mathrm{ref}}}_F/\norm{\Delta}_F$, with $\Delta_{\mathrm{ref}}$ the step of a reference method. Show that the grafted step has the degree of the reference, and that for Kronecker statistics it does not depend on the scalar $(\tr\Gamma)^{1/2}$ of \eqref{eq:kronecker:power}; with full-matrix AdaGrad as reference it is full-matrix AdaGrad's step.\label{ex:kronecker:graft}

\item \variation (Other noise, and $\epsilon$.) Let Shampoo's factors average gradients $\nabla\Loss(W)+\Xi$ of the Kronecker quadratic in the noise-dominated limit. (a) For $\vect(\Xi)\sim\cN(0,\sigma^2I)$, show that Shampoo is gradient descent ($q=0$). (b) For noise covariance $\sigma^2H^q$ with $q$ unknown in $[1,2]$, find the power that minimizes the worst-case exponent $\max_q\abs{1-qs}$. (c) Show that Adam with $\epsilon>0$ is not invariant under rescaling a coordinate's gradients by $c$: it tends to the $\epsilon=0$ update as $c\to\infty$ and to momentum gradient descent with step size $\eta/\epsilon$ as $c\to0$. Which modes lose the curvature-free property of \cref{prop:kronecker:soap} first? \hint{in (b), $s=2/3$ and the exponent is $1/3$.}\label{ex:kronecker:noise}

\item \simulation (The V and its floor.) Run \texttt{07-kronecker-ladder.py} (a) with $\kappa_A=900$ and $\kappa_C=1$, and check that the exponents still follow $\abs{1-qs}$ relative to $\kappa(H)=900$, although one-sided preconditioning is now exact; (b) with $100$, $1000$ and $10\,000$ sampled gradients, and check that the exponent at the bottom of each V falls like the inverse square root of their number, as the sampling error of the factors does (we measured $0.15$, $0.05$ and $0.016$ for $q=1$, $s=1$).\label{ex:kronecker:simV}

\item \simulation (Shampoo's transient.) Run plain Shampoo (sums, $\eps=10^{-12}$, power $\frac14$) on the Kronecker quadratic with $m=20$, $n=30$, $\kappa_A=\kappa_C=64$ and a standard Gaussian initial error, with step sizes $1$, $3$, $10$ and $30$. (a) Observe that for the two larger step sizes the loss first rises (in our runs by factors of about $4$ and $50$) and then falls faster than with the smaller ones. (b) Explain the transient: a step that is too long makes the gradients grow until the accumulated factors make it stable. (c) Compare the late-time rate with \cref{cor:kronecker:power} for $q=2$, and explain why the agreement is only rough.\label{ex:kronecker:transient}

\item \research (Are gradient statistics Kronecker?) For layers of a trained network, measure the relative distance $\min_{\Phi,\Psi}\norm{\Gamma-\Phi\otimes\Psi}_F/\norm{\Gamma}_F$ of the gradient second moment to the nearest Kronecker product \parencite{VanLoanPitsianis1993} and the constant $c^\star$. Find conditions on the joint law of inputs and back-propagated gradients, beyond independence, under which the distance is provably small. As far as we know no general theory exists.\label{ex:kronecker:kron}

\item \research (A schedule for the power.) Along a run the statistics pass from error-dominated ($q\approx2$) to noise-dominated ($q\approx1$) as the mean gradient shrinks (\cref{prop:kronecker:kfac}(iii)). In the noisy Kronecker quadratic started from an isotropic error, find the schedule of the power $s_t$ that minimizes the number of steps to a target loss, and compare it with a fixed power with grafting. This is open.\label{ex:kronecker:schedule}
\end{exercises}

\notesandrefs
Full-matrix and diagonal AdaGrad are due to \textcite{DuchiHazanSinger2011}. \Textcite{GuptaKorenSinger2018} introduced Shampoo for tensors of every order and proved \cref{thm:kronecker:bound} (their Lemma~8) and the regret bound; our proof of the bound passes through \cref{lem:kronecker:single}, which yields the middle term of \eqref{eq:kronecker:bound}. \Textcite{AnilEtAl2020} and \textcite{ShiEtAl2023} made Shampoo practical at scale, with grafting from \textcite{AgarwalEtAl2020Grafting}; \textcite{KasimbegEtAl2025} analyse the first AlgoPerf competition \parencite{DahlEtAl2023}.

K-FAC \parencite{MartensGrosse2015} approximates natural gradient descent \parencite{Amari1998}, and \textcite{GeorgeEtAl2018} corrected its eigenvalues; \textcite{KunstnerHennigBalles2019} and \textcite{Benzing2022} question the curvature reading of empirical Fishers and of K-FAC. \Textcite{MorwaniEtAl2025} relate Shampoo's factors to the optimal Kronecker approximation of full-matrix AdaGrad's statistics and of the Gauss--Newton matrix, and find the square of Shampoo's approximation, not the approximation itself, close to optimal. SOAP and the identity of \cref{prop:kronecker:adafactor} with Adafactor \parencite{ShazeerStern2018} are due to \textcite{VyasEtAl2025}, and \textcite{EschenhagenEtAl2025} take Shampoo's heuristics apart. That Shampoo without accumulation is the orthogonalized gradient, Muon's update without momentum (\cref{ch:spectral}), is observed by \textcite{BernsteinNewhouse2024}. We know no reference for the exponent law of \cref{cor:kronecker:power}, the exact constants of \cref{prop:kronecker:constant}, or \cref{prop:kronecker:soap}; all three are elementary consequences of \cref{prop:kronecker:partial,prop:kronecker:conditioning} and of Adam's invariance to diagonal rescaling \parencite{KingmaBa2015}.

\paragraph{Translating the literature.} \Textcite{GuptaKorenSinger2018} vectorize matrices by rows, so their Kronecker products appear in the reverse of our order (Exercise~\ref{ex:kronecker:vec}): the right side of their bound, $(\eps I_m+\sum_tG_tG_t^\top)^{1/2}\otimes(\eps I_n+\sum_tG_t^\top G_t)^{1/2}$, is our $R_T^{1/2}\otimes L_T^{1/2}$; their $L_t$, $R_t$, $r$ and $D$ are ours. \Textcite{MartensGrosse2015} write a layer's gradient as $g\bar a^\top$ ($g$ is our $\delta$, $\bar a$ our $a$ with a coordinate one appended for the bias) and approximate the Fisher block by $\bar A\otimes G$, where $\bar A$ is our $\E[aa^\top]$ and $G$ our $\E[\delta\delta^\top]$, not a gradient; they stack columns, as we do. \Textcite{VyasEtAl2025} write $L$, $R$, $Q_L$, $Q_R$ as we do, and their Shampoo with the $1/2$ power is our $s=\frac12$.

\printchapterbib
\end{refsection}


\casepart{Stability}{Practical step sizes sit at the edge of stability. Gradient descent drives the curvature up until the step is barely stable and then holds it there; the first steps of training meet curvature that is too high for the target step size, which is why the step size must start small; and curvature that grows with the gradient is tamed by clipping and by normalized steps.}{54}
% Chapter 8 of Optimizers in the Wild: the edge of stability.
\newcommand{\eosHmin}{\cH_{\min}}
\newcommand{\eosArea}{\mathcal{J}}

\casechapter{The Edge of Stability}{Thresholds, self-stabilization and central flows}{Learning rates at the stability limit}{ch:eos}
\begin{refsection}[chapters/08-eos.bib]
\nocite{HydrogenAtoms2026}

\begin{chaptersummary}
Learning rates are chosen as large as training tolerates, and full-batch training then runs with its curvature pinned at the stability threshold of the optimizer: the sharpness at $2/\eta$ for gradient descent, the preconditioned sharpness at $38/\eta$ for Adam with $\beta_1=0.9$. We derive these thresholds from one stability test, and then study the cubic model of Damian, Nichani and Lee, in which a single direction oscillates and the oscillation pushes the curvature down. In its continuous limit we prove that every cycle of oscillation conserves an energy and closes, and that over each cycle the curvature averages exactly $2/\eta$ while the squared oscillation averages exactly twice the rate of progressive sharpening; for the gradient-descent map itself we prove that the same energy increases at every step, so that the cycles spiral slowly outward, and we derive the rate. Averaging over the cycles gives the central flow, gradient flow projected onto the region where the step size is stable. The model takes progressive sharpening as an input and does not explain it.
\end{chaptersummary}

\section{In the wild: the largest learning rate that does not diverge}\label{sec:eos:wild}

Gradient descent on a quadratic multiplies the component of $\theta_t-\theta^\star$ along an eigenvector of the Hessian with eigenvalue $\lambda_i$ by $1-\eta\lambda_i$ at every step, so it converges from every initialization exactly when $\eta\lambda_i<2$ for every eigenvalue (\cref{ch:prologue}). The largest eigenvalue of the Hessian, $S(\theta)=\lmax(\nabla^2\Loss(\theta))$, is the \term{sharpness}\index{sharpness}, and $\eta S<2$ is the classical stability condition. In this chapter the learning rate and the step size are the same number $\eta$; we say learning rate when we describe practice and step size in the theory. Read as advice, the stability condition says: estimate the curvature, then choose the step. Networks are not trained that way.

\begin{practice}[folklore]
Choose the learning rate by a sweep on a logarithmic grid, with neighbouring values a factor of two or three apart, and keep the value with the lowest loss after a fixed budget; it is usually among the largest values that train without diverging, and a run that diverges is restarted with a smaller value. The curvature of the loss is not measured, neither at initialization nor later. Momentum, and Adam's average of the gradient with $\beta_1=0.9$ \parencite{KingmaBa2015}, stay at their defaults while the learning rate is swept.
\end{practice}

The rule is folklore in the sense of this book: it describes how sweeps are commonly run, not a measured regularity, and the \emph{Deep Learning Tuning Playbook} \parencite{GodboleEtAl2023} is the most systematic guide to running them. What the rule leaves unsaid is that the curvature is not a fixed property of the loss that the learning rate has to respect. It responds to the learning rate.\sidenote{That the step size limits where gradient descent can be stable, and so selects flat minima, is the dynamical-stability argument of \textcite{WuMaE2018}; \textcite{JastrzebskiEtAl2019} observed that the sharpness reached early in training adapts to the step size.}

\begin{observation}[Progressive sharpening and the edge of stability; \textcite{CohenEtAl2021}]\label[observation]{obs:eos:gd}
Train a network by full-batch gradient descent with a fixed step size $\eta$. (i) While $S(\theta_t)<2/\eta$, the sharpness tends to rise. (ii) Once it reaches $2/\eta$, it stops rising and hovers at or just above $2/\eta$, while the training loss keeps decreasing over long time scales, though not monotonically over short ones. (iii) With heavy-ball momentum $\beta$ the sharpness settles instead near $(2+2\beta)/\eta$, and with Nesterov momentum near $(2+2\beta)/((1+2\beta)\eta)$: in each case the largest sharpness at which the method is stable on a quadratic.
\end{observation}

Cohen and coauthors named the first stage \term{progressive sharpening}\index{progressive sharpening} and the second the \term{edge of stability}\index{edge of stability}, and documented both for fully connected, convolutional and transformer networks, with the square loss and with the cross-entropy loss. For momentum they derived the thresholds in (iii), which they call maximal stable sharpnesses, and observed the sharpness rise to them and level off in runs with $\eta=0.01$. The observation is robust for full-batch gradient descent, and its scope is limited in two ways that will matter later: progressive sharpening is weaker for wide or shallow networks and for easy tasks, and with the cross-entropy loss the sharpness typically falls at the end of training \parencite{CohenEtAl2021}.

\begin{observation}[The adaptive edge of stability; \textcite{CohenEtAl2022}]\label[observation]{obs:eos:adam}
Train a network full batch with an adaptive method, such as Adam or RMSProp, whose update is $\theta_{t+1}=\theta_t-\eta P_t^{-1}\hat m_t$ with a diagonal preconditioner $P_t$ ($\hat m_t$ is the gradient itself for RMSProp). The \term{preconditioned sharpness}\index{preconditioned sharpness} $S_P=\lmax(P_t^{-1}\nabla^2\Loss(\theta_t))$ rises until it reaches the stability threshold of the same method with $P_t$ frozen, which is $38/\eta$ for Adam with $\beta_1=0.9$, and then equilibrates there. The sharpness itself is not held at a threshold: the iterates keep moving into regions of higher curvature, and the preconditioner grows to compensate. Similar effects occur with minibatches, the more so the larger the batch.
\end{observation}

The source is an unrefereed preprint (arXiv 2022; a version appeared at a NeurIPS 2023 workshop), and most of its experiments are full batch, but the claim is easy to check: \cref{fig:eos:network} reproduces it, together with \cref{obs:eos:gd}, in a network small enough to train in half a minute. With minibatches the picture changes.

\begin{observation}[Minibatches; \textcite{JastrzebskiEtAl2020}, \textcite{AndreyevBeneventano2024}]\label[observation]{obs:eos:sgd}
With minibatch stochastic gradient descent the full-batch sharpness levels off below $2/\eta$, at a value set by the step size and the batch size. What settles at $2/\eta$ is a sharpness adapted to the minibatch: the expected curvature of a minibatch loss along its own minibatch gradient, which Andreyev and Beneventano call the \term{batch sharpness}\index{batch sharpness}.
\end{observation}

The first sentence is older and well replicated \parencite{JastrzebskiEtAl2019,JastrzebskiEtAl2020,LeeJang2023}; the second rests on an unrefereed preprint (arXiv 2024, revised in 2025). This chapter is about full batch, and \cref{ch:noise} takes up the noise.

The three observations have one shape: a stability threshold computed from a quadratic model with everything else frozen, and a training process that drives the relevant curvature up to the threshold and then holds it there. Two questions follow. Why does the iteration not diverge once the threshold is crossed, as the quadratic model says it must? And what does training do while it sits at the threshold: where does it go, and what does the oscillation cost? Volume 2 of this series, \emph{The Hydrogen Atoms of Deep Learning}, Chapter 15, answers the first question for a single excursion. In a two-parameter product model, a step size above the threshold makes the loss spike and throws the iterates to a flatter minimizer, a \term{catapult}\index{catapult} in the language of \textcite{LewkowyczEtAl2020}. That model has no progressive sharpening, so after one catapult it stays below the threshold. This chapter adds the sharpening.

\section{The idealization: one oscillating direction}\label{sec:eos:ideal}

The quadratic model freezes the Hessian. To see the curvature respond to the step size we need three ingredients and no more: a direction along which the iterates oscillate when the step is too large, a curvature along that direction that depends on the other parameters, and a training force that pushes the other parameters toward higher curvature. The following family of losses has exactly these.

\begin{assumption}[One oscillating direction]\label[assumption]{ass:eos:onedir}
The parameters split as $\theta=(x,w)\in\R\times\R^{d-1}$, and
\begin{equation}\label{eq:eos:onedir}
  \Loss(x,w)=\Loss_0(w)+\tfrac12\,s(w)\,x^2 ,
\end{equation}
where $\Loss_0$ and $s$ are twice continuously differentiable and $s(w)>\lmax(\nabla^2\Loss_0(w))$ for every $w$.
\end{assumption}

On the floor $x=0$ of the valley the Hessian is block diagonal, with the block $s(w)$ in the direction of $x$ and $\nabla^2\Loss_0(w)$ in the others, so the sharpness there is $s(w)$ and its eigenvector is the $x$-axis. Gradient descent with step size $\eta$ is, exactly,
\begin{equation}\label{eq:eos:gdonedir}
  x_{t+1}=\bigl(1-\eta s(w_t)\bigr)x_t,\qquad
  w_{t+1}=w_t-\eta\nabla\Loss_0(w_t)-\tfrac\eta2\,x_t^2\,\nabla s(w_t).
\end{equation}
The first equation is the quadratic picture: the oscillating coordinate is multiplied by $1-\eta s$, it decays while $\eta s<2$, and above the threshold it flips sign and grows. The second equation contains the new force. The oscillation pushes $w$ down the gradient of the curvature, with a strength proportional to $x^2$: the larger the oscillation, the faster the sharpness falls. Two numbers measure the competition between the training force and the push,
\begin{equation}\label{eq:eos:alphagamma}
  \alpha(w)=-\ip{\nabla\Loss_0(w),\nabla s(w)},\qquad \gamma(w)=\norm{\nabla s(w)}_2^2 .
\end{equation}
Along gradient flow $\dot w=-\nabla\Loss_0(w)$ the sharpness changes at the rate $\dot s=\alpha$, so $\alpha>0$ is progressive sharpening; along the push, a unit of $x^2$ lowers $s$ at the rate $\frac12\gamma$.

The smallest instance of the family linearizes $\Loss_0$ and $s$ near a point at the threshold. If $\Loss_0(w)=\Loss_0(0)+\ip{g,w}$ and $s(w)=2/\eta+\ip{c,w}$ with $c\ne0$, then $y=\ip{c,w}$ is the excess of the curvature over $2/\eta$, and \eqref{eq:eos:gdonedir} gives $x_{t+1}=-(1+\eta y_t)x_t$ and $y_{t+1}=y_t-\eta\ip{c,g}-\frac\eta2\norm c_2^2x_t^2$. The coupling $\norm{c}_2^2$ can be absorbed into $x$: the pair $(\norm c_2x_t,y_t)$ follows the recursion below with $\alpha=-\ip{g,c}$, while the component of $w$ orthogonal to $c$ moves by $-\eta$ times the component of $g$ orthogonal to $c$ at every step and never feels the oscillation. What remains is a two-dimensional map.

\begin{definition}[The cubic model]\label[definition]{def:eos:cubic}
Fix a step size $\eta>0$ and a rate $\alpha>0$. The \term{cubic model}\index{cubic model} is the loss
\[
  \Loss(x,y)=\tfrac12\Bigl(\frac2\eta+y\Bigr)x^2-\alpha y,\qquad (x,y)\in\R^2,
\]
trained by gradient descent with step size $\eta$:
\begin{equation}\label{eq:eos:cubicmap}
  x_{t+1}=-(1+\eta y_t)\,x_t,\qquad y_{t+1}=y_t+\eta\Bigl(\alpha-\tfrac12x_t^2\Bigr).
\end{equation}
\end{definition}

The name refers to the coupling $\frac12yx^2$, which is cubic. The coordinate $x$ is the displacement along the oscillating direction and $y$ the excess of its curvature over $2/\eta$; on the line $x=0$ the sharpness is exactly $2/\eta+y$, and elsewhere it is the largest eigenvalue of $\bigl(\begin{smallmatrix}2/\eta+y&x\\x&0\end{smallmatrix}\bigr)$, which is $2/\eta+y+O(\eta x^2)$. The loss is unbounded below in $y$: it is a model of a region of the landscape in which training raises the curvature at the constant rate $\alpha$, not of a whole landscape. With $x=0$ the curvature grows by $\eta\alpha$ per step forever; with $x\ne0$ it grows only while $x^2<2\alpha$.

\begin{dictionary}
the coordinate $x$, flipping sign at every step & the component of $\theta_t$ along the top eigenvector of the Hessian (of $P^{-1}\nabla^2\Loss$ for Adam), oscillating at period two\\
the other coordinates $w$, and $\Loss_0(w)$ & all other parameters, and the loss at the centre of the oscillation\\
the curvature $s(w)=2/\eta+y$ & the sharpness $S(\theta_t)$, and its distance $y$ from the threshold\\
the rate $\alpha=-\ip{\nabla\Loss_0,\nabla s}$ & the speed of progressive sharpening, $dS/d\tau$ along gradient flow\\
the coupling $\gamma=\norm{\nabla s}_2^2$ & how fast an oscillation lowers the sharpness, through the third derivative of the loss\\
the time $\tau=\eta t$ & the number of steps times the learning rate, the clock in which gradient flow has unit speed\\
the time average of $x^2$ & the variance $\sigma^2$ of the oscillation in a central flow\\
a frozen preconditioner $P$ & Adam's $\sqrt{\hat v_t}+\epsilon$, which moves on the time scale $1/(1-\beta_2)$
\end{dictionary}

\paragraph{What was thrown away.} The oscillating direction is fixed, while in a network the top eigenvector turns; only one direction is at the threshold, while in a network several can be; the loss is exactly quadratic and even in $x$, so there are no terms $x^3$ or $x^4$ along the oscillation; and in the cubic model the rates $\alpha$ and $\gamma$ are constant, so the model cannot leave the edge by itself. Gradients are exact, and there is no momentum and no preconditioner; for quadratics \cref{sec:eos:thresholds} puts both back. The choice keeps exactly the two forces whose balance is the phenomenon, the training force that raises the curvature and the push of the oscillation that lowers it, and it is what a third-order Taylor expansion of a general loss produces near the threshold (\cref{sec:eos:selfstab}).

\paragraph{What one had to notice.} In the quadratic theory an oscillation of growing amplitude is the beginning of divergence. The insight of \textcite{DamianNichaniLee2023} is that it is also the restoring force: through the third derivative of the loss, the oscillation produces the gradient $\frac12x^2\nabla s$, which lowers the curvature that caused it. This is \term{self-stabilization}\index{self-stabilization}. A second, technical observation makes the model solvable: the right variables are the logarithm of the squared amplitude, $A=\log x^2$, which forgets the sign flips, and the excess curvature $y$. In these variables a cycle of oscillation is a closed curve.

\section{The theory: stability thresholds}\label{sec:eos:thresholds}

All four thresholds of the in-the-wild section come from one stability test. A second-order linear recurrence $e_{t+1}=a e_t-be_{t-1}$ with real $a,b$ has the companion matrix $\bigl(\begin{smallmatrix}a&-b\\1&0\end{smallmatrix}\bigr)$, whose eigenvalues are the roots of $z^2-az+b$; its solutions tend to zero from every pair of initial values if and only if both roots lie in the open unit disk, and by the Schur--Cohn (Jury) test this happens if and only if $\abs b<1$ and $\abs a<1+b$ (\cref{lem:toolkit:recurrence}).

\begin{proposition}[Stability thresholds]\label[proposition]{prop:eos:thresholds}
Let $\Loss(\theta)=\frac12(\theta-\theta^\star)^\top H(\theta-\theta^\star)$ with $H$ symmetric positive semidefinite and $S=\lmax(H)>0$, and let $\eta>0$ and $0\le\beta<1$. Each of the following iterations converges to a minimizer of $\Loss$ from every initialization if and only if the stated inequality holds.
\begin{enumerate}
\item Gradient descent, $\theta_{t+1}=\theta_t-\eta\nabla\Loss(\theta_t)$: \ $\eta S<2$.
\item Heavy ball, $\theta_{t+1}=\theta_t-\eta\nabla\Loss(\theta_t)+\beta(\theta_t-\theta_{t-1})$: \ $\eta S<2+2\beta$.
\item Momentum as a moving average, $m_t=\beta m_{t-1}+(1-\beta)\nabla\Loss(\theta_t)$ and $\theta_{t+1}=\theta_t-\eta m_t$: \ $\eta S<(2+2\beta)/(1-\beta)$.
\item Adam with a frozen preconditioner: $m_t$ as in (iii) with $\beta=\beta_1$, and $\theta_{t+1}=\theta_t-\eta P^{-1}m_t$ with a fixed symmetric $P\succ0$: \ $\eta S_P<(2+2\beta_1)/(1-\beta_1)$, where $S_P=\lmax(P^{-1}H)$.
\end{enumerate}
For $\beta_1=0.9$ the threshold in (iv) is $\eta S_P<38$.
\end{proposition}

\begin{proof}
The plan is to split each iteration into scalar recursions, one for each eigenvector of the Hessian, and to apply the stability test to each. Let $Hq_i=\lambda_iq_i$ with orthonormal $q_1,\dots,q_d$, and write $e_{t,i}=\ip{q_i,\theta_t-\theta^\star}$. Since $\nabla\Loss(\theta)=H(\theta-\theta^\star)$, every iteration acts on each $e_{t,i}$ separately, and it converges to a minimizer exactly when every $e_{t,i}$ converges: the minimizers are the points $\theta^\star+v$ with $Hv=0$.

(i) Here $e_{t+1,i}=(1-\eta\lambda_i)e_{t,i}$. If $\lambda_i=0$ the component is constant. If $\lambda_i>0$, the only possible limit is $0$, and $(1-\eta\lambda_i)^te_{0,i}$ tends to $0$ for every $e_{0,i}$ if and only if $\abs{1-\eta\lambda_i}<1$, that is, $\eta\lambda_i<2$. All components converge if and only if this holds for the largest eigenvalue.

(ii) Here $e_{t+1,i}=(1+\beta-\eta\lambda_i)e_{t,i}-\beta e_{t-1,i}$. If $\lambda_i=0$, the roots of $z^2-(1+\beta)z+\beta$ are $1$ and $\beta$, so $e_{t,i}=a+b\beta^t$ converges. If $\lambda_i>0$, a limit $e$ would satisfy $\eta\lambda_ie=0$, so the component converges for every initialization if and only if it tends to zero, and by the stability test this happens if and only if $\beta<1$, which is assumed, and $\abs{1+\beta-\eta\lambda_i}<1+\beta$, which means $0<\eta\lambda_i<2+2\beta$.

(iii) For $t\ge1$, $\theta_{t+1}-\theta_t=-\eta m_t=-\eta\beta m_{t-1}-\eta(1-\beta)\nabla\Loss(\theta_t)=\beta(\theta_t-\theta_{t-1})-\eta(1-\beta)\nabla\Loss(\theta_t)$. This is the heavy ball (ii) with step size $\eta(1-\beta)$, started from $\theta_0$ and $\theta_1=\theta_0-\eta\beta m_{-1}-\eta(1-\beta)\nabla\Loss(\theta_0)$. If $\beta>0$, every $\theta_1$ arises from some $m_{-1}$, so (ii) gives convergence for every initialization if and only if $\eta(1-\beta)S<2+2\beta$; if $\beta=0$, the iteration is (i).

(iv) Put $\varphi_t=P^{1/2}(\theta_t-\theta^\star)$ and $\mu_t=P^{-1/2}m_t$. Multiplying the two updates by $P^{1/2}$ and $P^{-1/2}$ gives $\mu_t=\beta_1\mu_{t-1}+(1-\beta_1)H_P\varphi_t$ and $\varphi_{t+1}=\varphi_t-\eta\mu_t$, with $H_P=P^{-1/2}HP^{-1/2}$. This is iteration (iii) on the quadratic with Hessian $H_P$, which is symmetric positive semidefinite and similar to $P^{-1}H=P^{-1/2}H_PP^{1/2}$, so $\lmax(H_P)=S_P$, and $\varphi\mapsto P^{-1/2}\varphi+\theta^\star$ maps the minimizers of one problem onto those of the other. Finally $(2+1.8)/0.1=38$.
\end{proof}

Three remarks place the proposition. First, the moving-average form (iii), which Adam uses, is the heavy ball with the step size scaled by $1-\beta$, so at a fixed $\eta$ it tolerates a curvature $(1+\beta)/(1-\beta)$ times larger than gradient descent: $19$ times for $\beta=0.9$ and $199$ times for $\beta=0.99$. Second, Adam's bias corrections multiply its step by factors that tend to one, and its preconditioner is not frozen, so (iv) is a statement about a snapshot; \cref{obs:eos:adam} says that the snapshot's threshold is where Adam's preconditioned sharpness equilibrates. Third, a threshold says when the quadratic model of the loss at the current point is unstable, and nothing about what happens next. That is the content of the rest of the chapter. Nesterov's threshold, the third case of \cref{obs:eos:gd}, is \cref{ex:eos:nesterov}.

One more fact about snapshots is needed later. Adam's preconditioner can shrink, and its preconditioned sharpness rise, only at a speed that $\beta_2$ fixes.

\begin{proposition}[How fast the preconditioner can sharpen]\label[proposition]{prop:eos:speed}
Let $v_{t+1}=\beta_2v_t+(1-\beta_2)g_{t+1}^2$ coordinatewise, with $0<\beta_2<1$ and $v_t>0$, let $P_t=\diag(\sqrt{v_t})+\epsilon I$ with $\epsilon\ge0$, and let $H$ be symmetric with $\lmax(P_t^{-1}H)>0$. Then $P_{t+1}\succeq\sqrt{\beta_2}\,P_t$ and
\[
  \lmax(P_{t+1}^{-1}H)\le\beta_2^{-1/2}\,\lmax(P_t^{-1}H).
\]
\end{proposition}

\begin{proof}
Coordinatewise, $v_{t+1}\ge\beta_2v_t$, so $\sqrt{v_{t+1}}+\epsilon\ge\sqrt{\beta_2}\sqrt{v_t}+\epsilon\ge\sqrt{\beta_2}(\sqrt{v_t}+\epsilon)$, which is $P_{t+1}\succeq\sqrt{\beta_2}P_t$ for diagonal matrices. For $P\succ0$ the largest eigenvalue of $P^{-1}H$ is the largest value of $u^\top Hu/u^\top Pu$ over $u\ne0$, and when it is positive it is attained at a $u$ with $u^\top Hu>0$; for such $u$ the quotient for $P_{t+1}$ is at most $\beta_2^{-1/2}$ times the quotient for $P_t$.
\end{proof}

At a fixed Hessian the preconditioned sharpness can therefore grow by at most $\frac12\log(1/\beta_2)\approx\frac12(1-\beta_2)$ per step in logarithm: by a factor $e$ in no fewer than about $2000$ steps when $\beta_2=0.999$. The bound is attained when the new squared gradients are negligible next to $v_t$, which is the case once the iterates have nearly converged. The bias correction $\hat v_t=v_t/(1-\beta_2^{t+1})$ adds a factor $(1-\beta_2^{t+1})/(1-\beta_2^{t+2})$, which matters only for $t\lesssim1/(1-\beta_2)$.

\section{The theory: self-stabilization}\label{sec:eos:selfstab}

\subsection{Where the cubic model comes from}

The cubic model is what a general loss looks like near the threshold, to leading order. The derivation is due to \textcite{DamianNichaniLee2023}, and we give it in the form we need.

\begin{heuristic}[Heuristic derivation of the cubic model]
Let $\Loss$ be three times continuously differentiable and let $\theta^\star$ be a point with $S(\theta^\star)=2/\eta$, a simple top eigenvalue with unit eigenvector $u$, and $\ip{\nabla\Loss(\theta^\star),u}=0$. A simple eigenvalue is differentiable, with derivative $dS=u^\top(dH)u$, so $\nabla S(\theta^\star)=\nabla^3\Loss(\theta^\star)[u,u,\cdot\,]$. Write $\theta=\theta^\star+xu+\delta$ with $\delta\perp u$, put $y=\ip{\nabla S(\theta^\star),\delta}$, and assume $\ip{\nabla S(\theta^\star),u}=0$. The second-order Taylor expansion of the gradient,
\[
  \nabla\Loss(\theta)\approx\nabla\Loss(\theta^\star)+\nabla^2\Loss(\theta^\star)(xu+\delta)+\tfrac12\nabla^3\Loss(\theta^\star)[xu+\delta,xu+\delta,\cdot\,],
\]
has the component $\frac2\eta x+xy$ along $u$, after dropping $\frac12x^2\nabla^3\Loss[u,u,u]$ and $\frac12\nabla^3\Loss[u,\delta,\delta]$, and the component $-\alpha+\frac12\gamma x^2$ along $\nabla S$, with $\alpha=-\ip{\nabla\Loss(\theta^\star),\nabla S(\theta^\star)}$ and $\gamma=\norm{\nabla S(\theta^\star)}_2^2$, after dropping $\ip{\nabla S,\nabla^2\Loss\,\delta}$ and the other terms that involve $\delta$. A step of gradient descent therefore maps $(x,y)$ to approximately $\bigl(-(1+\eta y)x,\;y+\eta(\alpha-\frac12\gamma x^2)\bigr)$, which is \eqref{eq:eos:cubicmap} for $(\sqrt\gamma\,x,y)$. The motion of $\delta$ orthogonal to $\nabla S$ is gradient descent on $\Loss$, undisturbed by the oscillation.
\end{heuristic}

The dropped terms are of two kinds. Some describe how the reference point, its eigenvector and the rates $\alpha$ and $\gamma$ drift; Damian, Nichani and Lee follow them by moving the reference point along a constrained trajectory, and \cref{sec:eos:central} averages over them. Others are higher-order terms along the oscillation itself, such as $x^2\nabla^3\Loss[u,u,u]$ and a fourth-order term; \cref{sec:eos:misses} shows that they decide what happens over long times.

\subsection{Scales, and why two steps}

The cubic map flips the sign of $x$ at every step, so it is not close to the identity, and no differential equation can describe $x_t$ itself. The squared amplitude $x_t^2$ does not flip: by \eqref{eq:eos:cubicmap}, $x_{t+1}^2=(1+\eta y_t)^2x_t^2$, and every quantity that is even in $x$ evolves by a map close to the identity. Equivalently, the two-step map sends $x_t$ to $(1+\eta y_{t+1})(1+\eta y_t)x_t$, without a sign flip. A continuous limit can describe the envelope of the oscillation, not the oscillation.

Which quantities are of order one? In the variables $A=\log x^2$ and $y$, the map \eqref{eq:eos:cubicmap} reads
\begin{equation}\label{eq:eos:Amap}
  A_{t+1}=A_t+2\log(1+\eta y_t),\qquad y_{t+1}=y_t+\eta\Bigl(\alpha-\tfrac12e^{A_t}\Bigr).
\end{equation}
If $\alpha$, $y$ and $e^A=x^2$ are of order one, then each step changes $A$ and $y$ by amounts of order $\eta$, and the natural clock is $\tau=\eta t$, the time of gradient flow, in which the sharpness rises at the rate $\alpha$ when $x=0$. In this scaling the curvature along the oscillation is within $O(1)$ of $2/\eta$, so $\eta S-2=\eta y$ is of order $\eta$, and one cycle lasts a time of order $1/\sqrt\alpha$ in $\tau$, that is, of order $1/(\eta\sqrt\alpha)$ steps. The model has a single dimensionless parameter: with $a=A-\log(2\alpha)$, $Y=y/\sqrt\alpha$ and $s=\sqrt\alpha\,\tau$, the map becomes $a\mapsto a+2\log(1+\eps Y)$, $Y\mapsto Y+\eps(1-e^a)$, with
\[
  \eps=\eta\sqrt\alpha .
\]
Its square $\eps^2=\eta^2\alpha$ is the amount by which progressive sharpening raises $\eta S$ in one step. The continuous limit is $\eps\to0$: sharpening that is slow on the scale of a step, which is the regime of \cref{obs:eos:gd}.

\subsection{The theorem}

The limit of \eqref{eq:eos:Amap} is the differential equation
\begin{equation}\label{eq:eos:flow}
  \frac{dA}{d\tau}=2y,\qquad \frac{dy}{d\tau}=\alpha-\tfrac12e^{A}.
\end{equation}
Write $\phi(A)=\frac12e^A-\alpha A$ and
\begin{equation}\label{eq:eos:H}
  \cH(A,y)=y^2+\phi(A)=y^2+\tfrac12e^A-\alpha A .
\end{equation}
The function $\phi$ is strictly convex, with $\phi''=\frac12e^A>0$, and it tends to $+\infty$ as $A\to\pm\infty$; its minimum is at $A^*=\log(2\alpha)$, where $\phi(A^*)=\alpha(1-\log2\alpha)=:\eosHmin$. So $\cH$ has a unique minimum $\eosHmin$, at $(A^*,0)$, which corresponds to the orbit $x=\pm\sqrt{2\alpha}$, $y=0$ of period two of gradient descent. For $h>\eosHmin$ let $A_-<A^*<A_+$ be the two solutions of $\phi(A)=h$.

\begin{theorem}[Self-stabilization in the continuous limit]\label{thm:eos:selfstab}
Let $\alpha>0$, let $(A(\tau),y(\tau))$ solve \eqref{eq:eos:flow} from a point other than $(A^*,0)$, and let $h=\cH(A(0),y(0))$.
\begin{enumerate}
\item The energy $\cH$ is constant along the solution.
\item The solution is periodic, with period
\[
  T(h)=\int_{A_-}^{A_+}\frac{dA}{\sqrt{h-\phi(A)}}<\infty .
\]
\item Over each period the time averages are $\langle y\rangle=0$ and $\langle e^A\rangle=2\alpha$: the curvature along the oscillation averages exactly $2/\eta$, and the squared amplitude $x^2$ averages exactly $2\alpha$.
\item Gradient descent \eqref{eq:eos:cubicmap} started at a point $(x_0,y_0)$ with $\log x_0^2=A(0)$ and $y_0=y(0)$ follows the solution: for every $\tau_{\max}>0$ there are constants $C$ and $\eta_0>0$, which depend on $\alpha$, $A(0)$, $y(0)$ and $\tau_{\max}$ but not on $\eta$, such that for every $\eta\in(0,\eta_0]$ and every integer $t$ with $0\le t\le\tau_{\max}/\eta$,
\[
  \abs{\log x_t^2-A(\eta t)}+\abs{y_t-y(\eta t)}\le C\eta .
\]
\end{enumerate}
\end{theorem}

The theorem says that, to leading order in $\eps$, an oscillation at the edge is a closed cycle: the curvature rises above $2/\eta$, the oscillation grows and pushes it back down below $2/\eta$, the oscillation dies, and the cycle repeats with the same energy. Part (iii) is exact, not approximate. Over a cycle the curvature spends as much time above $2/\eta$ as it needs to spend below for its average to be exactly the threshold, and the oscillation carries exactly the energy that cancels progressive sharpening: $\frac12\langle x^2\rangle=\alpha$, or $\frac12\gamma\langle x^2\rangle=\alpha$ before the coupling is absorbed. Part (iv) is the continuous limit; we prove it after the rest.

\begin{proof}[Proof of (i)--(iii)]
The plan is to show that $\cH$ is a conserved energy, that its level sets are closed curves without equilibria, and that a periodic solution has the stated averages because the averages of derivatives over a period vanish.

(i) Along a solution, $\frac{d}{d\tau}\cH=2y\,\dot y+\phi'(A)\dot A=2y(\alpha-\frac12e^A)+(\frac12e^A-\alpha)\,2y=0$.

(ii) The sublevel sets of $\cH$ are compact, because $\cH\ge y^2+\eosHmin$ and $\phi$ tends to infinity at both ends, so by (i) the solution stays in a compact set and exists for all $\tau$. Since $h>\eosHmin$, the level set $\{\cH=h\}$ consists of the two arcs $y=\pm\sqrt{h-\phi(A)}$ over $A_-\le A\le A_+$, joined at the turning points $(A_\pm,0)$. The vector field of \eqref{eq:eos:flow} vanishes only at $(A^*,0)$, which is not on this curve. On the upper arc $\dot A=2y=2\sqrt{h-\phi(A)}>0$, so $A$ increases, and the time it takes to go from a point of the arc to $A_+$ is $\int dA/(2\sqrt{h-\phi(A)})$, which is at most $T(h)/2$. This integral is finite. By convexity, $\phi$ lies below its chord from $(A^*,\eosHmin)$ to $(A_+,h)$, so for $A^*\le A<A_+$,
\[
  h-\phi(A)\ge(h-\eosHmin)\frac{A_+-A}{A_+-A^*},
\]
and $\int_{A^*}^{A_+}dA/\sqrt{h-\phi(A)}\le2(A_+-A^*)/\sqrt{h-\eosHmin}$; the same argument bounds the integral over $[A_-,A^*]$. At the turning point $(A_+,0)$ we have $\dot y=-\phi'(A_+)<0$, because $A_+>A^*$, so the solution passes to the lower arc, along which $A$ decreases, reaches $(A_-,0)$ in the time $\int_{A_-}^{A_+}dA/(2\sqrt{h-\phi(A)})$, and there passes to the upper arc because $\dot y=-\phi'(A_-)>0$. A solution that starts on the lower arc does the same in the other order. After the time $T(h)$, the sum of the two transit times, the solution is back at its starting point, and by uniqueness of solutions it is periodic.

(iii) Over a period $[\tau_0,\tau_0+T]$ both $A$ and $y$ return to their initial values. Integrating the two equations of \eqref{eq:eos:flow} over the period gives $0=A(\tau_0+T)-A(\tau_0)=2\int y\,d\tau$ and $0=y(\tau_0+T)-y(\tau_0)=\int(\alpha-\frac12e^A)\,d\tau$, which are the two averages. The curvature along the oscillation is $2/\eta+y$, and $e^A=x^2$.
\end{proof}

The argument for (iii) is the one that gives the averages of the predator--prey cycles of Lotka and Volterra: over a closed cycle the average of a derivative is zero.\sidenote{The analogy is only partial. In the predator--prey system both populations are positive and enter through logarithms; here $y$ changes sign, and \eqref{eq:eos:flow} is the motion of a particle in the potential well $\phi$, since $d^2A/d\tau^2=-2\phi'(A)$.} The small cycles have the period of the linearization: near $(A^*,0)$, $d^2A/d\tau^2=2\alpha-e^A\approx-2\alpha(A-A^*)$, so $T(h)\to2\pi/\sqrt{2\alpha}$ as $h\downarrow\eosHmin$ (\cref{ex:eos:period}), about $2\pi/(\eta\sqrt{2\alpha})$ steps. Larger cycles are slower, because they spend a long time with $A$ very negative, an oscillation so small that nothing happens except that the curvature climbs back at the rate $\alpha$.

Part (iv) is an error estimate of the kind used for the Euler method, with one difference: the map \eqref{eq:eos:Amap} is not the Euler method of \eqref{eq:eos:flow}, because of the logarithm, but it agrees with it to second order in $\eta$.

\begin{lemma}[The continuous limit]\label[lemma]{lem:eos:limit}
Let $z(\tau)=(A(\tau),y(\tau))$ solve \eqref{eq:eos:flow} on $[0,\tau_{\max}]$, and let $K$ be the set of points within $\ell_1$-distance $1$ of the curve $\{z(\tau):0\le\tau\le\tau_{\max}\}$. Let $Y=\max_K\abs y$ and $E=\max_Ke^A$, and put $L=\max\{4,E/2\}$ and $C_1=\alpha+\frac12E+2Y^2+\frac12EY$. Let $z_t=(A_t,y_t)$ follow \eqref{eq:eos:Amap} from $z_0=z(0)$. If $\eta\le1/(2Y)$ and $C_1\eta\,(e^{L\tau_{\max}}-1)\le L$, then for every integer $t$ with $0\le t\le\tau_{\max}/\eta$ the iterate $z_t$ lies in $K$, $x_t\ne0$, and
\[
  \norm{z_t-z(\eta t)}_1\le\frac{C_1}{L}\bigl(e^{L\eta t}-1\bigr)\,\eta .
\]
\end{lemma}

\begin{proof}
The plan is to bound the error made in one step by the exact solution, to show that the map \eqref{eq:eos:Amap} is Lipschitz with constant $1+L\eta$ on $K$, and to add up the errors with the discrete Grönwall inequality (\cref{thm:toolkit:gronwall}). Write $F$ for the map \eqref{eq:eos:Amap}. The set $K$ is compact, and on it $\abs{\eta y}\le\frac12$, so $F$ is defined and smooth there.

\emph{One step of the exact solution.} Along the solution $\abs y\le Y$ and $e^A\le E$, so $\abs{\dot y}\le\alpha+\frac12E$ and $\abs{\dot A}\le2Y$. Fix $\tau\le\tau_{\max}-\eta$ and write $y=y(\tau)$, $A=A(\tau)$. The first component of $z(\tau+\eta)-F(z(\tau))$ is
\[
  \int_0^\eta2\bigl(y(\tau+r)-y\bigr)\,dr+2\bigl(\eta y-\log(1+\eta y)\bigr).
\]
The integral is at most $2\int_0^\eta r(\alpha+\frac12E)\,dr=(\alpha+\frac12E)\eta^2$ in absolute value. For the second term we use $0\le u-\log(1+u)\le u^2$ for $u\ge-\frac12$: the function $u^2-u+\log(1+u)$ vanishes at $0$ and has derivative $u(1+2u)/(1+u)$, which has the sign of $u$ there. So the second term lies in $[0,2\eta^2Y^2]$. The second component is $-\frac12\int_0^\eta(e^{A(\tau+r)}-e^{A})\,dr$, and $\abs{e^{A(\tau+r)}-e^A}\le E\abs{A(\tau+r)-A}\le2EYr$, so it is at most $\frac12EY\eta^2$. Together, $\norm{z(\tau+\eta)-F(z(\tau))}_1\le C_1\eta^2$.

\emph{Lipschitz bound.} Let $z,z'\in K$ with the segment between them in $K$. The derivative of $2\log(1+\eta y)$ in $y$ is $2\eta/(1+\eta y)\le4\eta$, and the derivative of $\frac12\eta e^A$ in $A$ is at most $\frac12\eta E$, so
\[
  \norm{F(z)-F(z')}_1\le\Bigl(1+\tfrac12\eta E\Bigr)\abs{A-A'}+(1+4\eta)\abs{y-y'}\le(1+L\eta)\norm{z-z'}_1 .
\]

\emph{Adding up.} Let $e_t=\norm{z_t-z(\eta t)}_1$ and suppose that $e_s\le1$ for $s\le t$, as holds for $t=0$. Then $z_t$ lies in the $\ell_1$-ball of radius $1$ around $z(\eta t)$, a convex subset of $K$, and the two bounds give $e_{t+1}\le\norm{F(z_t)-F(z(\eta t))}_1+\norm{F(z(\eta t))-z(\eta t+\eta)}_1\le(1+L\eta)e_t+C_1\eta^2$. With $e_0=0$ this gives, by induction,
\[
  e_{t+1}\le C_1\eta^2\sum_{s=0}^{t}(1+L\eta)^s=\frac{C_1\eta}{L}\bigl((1+L\eta)^{t+1}-1\bigr)\le\frac{C_1\eta}L\bigl(e^{L\eta(t+1)}-1\bigr),
\]
which is at most $1$ as long as $\eta(t+1)\le\tau_{\max}$, by the assumption on $\eta$. So the induction continues up to $t=\tau_{\max}/\eta$. Finally, $z_t\in K$ gives $1+\eta y_t\ge\frac12$, so no step sends $x$ to zero.
\end{proof}

\begin{proof}[Proof of \cref{thm:eos:selfstab}(iv)]
Apply \cref{lem:eos:limit} with $C=(C_1/L)(e^{L\tau_{\max}}-1)$ and $\eta_0=\min\{1/(2Y),\,1/C\}$, and note that $\log x_t^2=A_t$ by \eqref{eq:eos:Amap}.
\end{proof}

\subsection{What the discrete map does}

The theorem is about a limit. For the map itself, three identities hold exactly, at every step size.

\begin{lemma}[Exact identities]\label[lemma]{lem:eos:identities}
Let $(x_t,y_t)$ follow \eqref{eq:eos:cubicmap} from $x_0\ne0$, and suppose that $1+\eta y_t>0$ for $0\le t<T$. Write $A_t=\log x_t^2$. Then:
\begin{enumerate}
\item $\displaystyle\frac1T\sum_{t<T}x_t^2=2\alpha-\frac{2(y_T-y_0)}{\eta T}$;
\item $\displaystyle\frac1T\sum_{t<T}\log(1+\eta y_t)=\frac{A_T-A_0}{2T}$, and hence $\displaystyle\frac1T\sum_{t<T}y_t\ge\frac1\eta\bigl(e^{(A_T-A_0)/(2T)}-1\bigr)$;
\item for $t<T$, with $f_t=\alpha-\frac12x_t^2$,
\begin{equation}\label{eq:eos:dH}
  \cH(A_{t+1},y_{t+1})-\cH(A_t,y_t)=\eta^2f_t^2+\tfrac12\eta^2x_t^2y_t^2+2\alpha\bigl(\eta y_t-\log(1+\eta y_t)\bigr)\ \ge\ 0,
\end{equation}
with equality if and only if $x_t^2=2\alpha$ and $y_t=0$.
\end{enumerate}
\end{lemma}

\begin{proof}
By \eqref{eq:eos:cubicmap}, $x_{t+1}^2=(1+\eta y_t)^2x_t^2>0$ for $t<T$, so every $A_t$ with $t\le T$ is defined and \eqref{eq:eos:Amap} holds. Summing the $y$-update over $t<T$ gives $y_T-y_0=\eta\sum_{t<T}(\alpha-\frac12x_t^2)$, which is (i). Summing the $A$-update gives the identity in (ii), and since $\log$ is concave, $\log\bigl(\frac1T\sum_t(1+\eta y_t)\bigr)\ge\frac1T\sum_t\log(1+\eta y_t)$, which is the inequality. For (iii), drop the index $t$, and write $u=\eta y$, so that $e^{A_{t+1}}=(1+u)^2e^A$ and $A_{t+1}-A=2\log(1+u)$. Expanding,
\[
  \cH(A_{t+1},y_{t+1})-\cH(A,y)=2\eta yf+\eta^2f^2+\tfrac12e^A\bigl(2u+u^2\bigr)-2\alpha\log(1+u).
\]
The terms of first order combine: $2\eta yf+e^Au=2u(\alpha-\frac12e^A)+e^Au=2\alpha u$. What remains is $\eta^2f^2+\frac12u^2e^A+2\alpha(u-\log(1+u))$, which is \eqref{eq:eos:dH}. Each term is nonnegative, the last because $u-\log(1+u)\ge0$ for $u>-1$ with equality only at $u=0$. The sum vanishes if and only if $f=0$ and $u=0$, that is, $x^2=2\alpha$ and $y=0$.
\end{proof}

Identity (i) is the discrete form of $\langle x^2\rangle=2\alpha$, and it holds along every orbit with no limit taken: the squared oscillation averages $2\alpha$ up to a boundary term that is small whenever $y$ stays bounded. Identity (ii) says that the geometric mean of $\abs{x_{t+1}/x_t}=\eta S_t-1$, where $S_t=2/\eta+y_t$ is the curvature along the oscillation, is exactly $\abs{x_T/x_0}^{1/T}$; if the oscillation neither grows nor decays, the arithmetic mean of $\eta S_t$ is at least $2$, by the concavity of the logarithm. This is \cref{obs:eos:gd}(ii)'s ``at or just above $2/\eta$''. Identity (iii) is new information: the energy that the flow conserves increases at every step of the map. The cycles of gradient descent are not closed.

\begin{proposition}[Gradient descent spirals outward]\label[proposition]{prop:eos:drift}
Let $\eta>0$ and $\alpha>0$, and let $F$ be the map \eqref{eq:eos:Amap} on the half-plane $\{1+\eta y>0\}$.
\begin{enumerate}
\item The point $p=(A^*,0)$, the orbit $x=\pm\sqrt{2\alpha}$, $y=0$ of period two of \eqref{eq:eos:cubicmap}, is a fixed point of $F$, and the Jacobian of $F$ at $p$ has the eigenvalues $1\pm i\eta\sqrt{2\alpha}$, of modulus $\sqrt{1+2\alpha\eta^2}>1$.
\item Along every orbit of $F$, $\cH$ is nondecreasing, and it increases strictly at every step taken from a point other than $p$.
\item For every compact set $K$ in the half-plane, every orbit that starts in $K$ at a point other than $p$ leaves $K$.
\end{enumerate}
\end{proposition}

\begin{proof}
(i) At $p$ we have $1+\eta y=1$ and $e^A=2\alpha$, so $F(p)=p$. The Jacobian of $F$ is $\bigl(\begin{smallmatrix}1&2\eta/(1+\eta y)\\-\frac12\eta e^A&1\end{smallmatrix}\bigr)$, which at $p$ is $\bigl(\begin{smallmatrix}1&2\eta\\-\eta\alpha&1\end{smallmatrix}\bigr)$, with characteristic polynomial $(1-\lambda)^2+2\alpha\eta^2$. (ii) is \cref{lem:eos:identities}(iii). (iii) Suppose an orbit $z_t$ starts at $z_0\in K\setminus\{p\}$ and stays in $K$. The set $K'=K\cap\{\cH\ge\cH(z_0)\}$ is compact, contains every $z_t$ by (ii), and does not contain $p$, because $\cH(p)=\eosHmin<\cH(z_0)$. The increment $\cH(F(z))-\cH(z)$ is continuous and positive on $K'$, so it is at least some $c>0$ there, and $\cH(z_t)\ge\cH(z_0)+ct$ for all $t$, which contradicts the boundedness of $\cH$ on $K$.
\end{proof}

So the period-two orbit at the edge, the oscillation that exactly balances progressive sharpening, is unstable for every step size, and every other orbit drifts outward to larger cycles. The drift is slow, and its rate can be computed.\sidenote{Explicit time stepping does the same to every oscillator. For $\ddot q=-\omega^2q$ written as a first-order system, one Euler step multiplies the energy by exactly $1+\eta^2\omega^2$. In \eqref{eq:eos:Amap} the two variables are updated from the same old values, as in the Euler method.}

\begin{heuristic}[Heuristic derivation of the area law]
Let $\eosArea(h)$ be the area enclosed by the level curve $\{\cH=h\}$ in the $(A,y)$-plane. Over one cycle at level $h$, gradient descent stays within $O(\eta)$ of the cycle (\cref{lem:eos:limit}) and takes about $T(h)/\eta$ steps, each of which raises $\cH$ by $\eta^2P+O(\eta^3)$, with $P=f^2+y^2(\alpha+\frac12e^A)$, by expanding \eqref{eq:eos:dH} in $\eta$. Replacing the sum by an integral, the level rises by $\eta\oint P\,d\tau$ per cycle. On the cycle $dy=f\,d\tau$ and $dA=2y\,d\tau$, so $\oint P\,d\tau=\oint f\,dy+\frac12\oint(\alpha+\frac12e^A)\,y\,dA$. The cycle runs clockwise ($A$ increases where $y>0$), and Green's theorem turns the two line integrals into $\iint\frac12e^A$ and $\frac12\iint(\alpha+\frac12e^A)$ over the region $\{\cH\le h\}$. In that region $\iint(\alpha-\frac12e^A)\,dA\,dy=0$, because $\alpha-\frac12e^A=-\partial_A\cH$ and $\cH=h$ at both ends of every horizontal chord, so the average of $e^A$ over the region is $2\alpha$, as over the cycle, and $\oint P\,d\tau=2\alpha\eosArea(h)$. Finally $\eosArea(h)=\int_{A_-}^{A_+}2\sqrt{h-\phi(A)}\,dA$ has the derivative $\eosArea'(h)=T(h)$. So per cycle $\eosArea$ grows by $T(h)\cdot2\alpha\eta\eosArea(h)$, which is $2\alpha\eta^2\eosArea$ per step:
\begin{equation}\label{eq:eos:area}
  \eosArea\bigl(\cH(z_t)\bigr)\approx\eosArea\bigl(\cH(z_0)\bigr)\,e^{2\alpha\eta^2t}.\qedhere
\end{equation}
\end{heuristic}

Near the period-two orbit the area law is exact to first order, by \cref{prop:eos:drift}(i): the linearization multiplies the area of a small ellipse by its determinant $1+2\alpha\eta^2$ at every step. In the units of the theorem the drift takes the time $1/(2\alpha\eta)$ in $\tau$ to grow an area by the factor $e$, while a small cycle lasts $2\pi/\sqrt{2\alpha}$; the drift therefore needs about $1/(2\sqrt2\,\pi\eps)$ cycles, roughly $11$ cycles at $\eps=0.01$. It is invisible over a few cycles and dominant over many.\sidenote{The area $\eosArea$ is $2\pi$ times the action of the cycle in the sense of classical mechanics. Actions are conserved under slow changes of the parameters (adiabatic invariance), which is why the slow variation of $\alpha$ along a training run does not by itself change the size of the cycles, while the discretization does.}

\subsection{What Damian, Nichani and Lee prove}

Their theorem concerns gradient descent on a general loss, not the cubic model. Its assumptions include progressive sharpening along the constrained trajectory (their Assumption 1), a gap below the top eigenvalue, with the second eigenvalue of the Hessian below $c/\eta$ for an absolute constant $c<2$ (their Assumption 2), and bounds on the third derivative of the loss and on its Lipschitz constant. The constrained trajectory is projected gradient descent onto the stable set $\{\theta:S(\theta)\le2/\eta,\ \ip{\nabla\Loss(\theta),u(\theta)}=0\}$, with $u(\theta)$ the top eigenvector, and they show that gradient descent can be coupled with it for $\cT=O(\eps^{-1})$ steps, $\eps$ being the small parameter of their analysis, with the deviation captured by a two-dimensional system of the form \eqref{eq:eos:cubicmap} with time-dependent coefficients. The coupling assumes that the oscillation does not become too small: $\min_{t\le\cT}\abs{x_t}\ge c_1\delta$ for a constant $c_1$ and a scale $\delta$ of their analysis. They note that these conditions are not strictly necessary and that a local coupling of this kind is not sufficient to prove global convergence.

Two features of the statement matter here. The horizon is finite, and the lower bound on $\abs{x_t}$ excludes the long stretches of a large cycle in which the oscillation has almost died out. \Cref{prop:eos:drift} and \eqref{eq:eos:area} concern what lies beyond: a time of order $1/(\alpha\eta^2)$ steps, over which the reduced dynamics itself, exact for the cubic model, is not stable. A coupling theorem over a bounded number of cycles cannot decide whether real dynamics at the edge drifts outward, settles onto a stable period-two orbit, or does something else; that is decided by the terms that the reduction drops, and \cref{sec:eos:misses} shows how.

\subsection{Central flows}\label{sec:eos:central}

Average over the cycles. In the family \eqref{eq:eos:onedir}, suppose the curvature $s(w_t)$ has reached $2/\eta$ and the functions $\Loss_0$ and $s$ change little over one cycle. Then, over a cycle, the pair $(\sqrt{\gamma(w)}\,x,\ s(w)-2/\eta)$ follows the cubic model with the local rate $\alpha(w)$, and the theorem gives the cycle averages.

\begin{heuristic}[Heuristic derivation of the central flow]
If $\alpha(w)>0$, the squared oscillation averages $\langle x^2\rangle=2\alpha(w)/\gamma(w)$ over a cycle, by \cref{thm:eos:selfstab}(iii) for the rescaled pair. Averaging the $w$-update of \eqref{eq:eos:gdonedir} over the cycle, and writing $\tau=\eta t$, gives
\begin{equation}\label{eq:eos:central}
  \frac{dw}{d\tau}=-\nabla\Loss_0(w)-\tfrac12\sigma^2(w)\,\nabla s(w),\qquad
  \sigma^2(w)=\frac{2\max\{0,\,-\ip{\nabla\Loss_0(w),\nabla s(w)}\}}{\norm{\nabla s(w)}_2^2},
\end{equation}
while $s(w)=2/\eta$; when $\alpha(w)\le0$ the curvature falls below $2/\eta$, the oscillation decays, and $w$ follows gradient flow. With $\sigma^2=2\alpha/\gamma$, the velocity in \eqref{eq:eos:central} is $-\nabla\Loss_0$ minus its component along $\nabla s$: it is tangent to the surface $\{s=2/\eta\}$. So \eqref{eq:eos:central} is gradient flow projected onto the region $\{s\le2/\eta\}$, and the average of the loss over a cycle is $\Loss_0(w)+\frac12s\langle x^2\rangle=\Loss_0(w)+\sigma^2/\eta$.
\end{heuristic}

\Textcite{CohenEtAl2025} derived the same equation for a general loss and called it a \term{central flow}\index{central flow}: a differential equation for the time-averaged trajectory of gradient descent at the edge of stability, with $\Loss$ and $S$ in place of $\Loss_0$ and $s$. In their words, $\sigma^2$ is the unique value for which the downward force of oscillation-induced sharpness reduction cancels the upward force of progressive sharpening, and the flow projects the direction $\nabla S$ out of the gradient to keep the sharpness at $2/\eta$; they note that this re-derives a main result of \textcite{DamianNichaniLee2023} in a simpler but less rigorous way, and they extend the construction to several eigenvalues at the edge. The status is the same in their paper and here: derived, not proved. Its predictions are nevertheless quantitative: they report that central flows predict long-term trajectories of neural networks with high numerical accuracy. In the cubic model the central flow is the point $y=0$ with $\sigma^2=2\alpha$, which is \cref{thm:eos:selfstab}(iii).

For adaptive optimizers the central flow also moves the preconditioner, and the oscillation acts through it. \Textcite{CohenEtAl2025} derive central flows for RMSProp and Adam, and find what they call acceleration via regularization: the adaptive optimizers implicitly move toward regions of low curvature, in which they can take larger steps.

\section{Reading the theory}\label{sec:eos:reading}

\paragraph{Thresholds in numbers.} \Cref{prop:eos:thresholds} turns the defaults of practice into numbers. For the moving-average form of momentum, the largest stable $\eta S$, or $\eta S_P$ for Adam with a frozen preconditioner, is $(2+2\beta)/(1-\beta)$:
\begin{center}
\small
\begin{tabular}{@{}lccccc@{}}
$\beta$ (or $\beta_1$) & $0$ & $0.5$ & $0.9$ & $0.95$ & $0.99$\\
\hline
largest stable $\eta S$ & $2$ & $6$ & $38$ & $78$ & $398$
\end{tabular}
\end{center}
The left panel of \cref{fig:eos:network} shows full-batch gradient descent on a small network: the sharpness rises, overshoots, and then holds at $2/\eta$, within $1.4\%$ on average over the last $1500$ steps for $\eta=0.4$ and within $0.1\%$ for $\eta=0.25$. The right panel shows Adam on the same network: the preconditioned sharpness climbs to $38/\eta$ and stays there, with averages $0.98$ and $0.95$ of the threshold over the last $1500$ steps and excursions to about $1.03$ times it. The climb is close to the limit of \cref{prop:eos:speed} for a fixed Hessian, because the gradients are already small when it happens: for $\eta=0.003$ the preconditioned sharpness grows by the factor $e^{1.80}$ between steps $200$ and $1975$, where the limit allows $e^{1.67}$, and the slow growth of the Hessian itself makes up the difference. Meanwhile the sharpness of Adam's iterates stays between $7.4$ and $8.4$ from step $300$ on, in both runs, far below any threshold of its own: as \cref{obs:eos:adam} says, what Adam holds at its edge is the preconditioned sharpness.

\begin{figure}[t]
\begin{fullwidth}
\centering
\begin{tikzpicture}
\begin{groupplot}[group style={group size=2 by 1,horizontal sep=17mm},wplot,width=82mm,height=50mm,xmin=0,xmax=4000,xtick={0,1000,2000,3000,4000}]
\nextgroupplot[xlabel={step $t$},ylabel={sharpness $S(\theta_t)$},ymin=0,ymax=10,legend style={at={(0.98,0.04)},anchor=south east,font=\scriptsize},legend columns=2]
\addplot[theory,domain=0:4000,samples=2] {5};
\addplot[theory2,domain=0:4000,samples=2] {8};
\addplot[simulation,mark size=1.0pt] table[x=t,y=gdA] {figures/data/08-eos-network.dat};
\addplot[simulation2,mark size=0.9pt] table[x=t,y=gdB] {figures/data/08-eos-network.dat};
\legend{$2/\eta$\,($\eta=0.4$),$2/\eta$\,($\eta=0.25$),$\eta=0.4$,$\eta=0.25$}
\nextgroupplot[xlabel={step $t$},ylabel={preconditioned sharpness $S_P$},ymode=log,ymin=300,ymax=30000,legend style={at={(0.98,0.04)},anchor=south east,font=\scriptsize},legend columns=2]
\addplot[theory,domain=0:4000,samples=2] {3800};
\addplot[theory2,domain=0:4000,samples=2] {12666.67};
\addplot[simulation,mark size=1.0pt,restrict x to domain=25:4000] table[x=t,y=adA] {figures/data/08-eos-network.dat};
\addplot[simulation2,mark size=0.9pt,restrict x to domain=25:4000] table[x=t,y=adB] {figures/data/08-eos-network.dat};
\legend{$38/\eta$\,($\eta=0.01$),$38/\eta$\,($\eta=0.003$),$\eta=0.01$,$\eta=0.003$}
\end{groupplot}
\end{tikzpicture}
\end{fullwidth}
\caption{The thresholds of \cref{prop:eos:thresholds} (lines) and the measured curvature (points) in a tanh network $8\to32\to32\to1$ ($1377$ parameters) trained full batch on the square loss, with $256$ Gaussian inputs labelled by a random teacher (seeds fixed in \texttt{08-eos-network.py}). Left: gradient descent; the sharpness, estimated every $25$ steps by block power iteration with Hessian--vector products. Right: Adam with $\beta_1=0.9$, $\beta_2=0.999$, $\epsilon=10^{-8}$; the preconditioned sharpness $\lmax(P_t^{-1}H_t)$ of the preconditioner $P_t$ used in the update, from step $25$ on. From step $300$ on, the sharpness of Adam's own iterates stays between $7.4$ and $8.4$ in both runs.}
\label{fig:eos:network}
\end{figure}

\paragraph{The cubic model, cycle by cycle.} \Cref{fig:eos:cubic} compares the theorem with the map. On the left, gradient descent with $\alpha=1$ and $\eta=0.002$ follows the closed level curves of $\cH$ for two cycles each: a small one, started at the edge with $y_0=0.5$, and a large one, started below the edge ($y_0=-0.5$) with a small oscillation $x_0=0.01$, as at the onset of the edge of stability. The large cycle is a sequence of catapults: the curvature climbs to $2/\eta+3.0$, the oscillation grows from $0.01$ to about $5$, the curvature is pushed down to $2/\eta-3.0$, and the oscillation dies back to below $0.01$. The overshoot equals the undershoot because $\cH$ is even in $y$; with $\alpha=0$ the same conservation law turns a single excursion to $2/\eta+y_0$ into a landing at $2/\eta-y_0$, the reflection law that volume 2 of this series, Chapter 15, derives for the catapult. A small oscillation at onset means a large cycle: $\cH\approx y_0^2-\alpha\log x_0^2$ is logarithmic in the initial amplitude. The small cycle lasts $T=4.54$ in $\tau$, close to the $2\pi/\sqrt2=4.44$ of the linearization; the large one lasts $7.56$, that is, $3780$ steps. The second loop of the large cycle lies visibly outside the first: the drift.

In the middle, the running averages of $x^2$ and $y$ along the large cycle converge to $2\alpha=2$ and $0$, with oscillations that decay like $1/\tau$ and that gradient descent shares with the flow. After five periods gradient descent's average of $x^2$ is $2.04$, and identity (i) of \cref{lem:eos:identities} accounts for the difference exactly: it is the boundary term $-2(y_T-y_0)/(\eta T)$, which is not zero because the drifting orbit has fallen slightly out of phase with the flow. On the right, the area law \eqref{eq:eos:area}: three runs with different $(\alpha,\eta)$, started at $\cH=1$, collapse onto the line $\log(\eosArea_t/\eosArea_0)=2\alpha\eta^2t$ over the plotted range, in which the area grows by the factor $e^6$; beyond it, $\eta y$ is no longer small and the expansion in $\eta$ fails.

\begin{figure}[t]
\begin{fullwidth}
\centering
\begin{tikzpicture}
\begin{groupplot}[group style={group size=3 by 1,horizontal sep=14mm},wplot,width=58mm,height=50mm]
\nextgroupplot[xlabel={$A=\log x^2$},ylabel={$y=S-2/\eta$},xmin=-11,xmax=4.5,ymin=-3.6,ymax=3.6]
\addplot[theory] table[x=a2,y=y2] {figures/data/08-eos-cubic-levels.dat};
\addplot[theory] table[x=a1,y=y1] {figures/data/08-eos-cubic-levels.dat};
\addplot[simulation,mark size=0.6pt,unbounded coords=discard] table[x=a2,y=y2] {figures/data/08-eos-cubic-phase.dat};
\addplot[simulation2,mark size=0.6pt,unbounded coords=discard] table[x=a1,y=y1] {figures/data/08-eos-cubic-phase.dat};
\nextgroupplot[xlabel={time $\tau=\eta t$},ylabel={running averages},xmin=0,xmax=38,ymin=-1.2,ymax=4.2,legend style={at={(0.98,0.97)},anchor=north east,font=\scriptsize}]
\addplot[guide,domain=0:38,samples=2] {2};
\addplot[guide,domain=0:38,samples=2] {0};
\addplot[theory] table[x=tau,y=flowx2] {figures/data/08-eos-cubic-avg.dat};
\addplot[theory2] table[x=tau,y=flowy] {figures/data/08-eos-cubic-avg.dat};
\addplot[simulation,mark size=0.7pt] table[x=tau,y=gdx2] {figures/data/08-eos-cubic-avg.dat};
\addplot[simulation2,mark size=0.6pt] table[x=tau,y=gdy] {figures/data/08-eos-cubic-avg.dat};
\legend{,,{flow, $x^2$},{flow, $y$},{GD, $x^2$},{GD, $y$}}
\nextgroupplot[xlabel={$2\alpha\eta^2t$},ylabel={$\log(\eosArea_t/\eosArea_0)$},xmin=0,xmax=6.2,ymin=0,ymax=6.4,legend style={at={(0.03,0.97)},anchor=north west,font=\scriptsize}]
\addplot[theory,domain=0:6.2,samples=2] {x};
\addplot[simulation,mark size=1.1pt] table[x=s1,y=r1] {figures/data/08-eos-cubic-area.dat};
\addplot[simulation2,mark size=1.0pt] table[x=s2,y=r2] {figures/data/08-eos-cubic-area.dat};
\addplot[only marks,mark=triangle*,mark size=1.4pt,color=sea,mark options={fill=sea,draw=white,line width=0.2pt}] table[x=s3,y=r3] {figures/data/08-eos-cubic-area.dat};
\legend{area law,{$\alpha=1$, $\eta=0.01$},{$\alpha=1$, $\eta=0.005$},{$\alpha=\frac12$, $\eta=0.02$}}
\end{groupplot}
\end{tikzpicture}
\end{fullwidth}
\caption{The cubic model: theory (curves) against the gradient-descent map \eqref{eq:eos:cubicmap} (points); no randomness (script \texttt{08-eos-cubic.py}). Left: two cycles each of a small and a large orbit, $\alpha=1$, $\eta=0.002$, every $25$th iterate, on the level curves of $\cH$ through their starting points. Middle: running averages of $x^2$ and of $y$ along the large orbit over five periods, for the flow \eqref{eq:eos:flow} and for gradient descent. Right: the area enclosed by the level curve through the current iterate, for three runs started on $y=0$ at $\cH=1$, against the area law \eqref{eq:eos:area}.}
\label{fig:eos:cubic}
\end{figure}

\paragraph{Adam's edge without sharpening.} The preconditioned sharpness can reach its threshold even when the Hessian does not change at all, because the preconditioner shrinks as the gradients do.

\begin{example}[Adam on a parabola]\label[example]{ex:eos:parabola}
Run Adam with $\beta_1=0.9$, $\beta_2=0.999$ and $\epsilon=0$ on $\Loss(x)=\frac12hx^2$. The iterate converges toward zero, the gradients become small next to $\sqrt{v_t}$, and the preconditioned sharpness $h/P_t$ rises at the limit of \cref{prop:eos:speed}. In a run with $h=1$, $\eta=0.01$ and $x_0=1$ it grows from $\eta S_P=0.066$ at step $1000$ to the threshold $38$ at step $13{,}254$, which is the rate $\frac12\log(1/\beta_2)$ per step corrected for the bias correction of $\hat v$. It then oscillates between about $31$ and $40$, around an average of $35.4$, with $\abs{x}$ of order $0.03\,\eta$. Adam with $\epsilon=0$ is invariant under rescaling the loss, and rescaling $x$ rescales $\eta$; so the range and the average of $\eta S_P$ do not depend on $h$ or on $\eta$, which the runs confirm.
\end{example}

The example separates two mechanisms that the network run mixes. Adam can sharpen through its preconditioner alone, at a speed that $\beta_2$ fixes, and its self-stabilization then runs through $v_t$ rather than through the third derivative of the loss: a growing oscillation raises the squared gradients, $v_t$ grows, and $S_P$ falls below the threshold. The average sits about $7\%$ below the threshold, unlike the cubic model's exact $2/\eta$, because the oscillation grows fast above the threshold but decays slowly below it: the two roots of the momentum recursion have the product $\beta_1$, so the amplitude is multiplied by at least $\sqrt{\beta_1}$ per step however far below the threshold $\eta S_P$ is, and the cycle spends most of its time below (\cref{ex:eos:adam1d}).

\paragraph{Where the agreement must fail.} The theorem is a statement about $\eps\to0$ with the energy fixed. Large cycles fail first: their curvature excursion $\sqrt{h-\eosHmin}$ grows with the energy, and once $\eta y$ is no longer small the expansion behind \eqref{eq:eos:area} breaks down. In the run that the right panel of \cref{fig:eos:cubic} does not show, with $\eta=0.01$, $\alpha=1$ and $\cH_0=10$, the drift carries the orbit to cycles in which $x$ falls below $10^{-300}$; in floating point $x$ underflows to zero after about $4\cdot10^4$ steps, after which the iterates never oscillate again and the curvature grows without bound. In exact arithmetic the cycles keep growing until one of them bursts to an oscillation of astronomical size. Neither is what networks do, and \cref{sec:eos:misses} says why.

\section{What the theory misses}\label{sec:eos:misses}

\paragraph{What was idealized away.} The rate of progressive sharpening $\alpha$ is an input, constant in the cubic model and a given function of $w$ in \eqref{eq:eos:onedir}; why it is positive in networks is not addressed. One direction is at the edge, with a fixed eigenvector and a spectral gap below it, while in networks the top eigenvector turns and several eigenvalues can reach $2/\eta$ together. The loss is exactly quadratic and even in the oscillating coordinate: no terms $x^3$ or $x^4$, which are precisely the terms that decide the long-time behaviour below. The gradients are exact. And the thresholds for momentum and Adam are proved for quadratics with a frozen preconditioner, while the self-stabilization theory is proved for plain gradient descent only.

\paragraph{Where it breaks.} (1) \emph{Long times.} \Cref{prop:eos:drift} shows that the cubic model has no stable cycle: its period-two orbit repels at the rate $\alpha\eta^2$ per step. A fourth-order term $\frac1{24}\kappa x^4$ added to the loss changes the continuous limit at order one: the equations become $dA/d\tau=2y+\frac13\kappa e^A$ and $dy/d\tau=\alpha-\frac12e^A$, and for $\kappa<0$, a loss that grows more slowly than quadratically along the oscillation, the cycles spiral in to a stable period-two orbit with $y=-\kappa\alpha/3>0$, a sharpness slightly above $2/\eta$ (\cref{ex:eos:quartic}). Near the orbit the area rate becomes $2\alpha(\eta+\kappa/3)$ per unit of $\tau$, so the orbit is stable when $\kappa<-3\eta$, a prediction that the exact map confirms. A cubic term $\frac16\kappa_3x^3$ contributes $-\frac12\eta\kappa_3^2e^A$ to $dA/d\tau$, which also damps. In one dimension the combination is the condition $3\Loss'''^2-\Loss''\Loss''''>0$ for a stable period-two orbit near the threshold (\cref{ex:eos:quartic}); \textcite{ChenBruna2023} analyse such orbits through the two-step map. Whether a network's cycles are damped is therefore a property of its fourth derivative along the top eigenvector, which the cubic model sets to zero; the sustained edge of \cref{obs:eos:gd}, hovering ``at or just above'' $2/\eta$, is what damping predicts. (2) \emph{Progressive sharpening.} The model does not say why the sharpness rises, and with the cross-entropy loss it eventually falls \parencite{CohenEtAl2021}; for logistic regression on separable data, gradient descent converges at every step size after a phase in which the loss oscillates \parencite{WuBravermanLee2023,WuEtAl2024}, so there the edge is left for good. (3) \emph{Several eigenvalues.} When two eigenvalues are at the edge, the averaging identity of \cref{lem:eos:identities}(i) becomes a linear system for the two oscillation energies, which can have a negative solution; then one direction cannot stay at the edge (\cref{ex:eos:two}). The single-eigenvalue picture, and the eigengap assumption of \textcite{DamianNichaniLee2023}, fail exactly there. (4) \emph{Noise.} With minibatches the full-batch sharpness stays below $2/\eta$ (\cref{obs:eos:sgd}); the model has no noise and predicts the full-batch threshold. (5) \emph{Adam's moving preconditioner.} \Cref{ex:eos:parabola} shows that Adam reaches its edge and stabilizes there through $v_t$ alone, with an average $7\%$ below the frozen-preconditioner threshold; the cubic mechanism and this one act together in a network, and the theory of this chapter covers neither their interaction nor the observation that Adam keeps moving into regions of higher sharpness \parencite{CohenEtAl2022}.

\paragraph{Repairs and open problems.} Central flows \parencite{CohenEtAl2025} handle several eigenvalues at the edge and adaptive optimizers, and predict long training trajectories, at the level of a derivation. Rod flows \parencite{RegisChewi2026} model gradient descent as an extended object with a centre and an extent; their authors report that rod flow matches central flows on representative networks, contains central flows as a Taylor approximation, and exhibits self-stabilization in quartic potentials, and a second paper extends the construction to Adam, heavy ball, Nesterov and RMSProp \parencite{RegisChewi2026Adam}; both are unrefereed preprints of 2026. Minimalist models prove edge-of-stability behaviour for all time in low dimension: a two-parameter example \parencite{ZhuEtAl2023}, the learning of threshold neurons \parencite{AhnEtAl2023}, normalized gradient descent and a flow that reduces sharpness \parencite{AroraLiPanigrahi2022}, scalar networks \parencite{KreislerEtAl2023}, and a bifurcation analysis that aligns trajectories \parencite{SongYun2023}; \textcite{KalraHeBarkeshli2025} analyse the fixed points of a low-dimensional model and its route to chaos. Progressive sharpening has partial explanations, through the growth of the output layer \parencite{LiWangLi2022} and in second-order regression models \parencite{AgarwalaPedregosaPennington2023}. For minibatches, the batch sharpness of \textcite{AndreyevBeneventano2024} and the sharpness measure of \textcite{LeeJang2023} replace $S$. Open: a model in which progressive sharpening is derived and the long phase at the edge is proved; a global-in-time theorem for gradient descent at the edge beyond low-dimensional examples; the sign of the damping in networks, which can be measured (\cref{ex:eos:drift}); and a theory of Adam's edge with the preconditioner moving.

\section{Back in the wild}\label{sec:eos:back}

\fieldnotes{Learning rates chosen as large as training tolerates; full-batch gradient descent then trains with the sharpness at $2/\eta$, Adam with the preconditioned sharpness at $38/\eta$.}
{One oscillating direction whose curvature depends on the other parameters; its affine case, the cubic model $\frac12(2/\eta+y)x^2-\alpha y$.}
{Thresholds $2/\eta$, $(2+2\beta)/\eta$, $(2+2\beta)/((1-\beta)\eta)$, and $38/\eta$ for Adam with a frozen preconditioner. In the continuous limit, closed cycles with $\langle y\rangle=0$ and $\langle x^2\rangle=2\alpha$; for the map, exact averaging identities and an energy that increases at every step.}
{The sharpness averages $2/\eta$, or slightly more; the oscillation carries $\langle x^2\rangle=2\alpha/\norm{\nabla S}^2$; small cycles last $2\pi/(\eta\sqrt{2\alpha})$ steps; cycles grow by the factor $e$ in area every $1/(2\alpha\eta^2)$ steps unless higher-order terms damp them; the central flow (derived).}
{Why the sharpness rises; several eigenvalues at the edge; noise; a moving preconditioner; the fourth-order terms that decide whether the cycles are damped.}
{The step size sets the curvature, not the reverse (observed; explained in the model). Compare learning rates across momentum values through $(2+2\beta)/(1-\beta)$ (proved for quadratics). The oscillation costs about $\sigma^2/\eta$ of loss, which averaging two consecutive iterates removes (derived in the model; suggested).}

For a practitioner the first lesson is a change of viewpoint. The learning rate is not chosen to respect a curvature; the curvature adapts to the learning rate, rising until the step is barely stable and then staying there. A sweep that keeps the largest learning rate that does not diverge is therefore not choosing a point safely inside a fixed stability region: it is choosing the curvature at which training will run. Divergence happens when self-stabilization fails, for example when the step size grows faster than the curvature can follow, which is the subject of warmup (\cref{ch:warmup}), or when one cycle grows too large, which is one of the anatomies of a loss spike (\cref{ch:spikes}). The second lesson is quantitative but modest. The thresholds tell how momentum moves the edge: at $\beta=0.9$ the moving-average form tolerates a curvature nineteen times larger than plain gradient descent at the same $\eta$, so a learning rate tuned for one $\beta$ is not the right one for another; \cref{ch:momentum} proves the regions and their noise counterparts. And the oscillation is not a malfunction: it carries exactly the energy that progressive sharpening requires, and its cost in loss, about $\sigma^2/\eta$ in the model, is paid by the iterates and not by their average. In the model, the average of two consecutive iterates, $\frac12(x_t+x_{t+1})=-\frac12\eta y_tx_t$, has almost no oscillation, so its loss is lower by nearly $\sigma^2/\eta$; this is a reason to expect averaging of iterates, discussed in \cref{ch:momentum,ch:schedules}, to help at the edge, but it is a suggestion, not a result about networks.

What the theory does not tell a practitioner is equally clear. It does not say why the sharpness rises, so it cannot predict how long a run takes to reach the edge or where it goes along it, beyond the central flow's projection. With noise the edge moves (\cref{ch:noise}), and with Adam the edge involves the preconditioner, whose own dynamics the frozen threshold ignores; why Adam beats gradient descent on transformers is a separate question with several answers (\cref{ch:transformers}).

\begin{exercises}
\item \consolidation (Nesterov's threshold.) Nesterov's momentum is $\theta_{t+1}=\vartheta_t-\eta\nabla\Loss(\vartheta_t)$ with $\vartheta_t=\theta_t+\beta(\theta_t-\theta_{t-1})$. (a) On a quadratic with curvature $\lambda>0$ along an eigenvector, show that the component obeys $e_{t+1}=(1+\beta)(1-\eta\lambda)e_t-\beta(1-\eta\lambda)e_{t-1}$. (b) With the stability test, show that it is stable for every initialization if and only if $0<\eta\lambda<(2+2\beta)/(1+2\beta)$. (c) Compare with the heavy ball for $\beta=0.9$, and with the maximal stable sharpness of Nesterov's method reported by \textcite{CohenEtAl2021}.\label{ex:eos:nesterov}

\item \consolidation (Small cycles.) (a) Show that the linearization of \eqref{eq:eos:flow} at $(A^*,0)$ has the angular frequency $\sqrt{2\alpha}$. (b) Substitute $A=A^*+\sqrt{2(h-\eosHmin)/\alpha}\,\sin\vartheta$ in the formula for $T(h)$ and show that $T(h)\to2\pi/\sqrt{2\alpha}$ as $h\downarrow\eosHmin$. (c) For $\alpha=1$, compute $T$ numerically at $h=\cH(A^*,0.5)$ and compare with the value $4.54$ quoted in \cref{sec:eos:reading}.\label{ex:eos:period}

\item \consolidation (The reduction.) In the family \eqref{eq:eos:onedir} with $\Loss_0(w)=\Loss_0(0)+\ip{g,w}$ and $s(w)=2/\eta+\ip{c,w}$, show that $(\norm c_2x_t,\ip{c,w_t})$ follows \eqref{eq:eos:cubicmap} with $\alpha=-\ip{g,c}$, and that the component of $w_t$ orthogonal to $c$ performs gradient descent on $\Loss_0$ restricted to that subspace. Deduce that the central flow \eqref{eq:eos:central} is exact for this family, in the sense that $\langle y\rangle=0$ and the orthogonal component moves at the velocity $-\eta$ times the projected $g$ per step.\label{ex:eos:reduction}

\item \consolidation (Hellmann--Feynman.) Let $H(\theta)$ be a smooth family of symmetric matrices whose top eigenvalue $S(\theta)$ is simple, with unit eigenvector $u(\theta)$. Differentiate $Hu=Su$ and $u^\top u=1$ to show that $dS=u^\top(dH)\,u$, so that $\nabla S=\nabla^3\Loss[u,u,\cdot\,]$ when $H=\nabla^2\Loss$. Why does the formula fail when the top eigenvalue is not simple?\label{ex:eos:hellmann}

\item \consolidation (Two consecutive iterates.) In the cubic model, show that $\Loss(\frac12(x_t+x_{t+1}),y_t)=\frac18\eta^2y_t^2x_t^2(2/\eta+y_t)-\alpha y_t$, and compare with $\Loss(x_t,y_t)$. Averaging over a cycle in the continuous limit, show that the loss of the iterates exceeds that of the averages of consecutive iterates by $2\alpha/\eta+O(1)$.\label{ex:eos:average}

\item \variation (A quartic term.) Add $\frac1{24}\kappa x^4$ to the cubic model. (a) Show that gradient descent becomes $x_{t+1}=-(1+\eta y_t)x_t-\frac16\eta\kappa x_t^3$, with the same $y$-update. (b) Derive the continuous limit $dA/d\tau=2y+\frac13\kappa e^A$, $dy/d\tau=\alpha-\frac12e^A$, and show that its fixed point is $(\log2\alpha,-\kappa\alpha/3)$, a stable focus if $\kappa<0$ and an unstable one if $\kappa>0$. (c) Show that $d\cH/d\tau=\frac13\kappa e^A(\frac12e^A-\alpha)$, and use $\langle e^A\rangle=2\alpha$ and Jensen's inequality to show that for $\kappa<0$ the energy decreases on average over a cycle of the unperturbed flow. (d) Combine with the area law to predict stability of the period-two orbit of the map when $\kappa<-3\eta$, and check numerically with $\eta=0.01$, $\alpha=1$ and $\kappa\in\{-0.01,-0.05\}$. (e) For a one-dimensional loss with $\Loss''=2/\eta$ at a minimum, show that the two-step map of gradient descent is $x\mapsto x-2(k_2^2+k_3)x^3+O(x^4)$ with $k_2=-\frac12\eta\Loss'''$ and $k_3=-\frac16\eta\Loss''''$, so that the period-two orbit born at the threshold is stable when $3\Loss'''^2-\Loss''\Loss''''>0$.\label{ex:eos:quartic}

\item \variation (Momentum does not change the energy.) Run the heavy ball, $\theta_{t+1}=\theta_t-\eta\nabla\Loss(\theta_t)+\beta(\theta_t-\theta_{t-1})$, on the cubic model with $2/\eta$ replaced by $(2+2\beta)/\eta$. Show that $\frac1T\sum_{t<T}x_t^2=2\alpha-2\bigl((y_T-y_0)-\beta(y_{T-1}-y_{-1})\bigr)/(\eta T)$, so that $\langle x^2\rangle=2\alpha$ along every orbit on which $y$ stays bounded, and derive the same for the moving-average form. What does change with momentum?\label{ex:eos:momentum}

\item \variation (Two directions at the edge.) Let $\Loss=\frac12(2/\eta+y_1)x_1^2+\frac12(2/\eta+y_2)x_2^2-\alpha_1y_1-\alpha_2y_2$ arise, as in \cref{ex:eos:reduction}, from $s_j(w)=2/\eta+\ip{c_j,w}$ and $\Loss_0(w)=\ip{g,w}$, so that $y_j=\ip{c_j,w}$ and $\alpha_j=-\ip{g,c_j}$. (a) Show that $y_{j,t+1}=y_{j,t}+\eta\bigl(\alpha_j-\frac12\sum_kG_{jk}x_{k,t}^2\bigr)$ with $G_{jk}=\ip{c_j,c_k}$. (b) Show that if both oscillations stay bounded and $y$ stays bounded, then the time averages satisfy $G\langle x^2\rangle=2\alpha$. (c) Take unit vectors with $\ip{c_1,c_2}=0.8$ and $(\alpha_1,\alpha_2)=(1,0.5)$, for which $G\xi=2\alpha$ has the solution $\xi=(10/3,-5/3)$. With $\eta=0.002$, starting below the edge with small oscillations, show numerically that the second oscillation dies out while the first stays at the edge with $\langle x_1^2\rangle\approx2$; and that with $(\alpha_1,\alpha_2)=(1,0.9)$ both persist, with averages near $\xi=(1.56,0.56)$.\label{ex:eos:two}

\item \simulation (The area law.) Iterate \eqref{eq:eos:cubicmap} from $y_0=0$ and $\cH_0=1$ for $(\alpha,\eta)\in\{(1,0.01),(1,0.005),(\frac12,0.02)\}$, compute the enclosed area $\eosArea$ of the level curve through each iterate by quadrature, and reproduce the right panel of \cref{fig:eos:cubic}. Then start at $\cH_0=10$ with $\eta=0.01$ and $\alpha=1$, and observe that $x$ underflows to zero in floating point after about $4\cdot10^4$ steps. Repeat in the variables $(A,y)$ of \eqref{eq:eos:Amap}, which cannot underflow, and describe what happens to the cycles.\label{ex:eos:area}

\item \simulation (Adam on a parabola.) Reproduce \cref{ex:eos:parabola}. (a) Measure the rate at which $\log S_P$ grows before the threshold and compare it with \cref{prop:eos:speed}, including the bias correction. (b) Repeat with $\beta_2=0.99$ and with $\beta_1=0.95$, whose threshold is $78$, and report the average and range of $\eta S_P$. (c) Explain why the average lies below the threshold: compare the largest rate at which the oscillation can grow when $\eta S_P$ exceeds $(2+2\beta_1)/(1-\beta_1)$ by a few percent with the rate at which it decays when $\eta S_P$ is $20\%$ below.\label{ex:eos:adam1d}

\item \simulation (Momentum in the network.) Modify \texttt{08-eos-network.py} to run Adam with $\eta=0.01$ and $\beta_1\in\{0.8,0.95\}$, and gradient descent with the moving-average momentum of \cref{prop:eos:thresholds}(iii) at $\beta=0.5$ and $\eta=1$. Check that the preconditioned sharpness settles near $18/\eta$ and somewhat below $78/\eta$, and that the sharpness of the momentum run settles at $6$. Then run the moving average with $\beta=0.9$ and $\eta=4.5$, whose threshold $38/\eta\approx8.4$ is three times the initial sharpness $2.7$, and observe that training diverges within twenty steps. Which hypothesis of \cref{prop:eos:thresholds} fails in the first steps? (\Cref{ch:warmup} takes up the question.)\label{ex:eos:network}

\item \research\hard (Is the edge damped?) For a network at the edge of stability, estimate the third and fourth derivatives of the loss along the top eigenvector, $\nabla^3\Loss[u,u,u]$ and $\nabla^4\Loss[u,u,u,u]$, by finite differences of Hessian--vector products, and compute the damping coefficient of \cref{ex:eos:quartic}(e) together with the discretization's anti-damping $2\alpha\eta$, with $\alpha$ estimated from the sharpening rate before the edge. Does the sign of the total predict whether the oscillation settles onto a period-two orbit? Whether the cycles of networks are damped, and by which terms, is open.\hint{the effective coefficient must include the curvature of $S$ along $u$, since the fourth-order term enters through $\frac{d}{dx}$ of the curvature along the oscillation.}\label{ex:eos:drift}

\item \research (Sharpening as an output.) The cubic model is given progressive sharpening as the rate $\alpha$. Find a model in which $\alpha$ is derived from the data and the architecture, and in which the long phase at the edge, with a decreasing loss, can be proved. \Textcite{AgarwalaPedregosaPennington2023} derive progressive sharpening in second-order regression models; what is open is a model in which both the rate and the long phase are solved.\label{ex:eos:sharpening}
\end{exercises}

\notesandrefs
The threshold $2/\eta$ for gradient descent and the analysis of the heavy ball on quadratics are classical \parencite{Polyak1987}. That stability selects the minima that gradient descent can reach was proposed by \textcite{WuMaE2018}; \textcite{JastrzebskiEtAl2019,JastrzebskiEtAl2020} observed that early in training the sharpness adapts to the step size and the batch size; and \textcite{LewkowyczEtAl2020} described the catapult, solved in volume 2 of this series. \Textcite{CohenEtAl2021} named progressive sharpening and the edge of stability, documented both across architectures and losses, and derived the momentum thresholds; \textcite{CohenEtAl2022} found the adaptive edge of stability.

\Textcite{DamianNichaniLee2023} derived the cubic model by a third-order expansion, identified self-stabilization as the mechanism, and proved the coupling with the constrained trajectory discussed in \cref{sec:eos:selfstab}; an earlier version of their paper appeared at a NeurIPS 2022 workshop. \Textcite{CohenEtAl2025} derived central flows, for gradient descent and for adaptive optimizers. Our \cref{thm:eos:selfstab} is the continuous-limit analysis of the cubic model with the proofs written out. We know no reference for the exact identities of \cref{lem:eos:identities}, for the monotone increase of the energy in \cref{prop:eos:drift}, or for the area law \eqref{eq:eos:area}; they are elementary consequences of the model. Stable period-two orbits near the threshold are the flip bifurcation of dynamical systems \parencite[\S4.4--4.5]{Kuznetsov2004}, and \textcite{ChenBruna2023} study them through the two-step map of gradient descent.

Models with complete proofs of edge-of-stability behaviour are listed in \cref{sec:eos:misses}. Large step sizes for logistic regression, where the edge is eventually left, are analysed by \textcite{WuBravermanLee2023} and \textcite{WuEtAl2024}. Rod flows \parencite{RegisChewi2026,RegisChewi2026Adam} are a recent alternative to central flows.

\paragraph{Translating the literature.} \Textcite{DamianNichaniLee2023} write $x_t$ and $y_t$ as we do, for the displacement along the top eigenvector and its projection on $\nabla S$, and $\alpha=-\ip{\nabla\Loss,\nabla S}$ as we do; their $\beta=\norm{\nabla S}^2$ is our $\gamma$, renamed because $\beta$ is the momentum in this book; their constrained trajectory $\theta_t^\dagger$ is the discrete counterpart of the central flow. \Textcite{CohenEtAl2025} write $w$ for the parameters and count time in steps, so their flows carry a factor $\eta$ that our $\tau=\eta t$ absorbs; their $\sigma^2(t)$ is our $\sigma^2(w(\tau))$, the average of $x^2$ over a cycle. \Textcite{CohenEtAl2021} write $\mathrm{MSS}(\eta,\beta)$ for the maximal stable sharpnesses of \cref{obs:eos:gd}(iii). In \cref{fig:eos:network} we include Adam's bias corrections in the preconditioner, $P_t=(1-\beta_1^{t+1})(\sqrt{\hat v_t}+\epsilon)$, which matters only in the first few hundred steps. Volume 2 of this series, Chapter 15, writes the reduced dynamics as $dA/dt=2\eps A$, $d\eps/dt=\gamma-\nu A$ in steps; its $(A,\eps,\gamma,\nu)$ are our $(x^2,\eta y,\eta^2\alpha,\frac12\eta^2\gamma)$, so its $A$ is our $e^A$.

\printchapterbib
\end{refsection}

% Chapter 9 of Optimizers in the Wild: warmup.

\casechapter{Warmup}{Why the step size must start small}{Linear warmup}{ch:warmup}
\begin{refsection}[chapters/09-warmup.bib]

\begin{chaptersummary}
Large training runs raise the learning rate from nearly zero to its peak over the first thousands of steps. The smallest idealization that makes this necessary bounds the curvature by an affine function of the loss gap, $S\le K_0+K_1(\Loss-\Loss^\star)$, so that the step that is safe at a point grows as the loss falls. We prove a descent lemma for this class and show that, under the Polyak--Łojasiewicz inequality, a warmup indexed by the loss reaches accuracy $\eps$ in $(K_0\log(\Delta_0/\eps)+K_1\Delta_0)/\mu$ steps while a constant step pays the initial curvature again in every factor of $\eps$; the optimal profile is a rising hyperbola that saturates, given by the Wright omega function, and a model that saturates both hypotheses shows it optimal among stable schedules and a linear ramp too slow by the factor $u_0/(4e)$. We also prove that Adam's first step is a sign step and that its second moment averages only the few gradients seen so far, and we derive that a warmup at the edge of stability drives the sharpness down along $2/\eta_t$, as a small network confirms. The theory does not derive its curvature bound for networks, where it fails globally, and it misses the catapults through which shorter, unstable warmups succeed.
\end{chaptersummary}

\section{In the wild: warming up the learning rate}\label{sec:warmup:wild}

The learning rate of a modern training run is a curve, and the curve starts low. An early instance in deep learning is in the paper that introduced residual networks: to train a network with $110$ layers on CIFAR-10, \textcite{HeEtAl2016} used the learning rate $0.01$ until the training error fell below $80\%$, about $400$ iterations, and only then the rate $0.1$ of their other networks, which they found too large to start converging.\sidenote{In this chapter \emph{learning rate}, the word of the practice sections, and \emph{step size}, the word of the theory, mean the same thing.} \Textcite{GoyalEtAl2017} made the ramp gradual. To train a ResNet-50 on ImageNet with minibatches of $8192$ images, they multiplied the learning rate by the factor by which the minibatch had grown, the linear scaling rule of \cref{ch:noise}, and reached that rate by a linear increase over the first five epochs, because the rule fails early in training, when the network changes rapidly. The original Transformer \parencite{VaswaniEtAl2017} increased its learning rate linearly for $4000$ steps and then decayed it as the inverse square root of the step. Large language models inherited the ramp.

\begin{practice}[GPT-3, \textcite{BrownEtAl2020}; LLaMA, \textcite{TouvronEtAl2023}; Adam's rule of thumb, \textcite{MaYarats2021}]
Raise the learning rate linearly from (nearly) zero to its peak, over the first $375$ million tokens for GPT-3 and the first $2000$ steps for LLaMA, then decay it along a cosine to $10\%$ of the peak (\cref{ch:schedules}); both runs use AdamW with $\beta_2=0.95$. For Adam a rule of thumb sets the length of a linear warmup to $2/(1-\beta_2)$ steps: $2000$ steps for $\beta_2=0.999$, and $40$ for $\beta_2=0.95$. Otherwise the length is inherited from an earlier recipe or tuned by a sweep; we know of no derivation of these numbers.
\end{practice}

What has been observed about warmup falls into three groups.

\begin{observation}[Warmup and the placement of layer normalization; \textcite{XiongEtAl2020}]\label[observation]{obs:warmup:postln}
In the original Transformer, which normalizes after each residual addition (Post-LN), the gradients of the parameters near the output are large at initialization, so that a large initial learning rate makes training unstable, and the warmup stage avoids this. With the normalization inside the residual branches (Pre-LN) the gradients are well behaved at initialization, and the warmup can be removed: Pre-LN Transformers trained without it reach results comparable to the baselines with less training time and less tuning.
\end{observation}

The statement about the gradients is a theorem of \textcite{XiongEtAl2020}, proved with mean-field approximations at initialization; the rest is experimental, on translation and language-model pretraining. An initialization can remove the need as well: \textcite{HuangEtAl2020} train Transformers without warmup and without layer normalization by rescaling their initial weights.

\begin{observation}[Warmup lowers the sharpness; \textcite{GilmerEtAl2022}, \textcite{KalraBarkeshli2024}]\label[observation]{obs:warmup:sharpness}
Early instabilities coincide with a sharpness well above $2/\eta$, and warmup stabilizes training by letting the sharpness fall before the learning rate is high \parencite{GilmerEtAl2022}. When a network starts in a region sharper than the target learning rate tolerates, its sharpness falls during warmup and stays near $2/\eta_t$, and the main benefit of warmup is that larger target learning rates become trainable \parencite{KalraBarkeshli2024}.
\end{observation}

The \term{sharpness}\index{sharpness} is the largest eigenvalue of the Hessian, $S(\theta)=\lmax(\nabla^2\Loss(\theta))$, and $2/\eta$ is the stability threshold of gradient descent on a quadratic (\cref{ch:eos}). \Cref{fig:warmup:mlp} shows the second pattern in a network small enough to train in seconds: the sharpness falls from $132$ to $20$ while the learning rate rises to $0.1$, and without warmup the same learning rate diverges in six steps.

\begin{observation}[Explanations specific to Adam]\label[observation]{obs:warmup:adam}
\textcite{LiuEtAl2020RAdam} attributed Adam's need for warmup to the variance of its adaptive learning rate in the first steps, when the second-moment estimate has averaged few gradients, and proposed RAdam, which rectifies the step until that variance is tractable. \textcite{MaYarats2021} found that an untuned linear warmup over $2/(1-\beta_2)$ steps works as well as RAdam. \textcite{KossonMessmerJaggi2024} attributed the need for warmup in GPT training to large angular updates of the weights in the first steps, and reduced it by controlling them. \textcite{KalraBarkeshli2024} initialized Adam's second moment with the first squared gradient (GI-Adam) and reported benefits similar to those of warmup.
\end{observation}

The three groups are not rival theories. The first two are about curvature, the third about Adam's statistics and the size of its early updates, and a transformer trained with Adam has all of them. What is robust, across architectures and optimizers, is that some warmup prevents early divergence at the learning rates that practice wants; in small transformers a longer warmup also makes the final loss less sensitive to the learning rate \parencite{WortsmanEtAl2024}. What is folklore is the length. In \cref{obs:warmup:sharpness} the curvature falls because the schedule rises. The idealization of this chapter starts from the other end, a curvature that falls because the loss does, and asks two questions: why must the step size start small, and how should it grow?

\section{The idealization: curvature bounded by the loss}\label{sec:warmup:ideal}

Gradient descent on a function whose curvature is at most $L$ everywhere decreases the loss at every step if $\eta<2/L$ (\cref{lem:toolkit:descent}). A network has no useful global $L$: the sharpness of the network of \cref{fig:warmup:mlp} is $132$ at initialization and $20$ after a few hundred steps, and the sharpness of a network with a rescaling symmetry is not even bounded on its set of minimizers (\cref{rem:warmup:symmetry}). The idealization replaces the global bound by one that improves as the loss falls.

\begin{assumption}[Loss-dependent curvature]\label[assumption]{ass:warmup:curvature}
The loss $\Loss\colon\R^d\to\R$ is twice continuously differentiable and bounded below, with $\Loss^\star=\inf\Loss$ and gap $\Delta(\theta)=\Loss(\theta)-\Loss^\star$. For the initial point $\theta_0$ let $\Theta_0=\{\theta\colon\Loss(\theta)\le\Loss(\theta_0)\}$. There are constants $K_0,K_1\ge0$ such that
\begin{equation}\label{eq:warmup:curvature}
  S(\theta)\le K_0+K_1\Delta(\theta)\qquad\text{for every }\theta\in\Theta_0 .
\end{equation}
\end{assumption}

The \term{loss-dependent curvature}\index{loss-dependent curvature} condition gives the curvature $K_0$ that may remain at the bottom and the excess $K_1\Delta$ that a gap $\Delta$ allows. It restricts the Hessian from above only; nothing in this chapter uses a lower bound. \Textcite{AlimisisIslamovLucchi2026} introduced the condition with the operator norm $\norm{\nabla^2\Loss}_{\op}$ in place of $S$, which implies \eqref{eq:warmup:curvature}, and \textcite{RiabininEtAl2026} state it in a general norm with an extra term $K_\rho\Delta^\rho$ (\cref{sec:warmup:normalized}).\sidenote{Its relative for the gradient, $\norm{\nabla^2\Loss}\le L_0+L_1\norm{\nabla\Loss}$, is the $(L_0,L_1)$-smoothness of \textcite{ZhangHeSraJadbabaie2020} and of \cref{ch:spikes}. \Textcite{AlimisisIslamovLucchi2026} prefer the loss to the gradient because the gradient norm can rise early in training, whereas the loss falls along safe steps.} A quadratic satisfies it with $K_1=0$. The function
\begin{equation}\label{eq:warmup:cosh}
  \Loss_{\mathrm c}(z)=\frac{K_0}{K_1}\Bigl(\cosh\bigl(\sqrt{K_1}\,z\bigr)-1\Bigr),\qquad z\in\R,
\end{equation}
satisfies it with equality, $\Loss_{\mathrm c}''=K_0\cosh(\sqrt{K_1}z)=K_0+K_1\Loss_{\mathrm c}$: its curvature grows as fast as its loss allows, and $2\cosh z-2$ is the case $K_0=2$, $K_1=1$. The second hypothesis says that the gradient is large when the loss is.

\begin{assumption}[Polyak--Łojasiewicz inequality]\label[assumption]{ass:warmup:pl}
There is $\mu>0$ such that $\norm{\nabla\Loss(\theta)}_2^2\ge2\mu\Delta(\theta)$ for every $\theta\in\Theta_0$.
\end{assumption}

Under the \term{Polyak--Łojasiewicz inequality}\index{Polyak--Łojasiewicz inequality} every stationary point in $\Theta_0$ is a global minimizer, and with a global curvature bound $L$ gradient descent converges at the linear rate $1-\mu/L$ (\cref{lem:toolkit:pl}).\sidenote{The inequality is due to Polyak and to Łojasiewicz, both in 1963; see \textcite{Polyak1987} and, for its uses in machine learning, \textcite{KarimiNutiniSchmidt2016}.} Both hypotheses are statements about the loss, and that is the point of combining them: the first says how large a step the current loss allows, the second how much a step of that size achieves.

The extremal member of the class has two kinds of coordinates.

\begin{definition}[The warmup model]\label[definition]{def:warmup:model}
Fix $\mu>0$, $K_0\ge2\mu$ and $K_1>0$. The \term{warmup model}\index{warmup model} is the loss
\[
  \Loss(x,z)=q(x)+\tfrac12S(x)\norm{z}_2^2,\qquad q(x)=\tfrac\mu2\norm{x}_2^2,\qquad S(x)=K_0+K_1q(x),
\]
on $\theta=(x,z)\in\R^{d_x}\times\R^{d_z}$, trained by gradient descent from a point $(x_0,z_0)$ with $z_0$ small.
\end{definition}

The coordinate $x$ is slow: it carries the loss, with curvature $\mu$. The coordinate $z$ is stiff: its curvature $S(x)$ is set by the loss of the slow coordinate. The model caricatures the start of training, where the loss is high because the network has not yet fitted the data, and the curvature in the stiffest direction is high for the same reason.

\begin{proposition}[The warmup model saturates both hypotheses]\label[proposition]{prop:warmup:model}
In the warmup model $\Loss^\star=0$, and the following hold.
\begin{enumerate}
\item \Cref{ass:warmup:curvature,ass:warmup:pl} hold on all of $\R^{d_x+d_z}$ with the constants $K_0$, $K_1$ and $\mu$, and even $\norm{\nabla^2\Loss(\theta)}_{\op}\le K_0+K_1\Delta(\theta)$.
\item On the subspace $\{z=0\}$ both hold with equality: $S=K_0+K_1\Delta$ and $\norm{\nabla\Loss}_2^2=2\mu\Delta$.
\item Gradient descent with steps $\eta_t$ satisfies
\begin{equation}\label{eq:warmup:modelmap}
  x_{t+1}=\Bigl(1-\eta_t\mu\bigl(1+\tfrac12K_1\norm{z_t}_2^2\bigr)\Bigr)x_t,\qquad z_{t+1}=\bigl(1-\eta_tS(x_t)\bigr)z_t .
\end{equation}
So the subspace $\{z=0\}$ is invariant, on it $\Delta_{t+1}=(1-\mu\eta_t)^2\Delta_t$, and a perturbation of the stiff coordinate is amplified at step $t$ exactly when $\eta_tS(x_t)>2$.
\end{enumerate}
\end{proposition}

\begin{proof}
Write $\varepsilon=\frac12K_1\norm{z}_2^2$. Then $\nabla_x\Loss=\mu(1+\varepsilon)x$, $\nabla_z\Loss=S(x)z$, $\Loss=q(1+\varepsilon)+\frac12K_0\norm{z}_2^2$ and $K_0+K_1\Loss=S(1+\varepsilon)$. Part (iii) follows from the gradient, and part (ii) from setting $z=0$, where the Hessian is block diagonal with blocks $\mu I$ and $S(x)I$, and $S(x)\ge K_0\ge\mu$.

For (i), the Hessian has the diagonal blocks $\mu(1+\varepsilon)I$ and $SI$ and the off-diagonal block $K_1\mu\,xz^\top$. On vectors orthogonal to $x$ in the first block or to $z$ in the second it acts as $\mu(1+\varepsilon)$ or as $S$, and both lie in $[0,S(1+\varepsilon)]$. If $x$ and $z$ are nonzero, it remains to bound the $2\times2$ matrix $A$ that it induces on the span of $(x/\norm{x}_2,0)$ and $(0,z/\norm{z}_2)$, with diagonal entries $\mu(1+\varepsilon)$ and $S$ and off-diagonal entry $b=K_1\mu\norm{x}_2\norm{z}_2$, for which $b^2=4\mu K_1q\varepsilon$. The matrix $S(1+\varepsilon)I-A$ has the nonnegative diagonal entries $(S-\mu)(1+\varepsilon)$ and $S\varepsilon$, and its determinant is at least $(S-\mu)S\varepsilon-b^2=\varepsilon\bigl((S-\mu)S-4\mu K_1q\bigr)$. The bracket is nonnegative: $K_0\ge2\mu$ gives $S-\mu\ge K_0/2+K_1q$, and the inequality of arithmetic and geometric means gives $(K_0/2+K_1q)(K_0+K_1q)\ge2\sqrt2\,K_0K_1q\ge4\mu K_1q$. Hence $\lmax(A)\le S(1+\varepsilon)=K_0+K_1\Loss$. Also $\lmin(A)\ge-b$, and $b=2\sqrt{K_1q\cdot\mu\varepsilon}\le K_1q+\mu\varepsilon\le S(1+\varepsilon)$, which gives the bound on the operator norm. Finally, because $S^2\ge K_0^2\ge\mu K_0$,
\[
  \norm{\nabla\Loss}_2^2=2\mu q(1+\varepsilon)^2+S^2\norm{z}_2^2\ge2\mu q(1+\varepsilon)+\mu K_0\norm{z}_2^2=2\mu\Loss .\qedhere
\]
\end{proof}

So on the line $z=0$ the warmup model has the most curvature that \cref{ass:warmup:curvature} allows and the least gradient that \cref{ass:warmup:pl} allows. Gradient descent on it is a scalar recursion as long as the stiff coordinate stays quiet, and the instability of that coordinate is the only way in which a step can fail. The hypothesis $K_0\ge2\mu$ is a convenience (Exercise~\ref{ex:warmup:threshold}).

\begin{dictionary}
the gap $\Delta_t$ & the training loss minus its floor, which is unknown in practice\\
$S\le K_0+K_1\Delta$ & the sharpness falls as the loss falls early in training (measured, not derived)\\
$K_0$, and the final step $1/K_0$ & the curvature late in training, which limits the peak learning rate\\
the excess $u_0=K_1\Delta_0/K_0$ & how much sharper the initialization is than the trained network\\
the PL constant $\mu$, and $\kappa=K_0/\mu$ & the slowest direction of descent, and a condition number\\
the loss-indexed step $1/(K_0+K_1\Delta_t)$ & a warmup that follows the loss rather than the step count\\
slow coordinate $x$, stiff coordinate $z$ & the bulk of the parameters, and the top eigenvector of the Hessian\\
$\eta_tS_t>2$ & the edge of stability (\cref{ch:eos}): a catapult or a divergence\\
Adam's $n_t$ & the number of gradients that the second-moment estimate has effectively averaged
\end{dictionary}

\paragraph{What was thrown away.} Noise: the gradients are exact, so the gradient noise behind the warmup of \textcite{GoyalEtAl2017} is absent (\cref{ch:noise}). The response of the landscape: \cref{ass:warmup:curvature} is a fixed inequality, whereas in a network the sharpness responds to the step size (\cref{ch:eos}). The geometry: in a network the Polyak--Łojasiewicz inequality can hold at best near the trajectory, and \cref{ass:warmup:curvature} only on a bounded region (\cref{rem:warmup:symmetry}). The optimizer: gradient descent without momentum or preconditioning, until \cref{sec:warmup:more}. Each simplification removes a mechanism that practice may rely on, so that what remains is the one that the hypothesis isolates: the curvature at the start exceeds what the target step tolerates, and it falls as training proceeds. If warmup were not needed in this idealization, it would not be needed for this reason in a network either.

\paragraph{What one had to notice.} The safe step is a property of the current point, and it can be indexed by the loss, which, unlike the gradient norm, does not increase along safe steps. A step that is safe at the current loss is therefore safe at every later one, a bound in terms of the loss turns into a bound in terms of time, and an increasing schedule falls out of it.

\section{The theory: a safe step that grows as the loss falls}\label{sec:warmup:theory}

The classical descent lemma integrates a curvature bound along the segment of a step. Under \cref{ass:warmup:curvature} the bound is available only where the loss is below its starting value, and the segment is not known in advance to stay there. A continuity argument closes the gap.

\begin{lemma}[A local descent lemma]\label[lemma]{lem:warmup:descent}
Let $\Loss$ be twice continuously differentiable, let $\theta,u\in\R^d$ satisfy $\ip{\nabla\Loss(\theta),u}<0$, and let $M>0$ be such that $u^\top\nabla^2\Loss(\theta+su)\,u\le M$ for every $s\ge0$ with $\Loss(\theta+su)\le\Loss(\theta)$. Then for every $s\in\bigl[0,2\abs{\ip{\nabla\Loss(\theta),u}}/M\bigr]$ the segment from $\theta$ to $\theta+su$ lies in $\{\Loss\le\Loss(\theta)\}$, and
\[
  \Loss(\theta+su)\le\Loss(\theta)+s\ip{\nabla\Loss(\theta),u}+\tfrac12Ms^2 .
\]
\end{lemma}

\begin{proof}
Let $\varphi(r)=\Loss(\theta+ru)$, $\gamma=-\varphi'(0)=\abs{\ip{\nabla\Loss(\theta),u}}$ and $\bar s=2\gamma/M$, and fix $s\in(0,\bar s]$; the case $s=0$ is trivial. The plan is to show that $\varphi$ cannot climb back to $\varphi(0)$ before $s$. Let $A$ be the set of $r\in[0,s]$ such that $\varphi\le\varphi(0)$ on $[0,r]$. It contains $0$, it is an interval, and it is closed because $\varphi$ is continuous, so $A=[0,r^*]$. On $A$ the hypothesis gives $\varphi''\le M$, and Taylor's formula with integral remainder gives, for $r\in A$,
\begin{equation}\label{eq:warmup:taylor}
  \varphi(r)\le\varphi(0)-\gamma r+\tfrac12Mr^2=\varphi(0)-\gamma r\bigl(1-r/\bar s\bigr).
\end{equation}
Suppose that $r^*<s$. If $r^*=0$, then $\varphi'(0)<0$ gives $\varphi<\varphi(0)$ on some interval $(0,\delta)$. If $r^*>0$, then $0<r^*<\bar s$ and \eqref{eq:warmup:taylor} give $\varphi(r^*)<\varphi(0)$, and by continuity $\varphi<\varphi(0)$ on some interval $[r^*,r^*+\delta)$. Either way $A$ contains points beyond $r^*$, which is absurd. So $r^*=s$, the segment lies in the sublevel set, and \eqref{eq:warmup:taylor} at $r=s$ is the claim.
\end{proof}

For gradient descent at a point $\theta\in\Theta_0$ with gap $\Delta$ and gradient $g\ne0$, take $u=-g$. The points of the ray with $\Loss\le\Loss(\theta)$ lie in $\Theta_0$ and have gap at most $\Delta$, so \cref{ass:warmup:curvature} gives the hypothesis of the lemma with $M=(K_0+K_1\Delta)\norm{g}_2^2$. Hence, for $0\le\eta\le2/(K_0+K_1\Delta)$,
\begin{equation}\label{eq:warmup:gd}
  \Loss(\theta-\eta g)\le\Loss(\theta)-\eta\Bigl(1-\frac{\eta(K_0+K_1\Delta)}2\Bigr)\norm{g}_2^2,
\end{equation}
and in particular $\Loss(\theta-\eta g)\le\Loss(\theta)-\frac\eta2\norm{g}_2^2$ if $\eta\le1/(K_0+K_1\Delta)$. This is the descent lemma with the global constant replaced by the curvature that the current loss allows. The step does not increase the loss, so the step that \eqref{eq:warmup:gd} certifies at the next point is at least as large: the \term{local descent lemma}\index{local descent lemma} already contains a warmup.

\begin{theorem}[Warmup under loss-dependent curvature]\label{thm:warmup:schedule}
Let \cref{ass:warmup:curvature,ass:warmup:pl} hold with $K_0>0$, let $\Delta_0=\Delta(\theta_0)>0$, and define $r(D)=\mu D/(K_0+K_1D)$ and the scalar sequence
\[
  D_0=\Delta_0,\qquad D_{t+1}=D_t-r(D_t).
\]
Then $\mu\le K_0$, and the gaps $\Delta_t=\Delta(\theta_t)$ of gradient descent satisfy the following.
\begin{enumerate}
\item \emph{Loss-indexed warmup.} With $\eta_t=1/(K_0+K_1\Delta_t)$, $\Delta_{t+1}\le\Delta_t-r(\Delta_t)$ for every $t$.
\item \emph{Scheduled warmup.} With $\eta_t=1/(K_0+K_1D_t)$, $\Delta_t\le D_t$ for every $t$. This schedule is nondecreasing, from $1/(K_0+K_1\Delta_0)$ to the limit $1/K_0$.
\item In both cases, for every $t\ge0$,
\begin{equation}\label{eq:warmup:time}
  K_0\log\frac{\Delta_0}{\Delta_t}+K_1(\Delta_0-\Delta_t)\ge\mu t,
\end{equation}
so that $\Delta_T\le\eps$ for every $\eps\in(0,\Delta_0]$ and every $T\ge\bigl(K_0\log(\Delta_0/\eps)+K_1(\Delta_0-\eps)\bigr)/\mu$.
\item \emph{Constant step.} With $\eta_t=\eta\le1/(K_0+K_1\Delta_0)$, $\Delta_t\le(1-\mu\eta)^t\Delta_0$; at the largest such $\eta$ this guarantees $\Delta_T\le\eps$ for $T\ge(K_0+K_1\Delta_0)\log(\Delta_0/\eps)/\mu$.
\end{enumerate}
\end{theorem}

\begin{proof}
The plan is to apply \eqref{eq:warmup:gd} at every step, after checking that the step is admissible there, and to turn the one-step decrease into the time bound \eqref{eq:warmup:time} by comparison with an integral. A step from a point with $\Delta_t=0$ does not move, because such a point is a global minimizer; and at a point of $\Theta_0$ with $\Delta_t>0$ the gradient is nonzero by \cref{ass:warmup:pl}. We use the convention $\log(\Delta_0/0)=+\infty$.

\emph{Part (i).} Suppose that $\theta_t\in\Theta_0$ and $\Delta_t>0$. The step $\eta_t=1/(K_0+K_1\Delta_t)$ is admissible in \eqref{eq:warmup:gd}, and with \cref{ass:warmup:pl},
\[
  \Delta_{t+1}\le\Delta_t-\frac{\norm{\nabla\Loss(\theta_t)}_2^2}{2(K_0+K_1\Delta_t)}\le\Delta_t-\frac{\mu\Delta_t}{K_0+K_1\Delta_t}=\Delta_t-r(\Delta_t).
\]
In particular $\theta_{t+1}\in\Theta_0$, and (i) follows by induction.

\emph{The inequality $\mu\le K_0$.} The same computation at an arbitrary point of $\Theta_0$ with positive gap $\Delta$, together with $\Loss^\star\le\Loss(\theta-\eta g)$, gives $\mu\le K_0+K_1\Delta$ there. The iterates of (i) have $\Delta_{t+1}\le(1-\mu/(K_0+K_1\Delta_0))\Delta_t$, so their gaps tend to $0$. If a gap vanishes, at a first step $T$, then by \cref{lem:warmup:descent} the segment from $\theta_{T-1}$ to $\theta_T$ lies in $\Theta_0$, and $\Loss$ takes on it every value between $\Loss^\star$ and $\Loss(\theta_{T-1})$. Either way $\Theta_0$ contains points of arbitrarily small positive gap, and $\mu\le K_0$.

\emph{Part (ii).} Because $\mu\le K_0$, the map $D\mapsto D-r(D)=D\bigl(1-\mu/(K_0+K_1D)\bigr)$ sends $[0,\infty)$ into itself, so the $D_t$ are nonnegative and nonincreasing; they tend to $0$ because $D_{t+1}\le(1-\mu/(K_0+K_1\Delta_0))D_t$. Hence $\eta_t$ is nondecreasing with limit $1/K_0$. For the comparison, suppose that $\Delta_t\le D_t$, as holds at $t=0$. If $\Delta_t=0$ there is nothing to prove. Otherwise $\eta_t=1/(K_0+K_1D_t)\le1/(K_0+K_1\Delta_t)$ is admissible in \eqref{eq:warmup:gd}, and
\[
  \Delta_{t+1}\le\Delta_t-\frac{\eta_t}2\norm{\nabla\Loss(\theta_t)}_2^2\le(1-\mu\eta_t)\Delta_t\le(1-\mu\eta_t)D_t=D_{t+1},
\]
where the last inequality uses $1-\mu\eta_t\ge1-\mu/K_0\ge0$.

\emph{Part (iii).} Let $\Phi(D)=K_0\log D+K_1D$ for $D>0$ and $\Phi(0)=-\infty$. It is increasing, and its derivative $\Phi'(D)=(K_0+K_1D)/D$ is decreasing. Let $(E_t)$ be a nonnegative sequence with $E_0=\Delta_0$ and $E_{t+1}\le E_t-r(E_t)$, such as the gaps of (i) or the numbers $D_t$ of (ii). If $E_{t+1}>0$, then, because $\Phi'$ is smallest at the right end of $[E_{t+1},E_t]$,
\[
  \Phi(E_t)-\Phi(E_{t+1})=\int_{E_{t+1}}^{E_t}\Phi'(s)\,\dd s\ge(E_t-E_{t+1})\,\Phi'(E_t)\ge r(E_t)\,\frac{K_0+K_1E_t}{E_t}=\mu .
\]
Summing, $\Phi(E_t)\le\Phi(\Delta_0)-\mu t$, which also holds when some $E_s$ vanishes. In case (ii), $\Delta_t\le D_t$ and the monotonicity of $\Phi$ give $\Phi(\Delta_t)\le\Phi(D_t)$. So in both cases $\Phi(\Delta_t)\le\Phi(\Delta_0)-\mu t$, which is \eqref{eq:warmup:time}. If $T\ge(\Phi(\Delta_0)-\Phi(\eps))/\mu$, then $\Phi(\Delta_T)\le\Phi(\eps)$, so $\Delta_T\le\eps$.

\emph{Part (iv).} If $\Delta_t\le\Delta_0$, then $\eta\le1/(K_0+K_1\Delta_t)$ is admissible in \eqref{eq:warmup:gd}, which gives $\Delta_{t+1}\le(1-\mu\eta)\Delta_t$. By induction $\Delta_t\le(1-\mu\eta)^t\Delta_0\le e^{-\mu\eta t}\Delta_0$, which is at most $\eps$ once $t\ge\log(\Delta_0/\eps)/(\mu\eta)$.
\end{proof}

Parts (i) to (iii) are the warmup: a step size that starts at the value $1/(K_0+K_1\Delta_0)$ that the initial curvature allows and rises to $1/K_0$ as the loss falls. The schedule of (ii) is a fixed function of the step count, computed in advance from four numbers. The rule of (i) needs the gap, which in practice is unknown, because $\Loss^\star$ is. The bound \eqref{eq:warmup:time} is additive: the initial excess curvature costs $K_1\Delta_0/\mu$ steps once, and the final curvature $K_0/\mu$ steps per factor $e$ of accuracy. The constant step of (iv) pays $(K_0+K_1\Delta_0)/\mu$ steps per factor $e$. The shape of the schedule can be computed exactly.

\begin{corollary}[The warmup profile]\label[corollary]{cor:warmup:profile}
In the setting of \cref{thm:warmup:schedule} let $K_1>0$ and $u_0=K_1\Delta_0/K_0$, and for real $s$ let $\omega(s)$ be the positive solution $w$ of $w+\log w=s$. The function
\[
  D(t)=\frac{K_0}{K_1}\,\omega\Bigl(u_0+\log u_0-\frac{\mu t}{K_0}\Bigr),\qquad t\ge0,
\]
solves $\dot D=-r(D)$ with $D(0)=\Delta_0$, and the numbers $D_t$ of the scheduled warmup satisfy, if $\mu<K_0$,
\begin{equation}\label{eq:warmup:sandwich}
  D\Bigl(\frac{t}{1-\mu/K_0}\Bigr)\le D_t\le D(t)\qquad\text{for every }t\ge0
\end{equation}
(the right inequality also when $\mu=K_0$). So the scheduled step $1/(K_0+K_1D_t)$ lies between the \term{warmup profile}
\[
  \eta(t)=\frac{1}{K_0\bigl(1+\omega(u_0+\log u_0-\mu t/K_0)\bigr)}
\]
and the same profile at time $t/(1-\mu/K_0)$. If $u_0\ge1$, the profile reaches half its final value, $1/(2K_0)$, at the time $T_{1/2}=(K_0/\mu)(u_0-1+\log u_0)$.
\end{corollary}

\begin{proof}
With $u=K_1D/K_0$, the function $\Phi$ of the proof of \cref{thm:warmup:schedule} is $\Phi(D)=K_0(u+\log u)+K_0\log(K_0/K_1)$, so $D(t)$ is defined by $\Phi(D(t))=\Phi(\Delta_0)-\mu t$. Differentiating, $\Phi'(D)\dot D=-\mu$, that is, $\dot D=-\mu D/(K_0+K_1D)=-r(D)$. The right inequality of \eqref{eq:warmup:sandwich} is $\Phi(D_t)\le\Phi(\Delta_0)-\mu t=\Phi(D(t))$, proved in part (iii) of the theorem, with the monotonicity of $\Phi$. For the left one, $\Phi'$ is largest at the left end of $[D_{t+1},D_t]$, and $D_{t+1}/D_t=1-\mu/(K_0+K_1D_t)\ge1-\mu/K_0$, so
\[
  \Phi(D_t)-\Phi(D_{t+1})\le(D_t-D_{t+1})\frac{K_0+K_1D_{t+1}}{D_{t+1}}\le r(D_t)\,\frac{K_0+K_1D_t}{D_{t+1}}=\mu\frac{D_t}{D_{t+1}}\le\frac{\mu}{1-\mu/K_0} .
\]
Summing, $\Phi(D_t)\ge\Phi(\Delta_0)-\mu t/(1-\mu/K_0)=\Phi\bigl(D(t/(1-\mu/K_0))\bigr)$. The time $T_{1/2}$ is the solution of $\omega(u_0+\log u_0-\mu t/K_0)=1$, that is, of $1=u_0+\log u_0-\mu t/K_0$.
\end{proof}

The function $\omega$ is the Wright omega function, $\omega(s)=W(e^s)$ with $W$ the principal branch of the Lambert function.\sidenote{The Lambert function inverts $w\mapsto we^w$, so $w+\log w=s$ means $we^w=e^s$. It is available as \texttt{scipy.special.lambertw}, and $\omega$ as \texttt{scipy.special.wrightomega}.} The profile has two phases. While $u=K_1D/K_0$ is large, $u+\log u=u_0+\log u_0-\mu t/K_0$ gives $u\approx u_0-\mu t/K_0$: the gap falls by about $\mu/K_1$ per step, and the step size rises along the hyperbola $\eta(t)\approx1/(K_0+K_1\Delta_0-\mu t)$, slowly at first and then faster. Once $u$ is small, $u\approx u_0e^{u_0}e^{-\mu t/K_0}$: the gap decays at the rate $\mu/K_0$, and the step size approaches $1/K_0$ exponentially fast. The profile is convex while $\eta<2/(3K_0)$ and concave afterwards (Exercise~\ref{ex:warmup:profile}). The \term{warmup}\index{warmup!length} is the first phase, and its length $T_{1/2}$ is the initial excess curvature $K_1\Delta_0$ divided by $\mu$, up to a logarithm. The sandwich \eqref{eq:warmup:sandwich} says that the continuum describes the discrete schedule up to a relative error $1/\kappa$ in time, where $\kappa=K_0/\mu$.

How much does the warmup buy? A lower bound for constant steps needs a hard function, and the cosh function \eqref{eq:warmup:cosh} supplies one: a constant step large enough for the late phase overshoots it at the start, while a quadratic punishes every step small enough for the start.

\begin{proposition}[A constant step pays the initial curvature at every scale]\label[proposition]{prop:warmup:constant}
Let $0<\mu\le K_0$, $K_1>0$, $\Delta_0>0$ and $0<\eps<\Delta_0/2$, and let
\[
  \Loss(x,z)=\tfrac\mu2x^2+\Loss_{\mathrm c}(z),\qquad(x,z)\in\R^2,
\]
with $\Loss_{\mathrm c}$ as in \eqref{eq:warmup:cosh}, started at a point $(x_0,z_0)$ with $\frac\mu2x_0^2=\Loss_{\mathrm c}(z_0)=\Delta_0/2$ and $z_0>0$. Then $\Loss^\star=0$, \cref{ass:warmup:curvature,ass:warmup:pl} hold on $\R^2$ with the constants $K_0$, $K_1$ and $\mu$, and with $a=\sqrt{K_1}z_0$ and $\eta^*=2a/(K_0\sinh a)$ gradient descent with a constant step $\eta$
\begin{enumerate}
\item never reaches $\Delta_t\le\eps$ if $\eta\ge\eta^*$, and
\item needs at least $(1-\eta\mu)\log(\Delta_0/(2\eps))/(2\eta\mu)$ steps to reach it if $\eta<\eta^*$.
\end{enumerate}
Hence every constant step needs at least
\begin{equation}\label{eq:warmup:lowerconst}
  T\ge\Bigl(\frac{(K_0+K_1\Delta_0/2)\tanh a}{4\mu a}-\frac12\Bigr)\log\frac{\Delta_0}{2\eps}
\end{equation}
steps, where $a=\operatorname{arccosh}(1+K_1\Delta_0/(2K_0))\approx\log(2+K_1\Delta_0/K_0)$.
\end{proposition}

\begin{proof}
The Hessian is diagonal with entries $\mu\le K_0$ and $K_0\cosh(\sqrt{K_1}z)=K_0+K_1\Loss_{\mathrm c}(z)$, both at most $K_0+K_1\Loss$. For the gradient, $\Loss_{\mathrm c}'(z)^2=(K_0^2/K_1)\sinh^2(\sqrt{K_1}z)$ and $\sinh^2=(\cosh-1)(\cosh+1)\ge2(\cosh-1)$ give $\Loss_{\mathrm c}'^2\ge2K_0\Loss_{\mathrm c}\ge2\mu\Loss_{\mathrm c}$, and $(\mu x)^2=2\mu\cdot\frac\mu2x^2$; adding, $\norm{\nabla\Loss}_2^2\ge2\mu\Loss$.

Gradient descent acts on the two coordinates separately: $x_{t+1}=(1-\eta\mu)x_t$ and $z_{t+1}=\psi(z_t)$ with the odd map $\psi(z)=z-\eta(K_0/\sqrt{K_1})\sinh(\sqrt{K_1}z)$. For $z>0$, $\abs{\psi(z)}\ge z$ holds exactly when $\psi(z)\le-z$, that is, when $\eta\ge\eta^*(z)=2\sqrt{K_1}z/(K_0\sinh(\sqrt{K_1}z))$, and $\eta^*(z)$ decreases in $z$ because $b/\sinh b$ decreases on $(0,\infty)$. For (i), suppose that $\eta\ge\eta^*=\eta^*(z_0)$ and $\abs{z_t}\ge z_0$; then $\eta\ge\eta^*(\abs{z_t})$, so $\abs{z_{t+1}}=\abs{\psi(\abs{z_t})}\ge\abs{z_t}\ge z_0$. By induction $\Delta_t\ge\Loss_{\mathrm c}(z_t)\ge\Loss_{\mathrm c}(z_0)=\Delta_0/2>\eps$ for every $t$. For (ii), $\Delta_t\ge\frac\mu2x_t^2=(1-\eta\mu)^{2t}\Delta_0/2$. If $\eta\mu\ge1$ the claimed bound is not positive. Otherwise $\Delta_T\le\eps$ forces $2T\log(1/(1-\eta\mu))\ge\log(\Delta_0/(2\eps))$, and $\log(1/(1-s))\le s/(1-s)$ gives the claim. Finally, $(1-\eta\mu)/(2\eta\mu)=1/(2\eta\mu)-1/2$ decreases in $\eta$, so for $\eta<\eta^*$ it exceeds $1/(2\eta^*\mu)-1/2$, and $K_0\sinh a=K_0\cosh a\tanh a=(K_0+K_1\Delta_0/2)\tanh a$ gives \eqref{eq:warmup:lowerconst}.
\end{proof}

For $K_1\Delta_0\gg K_0$ the bound \eqref{eq:warmup:lowerconst} is about $K_1\Delta_0\log(\Delta_0/(2\eps))/(8\mu a)$, against the guarantee $(K_0\log(\Delta_0/\eps)+K_1\Delta_0)/\mu$ of the warmup: up to the logarithm $a$ and constants, a constant step pays the initial excess curvature once per factor $e$ of accuracy, and the warmup pays it once. The logarithm is real. The cheapest constant step that diverges on $\Loss_{\mathrm c}$ is $\eta^*=(2a\coth a)/(K_0+K_1\Delta_0/2)$, larger than the safe step by the factor $a\coth a$, because the curvature of $\Loss_{\mathrm c}$ falls quickly along a long step, and a step overshoots only if it exceeds the average curvature, not the curvature at its start. Constants matter as well: with $K_0=K_1=1$, $\mu=0.02$, $\Delta_0=50$ and $\eps=10^{-6}$, the best constant step on this instance needs $1398$ steps, \eqref{eq:warmup:lowerconst} gives $1392$, the loss-indexed warmup takes $1051$ steps, and \cref{thm:warmup:schedule} guarantees only $3386$. The comparison of worst cases exhibits the structure; the warmup model exhibits the size.

\Textcite{AlimisisIslamovLucchi2026}, who write $H_0,H_1$ for $K_0,K_1$, prove bounds of both kinds. For gradient descent with a warmup step size, under an aiming condition (a weakening of convexity) in place of \cref{ass:warmup:pl}, the term of their upper bound that grows as $\eps\to0$ involves $H_0$ alone, and $H_1$ enters only an additive term. For constant steps they prove lower bounds, among them, for functions satisfying \cref{ass:warmup:pl}, a number of iterations of order $(H_1/(4\mu))\,\Delta_0\log(\Delta_0/\eps)/(1+\log\Delta_0)$ in their normalization, with a hard function that is, like ours, quadratic near the minimum with an exponential tail; the logarithm of $\Delta_0$ in the denominator is the factor $a$ above. The next proposition makes the comparison sharp on the warmup model, where the stiff curvature is driven by the slow coordinate's loss.

\begin{proposition}[Warmup is optimal among stable schedules]\label[proposition]{prop:warmup:stable}
Consider the warmup model on the line $z=0$, with $\Delta_0=q(x_0)>0$, $S_t=K_0+K_1\Delta_t$ and $\Phi(D)=K_0\log D+K_1D$, and let $0<\eps<\Delta_0$.
\begin{enumerate}
\item Every schedule with $\eta_tS_t\le2$ for all $t$, so that no perturbation of the stiff coordinate is ever amplified, needs at least $(1-2\mu/K_0)^2\bigl(\Phi(\Delta_0)-\Phi(\eps)\bigr)/(4\mu)$ steps to reach $\Delta_T\le\eps$.
\item The loss-indexed schedule $\eta_t=c/S_t$ with $0<c\le2$ reaches it within $\bigl(\Phi(\Delta_0)-\Phi(\eps)\bigr)/\bigl(c\mu(2-c\mu/K_0)\bigr)$ steps.
\item Every constant step with $\eta S_0\le2$ needs at least $(1-2\mu/S_0)\,S_0\log(\Delta_0/\eps)/(4\mu)$ steps.
\end{enumerate}
\end{proposition}

\begin{proof}
By \cref{prop:warmup:model}(iii), $\Delta_{t+1}=(1-\mu\eta_t)^2\Delta_t$. For (i), $\mu\eta_t\le2\mu/S_t\le2\mu/K_0\le1$, and the inequality $\mu\eta_t<1$ is strict when $\Delta_t>0$, because then $S_t>K_0$; so the gaps stay positive, and $\Delta_t-\Delta_{t+1}=\Delta_t\mu\eta_t(2-\mu\eta_t)\le4\mu\Delta_t/S_t$. Since $\Phi'(D)=(K_0+K_1D)/D$ is decreasing,
\[
  \Phi(\Delta_t)-\Phi(\Delta_{t+1})\le(\Delta_t-\Delta_{t+1})\frac{S_{t+1}}{\Delta_{t+1}}\le\frac{4\mu\Delta_t}{\Delta_{t+1}}=\frac{4\mu}{(1-\mu\eta_t)^2}\le\frac{4\mu}{(1-2\mu/K_0)^2},
\]
and summing over $t<T$ with $\Phi(\Delta_T)\le\Phi(\eps)$ gives (i). For (ii), $c\mu/S_t\le2\mu/K_0\le1$ and
\[
  \Delta_t-\Delta_{t+1}=\Delta_t\frac{c\mu}{S_t}\Bigl(2-\frac{c\mu}{S_t}\Bigr)\ge c\mu\Bigl(2-\frac{c\mu}{K_0}\Bigr)\frac{\Delta_t}{S_t},
\]
so the integral comparison of part (iii) of \cref{thm:warmup:schedule} gives $\Phi(\Delta_t)\le\Phi(\Delta_0)-c\mu(2-c\mu/K_0)\,t$. For (iii), $\Delta_T=(1-\mu\eta)^{2T}\Delta_0\ge(1-2\mu/S_0)^{2T}\Delta_0$, and the inequality $\log(1/(1-s))\le s/(1-s)$ with $s=2\mu/S_0$ finishes as in \cref{prop:warmup:constant}.
\end{proof}

On the warmup model the theorem's schedule ($c=1$) is within a factor $2+O(1/\kappa)$ of every stable schedule, and riding the stability limit ($c=2$) is optimal among them up to $1+O(1/\kappa)$. The best stable constant step is slower by about $S_0\log(\Delta_0/\eps)/(\Phi(\Delta_0)-\Phi(\eps))$, which is close to $\min\{1+u_0,\log(\Delta_0/\eps)\}$. The guarantee of \cref{thm:warmup:schedule} is conservative here by the factor $2$ of the descent lemma, which certifies only half of the first-order decrease. Whether unstable schedules can do better is a different question, and \cref{sec:warmup:reading} answers it: they can.

\section{The theory: cross-entropy, normalized steps, Adam and the edge}\label{sec:warmup:more}

\subsection{Where loss-dependent curvature comes from}\label{sec:warmup:ce}

Losses built from exponentials have curvature bounded by their value. The exponential $\Loss(\theta)=e^{-\theta}$ on $\R$ has $\Loss''=\Loss$, so $K_0=0$ and $K_1=1$, and its infimum $0$ is not attained; \cref{ass:warmup:pl} fails, since $\Loss'^2=\Loss^2$. Still, the loss-indexed step $1/(K_1\Delta_t)=e^{\theta_t}$ gives $\theta_{t+1}=\theta_t+1$: the loss falls by the factor $e$ at every step, and the step size grows without bound, a warmup that never ends (Exercise~\ref{ex:warmup:exponential}). The cosh function \eqref{eq:warmup:cosh} has $\Loss_{\mathrm c}'^2=2K_0\Loss_{\mathrm c}+K_1\Loss_{\mathrm c}^2$, so it satisfies \cref{ass:warmup:pl} with $\mu=K_0$. For cross-entropy the bound holds with a constant that the plan for this chapter had wrong by a factor of two.

\begin{proposition}[Softmax regression]\label[proposition]{prop:warmup:softmax}
For an input $a\in\R^{d_{\mathrm{in}}}$, a class $c\in\{1,\dots,K\}$ and weights $W\in\R^{K\times d_{\mathrm{in}}}$, let $p=\softmax(Wa)$ and $\ell(W)=-\log p_c$. Then
\[
  \lmax\bigl(\nabla^2_W\ell\bigr)=\norm{a}_2^2\,\lmax\bigl(\diag(p)-pp^\top\bigr)\le2\norm{a}_2^2\,(1-p_c)\le2\norm{a}_2^2\,\ell(W),
\]
and the constant $2$ cannot be lowered. Consequently the average $\Loss$ of $\ell$ over a data set with $\norm{a_i}_2\le R$ satisfies \cref{ass:warmup:curvature} with $K_0=2R^2\Loss^\star$ and $K_1=2R^2$, and with $K_0=0$ when the data are separable.
\end{proposition}

\begin{proof}
With $s=Wa$, $\ell=-s_c+\log\sum_ke^{s_k}$ has gradient $p-e_c$ and Hessian $M=\diag(p)-pp^\top$ in $s$. Since $s$ is linear in $W$, the second derivative of $\ell$ in a direction $V\in\R^{K\times d_{\mathrm{in}}}$ is $(Va)^\top M(Va)$. As $V$ ranges over matrices with $\norm{V}_F=1$, the vector $Va$ ranges over the ball of radius $\norm{a}_2$, the extremes being $V=va^\top/\norm{a}_2$ with $\norm{v}_2=1$; so the largest eigenvalue is $\norm{a}_2^2\lmax(M)$. By Gershgorin's theorem every eigenvalue of $M$ is at most $\max_i(M_{ii}+\sum_{j\ne i}\abs{M_{ij}})=\max_i2p_i(1-p_i)$. For $i=c$ this is at most $2(1-p_c)$; for $i\ne c$ it is at most $2p_i\le2(1-p_c)$. Finally $1-p_c\le-\log p_c$. For $K=2$, $M$ has the eigenvalue $2p_c(1-p_c)$, whose ratio to $\ell=-\log p_c$ tends to $2$ as $p_c\to1$. Averaging over the data gives $\lmax(\nabla^2\Loss)\le2R^2\Loss=2R^2\Loss^\star+2R^2\Delta$; for separable data $\Loss^\star=0$.
\end{proof}

The plan for this chapter stated the bound with $1-p_c$ in place of $2(1-p_c)$, which fails for two classes as soon as $p_c>1/2$. For binary logistic regression with a single weight vector, $\ell(w)=\log(1+e^{-y\ip{w,a}})$, the Hessian is $\sigma(m)\sigma(-m)aa^\top$ with $m=y\ip{w,a}$ and $\sigma$ the logistic function, and the constant is $1$ (Exercise~\ref{ex:warmup:softmax}); the factor $2$ comes from the two rows of $W$ moving in opposite directions. So cross-entropy on a linear model, or on the last layer of a network with its features frozen, has curvature bounded by the loss, and when the data are separable, the step that \eqref{eq:warmup:gd} certifies grows without bound as the loss falls. The rest of a network is a different matter.

\begin{remark}[Rescaling symmetries break the global inequality]\label[remark]{rem:warmup:symmetry}
For the product model $\Loss(u,v)=\frac12(uv-y_\star)^2$ with $y_\star>0$, every point of the hyperbola $uv=y_\star$ is a global minimizer, with sharpness $u^2+v^2$, which is unbounded on it. Every sublevel set with positive gap contains the hyperbola, so no constants satisfy \cref{ass:warmup:curvature}. The same holds, with the square loss, for every network whose loss is invariant under multiplying one layer by $\alpha$ and the next by $1/\alpha$, as linear and ReLU networks are: their minimizers come in families along which the sharpness is unbounded. In such networks the assumption can hold only on a bounded region, with constants that depend on it, or along a trajectory.
\end{remark}

The product model is the subject of volume 2 of this series, \emph{The Hydrogen Atoms of Deep Learning}, Chapter 15 \parencite{HydrogenAtoms2026}, where gradient descent with a large step leaves a sharp minimizer through a catapult: the dynamics selects the curvature, which no inequality between curvature and loss can express. \Textcite{AlimisisIslamovLucchi2026} report that the condition is satisfied by common architectures, and \textcite{RiabininEtAl2026} check it along training trajectories; both are statements about the region that training visits.

\subsection{Normalized steps}\label{sec:warmup:normalized}

The theorem is about the step size of gradient descent, which multiplies the gradient. Optimizers whose update has a fixed length in some norm, such as sign descent, Muon and, roughly, Adam (\cref{ch:steepest,ch:sign,ch:spectral}), have a learning rate that is a length, and the question is whether that length must warm up. Let $\norm{\cdot}$ be a norm with dual norm $\norm{\cdot}_*$, and let normalized steepest descent take the step $\theta_{t+1}=\theta_t+\eta_t\lmo(\nabla\Loss(\theta_t))$ of length $\eta_t$ (\cref{ch:steepest}).

\begin{proposition}[When normalized steps need warmup]\label[proposition]{prop:warmup:normalized}
Suppose that on $\Theta_0$ the loss satisfies $u^\top\nabla^2\Loss(\theta)u\le K(\Delta(\theta))\norm{u}^2$ for every $u$, with $K\colon[0,\infty)\to(0,\infty)$ nondecreasing, and $\norm{\nabla\Loss(\theta)}_*\ge c\,\Delta(\theta)^\nu$ for some $c>0$ and $\nu\in(0,1]$. Then normalized steepest descent with the loss-indexed length $\eta_t=c\Delta_t^\nu/K(\Delta_t)$ satisfies
\[
  \Delta_{t+1}\le\Delta_t-\frac{c^2\Delta_t^{2\nu}}{2K(\Delta_t)},
\]
and since the gaps decrease, the length rises over a stretch of training exactly where $\Delta\mapsto\Delta^\nu/K(\Delta)$ decreases. For $K(\Delta)=K_0+K_1\Delta+K_\rho\Delta^\rho$ with $K_0>0$, $K_1,K_\rho\ge0$ and $\rho>1$:
\begin{enumerate}
\item if $\nu=1$ and $K_\rho=0$, the length only decreases: no warmup;
\item if $\nu=1$ and $K_\rho>0$, the length rises while $\Delta_t>\Delta'=\bigl(K_0/((\rho-1)K_\rho)\bigr)^{1/\rho}$ and decreases afterwards;
\item if $\nu=\frac12$ and $K_\rho=0$, the length rises while $K_1\Delta_t>K_0$ and decreases afterwards.
\end{enumerate}
\end{proposition}

\begin{proof}
Let $g=\nabla\Loss(\theta_t)$ and $u=\lmo(g)$, so that $\norm{u}=1$ and $\ip{g,u}=-\norm{g}_*$. If $\Delta_t=0$ the length is $0$; otherwise $g\ne0$. Points of the ray with $\Loss\le\Loss(\theta_t)$ have gap at most $\Delta_t$, and $K$ is nondecreasing, so \cref{lem:warmup:descent} applies with $M=K(\Delta_t)$. The length satisfies $\eta_t\le\norm{g}_*/K(\Delta_t)$, so
\[
  \Delta_{t+1}\le\Delta_t-\eta_t\norm{g}_*+\tfrac12K(\Delta_t)\eta_t^2\le\Delta_t-\tfrac12\eta_t\norm{g}_*\le\Delta_t-\tfrac12\eta_tc\Delta_t^\nu,
\]
which is the claim. The logarithmic derivative of $\Delta^\nu/K(\Delta)$ is $(\nu K-\Delta K')/(\Delta K)$, and $\nu K-\Delta K'$ equals $K_0$ in case (i), $K_0-(\rho-1)K_\rho\Delta^\rho$ in case (ii), and $\frac12(K_0-K_1\Delta)$ in case (iii).
\end{proof}

The gradient bound with $\nu=1$ holds for convex losses whose sublevel set $\Theta_0$ lies within distance $R$ of a minimizer, with $c=1/R$, because $\Delta\le\ip{\nabla\Loss(\theta),\theta-\theta^\star}\le R\norm{\nabla\Loss(\theta)}_*$; $\nu=\frac12$ is the Polyak--Łojasiewicz inequality in the dual norm. So under \cref{ass:warmup:curvature} a normalized method needs no warmup when the gradient's dual norm falls as fast as the loss, and it needs one when the curvature falls faster than the gradient. Case (ii) is the shape that \textcite{RiabininEtAl2026} derive for norm-constrained optimizers such as Muon and Lion: in their analysis the learning rate warms up and then decays only when $\rho>1$, with the turning point at the same gap $\Delta'$, which is where $\Delta/K(\Delta)$ is largest. In every case the length tends to $0$ with the gap: a normalized step must be decayed to converge, the thread of \cref{ch:momentum,ch:schedules}. The step length of gradient descent itself, $\eta_t\norm{\nabla\Loss(\theta_t)}_2$, follows case (iii) on the warmup model.

\subsection{Adam's first steps}\label{sec:warmup:adam}

Adam with bias correction \parencite{KingmaBa2015} keeps, from $m_{-1}=v_{-1}=0$, the moving averages $m_t=\beta_1m_{t-1}+(1-\beta_1)g_t$ and $v_t=\beta_2v_{t-1}+(1-\beta_2)g_t^2$ (coordinatewise), divides them by $1-\beta_1^{t+1}$ and $1-\beta_2^{t+1}$, and steps $\theta_{t+1}=\theta_t-\eta\,\hat m_t/(\sqrt{\hat v_t}+\epsilon)$. The \term{bias correction}\index{bias correction} removes the pull of the zero initial state, and it has a consequence for the first steps.

\begin{proposition}[Adam's first steps]\label[proposition]{prop:warmup:adamfirst}
Let $\epsilon=0$, with $0/0=0$, and fix a coordinate.
\begin{enumerate}
\item The first step is a sign step: $\theta_1-\theta_0=-\eta\sign(g_0)$, whatever the size of $g_0$.
\item For every $t\ge0$, $\hat v_t=\sum_{s=0}^tw_{t,s}g_s^2$ with weights $w_{t,s}=(1-\beta_2)\beta_2^{t-s}/(1-\beta_2^{t+1})$ that sum to $1$ and satisfy $\beta_2^t/(t+1)\le w_{t,s}\le\beta_2^{-t}/(t+1)$. If $t(1-\beta_2)\le\delta<1$, then $(1-\delta)\,\bar v_t\le\hat v_t\le\bar v_t/(1-\delta)$, where $\bar v_t=\frac1{t+1}\sum_{s\le t}g_s^2$ is the plain average of the squared gradients so far.
\item The effective number of averaged gradients is
\begin{equation}\label{eq:warmup:neff}
  n_t=\Bigl(\sum_{s=0}^tw_{t,s}^2\Bigr)^{-1}=\frac{1+\beta_2}{1-\beta_2}\cdot\frac{1-\beta_2^{t+1}}{1+\beta_2^{t+1}},
\end{equation}
which increases from $n_0=1$ to $(1+\beta_2)/(1-\beta_2)=2/(1-\beta_2)-1$. If the $g_s^2$ are independent with common mean $\sigma^2$ and variance $\varsigma^2$, then $\E\hat v_t=\sigma^2$ and $\Var(\hat v_t)=\varsigma^2/n_t$.
\end{enumerate}
\end{proposition}

\begin{proof}
At $t=0$, $\hat m_0=g_0$ and $\hat v_0=g_0^2$, so the step is $-\eta g_0/\abs{g_0}=-\eta\sign(g_0)$, which proves (i). Unrolling the recursion, $v_t=(1-\beta_2)\sum_{s=0}^t\beta_2^{t-s}g_s^2$ and $\sum_{k=0}^t\beta_2^k=(1-\beta_2^{t+1})/(1-\beta_2)$, which gives the weights and their sum. Each numerator $\beta_2^{t-s}$ lies in $[\beta_2^t,1]$ and the denominator $\sum_{k\le t}\beta_2^k$ in $[(t+1)\beta_2^t,t+1]$, which gives the bounds; Bernoulli's inequality $\beta_2^t\ge1-t(1-\beta_2)\ge1-\delta$ gives the comparison with $\bar v_t$. For (iii), $\sum_sw_{t,s}^2=(1-\beta_2)^2(1-\beta_2^{2t+2})/\bigl((1-\beta_2^2)(1-\beta_2^{t+1})^2\bigr)$, which simplifies to the reciprocal of \eqref{eq:warmup:neff}; the mean and variance follow from linearity and independence.
\end{proof}

So the first update has size $\eta$ in every coordinate, however small the gradient (\cref{ch:sign}), and for the first $\delta/(1-\beta_2)$ steps the preconditioner is the root mean square of the handful of gradients seen so far. With $\beta_1=\beta_2$ the bound of \cref{ch:sign} keeps every step at most $\eta$ per coordinate, but Adam's step is never small just because the gradient is: an optimizer that starts at its peak learning rate moves every coordinate by $\eta$ at once. RAdam replaces the first steps by a rectified version.

\begin{fieldguide}{RAdam}{\textcite{LiuEtAl2020RAdam}}
\fgrow{update}{Adam's $m,v$; after $k$ gradients, with $\rho_\infty=2/(1-\beta_2)-1$ and $\rho_k=\rho_\infty-2k\beta_2^k/(1-\beta_2^k)$, step $-\eta\,r_k\hat m/\sqrt{\hat v}$ if $\rho_k>4$ and $-\eta\hat m$ otherwise, $r_k=\bigl((\rho_k-4)(\rho_k-2)\rho_\infty/((\rho_\infty-4)(\rho_\infty-2)\rho_k)\bigr)^{1/2}$}
\fgrow{norm}{Euclidean for the first four steps (momentum SGD), then as Adam}
\fgrow{preconditioner}{$\diag(\sqrt{\hat v})$, damped by $r_k\approx(k(1-\beta_2)/2)^{1/2}$ while $5\ll k\ll1/(1-\beta_2)$}
\fgrow{averages}{as Adam; $\rho_k$ is the length of a simple moving average that matches the bias-corrected one, and $\rho_\infty=n_\infty$ of \eqref{eq:warmup:neff}}
\fgrow{state}{two numbers per parameter}
\fgrow{in the wild}{proposed to make warmup unnecessary; an untuned linear warmup over $2/(1-\beta_2)$ steps does as well \parencite{MaYarats2021}}
\end{fieldguide}

\begin{heuristic}[Heuristic derivation: the $2/(1-\beta_2)$ rule]
Suppose that the gradients of a coordinate in the first steps are independent draws from one distribution, and that a step is safe once the preconditioner is nearly as reliable as it will ever be. By \cref{prop:warmup:adamfirst}(iii) the relative variance of $\hat v_t$ is proportional to $1/n_t$, and
\[
  \frac{n_t}{n_\infty}=\tanh\Bigl(\frac{(t+1)\log(1/\beta_2)}2\Bigr),
\]
which reaches $\tanh1\approx0.76$ when $t+1=2/\log(1/\beta_2)\approx2/(1-\beta_2)$. A warmup that holds the step down until then lasts $2/(1-\beta_2)$ steps, the rule of thumb of \textcite{MaYarats2021}; RAdam's factor $r_k\approx\sqrt{k(1-\beta_2)/2}$ is a square-root warmup on the same time scale.
\end{heuristic}

The $2/(1-\beta_2)$ time scale is the \term{effective sample size}\index{effective sample size} of the second-moment estimate, the same number that sets the memory of Adam's preconditioner in \cref{ch:adaptive}. It is $2000$ steps for $\beta_2=0.999$ and $40$ steps for $\beta_2=0.95$. LLaMA uses $\beta_2=0.95$ and warms up for $2000$ steps: the long warmups of language-model recipes are not explained by the statistics of Adam's second moment alone.

\subsection{Warmup at the edge of stability}\label{sec:warmup:edge}

\Cref{ass:warmup:curvature} makes the curvature a function of the loss. In networks trained by gradient descent the causality also runs the other way: the sharpness rises to $2/\eta$ and stays there (\cref{ch:eos}). What does a rising $\eta_t$ do to it?

\begin{heuristic}[Heuristic derivation: the sharpness follows $2/\eta_t$]
Suppose that full-batch gradient descent with a slowly increasing schedule has reached the \term{edge of stability}\index{edge of stability}, $S(\theta_t)=2/\eta_t$, and that from then on the iterates follow the \term{central flow}\index{central flow} of \cref{ch:eos}: in the time $\tau=\sum_s\eta_s$,
\[
  \frac{d\theta}{d\tau}=-\nabla\Loss-\tfrac12\sigma^2\nabla S,
\]
with $\sigma^2$ the mean square of the oscillation along the top eigenvector of the Hessian, chosen so that the flow stays on the edge. With $\dot\eta=\eta_{t+1}-\eta_t$, the edge moves at the rate $d(2/\eta)/d\tau=-2\dot\eta/\eta^3$, and the flow moves the sharpness at the rate $\ip{\nabla S,d\theta/d\tau}=\alpha-\frac12\sigma^2\norm{\nabla S}_2^2$, with $\alpha=-\ip{\nabla\Loss,\nabla S}$ the rate of progressive sharpening. Equating the two,
\[
  \sigma^2=\frac{2\,(\alpha+2\dot\eta/\eta^3)}{\norm{\nabla S}_2^2}.
\]
The edge holds as long as $\sigma^2>0$, that is, as long as the threshold $2/\eta_t$ falls faster than the sharpness would fall under gradient flow, and then the sharpness follows $2/\eta_t$ down: the schedule, not the loss, sets the curvature. The oscillation carries an excess loss of about $\frac12S\sigma^2=\sigma^2/\eta$, which grows with the speed of the warmup. If $\eta_t$ rises so fast that $\sigma^2$ is not small, the cubic expansion behind the central flow fails, and the iterates catapult or diverge. On the warmup model $\alpha=-2\mu(S-K_0)$ and $\norm{\nabla S}_2^2=2\mu K_1(S-K_0)$, so on the edge $S=2/\eta$ the condition $\sigma^2>0$ reads $\dot\eta>\mu\eta^2(2-\eta K_0)$.
\end{heuristic}

If $\alpha>0$, the edge, once reached, holds for every nondecreasing schedule, and after the warmup the sharpness stays at $2/\eta_{\max}$ as in \cref{ch:eos}. If $\alpha<0$, as on the warmup model, it holds only while the warmup outruns the flattening that gradient descent produces on its own. In both cases the warmup lowers the curvature that the run will have at its peak learning rate, which is \cref{obs:warmup:sharpness}. The idealizations agree on the direction and disagree on the mechanism: in \cref{thm:warmup:schedule} the step waits for the loss to lower the curvature; at the edge the step lowers the curvature itself, through the oscillation, and the loss pays for it.

\section{Reading the theory}\label{sec:warmup:reading}

\paragraph{The shape and the length of the warmup.} \Cref{cor:warmup:profile} predicts a step size that starts at $1/(K_0+K_1\Delta_0)$, rises slowly along a hyperbola while the excess curvature dominates, turns concave at $2/(3K_0)$ and saturates exponentially at $1/K_0$; its length is $T_{1/2}=(K_0/\mu)(u_0-1+\log u_0)$. On the warmup model the rate $\mu$ is replaced by $2\mu$, because there the descent lemma is loose by its factor $2$. \Cref{fig:warmup:model} (right) compares this profile, for $K_0=K_1=1$, $\mu=0.02$ and $\Delta_0=50$, with gradient descent on the full model: the step reaches $1/(2K_0)$ at step $1324$, where the continuum predicts $1323$.

\begin{figure}[t]
\begin{fullwidth}
\centering
\begin{tikzpicture}
\begin{groupplot}[group style={group size=2 by 1,horizontal sep=17mm},wplot,width=82mm,height=52mm,xmin=0]
\nextgroupplot[xlabel={step $t$},ylabel={gap $\Delta_t$},ymode=log,ymin=1e-8,ymax=3e4,xmax=4000,ytick={1e-8,1e-4,1,1e4},legend style={at={(0.5,1.03)},anchor=south,font=\scriptsize},legend columns=2]
\addplot[theory] table[x=t,y=adrec] {figures/data/09-warmup-model.dat};
\addplot[simulation,each nth point=4,filter discard warning=false] table[x=t,y=adsim] {figures/data/09-warmup-model.dat};
\addplot[theory3] table[x=t,y=corec] {figures/data/09-warmup-model.dat};
\addplot[simulation2,each nth point=4,filter discard warning=false] table[x=t,y=cosim] {figures/data/09-warmup-model.dat};
\addplot[theory2] table[x=t,y=lcrec] {figures/data/09-warmup-model.dat};
\addplot[only marks,mark=triangle*,mark size=1.6pt,color=sea,mark options={fill=sea,draw=white,line width=0.2pt},each nth point=2,filter discard warning=false] table[x=t,y=lcsim] {figures/data/09-warmup-model.dat};
\addplot[guide] table[x=t,y=guarantee] {figures/data/09-warmup-model.dat};
\legend{loss-indexed: recursion,{\quad gradient descent},{constant $1/S_0$: recursion},{\quad gradient descent},{linear, $T_{\mathrm w}=2000$: recursion},{\quad gradient descent},{guarantee of \cref{thm:warmup:schedule}}}
\nextgroupplot[xlabel={step $t$},ylabel={step size $\eta_t$},ymode=log,ymin=5e-3,ymax=4,xmax=7000,legend style={at={(0.5,1.03)},anchor=south,font=\scriptsize},legend columns=2]
\addplot[theory] table[x=t,y=etacont] {figures/data/09-warmup-model.dat};
\addplot[simulation,each nth point=6,filter discard warning=false] table[x=t,y=etaad] {figures/data/09-warmup-model.dat};
\addplot[guide] table[x=t,y=limad] {figures/data/09-warmup-model.dat};
\addplot[theory3] table[x=t,y=etals] {figures/data/09-warmup-model.dat};
\addplot[theory2] table[x=t,y=limls] {figures/data/09-warmup-model.dat};
\legend{{profile, \cref{cor:warmup:profile} with $2\mu$},{loss-indexed, gradient descent},{its limit $2/S_t$},{linear, $T_{\mathrm w}=6150$},{its limit $2/S_t$}}
\end{groupplot}
\end{tikzpicture}
\end{fullwidth}
\caption{The warmup model with $K_0=K_1=1$, $\mu=0.02$ and $\Delta_0=50$ ($\kappa=u_0=50$), one slow and one stiff coordinate, the stiff one started at $z_0=10^{-6}$; no randomness (script \texttt{09-warmup-model.py}). Left: the gap under the loss-indexed warmup $\eta_t=1/(K_0+K_1\Loss(\theta_t))$, under the constant step $1/S_0=1/(K_0+K_1\Delta_0)$, and under a linear warmup to $1/K_0$ over $2000$ steps. Curves: the recursion $\Delta_{t+1}=(1-\mu\eta_t)^2\Delta_t$ of \cref{prop:warmup:model}; points: gradient descent on the full model. The linear warmup crosses the edge $\eta_tS_t=2$, its stiff coordinate grows, and a catapult lifts the loss to $6\cdot10^3$ before it converges, faster than the recursion. Dashed: the guarantee \eqref{eq:warmup:time}, conservative by the factor $2$ of the descent lemma. Without warmup, $\eta=1/K_0$, the loss exceeds $10^{13}$ after six steps. Right: the step size of the loss-indexed warmup against the profile of \cref{cor:warmup:profile} (with $2\mu$ for $\mu$), and the shortest linear warmup to $1/K_0$ that never crosses the edge, which touches its stability limit near $t=400$.}
\label{fig:warmup:model}
\end{figure}

\paragraph{What the warmup buys.} On the same model the loss-indexed warmup reaches $\Delta\le10^{-8}$ in $1805$ steps, against $1808$ predicted by \cref{prop:warmup:stable}(ii) with the exact factor $(1-\mu\eta_t)^2$ and $3617$ guaranteed by \cref{thm:warmup:schedule}. The constant step $1/S_0$ needs $28469$ steps, and the best stable constant step, just below $2/S_0$, needs $14304$, against the lower bound $14224$ of \cref{prop:warmup:stable}(iii). The ratio, about $8$, is the factor $\min\{1+u_0,\log(\Delta_0/\eps)\}$ of the theory with its constants: $(1+u_0)\log(\Delta_0/\eps)/(2(\log(\Delta_0/\eps)+u_0))=7.9$. It grows with the excess curvature at initialization and with the accuracy required, and it is $1$ when the network does not start sharp.

\paragraph{A linear warmup is the wrong shape.} Practice uses straight lines from zero, and the model says how long a straight line must be to stay stable.

\begin{heuristic}[Heuristic derivation: the shortest stable linear warmup]
On the line $z=0$ of the warmup model take $\eta_t=t/(K_0T_{\mathrm w})$ for $t\le T_{\mathrm w}$, and replace the recursion by its continuum $\dot\Delta=-2\mu\eta_t\Delta$, so that $\Delta_t=\Delta_0\exp(-\mu t^2/(K_0T_{\mathrm w}))$. With $s=t/T_{\mathrm w}$ and $\chi=\mu T_{\mathrm w}/K_0$, the condition $\eta_tS_t\le2$ reads $s+u_0se^{-\chi s^2}\le2$ for $s\in[0,1]$. The second term peaks at $s=1/\sqrt{2\chi}$ with the value $u_0/\sqrt{2e\chi}$, and evaluating the condition there gives $(1+u_0/\sqrt e)/\sqrt{2\chi}\le2$, that is,
\[
  T_{\mathrm w}\ge\frac{K_0}{8\mu}\Bigl(1+\frac{u_0}{\sqrt e}\Bigr)^2.
\]
After the warmup the step is constant and $S_t$ decreases, so the condition holds there too as soon as $u_0e^{-\chi}\le1$, which is weaker.
\end{heuristic}

For $\kappa=u_0=50$ the formula gives $6133$ steps, and bisection on the exact recursion gives $6150$; for $u_0$ from $5$ to $200$ the two agree within $6\%$. A line from the origin must pass below the limit $2/S_t$, which stays near $2/S_0$ until the loss has fallen (\cref{fig:warmup:model}, right), so the stable length is quadratic in the initial excess $u_0$, while the profile of \cref{cor:warmup:profile} starts at the limit and is linear in $u_0$. The ratio of the two lengths is about $u_0/(4e)$, here $4.6$, and the stable linear warmup reaches $10^{-8}$ in $2617$ steps instead of $1805$.

\paragraph{Shorter warmups and catapults.} The theory says nothing about schedules that cross the edge, and the model shows that they can win. Linear warmups to $1/K_0$ over $1400$ steps or fewer diverge. Between $1600$ and $2600$ steps the run crosses the edge, a catapult lifts the loss to between $2\cdot10^3$ and $10^4$, and the run then converges faster than the stable recursion predicts: $865$ steps at $T_{\mathrm w}=1624$, half the steps of the loss-indexed warmup. From $3000$ steps on, the run follows the recursion, although it crosses the edge for every $T_{\mathrm w}<6150$: the factor $\abs{1-\eta_tS_t}$ is close to $0$ when $\eta_tS_t\approx1$, which crushes the stiff coordinate, and it has no time to grow back before the edge is left. The recursion describes the run only while the stiff coordinate is negligible, and that is where the agreement must fail.

\paragraph{The edge in a network.} \Cref{fig:warmup:mlp} trains a small network from a sharp initialization, $S_0\approx132$, to the peak learning rate $0.1$, whose threshold is $2/\eta=20$. Before the edge is reached, the sharpness falls slowly on its own. Once $2/\eta_t$ meets it, after about $200$ steps of the long warmup and $100$ of the short one, the sharpness follows $2/\eta_t$ down, a few percent above it, as \cref{sec:warmup:edge} predicts; after the warmup it stays between $19.4$ and $23$ for the long warmup and settles near $20$ for the short one. The short warmup meets the edge with a catapult: its loss rises from $1.3$ to $6$ before the sharpness drops. Warmups of $50$ or $100$ steps diverge, and $150$ steps suffice. Unlike the warmup model, the network holds the edge rather than overshooting it, because every step of a network feeds its top eigendirection, whereas the stiff coordinate of the model is fed only by its initial value.

\begin{figure}[t]
\begin{fullwidth}
\centering
\begin{tikzpicture}
\begin{groupplot}[group style={group size=2 by 1,horizontal sep=17mm},wplot,width=82mm,height=50mm,xmin=0,xmax=2000]
\nextgroupplot[xlabel={step $t$},ylabel={sharpness $S(\theta_t)$},ymode=log,ymin=10,ymax=400,legend style={at={(0.5,1.03)},anchor=south,font=\scriptsize},legend columns=2]
\addplot[theory] table[x=t,y=edgeB] {figures/data/09-warmup-mlp.dat};
\addplot[simulation2] table[x=t,y=SB] {figures/data/09-warmup-mlp.dat};
\addplot[theory3] table[x=t,y=edgeA] {figures/data/09-warmup-mlp.dat};
\addplot[simulation] table[x=t,y=SA] {figures/data/09-warmup-mlp.dat};
\legend{{$2/\eta_t$, $T_{\mathrm w}=1000$},{measured},{$2/\eta_t$, $T_{\mathrm w}=250$},{measured}}
\nextgroupplot[xlabel={step $t$},ylabel={training loss},ymode=log,ymin=5e-4,ymax=30,legend style={at={(0.5,1.03)},anchor=south,font=\scriptsize},legend columns=3]
\addplot[simulation2] table[x=t,y=LB] {figures/data/09-warmup-mlp.dat};
\addplot[simulation] table[x=t,y=LA] {figures/data/09-warmup-mlp.dat};
\addplot[only marks,mark=triangle*,mark size=1.7pt,color=sea,mark options={fill=sea,draw=white,line width=0.2pt}] table[x=t,y=loss] {figures/data/09-warmup-mlp-nowarm.dat};
\legend{{$T_{\mathrm w}=1000$},{$T_{\mathrm w}=250$},{no warmup}}
\end{groupplot}
\end{tikzpicture}
\end{fullwidth}
\caption{Warmup at the edge in a small network: a tanh network with three hidden layers of width $64$ and output multiplier $5$ (a sharp initialization), trained by full-batch gradient descent on the square loss with $256$ Gaussian inputs in $\R^{10}$ labelled by a random tanh teacher, with a linear warmup to $\eta=0.1$ over $250$ or $1000$ steps (seeds fixed in \texttt{09-warmup-mlp.py}). Left: the sharpness, by power iteration with Hessian--vector products every $20$ steps (points), and the prediction $2/\eta_t$ of \cref{sec:warmup:edge} (curves), which holds once the edge is reached. Right: the training loss; without warmup it exceeds $10^9$ after six steps.}
\label{fig:warmup:mlp}
\end{figure}

\section{What the theory misses}\label{sec:warmup:misses}

\paragraph{What was idealized away.} A global inequality between curvature and loss: in networks it is measured along trajectories, and \cref{rem:warmup:symmetry} shows that it cannot hold on a sublevel set when the network has a rescaling symmetry. The Polyak--Łojasiewicz inequality, which makes the rates computable and is at best local in a network. Full batches: gradient noise adds its own reason to start small, the one \textcite{GoyalEtAl2017} gave, and at small batches the edge sits lower \parencite{AndreyevBeneventano2024}. A fixed preconditioner: for Adam the relevant curvature is the preconditioned sharpness, whose threshold is near $38/\eta$ for $\beta_1=0.9$ \parencite{CohenEtAl2022}, and the preconditioner itself changes fastest in the first steps (\cref{prop:warmup:adamfirst}). Normalization layers and weight decay, which make the effective learning rate an angular speed that is far from its equilibrium early in training (\cref{ch:decay}). Each idealization removes a reason for warmup, so the length that the theory prescribes is a lower bound on what practice needs only in the curvature mechanism, and not in general.

\paragraph{Where it breaks.} Three findings contradict the picture of \cref{thm:warmup:schedule}. First, the causality runs the other way in networks: the sharpness at the end of a warmup is $2/\eta_{\max}$, set by the schedule and not by the loss \parencite{CohenEtAl2021,KalraBarkeshli2024}; in \cref{fig:warmup:mlp} both warmups end at $S\approx20$ with different losses, so a bound $S\le K_0+K_1\Delta$ fitted to one run would be wrong for the other. Second, instability is not failure. Constant steps above $2/S_0$ converge through a catapult in wide networks \parencite{LewkowyczEtAl2020} and in the product model (volume 2 of this series, Chapter 15); for logistic regression, large steps with a non-monotone loss converge faster than every stable step \parencite{WuEtAl2024}; and on the warmup model, a linear warmup that crosses the edge converges in half the steps of the stable optimum of \cref{prop:warmup:stable}. The lower bounds of \cref{prop:warmup:constant,prop:warmup:stable} concern schedules that never cross the edge; the theory has nothing to say about the others. Third, the mechanism need not be curvature at all. With Pre-LN or a suitable initialization, transformers train without warmup \parencite{XiongEtAl2020,HuangEtAl2020}; for Adam, initializing the second moment with the first squared gradient or controlling the angular update recovers much of warmup's benefit \parencite{KalraBarkeshli2024,KossonMessmerJaggi2024}, which points to Adam's early statistics and update sizes rather than to a curvature that only time can lower. Finally, practice fixes the length by rules of thumb: $2000$ steps in LLaMA, where $2/(1-\beta_2)=40$ and where $K_1\Delta_0/\mu$ has never been measured; the theory reproduces the shape of the practice, not its numbers.

\paragraph{Repairs and open problems.} \Textcite{AlimisisIslamovLucchi2026} turn their analysis into a warmup schedule with one parameter, which matches a tuned linear warmup on language and vision models; \textcite{RiabininEtAl2026} derive a schedule for norm-constrained optimizers that sets the warmup length automatically and, in LLaMA pretraining, matches or beats the best tuned warmups; \textcite{LiuEtAl2025Warmup}, with curvature bounded by $K_0+K_\rho\Delta^\rho$, show that warmup can make gradient descent up to $\Theta(T)$ times faster than every non-increasing schedule in some cases (Exercise~\ref{ex:warmup:power}). On the side of the optimizer, GI-Adam, Pre-LN, rescaled initializations and the update controls of \textcite{KossonMessmerJaggi2024} remove much of the need, each for its own reason. Open: a derivation of the curvature--loss inequality along the trajectories of a network, even a two-layer one (Exercise~\ref{ex:warmup:twolayer}); lower bounds over all schedules, including those that cross the edge (Exercise~\ref{ex:warmup:catapult}); a theory of warmup with gradient noise, where the edge itself depends on the batch size; and the preconditioned version of the whole story for Adam, in which the preconditioner and the curvature evolve on different time scales.

\section{Back in the wild}\label{sec:warmup:back}

\fieldnotes{Linear warmup of the learning rate from nearly zero over a few thousand steps ($2000$ for LLaMA, $375$ million tokens for GPT-3, five epochs for large-batch ImageNet).}
{Curvature bounded by the loss, $S\le K_0+K_1\Delta$, with the Polyak--Łojasiewicz inequality; full-batch gradient descent; the warmup model; Adam's moving averages; the central flow at the edge.}
{A local descent lemma; a loss-indexed or scheduled warmup reaches $\eps$ in $(K_0\log(\Delta_0/\eps)+K_1\Delta_0)/\mu$ steps, while a constant step pays the initial curvature at every scale; the exact profile; optimality among stable schedules on the warmup model; softmax curvature at most $2R^2\ell$; Adam's first step is a sign step.}
{A hyperbolic rise that saturates, of length about $K_1\Delta_0/\mu$; a stable linear warmup must be $u_0/(4e)$ times longer; the sharpness follows $2/\eta_t$ at the edge (derived); Adam's statistics settle after $2/(1-\beta_2)$ steps (derived).}
{The inequality is measured, not derived, and fails globally under rescaling symmetries; catapults let unstable warmups succeed and even win; noise, preconditioning and normalization layers.}
{Warm up until the curvature has fallen below what the peak step tolerates, and expect the needed length to grow with how sharp the initialization is (proved in the idealization). Indexing the warmup by the loss or by a measured sharpness, rather than by a fixed count, is suggested by the theory and by two recent schedulers; long warmups at $\beta_2=0.95$ are not explained by Adam's statistics (derived).}

The theory gives a practitioner one firm statement and one suggestion. The firm statement, proved in the idealization, is that a constant step must be safe at the start and then pays for the initial curvature at every scale of accuracy, while a warmup pays once; the gain is large only when the initialization is much sharper than the trained network, which a power iteration on the first minibatch can check against $2/\eta_{\max}$. The suggestion is the shape: start at the largest step that the current curvature allows and rise as it falls, which is a hyperbola in time, not a line from zero. The theory does not give the length of a language-model warmup, and it should not be read as forbidding short ones; at the edge, and through catapults, networks lower their own curvature faster than the stable theory allows.

The chapter sits between the edge of stability of \cref{ch:eos}, which supplies the mechanism of \cref{sec:warmup:edge}, and the clipping of \cref{ch:spikes}, which tames a curvature that grows with the gradient rather than with the loss. Warmup is the first phase of the schedules of \cref{ch:schedules}, and its interaction with normalized updates continues in the effective learning rate of \cref{ch:decay}: with weight decay and normalization layers, the angular update that \textcite{KossonMessmerJaggi2024} control early in training is the quantity that equilibrates late. That Adam's first step is a sign step connects it to \cref{ch:sign}, and the role of the norm in \cref{prop:warmup:normalized} to \cref{ch:steepest}.

\begin{exercises}
\item \consolidation (The softmax constant.) (a) Show that $\lmax(\diag(p)-pp^\top)\le2\max_ip_i(1-p_i)\le2(1-p_c)$ for every probability vector $p$ and class $c$. (b) For $K=2$ show that the top eigenvalue is $2p_c(1-p_c)$, and conclude that the constant $2$ of \cref{prop:warmup:softmax} cannot be lowered and that $\lmax\le1-p_c$ fails when $p_c>\frac12$. (c) For binary logistic regression with one weight vector, $\ell(w)=\log(1+e^{-y\ip{w,a}})$, show that $\nabla^2\ell=\sigma(m)\sigma(-m)\,aa^\top\preceq\ell(w)\,aa^\top$ with $m=y\ip{w,a}$.\label{ex:warmup:softmax}

\item \consolidation (The extremal function.) (a) Show that \cref{ass:warmup:curvature} implies $\norm{\nabla\Loss}_2^2\le2(K_0+K_1\Delta)\Delta$ on $\Theta_0$. (b) For $\Loss_{\mathrm c}$ of \eqref{eq:warmup:cosh} show that $\Loss_{\mathrm c}'^2=2K_0\Loss_{\mathrm c}+K_1\Loss_{\mathrm c}^2$. (c) Show that gradient descent with the loss-indexed step on $\Loss_{\mathrm c}$ satisfies $a_{t+1}=a_t-\tanh a_t$ for $a_t=\sqrt{K_1}z_t$, so that the loss falls by a factor close to $e$ per step while $a_t$ is large, and the warmup on this function lasts about $\log(K_1\Delta_0/K_0)$ steps rather than $K_1\Delta_0/\mu$. Which hypothesis of \cref{thm:warmup:schedule} is far from tight here?\label{ex:warmup:extremal}

\item \consolidation (The profile.) (a) Derive the formula of \cref{cor:warmup:profile} directly from $\dot D=-r(D)$. (b) Show that $\eta(t)$ is convex while $\eta(t)<2/(3K_0)$ and concave afterwards. \hint{with $u=K_1D/K_0$, $\dot u=-(\mu/K_0)u/(1+u)$, and $\ddot\eta$ has the sign of $2u-1$.} (c) For $K_0=K_1=1$, $\mu=0.01$ and $\Delta_0=100$, compute $T_{1/2}$ and the time at which $\eta$ reaches $0.9/K_0$.\label{ex:warmup:profile}

\item \consolidation\hard (The hypothesis $K_0\ge2\mu$.) Show that the warmup model satisfies $\lmax(\nabla^2\Loss)\le K_0+K_1\Loss$ everywhere if and only if $K_0\ge25\mu/16$. \hint{in the proof of \cref{prop:warmup:model} let $\varepsilon\to0$: the condition becomes $(S-\mu)S\ge4\mu K_1q$ for all $q\ge0$, a quadratic inequality in $K_1q$ whose discriminant is $25\mu^2-16K_0\mu$.}\label{ex:warmup:threshold}

\item \consolidation (Adam's effective sample size.) (a) Prove \eqref{eq:warmup:neff} and the formula $n_t/n_\infty=\tanh((t+1)\log(1/\beta_2)/2)$. (b) Show that RAdam's $\rho_k$ and $n_{k-1}$ both equal $k$ to first order in $k(1-\beta_2)$, and that both tend to $(1+\beta_2)/(1-\beta_2)$. (c) Show that with $\beta_1>0$ the first step is still $-\eta\sign(g_0)$, and compute the second step when $g_1=-g_0$.\label{ex:warmup:neffex}

\item \variation (No curvature floor: the exponential loss.) Let $\Loss(w)=\frac1n\sum_i\exp(-y_i\ip{w,a_i})$ with $\norm{a_i}_2\le R$, and suppose the data are separable with margin $\gamma$: some unit vector $\bar w$ has $y_i\ip{\bar w,a_i}\ge\gamma$ for all $i$. (a) Show that $\lmax(\nabla^2\Loss(w))\le R^2\Loss(w)$ and $\norm{\nabla\Loss(w)}_2\ge\gamma\Loss(w)$. (b) Show that gradient descent with the loss-indexed step $\eta_t=1/(R^2\Loss(w_t))$ satisfies $\Loss(w_{t+1})\le(1-\gamma^2/(2R^2))\Loss(w_t)$: linear convergence without the Polyak--Łojasiewicz inequality, with step sizes that grow geometrically. (c) Contrast with a constant step, under which the loss decays only polynomially (\cref{ch:bias}), and relate the result to the large-step regime of \textcite{WuEtAl2024}.\label{ex:warmup:exponential}

\item \variation (Curvature that grows like a power of the gap.) Suppose that $S\le K_0+K_\rho\Delta^\rho$ on $\Theta_0$ with $0<\rho\le1$, as for \textcite{LiuEtAl2025Warmup}, and that \cref{ass:warmup:pl} holds. Show that the loss-indexed step $1/(K_0+K_\rho\Delta_t^\rho)$ reaches $\eps$ within $\bigl(K_0\log(\Delta_0/\eps)+K_\rho(\Delta_0^\rho-\eps^\rho)/\rho\bigr)/\mu$ steps. Where does the proof of \cref{thm:warmup:schedule} use $\rho\le1$?\label{ex:warmup:power}

\item \variation (Momentum on the warmup model.) With the moving-average momentum of \cref{ch:momentum}, $m_t=\beta m_{t-1}+(1-\beta)\nabla\Loss(\theta_t)$ and $\theta_{t+1}=\theta_t-\eta m_t$, the stiff coordinate with frozen curvature $S$ is stable if and only if $\eta S<2(1+\beta)/(1-\beta)$. Redo the heuristic derivation of the shortest stable linear warmup with this threshold in place of $2$, assuming that the slow coordinate still obeys $\dot\Delta=-2\mu\eta_t\Delta$, and evaluate the result for $\beta=0.9$.\label{ex:warmup:momentum}

\item \variation (The step length of gradient descent.) On the line $z=0$ of the warmup model, show that the step length $\eta_t\norm{\nabla\Loss(\theta_t)}_2$ of the loss-indexed warmup is $\sqrt{2\mu\Delta_t}/(K_0+K_1\Delta_t)$, that it rises until $K_1\Delta_t=K_0$ and falls afterwards, and that this is case (iii) of \cref{prop:warmup:normalized}. What does a sign-descent run with a fixed step length do on this model when that length exceeds the peak?\label{ex:warmup:length}

\item \simulation (The warmup model.) With \texttt{09-warmup-model.py}: (a) check the shortest stable linear warmup against the heuristic formula for $u_0\in\{5,10,20,50,100,200\}$; (b) scan the warmup length and locate the window of catapults ($T_{\mathrm w}$ between about $1600$ and $2600$ for the default parameters), recording the largest loss of each run; (c) run the scheduled warmup of \cref{thm:warmup:schedule}(ii) and the loss-indexed one with $c=2$, and explain their step counts by \cref{prop:warmup:stable}.\label{ex:warmup:simmodel}

\item \simulation (A network at the edge.) With \texttt{09-warmup-mlp.py}: (a) measure $\eta_tS(\theta_t)$ during and after the warmup and check that it stays close to $2$ once the edge is reached; (b) find the shortest warmup that trains for the peak learning rates $0.05$ and $0.2$ (it is between $100$ and $150$ steps for $0.1$), and compare how it scales with the peak; (c) set the output multiplier to $1$, where $S_0\approx5.7$, and observe that with a peak whose threshold $2/\eta$ exceeds $S_0$ the sharpness rises toward $2/\eta_t$ instead of falling.\label{ex:warmup:simmlp}

\item \research (The inequality for a two-layer network.) \Cref{rem:warmup:symmetry} rules out \cref{ass:warmup:curvature} on sublevel sets of networks with rescaling symmetries. Find conditions, such as balanced or small initializations, bounded weight norms, or a restriction to gradient-descent trajectories, under which a two-layer network with the square or cross-entropy loss satisfies $S\le K_0+K_1\Delta$ with explicit constants, and compute $u_0=K_1\Delta_0/K_0$ for standard initializations. As far as we know, no such derivation exists.\label{ex:warmup:twolayer}

\item \research (Lower bounds that allow catapults.) On the warmup model, catapulting warmups beat every stable schedule (\cref{sec:warmup:reading}). Is there a lower bound over all schedules, stable or not, for gradient descent on the class of \cref{ass:warmup:curvature,ass:warmup:pl}? What is the fastest schedule on the warmup model, and does it cross the edge? Both questions are open.\label{ex:warmup:catapult}
\end{exercises}

\notesandrefs
Warmup entered deep learning as a fix for runs that would not start: \textcite{HeEtAl2016} held a $110$-layer residual network at a tenth of its learning rate for its first $400$ iterations, \textcite{GoyalEtAl2017} made the ramp linear to train with large minibatches, and \textcite{VaswaniEtAl2017} paired a linear warmup with an inverse-square-root decay. The recipes of GPT-3 and LLaMA \parencite{BrownEtAl2020,TouvronEtAl2023} fixed the modern form, a linear ramp followed by cosine decay. \Textcite{XiongEtAl2020} traced the need in Post-LN Transformers to large gradients near the output at initialization, and \textcite{HuangEtAl2020} removed it by initialization.

The curvature view of warmup is empirical: \textcite{GilmerEtAl2022} tied early instabilities to the sharpness, and \textcite{KalraBarkeshli2024} described the sharpness-reduction and progressive-sharpening regimes of warmup and proposed GI-Adam. The edge of stability is due to \textcite{CohenEtAl2021}, its self-stabilization to \textcite{DamianNichaniLee2023}, the central flow to \textcite{CohenEtAl2025}, the catapult to \textcite{LewkowyczEtAl2020}, and the solvable product model with its catapult to volume 2 of this series \parencite{HydrogenAtoms2026}. Large steps with a non-monotone loss on logistic regression are analysed by \textcite{WuEtAl2024}.

The theory of warmup under loss-dependent curvature is recent. \Textcite{AlimisisIslamovLucchi2026} introduced $(H_0,H_1)$-smoothness, proved upper bounds for gradient descent with a warmup step and lower bounds for constant steps, and turned the analysis into a one-parameter schedule. \Textcite{LiuEtAl2025Warmup} proved acceleration by warmup under a power-law version of the condition, and \textcite{RiabininEtAl2026}, in a preprint, extended the analysis to norm-constrained optimizers and derived an adaptive scheduler. The loss-dependent condition is a relative of the gradient-dependent $(L_0,L_1)$-smoothness of \textcite{ZhangHeSraJadbabaie2020}. The Polyak--Łojasiewicz inequality is surveyed by \textcite{KarimiNutiniSchmidt2016}. As far as we know, the warmup model, the profile of \cref{cor:warmup:profile} in terms of the Wright omega function, \cref{prop:warmup:stable} and the shortest stable linear warmup are new here.

For Adam, the bias correction is in \textcite{KingmaBa2015}; the variance argument and RAdam are due to \textcite{LiuEtAl2020RAdam}, the rule of thumb $2/(1-\beta_2)$ to \textcite{MaYarats2021}, and the angular-update analysis of warmup in GPT training to \textcite{KossonMessmerJaggi2024}. Adam's edge, with its threshold near $38/\eta$, is in \textcite{CohenEtAl2022}; the minibatch edge in \textcite{AndreyevBeneventano2024}. \Textcite{WortsmanEtAl2024} study warmup among the knobs of small proxy models for large-scale instabilities.

\paragraph{Translating the literature.} \Textcite{AlimisisIslamovLucchi2026} write $f$ for $\Loss$, $w$ for $\theta$, $f^*$ for $\Loss^\star$ and $H_0,H_1$ for $K_0,K_1$, and bound $\norm{\nabla^2f(w)}$ by $H_0+H_1(f(w)-f^*)$; their warmup step is a multiple $\theta\eta_k$ of a base step $\eta_k$, so their $\theta$ is a constant and not the parameter vector, and they count iterations by $K$. \Textcite{RiabininEtAl2026} write $x$ for $\theta$ and $\Delta^t$ for our $\Delta_t$, and bound the dual norm of gradient differences by $K(x)\norm{x-y}$ with $K(x)=K_0+K_1(f(x)-f^\star)+K_\rho(f(x)-f^\star)^\rho$; with $K_\rho=0$ and the Euclidean norm their assumption reduces to that of Alimisis et al. \Textcite{LiuEtAl2025Warmup} call their condition $(\rho,K_0,K_\rho)$-smoothness, a bound $K_0+K_\rho(f(w)-f^*)^\rho$ on the local curvature. In RAdam, $\rho_t$ is indexed by the number of gradients, our $t+1$.

\printchapterbib
\end{refsection}

% Chapter 10 of Optimizers in the Wild: gradient clipping, generalized smoothness and loss spikes.
\newcommand{\spikescliff}{\Loss_{\mathrm{cliff}}}
\newcommand{\spikesplat}{\Loss_{\mathrm{plat}}}
\newcommand{\spikesclip}{\operatorname{clip}}

\casechapter{Clipping and Spikes}{Generalized smoothness and the anatomy of an instability}{Gradient clipping, z-loss, QK-norm}{ch:spikes}
\begin{refsection}[chapters/10-spikes.bib]

\begin{chaptersummary}
Large language models are trained with the norm of the gradient clipped at $1$, a penalty on the normalizer of the output softmax and normalized queries and keys, and their losses still spike. The idealization is $(L_0,L_1)$-smoothness, a bound on the curvature by $L_0$ plus $L_1$ times the norm of the gradient. We prove a sharp local smoothness lemma, show that clipped steepest descent reaches a gradient of size $\eps$ within $\max\{4L_0\Delta_0/\eps^2,6L_1\Delta_0/\eps\}$ steps while every fixed step size pays almost linearly for the largest gradient met on the way, and show that normalized updates are clipped by construction. We then take three instabilities apart exactly: a block of weights that wakes up after a quiet period receives Adam steps of up to $0.21/\sqrt{1-\beta_2}$ times the learning rate; the cross-entropy has a flat direction along which adaptive updates drift and which the z-loss removes; and attention logits grow with the square of the weight scale unless queries and keys are normalized. The theory explains why each defence helps; it does not predict when a spike will come, and real spikes have causes it leaves out.
\end{chaptersummary}

\section{In the wild: clipped gradients and spiking losses}\label{sec:spikes:wild}

Almost every large training run rescales its gradient before using it. The Euclidean norm of the gradient is computed over all parameters at once; if it exceeds a threshold, the whole gradient is multiplied by the threshold divided by the norm, and otherwise it is left alone. The rule comes from recurrent networks, whose losses have \term{cliffs}\index{cliff}: thin regions in which the gradient and the curvature are both enormous, so that a single gradient step throws the parameters far away.\sidenote{Clipping entered deep learning with recurrent language models \parencite{Mikolov2012}. \Textcite{PascanuMikolovBengio2013} proposed clipping the norm of the gradient and explained the need with the geometry of cliffs, building on the analysis of exploding and vanishing gradients by \textcite{BengioSimardFrasconi1994}.} Transformers inherited the rule, and with it a growing list of defences against the instabilities that clipping does not prevent.

\begin{practice}[GPT-3, \textcite{BrownEtAl2020}; LLaMA, \textcite{TouvronEtAl2023}; PaLM, \textcite{ChowdheryEtAl2023}; ViT-22B, \textcite{DehghaniEtAl2023}; Kimi K2, \textcite{KimiTeam2025}]
Clip the global norm of the gradient at $1.0$ before the optimizer step (GPT-3, LLaMA, PaLM). Use Adam with $\beta_2=0.95$ rather than the default $0.999$ (GPT-3, LLaMA). Add $10^{-4}(\log Z)^2$ to the cross-entropy, where $Z$ is the normalizer of the softmax over the vocabulary (the \term{z-loss}\index{z-loss} of PaLM). Apply a LayerNorm to the queries and the keys before their dot product (\term{QK-norm}\index{QK-norm}, ViT-22B). With Muon, after each step rescale the query and key weights of every attention head whose largest logit on the batch exceeds a threshold, $100$ for Kimi K2 (\term{MuonClip}\index{MuonClip}). When the loss spikes anyway, restart from a checkpoint shortly before the spike and skip the batches around it (PaLM).
\end{practice}

The values travel from recipe to recipe; we know of no derivation of the threshold $1.0$ or of the coefficient $10^{-4}$, and in that sense both are folklore. What has been observed is the following.

\begin{observation}[Loss spikes; \textcite{ChowdheryEtAl2023}]\label[observation]{obs:spikes:palm}
While the $540$-billion-parameter PaLM model was trained, with the gradient norm clipped at $1.0$, the loss spiked roughly $20$ times, at irregular intervals and sometimes late in training. Restarting from a checkpoint about $100$ steps before a spike and skipping about $200$ to $500$ batches avoided it. Replaying the batches around a spike from a different, earlier checkpoint did not produce a spike.
\end{observation}

The last sentence carries the weight: the spike was in neither the data nor the parameters alone, but in their combination. The report is qualitative and concerns one run. The OPT team describes the same experience \parencite{ZhangEtAl2022OPT}: their divergences were cured by lowering the learning rate and restarting from an earlier checkpoint, and they coincided with a collapse of the dynamic loss scale used for mixed precision and with spikes in the norm of the activations of the last layer.

\begin{observation}[Curvature grows with the gradient; \textcite{ZhangHeSraJadbabaie2020}]\label[observation]{obs:spikes:L0L1}
Along the training trajectory of an LSTM language model, a local smoothness estimated from the change of the gradient between nearby iterates was strongly correlated with the norm of the gradient: where the gradient was large, so was the curvature.
\end{observation}

If the observation holds, a global bound on the curvature is set by the steepest part of the trajectory and is far larger than the curvature met during most of the run, so the step size $1/L$ that the classical theory prescribes is wasted almost everywhere. The evidence is a correlation along trajectories of a few networks, not a law.

\begin{observation}[Two instabilities in small models; \textcite{WortsmanEtAl2024}]\label[observation]{obs:spikes:proxies}
Small transformers trained at high learning rates reproduce two instabilities first reported for large runs: the growth of the attention logits, which makes the attention weights nearly one-hot \parencite{DehghaniEtAl2023}, and the divergence of the output logits from log-probabilities \parencite{ChowdheryEtAl2023}. Normalizing queries and keys removes the first, and the z-loss removes the second.
\end{observation}

\begin{observation}[Spikes follow a stale second moment; \textcite{WortsmanEtAl2023}]\label[observation]{obs:spikes:stale}
In training large CLIP models with AdamW, loss spikes followed, within a few iterations, steps at which the squared gradients of some layers were much larger than AdamW's running estimate of them: the second-moment estimate was out of date.
\end{observation}

From runs of language models with $7$ to $546$ billion parameters, \textcite{MolybogEtAl2023} (unrefereed) argue for a related account: Adam can enter a state in which its update has a large norm and is nearly uncorrelated with the direction of descent, and this is more likely in deep models trained with large batches.

The reports share a pattern. Each instability is a step too long for the curvature where it lands, and in each the step became too long because some quantity grew that the loss did not register in time: the gradient near a cliff, the gap between Adam's estimate of the gradient's scale and the scale itself, a common shift of the output logits, the size of queries and keys. Why can a fixed step size not cope with curvature that grows with the gradient, and what exactly does clipping buy? Through which mechanisms does a well-tuned run produce a step that is far too long? The idealization of this chapter answers the first two questions with theorems and the third with three exact mechanisms. It does not predict when a spike happens.

\section{The idealization: curvature that grows with the gradient}\label{sec:spikes:ideal}

Classical smoothness bounds the curvature by one number everywhere. On a cliff that number is huge, and the step size it allows is wasted on the rest of the landscape. \Textcite{ZhangHeSraJadbabaie2020} let the bound grow with the gradient.\sidenote{They wrote the condition in the Euclidean norm. The general norm costs nothing in the proofs and is what \cref{prop:spikes:normalized} needs: sign descent and Muon are steepest descent in the $\ell_\infty$ and spectral norms (\cref{ch:sign,ch:spectral}).}

\begin{assumption}[$(L_0,L_1)$-smoothness]\label[assumption]{ass:spikes:L0L1}
Let $\norm{\cdot}$ be a norm on $\R^d$ with dual norm $\norm{\cdot}_*$, and for a symmetric matrix $H$ let $\norm{H}=\max_{\norm u\le1}\norm{Hu}_*$. The objective $\Loss\colon\R^d\to\R$ is twice continuously differentiable and bounded below, $\Loss^\star=\inf\Loss>-\infty$, and there are constants $L_0\ge0$ and $L_1>0$ such that
\begin{equation}\label{eq:spikes:L0L1}
  \norm{\nabla^2\Loss(\theta)}\le L_0+L_1\norm{\nabla\Loss(\theta)}_*\qquad\text{for every }\theta\in\R^d .
\end{equation}
\end{assumption}

\term{$(L_0,L_1)$-smoothness}\index{generalized smoothness} with $L_1=0$ would be ordinary $L_0$-smoothness; the point of $L_1>0$ is to allow curvature that grows without bound, provided the gradient grows with it. The exponential $e^x$ is $(0,1)$-smooth and $\cosh x-1$ is $(1,1)$-smooth; the logistic loss $\log(1+e^{-x})$ is both $(\frac14,0)$- and $(0,1)$-smooth; $x^4$ is $(L_0,L_1)$-smooth exactly when $L_0L_1^2\ge16$; and $\sum_ia_i(\cosh\theta_i-1)$ is $(\max_ia_i,1)$-smooth in the Euclidean norm and $(\sum_ia_i,1)$-smooth in $\ell_\infty$ (\cref{ex:spikes:constants}). A quadratic with Hessian $H$ is $(\norm H,0)$-smooth and no better, because at its minimizer the gradient vanishes. The function that saturates \eqref{eq:spikes:L0L1} everywhere is the \term{exponential cliff}\index{cliff!exponential}
\begin{equation}\label{eq:spikes:cliff}
  \spikescliff(x)=\frac{L_0}{L_1^2}\bigl(e^{L_1\abs x}-1-L_1\abs x\bigr),\qquad
  \abs{\spikescliff'(x)}=\frac{L_0}{L_1}\bigl(e^{L_1\abs x}-1\bigr),\qquad
  \spikescliff''(x)=L_0e^{L_1\abs x},
\end{equation}
a valley of curvature $L_0$ whose walls steepen exponentially; it is twice continuously differentiable, and $\spikescliff''=L_0+L_1\abs{\spikescliff'}$. Exploding gradients have this shape: in a recurrent network the loss depends on a recurrent weight through high powers of it \parencite{BengioSimardFrasconi1994,PascanuMikolovBengio2013}, and a polynomial of high degree has derivatives that grow together.

\paragraph{The method.} Write $g^\sharp=-\norm g_*\lmo(g)$ for the sharp operator of the norm (\cref{ch:steepest}; \cref{lem:toolkit:conjugate}), so that $\ip{g,g^\sharp}=\norm g_*^2$ and $\norm{g^\sharp}=\norm g_*$. \term{Clipped steepest descent}\index{gradient clipping} with step size $\eta$ and largest step length $\gamma$ is
\begin{equation}\label{eq:spikes:clipped}
  \theta_{t+1}=\theta_t-\eta_t\,g_t^\sharp,\qquad g_t=\nabla\Loss(\theta_t),\qquad \eta_t=\min\Bigl\{\eta,\frac{\gamma}{\norm{g_t}_*}\Bigr\},
\end{equation}
with $\eta_t=\eta$ when $g_t=0$. Its step has length $\min\{\eta\norm{g_t}_*,\gamma\}$. Below the threshold $\norm{g_t}_*=\gamma/\eta$ it is unnormalized steepest descent with step size $\eta$; above it, it is normalized steepest descent $\theta_{t+1}=\theta_t+\gamma\lmo(g_t)$ with step length $\gamma$. In the Euclidean norm $g^\sharp=g$, and \eqref{eq:spikes:clipped} is the rule of the practice box: $\theta_{t+1}=\theta_t-\eta\spikesclip_\tau(g_t)$ with $\spikesclip_\tau(g)=\min\{1,\tau/\norm g_2\}g$ and threshold $\tau=\gamma/\eta$. The norm is taken over all parameters of the network at once, as in PyTorch's \texttt{clip\_grad\_norm\_}; clipping each layer separately is clipped steepest descent in a product norm (\cref{ch:modular}).\sidenote{Clipping \emph{by value} instead clips each coordinate at the threshold; it is the trust-region step for a box rather than a ball (\cref{ex:spikes:trust}), and it changes the direction of the step.} Clipping is a trust region\index{trust region}: the clipped step minimizes $\ip{g,u}+\norm u_2^2/(2\eta)$ over the ball $\norm u_2\le\gamma$ (\cref{ex:spikes:trust}).

\begin{fieldguide}{Clipped steepest descent}{\textcite{PascanuMikolovBengio2013}, Euclidean norm}
\fgrow{update}{$\theta_{t+1}=\theta_t-\min\{\eta,\gamma/\norm{g_t}_*\}\,g_t^\sharp$; in $\ell_2$, $\theta_{t+1}=\theta_t-\eta\min\{1,\tau/\norm{g_t}_2\}\,g_t$, threshold $\tau=\gamma/\eta$}
\fgrow{norm}{every norm; the step is unnormalized below the threshold and normalized, of length $\gamma$, above it}
\fgrow{preconditioner}{none; in practice the clipped gradient is passed to Adam, whose moments see the clipped values}
\fgrow{averages}{none}
\fgrow{state}{none; one norm per step, over all parameters}
\fgrow{in the wild}{$\tau=1.0$ on the global $\ell_2$ norm, in front of AdamW (GPT-3, LLaMA, PaLM)}
\fgrow{variants}{clipping by value, an $\ell_\infty$ trust region (\cref{ex:spikes:trust}); clipping the update rather than the gradient \parencite{ShazeerStern2018}}
\end{fieldguide}

\begin{dictionary}
$(L_0,L_1)$-smoothness \eqref{eq:spikes:L0L1} & curvature that is largest where the gradient is largest: cliffs, the first steps of training, a spike\\
$L_0$ & the curvature of the flat part of the landscape, which sets the useful learning rate\\
$L_1\norm{\nabla\Loss}_*$ & the extra curvature near a cliff or during a spike\\
$M$, the largest gradient on the sublevel set & the largest gradient norm the run will meet, including its spikes\\
step size $\eta$; step length $\gamma$; $\gamma/\eta$ & learning rate; learning rate times threshold; the threshold \texttt{max\_norm}\\
a coordinate that is silent, then loud & a token's embedding, a head or an expert whose gradient was near zero and suddenly is not\\
$a_k$ of \eqref{eq:spikes:burst} & Adam's update, per coordinate and in units of the learning rate, $k$ steps after waking up\\
the flat direction $\mathbf1$ of the logits & a common shift of all output logits, which moves $\log Z$\\
$\alpha(\log Z)^2$ & the z-loss, $\alpha=10^{-4}$\\
$s=\ip{W_Qx,W_Kx'}/\sqrt{d_{\mathrm h}}$ & an attention logit; MuonClip watches its maximum per head\\
the gains $\omega_q,\omega_k$ & the scale parameters of the normalizations of queries and keys
\end{dictionary}

\paragraph{What was thrown away.} The gradient is exact in the main theorems; noise returns only in \cref{thm:spikes:heavy} and in the anatomy of \cref{sec:spikes:anatomy}. The network is reduced to one relation between its curvature and its gradient, with constants that do not change during training; the constants of a real run do change, and we cannot compute them. Finite precision is absent. These are the right first simplifications because the failure of fixed step sizes is already present without noise or architecture: a step size that is safe where the gradient is largest is too small everywhere else, and \cref{thm:spikes:lower} shows that this costs a factor that grows with the largest gradient.

\paragraph{What one had to notice.} Classical analyses take a global smoothness constant as given and derive a step size from it. \Textcite{ZhangHeSraJadbabaie2020} measured the curvature along trajectories and saw that it moved with the gradient norm. Once curvature is tied to the gradient, a step whose \emph{length} is bounded is safe everywhere, and bounding the length of the step is exactly what clipping does. For spikes, the thing to notice is that each defence in the practice box targets a quantity whose growth the loss does not see until it is too late: the optimizer's state, a flat direction of the loss, a scale.

\section{The theory: clipping under generalized smoothness}\label{sec:spikes:clipping}

Everything rests on one lemma: under \eqref{eq:spikes:L0L1}, the gradient cannot change faster than exponentially along a segment, at a rate set by $L_1$.

\begin{lemma}[Local smoothness]\label[lemma]{lem:spikes:local}
Let $\Loss$ satisfy \cref{ass:spikes:L0L1}, let $\theta,u\in\R^d$, and put $g=\nabla\Loss(\theta)$ and $r=L_1\norm u$. Then
\begin{align}
  \norm{\nabla\Loss(\theta+u)-g}_*&\le\Bigl(\frac{L_0}{L_1}+\norm g_*\Bigr)\bigl(e^{r}-1\bigr),\label{eq:spikes:gradgrowth}\\
  \bigl\lvert\Loss(\theta+u)-\Loss(\theta)-\ip{g,u}\bigr\rvert&\le\Bigl(\frac{L_0}{L_1}+\norm g_*\Bigr)\frac{e^{r}-1-r}{L_1}.\label{eq:spikes:taylor}
\end{align}
Consequently, with $\varphi_1(c)=(e^c-1)/c$ and $\varphi_2(c)=(e^c-1-c)/c^2$, if $r\le c$ then
\begin{equation}\label{eq:spikes:local}
  \norm{\nabla\Loss(\theta+u)-g}_*\le\varphi_1(c)\,(L_0+L_1\norm g_*)\norm u,\qquad
  \Loss(\theta+u)\le\Loss(\theta)+\ip{g,u}+\varphi_2(c)\,(L_0+L_1\norm g_*)\norm u^2 .
\end{equation}
For $c=\frac13$, $\varphi_1(c)<1.19$ and $\varphi_2(c)<\frac35$. Both \eqref{eq:spikes:gradgrowth} and \eqref{eq:spikes:taylor} hold with equality for the cliff \eqref{eq:spikes:cliff} at $\theta\ge0$ and $u\ge0$.
\end{lemma}

As $c\to0$, $\varphi_1(c)\to1$ and $\varphi_2(c)\to\frac12$, and \eqref{eq:spikes:local} becomes the descent lemma for the local smoothness constant $L_0+L_1\norm g_*$ (\cref{lem:toolkit:descent}). The lemma says that this constant may be used as long as the step is short compared with $1/L_1$, at the price of the factors $\varphi_1,\varphi_2$.\sidenote{These are the functions $\varphi_1$ and $\varphi_2$ of exponential integrators for stiff differential equations, where they arise for the same reason: they integrate the exponential growth of a linear equation over one step.}

\begin{proof}
The plan is to bound the growth of the gradient along the segment from $\theta$ to $\theta+u$ by a linear differential inequality, solve it with Grönwall's inequality, and integrate. If $u=0$ there is nothing to prove, so let $u\ne0$. For $s\in[0,1]$ let $\phi(s)=\nabla\Loss(\theta+su)$, whose derivative $\phi'(s)=\nabla^2\Loss(\theta+su)u$ is continuous, and let $\psi(s)=\int_0^s\norm{\phi'(\sigma)}_*\dd\sigma$, a continuously differentiable function with $\psi(0)=0$. By the fundamental theorem of calculus and the triangle inequality for integrals, $\norm{\phi(s)-g}_*\le\psi(s)$. The assumption \eqref{eq:spikes:L0L1}, applied at $\theta+su$, bounds the derivative of $\psi$ by the current gradient, and the current gradient by $\norm g_*+\psi(s)$:
\[
  \psi'(s)=\norm{\nabla^2\Loss(\theta+su)u}_*\le\bigl(L_0+L_1\norm{\phi(s)}_*\bigr)\norm u\le a+r\,\psi(s),\qquad a=(L_0+L_1\norm g_*)\norm u .
\]
Hence $(e^{-rs}\psi(s))'=e^{-rs}(\psi'(s)-r\psi(s))\le ae^{-rs}$, and integrating from $0$ to $s$ gives $\psi(s)\le\frac ar(e^{rs}-1)$, the differential form of Grönwall's inequality (\cref{thm:toolkit:gronwall}). Since $a/r=L_0/L_1+\norm g_*$, the case $s=1$ is \eqref{eq:spikes:gradgrowth}. For \eqref{eq:spikes:taylor}, write the remainder of the first-order expansion as an integral and insert the bound on $\psi$:
\[
  \bigl\lvert\Loss(\theta+u)-\Loss(\theta)-\ip{g,u}\bigr\rvert=\Bigl\lvert\int_0^1\ip{\phi(s)-g,u}\dd s\Bigr\rvert\le\norm u\int_0^1\psi(s)\dd s\le\frac ar\norm u\Bigl(\frac{e^r-1}r-1\Bigr),
\]
which equals the right side of \eqref{eq:spikes:taylor}. The power series $\varphi_1(c)=\sum_{k\ge0}c^k/(k+1)!$ and $\varphi_2(c)=\sum_{k\ge0}c^k/(k+2)!$ have positive coefficients, so both functions increase on $(0,\infty)$; for $r\le c$ this gives $e^r-1=r\varphi_1(r)\le r\varphi_1(c)$ and $e^r-1-r=r^2\varphi_2(r)\le r^2\varphi_2(c)$, and substituting $r=L_1\norm u$ gives \eqref{eq:spikes:local}. The numbers: $\varphi_1(\frac13)=3(e^{1/3}-1)=1.187$, and $\varphi_2(\frac13)=9(e^{1/3}-\frac43)<9(\frac75-\frac43)=\frac35$ because $e<(\frac75)^3=2.744$. Finally, for the cliff at $\theta\ge0$ and $u\ge0$, both sides of \eqref{eq:spikes:gradgrowth} equal $\frac{L_0}{L_1}e^{L_1\theta}(e^{L_1u}-1)$, and integrating this in $u$ gives equality in \eqref{eq:spikes:taylor}.
\end{proof}

With the lemma, the main theorem is one application per step.

\begin{theorem}[Clipped steepest descent]\label{thm:spikes:clipped}
Let $\Loss$ satisfy \cref{ass:spikes:L0L1}, let $\theta_0\in\R^d$ and $\Delta_0=\Loss(\theta_0)-\Loss^\star$, and run \eqref{eq:spikes:clipped} with $0<\gamma\le1/(3L_1)$ and $0<\eta\le1/(2L_0)$ (every $\eta>0$ is allowed if $L_0=0$). Then:
\begin{enumerate}
\item $\Loss(\theta_{t+1})\le\Loss(\theta_t)-\frac12\min\bigl\{\eta\norm{g_t}_*^2,\ \gamma\norm{g_t}_*\bigr\}$ for every $t$;
\item for every $\eps>0$, the number of steps with $\norm{g_t}_*>\eps$ is at most $2\Delta_0\max\{1/(\eta\eps^2),1/(\gamma\eps)\}$. If $L_0>0$, then with $\eta=1/(2L_0)$ and $\gamma=1/(3L_1)$ it is at most
\begin{equation}\label{eq:spikes:clippedrate}
  \max\Bigl\{\frac{4L_0\Delta_0}{\eps^2},\ \frac{6L_1\Delta_0}{\eps}\Bigr\},
\end{equation}
which is $4L_0\Delta_0/\eps^2$ whenever $\eps\le2L_0/(3L_1)$.
\end{enumerate}
\end{theorem}

\begin{proof}
The plan is to apply \cref{lem:spikes:local} to one step, whose length is at most $\gamma$, and to check that the quadratic remainder costs at most half of the first-order decrease. If $g_t=0$, the step is zero and (i) holds with equality. Otherwise put $u=-\eta_tg_t^\sharp$, so that $\ip{g_t,u}=-\eta_t\norm{g_t}_*^2$ and $\norm u=\eta_t\norm{g_t}_*=\min\{\eta\norm{g_t}_*,\gamma\}$. Since $L_1\norm u\le L_1\gamma\le\frac13$, \eqref{eq:spikes:local} with $c=\frac13$ and $\varphi_2(\frac13)<\frac35$ gives
\[
  \Loss(\theta_{t+1})\le\Loss(\theta_t)-\eta_t\norm{g_t}_*^2+\tfrac35\bigl(L_0+L_1\norm{g_t}_*\bigr)\eta_t^2\norm{g_t}_*^2
  =\Loss(\theta_t)-\eta_t\norm{g_t}_*^2\Bigl(1-\tfrac35\bigl(L_0\eta_t+L_1\norm u\bigr)\Bigr).
\]
The two terms in the inner parenthesis are controlled by the two parameters: $L_0\eta_t\le L_0\eta\le\frac12$ and $L_1\norm u\le\frac13$. The large parenthesis is therefore at least $1-\frac35\cdot\frac56=\frac12$, and the decrease is at least $\frac12\eta_t\norm{g_t}_*^2=\frac12\min\{\eta\norm{g_t}_*^2,\gamma\norm{g_t}_*\}$, which is (i).

For (ii), sum (i) over $t<T$: the sum of the decreases is at most $\Loss(\theta_0)-\Loss(\theta_T)\le\Delta_0$. The functions $x\mapsto\eta x^2$ and $x\mapsto\gamma x$ increase, so every step with $\norm{g_t}_*>\eps$ decreases the loss by more than $\frac12\min\{\eta\eps^2,\gamma\eps\}$, and there are at most $2\Delta_0/\min\{\eta\eps^2,\gamma\eps\}=2\Delta_0\max\{1/(\eta\eps^2),1/(\gamma\eps)\}$ of them, for every $T$. Substituting $\eta=1/(2L_0)$ and $\gamma=1/(3L_1)$ gives \eqref{eq:spikes:clippedrate}; its first entry is the larger exactly when $\eps\le2L_0/(3L_1)$.
\end{proof}

The bound does not contain the largest gradient met along the way. A clipped step on a cliff is short, but the gradient there is large, and it decreases the loss by $\frac12\gamma\norm{g_t}_*$, more than an unclipped step does near the target; so the cliff costs nothing in \eqref{eq:spikes:clippedrate} beyond its share of $\Delta_0$. For gradient descent with a fixed step the largest gradient is the whole story. Let
\[
  M=\sup\bigl\{\norm{\nabla\Loss(\theta)}_*:\Loss(\theta)\le\Loss(\theta_0)\bigr\}
\]
be the largest gradient on the initial sublevel set, assumed finite.

\begin{corollary}[Gradient descent with a step set by $M$]\label[corollary]{cor:spikes:gd}
Let $\Loss$ satisfy \cref{ass:spikes:L0L1} with $M<\infty$, and let $0<\eta\le\min\{1/(2L_0),1/(3L_1M)\}$. Then unnormalized steepest descent $\theta_{t+1}=\theta_t-\eta g_t^\sharp$ keeps every iterate in the sublevel set, decreases the loss by at least $\frac\eta2\norm{g_t}_*^2$ at every step, and has $\norm{g_t}_*>\eps$ for at most $2\Delta_0/(\eta\eps^2)$ steps, for every $\eps\le M$. With the largest allowed $\eta$ this is $2\Delta_0\max\{2L_0,3L_1M\}/\eps^2$.
\end{corollary}

\begin{proof}
Put $\gamma=\eta M\le1/(3L_1)$ and argue by induction. If $\Loss(\theta_t)\le\Loss(\theta_0)$, then $\norm{g_t}_*\le M=\gamma/\eta$, so $\eta_t=\eta$ in \eqref{eq:spikes:clipped} and the step of gradient descent is the clipped step. \Cref{thm:spikes:clipped}(i) then gives a decrease of at least $\frac12\min\{\eta\norm{g_t}_*^2,\eta M\norm{g_t}_*\}=\frac\eta2\norm{g_t}_*^2$, so $\theta_{t+1}$ is in the sublevel set too. The count follows as in \cref{thm:spikes:clipped}(ii), since $\gamma\eps=\eta M\eps\ge\eta\eps^2$.
\end{proof}

So gradient descent works on the class if its step size is of order $1/(L_0+L_1M)$, and then it needs of order $(L_0+L_1M)\Delta_0/\eps^2$ steps. The next theorem shows that no fixed step size does much better: up to a logarithm, gradient descent must pay for $M$.

\begin{theorem}[A fixed step pays for the largest gradient]\label{thm:spikes:lower}
Let $L_0,L_1,M,\Delta_0>0$ and $0<\eps\le M/2$, with $\Delta_0\ge M/L_1$, and fix $x_0\in\R$. Let $\cF$ be the set of twice continuously differentiable functions $f\colon\R\to\R$ that are bounded below and $(L_0,L_1)$-smooth, with $f(x_0)-\inf f\le\Delta_0$ and $\abs{f'(x)}\le M$ whenever $f(x)\le f(x_0)$. For every step size $\eta>0$ there is an $f\in\cF$ on which gradient descent $x_{t+1}=x_t-\eta f'(x_t)$ has $\abs{f'(x_t)}>\eps$ for every $t<N$, with
\begin{equation}\label{eq:spikes:lower}
  N=\frac{L_1M}{\log(1+L_1M/L_0)}\cdot\frac{\Delta_0-4\eps^2/(\sqrt3L_0)}{8\eps^2}.
\end{equation}
If $\eta\ge\eta^*=2\log(1+L_1M/L_0)/(L_1M)$, then $f$ can be chosen so that $\abs{f'(x_t)}\ge M$ for every $t$.
\end{theorem}

When $\eps^2\le L_0\Delta_0/4$, the second factor of \eqref{eq:spikes:lower} is at least $\Delta_0/(20\eps^2)$. Because $\log(1+P)<P$ for $P>0$, the first factor exceeds $L_0$, so the theorem also contains the familiar lower bound of order $L_0\Delta_0/\eps^2$; what is new is the factor $L_1M/\log(1+L_1M/L_0)$, which is unbounded in $M$. The functions of the proof are functions of one variable; as functions of the first coordinate they give the same bound in $\R^d$ with the Euclidean norm.

\begin{proof}
The plan is a dichotomy: a large step diverges on a cliff, and a small step crawls on a plateau of slope just above $\eps$. Put $P=L_1M/L_0$ and $y_0=\log(1+P)$, so that $\eta^*=2y_0/(L_1M)$.

\emph{A large step diverges on the cliff.} Let $\eta\ge\eta^*$ and $f(x)=\spikescliff(x-c)$ with $c=x_0-y_0/L_1$. By \eqref{eq:spikes:cliff}, $f$ is $(L_0,L_1)$-smooth, $\abs{f'(x_0)}=\frac{L_0}{L_1}(e^{y_0}-1)=M$, the sublevel set $\{f\le f(x_0)\}$ is the interval $\abs{x-c}\le y_0/L_1$, on which $\abs{f'}\le M$, and $f(x_0)-\inf f=\frac{L_0}{L_1^2}(P-\log(1+P))<M/L_1\le\Delta_0$; so $f\in\cF$. In the variable $z_t=L_1(x_t-c)$ gradient descent reads $z_{t+1}=z_t\bigl(1-\eta L_0\varphi_1(\abs{z_t})\bigr)$, with the increasing function $\varphi_1$ of \cref{lem:spikes:local}. At $t=0$ we have $\eta L_0\varphi_1(y_0)=\eta L_0P/y_0=\eta L_1M/y_0\ge2$, hence $\abs{z_1}=\abs{z_0}(\eta L_0\varphi_1(\abs{z_0})-1)\ge\abs{z_0}$. Since $\varphi_1$ increases, the inequality $\eta L_0\varphi_1(\abs{z_t})\ge2$ propagates, and by induction $\abs{z_t}\ge y_0$ for every $t$, that is, $\abs{f'(x_t)}\ge M>\eps$.

\emph{A small step crawls on the plateau.} Let $\eta<\eta^*$. If $\Delta_0\le4\eps^2/(\sqrt3L_0)$, then $N\le0$ and there is nothing to prove, so assume the opposite, and let
\[
  \spikesplat(x)=\frac{4\eps^2}{L_0}\Bigl(\sqrt{1+\bigl(\tfrac{L_0x}{2\eps}\bigr)^2}-1\Bigr),\qquad
  \spikesplat'(x)=\frac{2\eps\,\sigma}{\sqrt{1+\sigma^2}},\qquad
  \spikesplat''(x)=\frac{L_0}{(1+\sigma^2)^{3/2}},\qquad \sigma=\frac{L_0x}{2\eps}.
\]
Then $\abs{\spikesplat'}<2\eps\le M$ and $0<\spikesplat''\le L_0$, so $f(x)=\spikesplat(x-c)$ is $(L_0,L_1)$-smooth for every $c$; choose $c<x_0$ with $f(x_0)=\Delta_0$, which puts $f$ in $\cF$. Solving $f(x_0)=\Delta_0$ with $a=L_0\Delta_0/(4\eps^2)$ gives $x_0-c=\frac{2\eps}{L_0}\sqrt{(1+a)^2-1}\ge\frac{2\eps}{L_0}a=\frac{\Delta_0}{2\eps}$. On the other hand $\abs{f'(x)}>\eps$ exactly when $\abs{x-c}>r_\eps=2\eps/(\sqrt3L_0)$. Each step moves the iterate by $\eta\abs{f'(x_t)}<2\eta\eps$, so $\abs{f'(x_t)}>\eps$ for every
\[
  t<\frac{x_0-c-r_\eps}{2\eta\eps}\ ,\qquad\text{and}\qquad
  \frac{x_0-c-r_\eps}{2\eta\eps}\ge\frac{\Delta_0-4\eps^2/(\sqrt3L_0)}{4\eta\eps^2}>\frac{\Delta_0-4\eps^2/(\sqrt3L_0)}{4\eta^*\eps^2}=N.\qedhere
\]
\end{proof}

\Textcite{ZhangHeSraJadbabaie2020} proved the first results of this kind: the lemma, with the factor $\varphi_1(c)$ in front of $L_1\norm g_*$ and a slightly larger one in front of $L_0$; a bound of the form $O(L_0\Delta_0/\eps^2+L_1^2\Delta_0/L_0)$ for clipped gradient descent, with larger absolute constants, under smoothness assumed on a neighbourhood of the sublevel set; and a lower bound of the form \eqref{eq:spikes:lower} for fixed steps, with $\log M+1$ in place of our logarithm. Our constants are smaller for two reasons: integrating the exact bound \eqref{eq:spikes:gradgrowth} gives the factor $\varphi_2(c)$, which is below $\frac12\varphi_1(c)$, in the descent inequality; and counting only the steps with $\norm{g_t}_*>\eps$ turns their sum of two terms into a maximum, which is never larger up to a constant, since $6L_1/\eps\le4L_0/\eps^2+\frac94L_1^2/L_0$ by the inequality of arithmetic and geometric means.

Normalized updates bound their step length by definition, so the same analysis applies to them, with the same weak dependence on $L_1$.

\begin{proposition}[Normalized steps are clipped by construction]\label[proposition]{prop:spikes:normalized}
Let $\Loss$ satisfy \cref{ass:spikes:L0L1} and $0<\gamma\le1/(3L_1)$.
\begin{enumerate}
\item Normalized steepest descent $\theta_{t+1}=\theta_t+\gamma\lmo(g_t)$ satisfies $\Loss(\theta_{t+1})\le\Loss(\theta_t)-\frac45\gamma\norm{g_t}_*+\frac35L_0\gamma^2$. Hence
\[
  \frac1T\sum_{t<T}\norm{\nabla\Loss(\theta_t)}_*\le\frac{5\Delta_0}{4\gamma T}+\frac{3L_0\gamma}4,
\]
and with $\gamma=\min\{1/(3L_1),2\eps/(3L_0)\}$ some $t<T$ has $\norm{\nabla\Loss(\theta_t)}_*\le\eps$ as soon as $T\ge\max\{15L_0\Delta_0/(4\eps^2),15L_1\Delta_0/(2\eps)\}$.
\item This holds for normalized gradient descent ($\ell_2$), for sign descent ($\ell_\infty$, with \eqref{eq:spikes:L0L1} read as $\norm{\nabla^2\Loss}_{\infty\to1}\le L_0+L_1\norm{\nabla\Loss}_1$), and for orthogonalized gradient descent on a weight matrix (the spectral norm, with the nuclear norm as dual).
\item Adam without bias correction, with $\epsilon=0$ and $\beta_1^2<\beta_2$, takes steps with $\norm{\theta_{t+1}-\theta_t}_\infty\le C\eta$, where $C=(1-\beta_1)/\sqrt{(1-\beta_2)(1-\beta_1^2/\beta_2)}$. If $\Loss$ satisfies \cref{ass:spikes:L0L1} in $\ell_\infty$ and $C\eta\le c/L_1$, every Adam step therefore satisfies \eqref{eq:spikes:local}.
\end{enumerate}
\end{proposition}

\begin{proof}
The plan is to repeat the proof of \cref{thm:spikes:clipped}(i) for a step of fixed length. For (i), $u=\gamma\lmo(g_t)$ has $\norm u\le\gamma$ and $\ip{g_t,u}=-\gamma\norm{g_t}_*$, also when $g_t=0$ and $\lmo(g_t)$ is an arbitrary point of the unit ball. \Cref{lem:spikes:local} with $c=\frac13$ gives
\[
  \Loss(\theta_{t+1})\le\Loss(\theta_t)-\gamma\norm{g_t}_*+\tfrac35(L_0+L_1\norm{g_t}_*)\gamma^2\le\Loss(\theta_t)-\tfrac45\gamma\norm{g_t}_*+\tfrac35L_0\gamma^2,
\]
using $\frac35L_1\gamma\le\frac15$. Summing over $t<T$ and dividing by $\frac45\gamma T$ gives the average. With the stated $\gamma$ the second term is at most $\eps/2$, and the first is at most $\eps/2$ once $T\ge5\Delta_0/(2\gamma\eps)$, which is the stated condition; the smallest of the $T$ gradients is at most their average. Part (ii) consists of instances of (i): the $\lmo$ of the three norms is $-g/\norm g_2$, $-\sign(g)$ and $-UV^\top$ (\cref{ch:steepest,ch:spectral}). For (iii), $\abs{m_{t,i}}/\sqrt{v_{t,i}}\le C$ for every coordinate by Chapter~2's bound on Adam's update (\cref{prop:sign:adambound}), so every step has $\ell_\infty$-length at most $C\eta$, and \cref{lem:spikes:local} applies to it.
\end{proof}

Part (iii) does not prove that Adam converges: its direction need not be a descent direction, and a rate needs the analyses of \cref{ch:adaptive}. It proves that Adam, like clipped descent, never takes a step that leaves the region where the local curvature bound holds, provided $C\eta\lesssim1/L_1$; for $(\beta_1,\beta_2)=(0.9,0.999)$ the constant is $C=7.27$, for $(0.9,0.95)$ it is $1.17$. \Cref{sec:spikes:anatomy} shows that Adam comes close to this bound exactly when a quiet block of weights wakes up.

The last result of the section is about noise. \Textcite{SimsekliSagunGurbuzbalaban2019} measured \term{heavy tails}\index{heavy-tailed noise} in the gradient noise of deep networks, and \textcite{ZhangEtAl2020Attention} found heavier tails for attention models than for convolutional networks and proposed clipping as the reason adaptive methods train them faster. If the noise has an infinite variance, stochastic gradient descent fails at the first step: on $\Loss(\theta)=\frac12\norm\theta_2^2$ with noise $\zeta_0$ of infinite second moment, $\E\Loss(\theta_1)=\frac12\E\norm{(1-\eta)\theta_0-\eta\zeta_0}_2^2=\infty$ for every $\eta>0$. Clipping restores a rate.

\begin{theorem}[Clipped SGD with heavy-tailed noise]\label{thm:spikes:heavy}
Let $\Loss$ be $L$-smooth in the Euclidean norm and bounded below, with $\Delta_0=\Loss(\theta_0)-\Loss^\star$, and let the stochastic gradients be unbiased, $\E[g_t\mid\theta_t]=\nabla\Loss(\theta_t)$, with a bounded moment of order $p\in(1,2]$: $\E[\norm{g_t}_2^p\mid\theta_t]\le\Gamma^p$. Clipped SGD $\theta_{t+1}=\theta_t-\eta\spikesclip_\tau(g_t)$ satisfies
\begin{equation}\label{eq:spikes:heavy}
  \frac1T\sum_{t<T}\E\norm{\nabla\Loss(\theta_t)}_2^2\le\frac{2\Delta_0}{\eta T}+\frac{\Gamma^{2p}}{\tau^{2(p-1)}}+L\eta\Gamma^p\tau^{2-p}.
\end{equation}
With $\tau=\Gamma T^{1/(3p-2)}$ and $\eta=(2\Delta_0/(L\Gamma^2))^{1/2}T^{-p/(3p-2)}$, the right side equals $(2\Gamma\sqrt{2L\Delta_0}+\Gamma^2)\,T^{-2(p-1)/(3p-2)}$, so $\min_{t<T}\E\norm{\nabla\Loss(\theta_t)}_2\lesssim T^{-(p-1)/(3p-2)}$.
\end{theorem}

\begin{proof}
The plan is to bound the bias and the second moment of a clipped gradient by the $p$th moment, and then to run the descent lemma. Fix $t$ and condition on $\theta_t$; write $g=g_t$ and $\hat g=\spikesclip_\tau(g)$. The bias: $\hat g$ and $g$ differ only where $\norm g_2>\tau$, by $\norm g_2-\tau<\norm g_2\le\norm g_2^p\tau^{1-p}$, so
\[
  \bigl\lVert\E\hat g-\nabla\Loss(\theta_t)\bigr\rVert_2\le\E\norm{\hat g-g}_2\le\frac{\Gamma^p}{\tau^{p-1}}\eqqcolon b .
\]
The second moment: $\norm{\hat g}_2^2=\min\{\norm g_2,\tau\}^2\le\norm g_2^p\tau^{2-p}$, by checking the cases $\norm g_2\le\tau$ and $\norm g_2>\tau$, so $\E\norm{\hat g}_2^2\le\Gamma^p\tau^{2-p}$. The descent lemma (\cref{lem:toolkit:descent}) gives $\Loss(\theta_{t+1})\le\Loss(\theta_t)-\eta\ip{\nabla\Loss(\theta_t),\hat g}+\frac L2\eta^2\norm{\hat g}_2^2$. The bias enters only the middle term, through $\ip{\nabla\Loss,\E\hat g}\ge\norm{\nabla\Loss}_2^2-\norm{\nabla\Loss}_2b\ge\frac12\norm{\nabla\Loss}_2^2-\frac12b^2$. Taking expectations,
\[
  \E\Loss(\theta_{t+1})\le\E\Loss(\theta_t)-\frac\eta2\E\norm{\nabla\Loss(\theta_t)}_2^2+\frac\eta2b^2+\frac L2\eta^2\Gamma^p\tau^{2-p}.
\]
Summing over $t<T$, using $\E\Loss(\theta_T)\ge\Loss^\star$ and dividing by $\eta T/2$ gives \eqref{eq:spikes:heavy}. With the stated $\tau$ and $\eta$, the first and third terms both equal $\Gamma\sqrt{2L\Delta_0}\,T^{-2(p-1)/(3p-2)}$ and the second equals $\Gamma^2T^{-2(p-1)/(3p-2)}$. The last claim follows because the smallest of the $T$ numbers $(\E\norm{\nabla\Loss(\theta_t)}_2)^2\le\E\norm{\nabla\Loss(\theta_t)}_2^2$ is at most their average.
\end{proof}

For $p=2$ the rate is $T^{-1/4}$, the rate of stochastic gradient descent with bounded variance; for $p\to1$ it degrades to nothing. The threshold must grow with the length of the run, like $T^{1/(3p-2)}$: a fixed threshold leaves the bias $\Gamma^p/\tau^{p-1}$, and clipping with noise is biased (\cref{ex:spikes:bias}). \Textcite{ZhangEtAl2020Attention} proved this rate and a matching lower bound, and \textcite{KoloskovaHendrikxStich2023} analysed the bias of a fixed threshold.

\section{The theory: the anatomy of a spike}\label{sec:spikes:anatomy}

A \term{loss spike}\index{loss spike} is a step too long for the curvature it lands in. \Cref{thm:spikes:clipped} says that a step of bounded length is safe; so the question is how a clipped, well-tuned run produces steps that are too long, or curvature that is too high. We take apart three mechanisms, one for each further defence in the practice box: Adam's state ($\beta_2$), the output logits (z-loss) and the attention logits (QK-norm, MuonClip).

\paragraph{The optimizer's state.} \Cref{ch:adaptive} showed that after a long stretch of tiny gradients a coordinate whose gradient jumps receives an Adam step of about $(1-\beta_1)/\sqrt{1-\beta_2}$ times the learning rate (\cref{prop:adaptive:quietloud}). That is the first step of a burst, not its peak.\sidenote{\Textcite{KingmaBa2015} observed that Adam's step is roughly bounded by the learning rate, and singled out the most severe case of sparsity, a gradient that was zero at all earlier steps, in which the step is $(1-\beta_1)/\sqrt{1-\beta_2}$ times the learning rate.}

\begin{proposition}[A quiet block wakes up]\label[proposition]{prop:spikes:burst}
Consider one coordinate of Adam with $\beta_1,\beta_2\in[0,1)$ and $\epsilon\ge0$, moments $m_t=\beta_1m_{t-1}+(1-\beta_1)g_t$ and $v_t=\beta_2v_{t-1}+(1-\beta_2)g_t^2$, and step $-\eta u_t$ with $u_t=m_t/(\sqrt{v_t}+\epsilon)$. Suppose that $\abs{m_{t_0-1}}\le\delta$ and $0\le v_{t_0-1}\le\delta^2$, as holds when $\abs{g_s}\le\delta$ for every $s<t_0$ and $m_{-1}=v_{-1}=0$, and that $g_s=\bar g\ne0$ for $t_0\le s<t_0+k$. At the $k$th loud step $t=t_0+k-1$,
\begin{equation}\label{eq:spikes:burstbounds}
  \sign(\bar g)\,u_t\le\frac{(1-\beta_1^k)\abs{\bar g}+\beta_1^k\delta}{\sqrt{1-\beta_2^k}\,\abs{\bar g}+\epsilon},
\end{equation}
and $\sign(\bar g)u_t\ge\bigl((1-\beta_1^k)\abs{\bar g}-\beta_1^k\delta\bigr)/\bigl(\sqrt{1-\beta_2^k}\,\abs{\bar g}+\beta_2^{k/2}\delta+\epsilon\bigr)$ whenever the numerator is nonnegative. If the coordinate was silent, $\delta=0$, then $u_t=\sign(\bar g)\,a_k$ with
\begin{equation}\label{eq:spikes:burst}
  a_k=\frac{1-\beta_1^k}{\sqrt{1-\beta_2^k}+\epsilon/\abs{\bar g}} .
\end{equation}
With $\epsilon=0$: $a_1=(1-\beta_1)/\sqrt{1-\beta_2}$; $a_k\to1$ as $k\to\infty$; $a_k\le C(\beta_1,\beta_2)=(1-\beta_1)/\sqrt{(1-\beta_2)(1-\beta_1^2/\beta_2)}$ for every $k$ when $\beta_1^2<\beta_2$; and $a_k\ge(1-\beta_1^k)/\sqrt{k(1-\beta_2)}$, so that for $\beta_1=0.9$, $\max_ka_k\ge a_{12}\ge0.207/\sqrt{1-\beta_2}$.
\end{proposition}

\begin{proof}
The plan is to unroll the two recursions over the loud steps. Over $k$ steps with the constant gradient $\bar g$, the first moment forgets its initial value at the rate $\beta_1$ and the second at the rate $\beta_2$:
\[
  m_t=\beta_1^km_{t_0-1}+(1-\beta_1^k)\bar g,\qquad v_t=\beta_2^kv_{t_0-1}+(1-\beta_2^k)\bar g^2,
\]
by induction on $k$, using $(1-\beta)\sum_{j<k}\beta^j=1-\beta^k$. Hence $\abs{\sign(\bar g)m_t-(1-\beta_1^k)\abs{\bar g}}\le\beta_1^k\delta$ and $(1-\beta_2^k)\bar g^2\le v_t\le(1-\beta_2^k)\bar g^2+\beta_2^k\delta^2$. For the upper bound, if $\sign(\bar g)m_t<0$ there is nothing to prove; otherwise divide the upper bound on $\sign(\bar g)m_t$ by the lower bound $\sqrt{1-\beta_2^k}\abs{\bar g}+\epsilon$ on $\sqrt{v_t}+\epsilon$. For the lower bound, divide the nonnegative lower bound on $\sign(\bar g)m_t$ by the upper bound on $\sqrt{v_t}+\epsilon$, using $\sqrt{x+y}\le\sqrt x+\sqrt y$. With $\delta=0$ the two bounds coincide and give \eqref{eq:spikes:burst}. The limit $a_k\to1$ is immediate. The bound by $C(\beta_1,\beta_2)$ is Chapter~2's bound on $\abs{m_t}/\sqrt{v_t}$ (\cref{prop:sign:adambound}), which holds for every sequence of gradients, here the sequence that is zero up to $t_0-1$ and $\bar g$ afterwards. Finally $1-\beta_2^k=(1-\beta_2)\sum_{j<k}\beta_2^j\le k(1-\beta_2)$, and $(1-0.9^{12})/\sqrt{12}=0.2071$.
\end{proof}

In words: after a long silence both moments are empty, and the coordinate goes through a \term{burst}\index{Adam!burst after a quiet period}. The first moment refills in about $1/(1-\beta_1)=10$ steps, the second in about $1/(1-\beta_2)$ steps, and in between $m$ has nearly its full size while $\sqrt v$ does not. The step grows for about $1.26/\log(1/\beta_1)\approx12$ steps (the maximizer of $(1-\beta_1^k)/\sqrt k$), and then decays to $1$ over about $1/(1-\beta_2)$ steps, as the second moment catches up:

\begin{center}\small
\begin{tabular}{@{}lccc@{}}
\toprule
$\beta_2$ & first step $a_1$ & peak $\max_ka_k$ (at $k$) & worst case $C(0.9,\beta_2)$\\
\midrule
$0.95$ & $0.45$ & $1.10$ ($k=20$) & $1.17$\\
$0.99$ & $1.00$ & $2.13$ ($k=13$) & $2.35$\\
$0.999$ & $3.16$ & $6.57$ ($k=12$) & $7.27$\\
\bottomrule
\end{tabular}
\end{center}

\noindent The burst reaches about nine tenths of the worst case over all gradient sequences, which is approached by gradients that grow geometrically. Bias correction does not help: it divides by $1-\beta_1^{t+1}$ and $1-\beta_2^{t+1}$ with $t$ counted from the start of training, and for $t\gg1/(1-\beta_2)$ both are $1$. The burst is precisely the transient that bias correction removes at the start of training (\cref{ch:warmup}), recreated for a coordinate whose memory has been emptied at a time when bias correction has long stopped acting; a bias correction that counted the age $k$ of the coordinate would remove it (\cref{ex:spikes:burst}). Without momentum the burst is a single step: for RMSProp, $\beta_1=0$, the largest $a_k$ is $a_1=1/\sqrt{1-\beta_2}$, which is $31.6$ at $\beta_2=0.999$. And $\epsilon$ caps the burst only for gradients with $\abs{\bar g}\lesssim\epsilon/\sqrt{1-\beta_2^k}$.

Which blocks go silent? Rows of an embedding matrix for tokens absent from recent batches, experts that a router stops choosing, units whose inputs vanish; in a milder form, every block whose gradient grows by a large factor within fewer than $1/(1-\beta_2)$ steps (\cref{ex:spikes:jump}). A whole block that wakes up together moves, coordinate by coordinate, by up to $6.6$ learning rates per step at $\beta_2=0.999$, against about one for steady gradients; whether the loss then spikes depends on the curvature of the block, which the proposition does not see.

\paragraph{The output logits.} The second mechanism is a direction in which the loss does not change at all.

\begin{proposition}[A flat direction and the z-loss]\label[proposition]{prop:spikes:zloss}
For logits $z\in\R^K$ let $Z(z)=\sum_ke^{z_k}$ and $\pi=\softmax(z)$, and let $\ell(z)=-\sum_ky_kz_k+\log Z(z)$ be the cross-entropy against a probability vector $y$ (a one-hot label, or the label frequencies of a batch).
\begin{enumerate}
\item $\ell(z+c\mathbf1)=\ell(z)$ for every $c\in\R$; $\ip{\nabla\ell(z),\mathbf1}=0$ and $\nabla^2\ell(z)\mathbf1=0$.
\item For an output layer $z=Wh+b$ with $W\in\R^{K\times d_{\mathrm{in}}}$, every gradient of the loss of a batch satisfies $\mathbf1^\top\nabla_W\ell=0$ and $\ip{\mathbf1,\nabla_b\ell}=0$. Hence an update that is a linear combination of gradients, multiplied on the right by matrices (gradient descent, momentum, right preconditioning), and also an orthogonalized update $-\eta\polar(G)$, leaves $\mathbf1^\top W$ and $\ip{\mathbf1,b}$ unchanged; a sign step on $b$ changes $\ip{\mathbf1,b}$ by $-\eta\sum_k\sign(\pi_k-y_k)$.
\item With the z-loss, $\ell_\alpha(z)=\ell(z)+\alpha(\log Z(z))^2$ with $\alpha>0$, the restriction to the flat line is a parabola: $\ell_\alpha(z+c\mathbf1)=\ell(z)+\alpha(\log Z(z)+c)^2$, with second derivative $2\alpha$ in $c$ and minimum at $\log Z=0$. If $y_k>0$ for every $k$, the only stationary point of $\ell_\alpha$ is $z=\log y$ (coordinatewise), where $\log Z=0$ and $\pi=y$.
\end{enumerate}
\end{proposition}

\begin{proof}
The plan is to compute the gradient and the Hessian of $\ell$ and to read everything off them. (i) $Z(z+c\mathbf1)=e^cZ(z)$, so $\log Z$ increases by $c$, exactly as much as $\sum_ky_kz_k$ does. Differentiating, $\nabla\ell=\pi-y$, whose coordinates sum to $1-1=0$, and $\nabla^2\ell=\diag(\pi)-\pi\pi^\top$, which maps $\mathbf1$ to $\pi-\pi=0$. (ii) For one example the chain rule gives $\nabla_W\ell=(\pi-y)h^\top$ and $\nabla_b\ell=\pi-y$, so $\mathbf1^\top\nabla_W\ell=0$ by (i); a batch gradient is an average of such terms. A linear combination $\sum_sc_sG_sR_s$ inherits $\mathbf1^\top G_s=0$. If $\mathbf1^\top G=0$, then $\mathbf1$ is orthogonal to the column space of $G$, which is spanned by the left singular vectors $U$, so $\mathbf1^\top\polar(G)=\mathbf1^\top UV^\top=0$. The sign step is a computation. (iii) The first claim follows from $\log Z(z+c\mathbf1)=\log Z(z)+c$ and (i). A stationary point satisfies $\nabla\ell_\alpha=\pi-y+2\alpha\log Z\,\pi=0$, that is, $(1+2\alpha\log Z)\pi=y$; summing the coordinates gives $1+2\alpha\log Z=1$, hence $\log Z=0$ and $\pi=y$, and then $z_k=\log\pi_k+\log Z=\log y_k$.
\end{proof}

Gradient descent therefore never moves the logits along this \term{flat direction}\index{flat direction of the softmax}, and neither does Muon (\cref{ex:spikes:muonmean}), whereas coordinatewise normalization does:\sidenote{In the recipes of \textcite{JordanEtAl2024Muon} and \textcite{LiuEtAl2025Muon}, Muon trains the hidden matrices and AdamW trains the embedding and output layers, so the output logits are still exposed to the drift.} the coordinates of $\pi-y$ sum to zero, their signs need not. Decoupled weight decay shrinks the flat component, and the z-loss adds a spring of stiffness $2\alpha$ that holds $\log Z$ near $0$. In exact arithmetic the drift is harmless, since the loss does not see it; it matters through finite precision, for instance in bfloat16, which keeps $8$ significant bits, so that logits near $100$ are rounded to multiples of $0.5$ and the probabilities computed from them are off by factors up to $e^{\pm0.25}$.\sidenote{\Textcite{ZophEtAl2022} introduced a z-loss on the logits of the router of a mixture of experts, motivated by roundoff in the exponentials of large logits. Whether the output-logit divergence of \cref{obs:spikes:proxies} is a precision effect has, as far as we know, not been established.} How strong a spring is needed depends on how fast the drift is, which the proposition does not say; \cref{sec:spikes:reading} measures it in a small model.

\paragraph{The attention logits.} The third mechanism is a scale, that of the \term{attention logits}\index{attention logits}.

\begin{proposition}[Attention logits, QK-norm and QK-clip]\label[proposition]{prop:spikes:qknorm}
For an attention head with query and key weights $W_Q,W_K\in\R^{d_{\mathrm h}\times d_{\mathrm{in}}}$ and inputs $x,x'\in\R^{d_{\mathrm{in}}}$, let $s(x,x')=\ip{W_Qx,W_Kx'}/\sqrt{d_{\mathrm h}}$.
\begin{enumerate}
\item The logit is bilinear in $(W_Q,W_K)$: replacing $(W_Q,W_K)$ by $(\rho W_Q,\rho W_K)$ multiplies every logit by $\rho^2$, and $\abs{s(x,x')}\le\norm{W_Q}_{\op}\norm{W_K}_{\op}\norm x_2\norm{x'}_2/\sqrt{d_{\mathrm h}}$.
\item With QK-norm, $\tilde s(x,x')=\ip{\omega_q\odot\mathcal N(W_Qx),\,\omega_k\odot\mathcal N(W_Kx')}/\sqrt{d_{\mathrm h}}$, where $\mathcal N$ is RMS normalization or a LayerNorm without bias and without $\epsilon$, and $\omega_q,\omega_k\in\R^{d_{\mathrm h}}$ are gains,
\[
  \abs{\tilde s(x,x')}\le\sqrt{d_{\mathrm h}}\,\norm{\omega_q}_\infty\norm{\omega_k}_\infty
\]
for all $W_Q,W_K,x,x'$ for which the normalizations are defined, with equality when $\mathcal N(W_Qx)=\mathcal N(W_Kx')$ and both gains are the same constant vector.
\item If the logits of one row of attention lie in an interval of length $R$, then each of the $n_{\mathrm{key}}$ attention weights of the row is at most $1/(1+(n_{\mathrm{key}}-1)e^{-R})$, and the row's entropy is at least $\log(1+(n_{\mathrm{key}}-1)e^{-R})$.
\item (QK-clip.) If the largest logit of a head on a batch is $s_{\max}>\bar s>0$, then replacing $(W_Q,W_K)$ by $(\rho W_Q,\rho W_K)$ with $\rho=\sqrt{\bar s/s_{\max}}$ multiplies every logit of the head by $\bar s/s_{\max}$, so that the largest logit on that batch becomes exactly $\bar s$.
\end{enumerate}
\end{proposition}

\begin{proof}
The plan is to treat the four parts in turn; only (ii) and (iii) need a computation. (i) Bilinearity is immediate, and Cauchy--Schwarz gives $\abs{\ip{W_Qx,W_Kx'}}\le\norm{W_Qx}_2\norm{W_Kx'}_2\le\norm{W_Q}_{\op}\norm x_2\norm{W_K}_{\op}\norm{x'}_2$. (ii) RMS normalization maps $v$ to $v/(\norm v_2/\sqrt{d_{\mathrm h}})$, and a LayerNorm without bias and $\epsilon$ maps it to $(v-\bar v\mathbf1)/(\norm{v-\bar v\mathbf1}_2/\sqrt{d_{\mathrm h}})$, with $\bar v$ the mean of the coordinates; both images have Euclidean norm $\sqrt{d_{\mathrm h}}$. Hence, with $a=\mathcal N(W_Qx)$ and $b=\mathcal N(W_Kx')$,
\[
  \Bigl\lvert\sum_i\omega_{q,i}\omega_{k,i}a_ib_i\Bigr\rvert\le\norm{\omega_q}_\infty\norm{\omega_k}_\infty\sum_i\abs{a_ib_i}\le\norm{\omega_q}_\infty\norm{\omega_k}_\infty\norm a_2\norm b_2=\norm{\omega_q}_\infty\norm{\omega_k}_\infty\,d_{\mathrm h},
\]
and dividing by $\sqrt{d_{\mathrm h}}$ gives the bound; in the stated case every inequality is an equality. (iii) Let the logits be $s_1,\dots,s_{n_{\mathrm{key}}}$. Every $s_i$ is at least $s_j-R$, so the weight of key $j$ is $e^{s_j}/\sum_ie^{s_i}\le e^{s_j}/(e^{s_j}+(n_{\mathrm{key}}-1)e^{s_j-R})$, which is the bound; the entropy $\sum_jp_j\log(1/p_j)$ is at least $\log(1/\max_jp_j)$. (iv) follows from (i).
\end{proof}

Part (i) is the source of the growth: the logit is quadratic in the scale of $(W_Q,W_K)$, and under Adam the scale can grow linearly in time. Every coordinate moves by at most $C\eta$ per step (\cref{prop:spikes:normalized}(iii)), so after $T$ steps $\norm{W_Q-W_Q^{(0)}}_{\op}\le\norm{W_Q-W_Q^{(0)}}_F\le C\eta T\sqrt{d_{\mathrm h}d_{\mathrm{in}}}$, and for inputs normalized to $\norm x_2=\sqrt{d_{\mathrm{in}}}$ the bound of (i) grows like $(C\eta T)^2d_{\mathrm{in}}^2\sqrt{d_{\mathrm h}}$. The Frobenius bound is attained by rank-one sign updates $-\eta\sign(\delta)\sign(x)^\top$, which is what Adam does with a consistent rank-one gradient. Under QK-norm the same argument applied to the gains gives $\sqrt{d_{\mathrm h}}(\norm{\omega^{(0)}}_\infty+C\eta T)^2$: the growth rate of the worst case falls by the factor $d_{\mathrm{in}}^2$. QK-norm does not forbid sharp attention; with unit gains and $d_{\mathrm h}=128$ the logits may range over an interval of length $2\sqrt{128}\approx23$, enough for nearly one-hot rows. What it removes is growth driven by the scale of the weight matrices, and part (iii) says that sharp attention needs large logits: entropy collapse in the sense of \textcite{ZhaiEtAl2023} is a growth of the logit range.\sidenote{\Textcite{HenryEtAl2020} normalized queries and keys to unit Euclidean norm with a learned temperature for machine translation; \textcite{DehghaniEtAl2023} used LayerNorm on queries and keys to remove the divergence of ViT-22B.} QK-clip acts on the same scale after the fact: by (iv) it is the smallest rescaling of $(W_Q,W_K)$, split evenly between them, that brings the largest logit of the batch back to the threshold.

\begin{fieldguide}{MuonClip}{\textcite{KimiTeam2025}}
\fgrow{update}{a Muon step on every hidden matrix (\cref{ch:spectral}); then, for each attention head whose largest logit $s_{\max}$ on the batch exceeds $\bar s$, $W_Q\leftarrow\rho W_Q$ and $W_K\leftarrow\rho W_K$ with $\rho=\sqrt{\bar s/s_{\max}}$}
\fgrow{norm}{spectral, for the Muon step}
\fgrow{preconditioner}{none beyond orthogonalization}
\fgrow{averages}{Muon's momentum}
\fgrow{state}{one momentum matrix per weight; the largest logit of each head, read off the forward pass}
\fgrow{in the wild}{$\bar s=100$ in Kimi K2; Muon itself as scaled by \textcite{LiuEtAl2025Muon}}
\fgrow{caution}{acts on the weights after the step, not on the gradient, and only on heads above the threshold}
\end{fieldguide}

\section{Reading the theory}\label{sec:spikes:reading}

\paragraph{The price of a fixed step.} \Cref{fig:spikes:separation} puts the two theorems side by side on the simplest function that has both a cliff and a plateau, $f(x,y)=\spikescliff(x)+\spikesplat(y)$ with $L_0=L_1=1$, which is $(1,1)$-smooth in the Euclidean norm because its Hessian is diagonal. The largest step with which gradient descent does not diverge is exactly $\eta^*=2\log(1+L_1M/L_0)/(L_1M)$, the cliff's threshold from the proof of \cref{thm:spikes:lower}: $1.39$ at $M=1$, $0.092$ at $M=100$ and $2.8\cdot10^{-5}$ at $M=10^6$. Gradient descent at $0.9\eta^*$ then crawls along the plateau, and its number of steps equals the gradient-flow time of the plateau divided by the step size to within one step, because the steps are small compared with $1/L_0$: from $81$ steps at $M=1$ to $4.07\cdot10^6$ at $M=10^6$, growing like $M/\log M$. Clipped descent with the parameters of \cref{thm:spikes:clipped} needs $203$ steps at $M=1$ and $241$ at $M=10^6$: the plateau costs it $2\cdot101$ steps at the step size $1/(2L_0)$, and the cliff $3\log(1+M)$ normalized steps of length $\gamma=1/3$. Its largest convergent step size does not depend on $M$ at all. For every start it lies between $L_1\gamma/(L_0(e^{L_1\gamma/2}-1))=1.838$, below which clipped descent cannot get caught in a two-cycle across the bottom of the valley, and $2/L_0$, above which it cannot converge at the bottom; the measured values fill this band.

\begin{figure}[t]
\begin{fullwidth}
\centering
\begin{tikzpicture}
\begin{groupplot}[group style={group size=2 by 1,horizontal sep=18mm},wplot,width=82mm,height=52mm,xmode=log,xmin=0.7,xmax=1.5e6]
\nextgroupplot[ymode=log,xlabel={largest gradient $M$},ylabel={steps to $\norm{\nabla f}_2\le\eps$},ymin=40,ymax=1e7,legend style={at={(0.03,0.97)},anchor=north west,font=\scriptsize}]
\addplot[theory] table[x=M,y=gdtheory] {figures/data/10-spikes-separation.dat};
\addplot[simulation] table[x=M,y=gdsim] {figures/data/10-spikes-separation.dat};
\addplot[theory2] table[x=M,y=cliptheory] {figures/data/10-spikes-separation.dat};
\addplot[simulation2] table[x=M,y=clipsim] {figures/data/10-spikes-separation.dat};
\legend{gradient flow$/\eta$,{GD, $\eta=0.9\eta^*$},flow$/\eta$ $+$ cliff steps,clipped GD}
\nextgroupplot[ymode=log,xlabel={largest gradient $M$},ylabel={largest stable step size},ymin=1e-5,ymax=5,legend style={at={(0.03,0.03)},anchor=south west,font=\scriptsize}]
\fill[seamist] (axis cs:0.7,1.83796) rectangle (axis cs:1.5e6,2);
\addplot[theory] table[x=M,y=thrtheory] {figures/data/10-spikes-separation.dat};
\addplot[simulation] table[x=M,y=thrsim] {figures/data/10-spikes-separation.dat};
\addplot[simulation2] table[x=M,y=clipthr] {figures/data/10-spikes-separation.dat};
\legend{$\eta^*=2\log(1+M)/M$,GD (bisection),clipped GD (scan)}
\end{groupplot}
\end{tikzpicture}
\end{fullwidth}
\caption{Clipped and fixed-step gradient descent on a cliff followed by a plateau, $f(x,y)=\spikescliff(x)+\spikesplat(y)$ with $L_0=L_1=1$ and $\eps=0.05$, started where the slope of the cliff is $M$ and the height of the plateau is $1$. Left: steps until the gradient norm is at most $\eps$, for gradient descent at $0.9$ times its largest stable step (circles) and clipped descent with $\eta=1/(2L_0)$, $\gamma=1/(3L_1)$ (squares); curves: the gradient-flow time of the plateau, $101.15$, divided by the step size, plus for clipped descent the $3\log(1+M)-3\log\frac53$ normalized steps down the cliff. Right: the largest step size with which gradient descent does not diverge (circles, bisection) against $\eta^*$ (curve), and the largest with which clipped descent converges (squares, scan in steps of $0.0025$), against the band $[1.838,2]$ predicted for every start (shaded). Deterministic; script \texttt{10-spikes-separation.py}.}
\label{fig:spikes:separation}
\end{figure}

The agreement is exact because the functions are those of the proofs; it is a check of the arithmetic, not of the idealization. It fails as soon as the step sizes stop being small compared with $1/L_0$ on the plateau, where the flow approximation of the step count breaks down, and it says nothing about functions that are not separable.

\paragraph{Where to clip.} In units of the gradient norm, the threshold of \cref{thm:spikes:clipped} is $\gamma/\eta=2L_0/(3L_1)$: clipping should begin where the term $L_1\norm{\nabla\Loss}_*$ of the curvature bound starts to dominate $L_0$. That is a property of the loss, not of the optimizer, and the theory gives no reason for a universal value such as $1.0$; global gradient norms change with the normalization of the loss and with the number of parameters. What it predicts is that a good threshold is inactive where the landscape is flat and active on cliffs and during spikes, and that the price of a threshold set too low is the factor $1/(\gamma\eps)$ of \cref{thm:spikes:clipped}(ii): the run is slowed wherever the gradient exceeds the threshold.

\paragraph{Clipping in front of Adam.} Adam's update is unchanged when all gradients are multiplied by the same factor (up to $\epsilon$), and \cref{prop:spikes:normalized}(iii) bounds its step whatever the gradient. So clipping in front of Adam does not bound Adam's step; it bounds the weight that one large gradient receives in $m$ and $v$. When the global gradient norm jumps by a large factor at one step and the threshold sits near its old value, the clipped gradient keeps the old norm, and the loud event that would start a burst is suppressed. A block that wakes up while the global norm stays below the threshold is not affected: the mechanism of \cref{prop:spikes:burst} survives global clipping.

\begin{figure}[t]
\begin{fullwidth}
\centering
\begin{tikzpicture}
\begin{groupplot}[group style={group size=2 by 1,horizontal sep=18mm},wplot,width=82mm,height=52mm]
\nextgroupplot[xmode=log,xmin=1,xmax=3000,ymin=0,ymax=7.8,xlabel={steps $k$ since the block woke up},ylabel={Adam step $\abs{\Delta\theta_i}/\eta$},legend style={at={(0.97,0.97)},anchor=north east,font=\scriptsize},legend columns=2]
\addplot[guide,domain=1:3000] {7.2703};
\addplot[guide,domain=1:3000] {2.3452};
\addplot[guide,domain=1:3000] {1.1650};
\addplot[theory3] table[x=k,y=th999] {figures/data/10-spikes-burst-k.dat};
\addplot[simulation,each nth point=2] table[x=k,y=sim999] {figures/data/10-spikes-burst-k.dat};
\addplot[theory2] table[x=k,y=th99] {figures/data/10-spikes-burst-k.dat};
\addplot[simulation2,each nth point=2] table[x=k,y=sim99] {figures/data/10-spikes-burst-k.dat};
\addplot[theory] table[x=k,y=th95] {figures/data/10-spikes-burst-k.dat};
\addplot[simulation,mark=triangle*,mark size=1.8pt,each nth point=2] table[x=k,y=sim95] {figures/data/10-spikes-burst-k.dat};
\legend{,,,{$a_k$, $\beta_2=0.999$},Adam,{$a_k$, $\beta_2=0.99$},Adam,{$a_k$, $\beta_2=0.95$},Adam}
\nextgroupplot[xmode=log,ymode=log,x dir=reverse,xmin=7e-5,xmax=0.14,ymin=0.25,ymax=30,xlabel={$1-\beta_2$},ylabel={largest step $/\eta$},legend style={at={(0.03,0.97)},anchor=north west,font=\scriptsize}]
\addplot[guide] table[x=omb2,y=worst] {figures/data/10-spikes-burst-peak.dat};
\addplot[theory] table[x=omb2,y=peak] {figures/data/10-spikes-burst-peak.dat};
\addplot[theory3,densely dotted] table[x=omb2,y=a1] {figures/data/10-spikes-burst-peak.dat};
\addplot[theory2] table[x=omb2,y=peakeps] {figures/data/10-spikes-burst-peak.dat};
\addplot[simulation] table[x=omb2,y=sim] {figures/data/10-spikes-burst-peak.dat};
\addplot[simulation2] table[x=omb2,y=simeps] {figures/data/10-spikes-burst-peak.dat};
\legend{{worst case $C(\beta_1,\beta_2)$},peak $\max_ka_k$,first step $a_1$,{peak, $\abs{\bar g}=10\epsilon$},Adam,{Adam, $\abs{\bar g}=10\epsilon$}}
\end{groupplot}
\end{tikzpicture}
\end{fullwidth}
\caption{Adam's step on a block of $32$ weights of a linear regression whose input features are absent for the first $10^5$ steps of training, so that the block's moments are zero, and present afterwards ($\beta_1=0.9$, $\epsilon=10^{-8}$, bias correction counted from the start of training, batches of $4096$, step size $10^{-6}$, labels $\ip{\theta_\star,x}$ plus unit Gaussian noise with $\theta_\star\in\{\pm1\}^{32}$). Left: the median over the block of $\abs{\Delta\theta_i}/\eta$ at the $k$th step after the features appear (points, every other measured $k$ shown), against $a_k$ of \eqref{eq:spikes:burst} (curves); dashed lines: the worst case $C(0.9,\beta_2)$. Right: the largest such step over the first $80$ steps against $1-\beta_2$ (points), with the burst peak $\max_ka_k$, the first step $a_1$ and the worst case; the lower curve and squares repeat the computation with the loss multiplied by $10^{-7}$, so that $\abs{\bar g}=10\epsilon$. One run per point; seeds in \texttt{10-spikes-burst.py}.}
\label{fig:spikes:burst}
\end{figure}

\paragraph{The burst.} \Cref{fig:spikes:burst} shows the anatomy of \cref{prop:spikes:burst} in a training run. The measured steps follow $a_k$ to within one percent: they rise for about twelve steps, peak at $6.59$, $2.13$ and $1.09$ learning rates for $\beta_2=0.999,0.99,0.95$ (predicted $6.57$, $2.13$, $1.10$), and decay to $1$ over about $1/(1-\beta_2)$ steps; the small differences come from minibatch noise, which enters $v$ but not, on average, $m$. Across $\beta_2$ the peak follows $0.207/\sqrt{1-\beta_2}$ for long memories, stays below the worst case, and is roughly twice the single-jump value $a_1$ of \cref{ch:adaptive}. Lowering $\beta_2$ from $0.999$ to $0.95$ shrinks the largest burst six-fold, from $6.6$ to $1.1$ learning rates per coordinate. This is a mechanism behind the choice $\beta_2=0.95$ in large language models and behind \cref{obs:spikes:stale}; that it is \emph{the} reason for the choice is a suggestion, not a result, and the price of a short memory is a noisier preconditioner (\cref{ch:adaptive}). A larger $\epsilon$ also damps the burst, but only for coordinates whose gradient is not much larger than $\epsilon/\sqrt{1-\beta_2^k}$: at $\abs{\bar g}=10\epsilon$ the peak is roughly halved near $\beta_2=0.999$ (right panel), and wherever gradients are smaller than $\epsilon$ Adam becomes gradient descent with step size $\eta/\epsilon$.

\paragraph{How fast the logits drift.} \Cref{prop:spikes:zloss} says that the z-loss removes the flat direction; it does not say whether $\alpha=10^{-4}$ is enough. In a unigram model, a softmax over $K=1000$ classes with Zipf frequencies and logits as the only parameters, trained from zero with minibatches of $64$ labels, stochastic gradient descent keeps the mean logit fixed exactly, as (ii) requires. Adam with $\beta_2=0.95$ does not: $\log Z$ falls steadily, by $0.12$ step sizes per step, while the cross-entropy stays within $0.14$ of its minimum (script \texttt{10-spikes-drift.py}). The drift is a minibatch effect of coordinatewise normalization: it vanishes in full batch, it falls with the batch size ($0.21$, $0.12$, $0.05$, $0.02$, $0.009$ step sizes per step at batches of $16$ to $4096$) and with $\beta_2$ ($0.02$ at $0.99$, $0.001$ at $0.999$), and it comes from the skewed distribution of a rare class's gradient, small and positive in most batches and large and negative in the few that contain it. A z-loss of $10^{-4}$ leaves the drift unchanged; a z-loss of $0.1$ stops it at the price of an excess cross-entropy of $0.38$, and decoupled weight decay of $0.1$ stops it at a price of $0.03$. In this model, then, the coefficient used in practice is far too weak to pin $\log Z$, and weight decay does the pinning; why $10^{-4}$ suffices in large runs is not explained by this mechanism.

\section{What the theory misses}\label{sec:spikes:misses}

\paragraph{What was idealized away.} The main theorems have exact gradients. With noise, a fixed threshold biases the clipped gradient (\cref{ex:spikes:bias}), and \cref{thm:spikes:heavy} needs a threshold that grows with the length of the run and only $L$-smoothness; a theorem that combines heavy tails with $(L_0,L_1)$-smoothness and a fixed practical threshold is more delicate. The constants $L_0$ and $L_1$ are global and fixed, whereas the curvature of a network changes during training, and the assumption of two continuous derivatives excludes ReLU networks. The lower bound is a statement about a class of functions, not about a given network. The burst proposition assumes an exactly constant loud gradient after an exactly silent period; real blocks are rarely silent, and \eqref{eq:spikes:burstbounds} quantifies what a residual state does. The propositions on logits are about invariances and bounds, not dynamics. Finite precision appears nowhere, although it is a candidate for the step from a drifted logit to a divergence.

\paragraph{Where it breaks.} (1) \Cref{ass:spikes:L0L1} fails globally for networks with rescaling symmetries. In the product model $\frac12(uv-1)^2$ of volume 2 of this series, \emph{The Hydrogen Atoms of Deep Learning}, Chapter 15, the minimizers $(u,1/u)$ have zero gradient and sharpness $u^2+u^{-2}$, so no constants $L_0,L_1$ work (\cref{ex:spikes:product}); the assumption holds only on bounded sets, with an $L_0$ that grows with the bound, and then the guarantee depends on how far the weights travel. \Cref{obs:spikes:L0L1} is a correlation along trajectories, not a derivation. (2) Heavy-tailed noise is not the main reason that adaptive methods train transformers faster: \textcite{KunstnerEtAl2023} found that the gap between SGD and Adam persists, and even grows, in full batch, where there is no noise, while sign descent with momentum closes most of it (\cref{ch:transformers}). The case for clipping via \cref{thm:spikes:heavy} is therefore weaker than the case via \cref{thm:spikes:clipped}. (3) Real spikes have interacting causes. In PaLM, neither the batches nor the parameter state alone produced a spike (\cref{obs:spikes:palm}); in OPT, divergences coincided with a collapse of the loss scale of mixed precision \parencite{ZhangEtAl2022OPT}; \textcite{TakaseEtAl2023} (unrefereed) relate spikes to the norms of the Jacobians of the layers and remove them with changes to initialization and to the scale of embeddings, a mechanism of architecture outside all three propositions. (4) In our unigram model the z-loss coefficient of practice does not pin the logits, so the observed benefit of $10^{-4}$ in large runs is unexplained by \cref{prop:spikes:zloss}. (5) MuonClip was needed because, in Kimi K2's experiments, attention logits grew more often with Muon than with AdamW \parencite{KimiTeam2025}; \cref{prop:spikes:qknorm} explains why rescaling removes the symptom, not why an orthogonalized optimizer produces it. (6) That a threshold of $1.0$ is rarely active late in training, except during spikes, is folklore; we know no systematic study.

\paragraph{Repairs and open problems.} The theory of generalized smoothness has grown quickly: clipping with momentum and sharper constants \parencite{ZhangJinFangWang2020}; a symmetric form of \eqref{eq:spikes:L0L1} that needs no second derivatives \parencite{ChenZhouLiangLu2023}; curvature bounded by a general function of the gradient norm, under which gradient descent with a small fixed step converges on sublevel sets \parencite{LiQianTianRakhlinJadbabaie2023}; Adam under such assumptions \parencite{LiRakhlinJadbabaie2023}; generalized sign descent \parencite{CrawshawEtAl2022}; and sharper complexities for gradient methods in the nonconvex and convex cases \parencite{VankovEtAl2025,GorbunovEtAl2025}. For heavy tails there are accelerated and high-probability versions of clipping \parencite{GorbunovDanilovaGasnikov2020} and an analysis of the bias of fixed thresholds \parencite{KoloskovaHendrikxStich2023}. Against the burst, the repairs act on the update: Adafactor-style clipping of the update by its root mean square \parencite{ShazeerStern2018}, adopted for AdamW as StableAdamW \parencite{WortsmanEtAl2023}; spike-aware variants of Adam that clip gradient entries which are large relative to the second-moment estimate and reset the moments \parencite{HuangEtAl2025SPAM}; and Adam-atan2, which replaces $m/(\sqrt v+\epsilon)$ by a bounded function without $\epsilon$ \parencite{EverettEtAl2024}. Against the logits, z-loss, QK-norm, QK-clip, and reparametrizations that control the spectral norm of the attention logits \parencite{ZhaiEtAl2023}. Open: a solvable model in which loss spikes arise and recur without being put in by hand (\cref{ex:spikes:spike}); a derivation of $(L_0,L_1)$-type bounds for a transformer, with constants that depend on the gains of its normalizations; a principled clipping threshold (\cref{ex:spikes:threshold}); and why orthogonalized updates grow attention logits.

\section{Back in the wild}\label{sec:spikes:back}

\fieldnotes{Clip the global gradient norm at $1.0$; $\beta_2=0.95$; z-loss $10^{-4}(\log Z)^2$; QK-norm; MuonClip; after a spike, restart and skip batches.}
{$(L_0,L_1)$-smoothness $\norm{\nabla^2\Loss}\le L_0+L_1\norm{\nabla\Loss}_*$; Adam's moments after a quiet period; the flat direction of the softmax; the bilinear attention logit.}
{A sharp local smoothness lemma. Clipped steepest descent: at most $\max\{4L_0\Delta_0/\eps^2,6L_1\Delta_0/\eps\}$ steps with $\norm g_*>\eps$. Every fixed step: $\gtrsim L_1M\Delta_0/(\eps^2\log(1+L_1M/L_0))$ on some function of the class. Normalized steps are clipped by construction. Clipped SGD at $T^{-(p-1)/(3p-2)}$ with $p$th moments. Exactly: the burst $a_k$, the pinning by the z-loss, the bound of QK-norm by its gains.}
{GD's safe step falls like $\log M/M$, clipped descent's does not; clip where $L_1\norm{\nabla\Loss}_*\approx L_0$; Adam's burst peaks near $0.21/\sqrt{1-\beta_2}$ learning rates about $12$ steps after a block wakes up, six times smaller at $\beta_2=0.95$ than at $0.999$; global clipping does not stop a block's burst.}
{When spikes happen; interacting causes (data, state, precision); $(L_0,L_1)$-smoothness fails under rescaling symmetries; why $10^{-4}$ and $1.0$.}
{Bound the length of every step, by clipping or by normalized updates (proved in the idealization). Prefer a short $\beta_2$ or update clipping when blocks can wake up (suggested). Use z-loss, weight decay and QK-norm to remove directions of growth that the loss cannot see (proved as mechanisms; their effect on spikes is observed). The values $1.0$ and $10^{-4}$ are folklore.}

The theory of this chapter says one thing with confidence: a step whose length is bounded is safe on landscapes whose curvature grows with the gradient, and a fixed step size is not. That is why clipping works, why normalized updates (\cref{ch:steepest,ch:sign,ch:spectral}) do not need it, and why Adam, whose step is bounded by Chapter~2's constant, is less exposed to cliffs than gradient descent; in front of Adam, clipping acts on the moments rather than on the step. It does not tell a practitioner where to put the threshold beyond ``where the curvature starts to grow'', and it predicts no spike. What it adds to the folklore of spikes is a set of exact mechanisms, each of which makes a step too long: an optimizer state that lags its gradient, a direction the loss cannot see, a scale that enters the logits quadratically. Each defence in the practice box removes one of them, and none removes all.

The chapter closes Part III. The stability threshold $2/\eta$ of \cref{ch:eos} bounds the curvature that a fixed step can stand; warmup (\cref{ch:warmup}) grows the step as the curvature falls with the loss; clipping shrinks the step when the curvature rises with the gradient. The loss-dependent curvature of \cref{ch:warmup} and the gradient-dependent curvature of this chapter are the same idea read two ways, and both hold for exponential-type losses. Weight decay (\cref{ch:decay}) is a further defence of the kind studied here: it pins the flat direction of the logits and bounds the scales that enter the attention logits. And the heavy tails of \cref{thm:spikes:heavy} return in \cref{ch:transformers}, where they turn out not to be the main reason that Adam beats gradient descent on transformers.

\begin{exercises}
\item \consolidation (Constants.) (a) Show that $x^4$ is $(L_0,L_1)$-smooth if and only if $L_0L_1^2\ge16$. (b) Show that $e^x$ is $(0,1)$-smooth and $L$-smooth for no $L$. (c) Show that $\log(1+e^{-x})$ is $(\frac14,0)$-smooth and $(0,1)$-smooth. (d) Show that $\sum_ia_i(\cosh\theta_i-1)$ with $a_i>0$ is $(\max_ia_i,1)$-smooth in the Euclidean norm and $(\sum_ia_i,1)$-smooth in $\ell_\infty$. (e) Show that $\frac12\theta^\top H\theta$ with $H\ne0$ is $(L_0,L_1)$-smooth only if $L_0\ge\norm H$. \hint{in (a), minimize $L_0+4L_1x^3-12x^2$ over $x\ge0$; in (c), $\sigma(x)(1-\sigma(x))\le\min\{\frac14,1-\sigma(x)\}$ for the logistic function $\sigma$.}\label{ex:spikes:constants}

\item \consolidation (The lemma is sharp.) (a) Verify the equality case of \cref{lem:spikes:local} for the cliff. (b) Deduce that neither $\varphi_1(c)$ nor $\varphi_2(c)$ in \eqref{eq:spikes:local} can be replaced by a smaller constant. (c) Show that for $\cosh x-1$, which is $(1,1)$-smooth, the bound \eqref{eq:spikes:gradgrowth} is not attained at $\theta=0$, and find the ratio of the two sides as $u\to\infty$.\label{ex:spikes:sharp}

\item \consolidation (Clipping is a trust region.) (a) Show that the clipped step $-\min\{\eta,\gamma/\norm g_2\}g$ minimizes $\ip{g,u}+\norm u_2^2/(2\eta)$ over the ball $\norm u_2\le\gamma$. (b) Show that clipping by value, $u_i=-\sign(g_i)\min\{\eta\abs{g_i},\gamma\}$, minimizes the same function over the box $\norm u_\infty\le\gamma$. (c) Show that clipping by value can change the direction of the step, and global clipping cannot.\label{ex:spikes:trust}

\item \consolidation (The burst.) (a) Show that if the moments were bias-corrected with the age of the coordinate, $\hat m_t=m_t/(1-\beta_1^k)$ and $\hat v_t=v_t/(1-\beta_2^k)$, the step in \cref{prop:spikes:burst} with $\delta=0$ and $\epsilon=0$ would be $\sign(\bar g)$ at every loud step. (b) For $\beta_1=0.9$, show that $(1-\beta_1^k)/\sqrt k$ is largest at $k=12$, and deduce $\max_ka_k\ge6.55$ for $\beta_2=0.999$. (c) For RMSProp, $\beta_1=0$, show that the largest $a_k$ is $a_1=1/\sqrt{1-\beta_2}$. (d) Show that for $\beta_2\to1$ with $\beta_1$ fixed, $\sqrt{1-\beta_2}\,\max_ka_k$ converges to $\max_{k\ge1}(1-\beta_1^k)/\sqrt k$ and $\sqrt{1-\beta_2}\,C(\beta_1,\beta_2)$ to $\sqrt{(1-\beta_1)/(1+\beta_1)}$; for $\beta_1=0.9$ the ratio is $0.903$.\label{ex:spikes:burst}

\item \consolidation (The z-loss.) (a) Show that $\nabla(\log Z)^2=2(\log Z)\pi$ and $\nabla^2(\log Z)^2=2\pi\pi^\top+2(\log Z)(\diag(\pi)-\pi\pi^\top)$. (b) Show that gradient descent on the bias $b$ of a softmax with the loss $\ell+\alpha(\log Z)^2$ changes $\ip{\mathbf1,b}$ by exactly $-2\alpha\eta\log Z$ per step. (c) With decoupled weight decay $\lambda$ instead of the z-loss, show that $\ip{\mathbf1,b}$ is multiplied by $1-\eta\lambda$ per step of gradient descent.\label{ex:spikes:zgrad}

\item \variation (No $(L_0,L_1)$ for the product model.) Let $\Loss(u,v)=\frac12(uv-1)^2$. (a) Show that no constants $L_0,L_1$ make $\Loss$ $(L_0,L_1)$-smooth in the Euclidean norm. (b) Show that $\norm{\nabla^2\Loss(u,v)}_{\op}\le2R^2+1$ on the disk $u^2+v^2\le R^2$. (c) Explain what (b) means for \cref{thm:spikes:clipped} and \cref{cor:spikes:gd} on this model.\label{ex:spikes:product}

\item \variation (Clipping with noise is biased.) Let $g=\mu+\zeta$ with $\mu>0$ and $\zeta=\pm s$ with probability $\frac12$ each, $s>\mu$. (a) Compute $\E\spikesclip_\tau(g)$ for every $\tau>0$, and show that it is $0$ when $\tau\le s-\mu$ and lies in $[0,\mu]$ always. (b) Deduce that clipped SGD with a fixed threshold can stall at a point where the gradient is not zero, and compare with the bias term of \eqref{eq:spikes:heavy}.\label{ex:spikes:bias}

\item \variation (A loud jump.) Suppose that the gradient of a coordinate has been $\delta'\ne0$ for so long that $m=\delta'$ and $v=\delta'^2$ (take $\epsilon=0$), and then jumps to $F\delta'$ with $F>1$. (a) Show that the $k$th step after the jump is $\bigl((1-\beta_1^k)+\beta_1^k/F\bigr)/\sqrt{(1-\beta_2^k)+\beta_2^k/F^2}$, which tends to $a_k$ as $F\to\infty$. (b) Show that the first step after the jump exceeds the steady-state value $1$ for every $F>1$ if $(1-\beta_1)^2\ge1-\beta_2$, and only for $F<(\beta_2-\beta_1^2)/((1-\beta_2)-(1-\beta_1)^2)$ otherwise; evaluate for $(0.9,0.95)$. (c) For $(0.9,0.999)$ find the jump factor at which the first step equals $2$.\label{ex:spikes:jump}

\item \variation (Orthogonalized updates keep the logits centred.) (a) Show that every Newton--Schulz iteration $X\mapsto aX+b(XX^\top)X+c(XX^\top)^2X$ of \cref{ch:spectral} preserves $\mathbf1^\top X=0$, so that Muon, applied to an output layer, leaves $\mathbf1^\top W$ unchanged. (b) Find a gradient $G\in\R^{3\times1}$ with $\mathbf1^\top G=0$ and $\mathbf1^\top\sign(G)\ne0$. (c) Show that the mean logit $\frac1K\mathbf1^\top Wh$ can still change during training when $\mathbf1^\top W$ is fixed, and say through which parameters.\label{ex:spikes:muonmean}

\item \simulation (The band of clipped step sizes.) Rerun \texttt{10-spikes-separation.py} with $\gamma\in\{\frac16,\frac13,\frac23\}$ and $\eta=1/(2L_0)$. Check that the largest convergent step size of clipped descent lies in $[L_1\gamma/(L_0(e^{L_1\gamma/2}-1)),2/L_0]$ for every $M$, that the band shrinks to $2/L_0$ as $\gamma\to0$, and that the number of steps down the cliff grows like $1/\gamma$. Explain why the two-cycle $\pm\gamma/2$ exists exactly when the lower edge is exceeded.\label{ex:spikes:simsep}

\item \simulation (Heavy tails.) Run SGD and clipped SGD on $\Loss(\theta)=\frac12\norm\theta_2^2$ in $\R^{10}$ with noise whose coordinates are independent and symmetric with $\Pr[\abs{\zeta_i}>x]=(1+x)^{-3/2}$, so that $p$th moments are finite exactly for $p<\frac32$. (a) Observe that SGD with a constant step makes rare huge excursions. (b) With the schedule of \cref{thm:spikes:heavy} for $p=1.4$, measure the average of $\norm{\nabla\Loss(\theta_t)}_2^2$ for $T=10^3,\dots,10^6$ and fit a slope in $\log T$; compare it with $-2(p-1)/(3p-2)=-0.36$ (in our runs, about $-0.32$). (c) Explain why the bias term of \eqref{eq:spikes:heavy} is much smaller here than its bound, and which term sets the slope.\label{ex:spikes:simheavy}

\item \simulation (Logit drift.) Run \texttt{10-spikes-drift.py}. (a) Check that SGD keeps the mean logit fixed to rounding error, and that Adam's drift vanishes in full batch. (b) Explain the sign of the drift from the distribution of the gradient of a rare class, small and positive in most batches and large and negative in a few, and the coordinatewise normalization of Adam. (c) Find, by bisection on $\alpha$, the smallest z-loss coefficient that stops the drift at batch size $64$, and compare it with $10^{-4}$.\label{ex:spikes:simdrift}

\item \research (A solvable spike.) Find a low-dimensional model in which Adam with fixed hyperparameters produces recurring loss spikes without a switch put in by hand, and compute their period and height as functions of $\eta$ and $\beta_2$. The burst of \cref{prop:spikes:burst} and the stability threshold of Adam in \cref{ch:eos} are natural ingredients; what is open is a model in which the quiet period itself is produced by the dynamics.\label{ex:spikes:spike}

\item \research (A threshold from data.) Estimate $L_0$ and $L_1$ along a training run by regressing a local smoothness estimate on the gradient norm, as \textcite{ZhangHeSraJadbabaie2020} did, and compare $2L_0/(3L_1)$ with the threshold $1.0$. Is a threshold set from such an estimate better than $1.0$, and how should it change with the size of the model? As far as we know, both questions are open.\label{ex:spikes:threshold}
\end{exercises}

\notesandrefs
Gradient clipping entered deep learning with recurrent language models \parencite{Mikolov2012}, and \textcite{PascanuMikolovBengio2013} gave the norm-clipping rule and the picture of cliffs, building on the analysis of exploding and vanishing gradients by \textcite{BengioSimardFrasconi1994}; the textbook of \textcite{GoodfellowBengioCourville2016} gives both. The rule passed into transformer recipes with a threshold of $1.0$ on the global norm \parencite{BrownEtAl2020,TouvronEtAl2023,ChowdheryEtAl2023}. Adam's step, bounded by construction \parencite{KingmaBa2015}, is clipped before it is computed: the clipped gradient enters its moments.

\Textcite{ZhangHeSraJadbabaie2020} introduced $(L_0,L_1)$-smoothness, measured it along trajectories, and proved the first separation between clipped and fixed-step gradient descent; our proofs follow their strategy, with the exact exponential bound of \cref{lem:spikes:local} in place of its linearization, a maximum in place of a sum in \cref{thm:spikes:clipped}, and the cliff and the plateau as the two functions of \cref{thm:spikes:lower}. The extensions are listed in \cref{sec:spikes:misses}. Heavy-tailed gradient noise was measured by \textcite{SimsekliSagunGurbuzbalaban2019}; \textcite{ZhangEtAl2020Attention} proved the rate of \cref{thm:spikes:heavy} with a matching lower bound and proposed a coordinatewise adaptive clipping; \textcite{KunstnerEtAl2023} showed that noise is not what separates SGD from Adam on transformers.

The record of spikes in large runs is mostly in technical reports: PaLM's restarts and skipped batches and its z-loss \parencite{ChowdheryEtAl2023}, OPT's restarts \parencite{ZhangEtAl2022OPT}, the attention-logit divergence of ViT-22B and its cure by QK-norm \parencite{DehghaniEtAl2023}, and Kimi K2's QK-clip \parencite{KimiTeam2025}. \Textcite{WortsmanEtAl2024} reproduced the two logit instabilities in small models and introduced the sensitivity of the final loss to the learning rate as a measure of stability; \textcite{WortsmanEtAl2023} tied spikes to an out-of-date second moment; \textcite{MolybogEtAl2023} proposed a theory of Adam's instability in large runs; \textcite{ZhaiEtAl2023} tied instability to the collapse of attention entropy. We know no reference for the burst formula \eqref{eq:spikes:burst} beyond its first step, which is the quiet-then-loud step of \cref{ch:adaptive}, nor for the drift of logits under Adam in the unigram model of \cref{sec:spikes:reading}; both are elementary.

\paragraph{Translating the literature.} \Textcite{ZhangHeSraJadbabaie2020} write $f$ and $x$ for our $\Loss$ and $\theta$, and clipped gradient descent as $x_{k+1}=x_k-h_c\nabla f(x_k)$ with $h_c=\min\{\eta_c,\gamma\eta_c/\norm{\nabla f(x_k)}\}$: their $\eta_c$ is our $\eta$, their $\gamma$ is a threshold on the gradient norm, our $\gamma/\eta$, and their $\gamma\eta_c$ is our $\gamma$; their $M$ is ours. \Textcite{ZhangEtAl2020Attention} write $\alpha$ for our moment order $p$ and $G$ for our $\Gamma$. PaLM's z-loss is our $\alpha(\log Z)^2$ with $\alpha=10^{-4}$. Kimi K2's report writes $\tau$ for our logit threshold $\bar s$ and $S_{\max}$ for the largest logit of a head, and multiplies the product of the query and key weights by $\min\{1,\tau/S_{\max}\}$, which is our $\rho^2$. Practitioners' \texttt{max\_norm} is our $\tau=\gamma/\eta$.

\printchapterbib
\end{refsection}


\casepart{Noise, Schedules and Equilibria}{Stochastic gradients set a floor under the loss that rises with the step size. The batch size trades noise for steps, the learning-rate schedule removes the floor at the end of training, and weight decay drives normalized layers to an equilibrium whose angular speed is the effective learning rate.}{72}
% Chapter 11 of Optimizers in the Wild: noise and batch size.
\newcommand{\noiseEmin}{E_{\min}}
\newcommand{\noiseTmin}{T_{\min}}
\newcommand{\noiseBcrit}{B_{\mathrm{crit}}}
\newcommand{\noiseBnoise}{B_{\mathrm{noise}}}
\newcommand{\noiseBsimple}{B_{\mathrm{simple}}}
\newcommand{\noiseBeps}{B_{\eps}}
\newcommand{\noisesl}{\mathrm{sl}}

\casechapter{Noise and Batch Size}{The noisy quadratic model and the critical batch size}{Batch-size scaling rules}{ch:noise}
\begin{refsection}[chapters/11-noise.bib]

\begin{chaptersummary}
Up to a critical batch size, doubling the batch and the learning rate together halves the number of steps of a training run; beyond it, larger batches buy almost nothing, and the critical batch size is larger for momentum and for preconditioned optimizers than for plain SGD. The noisy quadratic model, a quadratic loss whose gradients carry noise of covariance $\Sigma/B$, is the smallest idealization with both regimes. We solve its dynamics exactly and prove that the number of steps to a target loss, at the best constant step size, is asymptotic to $\noiseEmin/B$ at small batch sizes, where the step size is proportional to $B$, and equals a constant $\noiseTmin$ at large ones, where the step size sits near the stability limit $2/S$; both constants are explicit, and the critical batch size $\noiseEmin/\noiseTmin$ is of order $\tr\Sigma/(S\eps)$ for a target $\eps$. Preconditioning raises this scale, and when the noise lives in the bulk of the spectrum it leaves the small-batch cost unchanged. The model has a fixed landscape and fixed noise, so it cannot say how the critical batch size moves with the data and the model.
\end{chaptersummary}

\section{In the wild: batch sizes and scaling rules}\label{sec:noise:wild}

A data-parallel training run splits each batch of $B$ examples across devices, computes the per-example gradients in parallel, and averages them. The average is an unbiased estimate of the gradient, with the per-example covariance $\Sigma$ divided by $B$ as its covariance. As long as the batch fits the hardware, a step takes about the same time whatever the batch size, so the number of steps $T$ measures the wall-clock time of a run and the number of examples $E=BT$ measures its cost in data and in computation.\sidenote{For a language model an example is a sequence or a token, and batch sizes are quoted in tokens. Nothing in this chapter depends on the unit, as long as $B$ and $E$ are counted in the same one.} A larger batch can pay off only through fewer steps, and two questions follow. When the batch size changes, how should the learning rate change? And how large can the batch be before it stops buying steps? In the practice sections we say \emph{learning rate}, in the theory \emph{step size}; they are the same number $\eta$.

\begin{practice}[ResNet-50 on ImageNet, \textcite{GoyalEtAl2017}; Adam, \textcite{MalladiEtAl2022}; GPT-3, \textcite{BrownEtAl2020}]
For SGD, the \emph{linear scaling rule}: when the batch size is multiplied by $k$, multiply the learning rate by $k$. Goyal et al.\ trained ResNet-50 with the learning rate $0.1\cdot B/256$ for batch sizes up to $8192$, after a linear warmup over the first five epochs. For Adam, the \emph{square-root rule}: multiply the learning rate by $\sqrt k$, and $1-\beta_1$ and $1-\beta_2$ by $k$. Both rules are used up to a \emph{critical batch size}, beyond which a larger batch no longer reduces the number of steps much, and batch sizes are chosen near it. Language models train on batches of millions of tokens: $3.2$ million for the largest GPT-3 model, reached by increasing the batch size gradually at the start of training.
\end{practice}

The linear rule is discussed by \textcite{Krizhevsky2014}, and \textcite{GoyalEtAl2017} made it into a recipe that trained ResNet-50 on ImageNet in one hour on $256$ GPUs. \Textcite{HofferHubaraSoudry2017} argued instead for a learning rate proportional to $\sqrt B$ for SGD, which keeps the covariance of the update fixed, and \textcite{MalladiEtAl2022} derived the square-root rule for adaptive methods from stochastic differential equations. \Textcite{SmithEtAl2018} turned the rules around: increasing the batch size during training can replace a decay of the learning rate. What happens when the batch size alone is varied, with everything else retuned, was measured systematically only later.

\begin{observation}[Three regimes of data parallelism; \textcite{ShallueEtAl2019}]\label[observation]{obs:noise:regimes}
With the learning rate and the other optimizer hyperparameters retuned for every batch size, the number of steps needed to reach a fixed goal for the out-of-sample error falls in proportion to $1/B$ at small batch sizes (\term{perfect scaling}\index{perfect scaling}), then more slowly (diminishing returns), and finally not at all (maximal data parallelism). The shape is the same across the image and text workloads studied; the batch size at which perfect scaling ends differs widely between workloads and depends on the model, the data set and the optimizer.
\end{observation}

\begin{observation}[The gradient noise scale; \textcite{McCandlishEtAl2018}]\label[observation]{obs:noise:scale}
Across tasks from image classification to reinforcement learning, the critical batch size is predicted, roughly, by the gradient noise scale $\tr\Sigma/\norm{\nabla\Loss}_2^2$ measured during training. The noise scale grows over a run as the loss falls, and it is larger for harder tasks. The steps $T$ and examples $E$ needed to reach a target fit the relation $(T/\noiseTmin-1)(E/\noiseEmin-1)=1$.
\end{observation}

\begin{observation}[Optimizers and the critical batch size; \textcite{ZhangEtAl2019NQM}]\label[observation]{obs:noise:optimizers}
On image classification and language modeling workloads, the optimizers that precondition, Adam and K-FAC \parencite{MartensGrosse2015}, reach much larger critical batch sizes than SGD with momentum.
\end{observation}

\begin{observation}[Language models; \textcite{KaplanEtAl2020}, \textcite{ZhangEtAl2025CBS}, \textcite{MarekEtAl2025}]\label[observation]{obs:noise:lm}
In language-model pretraining the critical batch size, counted in tokens, grows as the loss falls, and it depends on the loss reached rather than directly on the model size (Kaplan et al.); in controlled sweeps it grows with the amount of training data more than with the model size (Zhang et al.). At the other end, very small batches, down to a single sequence, train stably when Adam's $\beta_2$ is adjusted to keep its averaging horizon fixed in tokens, and at such batch sizes plain SGD without momentum comes close to Adam; the gap between the two grows with the batch size (Marek et al.).
\end{observation}

The first three observations are robust in the settings where they were made, but they rest on expensive sweeps over a few workloads each, and the noise-scale prediction of \cref{obs:noise:scale} is accurate only to within a factor of a few; McCandlish et al.'s paper is an unrefereed preprint. Together they pose the question of this chapter. Why does doubling the batch halve the steps up to some batch size and not beyond it, what sets that batch size, and why do momentum and preconditioners move it, while leaving the small-batch end almost where it was? A model that answers needs two ingredients and nothing more: noise, whose effect falls with the batch size, and curvature, which caps the step size whatever the batch size.

\section{The idealization: the noisy quadratic model}\label{sec:noise:model}

\begin{assumption}[The noisy quadratic model]\label[assumption]{ass:noise:nqm}
\begin{enumerate}
\item The loss is $\Loss(\theta)=\frac12\theta^\top H\theta$ on $\R^d$, where $H$ is symmetric positive definite with eigenvalues $h_1\ge h_2\ge\dots\ge h_d>0$, and $S=h_1=\lmax(H)$ is the sharpness. We work in an orthonormal eigenbasis of $H$, so that $H=\diag(h_1,\dots,h_d)$, $\theta^\star=0$ and $\Loss^\star=0$.
\item At step $t$ the optimizer receives the stochastic gradient $g_t=H\theta_t+\zeta_t$. Conditionally on the past, that is, on the $\sigma$-field $\cF_t$ generated by $\zeta_0,\dots,\zeta_{t-1}$, the noise has mean zero and covariance $\Sigma/B$: $\E[\zeta_t\mid\cF_t]=0$ and $\E[\zeta_t\zeta_t^\top\mid\cF_t]=\Sigma/B$, for a fixed positive semidefinite matrix $\Sigma\ne0$ and a batch size $B>0$. We write $c_i=\Sigma_{ii}$ for the diagonal of $\Sigma$ in the eigenbasis of $H$.
\item The initialization $\theta_0$ is deterministic; $\Delta_{0,i}=\frac12h_i\theta_{0,i}^2$ is the initial loss of mode $i$, and $\Delta_0=\Loss(\theta_0)=\sum_i\Delta_{0,i}$.
\end{enumerate}
\end{assumption}

The average of $B$ independent per-example gradients with covariance $\Sigma$ satisfies (ii). In the theory the batch size is a positive real number, and it enters only through the noise level $\Sigma/B$. Nothing else about the noise will matter: it need not be Gaussian, and its covariance need not be diagonal in the eigenbasis of $H$. The standard choice is noise proportional to curvature, $\Sigma=\sigma^2H$, which is what \textcite{ZhangEtAl2019NQM} use in their noisy quadratic model: for a likelihood loss at the minimizer of a well-specified model, the covariance of the per-example gradients is the Fisher information, and it equals the Hessian of the population loss there.\sidenote{The Hessian spectra of trained networks have a few outlier eigenvalues above a bulk \parencite{GhorbaniKrishnanXiao2019,Papyan2020}. \Cref{prop:noise:spiked} uses this shape.} Our running example, which we call the \emph{harmonic model}, has $d=1000$, $h_i=1/i$, $\Sigma=H$, $\theta_0=(1,\dots,1)$, and the target $\eps=0.05\,\Delta_0$. Its initial loss is $\Delta_0=\frac12\sum_{i\le1000}1/i\approx3.74$ and its effective rank is $\tr H/S\approx7.49$.

\begin{dictionary}
the quadratic $\frac12\theta^\top H\theta$, its eigenvalues $h_i$ & the loss near the region the run is in, and its Hessian spectrum: a bulk and a few outliers\\
the sharpness $S=h_1$ and the limit $\eta<2/S$ & the largest stable learning rate; training near it is the edge of stability (\cref{ch:eos})\\
the covariance $\Sigma$, often $\sigma^2H$ & the covariance of per-example gradients, close to the Fisher information and to the Hessian near a minimum\\
the batch size $B$, the noise level $\Sigma/B$ & the batch size, in sequences or tokens\\
the target $\eps$ & a target value of the loss, above its minimum\\
$T(B)$ and $E(B)=B\,T(B)$ & steps and examples to the target, at a tuned constant learning rate\\
the floor $\cR_\infty\approx\eta\tr\Sigma/(4B)$ & the plateau of the loss at a constant learning rate, removed by a decay (\cref{ch:schedules})\\
the noise-limited regime, $\eta\approx\gamma^\star B$ & perfect scaling and the linear scaling rule\\
the curvature-limited regime, $\eta$ near $2/S$ & maximal data parallelism\\
$\noiseBcrit=\noiseEmin/\noiseTmin$ & the critical batch size\\
a preconditioner $P$, $S_P=\lmax(P^{-1}H)$ & Adam's $\sqrt{\hat v}$, Shampoo or K-FAC, frozen and idealized as a function of $H$
\end{dictionary}

\paragraph{What was thrown away.} The landscape is fixed: the curvature and the noise do not change during training, so there is no feature learning and no progressive sharpening. The noise does not depend on the parameters, while in an interpolating model it shrinks with the loss (Exercise~\ref{ex:noise:interp}). Every step sees fresh examples, so there are no epochs and no gap between training and test loss. The step size is constant, and the object of study is the expected loss, not the loss of one run. Each simplification is the right one to make first because it removes something the question does not need. The question is a trade-off between two effects of the step size, and a quadratic with additive noise is the smallest model that has both: curvature, which limits the step size through stability, and noise, whose effect on the loss grows with the step size and falls with the batch size. A model with less, a deterministic quadratic or pure noise, has one regime only.

\paragraph{What one had to notice.} The batch size enters the model only through the noise level $\Sigma/B$, while the step size enters twice: in the speed $\eta h_i$ of every mode, and in a noise floor proportional to $\eta/B$. So the step size meets two constraints. Stability, $\eta<2/S$, does not depend on $B$; the floor, about $\eta\tr\Sigma/(4B)$, must lie below the target, which caps $\eta$ at a multiple of $B$. When the floor binds, the best step size is proportional to $B$ and the number of steps to $1/B$; when stability binds, the number of steps no longer depends on $B$. The critical batch size is where the two constraints meet.

\section{The theory: exact dynamics}\label{sec:noise:exact}

The first theorem solves the model. It is elementary, and everything else in the chapter rests on it.

\begin{theorem}[Exact dynamics of SGD in the noisy quadratic model]\label{thm:noise:nqm}
Let \cref{ass:noise:nqm} hold, let $\theta_{t+1}=\theta_t-\eta g_t$ with a constant step size $\eta>0$, and write $m_{t,i}=\E\theta_{t,i}^2$ and $\cR_t(\eta,B)=\E\Loss(\theta_t)$.
\begin{enumerate}
\item For every $t\ge0$ and every $i$, $m_{t+1,i}=(1-\eta h_i)^2m_{t,i}+\eta^2c_i/B$.
\item If $0<\eta<2/S$, then for every $t\ge0$
\begin{equation}\label{eq:noise:loss}
  \cR_t(\eta,B)=\sum_{i=1}^d\Bigl[F_i+\bigl(\Delta_{0,i}-F_i\bigr)(1-\eta h_i)^{2t}\Bigr],\qquad F_i=\frac{\eta c_i}{2B(2-\eta h_i)},
\end{equation}
and $\cR_t(\eta,B)$ converges geometrically to the \term{noise floor}\index{noise floor}
\begin{equation}\label{eq:noise:floor}
  \cR_\infty(\eta,B)=\frac{\eta}{2B}\sum_{i=1}^d\frac{c_i}{2-\eta h_i},\qquad\text{with}\qquad\frac{\eta\tr\Sigma}{4B}\le\cR_\infty(\eta,B)\le\frac{\eta\tr\Sigma}{2B(2-\eta S)}.
\end{equation}
\item If $\eta S\ge2$ and $c_1>0$, or if $\eta S>2$ and $\theta_{0,1}\ne0$, then $\cR_t(\eta,B)\to\infty$.
\item The expected loss depends on the noise only through the conditional means and variances of its coordinates in the eigenbasis of $H$; neither Gaussian noise nor a diagonal $\Sigma$ is needed.
\end{enumerate}
\end{theorem}

\begin{proof}
The plan is to show that each coordinate is a scalar linear recursion driven by noise that is uncorrelated with its past, and to solve that recursion. In the eigenbasis the update reads $\theta_{t+1,i}=(1-\eta h_i)\theta_{t,i}-\eta\zeta_{t,i}$. The coordinate $\theta_{t,i}$ is a function of $\zeta_0,\dots,\zeta_{t-1}$, so it is $\cF_t$-measurable, and \cref{ass:noise:nqm}(ii) gives $\E[\zeta_{t,i}\mid\cF_t]=0$ and $\E[\zeta_{t,i}^2\mid\cF_t]=c_i/B$. Squaring the update and conditioning on $\cF_t$, the cross term vanishes:
\[
  \E\bigl[\theta_{t+1,i}^2\bigm|\cF_t\bigr]=(1-\eta h_i)^2\theta_{t,i}^2+\eta^2c_i/B .
\]
Taking expectations gives (i), and by induction every $m_{t,i}$ is finite.

For (ii), put $\rho_i=(1-\eta h_i)^2$. Since $0<\eta h_i\le\eta S<2$, we have $0\le\rho_i<1$ and $1-\rho_i=\eta h_i(2-\eta h_i)>0$. The recursion (i) has the fixed point $m_i^\infty=\eta^2c_i/(B(1-\rho_i))=\eta c_i/(Bh_i(2-\eta h_i))$, and subtracting it gives $m_{t,i}-m_i^\infty=\rho_i^t(m_{0,i}-m_i^\infty)$. Multiplying by $h_i/2$, and noting that $\frac12h_im_i^\infty=F_i$ and $\frac12h_im_{0,i}=\Delta_{0,i}$, we obtain the $i$th term of \eqref{eq:noise:loss}; the convergence is geometric with ratio $\max_i\rho_i<1$. The bounds in \eqref{eq:noise:floor} follow from $2-\eta S\le2-\eta h_i\le2$.

For (iii), $\rho_1\ge1$ when $\eta S\ge2$, so (i) gives $m_{t+1,1}\ge m_{t,1}+\eta^2c_1/B$ and hence $m_{t,1}\ge t\eta^2c_1/B$; it also gives $m_{t,1}\ge\rho_1^t\theta_{0,1}^2$, which tends to infinity when $\rho_1>1$ and $\theta_{0,1}\ne0$. In both cases $\cR_t\ge\frac12Sm_{t,1}\to\infty$. Part (iv) holds because the argument used nothing else about $\zeta_t$.
\end{proof}

Formula \eqref{eq:noise:loss} says that each mode relaxes geometrically from its initial loss to its share $F_i$ of the floor, at the rate $(1-\eta h_i)^2\approx e^{-2\eta h_i}$ per step. When $\Delta_{0,i}>F_i$ for every $i$, as at all the step sizes that matter in our examples, the expected loss decreases strictly in $t$. The step size enters twice: in the speed of every mode, and in the floor, which for $\eta S\ll1$ is about $\eta\tr\Sigma/(4B)$ and depends on $\eta$ and $B$ only through $\eta/B$. The batch size enters once, in the floor. \Cref{fig:noise:curves} shows both effects in the harmonic model. At a fixed batch size a larger step size reaches its floor sooner, but the floor is higher; at a fixed step size a larger batch lowers the floor and does not make the early steps faster. To reach a target $\eps$, the floor must lie below $\eps$, which caps $\eta$ at about $4B\eps/\tr\Sigma$, and the iteration must be stable, which caps $\eta$ below $2/S$ whatever $B$ is. The next section makes this sentence precise. The floor \eqref{eq:noise:floor} is also where \cref{ch:schedules} begins: a decaying step size removes it, and the river-valley model adds a slow direction to it.

\begin{figure}[t]
\begin{fullwidth}
\centering
\begin{tikzpicture}
\begin{groupplot}[group style={group size=2 by 1,horizontal sep=17mm},wplot,width=82mm,height=56mm,xmode=log,ymode=log,xmin=1,xmax=4000,ymin=0.02,ymax=5,xlabel={step $t$}]
\nextgroupplot[ylabel={expected loss $\cR_t(\eta,B)$},title={batch size $B=4$},legend style={at={(0.03,0.03)},anchor=south west,font=\scriptsize}]
\addplot[theory3] table[x=t,y=a1] {figures/data/11-noise-curves.dat};
\addplot[theory] table[x=t,y=a2] {figures/data/11-noise-curves.dat};
\addplot[theory2] table[x=t,y=a3] {figures/data/11-noise-curves.dat};
\addplot[simulation] table[x=t,y=a1] {figures/data/11-noise-curves-mc.dat};
\addplot[simulation,color=sea,mark options={fill=sea,draw=white,line width=0.2pt}] table[x=t,y=a2] {figures/data/11-noise-curves-mc.dat};
\addplot[simulation2] table[x=t,y=a3] {figures/data/11-noise-curves-mc.dat};
\addplot[guide] table[x=t,y=fa1] {figures/data/11-noise-curves.dat};
\addplot[guide] table[x=t,y=fa2] {figures/data/11-noise-curves.dat};
\addplot[guide] table[x=t,y=fa3] {figures/data/11-noise-curves.dat};
\addplot[ink,line width=0.6pt,densely dotted] table[x=t,y=eps] {figures/data/11-noise-curves.dat};
\legend{$\eta=0.10$,$\eta=0.26$,$\eta=1.00$}
\nextgroupplot[ylabel={expected loss $\cR_t(\eta,B)$},title={step size $\eta=0.26$},legend style={at={(0.03,0.03)},anchor=south west,font=\scriptsize}]
\addplot[theory3] table[x=t,y=b1] {figures/data/11-noise-curves.dat};
\addplot[theory] table[x=t,y=b2] {figures/data/11-noise-curves.dat};
\addplot[theory2] table[x=t,y=b3] {figures/data/11-noise-curves.dat};
\addplot[simulation] table[x=t,y=b1] {figures/data/11-noise-curves-mc.dat};
\addplot[simulation,color=sea,mark options={fill=sea,draw=white,line width=0.2pt}] table[x=t,y=b2] {figures/data/11-noise-curves-mc.dat};
\addplot[simulation2] table[x=t,y=b3] {figures/data/11-noise-curves-mc.dat};
\addplot[guide] table[x=t,y=fb1] {figures/data/11-noise-curves.dat};
\addplot[guide] table[x=t,y=fb2] {figures/data/11-noise-curves.dat};
\addplot[guide] table[x=t,y=fb3] {figures/data/11-noise-curves.dat};
\addplot[ink,line width=0.6pt,densely dotted] table[x=t,y=eps] {figures/data/11-noise-curves.dat};
\legend{$B=1$,$B=4$,$B=16$}
\end{groupplot}
\end{tikzpicture}
\end{fullwidth}
\caption{The expected loss of SGD in the harmonic model ($d=1000$, $h_i=1/i$, $\Sigma=H$, $\theta_0=(1,\dots,1)$). Curves: the exact formula \eqref{eq:noise:loss}; dashed lines: the floors \eqref{eq:noise:floor}; dotted line: the target $\eps=0.05\,\Delta_0$. Points: the average of $64$ independent runs with Gaussian noise of covariance $\Sigma/B$ (seed fixed in \texttt{11-noise-curves.py}); they agree with the curves to within $4.4\%$, the Monte Carlo error. Left: at a fixed batch size a larger step is faster and has a higher floor; $\eta=0.26$ is close to the step size that reaches the target first. Right: at a fixed step size the batch size changes only the floor.}
\label{fig:noise:curves}
\end{figure}

Preconditioning does not need a new theorem. A preconditioner that is a function of $H$ only changes the coordinates.

\begin{lemma}[Preconditioning is a change of variables]\label[lemma]{lem:noise:change}
Let $P=\diag(1/\pi_1,\dots,1/\pi_d)$ with every $\pi_i>0$, so that $P$ is positive definite and commutes with $H$, and let $\theta_{t+1}=\theta_t-\eta P^{-1}g_t$. Then $\theta'_t=P^{1/2}\theta_t$ is SGD with step size $\eta$ on the noisy quadratic model with Hessian $H'=P^{-1/2}HP^{-1/2}$ and noise covariance $\Sigma'=P^{-1/2}\Sigma P^{-1/2}$, and $\Loss(\theta_t)=\frac12\theta_t'^\top H'\theta'_t$. In the common eigenbasis, $h'_i=\pi_ih_i$, $c'_i=\pi_ic_i$ and $\Delta'_{0,i}=\Delta_{0,i}$. In particular, \cref{thm:noise:nqm} holds for preconditioned SGD with these substitutions: it is stable for $\eta<2/S_P$, where $S_P=\lmax(P^{-1}H)=\max_i\pi_ih_i$ is the preconditioned sharpness, and its noise floor $\frac{\eta}{2B}\sum_i\pi_ic_i/(2-\eta\pi_ih_i)$ lies between $\eta\tr(P^{-1}\Sigma)/(4B)$ and $\eta\tr(P^{-1}\Sigma)/(2B(2-\eta S_P))$.
\end{lemma}

\begin{proof}
Multiplying the update by $P^{1/2}$ and writing $\theta_t=P^{-1/2}\theta'_t$ gives
\[
  \theta'_{t+1}=\theta'_t-\eta P^{-1/2}\bigl(HP^{-1/2}\theta'_t+\zeta_t\bigr)=\theta'_t-\eta\bigl(H'\theta'_t+\zeta'_t\bigr),\qquad\zeta'_t=P^{-1/2}\zeta_t,
\]
and $\zeta'_t$ has conditional mean $0$ and conditional covariance $P^{-1/2}\Sigma P^{-1/2}/B=\Sigma'/B$. The loss is $\frac12\theta_t^\top H\theta_t=\frac12\theta_t'^\top H'\theta'_t$. Since $P^{-1/2}=\diag(\sqrt{\pi_i})$, we get $h'_i=\pi_ih_i$, $c'_i=\pi_ic_i$, $\theta'_{0,i}=\theta_{0,i}/\sqrt{\pi_i}$, hence $\Delta'_{0,i}=\frac12\pi_ih_i\theta_{0,i}^2/\pi_i=\Delta_{0,i}$, and $\tr\Sigma'=\sum_i\pi_ic_i=\tr(P^{-1}\Sigma)$.
\end{proof}

The powers $P=H^p$ with $0\le p\le1$ interpolate between SGD ($p=0$) and Newton's method ($p=1$); they give $h'_i=h_i^{1-p}$ and, for $\Sigma=\sigma^2H$, $c'_i=\sigma^2h_i^{1-p}$. The power $p=\frac12$ is what Adam's denominator becomes when the noise dominates and $\Sigma\propto H$ (the heuristic derivation in \cref{sec:noise:critical}). A preconditioner that does not commute with $H$ also produces a noisy quadratic model after the change of variables, with a Hessian that is no longer diagonal (Exercise~\ref{ex:noise:change}).

Momentum, the second way to modify SGD, keeps the model linear but couples each coordinate to its previous value. \Cref{ch:momentum} computes the stationary variance of the heavy ball on one quadratic direction; summed over the modes it gives the floor of the model.

\begin{proposition}[Momentum in the noisy quadratic model]\label[proposition]{prop:noise:momentum}
Let \cref{ass:noise:nqm} hold, and let $\theta_{t+1}=\theta_t-\eta g_t+\beta(\theta_t-\theta_{t-1})$ with $\eta>0$, $0\le\beta<1$ and $\theta_{-1}=\theta_0$.
\begin{enumerate}
\item The expected loss converges for every $\theta_0$ and every $\Sigma$ if and only if $\eta S<2+2\beta$, and then its limit is
\[
  \cR_\infty^\beta(\eta,B)=\sum_{i=1}^d\frac{\eta c_i(1+\beta)}{2B(1-\beta)\bigl(2(1+\beta)-\eta h_i\bigr)}=\sum_{i=1}^d\frac{\eta_{\mathrm e}c_i}{2B\bigl(2-\frac{1-\beta}{1+\beta}\eta_{\mathrm e}h_i\bigr)},\qquad\eta_{\mathrm e}=\frac{\eta}{1-\beta}.
\]
Each term is at most the corresponding term of the floor \eqref{eq:noise:floor} of SGD with the effective step size $\eta_{\mathrm e}$, when that step size is stable, and the two agree to first order in $\eta_{\mathrm e}$.
\item Fix $\beta$, and let $T_\beta(B)$ be the number of steps to the target $\eps$ at the best $\eta\in(0,(2+2\beta)/S)$. Then $B\,T_\beta(B)\to\noiseEmin$ as $B\to0$, the same limit as for SGD in \cref{thm:noise:critical} below.
\end{enumerate}
\end{proposition}

\begin{proof}
The plan for (i) is to treat each mode as a second-order linear recursion, to test its stability with the Jury criterion, and to solve the two-dimensional Lyapunov equation for its stationary variance. Mode $i$ evolves by $x_{t+1}=a_ix_t-\beta x_{t-1}-\eta\zeta_{t,i}$ with $a_i=1+\beta-\eta h_i$. By the Jury test (\cref{lem:toolkit:recurrence}), both roots of $z^2-a_iz+\beta$ lie in the open unit disk if and only if $\beta<1$ and $\abs{a_i}<1+\beta$, that is, $0<\eta h_i<2+2\beta$. If this holds for every $i$, the mean of every mode tends to zero and the second moments converge. If it fails for $i=1$, then a root $\mu$ of modulus at least $1$ exists, the other root $\beta/\mu$ has modulus less than $1$, the impulse response of mode $1$ does not tend to zero, and the variance of mode $1$, which is $\eta^2c_1/B$ times the sum of the squared impulse response, diverges whenever $c_1>0$.

For the stationary second moments $\gamma_0=\E x_t^2$ and $\gamma_1=\E x_tx_{t-1}$ of a stable mode (\cref{lem:toolkit:lyapunov}), multiply the recursion by $x_t$ and square it, using that $\zeta_{t,i}$ is uncorrelated with $x_t$ and $x_{t-1}$: $\gamma_1=a_i\gamma_0-\beta\gamma_1$ and $\gamma_0=(a_i^2+\beta^2)\gamma_0-2a_i\beta\gamma_1+\eta^2c_i/B$. Eliminating $\gamma_1=a_i\gamma_0/(1+\beta)$,
\[
  \gamma_0\,\frac{(1-\beta)\bigl((1+\beta)^2-a_i^2\bigr)}{1+\beta}=\frac{\eta^2c_i}{B},\qquad(1+\beta)^2-a_i^2=\eta h_i\bigl(2(1+\beta)-\eta h_i\bigr),
\]
and $\frac12h_i\gamma_0$ is the $i$th term of $\cR^\beta_\infty$. Substituting $\eta=(1-\beta)\eta_{\mathrm e}$ gives the second form. Since $(1-\beta)/(1+\beta)\le1$, its denominator is at least $2B(2-\eta_{\mathrm e}h_i)$, and both forms equal $\eta_{\mathrm e}c_i/(4B)$ to first order in $\eta_{\mathrm e}$.
\end{proof}

\begin{proofsketch}[Sketch of the proof of (ii)]
For fixed $\beta<1$ and $\eta\to0$, the roots of $z^2-(1+\beta-\eta h_i)z+\beta$ are $\mu_+=1-\eta h_i/(1-\beta)+O(\eta^2)$ and $\mu_-=\beta+O(\eta)$. With $\theta_{-1}=\theta_0$ the mean of mode $i$ is $\theta_{0,i}\mu_+^t$ up to a relative error $O(\eta)$, and its variance is $\eta^2c_i/B$ times the sum of the squared impulse response, which is $\frac{1}{2\eta h_i(1-\beta)}(1-\mu_+^{2t})$ up to the same relative error. With $\eta=(1-\beta)\gamma B$ and $t=E/B$, the expected loss therefore converges, uniformly for $(\gamma,E)$ in compact sets, to the limit $D(\gamma E)+\frac\gamma4Q(\gamma E)$ that SGD has with step size $\gamma B$ (\cref{def:noise:steps}), and the upper- and lower-bound arguments of \cref{thm:noise:critical}(iii) go through. What is omitted is the bookkeeping of the $O(\eta)$ errors, uniformly in $t$; it parallels the proof of \cref{thm:noise:critical}.
\end{proofsketch}

So momentum with a fixed $\beta$ needs, in the limit of small batches, exactly as many examples as SGD: in the harmonic model with $\beta=0.9$ it needs $9989$ steps at $B=1$, against $10063$ for SGD and the common limit $\noiseEmin=9987$. Its benefit is all at large batch sizes, where the region $\eta S<2+2\beta$ allows effective step sizes up to $(2+2\beta)/((1-\beta)S)=38/S$, and the iteration accelerates (\cref{ch:momentum}).

The last result of this section concerns one step. It is the model from which \textcite{McCandlishEtAl2018} derived the gradient noise scale and their relation between steps and examples; in the noisy quadratic model it is exact.

\begin{proposition}[The one-step model is exact]\label[proposition]{prop:noise:onestep}
Let \cref{ass:noise:nqm} hold, fix $\theta$ with $G=\nabla\Loss(\theta)=H\theta\ne0$, and let $g=G+\zeta$ with $\E\zeta=0$ and $\E\zeta\zeta^\top=\Sigma/B$. Then for every $\eta$,
\begin{equation}\label{eq:noise:onestep}
  \E\Loss(\theta-\eta g)=\Loss(\theta)-\eta\norm{G}_2^2+\frac{\eta^2}2\Bigl(G^\top HG+\frac{\tr(H\Sigma)}B\Bigr).
\end{equation}
The right side is smallest at $\eta=\eta_{\max}/(1+\noiseBnoise/B)$, where it equals $\Loss(\theta)-\Delta_{\max}/(1+\noiseBnoise/B)$, with
\[
  \eta_{\max}=\frac{\norm G_2^2}{G^\top HG},\qquad\Delta_{\max}=\frac{\norm G_2^4}{2G^\top HG},\qquad\noiseBnoise=\frac{\tr(H\Sigma)}{G^\top HG}.
\]
\begin{enumerate}
\item The \term{gradient noise scale}\index{gradient noise scale} $\noiseBsimple=\tr\Sigma/\norm G_2^2$ satisfies $\noiseBsimple\ge\tr\Sigma/(2S\Loss(\theta))$, with equality if and only if $\theta$ lies in the top eigenspace of $H$. If $H=hI$, then $\noiseBnoise=\noiseBsimple=\tr\Sigma/(2h\Loss(\theta))$.
\item $1/S\le\eta_{\max}\le1/h_d$, so the one-step optimal step size can exceed the stability limit $2/S$.
\item Let $\theta$ be distributed according to the stationary second moments of SGD with step size $\eta<2/S$ and batch size $B$ (\cref{thm:noise:nqm}(ii)). Then the noise scales measured with $\bar g_1=\E\norm G_2^2$ and $\bar g_2=\E G^\top HG$ satisfy
\[
  (2-\eta S)\frac{B\tr\Sigma}{\eta\tr(H\Sigma)}\le\frac{\tr\Sigma}{\bar g_1}\le\frac{2B\tr\Sigma}{\eta\tr(H\Sigma)},\qquad(2-\eta S)\frac{B\tr(H\Sigma)}{\eta\tr(H^2\Sigma)}\le\frac{\tr(H\Sigma)}{\bar g_2}\le\frac{2B\tr(H\Sigma)}{\eta\tr(H^2\Sigma)}.
\]
\end{enumerate}
\end{proposition}

\begin{proof}
The loss is quadratic, so $\Loss(\theta-\eta g)=\Loss(\theta)-\eta G^\top g+\frac12\eta^2g^\top Hg$ exactly. Taking expectations, $\E G^\top g=\norm G_2^2$ and $\E g^\top Hg=G^\top HG+\E\zeta^\top H\zeta=G^\top HG+\tr(H\Sigma)/B$, the cross term having mean zero; this is \eqref{eq:noise:onestep}. Its right side is a quadratic in $\eta$ with positive leading coefficient, minimized at $\norm G_2^2/(G^\top HG+\tr(H\Sigma)/B)$ with minimum $\Loss(\theta)-\norm G_2^4/(2G^\top HG+2\tr(H\Sigma)/B)$, which are the stated expressions.

(i) Since $H^2\preceq SH$, $\norm G_2^2=\theta^\top H^2\theta\le S\theta^\top H\theta=2S\Loss(\theta)$, with equality if and only if $(SH-H^2)^{1/2}\theta=0$, that is, $H(S-H)\theta=0$; since $H$ is invertible, this means $H\theta=S\theta$. If $H=hI$, then $G=h\theta$, $\norm G_2^2=2h\Loss(\theta)$, $G^\top HG=h\norm G_2^2$ and $\tr(H\Sigma)=h\tr\Sigma$. (ii) The number $1/\eta_{\max}$ is the Rayleigh quotient of $H$ at $G$, which lies in $[h_d,S]$. (iii) At stationarity $\E\theta_i^2=m_i^\infty=\eta c_i/(Bh_i(2-\eta h_i))$, so $\bar g_1=\sum_ih_i^2m_i^\infty=\frac\eta B\sum_ih_ic_i/(2-\eta h_i)$ and $\bar g_2=\frac\eta B\sum_ih_i^2c_i/(2-\eta h_i)$; the bounds follow from $2-\eta S\le2-\eta h_i\le2$, with $\tr(H\Sigma)=\sum_ih_ic_i$ and $\tr(H^2\Sigma)=\sum_ih_i^2c_i$.
\end{proof}

Part (iii) has a consequence that matters for practice. A noise scale measured on the iterates of a run that has reached its floor tells us the run's own ratio $B/\eta$, up to factors that depend on the spectrum, and not a property of the problem. We return to this in \cref{sec:noise:reading}.

\section{The theory: the critical batch size}\label{sec:noise:critical}

We now count steps. The target is a loss value $\eps$, and the step size is chosen, for each batch size, to reach it in the fewest steps.

\begin{definition}[Steps and examples to a target]\label[definition]{def:noise:steps}
Let \cref{ass:noise:nqm} hold and fix a target $0<\eps<\Delta_0$. For $B>0$, the \term{steps to target} $T(B)$ is the smallest $t\in\N$ for which some step size $\eta\in(0,2/S)$ gives $\cR_t(\eta,B)<\eps$, and $E(B)=B\,T(B)$ is the number of examples. Without noise, let $D_t(\eta)=\sum_i\Delta_{0,i}(1-\eta h_i)^{2t}$ be the loss of gradient descent, and let $\noiseTmin$ be the smallest $t$ for which some $\eta\in(0,2/S)$ gives $D_t(\eta)<\eps$. In continuous time let
\[
  D(\tau)=\sum_{i=1}^d\Delta_{0,i}e^{-2h_i\tau},\qquad Q(\tau)=\sum_{i=1}^dc_i\bigl(1-e^{-2h_i\tau}\bigr),
\]
let $\tau_\eps$ be the solution of $D(\tau_\eps)=\eps$, and let
\begin{equation}\label{eq:noise:Emin}
  \noiseEmin=\min_{\tau>\tau_\eps}\frac{\tau\,Q(\tau)}{4\bigl(\eps-D(\tau)\bigr)},
\end{equation}
which is attained at some $\tau^\star$; put $\gamma^\star=4(\eps-D(\tau^\star))/Q(\tau^\star)$, so that $\noiseEmin=\tau^\star/\gamma^\star$. The \term{critical batch size}\index{critical batch size} is $\noiseBcrit=\noiseEmin/\noiseTmin$.
\end{definition}

The function $D$ is the loss of gradient flow at time $\tau$, and $Q(\tau)$ measures the noise that has accumulated by then. Since $D$ is continuous and strictly decreasing from $\Delta_0$ to $0$, the time $\tau_\eps$ is well defined. The objective in \eqref{eq:noise:Emin} tends to infinity as $\tau\downarrow\tau_\eps$, because its numerator tends to $\tau_\eps Q(\tau_\eps)>0$ while its denominator tends to $0$, and as $\tau\to\infty$, where it is at least $\tau Q(\tau_\eps)/(4\eps)$; so the minimum exists. The formula has a reading as a budget. Spend the flow time $\tau$; the deterministic part of the loss is then $D(\tau)$, which leaves $\eps-D(\tau)$ for the noise; the largest step size per example whose accumulated noise $\frac\gamma4Q(\tau)$ fits in that part is $4(\eps-D(\tau))/Q(\tau)$, and the examples needed are $\tau$ divided by it. In the harmonic model, $\tau_\eps=350$, $\tau^\star=677$, $\gamma^\star=0.0678$ and $\noiseEmin=9987$; at the optimum the noise uses two thirds of the budget. Restricting the step sizes to $(0,2/S)$ is a modeling choice: a larger step makes the stiffest mode grow, which a run can tolerate for a few steps when that mode starts with little loss, but not for long.

\begin{theorem}[Two regimes and the critical batch size]\label{thm:noise:critical}
Let \cref{ass:noise:nqm} hold, fix a target $0<\eps<\Delta_0$, and use the notation of \cref{def:noise:steps}.
\begin{enumerate}
\item \emph{(Monotonicity.)} $T(B)$ is finite and nonincreasing in $B$, and $T(B)\ge\noiseTmin$ for every $B$.
\item \emph{(The curvature-limited regime.)} There is $B_1<\infty$ such that $T(B)=\noiseTmin$ for every $B\ge B_1$.
\item \emph{(The noise-limited regime.)} $B\,T(B)\to\noiseEmin$ as $B\to0$. Quantitatively,
\[
  B\,T(B)\le\noiseEmin(1+\gamma^\star BS)^2\Bigl(1+\frac{\gamma^\star B}{\tau^\star}\Bigr)+B\quad\text{if }\gamma^\star BS\le\tfrac12,\qquad B\,T(B)\ge\frac{\noiseEmin}{\bigl(1+S\sqrt{2B\eps/\tr(H\Sigma)}\bigr)^2}\quad\text{if }8BS^2\eps\le\tr(H\Sigma).
\]
\item \emph{(Linear scaling.)} For each $B$ let $\eta_B\in(0,2/S)$ satisfy $\cR_{T(B)}(\eta_B,B)<\eps$. Every limit point of $\eta_B/B$ as $B\to0$ equals $4(\eps-D(\tau))/Q(\tau)$ for a minimizer $\tau$ of \eqref{eq:noise:Emin}; if the minimizer is unique, then $\eta_B/B\to\gamma^\star$.
\end{enumerate}
\end{theorem}

\begin{proof}
The plan is to split the expected loss into a deterministic part and a noise part proportional to $1/B$. At large batch sizes the noise part disappears; at small ones the step size is forced to zero, so that the discrete recursion is close to its continuous-time version, which is the minimization \eqref{eq:noise:Emin}. For $0<\eta<2/S$, \cref{thm:noise:nqm}(ii) gives
\[
  \cR_t(\eta,B)=D_t(\eta)+\frac{M_t(\eta)}B,\qquad M_t(\eta)=\sum_{i=1}^dBF_i\bigl(1-(1-\eta h_i)^{2t}\bigr)\ge0,
\]
and $BF_i=\eta c_i/(2(2-\eta h_i))$ does not depend on $B$. We use three elementary facts. First, $(1-x)^2\le e^{-2x}$ for $0\le x\le1$. Second, $(1-x)^2\ge e^{-2x(1+x)}$ for $0\le x\le\frac12$, because $\log(1-x)\ge-x-x^2$ there. Third, $Q$ is nondecreasing and concave with $Q(0)=0$, so that $Q(a\tau)\le aQ(\tau)$ and $Q(\tau/a)\ge Q(\tau)/a$ for $a\ge1$ and $\tau\ge0$.

(i) As $\eta\to0$ the floor $\cR_\infty(\eta,B)\le\eta\tr\Sigma/(2B(2-\eta S))$ tends to $0$; so some $\eta$ has a floor below $\eps$, and since $\cR_t$ converges to the floor, $T(B)$ is finite. Since $M_t\ge0$, $\cR_t(\eta,B)$ is nonincreasing in $B$, so a pair $(\eta,t)$ that reaches the target at batch size $B$ reaches it at every larger batch size, and $T$ is nonincreasing. Since $\cR_t(\eta,B)\ge D_t(\eta)$, every such pair has $D_t(\eta)<\eps$, so $T(B)\ge\noiseTmin$. The number $\noiseTmin$ is finite: by the first fact $D_t(1/S)\le D(t/S)$, which is below $\eps$ once $t>S\tau_\eps$.

(ii) Choose $\eta_\infty\in(0,2/S)$ with $D_{\noiseTmin}(\eta_\infty)<\eps$, and let $B_1=1+M_{\noiseTmin}(\eta_\infty)/(\eps-D_{\noiseTmin}(\eta_\infty))$. For $B\ge B_1$ we get $\cR_{\noiseTmin}(\eta_\infty,B)<\eps$, so $T(B)\le\noiseTmin$, and (i) gives equality.

(iii) \emph{Upper bound.} Let $\gamma^\star BS\le\frac12$, let $0<\gamma\le\gamma^\star$, put $\eta=\gamma B$, so that $\eta S\le\frac12$, and let $t=\lceil\tau^\star/\eta\rceil$, so that $\tau^\star\le\eta t<\tau^\star+\eta$. By the first fact, and since $D$ is decreasing, $D_t(\eta)\le D(\eta t)\le D(\tau^\star)$. For the noise part we bound both factors of each term. Since $(1+x)(2-x)\ge2$ for $0\le x\le1$, we have $BF_i\le\frac14\eta c_i(1+\eta h_i)$; and by the second fact and the concavity of $u\mapsto1-e^{-u}$,
\[
  1-(1-\eta h_i)^{2t}\le1-e^{-2\eta h_it(1+\eta h_i)}\le(1+\eta h_i)\bigl(1-e^{-2h_i\eta t}\bigr).
\]
Summing over $i$, and then using the third fact in the form $Q(\eta t)\le Q(\tau^\star+\eta)\le(1+\eta/\tau^\star)Q(\tau^\star)$,
\[
  \frac{M_t(\eta)}B\le(1+\eta S)^2\,\frac\gamma4\,Q(\eta t)\le(1+\gamma^\star BS)^2\Bigl(1+\frac{\gamma^\star B}{\tau^\star}\Bigr)\frac\gamma4\,Q(\tau^\star).
\]
Since $\eps-D(\tau^\star)=\frac14\gamma^\star Q(\tau^\star)$, the target is reached at step $t$ as soon as $\gamma(1+\gamma^\star BS)^2(1+\gamma^\star B/\tau^\star)<\gamma^\star$. Hence $T(B)\le\lceil\tau^\star/(\gamma B)\rceil\le\tau^\star/(\gamma B)+1$ for every such $\gamma$, and letting $\gamma$ increase to $\gamma^\star(1+\gamma^\star BS)^{-2}(1+\gamma^\star B/\tau^\star)^{-1}$ gives the upper bound, because $\tau^\star/\gamma^\star=\noiseEmin$.

\emph{Lower bound.} Let $t=T(B)$ and $\eta\in(0,2/S)$ with $\cR_t(\eta,B)<\eps$, and put $\gamma=\eta/B$ and $\tau=\eta t$. Since $\cR_0=\Delta_0>\eps$, $t\ge1$, and \cref{thm:noise:nqm}(i) gives $m_{t,i}\ge\eta^2c_i/B$; therefore $\eps>\cR_t\ge\eta^2\tr(H\Sigma)/(2B)$, where $\tr(H\Sigma)=\sum_ih_ic_i>0$ because $\Sigma\ne0$. If $8BS^2\eps\le\tr(H\Sigma)$, this gives $\eta S\le S\sqrt{2B\eps/\tr(H\Sigma)}\le\frac12$. Then the second fact gives $(1-\eta h_i)^{2t}\ge e^{-2h_i\tau(1+\eta S)}$, the first fact gives $1-(1-\eta h_i)^{2t}\ge1-e^{-2h_i\tau}$, and $BF_i\ge\frac14\eta c_i$; so, with $\tau'=\tau(1+\eta S)$,
\[
  \eps>\cR_t(\eta,B)\ge D(\tau')+\frac\gamma4\,Q(\tau).
\]
Hence $D(\tau')<\eps$, that is $\tau'>\tau_\eps$, and $\gamma<4(\eps-D(\tau'))/Q(\tau)$. Using $Q(\tau)=Q(\tau'/(1+\eta S))\ge Q(\tau')/(1+\eta S)$,
\begin{equation}\label{eq:noise:lower}
  B\,T(B)=\frac\tau\gamma>\frac{\tau\,Q(\tau)}{4\bigl(\eps-D(\tau')\bigr)}\ge\frac{\tau'\,Q(\tau')}{4(1+\eta S)^2\bigl(\eps-D(\tau')\bigr)}\ge\frac{\noiseEmin}{(1+\eta S)^2},
\end{equation}
which is the lower bound. The two bounds give $B\,T(B)\to\noiseEmin$.

(iv) Let $B_k\to0$, and write $\eta_k=\eta_{B_k}$, $t_k=T(B_k)$, $\gamma_k=\eta_k/B_k$, $\tau_k=\eta_kt_k$ and $\tau'_k=\tau_k(1+\eta_kS)$. For large $k$ the lower-bound argument applies: $\eta_k\to0$, $\tau'_k>\tau_\eps$, $\gamma_k<4(\eps-D(\tau'_k))/Q(\tau_k)$, and $B_kt_k=\tau_k/\gamma_k\to\noiseEmin$. Since $\tau_k\ge\tau_\eps/(1+\eta_kS)\ge\frac23\tau_\eps$ and $B_kt_k>\tau_kQ(\tau_k)/(4\eps)\ge\tau_kQ(\frac23\tau_\eps)/(4\eps)$, the sequence $\tau_k$ is bounded, and so are $\gamma_k=\tau_k/(B_kt_k)$ and $1/\gamma_k$. Along a subsequence $\gamma_k\to\bar\gamma>0$ and $\tau_k\to\bar\tau=\bar\gamma\noiseEmin\ge\tau_\eps$, and in the limit $\bar\gamma\le4(\eps-D(\bar\tau))/Q(\bar\tau)$. Because $\bar\gamma>0$, this forces $D(\bar\tau)<\eps$, so $\bar\tau>\tau_\eps$ and
\[
  \noiseEmin=\frac{\bar\tau}{\bar\gamma}\ge\frac{\bar\tau\,Q(\bar\tau)}{4\bigl(\eps-D(\bar\tau)\bigr)}\ge\noiseEmin .
\]
Equality holds throughout, so $\bar\tau$ is a minimizer and $\bar\gamma=4(\eps-D(\bar\tau))/Q(\bar\tau)$. If the minimizer is unique, every subsequence has the limit $\gamma^\star$.
\end{proof}

Parts (iii) and (iv) are the \term{noise-limited regime}\index{noise-limited regime}: the number of examples is asymptotically independent of the batch size, and the best step size is asymptotically $\gamma^\star B$, the \term{linear scaling rule}\index{linear scaling rule}. The quantitative bounds say that the regime lasts while $\gamma^\star BS\ll1$, that is, while the step size of the linear rule is small compared with $1/S$. Part (ii) is the \term{curvature-limited regime}\index{curvature-limited regime}: the noise has become irrelevant, and the step size is limited by stability alone. The critical batch size is the ratio of the two asymptotes, $T(B)\approx\noiseEmin/B$ and $T(B)=\noiseTmin$. Its size is set by three numbers.

\begin{corollary}[The scale of the critical batch size]\label[corollary]{cor:noise:bcrit}
In the setting of \cref{thm:noise:critical}, for $0<\delta<\Delta_0$ let $\tau_\delta$ solve $D(\tau_\delta)=\delta$; let $D^\noisesl$ be the sum defining $D$ restricted to the modes with $h_i\le S/4$, and let $\tau^\noisesl_\eps=\inf\{\tau\ge0:D^\noisesl(\tau)<\eps\}$. Then
\[
  \frac{\tau_\eps\,Q(\tau_\eps)}{4\eps}\le\noiseEmin\le\frac{\tau_{\eps/2}\,Q(\tau_{\eps/2})}{2\eps},\qquad\frac{S\tau^\noisesl_\eps}4\le\noiseTmin\le S\tau_\eps+1 .
\]
Consequently, with $\noiseBeps=\tr\Sigma/(S\eps)$,
\[
  \frac{Q(\tau_\eps)}{4\tr\Sigma}\cdot\frac{S\tau_\eps}{S\tau_\eps+1}\le\frac{\noiseBcrit}{\noiseBeps}\le\frac{2\tau_{\eps/2}}{\tau^\noisesl_\eps},
\]
the right inequality when $\tau^\noisesl_\eps>0$.
\end{corollary}

\begin{proof}
For $\tau>\tau_\eps$ we have $Q(\tau)\ge Q(\tau_\eps)$ and $0<\eps-D(\tau)<\eps$, so the objective of \eqref{eq:noise:Emin} is at least $\tau_\eps Q(\tau_\eps)/(4\eps)$; and at $\tau=\tau_{\eps/2}>\tau_\eps$ its denominator is $2\eps$. For $\noiseTmin$, the step size $1/S$ gives $D_t(1/S)\le D(t/S)$ by the first fact in the proof of \cref{thm:noise:critical}, which is below $\eps$ for the smallest integer $t>S\tau_\eps$, and that integer is at most $S\tau_\eps+1$. Conversely, let $\eta\in(0,2/S)$ and $t$ satisfy $D_t(\eta)<\eps$. For the modes with $h_i\le S/4$ we have $x=\eta h_i<\frac12$, and $\log(1-x)\ge-2x$ on $[0,\frac12]$ gives $(1-\eta h_i)^{2t}\ge e^{-4\eta h_it}\ge e^{-2h_i(4t/S)}$; therefore $\eps>D_t(\eta)\ge D^\noisesl(4t/S)$, and $4t/S\ge\tau^\noisesl_\eps$. The bounds on $\noiseBcrit$ follow by division, with $Q\le\tr\Sigma$ in the upper one.
\end{proof}

Up to the shape factors on the outside, then, the critical batch size is the total noise divided by the product of the sharpness and the target:
\[
  \noiseBcrit\asymp\frac{\tr\Sigma}{S\,\eps}.
\]
It grows like $1/\eps$ as the target falls, it is proportional to the total noise, and it is inversely proportional to the sharpness, which caps the step size of the curvature-limited regime. For $\Sigma=\sigma^2H$ it is $\sigma^2/\eps$ times the effective rank $\tr H/S$. The shape factors depend on the spectrum and on the initialization through the profile of $D$ near the target: in the harmonic model $\noiseBeps=40$, the corollary places $\noiseBcrit$ between $9.5$ and $125$, and the exact value is $\noiseBcrit=9987/177=56$.

\begin{theorem}[Preconditioning raises the critical batch size]\label{thm:noise:precond}
Let \cref{ass:noise:nqm} hold, let $P=\diag(1/\pi_1,\dots,1/\pi_d)$ with every $\pi_i>0$, and define $T^P(B)$, $\noiseEmin^P$, $\noiseTmin^P$ and $\noiseBcrit^P=\noiseEmin^P/\noiseTmin^P$ for $\theta_{t+1}=\theta_t-\eta P^{-1}g_t$ as in \cref{def:noise:steps}, with step sizes $\eta\in(0,2/S_P)$.
\begin{enumerate}
\item \Cref{thm:noise:critical} and \cref{cor:noise:bcrit} hold for preconditioned SGD after the substitutions $h_i\to\pi_ih_i$, $c_i\to\pi_ic_i$, $S\to S_P=\max_i\pi_ih_i$ and $\tr\Sigma\to\tr(P^{-1}\Sigma)$. In particular, $D$ becomes $D_P(\tau)=\sum_i\Delta_{0,i}e^{-2\pi_ih_i\tau}$ and $Q$ becomes $Q_P(\tau)=\sum_i\pi_ic_i(1-e^{-2\pi_ih_i\tau})$, the linear scaling rule holds at small batch sizes, and $\noiseBcrit^P$ is within the shape factors of \cref{cor:noise:bcrit} of $\tr(P^{-1}\Sigma)/(S_P\eps)$.
\item The advantage of preconditioning, $T(B)/T^P(B)$, tends to $\noiseEmin/\noiseEmin^P$ as $B\to0$, and it equals $\noiseTmin/\noiseTmin^P$ for every large $B$.
\item For $P=H^p$ with $0\le p\le1$ and $\Sigma=\sigma^2H$,
\[
  \frac{\tr(P^{-1}\Sigma)}{S_P\,\eps}=\frac{\sigma^2}{\eps}\sum_{i=1}^d\Bigl(\frac{h_i}{S}\Bigr)^{1-p},
\]
which is nondecreasing in $p$; it is $\sigma^2/\eps$ times the effective rank $\tr H/S$ at $p=0$ and $\sigma^2d/\eps$ at $p=1$.
\end{enumerate}
\end{theorem}

\begin{proof}
By \cref{lem:noise:change}, preconditioned SGD in the coordinates $P^{1/2}\theta$ is SGD on a noisy quadratic model with $h'_i=\pi_ih_i$, $c'_i=\pi_ic_i$, the same initial losses $\Delta_{0,i}$ and the same loss values, so the same target; its stable step sizes are those in $(0,2/S_P)$. \Cref{thm:noise:critical} and \cref{cor:noise:bcrit} apply to this model, which is (i). For (ii), $T(B)/T^P(B)=(B\,T(B))/(B\,T^P(B))$ tends to $\noiseEmin/\noiseEmin^P$ by \cref{thm:noise:critical}(iii), applied to both models, and both step counts equal their limits for every $B$ beyond the two thresholds of \cref{thm:noise:critical}(ii). For (iii), $\pi_i=h_i^{-p}$ gives $\pi_ih_i=h_i^{1-p}$, whose maximum is $S^{1-p}$, and $\pi_ic_i=\sigma^2h_i^{1-p}$; each $(h_i/S)^{1-p}$ with $0<h_i/S\le1$ is nondecreasing in $p$.
\end{proof}

Preconditioning by $H^{-p}$ thus divides the stiffest curvature by more than the others, which raises the stable step size of the curvature-limited regime, and it raises the total noise in the preconditioned metric relative to that curvature; both push the critical batch size up. What it does to the small-batch cost $\noiseEmin^P$ depends on the problem. In the harmonic model, $p=\frac12$ gives $\noiseEmin^P=3390$ against $9987$ for SGD, an advantage of $2.95$ that persists as $B\to0$, while at large batch sizes it needs $10$ steps instead of $177$. A preconditioner helps at small batch sizes only by reallocating time among the modes: by \cref{lem:noise:change}, mode $i$ of preconditioned SGD is SGD on that mode with step size $\eta\pi_i$, and with $\Sigma=\sigma^2H$ the noise a mode accumulates per unit of its own relaxation time is the same in every mode (Exercise~\ref{ex:noise:isotropic} treats isotropic noise). SGD gives the stiff modes more relaxation time than they need, and pays for it in noise; when there is little noise in the stiff modes there is nothing to gain. This is the case of a spiked spectrum.

\begin{proposition}[A spiked spectrum: no gain at small batch sizes]\label[proposition]{prop:noise:spiked}
Let $H$ have $k$ eigenvalues $\kappa h$, the outliers, and $m$ eigenvalues $h$, the bulk, with $\kappa\ge4$. Let $\Sigma=\sigma^2H$ and $\theta_0=\bar\theta\mathbf1$, so that the bulk starts with the loss $L_{\mathrm b}=\frac12mh\bar\theta^2$ and the outliers with $\rho L_{\mathrm b}$, where $\rho=k\kappa/m$; and let the target be $\eps=\phi L_{\mathrm b}$ with $0<\phi<1$. Put $c(x)=-x\log(1-x)$, $\ell(x)=-\log(1-x)$,
\[
  V(\phi)=\inf_{1-\phi<x<1}\frac{c(x)}{x-(1-\phi)},\qquad W_\rho(\phi)=\inf_{1-\phi/(1+\rho)<x<1}\frac{c(x)+\rho\,\ell(x)}{\phi-(1+\rho)(1-x)},
\]
and $R(\rho,\phi)=W_\rho(\phi)/V(\phi)$, which is at least $1$ and tends to $1$ as $\rho\to0$.
\begin{enumerate}
\item \emph{(Small batch sizes.)} Every diagonal preconditioner $P$ has $\noiseEmin^P\ge\frac{\sigma^2}{4h\bar\theta^2}V(\phi)$, and every power $P=H^p$ with $0\le p\le1$, SGD included, has $\noiseEmin^{H^p}\le\frac{\sigma^2}{4h\bar\theta^2}W_\rho(\phi)$. Hence $\noiseEmin/\noiseEmin^P\le R(\rho,\phi)$ for every diagonal $P$, and $\noiseEmin/\noiseEmin^{H^p}\ge1/R(\rho,\phi)$.
\item \emph{(Large batch sizes.)} For $0\le p\le1$,
\[
  \frac{\noiseTmin}{\noiseTmin^{H^p}}\ge\frac{\kappa\log(1/\phi)}{8+4\kappa^{1-p}\log((1+\rho)/\phi)},
\]
which for large $\kappa$ is about $\frac14\kappa^p\log(1/\phi)/\log((1+\rho)/\phi)$.
\end{enumerate}
\end{proposition}

\begin{proof}
The plan is to write the objective of \eqref{eq:noise:Emin}, for the transformed model of \cref{thm:noise:precond}(i), in terms of the fraction of its initial loss that each mode has lost, and to compare the bulk with the outliers. For a diagonal $P$ and a time $\tau$, put $u_i=2\pi_ih_i\tau$ and $x_i=1-e^{-u_i}$. Because $c_i=\sigma^2h_i$,
\[
  \tau\pi_ic_i\bigl(1-e^{-u_i}\bigr)=\frac{\sigma^2}2\,u_i\bigl(1-e^{-u_i}\bigr)=\frac{\sigma^2}2\,c(x_i),\qquad\Delta_{0,i}e^{-u_i}=\tfrac12h_i\bar\theta^2(1-x_i),
\]
so $\tau Q_P(\tau)=\frac{\sigma^2}2\sum_ic(x_i)$ and $D_P(\tau)=\frac12\bar\theta^2\sum_ih_i(1-x_i)$.

(i) \emph{Lower bound.} Dropping the outliers from both sums decreases $\tau Q_P$ and $D_P$, and so decreases the objective. The function $c$ is convex on $[0,1)$, since $c''(x)=(1-x)^{-1}+(1-x)^{-2}>0$, so by Jensen's inequality the bulk contributes at least $m\,c(\bar x)$ to $\sum_ic(x_i)$, where $\bar x$ is the mean of the bulk $x_i$, while it contributes exactly $mh(1-\bar x)$ to $\sum_ih_i(1-x_i)$. If $D_P(\tau)<\eps$, then $L_{\mathrm b}(1-\bar x)<\phi L_{\mathrm b}$ and
\[
  \frac{\tau Q_P(\tau)}{4(\eps-D_P(\tau))}\ge\frac{\frac{\sigma^2}2m\,c(\bar x)}{4L_{\mathrm b}\bigl(\bar x-(1-\phi)\bigr)}=\frac{\sigma^2}{4h\bar\theta^2}\cdot\frac{c(\bar x)}{\bar x-(1-\phi)}\ge\frac{\sigma^2}{4h\bar\theta^2}V(\phi).
\]
\emph{Upper bound.} For $P=H^p$, $u_i=2h_i^{1-p}\tau$, so the bulk modes share one value $u=2h^{1-p}\tau$, with $x=1-e^{-u}$ and $u=\ell(x)$, and the outliers have $u_{\mathrm o}=\kappa^{1-p}u\ge u$. Then $c(x_{\mathrm o})\le u_{\mathrm o}\le\kappa u=\kappa\,\ell(x)$ and $1-x_{\mathrm o}\le1-x$, so, using $k\kappa=\rho m$,
\[
  \tau Q_P(\tau)\le\frac{\sigma^2}2\,m\bigl(c(x)+\rho\,\ell(x)\bigr),\qquad D_P(\tau)\le(1+\rho)L_{\mathrm b}(1-x).
\]
As $\tau$ runs over $(0,\infty)$, $x$ runs over $(0,1)$, and therefore
\[
  \noiseEmin^{H^p}\le\inf_x\frac{\frac{\sigma^2}2m\bigl(c(x)+\rho\,\ell(x)\bigr)}{4L_{\mathrm b}\bigl(\phi-(1+\rho)(1-x)\bigr)}=\frac{\sigma^2}{4h\bar\theta^2}W_\rho(\phi).
\]
SGD is the case $p=0$, so both ratios follow. For the properties of $R$: the numerator in $W_\rho$ is at least $c(x)$ and its denominator at most $x-(1-\phi)$, so $W_\rho(\phi)\ge V(\phi)$. The infimum defining $V(\phi)$ is attained at some $x_\phi\in(1-\phi,1)$, because the quotient tends to infinity at both ends of the interval; for $\rho$ small enough, $(1+\rho)(1-x_\phi)<\phi$, and $W_\rho(\phi)\le(c(x_\phi)+\rho\,\ell(x_\phi))/(\phi-(1+\rho)(1-x_\phi))$, which tends to $V(\phi)$ as $\rho\to0$.

(ii) We apply \cref{cor:noise:bcrit} to SGD and, through \cref{thm:noise:precond}(i), to $P=H^p$. For SGD, $S=\kappa h$, and the modes with $h_i\le S/4$ are the bulk, because $\kappa\ge4$; so $D^\noisesl(\tau)=L_{\mathrm b}e^{-2h\tau}$, $\tau^\noisesl_\eps=\log(1/\phi)/(2h)$, and $\noiseTmin\ge\frac18\kappa\log(1/\phi)$. For $H^p$, $S_P=(\kappa h)^{1-p}$, and $D_P(\tau)\le(1+\rho)L_{\mathrm b}e^{-2h^{1-p}\tau}$, so that $D_P$ falls below $\eps$ by the time $\log((1+\rho)/\phi)/(2h^{1-p})$; hence $\noiseTmin^{H^p}\le\frac12\kappa^{1-p}\log((1+\rho)/\phi)+1$. Dividing gives the bound.
\end{proof}

For the spiked spectra of \cref{fig:noise:precond} ($\kappa=100$, $k=10$, target $5\%$ of the initial loss), the bound $R$ is $2.39$ for $\rho=1$ and $1.12$ for $\rho=0.1$, and the observed small-batch advantages of $p=\frac12$ are $1.83$ and $1.09$; the large-batch advantages are $9.1$ and $9.8$, close to $\kappa^{1/2}=10$, so the lower bound (ii) is pessimistic by a factor of four to five. In words: when the noise lives in the bulk, no diagonal preconditioner can save examples at small batch sizes, because the bulk is already well conditioned and the outliers carry too little noise to matter; at large batch sizes the outliers set the stable step size, and preconditioning gains a factor of order $\kappa^p$.

Finally, the scaling rules. The linear rule is proved in the model; the square-root rule for Adam needs one more idealization.

\begin{heuristic}[Heuristic derivation of the scaling rules]
\emph{SGD.} In the noise-limited regime, the expected loss after $E$ examples at step size $\eta=\gamma B$ approaches $D(\gamma E)+\frac\gamma4Q(\gamma E)$, a function of $\gamma=\eta/B$ and $E$ only. It is the expected loss of the Ornstein--Uhlenbeck process $d\theta=-H\theta\,d\tau+\sqrt\gamma\,\Sigma^{1/2}dW_\tau$ at time $\tau=\gamma E$, which is the \term{stochastic differential equation}\index{stochastic differential equation} of \textcite{LiTaiE2017} for this model. Runs with the same $\eta/B$ have the same loss per example; that is the linear rule, and \cref{thm:noise:critical}(iv) shows that the tuned step size follows it. The rule fails when $\eta S$ is no longer small, because the discrete factor $(1-\eta h_i)^2$ then departs from $e^{-2\eta h_i}$ and stability caps $\eta$ below $2/S$.

\emph{Adam.} Take Adam with $\beta_1=0$, $\epsilon=0$ and $\beta_2$ close to $1$, in a stretch of training over which the gradient statistics change slowly, so that $\hat v_{t,i}\approx\E g_{t,i}^2=(\partial_i\Loss)^2+c_i/B$. While the noise dominates the signal in coordinate $i$, that is, while $B(\partial_i\Loss)^2\ll c_i$, the step on that coordinate is $\eta g_{t,i}/\sqrt{\hat v_{t,i}}\approx\eta\sqrt B\,g_{t,i}/\sqrt{c_i}$: preconditioned SGD with $P^{-1}=\diag(c_i^{-1/2})$ and step size $\eta\sqrt B$. For $\Sigma=\sigma^2H$ this is the power $p=\frac12$ of \cref{thm:noise:precond}, up to the factor $\sigma$. The linear rule for this preconditioned SGD asks for $\eta\sqrt B\propto B$, that is, $\eta\propto\sqrt B$: the \term{square-root scaling rule}\index{square-root scaling rule}. For $\hat v$ to average over a fixed number of examples, its horizon of $1/(1-\beta_2)$ steps must shrink like $1/B$, so $1-\beta_2\propto B$, and likewise $1-\beta_1\propto B$; and since $\sqrt{\hat v}$ scales like $B^{-1/2}$, Adam's $\epsilon$ must scale the same way to keep its relative size. The rule fails when $B(\partial_i\Loss)^2$ is no longer small compared with $c_i$: then $\hat v_{t,i}\approx(\partial_i\Loss)^2$, the step approaches the sign step $\eta\sign(\partial_i\Loss)$ of \cref{ch:sign}, whose size does not depend on $B$. It also fails, as the linear rule does, when the effective step size approaches the stability limit $2/S_P$.
\end{heuristic}

\Textcite{MalladiEtAl2022} derived the square-root rule from a stochastic differential equation for adaptive methods, and \textcite{LiMalladiArora2021} gave conditions under which such equations, and with them the scaling rules, describe SGD.

\section{Reading the theory}\label{sec:noise:reading}

\paragraph{Loss curves.} \Cref{fig:noise:curves} checks \cref{thm:noise:nqm} against simulation and shows the trade-off behind everything else. At $B=4$, the step size $0.10$ has a floor of $0.047$ but is still above the target after $4000$ steps; the step size $1.00$ is fast but levels off at $0.55$, above the target; $0.26$ is close to the best compromise and reaches the target after about $2600$ steps. At $\eta=0.26$ the three batch sizes follow one curve until each reaches its own floor, $0.50$, $0.13$ and $0.031$: larger batches at a fixed step do not train faster, they train lower.

\begin{figure}[t]
\begin{fullwidth}
\centering
\begin{tikzpicture}
\begin{groupplot}[group style={group size=2 by 1,horizontal sep=17mm},wplot,width=82mm,height=60mm,xmode=log,xmin=1,xmax=16384,xlabel={batch size $B$},xtick={1,10,100,1000,10000}]
\nextgroupplot[ymode=log,ymin=5,ymax=20000,ylabel={steps to target $T(B)$},legend style={at={(0.97,0.97)},anchor=north east,font=\scriptsize}]
\addplot[theory] table[x=B,y=sgd] {figures/data/11-noise-steps.dat};
\addplot[theory3] table[x=B,y=mom] {figures/data/11-noise-steps-mom.dat};
\addplot[theory2] table[x=B,y=pre] {figures/data/11-noise-steps.dat};
\addplot[guide] table[x=B,y=hypsgd] {figures/data/11-noise-steps.dat};
\addplot[simulation,color=sea,mark options={fill=sea,draw=white,line width=0.2pt}] table[x=B,y=sgd] {figures/data/11-noise-steps-mc.dat};
\addplot[simulation] table[x=B,y=mom] {figures/data/11-noise-steps-mc.dat};
\addplot[simulation2] table[x=B,y=pre] {figures/data/11-noise-steps-mc.dat};
\addplot[guide] table[x=B,y=hyppre] {figures/data/11-noise-steps.dat};
\legend{SGD,{momentum, $\beta=0.9$},{preconditioned, $H^{-1/2}$},hyperbola $\noiseTmin+\noiseEmin/B$}
\nextgroupplot[ymode=log,ymin=0.005,ymax=60,ylabel={best step size},legend style={at={(0.97,0.03)},anchor=south east,font=\scriptsize}]
\addplot[theory] table[x=B,y=etasgd] {figures/data/11-noise-steps.dat};
\addplot[theory3] table[x=B,y expr=\thisrow{etamom}/0.1] {figures/data/11-noise-steps-mom.dat};
\addplot[theory2] table[x=B,y=etapre] {figures/data/11-noise-steps.dat};
\addplot[guide] table[x=B,y=linsgd] {figures/data/11-noise-steps.dat};
\addplot[guide] table[x=B,y=linpre] {figures/data/11-noise-steps.dat};
\addplot[guide,domain=1:16384,samples=2] {2};
\addplot[guide,domain=1:16384,samples=2] {38};
\legend{$\eta$ for SGD,$\eta/(1-\beta)$ for momentum,$\eta$ for $H^{-1/2}$}
\end{groupplot}
\end{tikzpicture}
\end{fullwidth}
\caption{Steps to the target $\eps=0.05\,\Delta_0$ in the harmonic model, at the best constant step size. Left: curves are exact (\cref{thm:noise:nqm}, minimized over the stable step sizes; heavy ball by its exact second-moment recursion); points are simulations: the first step at which the average loss of $64$ runs ($256$ runs for $B\ge128$), at the step size chosen by the theory and smoothed over $2\%$ of the predicted time, is below the target (seed fixed in \texttt{11-noise-steps.py}). All agree to within $1.5\%$. Dashed: the hyperbolas $\noiseTmin+\noiseEmin/B$ built from the exact limits of \cref{thm:noise:critical}. Right: the best step sizes. Dashed: the linear rules $\gamma^\star B$ of \cref{thm:noise:critical}(iv), the limit $2/S=2/S_P=2$, and momentum's limit $(2+2\beta)/((1-\beta)S)=38$ for the effective step size.}
\label{fig:noise:steps}
\end{figure}

\paragraph{Steps against batch size.} \Cref{fig:noise:steps} is the theory's version of \cref{obs:noise:regimes}. For SGD the steps fall like $1/B$ up to a batch size of about $10$, bend between $10$ and $100$, and level off near $\noiseTmin=177$; the right panel shows the reason, the best step size growing like $\gamma^\star B$ with $\gamma^\star=0.0678$ until it approaches the stability limit, after which it creeps up to $1.98$. The linear rule reaches $1/S$ at $B=1/(\gamma^\star S)\approx15$ and the large-batch step size at $B\approx29$; the critical batch size $\noiseEmin/\noiseTmin=56$ lies a factor of two further out. The best large-batch step size is within one percent of $2/S$: in the curvature-limited regime the tuned learning rate sits at the edge of stability, as it does in practice (\cref{ch:eos}). Momentum with $\beta=0.9$ follows SGD at small batch sizes, as \cref{prop:noise:momentum} requires, and levels off at $23$ steps, eight times fewer, so its critical batch size is $434$. Preconditioning by $H^{-1/2}$ moves the whole small-batch line down by the factor $2.95$ and the large-batch level from $177$ to $10$ steps; its critical batch size is $339$. Defined instead as the batch size at which $T(B)=2\noiseTmin$, the knee moves from about $46$ to about $190$, a factor of four, while $\noiseBcrit$ moves by a factor of six.

\paragraph{The hyperbola.} If $T(B)=\noiseTmin+\noiseEmin/B$ held exactly, then $E=B\,T$ would satisfy McCandlish et al.'s relation $(T/\noiseTmin-1)(E/\noiseEmin-1)=1$. \Cref{thm:noise:critical} proves the two asymptotes, and in between the hyperbola is a model. For SGD the exact $T(B)$ lies between $0.89$ and $1.04$ times the hyperbola built from the exact limits; for $H^{-1/2}$, between $0.69$ and $1.00$. The product form is a much more sensitive test: computed from the exact first-passage times it rises from $0.43$ as $B\to0$ to $3.5$ at $B=16384$ for SGD, and stays between $0.11$ and $0.25$ for $H^{-1/2}$, instead of $1$. The reason is that the product compares corrections: as $B\to0$ it tends to $E_1/\noiseTmin$, where $E(B)=\noiseEmin+E_1B+o(B)$, and McCandlish et al.'s derivation assumes $E_1=\noiseTmin$, which nothing in the model enforces. A least-squares fit of the hyperbola to $T(B)$ is accurate to $8\%$ for SGD and gives $\noiseBcrit\approx52$ (Exercise~\ref{ex:noise:simfit}).

\paragraph{Noise scales.} McCandlish et al.'s noise scales, evaluated on the noise-free path of gradient flow at the moment it reaches the target, are $\noiseBnoise\approx9.5\times10^5$ and $\noiseBsimple\approx1.1\times10^4$, and they overestimate $\noiseBcrit=56$ by orders of magnitude. The one-step model behind them lets the step size grow to $\eta_{\max}$, which near the target is close to $1/h$ for the slow modes that still carry loss, far beyond the stability limit (\cref{prop:noise:onestep}(ii)). Measured instead, as in practice, on the noisy iterates of the run, at the step size the tuning chose, they are $\noiseBsimple\approx132$ and $\noiseBnoise\approx40$ at $B=1$, the right order of magnitude. \Cref{prop:noise:onestep}(iii) explains why: at the floor the full-batch gradient is dominated by the fluctuations of the stiff modes, and the measured scales are $2\tr\Sigma/(\gamma\tr(H\Sigma))$ and $2\tr(H\Sigma)/(\gamma\tr(H^2\Sigma))$ as $\eta=\gamma B\to0$. They measure the inverse of the step size per example that the tuning found, and through $\gamma^\star$ they inherit the scale $\tr\Sigma/\eps$ of the critical batch size. For an isotropic Hessian the agreement is close: then $\noiseTmin=1$, $\noiseBcrit=\noiseEmin$, and with $\Delta_0/\eps=20$ the critical batch size is $1.42$ times $\noiseBsimple$ evaluated at loss $\eps$ (Exercise~\ref{ex:noise:onemode}).

\begin{figure}[t]
\centering
\begin{tikzpicture}
\begin{axis}[wplot,xmode=log,ymode=log,xmin=0.5,xmax=32768,ymin=0.8,ymax=30,xlabel={batch size $B$},ylabel={advantage $T(B)/T^{H^{1/2}}(B)$},xtick={1,10,100,1000,10000},ytick={1,2,5,10,20},yticklabels={1,2,5,10,20},legend style={at={(0.03,0.97)},anchor=north west,font=\scriptsize}]
\addplot[theory] table[x=B,y=harmonic] {figures/data/11-noise-precond.dat};
\addplot[theory3] table[x=B,y=spike1] {figures/data/11-noise-precond.dat};
\addplot[theory2] table[x=B,y=spike01] {figures/data/11-noise-precond.dat};
\addplot[simulation,color=sea,mark options={fill=sea,draw=white,line width=0.2pt}] table[x=B,y=harmonic] {figures/data/11-noise-precond-mc.dat};
\addplot[simulation] table[x=B,y=spike1] {figures/data/11-noise-precond-mc.dat};
\addplot[simulation2] table[x=B,y=spike01] {figures/data/11-noise-precond-mc.dat};
\addplot[guide,domain=0.5:32768,samples=2] {1};
\addplot[guide,line width=1pt] table[x=B,y=harmonic] {figures/data/11-noise-precond-lim.dat};
\addplot[guide,line width=1pt] table[x=B,y=spike1] {figures/data/11-noise-precond-lim.dat};
\addplot[guide,line width=1pt] table[x=B,y=spike01] {figures/data/11-noise-precond-lim.dat};
\legend{{harmonic, $h_i=1/i$},{spiked, $\rho=1$},{spiked, $\rho=0.1$}}
\end{axis}
\end{tikzpicture}
\caption{The advantage of preconditioning by $H^{-1/2}$ over SGD, the ratio of their steps to target, against the batch size, with $\Sigma=H$, $\theta_0=(1,\dots,1)$ and the target $5\%$ of the initial loss. Harmonic model as in \cref{fig:noise:curves}; spiked models with $10$ outliers of curvature $1$ and $m=1000$ ($\rho=1$) or $m=10^4$ ($\rho=0.1$) bulk modes of curvature $0.01$. Curves: exact real-valued first-passage times; short dashes at the ends: the limits $\noiseEmin/\noiseEmin^P$ and $\noiseTmin/\noiseTmin^P$ of \cref{thm:noise:precond}(ii). Points: the same ratio from simulations as in \cref{fig:noise:steps}, with the crossing of the target located by linear interpolation; bulk modes are simulated by $500$ representative modes (\texttt{11-noise-precond.py}).}
\label{fig:noise:precond}
\end{figure}

\paragraph{The advantage of preconditioning.} \Cref{fig:noise:precond} is the batch-size dependence that \cref{ch:transformers} uses to compare Adam with SGD. In all three spectra the advantage grows with the batch size, from its small-batch limit to about its large-batch limit, which two of the curves overshoot slightly near the knee of the preconditioned curve. The small-batch limit is a property of the problem: $2.95$ for the harmonic spectrum, where SGD wastes noise on stiff modes that carry a large share of the trace, $1.83$ for the spiked spectrum with $\rho=1$, and $1.09$ for $\rho=0.1$, where the bulk carries $91\%$ of the noise. So the advantage does not vanish at small batch sizes in general: it tends to a constant set by the spectrum, the noise and the initialization through \eqref{eq:noise:Emin}, and that constant tends to $1$ when the outliers' share of the noise does (\cref{prop:noise:spiked}).

\paragraph{Where the agreement must fail.} The comparisons in the figures are between the exact expected loss and averages of independent runs; the residual differences are Monte Carlo error and the bias of locating a first passage on a noisy curve, which is largest where the loss approaches the target slowly, at small batch sizes for SGD. A single run fluctuates around the expected loss: at the target of the harmonic model with $B=4$ the loss of one run has a standard deviation of $17\%$ of the target, and since the fluctuations of the stiff modes decorrelate within a few steps, a single run first touches the target well before the expected loss reaches it (median $1981$ steps in $200$ runs, against $T(4)=2577$). The theory describes the expected loss, or a loss averaged over many runs or over a window of steps; it does not describe the first touch of one run. And the steps are integers, which produces the small steps in the curves at large batch sizes.

\section{What the theory misses}\label{sec:noise:misses}

\paragraph{What was idealized away.} The landscape and the noise are fixed. In a network both change during training: the sharpness rises toward the stability threshold (\cref{ch:eos}), the noise covariance changes with the parameters, and near interpolation it shrinks with the loss, which changes the critical batch size from $\tr\Sigma/(S\eps)$ to the target-independent effective rank $\tr H/S$ (Exercise~\ref{ex:noise:interp}). The model is online, so it has no epochs and no generalization gap, and it says nothing about how the critical batch size depends on the data set or on the model size, which enter only through $\Sigma$, $H$ and $\eps$. The step size is constant: an optimal schedule, or averaging of the iterates, changes the small-batch regime, and with averaging every diagonal preconditioner reaches the same asymptotic cost $\tr(H^{-1}\Sigma)/(2E)$ \parencite{PolyakJuditsky1992} (Exercise~\ref{ex:noise:averaging}). The preconditioner is a fixed function of $H$, whereas Adam's is estimated from noisy gradients and lags behind them (\cref{ch:adaptive}). The object is the expected loss; a practitioner sees one run.

\paragraph{Where it breaks.} (1) The critical batch size of a real run moves. \Textcite{McCandlishEtAl2018} found the noise scale growing during training, and \textcite{KaplanEtAl2020} fitted the critical batch size of language models as a power of the loss; the model's $1/\eps$ captures the direction of this growth, but not its rate, and \textcite{ZhangEtAl2025CBS} found that in pretraining it grows mainly with the amount of training data, not with the model size, a dependence the model cannot express because it has no data set. (2) The best learning rate does not follow the model's rule, linear and then constant, in every workload: \textcite{ShallueEtAl2019} found that the optimal learning rate followed neither the linear nor the square-root rule consistently across workloads. (3) Large batches can generalize worse at a fixed number of epochs \parencite{KeskarEtAl2017,HofferHubaraSoudry2017}, an effect the online model cannot show. (4) At very small batch sizes, \textcite{MarekEtAl2025} found plain SGD close to Adam on language models. In the harmonic model the advantage of preconditioning stays near $3$ as $B\to0$; it vanishes only if the noise lives in the bulk of a spiked spectrum (\cref{prop:noise:spiked}). Which spectrum describes a transformer is not settled: its Hessian is heterogeneous across blocks \parencite{ZhangEtAl2024Hessian}, and \cref{ch:transformers} weighs this against the other explanations of the gap. (5) The stochastic-differential-equation approximation behind the scaling rules has its own conditions of validity, and \textcite{LiMalladiArora2021} show experimentally where it fails; heavy-tailed gradient noise (\cref{ch:spikes}) is outside the model altogether.

\paragraph{Repairs and open problems.} Batch-size schedules exploit the growth of the critical batch size: \textcite{SmithEtAl2018} increase the batch size in place of decaying the learning rate, GPT-3 ramped its batch size up \parencite{BrownEtAl2020}, and \textcite{MerrillEtAl2025}, in an unrefereed preprint, measure the critical batch size of language models over a run and propose a batch-size warmup based on it. Noise-scale estimators make the critical batch size measurable during training \parencite{McCandlishEtAl2018}, with the caveat of \cref{prop:noise:onestep}(iii) that they partly measure the run's own step size. Layerwise normalized updates, LARS \parencite{YouGitmanGinsburg2017} and LAMB \parencite{YouEtAl2020}, were proposed for very large batches; \textcite{NadoEtAl2021} found that well-tuned standard optimizers match them. Models with parameter-dependent noise give the interpolation regime its own critical batch size \parencite{MaBassilyBelkin2018,JainEtAl2018}. Open: a solvable model in which the dependence of the critical batch size on the data-set size and the model size is derived rather than fitted; a theory of the critical batch size for normalized optimizers, sign descent and Muon, whose step length does not shrink with the noise (\cref{ch:momentum}); and a measurement of where transformer spectra and noise covariances lie between the harmonic and the spiked cases.

\section{Back in the wild}\label{sec:noise:back}

\fieldnotes
{Scale the learning rate linearly (SGD) or as the square root (Adam) with the batch size, up to a critical batch size; train language models with batches of millions of tokens, ramped up early in training.}
{The noisy quadratic model: a fixed quadratic, gradient noise of covariance $\Sigma/B$, a constant step size, and the expected loss.}
{Exact loss curves and floors; $T(B)=\noiseTmin$ for large $B$, and $B\,T(B)\to\noiseEmin$ with $\eta\approx\gamma^\star B$ as $B\to0$; $\noiseBcrit\asymp\tr\Sigma/(S\eps)$ up to shape factors; preconditioning raises this scale; in a spiked spectrum no diagonal preconditioner saves examples at small batch sizes.}
{Perfect scaling below $\noiseBcrit$, saturation above it with the step size near $2/S$; a critical batch size growing like $1/\eps$ as the target falls; momentum and preconditioners help mainly above the knee; noise scales measured on a run reflect its own $\eta/B$.}
{Curvature and noise that change during training, noise that depends on the parameters, finite data and epochs, schedules, heavy tails, and generalization.}
{Scale $\eta\propto B$ below the knee and expect diminishing returns above it (proved in the idealization); expect a larger critical batch size late in training and for preconditioned optimizers (suggested); do not expect a large gain from preconditioning at very small batch sizes when the noise lives in the bulk of the spectrum (suggested).}

For a practitioner the theory turns the scaling rules from recipes into statements with a range of validity. The linear rule is a theorem of the model, but only while the step size it prescribes is small compared with $1/S$; once the tuned learning rate approaches the stability limit, the rule stops, and so do the returns of larger batches. The critical batch size is of order $\tr\Sigma/(S\eps)$: it grows as the loss falls, which supports ramping the batch size up during training; it grows with the noise; and it falls with the sharpness. What the theory does not provide is the constant in front, the shape factors of \cref{cor:noise:bcrit}, or the way $\Sigma$ and $S$ evolve in a network; the batch size is still chosen by measuring. The common advice to train near the critical batch size is a sound trade-off between steps and examples in the model, but where exactly to sit depends on the relative cost of the two, which is folklore territory.

The chapter connects to others through its two constraints. The stability constraint is the threshold $2/S$ of \cref{ch:eos}, with momentum's region $2+2\beta$ from \cref{ch:momentum}; the floor is the starting point of \cref{ch:schedules}, where a decaying step size removes it and the river-valley model adds a slow direction to it, and of \cref{ch:decay}, whose equilibrium is another floor. The batch-size dependence of the advantage of preconditioning, \cref{thm:noise:precond} and \cref{prop:noise:spiked}, is the third explanation that \cref{ch:transformers} offers for why Adam beats SGD on transformers, and the heuristic for the square-root rule connects to the choice of Adam's $\beta_2$ in \cref{ch:adaptive}.

\begin{exercises}
\item \consolidation (Momentum's floor.) For the heavy ball in one mode, $x_{t+1}=(1+\beta-\eta h)x_t-\beta x_{t-1}+w_t$ with uncorrelated $w_t$ of variance $s^2$, derive $\E x^2=s^2(1+\beta)/((1-\beta)\eta h(2(1+\beta)-\eta h))$ at stationarity. Show that with $\eta_{\mathrm e}=\eta/(1-\beta)$ fixed the floor of \cref{prop:noise:momentum} decreases as $\beta$ increases, and that, when $\eta_{\mathrm e}S<2$, its relative difference from SGD's floor at $\eta_{\mathrm e}$ is at most $\eta_{\mathrm e}S/(2-\eta_{\mathrm e}S)$.\label{ex:noise:momfloor}

\item \consolidation (One mode.) Let $d=1$, with curvature $h$, noise $c$ and $r=\Delta_0/\eps>1$. (a) Show that $\noiseTmin=1$, attained at $\eta=1/h$, and that $T(B)=1$ if and only if $B>(1-1/r)\,c/(2h\eps)$. (b) Show that $\noiseEmin=\frac{c}{8h\eps}V_1(r)$ with $V_1(r)=\min_{x>\log r}x(1-e^{-x})/(1-re^{-x})$, and that $V_1(r)=\log r+\log\log r+O(1)$ as $r\to\infty$. (c) Deduce that $E(B)$ can be smaller than $\noiseEmin$: for $r=20$, where $V_1(20)=5.69$, a batch size just above $(1-1/r)c/(2h\eps)$ reaches the target in one step, with about two thirds of $\noiseEmin$ examples. So the assumption $E\ge\noiseEmin$ behind McCandlish et al.'s relation can fail in the curvature-limited regime. (d) For $H=hI$ in $d$ dimensions, show that $\noiseBcrit=\noiseEmin$ is $V_1(r)/4$ times $\noiseBsimple$ evaluated at loss $\eps$. \hint{in (b), put $x=\log r+s$ and minimize over $s$.}\label{ex:noise:onemode}

\item \consolidation (Noise scales at a floor.) (a) Derive \eqref{eq:noise:onestep} without assuming that the noise is Gaussian. (b) Using \cref{prop:noise:onestep}(iii) and \cref{thm:noise:critical}(iv), show that for a run at the best step size the measured scale $\tr\Sigma/\bar g_1$ tends to $2\tr\Sigma/(\gamma^\star\tr(H\Sigma))$ as $B\to0$, if the run is at its floor. (c) Evaluate it for the harmonic model, where $\tr\Sigma=7.49$, $\tr(H\Sigma)=1.64$ and $\gamma^\star=0.0678$, and compare with $\noiseBcrit=56$.\label{ex:noise:noisescale}

\item \consolidation (A preconditioner that does not commute with $H$.) Let $P\succ0$ be arbitrary. Show that $\theta'=P^{1/2}\theta$ performs SGD on a noisy quadratic model with Hessian $H'=P^{-1/2}HP^{-1/2}$ and noise covariance $P^{-1/2}\Sigma P^{-1/2}$, that $\lmax(H')=\lmax(P^{-1}H)=S_P$, and that \cref{thm:noise:nqm} applies in an eigenbasis of $H'$, where the relevant noise entries are the diagonal entries of $P^{-1/2}\Sigma P^{-1/2}$ in that basis.\label{ex:noise:change}

\item \consolidation (The diffusion limit.) Show that the Ornstein--Uhlenbeck process $d\theta=-H\theta\,d\tau+\sqrt\gamma\,\Sigma^{1/2}dW_\tau$ started at $\theta_0$ has $\E\Loss(\theta_\tau)=D(\tau)+\frac\gamma4Q(\tau)$, and that the infimum, over all $\gamma>0$, of the number of examples $E=\tau/\gamma$ after which this is below $\eps$ is $\noiseEmin$. Explain why in this limit the best step size is proportional to the batch size, and why the limit says nothing about step sizes with $\eta S$ close to $2$.\label{ex:noise:ou}

\item \variation (Interpolation.) Least squares with Gaussian inputs $x\sim\cN(0,H)$ and noiseless labels has per-example gradient covariance $\Sigma(\theta)=2\Loss(\theta)H+H\theta\theta^\top H$. (a) Show that the vector $m_t=(\E\theta_{t,i}^2)_i$ of SGD satisfies $m_{t+1}=Am_t$, with $A=\diag(a_i)+(\eta^2/B)hh^\top$, $a_i=(1-\eta h_i)^2+\eta^2h_i^2/B$ and $h=(h_1,\dots,h_d)$. (b) Show that the expected loss tends to $0$ from every $\theta_0$ if and only if $\max_ia_i<1$ and $\frac\eta B\sum_ih_i/(2-\eta h_i(1+1/B))<1$. (c) Deduce that the largest stable step size satisfies $\min\{B/\tr H,\,B/((B+1)S)\}\le\eta_{\max}(B)\le\min\{2B/\tr H,\,2B/((B+1)S)\}$: it grows linearly in $B$ up to about the effective rank $\tr H/S$ and saturates near $2/S$ beyond. Here the critical batch size is $\tr H/S$, whatever the target. \hint{$A$ is entrywise positive; its largest eigenvalue exceeds $\max_ia_i$ and solves $1=(\eta^2/B)\sum_ih_i^2/(\mu-a_i)$.}\label{ex:noise:interp}

\item \variation (Isotropic noise.) Repeat the proof of \cref{prop:noise:spiked}(i) for $\Sigma=\sigma^2I$. Show that now $\tau\pi_ic_i(1-e^{-u_i})=\frac{\sigma^2}{2h_i}c(x_i)$, that the outliers contribute at most $\frac{\sigma^2}{2h}k\kappa^{-p}\ell(x)$ to $\tau Q_P$ for $P=H^p$, and that $\noiseEmin/\noiseEmin^P$ again tends to $1$ for every diagonal $P$ as $k\kappa/m\to0$. In which modes does preconditioning by $H^{-p}$ raise the floor when the noise is isotropic?\label{ex:noise:isotropic}

\item \variation\hard (Averaging.) Let $\theta_1,\dots,\theta_n$ be consecutive iterates of SGD with step size $\eta<2/S$ in its stationary regime, and let $\bar\theta_n=\frac1n\sum_t\theta_t$. (a) Show that $\E\Loss(\bar\theta_n)=\frac1{2nB}\tr(H^{-1}\Sigma)\,(1+o(1))$ as $n\to\infty$, whatever $\eta\in(0,2/S)$ is. (b) Show that the same holds for SGD preconditioned by every diagonal $P$. (c) Conclude that with averaging, and for targets small enough that the stationary regime dominates, the number of examples to the target is about $\tr(H^{-1}\Sigma)/(2\eps)$ for every diagonal preconditioner, as in the theorem of \textcite{PolyakJuditsky1992}. What, then, is left for a preconditioner to improve? \hint{each mode is an autoregression of order one; sum its autocovariances.}\label{ex:noise:averaging}

\item \variation (The square-root rule for sign descent.) Let $\theta_{t+1,i}=\theta_{t,i}-\eta\sign(g_{t,i})$ with Gaussian noise, and write $\mu=\partial_i\Loss(\theta_t)$ and $s^2=c_i/B$. (a) Show that $\E[\sign g_{t,i}\mid\cF_t]=\erf(\mu/(s\sqrt2))=\sqrt{2/\pi}\,\mu/s+O((\mu/s)^3)$. (b) Deduce that, to leading order in $\mu/s$, the conditional mean of the sign step is that of SGD preconditioned by $\pi_i=\sqrt{2B/(\pi c_i)}$, and its conditional variance is $\pi/2$ times that of this preconditioned SGD. (c) For $\Sigma=\sigma^2H$, identify the power $p$ of \cref{thm:noise:precond}, and derive the square-root rule $\eta\propto\sqrt B$ in the noise-dominated regime. Why does it fail once $B\mu^2\gtrsim c_i$?\label{ex:noise:sign}

\item \simulation (The momentum coefficient.) In \texttt{11-noise-steps.py}, replace $\beta=0.9$ by $0.5$, $0.8$ and $0.95$. Check that the small-batch part of the curve does not move, as \cref{prop:noise:momentum}(ii) requires, while the large-batch level becomes $59$, $23$ and $41$ steps (against $177$ for SGD and $23$ for $\beta=0.9$). Compare the best $\beta$ with Polyak's choice $(\sqrt\kappa-1)^2/(\sqrt\kappa+1)^2$ for the condition number $\kappa\approx700$ of the modes that must converge (\cref{ch:momentum}), and explain why a $\beta$ that is too large hurts.\label{ex:noise:simbeta}

\item \simulation (The outliers' share.) Using \texttt{11-noise-precond.py}, compute the small-batch advantage of $H^{-1/2}$ for spiked spectra with $\kappa=100$, $k=10$ and $\rho\in\{1,0.1,0.01\}$, and compare with the bound $R(\rho,\phi)$ of \cref{prop:noise:spiked}, with $\phi=0.05(1+\rho)$. You should find $1.83$, $1.09$ and $1.009$ against $2.39$, $1.12$ and $1.012$. Then raise $\kappa$ to $1000$ at $\rho=0.1$ and check that the small-batch advantage stays near $1$ (it is $1.10$) while the large-batch one grows to about $\kappa^{1/2}$ (it is $31$).\label{ex:noise:simspike}

\item \simulation (Fitting the hyperbola.) Fit $T(B)=a+b/B$ to the exact SGD curve of \cref{fig:noise:steps} by least squares on $\log T$. You should find $a\approx180$, $b\approx9300$ and $b/a\approx52$, within $8\%$ of the curve everywhere, against the exact limits $\noiseTmin=177$ and $\noiseEmin=9987$. Repeat for $H^{-1/2}$, where the best fit is off by $31\%$ near the knee, and explain from \cref{fig:noise:steps} why the hyperbola fits the preconditioned curve worse.\label{ex:noise:simfit}

\item \research (The critical batch size and the data.) \Textcite{ZhangEtAl2025CBS} report that in pretraining the critical batch size grows with the amount of data rather than with the model size. The noisy quadratic model has no data set. Formulate a solvable model, for instance linear regression on a finite sample trained for several epochs, or a power-law random-feature model in which the target $\eps$ is tied to the data and model sizes by a scaling law, and derive how $\noiseBcrit$ depends on both. What is open is a model in which this dependence is derived and matches the measurements.\label{ex:noise:data}

\item \research (Normalized steps.) For sign descent and Muon the step length does not shrink with the noise, and without momentum they have a floor that does not vanish as $\eta\to0$ at fixed $B$ (\cref{ch:momentum}). Define a critical batch size for normalized steepest descent with momentum on the noisy quadratic model and compute it. Is it again of order $\tr\Sigma/(S\eps)$ in a suitable geometry? As far as we know, this is open.\label{ex:noise:normalized}
\end{exercises}

\notesandrefs
Quadratic models with additive gradient noise have long served to reason about step sizes: \textcite{SchaulZhangLeCun2013} derived adaptive learning rates from one, and \textcite{WuEtAl2018} used one to explain the short-horizon bias of learned learning rates. \Textcite{ZhangEtAl2019NQM} used the noisy quadratic model, with the per-mode recursion of \cref{thm:noise:nqm} and its floor, to predict how the critical batch size depends on the optimizer, including preconditioning by powers of $H$ and momentum; their solutions of the model show perfect scaling, its end, and how preconditioning and momentum move the end. Constant-step SGD on a quadratic is an Ornstein--Uhlenbeck process in the diffusion limit, and \textcite{MandtHoffmanBlei2017} studied its stationary distribution; the floor proportional to the step size is a standard part of SGD theory \parencite{BottouCurtisNocedal2018}, and averaging removes it \parencite{PolyakJuditsky1992}, with asymptotic optimality in the sense of local minimax theory \parencite{DuchiRuan2021}.

The scaling rules have two sources. Empirically, \textcite{Krizhevsky2014}, \textcite{GoyalEtAl2017} and \textcite{HofferHubaraSoudry2017} studied how the learning rate should follow the batch size, and \textcite{SmithLe2018} and \textcite{SmithEtAl2018} argued that the ratio of learning rate to batch size controls the noise of SGD. Theoretically, \textcite{LiTaiE2017,LiTaiE2019} introduced stochastic modified equations, \textcite{LiMalladiArora2021} examined when they are valid, and \textcite{MalladiEtAl2022} derived the square-root rule for adaptive methods. \Textcite{KidambiEtAl2018} showed that heavy-ball and Nesterov momentum do not improve on SGD for stochastic least squares, the counterpart of \cref{prop:noise:momentum}(ii).

The critical batch size was named and measured by \textcite{McCandlishEtAl2018}, who also derived the one-step model of \cref{prop:noise:onestep}; \textcite{ShallueEtAl2019} mapped the three regimes across workloads, and \textcite{KaplanEtAl2020} folded the critical batch size into the scaling laws of language models. In the interpolation regime of least squares, \textcite{MaBassilyBelkin2018} and \textcite{JainEtAl2018} found a critical batch size set by the squared norm of the inputs, a trace for Gaussian inputs, divided by the top eigenvalue of the Hessian (Exercise~\ref{ex:noise:interp}). Recent work on language models measures how the critical batch size grows with data \parencite{ZhangEtAl2025CBS} and over a run \parencite{MerrillEtAl2025}, and how small batches can be \parencite{MarekEtAl2025}.

The two limits of \cref{thm:noise:critical} with the variational formula \eqref{eq:noise:Emin}, the bounds of \cref{cor:noise:bcrit}, the spiked-spectrum proposition and the stationary noise scales of \cref{prop:noise:onestep}(iii) are, as far as we know, not stated in this form elsewhere; the regimes themselves, and the numerical evidence that preconditioning moves the knee, are in \textcite{ZhangEtAl2019NQM}.

\paragraph{Translating the literature.} \Textcite{McCandlishEtAl2018} write $\epsilon$ for the learning rate (our $\eta$), $S$ and $E$ for the steps and examples to a target (our $T$ and $E$; our $S$ is the sharpness), $G$ for the true gradient (our $\nabla\Loss$) and $S_{\min}$, $E_{\min}$ for the limits (our $\noiseTmin$, $\noiseEmin$); their $B_{\mathrm{crit}}=E_{\min}/S_{\min}$, $B_{\mathrm{noise}}$ and $B_{\mathrm{simple}}$ are ours. \Textcite{ZhangEtAl2019NQM} call the noise covariance $C$ (our $\Sigma$), with diagonal entries $c_i$ in the eigenbasis of $H$ as here, and precondition by $P=H^p$, as in \cref{thm:noise:precond}. \Textcite{KaplanEtAl2020} express the critical batch size in tokens as a function of the loss, which in the noisy quadratic model corresponds to $\noiseBcrit$ as a function of the target $\eps$.

\printchapterbib
\end{refsection}

% Chapter 12 of Optimizers in the Wild: learning-rate schedules.
\newcommand{\schedulesG}{\mathcal{G}}
\newcommand{\schedulesK}{\mathcal{K}}

\casechapter{Schedules}{Cosine, warmup--stable--decay and the river valley}{Cosine and WSD schedules}{ch:schedules}
\begin{refsection}[chapters/12-schedules.bib]

\begin{chaptersummary}
Large models are trained with a learning rate that rises briefly and then decays, along a cosine, along a straight line, or in a short cooldown after a long plateau, and the loss falls fastest exactly when the learning rate falls. For convex Lipschitz losses we prove a bound on the last iterate of stochastic gradient descent under an arbitrary schedule, show that linear decay to zero attains the rate of iterate averaging exactly, and evaluate the bound for the schedules in use. In a river valley, a slowly descending direction flanked by stiff and noisy ones, we compute the expected loss exactly for every schedule: it is a term that depends on the area under the schedule plus a term that remembers only the recent schedule, which is the form of the empirical laws that predict loss curves across schedules, and it predicts how the best cooldown scales with the length of the run. Neither idealization has a landscape that changes during training, and the river valley's predictions depend on constants that must be measured.
\end{chaptersummary}

\section{In the wild: the shape of the learning rate}\label{sec:schedules:wild}

Ask what learning rate a large language model was trained with and the answer is a curve. Most published pretraining recipes since GPT-3 warm the learning rate up over a small part of the run and then decay it to a small value at the planned end, and the shape of the decay is one of the few hyperparameters still argued about in public. As of 2025 three shapes dominate the open recipes.

\begin{practice}[GPT-3, \textcite{BrownEtAl2020}; LLaMA, \textcite{TouvronEtAl2023}; Chinchilla, \textcite{HoffmannEtAl2022}; MiniCPM, \textcite{HuEtAl2024}; \textcite{HageleEtAl2024}; \textcite{BergsmaEtAl2025}]
Warm the learning rate up linearly from zero to its peak $\eta_{\max}$, over the first $375$ million tokens (GPT-3) or $2000$ steps (LLaMA). Then do one of three things. \emph{Cosine}: decay along half a period of a cosine to a tenth of the peak, $\eta_t=\eta_{\max}\bigl(0.1+0.45(1+\cos(\pi t/T))\bigr)$; LLaMA and Chinchilla reach the tenth at the end of the run, GPT-3 after $260$ of its $300$ billion tokens. \emph{Linear decay to zero}, $\eta_t=\eta_{\max}(1-t/T)$ \parencite{DefazioEtAl2023,BergsmaEtAl2025}. \emph{Warmup--stable--decay} (WSD): hold the peak, then decay to zero or nearly zero over the last $10$--$20\%$ of the steps (MiniCPM; \textcite{HageleEtAl2024} recommend about $20\%$, with the shape $1-\sqrt{\cdot}$). The original transformer used the inverse square root, $\eta_t\propto\min(t^{-1/2},\,t\,T_{\mathrm w}^{-3/2})$ with $T_{\mathrm w}=4000$ \parencite{VaswaniEtAl2017}.
\end{practice}

The cosine\index{cosine schedule} came from vision, where \textcite{LoshchilovHutter2017} introduced cosine annealing with periodic restarts; language models kept a single half period.\sidenote{Earlier vision recipes divided the learning rate by $10$ at two or three fixed epochs, a \emph{step decay}. The classical rule of stochastic approximation is $\eta_t\propto1/t$, which satisfies the conditions $\sum_t\eta_t=\infty$ and $\sum_t\eta_t^2<\infty$ of \textcite{RobbinsMonro1951}. Every schedule in the box violates the second: at a fixed peak, $\sum_{t\le T}\eta_t^2$ grows without bound with the horizon $T$, linearly for the decaying shapes and like $\log T$ for the inverse square root.} The shapes look very different, yet their final losses in tuned runs are close, and which is best is still debated.

\begin{observation}[The cooldown drop; \textcite{HuEtAl2024}, \textcite{HageleEtAl2024}]\label[observation]{obs:schedules:cooldown}
Train with WSD\index{warmup--stable--decay (WSD)}\index{cooldown}. During the stable phase the loss decreases slowly; when the decay begins it falls much faster, and after a decay over $10$--$20\%$ of the steps it matches, or slightly beats, a cosine schedule of the same length. Longer decays help a little more, with diminishing returns.
\end{observation}

\begin{observation}[The horizon; \textcite{HoffmannEtAl2022}]\label[observation]{obs:schedules:horizon}
A cosine schedule planned for more steps than the run takes, by more than about $25\%$, gives a worse model than one planned for exactly the steps taken. WSD runs of different lengths coincide until their cooldowns, so one stable phase can be branched into cooldowns of several lengths, which makes scaling experiments cheaper \parencite{HageleEtAl2024}.
\end{observation}

The next two observations are of a different kind: fits and analogies that describe whole loss curves.

\begin{observation}[Loss curves across schedules; \textcite{TissueEtAl2025}, \textcite{LuoEtAl2025}]\label[observation]{obs:schedules:laws}
Fitted to the loss curves of one or two schedules, a law of the form
\[
  \Loss(t)\approx L_\infty+C\,S_1(t)^{-\alpha}-\mathrm{LD}(t),\qquad S_1(t)=\sum_{s\le t}\eta_s ,
\]
predicts the loss curves of the same model under other schedules\index{loss-curve laws}. The \emph{forward area} $S_1(t)$ is the area under the schedule so far; the \emph{loss drop} $\mathrm{LD}(t)$ adds, for each decrease of the learning rate, a term proportional to the decrease that grows and then saturates with the number of steps \parencite{TissueEtAl2025} or with the area \parencite{LuoEtAl2025} since the decrease.
\end{observation}

\begin{observation}[Convex theory agrees; \textcite{SchaippEtAl2025}]\label[observation]{obs:schedules:convex}
Evaluate, for each schedule and each step $t$, the bound on $\E\Loss(\theta_t)-\Loss^\star$ that convex optimization gives for stochastic gradient descent on a convex Lipschitz loss. The curves look like language-model loss curves: WSD plateaus and then drops suddenly, and cosine and WSD end close together. \Textcite{SchaippEtAl2025} report that the bound reproduces the matching performance of cosine and WSD and the necessity and length of the cooldown, and that learning rates chosen with it improved Llama-type models with $124$M and $210$M parameters.
\end{observation}

Language-model training is far from convex, so \cref{obs:schedules:convex} suggests that the schedule acts through something simpler than the landscape; in the experiments of \textcite{DefazioEtAl2023}, too, warmup followed by linear decay matched or beat cosine almost everywhere, as their worst-case theory predicts. Why does the loss fall when the learning rate falls, why do such different shapes end so close together, and how long must the decay be?

\section{The idealization: two pictures of a schedule}\label{sec:schedules:ideal}

Throughout, $\theta_1$ is the initial point and stochastic gradient descent with step sizes $\eta_1,\eta_2,\dots$ is
\begin{equation}\label{eq:schedules:sgd}
  \theta_{t+1}=\theta_t-\eta_t\,g_t ,\qquad t=1,2,\dots,
\end{equation}
with $g_t$ a stochastic gradient at $\theta_t$; ``learning rate'' (in practice) and ``step size'' (in the theory) are synonyms. Let $\cF_t$ be the $\sigma$-algebra generated by $g_1,\dots,g_{t-1}$, so that $\theta_t$ is $\cF_t$-measurable. A \term{schedule}\index{schedule} is the sequence $(\eta_t)$, often written $\eta_t=\eta\,\varphi(t/T)$ with a peak $\eta$ and a \emph{shape} $\varphi\colon[0,1]\to[0,1]$. The \term{area}\index{area under the schedule} before step $t$ is $A_t=\sum_{s=1}^{t-1}\eta_s$, the time of the gradient flow that the run follows when its steps are small.

\paragraph{Convex Lipschitz losses.} The first picture keeps nothing of the landscape except convexity.

\begin{assumption}[Convex Lipschitz stochastic optimization]\label[assumption]{ass:schedules:convex}
The loss $\Loss\colon\R^d\to\R$ is convex and has a minimizer $\theta^\star$, and $D=\norm{\theta_1-\theta^\star}_2$. The stochastic gradients satisfy $\E[g_t\mid\cF_t]\in\partial\Loss(\theta_t)$ and $\E\bigl[\norm{g_t}_2^2\mid\cF_t\bigr]\le G^2$.
\end{assumption}

The constant $G$ bounds subgradients and noise together, so the subgradient method on a $G$-Lipschitz function and stochastic gradient descent on a smooth function with $\norm{\nabla\Loss}_2^2+\E\norm{\zeta_t}_2^2\le G^2$ are both covered; the norm is Euclidean because stochastic gradient descent is steepest descent in $\ell_2$ (\cref{ch:steepest}). A statement about the whole class cannot see the scales of a particular landscape, and this will show in the numbers.

\paragraph{The river valley.} The second picture keeps the simplest landscape with two time scales. \Textcite{WenEtAl2025} proposed that the pretraining loss of a language model is shaped like a \term{river valley}\index{river valley}: a deep valley whose floor descends slowly along a river, between steep walls. With a large learning rate the iterate oscillates between the walls while it moves quickly along the river; when the learning rate decays the oscillation dies out, and the progress made along the river shows in the loss.\sidenote{\Textcite{WenEtAl2025} derive the picture in a toy model trained on bigram data and use it to design a WSD variant that reuses decayed checkpoints. The quadratic walls and the noise below are our simplification.} We make the walls quadratic, the oscillation noise-driven, and the river straight or weakly curved.

\begin{assumption}[The river valley]\label[assumption]{ass:schedules:river}
The parameter is $\theta=(x,y)$ with river coordinates $x\in\R^p$ and hill coordinates $y\in\R^m$, and
\[
  \Loss(x,y)=r(x)+\frac12\sum_{i=1}^m h_i\,y_i^2,\qquad h_i>0 .
\]
The stochastic gradient is $g_t=\nabla\Loss(\theta_t)+\zeta_t$ with $\E[\zeta_t\mid\cF_t]=0$ and $\E[\zeta_{t,i}^2\mid\cF_t]=\sigma_i^2$ for each hill coordinate $y_i$. The river $r$ is linear, with a constant gradient $\nabla r$, or a quadratic whose curvatures are much smaller than the $h_i$, with noise of fixed variance in each of its eigendirections.
\end{assumption}

With a quadratic river the model is the noisy quadratic model\index{noisy quadratic model} of \cref{ch:noise}, its spectrum split into hills, whose curvature times the area of the run is large, and river directions, for which it is small. In a hill the iterate soon forgets its past and sits in a noise-driven equilibrium; along the river it never forgets, and it is still descending when the run ends.

\begin{dictionary}
step size $\eta_t$, peak $\eta$, shape $\varphi$ & the schedule $\eta_{\max}\varphi(t/T)$\\
area $A_t=\sum_{s<t}\eta_s$ & the forward area $S_1$ of the empirical laws; how far the run has travelled along the river\\
river $r(x)$ & the part of the loss that keeps falling at a constant learning rate, roughly a power of the area\\
hill $y_i$, curvature $h_i$ & a sharp direction of the Hessian, along which the iterate oscillates at a high learning rate\\
noise variance $\sigma_i^2$ & gradient noise along that direction, $\Sigma_{ii}/B$ at batch size $B$\\
noise floor $\eta\sigma_i^2/(2(2-\eta h_i))$ & the excess loss of the stable phase, removed by the cooldown\\
relaxation time $1/(2\eta h_i)$ steps & how long the loss takes to respond to a change of the learning rate\\
annealing term $\sum_k(\eta_{k-1}-\eta_k)\,\schedulesG(A_T-A_k)$ & the loss drop $\mathrm{LD}(t)$ of the empirical laws\\
$D$ and $G^2$ & the distance to a good solution; the mean squared gradient
\end{dictionary}

\paragraph{What was thrown away.} The convex picture throws away curvature, and with it every distinction between fast and slow directions, as well as nonconvexity and every optimizer but plain stochastic gradient descent; it comes first because in it a statement about every schedule and every loss of a class can be proved. The river valley throws away the coupling of river and hills (in the picture of \textcite{WenEtAl2025} the river bends and the walls turn with it), changes of curvature, noise that depends on the position, and momentum, Adam and weight decay. It keeps what the convex picture lacks, two time scales, and it is solvable exactly.

\paragraph{What one had to notice.} The loss at the end of a run is a functional of the whole schedule, and in both pictures it depends on the schedule through two features with very different memories: the area, which remembers everything, and an average of the step sizes over the last few relaxation times, which remembers almost nothing. Large steps buy area; small steps at the end erase the short memory. A schedule is a compromise between two terms, and the shapes in use are different compromises between the same two terms.

\section{The theory: the last iterate under an arbitrary schedule}\label{sec:schedules:convex}

The classical analysis bounds the average iterate: by the online-gradient-descent bound of \cref{thm:toolkit:ogd} and Jensen's inequality, the $\eta$-weighted average of the iterates is within $(D^2+G^2\sum_t\eta_t^2)/(2\sum_t\eta_t)$ of the optimum, which is $DG/\sqrt T$ for the constant step $D/(G\sqrt T)$ \parencite{NemirovskiEtAl2009}. Practitioners use the \term{last iterate}\index{last iterate}.\sidenote{Averaging in stochastic approximation goes back to \textcite{PolyakJuditsky1992}; in deep learning it survives as the moving average of the weights and in schedule-free methods (\cref{ch:momentum}).} The tool for $\theta_T$ is the same expansion with the minimizer replaced by an earlier iterate.\sidenote{The proofs use one more stochastic gradient $g_T$ at $\theta_T$ and the point $\theta_{T+1}$. Neither affects $\theta_1,\dots,\theta_T$, so the step size $\eta_T$ is a free parameter of the bounds.}

\begin{lemma}[Anchored regret]\label[lemma]{lem:schedules:anchor}
Let \cref{ass:schedules:convex} hold and let $\eta_1,\dots,\eta_T>0$. For each $k\in\{1,\dots,T\}$ and for $u=\theta^\star$ or $u=\theta_k$,
\[
  \sum_{t=k}^{T}\eta_t\,\E\bigl[\Loss(\theta_t)-\Loss(u)\bigr]\le\frac12\,\E\norm{\theta_k-u}_2^2+\frac{G^2}2\sum_{t=k}^{T}\eta_t^2 .
\]
\end{lemma}

\begin{proof}
The plan is to expand the distance of $\theta_{t+1}$ to $u$ and telescope from $t=k$ on. Let $s_t=\E[g_t\mid\cF_t]\in\partial\Loss(\theta_t)$; by Jensen's inequality $\norm{s_t}_2\le G$. All expectations are finite: $(\E\norm{\theta_t}_2^2)^{1/2}\le\norm{\theta_1}_2+G\sum_{s<t}\eta_s$ by Minkowski's inequality, and $0\le\Loss(\theta_t)-\Loss^\star\le\Loss(\theta_1)-\Loss^\star+G\norm{\theta_t-\theta_1}_2$ by the subgradient inequality at $\theta_t$. By \eqref{eq:schedules:sgd},
\[
  \norm{\theta_{t+1}-u}_2^2=\norm{\theta_t-u}_2^2-2\eta_t\ip{g_t,\theta_t-u}+\eta_t^2\norm{g_t}_2^2 .
\]
For $t\ge k$ both $\theta_t$ and $u$ are $\cF_t$-measurable, so conditioning on $\cF_t$ replaces $g_t$ by $s_t$ in the middle term, and $\ip{s_t,\theta_t-u}\ge\Loss(\theta_t)-\Loss(u)$. Hence $2\eta_t\E[\Loss(\theta_t)-\Loss(u)]\le\E\norm{\theta_t-u}_2^2-\E\norm{\theta_{t+1}-u}_2^2+\eta_t^2G^2$. Summing over $t=k,\dots,T$ and dropping $-\E\norm{\theta_{T+1}-u}_2^2\le0$ gives the claim.
\end{proof}

\begin{theorem}[Last iterate under an arbitrary schedule]\label{thm:schedules:lastiterate}
Let \cref{ass:schedules:convex} hold, let $\eta_1,\dots,\eta_T>0$, and let $\theta_1,\dots,\theta_T$ be generated by \eqref{eq:schedules:sgd}. Write $R_k=\sum_{t=k}^{T}\eta_t$ and $V_k=\sum_{t=k}^{T}\eta_t^2$. Then
\begin{equation}\label{eq:schedules:lastiterate}
  \E\Loss(\theta_T)-\Loss^\star\ \le\ \frac{D^2+G^2V_1}{2R_1}\;+\;\frac{G^2}{2}\sum_{k=1}^{T-1}\frac{\eta_k}{R_{k+1}}\cdot\frac{V_k}{R_k} .
\end{equation}
\end{theorem}

\begin{proof}
The plan is to compare, for each $k$, the $\eta$-weighted mean of the expected losses from step $k$ on with the loss at step $k$, using $\theta_k$ as the comparator in \cref{lem:schedules:anchor}, and to telescope from the mean over the whole run down to the last iterate. Write $\ell_t=\E\Loss(\theta_t)$ and $\bar\ell_k=R_k^{-1}\sum_{t=k}^{T}\eta_t\ell_t$, so that $\bar\ell_T=\ell_T$.

\emph{The whole run.} \Cref{lem:schedules:anchor} with $k=1$ and $u=\theta^\star$ gives $R_1(\bar\ell_1-\Loss^\star)\le\frac12(D^2+G^2V_1)$.

\emph{One step of the telescope.} Let $1\le k\le T-1$. \Cref{lem:schedules:anchor} with $u=\theta_k$ gives $R_k(\bar\ell_k-\ell_k)\le\frac12G^2V_k$. By definition $R_k\bar\ell_k=\eta_k\ell_k+R_{k+1}\bar\ell_{k+1}$; subtracting $R_{k+1}\bar\ell_k=(R_k-\eta_k)\bar\ell_k$ from both sides,
\[
  R_{k+1}\bigl(\bar\ell_{k+1}-\bar\ell_k\bigr)=\eta_k\bigl(\bar\ell_k-\ell_k\bigr)\le\frac{G^2}2\,\frac{\eta_kV_k}{R_k} .
\]
\emph{Summation.} Divide by $R_{k+1}>0$ and add the identity $\ell_T-\Loss^\star=(\bar\ell_1-\Loss^\star)+\sum_{k=1}^{T-1}(\bar\ell_{k+1}-\bar\ell_k)$.
\end{proof}

The first term of \eqref{eq:schedules:lastiterate} is the averaging bound. The second, the price of the last iterate, averages the step sizes near the end: $V_k/R_k$ is the mean step size from step $k$ on, and its weight $\eta_k/R_{k+1}$, the current step over the area still to come, is negligible early and of order one at the end. The two terms are the convex counterparts of the river and the hills of \cref{sec:schedules:river}: one decreases with the area, the other remembers the recent schedule. Constant steps make the weights $1/(T-k)$, whose harmonic sum costs a logarithm; a decay to zero makes the step sizes small where the weights are large.\sidenote{The comparison with an earlier iterate is the argument of \textcite{ShamirZhang2013}; \textcite{Orabona2020} gave a version without a bounded domain, which we follow. \Textcite{DefazioEtAl2023} prove a last-iterate bound for general schedules by a reduction from regret, and \textcite{SchaippEtAl2025} evaluate such a bound along training runs.}

For linear decay the theorem gives about $2DG/\sqrt T$ (\cref{cor:schedules:shapes}), twice the averaging bound. A different argument closes the gap.

\begin{proposition}[Linear decay matches averaging]\label[proposition]{prop:schedules:linear}\index{linear decay}
Let \cref{ass:schedules:convex} hold, let $\gamma>0$, and run \eqref{eq:schedules:sgd} with $\eta_t=\gamma\,(T-t)/T$ for $t=1,\dots,T-1$. Then
\[
  \E\Loss(\theta_T)-\Loss^\star\le\frac{D^2}{2\gamma T}+\frac{\gamma G^2}2 ,
\]
which is $DG/\sqrt T$ for $\gamma=D/(G\sqrt T)$, the bound for the average iterate of stochastic gradient descent with the constant step $\gamma$.
\end{proposition}

\begin{proof}
The plan is to write the last iterate as the average of a sequence that takes constant steps along the same stochastic gradients, and to show by convexity that the regret of that sequence controls the last iterate. Let $z_1=\theta_1$ and $z_{t+1}=z_t-\gamma g_t$ for $t=1,\dots,T$, with $g_t$ the stochastic gradient at $\theta_t$.

\emph{The last iterate is an average.} We claim that $T\theta_t=\sum_{s<t}z_s+(T-t+1)z_t$ for $1\le t\le T$. It holds for $t=1$, and if it holds for $t$, the claimed values at $t+1$ and $t$ differ by $T^{-1}\bigl(z_t+(T-t)z_{t+1}-(T-t+1)z_t\bigr)=\frac{T-t}T(z_{t+1}-z_t)=-\eta_tg_t$, which is the step \eqref{eq:schedules:sgd}. In particular $\theta_T=T^{-1}\sum_{s=1}^Tz_s$.

\emph{A convexity inequality.} Let $p_t=t^{-1}\sum_{s\le t}z_s$, so that $p_1=\theta_1$ and $p_T=\theta_T$. For $2\le t\le T$ the claim reads $T\theta_t=(t-1)p_{t-1}+(T-t+1)z_t$; eliminating $z_t$ with $tp_t=(t-1)p_{t-1}+z_t$ gives
\[
  z_t-\theta_t=\frac{t-1}{T-t+1}\,(\theta_t-p_{t-1}),\qquad p_t=\frac{(t-1)(T-t)\,p_{t-1}+T\theta_t}{t\,(T-t+1)} .
\]
The second identity makes $p_t$ a convex combination of $p_{t-1}$ and $\theta_t$, so with $F_1=\Loss(\theta_1)$ and $F_t=\bigl((t-1)(T-t)F_{t-1}+T\Loss(\theta_t)\bigr)/(t(T-t+1))$, convexity and induction give $\Loss(p_t)\le F_t$; note that $F_T=\Loss(\theta_T)$. Let $s_t\in\partial\Loss(\theta_t)$ and $u\in\R^d$. The first identity and the subgradient inequality at $\theta_t$, applied to $u$ and to $p_{t-1}$, give
\[
  \ip{s_t,z_t-u}\ge\Loss(\theta_t)-\Loss(u)+\frac{t-1}{T-t+1}\bigl(\Loss(\theta_t)-F_{t-1}\bigr)=tF_t-(t-1)F_{t-1}-\Loss(u),
\]
the equality being the definition of $F_t$ rearranged (for $t=1$ read $(t-1)F_{t-1}=0$). The right side telescopes, so $\sum_{t=1}^{T}\ip{s_t,z_t-u}\ge T\bigl(\Loss(\theta_T)-\Loss(u)\bigr)$.

\emph{Conclusion.} Take $u=\theta^\star$ and $s_t=\E[g_t\mid\cF_t]$. Since $z_t$ is $\cF_t$-measurable, $\E\ip{g_t,z_t-\theta^\star}=\E\ip{s_t,z_t-\theta^\star}$. Since $z_{t+1}=z_t-\gamma g_t$, the expansion in the proof of \cref{lem:schedules:anchor}, with $z_t$ in place of $\theta_t$ (it uses only the form of the step, not where $g_t$ was computed), gives $\sum_{t=1}^T\E\ip{g_t,z_t-\theta^\star}\le D^2/(2\gamma)+\gamma TG^2/2$. Divide by $T$.
\end{proof}

The gradient $g_t$ enters $\theta_T$ with the weight $\gamma(T-t)/T$, exactly its total weight in the average of $z_1,\dots,z_T$: the last iterate of linear decay \emph{is} an average, computed without storing one. \Textcite{DefazioEtAl2023} derive the schedule from a conversion of this kind with general weights (Exercise~\ref{ex:schedules:weights}).

\begin{corollary}[Constant steps and linear decay]\label[corollary]{cor:schedules:shapes}
Under the hypotheses of \cref{thm:schedules:lastiterate}, let $H_n=\sum_{j=1}^n1/j$.
\begin{enumerate}
\item For $\eta_t=\eta$ the right side of \eqref{eq:schedules:lastiterate} is $D^2/(2\eta T)+\frac12G^2\eta\,(1+H_{T-1})$. Its minimum over $\eta$ is $DG\sqrt{(1+H_{T-1})/T}$; the step $\eta=D/(G\sqrt T)$ gives $(DG/\sqrt T)(1+\frac12H_{T-1})$.
\item For $\eta_t=\eta\,(T+1-t)/(T+1)$ it is $D^2/(\eta T)+G^2\eta\,(T-\frac12+H_{T-1})/(T+1)$, with minimum $(2DG/\sqrt T)(1+O(T^{-1}\log T))$.
\end{enumerate}
\end{corollary}

\begin{proof}
(i) Here $V_k/R_k=\eta$ and $\eta_k/R_{k+1}=1/(T-k)$, so the sum is $\eta H_{T-1}$, and the first term is $D^2/(2\eta T)+G^2\eta/2$; the minimum of $a/\eta+b\eta$ is $2\sqrt{ab}$. (ii) Put $m=T+1-k$ and $\delta=\eta/(T+1)$. Then $\eta_k=\delta m$, $R_k=\delta m(m+1)/2$, $V_k=\delta^2m(m+1)(2m+1)/6$ and $R_{k+1}=\delta(m-1)m/2$, so the $k$th term of the sum is $\frac{2\delta}{3}\cdot\frac{2m+1}{m-1}$, and $\sum_{m=2}^{T}\frac{2m+1}{m-1}=2(T-1)+3H_{T-1}$. With $R_1=\eta T/2$ and $V_1=\eta^2T(2T+1)/(6(T+1))$, collecting terms gives the formula.
\end{proof}

With the best constant step the bound exceeds the averaging bound by the factor $\sqrt{1+H_{T-1}}\approx\sqrt{\ln T}$, and with the step $D/(G\sqrt T)$ by about $\frac12\ln T$. Its constant is not sharp: \textcite{ZamaniGlineur2025} computed the exact worst case of the last iterate of the subgradient method with constant steps, and for the best constant step it is at most $DG\sqrt{1+\log(N)/4}/\sqrt{N+1}$ after $N$ steps.\sidenote{\Textcite{ZamaniGlineur2025} also show that linear decay attains the lower bound $DG/\sqrt{N+1}$ that holds for all first-order methods, and that no single sequence of step sizes attains it at every $N$: the optimal schedule must know its horizon, which is \cref{obs:schedules:horizon} in the worst case. For $\eta_t\propto1/\sqrt t$, which does not, the last iterate loses a logarithmic factor that cannot be removed \parencite{HarveyEtAl2019}.} For other shapes the sums in \eqref{eq:schedules:lastiterate} are Riemann sums. For a shape $\varphi$ of bounded variation on $[0,1]$ and positive on $[0,1)$, put $\Phi(s)=\int_s^1\varphi$, $\Psi(s)=\int_s^1\varphi^2$ and
\begin{equation}\label{eq:schedules:J}
  J[\varphi]=\frac{\Psi(0)}{\Phi(0)}+\int_0^1\frac{\varphi(s)\,\Psi(s)}{\Phi(s)^2}\,ds .
\end{equation}

\begin{proposition}[The bound for a shape]\label[proposition]{prop:schedules:shape}
Let $\varphi$ be a shape as above such that $\varphi(s)/(1-s)^p$ has a positive limit as $s\to1$, for some $p>0$. Let $\eta_t=\eta\,\varphi(t/(T+1))$ for $t\le T$, with $\eta=\gamma D/(G\sqrt T)$ for a fixed $\gamma>0$. Then $J[\varphi]<\infty$; as $T\to\infty$ the right side of \eqref{eq:schedules:lastiterate} is $(DG/\sqrt T)\bigl(1/(2\gamma\Phi(0))+\gamma J[\varphi]/2+o(1)\bigr)$, and its minimum over $\gamma$ is $DG\sqrt{J[\varphi]/(\Phi(0)T)}\,(1+o(1))$.
\end{proposition}

\begin{proofsketch}
Near $s=1$ the hypothesis gives $\Phi\asymp(1-s)^{p+1}$, $\Psi\asymp(1-s)^{2p+1}$ and $\varphi\Psi/\Phi^2\asymp(1-s)^{p-1}$, which is integrable, so $J[\varphi]<\infty$. With $s=k/(T+1)$, the sums $R_k$ and $V_k$ are Riemann sums of $\eta T\Phi(s)$ and $\eta^2T\Psi(s)$, so for $1-s\gg1/T$ the $k$th term of the sum in \eqref{eq:schedules:lastiterate} is $(\eta/T)\,\varphi\Psi/\Phi^2$ at $s$ up to a factor $1+o(1)$: the sum is a Riemann sum of the integral in \eqref{eq:schedules:J}. The power law at $s=1$ makes the terms with $1-s\le\eps$ contribute $O(\eta\eps^p)$ in total, uniformly in $T$; we omit this estimate, which for the shapes of \cref{tab:schedules:bound} can be done by exact summation as in \cref{cor:schedules:shapes}(ii). The hypothesis cannot be dropped: if $\varphi$ falls faster than every power of $1-s$, as $e^{-0.02/(1-s)}$ does, the last term alone is close to $\eta$, and the bound converges to the formula with $J[\varphi]+1$ in place of $J[\varphi]$. At $T=10^5$ the discrete bound agrees with the limit within $0.02\%$ for the shapes of \cref{tab:schedules:bound}.
\end{proofsketch}

\begin{table}[t]
\begin{fullwidth}
\centering
\small
\begin{tabular}{@{}lccl@{}}
\toprule
shape & $\Phi(0)$ & $J[\varphi]$ & best bound $\times\sqrt T/(DG)$\\
\midrule
constant & $1$ & $1+H_{T-1}$ (exact) & $\sqrt{1+H_{T-1}}$, which is $3.28$ at $T=10^4$\\
linear decay to zero & $1/2$ & $2$ & $2$\\
cosine to zero & $1/2$ & $2.122$ & $2.060$\\
cosine to $\rho=0.1$ & $0.55$ & $1.775+0.1\ln T$ & $2.21$ at $T=10^4$; $2.31$ at $T=10^5$\\
WSD, linear cooldown of fraction $c$ & $1-c/2$ & $2+\ln\frac{2-c}c$ & $2.28$, $2.16$, $2.03$ for $c=0.1$, $0.2$, $0.5$\\
WSD, $1-\sqrt\cdot$ cooldown, $c=0.2$ & $0.867$ & $3.92$ & $2.13$\\
\bottomrule
\end{tabular}
\caption{The last-iterate bound \eqref{eq:schedules:lastiterate} at its best peak (\cref{prop:schedules:shape}); the constant and floor rows are exact discrete values, whose $J$ grows with $T$. A floor $\rho$ adds $\rho\ln T$ to $J$ (Exercise~\ref{ex:schedules:floor}).}\label{tab:schedules:bound}
\end{fullwidth}
\end{table}

\begin{corollary}[The cooldown in the bound]\label[corollary]{cor:schedules:cooldown}
For WSD with peak $\eta$ and a linear cooldown to zero over the last fraction $c\in(0,1]$ of the steps, the right side of \eqref{eq:schedules:lastiterate} is, as $T\to\infty$ with $\eta\sqrt T$ fixed, $U(c)\,(1+o(1))$ with
\begin{equation}\label{eq:schedules:wsdbound}
  U(c)=\frac{D^2}{(2-c)\,\eta T}+\frac{G^2\eta}2\Bigl(2+\ln\frac{2-c}{c}\Bigr).
\end{equation}
\begin{enumerate}
\item With the peak tuned for each $c$, the bound is
\[
  \min_{\eta>0}U(c)=DG\,\sqrt{\frac{2\,\bigl(2+\ln\frac{2-c}{c}\bigr)}{(2-c)\,T}} ,
\]
which decreases in $c$: the bound prefers to decay over the whole run.
\item With the peak fixed, let $\mu=G\eta\sqrt T/D$ be the peak in units of the best peak for linear decay. Then $U$ is minimized at $c^\star=\min\{1,\,2\mu^2/(1+\mu^2)\}$: a peak below the best one asks for a shorter cooldown.
\end{enumerate}
\end{corollary}

\begin{proof}
The plan is to compute $J$ for this shape and apply \cref{prop:schedules:shape}. Here $\Phi(0)=1-c/2$ and $\Psi(0)=1-2c/3$. On the cooldown $\Phi(s)=(1-s)^2/(2c)$ and $\Psi(s)=(1-s)^3/(3c^2)$, so the integrand of \eqref{eq:schedules:J} is $4/(3c)$ and contributes $4/3$. On the stable phase, with $w=1-s-c/2$, the integrand is $(w-c/6)/w^2$, with integral $\ln\frac{2-c}c-\frac13+\frac{c}{3(2-c)}$ over $w\in[c/2,1-c/2]$. Adding $\Psi(0)/\Phi(0)=\frac{2(3-2c)}{3(2-c)}$ gives $J=2+\ln\frac{2-c}c$.
(i) The tuned value is $DG\sqrt{2F(c)/T}$ with $F(c)=(2+\ln r)/(2-c)$ and $r=(2-c)/c\ge1$, and $F'(c)=(1-r+\ln r)/(2-c)^2\le0$ because $\ln r\le r-1$.
(ii) The derivative $U'(c)=D^2/((2-c)^2\eta T)-G^2\eta/(c(2-c))$, multiplied by $c(2-c)^2\eta T/D^2>0$, is $c(1+\mu^2)-2\mu^2$, which changes sign from negative to positive at $2\mu^2/(1+\mu^2)$; if that exceeds $1$, then $U$ decreases on $(0,1]$.
\end{proof}

So the bound reproduces \cref{obs:schedules:cooldown} in a weak form: a cooldown removes the logarithm of the constant schedule, and its benefit saturates. It does not single out $20\%$; with tuned peaks it prefers a decay over the whole run, as in the experiments of \textcite{DefazioEtAl2023} and \textcite{BergsmaEtAl2025}, and it prefers the $1-\sqrt\cdot$ cooldown of \textcite{HageleEtAl2024} to a linear one.

\section{The theory: the river valley}\label{sec:schedules:river}

In the river valley the expected loss can be computed exactly for every schedule. For a hill $i$ write $q_{i,t}=(1-\eta_th_i)^2$ and $P_i(k,T)=\prod_{t=k}^{T-1}q_{i,t}$ (an empty product is $1$), the contraction of its second moment between steps $k$ and $T$. When $0\le\eta h_i<2$, let
\[
  v_i(\eta)=\frac{\eta\,\sigma_i^2}{h_i\,(2-\eta h_i)} ,
\]
the stationary second moment of the hill at the constant step $\eta$ (\cref{lem:toolkit:lyapunov}).

\begin{theorem}[Stochastic gradient descent in the river valley]\label{thm:schedules:river}
Let \cref{ass:schedules:river} hold and let $(\eta_t)$ be a schedule. For each hill $i$ and each $T\ge1$:
\begin{enumerate}
\item $\E y_{T,i}^2=P_i(1,T)\,y_{1,i}^2+\sigma_i^2\sum_{k=1}^{T-1}\eta_k^2\,P_i(k+1,T)$.
\item If $0\le\eta_th_i<2$ for $t<T$, then
\begin{multline}\label{eq:schedules:annealing}
  \E y_{T,i}^2=v_i(\eta_1)+P_i(1,T)\bigl(y_{1,i}^2-v_i(\eta_1)\bigr)\\
  -\sum_{k=2}^{T-1}\bigl(v_i(\eta_{k-1})-v_i(\eta_k)\bigr)\bigl(1-P_i(k,T)\bigr).
\end{multline}
\item If $0\le\eta_th_i\le\frac12$ for $k\le t<T$, then $e^{-2h_i(A_T-A_k)-2h_i^2\sum_{t=k}^{T-1}\eta_t^2}\le P_i(k,T)\le e^{-2h_i(A_T-A_k)}$.
\item If the river is linear, then $\E r(x_T)=r(x_1)-\norm{\nabla r}_2^2\,A_T$. If it is quadratic, each of its eigendirections satisfies (i)--(iii) with its own curvature and noise.
\end{enumerate}
Consequently $\E\Loss(\theta_T)=\E r(x_T)+\frac12\sum_ih_i\,\E y_{T,i}^2$ is given in closed form by the schedule.
\end{theorem}

\begin{proof}
The plan is to derive the second-moment recursion of each coordinate and to rewrite it around the moving equilibrium $v_i(\eta_t)$. (i) The hill coordinate moves by $y_{t+1,i}=(1-\eta_th_i)y_{t,i}-\eta_t\zeta_{t,i}$. Square it and condition on $\cF_t$; the cross term vanishes because $y_{t,i}$ is $\cF_t$-measurable and $\E[\zeta_{t,i}\mid\cF_t]=0$, so $\E[y_{t+1,i}^2\mid\cF_t]=q_{i,t}y_{t,i}^2+\eta_t^2\sigma_i^2$. Take expectations and unroll from $t=1$.
(ii) Write $u_t=\E y_{t,i}^2$ and $v_t=v_i(\eta_t)$; the step size $\eta_T$ does not affect $y_{T,i}$, so we may take $\eta_T=\eta_{T-1}$ and $v_T$ is defined. A direct computation gives $q_{i,t}v_t+\eta_t^2\sigma_i^2=v_t$, so subtracting it from the recursion of (i) gives $u_{t+1}-v_t=q_{i,t}(u_t-v_t)$, that is, $u_{t+1}-v_{t+1}=q_{i,t}(u_t-v_t)+(v_t-v_{t+1})$. Unrolling, $u_T-v_T=P_i(1,T)(u_1-v_1)+\sum_{k=2}^{T}(v_{k-1}-v_k)P_i(k,T)$. The term $k=T$ has $P_i(T,T)=1$ and cancels the unused $v_T$, and $v_{T-1}=v_1-\sum_{k=2}^{T-1}(v_{k-1}-v_k)$ gives \eqref{eq:schedules:annealing}.
(iii) For $0\le z\le\frac12$, $e^{-z-z^2}\le1-z\le e^{-z}$: the right inequality is convexity of the exponential, and $\ln(1-z)+z+z^2$ vanishes at $0$ and has derivative $z(1-2z)/(1-z)\ge0$. Apply it with $z=\eta_th_i$, square, and multiply.
(iv) A linear river has $x_{t+1}=x_t-\eta_t(\nabla r+\zeta_{t,x})$, so $\E x_T=x_1-A_T\nabla r$, and $r$ is affine. The eigendirections of a quadratic river are decoupled coordinates of the same form as the hills.
\end{proof}

Part (ii), the \term{annealing decomposition}\index{annealing decomposition}, is the structure behind the chapter: the second moment of a hill at step $T$ is its equilibrium at the initial step size, plus the fading memory of the start, minus one term for each decrease of the step size, the decrease of the equilibrium times the fraction $1-P_i(k,T)$ of it realized since. By (iii) that fraction depends on the area since the decrease, $1-e^{-2h_i(A_T-A_k)}$ up to relative corrections of order $\eta h_i$; an increase of the step size, as in warmup, enters with the opposite sign.

\begin{corollary}[Noise floor and relaxation]\label[corollary]{cor:schedules:floor}
Under \cref{ass:schedules:river}, let $0<\eta h_i<2$ for every hill.
\begin{enumerate}
\item At the constant step $\eta$ the hill part of the expected loss converges, at the rate $(1-\eta h_i)^2$ per step in hill $i$, to the \term{noise floor}\index{noise floor} $\Loss_{\mathrm{hill}}(\eta)=\sum_{i}\eta\sigma_i^2/(2(2-\eta h_i))$, which lies between $\frac\eta4\sum_i\sigma_i^2$ and $\eta\sum_i\sigma_i^2/(4-2\eta\max_ih_i)$.
\item If $\E y_{k,i}^2=v_i(\eta)$ and the step size is $\eta'\in[0,\eta)$ from step $k$ on, then $\E y_{T,i}^2=v_i(\eta')+\bigl(v_i(\eta)-v_i(\eta')\bigr)(1-\eta'h_i)^{2(T-k)}$ for $T\ge k$.
\end{enumerate}
\end{corollary}

\begin{proof}
(i) In \eqref{eq:schedules:annealing} all differences vanish and $P_i(1,T)=(1-\eta h_i)^{2(T-1)}$; multiply by $h_i/2$ and sum, and use $2-\eta\max_ih_i\le2-\eta h_i\le2$. (ii) The recursion $u_{t+1}-v_i(\eta')=(1-\eta'h_i)^2(u_t-v_i(\eta'))$ of the proof of \cref{thm:schedules:river}(ii) holds for $t\ge k$.
\end{proof}

So at a constant learning rate each hill adds about $\eta\sigma_i^2/4$ to the loss, a floor proportional to the learning rate and to the noise, hence to $1/B$. A drop of the learning rate removes the excess on the \term{relaxation time}\index{relaxation time} $1/(\eta'h_i(2-\eta'h_i))\approx1/(2\eta'h_i)$ of the \emph{new} step size, so a decay to zero never quite finishes.

\begin{heuristic}[The cooldown law]\index{cooldown!law}
Take one hill in equilibrium at the peak $\eta$, and a linear cooldown to zero over the last $K$ steps. Measure time by the area: to first order in $\eta h$ the recursion of \cref{thm:schedules:river}(i) is $dv/dA=-2hv+\eta(A)\sigma^2$, because a step changes $v$ by $-(2\eta_th-\eta_t^2h^2)v+\eta_t^2\sigma^2$ and the area by $\eta_t$. The cooldown has area $\eta K/2$, and with $\varrho$ the area that remains before the end the step size is $\eta(\varrho)=\sqrt{2\eta\varrho/K}$. Solving the linear equation from $v=\eta\sigma^2/(2h)$,
\[
  v_{\mathrm{end}}=\frac{\eta\sigma^2}{2h}e^{-\eta hK}+\sigma^2\int_0^{\eta K/2}e^{-2h\varrho}\sqrt{2\eta\varrho/K}\,d\varrho ,
\]
an incomplete gamma function. With $\kappa=\eta hK$ the two terms combine into
\begin{equation}\label{eq:schedules:cooldownlaw}
  \frac h2\,v_{\mathrm{end}}=\frac{\eta\sigma^2}4\,\chi(\kappa),\qquad \chi(\kappa)=\frac{\sqrt\pi}2\,\frac{\erf\sqrt\kappa}{\sqrt\kappa} .
\end{equation}
A short cooldown has $\chi(\kappa)=1-\kappa/3+O(\kappa^2)$ and removes almost nothing; a long one leaves $\chi(\kappa)\approx\frac{\sqrt\pi}2\kappa^{-1/2}$. The residual noise after a linear cooldown of $K$ steps decays like $K^{-1/2}$, not exponentially: the noise injected within the last area $1/(2h)$, at step sizes of order $(\eta/(hK))^{1/2}$, has no area left in which to be forgotten.\qedhere
\end{heuristic}

The exact recursion confirms \eqref{eq:schedules:cooldownlaw} within relative errors of order $\eta h$: started at the equilibrium $v(\eta)$, with the steps $\eta(1-j/K)$ for $j=1,\dots,K$, for $\eta h=0.01$ and $K=100$, $1000$, $10^4$ it leaves $0.746$, $0.279$ and $0.0882$ of the floor, against $\chi=0.747$, $0.280$ and $0.0886$. The cooldown that removes the most noise in $K$ steps, ignoring the river, is a $1/t$ decay, with a residual of order $1/K$ (Exercise~\ref{ex:schedules:greedy}).\sidenote{The error function enters as $\int_0^x\sqrt u\,e^{-u}du=\frac{\sqrt\pi}2\erf\sqrt x-\sqrt x\,e^{-x}$. A cooldown whose step size vanishes like the $p$th power of the remaining steps leaves a tail $\kappa^{-p/(p+1)}$: $\kappa^{-1/2}$ for the linear and $1-\sqrt\cdot$ shapes, $\kappa^{-2/3}$ for a cosine.}

\begin{heuristic}[The best cooldown]\index{cooldown!optimal length}
Let the river be linear, so that by \cref{thm:schedules:river}(iv) it loses $\omega=\norm{\nabla r}_2^2$ per unit of area, and replace the last $K$ steps of a constant schedule by a linear cooldown. The cooldown gives up the area $\eta K/2$ and removes noise by \eqref{eq:schedules:cooldownlaw}, so the final loss is, up to a constant, $\omega\eta K/2+\sum_i\frac14\eta\sigma_i^2\chi(\eta h_iK)$. For one hill with $\eta hK\gg1$ the derivative vanishes at
\begin{equation}\label{eq:schedules:Kstar}
  K^\star=\Bigl(\frac{\sqrt\pi\,\sigma^2}{8\,\omega}\Bigr)^{2/3}(\eta h)^{-1/3},
\end{equation}
where the residual noise is twice the river progress given up. The best cooldown is a fixed number of steps, and its \emph{fraction} of the run falls like $1/T$. A river that flattens changes this. If the river loss is a power of the area, $\bar r(A)=L_\infty+CA^{-\alpha}$, then its slope at the end of a run of area $\eta T$ is $\alpha C(\eta T)^{-1-\alpha}$, and the same balance, with $1-c/2\approx1$, gives the best cooldown fraction
\begin{equation}\label{eq:schedules:cstar}
  c^\star\approx\Bigl(\frac{\sqrt\pi}{8\alpha C}\sum_i\frac{\sigma_i^2}{\sqrt{h_i}}\Bigr)^{2/3}\,\eta^{(1+2\alpha)/3}\,T^{(2\alpha-1)/3} .
\end{equation}
A fixed fraction is optimal at every horizon exactly when $\alpha=\frac12$; a river that falls faster than the inverse square root of the area asks for a longer cooldown in a longer run, a slower one for a shorter cooldown.\qedhere
\end{heuristic}

For one hill with $\eta h=0.01$, $\sigma^2=1$ and $\omega=10^{-4}$, the exact recursion puts the best linear cooldown at $790$ steps, \eqref{eq:schedules:Kstar} at $789$; the asymptotic form fails, as it must, when $\eta hK^\star$ is of order one. The decomposition also explains the laws of \cref{obs:schedules:laws}.

\begin{heuristic}[From the river valley to the empirical laws]
Multiply \eqref{eq:schedules:annealing} by $h_i/2$, sum over the hills, and keep the first order in $\eta h_i$, where $\frac{h_i}2v_i(\eta)\approx\frac14\eta\sigma_i^2$ and $1-P_i(k,T)\approx1-e^{-2h_i(A_T-A_k)}$. Let $\schedulesG(u)=\sum_i\frac14\sigma_i^2\bigl(1-e^{-2h_iu}\bigr)$, a saturation function built from the Laplace transform of the noise-weighted hill spectrum. Once the initial transient has faded,
\begin{equation}\label{eq:schedules:law}
  \E\Loss(\theta_T)\approx\E r(x_T)+\frac{\eta_1}4\sum_i\sigma_i^2-\sum_{k=2}^{T-1}(\eta_{k-1}-\eta_k)\,\schedulesG(A_T-A_k).
\end{equation}
This is the form of \cref{obs:schedules:laws}: an area term, close to $L_\infty+CA_T^{-\alpha}$ for a quadratic river with a power-law spectrum, and a loss drop that adds, for each decrease of the learning rate, the decrease times a saturating function of the area since. If the noise is spread over many hills and its distribution over their curvatures, with weights $\sigma_i^2$, is a gamma distribution with shape $\nu$ and scale $\bar h$, then $\schedulesG(u)=\frac14\bigl(\sum_i\sigma_i^2\bigr)\bigl(1-(1+2\bar hu)^{-\nu}\bigr)$, the power-law saturation of \textcite{LuoEtAl2025}. A single hill at a nearly constant step $\bar\eta$ after the decreases gives the saturation $1-e^{-2\bar\eta h(T-k)}$, counted in steps, which is the annealing area of \textcite{TissueEtAl2025}. Finally, the continuum form of \cref{thm:schedules:river}(i) writes the noise term as $\sum_k\eta_k^2\schedulesK(A_T-A_{k+1})$ with the kernel $\schedulesK=\schedulesG'$, a convolution of the squared schedule with a kernel in area-time; \textcite{LiEtAl2025} derive a functional scaling law of this kind for stochastic gradient descent on kernel regression with power-law spectra.\qedhere
\end{heuristic}

The model predicts what the laws leave open: the loss relaxes after a decrease in \emph{area}, not in steps, because a quadratic hill contracts by $(1-\eta h)^2\approx e^{-2h\eta}$ per step, that is, by $e^{-2h}$ per unit of area. \Textcite{TissueEtAl2025} count the relaxation in steps, \textcite{LuoEtAl2025} in area with a fitted correction that depends on the learning rate, which measures a departure from a fixed quadratic landscape.

\begin{remark}[Branching, rewarming and averaging]\label[remark]{rem:schedules:wsd}
After a long stable phase the hills have forgotten everything but the peak, so by \eqref{eq:schedules:annealing} the loss after a cooldown depends on the stable phase only through its area (\cref{obs:schedules:horizon}); rewarming a cooled checkpoint refills the hills within a relaxation time and keeps the river progress. Averaging can replace the decay, as in schedule-free methods \parencite{DefazioEtAl2024} (\cref{ch:momentum}): the mean of the last $K$ iterates at a constant step gives up the same area as a linear cooldown of $K$ steps and leaves a residual of order $1/(\eta hK)$, not $(\eta hK)^{-1/2}$ (Exercise~\ref{ex:schedules:averaging}).
\end{remark}

\section{Reading the theory}\label{sec:schedules:reading}

\paragraph{Loss curves.} \Cref{fig:schedules:river} runs the four schedules in a river valley built to look like a language model: a quadratic river in $200$ directions whose loss falls like the inverse square root of the area, and $40$ noisy hills with curvatures between $0.2$ and $2$. The points, means of $100$ runs, agree with the exact curves within their sampling error (final values $0.203$, $0.156$, $0.158$, $0.131$ against $0.207$, $0.157$, $0.158$, $0.131$). The constant schedule reaches the floor $\Loss_{\mathrm{hill}}(0.04)=0.102$ within a few hundred steps and follows the river down at that height. In the cooldown of WSD the hill part falls from $0.102$ to $0.021$ while the river part continues from $0.117$ to $0.110$: four fifths of the floor removed for a tenth of the area. Cosine and linear decay remove almost all of the floor but cover half the area, and their river part ends at $0.149$. At a common peak the model ranks WSD first, cosine and linear decay together, the constant schedule last.

\begin{figure}[t]
\begin{fullwidth}
\centering
\begin{tikzpicture}
\begin{groupplot}[group style={group size=2 by 1,horizontal sep=17mm},wplot,width=82mm,height=48mm,xmin=0,xmax=4000,xtick={0,1000,2000,3000,4000},scaled x ticks=false,xticklabel style={/pgf/number format/1000 sep={}}]
\nextgroupplot[xlabel={step $t$},ylabel={expected loss},ymode=log,ymin=0.1,ymax=0.65,ytick={0.1,0.15,0.2,0.3,0.4,0.6},yticklabels={$0.1$,$0.15$,$0.2$,$0.3$,$0.4$,$0.6$},legend style={at={(0.5,1.03)},anchor=south},legend columns=4]
\addplot[theory] table[x=t,y=const] {figures/data/12-schedules-river.dat};
\addplot[theory3] table[x=t,y=cosine] {figures/data/12-schedules-river.dat};
\addplot[theory2] table[x=t,y=linear] {figures/data/12-schedules-river.dat};
\addplot[muted,line width=1.0pt,dashdotted] table[x=t,y=wsd] {figures/data/12-schedules-river.dat};
\addplot[only marks,mark=*,mark size=1.1pt,sea,each nth point=2,filter discard warning=false] table[x=t,y=const] {figures/data/12-schedules-river-sgd.dat};
\addplot[only marks,mark=*,mark size=1.1pt,sandline!55!black,each nth point=2,filter discard warning=false] table[x=t,y=cosine] {figures/data/12-schedules-river-sgd.dat};
\addplot[only marks,mark=square*,mark size=1.0pt,seadeep,each nth point=2,filter discard warning=false] table[x=t,y=linear] {figures/data/12-schedules-river-sgd.dat};
\addplot[only marks,mark=triangle*,mark size=1.3pt,muted,each nth point=2,filter discard warning=false] table[x=t,y=wsd] {figures/data/12-schedules-river-sgd.dat};
\legend{constant,cosine,linear,WSD}
\nextgroupplot[xlabel={step $t$},ylabel={river and hill parts},ymode=log,ymin=0.004,ymax=0.65,legend style={at={(0.5,1.03)},anchor=south},legend columns=2]
\addplot[muted,line width=1.0pt,dashdotted] table[x=t,y=wsdriver] {figures/data/12-schedules-river.dat};
\addplot[muted,line width=0.8pt,densely dotted] table[x=t,y=wsdhill] {figures/data/12-schedules-river.dat};
\addplot[theory3] table[x=t,y=cosriver] {figures/data/12-schedules-river.dat};
\addplot[sandline!55!black,line width=0.8pt,densely dotted] table[x=t,y=coshill] {figures/data/12-schedules-river.dat};
\addplot[only marks,mark=triangle*,mark size=1.3pt,muted,each nth point=2,filter discard warning=false] table[x=t,y=wsdriver] {figures/data/12-schedules-river-sgd.dat};
\addplot[only marks,mark=triangle*,mark size=1.3pt,muted,each nth point=2,filter discard warning=false] table[x=t,y=wsdhill] {figures/data/12-schedules-river-sgd.dat};
\addplot[only marks,mark=*,mark size=1.1pt,sandline!55!black,each nth point=2,filter discard warning=false] table[x=t,y=cosriver] {figures/data/12-schedules-river-sgd.dat};
\addplot[only marks,mark=*,mark size=1.1pt,sandline!55!black,each nth point=2,filter discard warning=false] table[x=t,y=coshill] {figures/data/12-schedules-river-sgd.dat};
\legend{WSD river,WSD hills,cosine river,cosine hills}
\end{groupplot}
\end{tikzpicture}
\end{fullwidth}
\caption{Four schedules in a river valley: curves are the exact expected loss (\cref{thm:schedules:river}), points the mean of $100$ runs of stochastic gradient descent with Gaussian noise (seed $12$), every $200$ steps. River: $200$ directions with curvatures $0.02\,j^{-2}$, initial losses proportional to $j^{-2}$ (total $0.5$), no noise. Hills: $40$, curvatures log-spaced in $[0.2,2]$, started at $0$, noise variance $1/4$ each. Peak $0.04$, $T=4000$, no warmup; WSD cools down linearly over the last $800$ steps. Left: total loss. Right: river part (upper curves) and hill part (lower curves). The points agree with the curves within sampling error (at most $2.7\%$ apart, less than three standard errors); the river carries no noise, so its points coincide with its curves. Script: \texttt{12-schedules-river.py}.}
\label{fig:schedules:river}
\end{figure}

That ranking is not a prediction about tuned runs. With each peak tuned separately among the peaks at which every step is stable, $\eta\max_ih_i<2$, the model raises the peaks of the decaying schedules far above the common one (cosine to that limit, linear decay to $0.64$, WSD to $0.21$) and ranks cosine first ($0.052$, against $0.074$ for linear decay, $0.096$ for WSD and $0.195$ for constant steps): a larger peak buys area, and a long decay removes the larger floor. Cosine even gains from peaks up to nearly twice that limit ($0.050$ at the peak $1.9$): the stiff hills grow during its first, unstable steps, and the contraction during the later steps undoes the growth, which only exactly quadratic hills allow. Nothing in the model penalizes a large peak short of divergence (\cref{ch:eos,ch:spikes}). What is robust are the decomposition and its time scales: the cooldown drop is the noise floor at the peak, about $\eta\tr(\Sigma_{\mathrm{hill}})/(4B)$, so it should double when the peak doubles or the batch size halves; it takes a few relaxation times of the slowest noisy direction; and it leaves a residual decaying only like $(\eta hK)^{-1/2}$. The curves are exact for this model, so their agreement with the points is guaranteed up to sampling error; it must fail once the curvature changes along the run, the noise depends on the position, or the optimizer is not plain stochastic gradient descent.

\paragraph{The convex bound along a run.} \Cref{fig:schedules:bound} (left) evaluates \eqref{eq:schedules:lastiterate} at every step on the \term{cone}\index{cone (test problem)} $\Loss(\theta)=\E_a\abs{\ip{a,\theta-\theta^\star}}=\sqrt{2/\pi}\,\norm{\theta-\theta^\star}_2$ with $a\sim\cN(0,I_{100})$, which is convex with a nonsmooth minimum and has $G^2=d$ exactly. The shapes agree with the simulation and with \cref{obs:schedules:convex}: constant steps plateau, WSD plateaus and then drops, cosine and linear decay descend smoothly, and both rank the final values linear decay, cosine, WSD, constant. The values do not agree: the bound is $7.6$ to $9.5$ times the simulated suboptimality at the end. On the cone the last iterate of constant steps settles exactly at the averaging floor $\eta G^2/2$, with no logarithm, and it forgets a change of the step size within about $130$ steps whatever the step size, so its final loss is set by the last steps. Tuned comparisons (right panel) suffer: the bound charges a $20\%$ cooldown $8\%$ over linear decay, the cone $68\%$; cosine $3\%$, the cone $31\%$. The bound covers every convex Lipschitz function, and the functions on which the last iterate comes within a constant factor of it are specially constructed \parencite{HarveyEtAl2019,ZamaniGlineur2025}; it agrees with training curves in shapes and rankings, not in values.

\begin{figure}[t]
\begin{fullwidth}
\centering
\begin{tikzpicture}
\begin{groupplot}[group style={group size=2 by 1,horizontal sep=17mm},wplot,width=82mm,height=48mm]
\nextgroupplot[xlabel={step $t$},ylabel={$\E\Loss(\theta_t)-\Loss^\star$},ymode=log,xmin=0,xmax=1000,ymin=0.04,ymax=4,legend style={at={(0.5,1.03)},anchor=south},legend columns=4,xticklabel style={/pgf/number format/1000 sep={}}]
\addplot[theory,restrict x to domain=10:1000] table[x=t,y=const] {figures/data/12-schedules-bound.dat};
\addplot[theory3,restrict x to domain=10:1000] table[x=t,y=cosine] {figures/data/12-schedules-bound.dat};
\addplot[theory2,restrict x to domain=10:1000] table[x=t,y=linear] {figures/data/12-schedules-bound.dat};
\addplot[muted,line width=1.0pt,dashdotted,restrict x to domain=10:1000] table[x=t,y=wsd] {figures/data/12-schedules-bound.dat};
\addplot[only marks,mark=*,mark size=1.0pt,sea,restrict x to domain=10:1000,each nth point=2,filter discard warning=false] table[x=t,y=const] {figures/data/12-schedules-bound-sgd.dat};
\addplot[only marks,mark=*,mark size=1.0pt,sandline!55!black,restrict x to domain=10:1000,each nth point=2,filter discard warning=false] table[x=t,y=cosine] {figures/data/12-schedules-bound-sgd.dat};
\addplot[only marks,mark=square*,mark size=0.9pt,seadeep,restrict x to domain=10:1000,each nth point=2,filter discard warning=false] table[x=t,y=linear] {figures/data/12-schedules-bound-sgd.dat};
\addplot[only marks,mark=triangle*,mark size=1.2pt,muted,restrict x to domain=10:1000,each nth point=2,filter discard warning=false] table[x=t,y=wsd] {figures/data/12-schedules-bound-sgd.dat};
\legend{constant,cosine,linear,WSD}
\nextgroupplot[xlabel={cooldown fraction $c$},ylabel={tuned final value / linear decay},xmode=log,xmin=0.02,xmax=1.05,ymin=0.9,ymax=2.2,xtick={0.02,0.05,0.1,0.2,0.5,1},xticklabels={$0.02$,$0.05$,$0.1$,$0.2$,$0.5$,$1$},legend style={at={(0.5,1.03)},anchor=south},legend columns=2]
\addplot[theory] table[x=c,y=boundrel] {figures/data/12-schedules-bound-tunedcurve.dat};
\addplot[simulation] table[x=c,y=sgdrel] {figures/data/12-schedules-bound-tuned.dat};
\addplot[guide,domain=0.02:1.05] {1};
\legend{bound \eqref{eq:schedules:lastiterate},SGD on the cone}
\end{groupplot}
\end{tikzpicture}
\end{fullwidth}
\caption{The last-iterate bound against stochastic gradient descent on the cone $\sqrt{2/\pi}\,\norm{\theta-\theta^\star}_2$ in $d=100$ dimensions, $D=1$, $G^2=d$. Left: the bound \eqref{eq:schedules:lastiterate} for $\theta_t$, computed from $\eta_1,\dots,\eta_t$ (curves), and the mean of $200$ runs (points, seed $21$), $T=1000$, common peak $D/(G\sqrt T)$; WSD cools down over the last $200$ steps. At $t=T$ the bound is $7.6$ to $9.5$ times the simulated value. Right: WSD with cooldown fraction $c$, peak tuned for each $c$, final value relative to linear decay; the bound is minimized exactly, the simulation takes the best of $17$ peaks between a quarter of and four times $D/(G\sqrt T)$, $40$ runs each (seed $22$). Script: \texttt{12-schedules-bound.py}.}
\label{fig:schedules:bound}
\end{figure}

\paragraph{How long to cool down.} \Cref{fig:schedules:cooldown} tests the cooldown heuristics in the valley of \cref{fig:schedules:river} and in one whose river falls like $A^{-1}$. For $\alpha=\frac12$ the best cooldown fraction is $0.24$, $0.28$ and $0.27$ at $T=1000$, $4000$ and $16000$, fixed as \eqref{eq:schedules:cstar} predicts; for $\alpha=1$ it is $0.21$, $0.35$ and $0.48$, a factor $2.3$ over a sixteenfold increase of $T$, against $16^{1/3}=2.5$. The continuum minimizer coincides with the exact one within the grid spacing, and the minimum is flat, as in \cref{obs:schedules:cooldown}: at $\alpha=\frac12$ and $T=4000$ every fraction between $0.18$ and $0.41$ is within $1\%$ of the best. At fractions this large the peak dependence of \eqref{eq:schedules:cstar} is approximate: doubling the peak multiplies $c^\star$ by $1.42$ to $1.55$ for $\alpha=\frac12$ and by $1.56$ to $1.72$ for $\alpha=1$, against $2^{2/3}\approx1.59$ and $2$. The value $c^\star\approx0.27$ depends on the noise, the curvatures and the river; the scaling with $T$ and its dependence on $\alpha$ are the prediction.

\begin{figure}[t]
\begin{fullwidth}
\centering
\begin{tikzpicture}
\begin{groupplot}[group style={group size=2 by 1,horizontal sep=17mm},wplot,width=82mm,height=48mm,xmode=log,xmin=0.01,xmax=1.05,xtick={0.01,0.03,0.1,0.3,1},xticklabels={$0.01$,$0.03$,$0.1$,$0.3$,$1$},ymode=log]
\nextgroupplot[xlabel={cooldown fraction $c$},ylabel={final expected loss},title={river $\propto A^{-1/2}$},ymin=0.055,ymax=0.36,ytick={0.06,0.1,0.15,0.2,0.3},yticklabels={$0.06$,$0.1$,$0.15$,$0.2$,$0.3$}]
\addplot[theory] table[x=c,y=a05T1000] {figures/data/12-schedules-cooldown.dat};
\addplot[theory3] table[x=c,y=a05T4000] {figures/data/12-schedules-cooldown.dat};
\addplot[theory2] table[x=c,y=a05T16000] {figures/data/12-schedules-cooldown.dat};
\addplot[simulation] table[x=c,y=a05T1000] {figures/data/12-schedules-cooldown-sgd.dat};
\addplot[simulation] table[x=c,y=a05T4000] {figures/data/12-schedules-cooldown-sgd.dat};
\addplot[simulation] table[x=c,y=a05T16000] {figures/data/12-schedules-cooldown-sgd.dat};
\addplot[only marks,mark=diamond*,mark size=2.6pt,ink,mark options={fill=white}] table[x=cstara05,y=lossa05] {figures/data/12-schedules-cooldown-opt.dat};
\nextgroupplot[xlabel={cooldown fraction $c$},ylabel={final expected loss},title={river $\propto A^{-1}$},ymin=0.015,ymax=0.3,ytick={0.02,0.05,0.1,0.2},yticklabels={$0.02$,$0.05$,$0.1$,$0.2$},legend style={at={(0.97,0.63)},anchor=east,font=\scriptsize},legend columns=1]
\addplot[theory] table[x=c,y=a1T1000] {figures/data/12-schedules-cooldown.dat};
\addplot[theory3] table[x=c,y=a1T4000] {figures/data/12-schedules-cooldown.dat};
\addplot[theory2] table[x=c,y=a1T16000] {figures/data/12-schedules-cooldown.dat};
\addplot[simulation] table[x=c,y=a1T1000] {figures/data/12-schedules-cooldown-sgd.dat};
\addplot[simulation] table[x=c,y=a1T4000] {figures/data/12-schedules-cooldown-sgd.dat};
\addplot[simulation] table[x=c,y=a1T16000] {figures/data/12-schedules-cooldown-sgd.dat};
\addplot[only marks,mark=diamond*,mark size=2.6pt,ink,mark options={fill=white}] table[x=cstara1,y=lossa1] {figures/data/12-schedules-cooldown-opt.dat};
\legend{$T=1000$,$T=4000$,$T=16000$}
\end{groupplot}
\end{tikzpicture}
\end{fullwidth}
\caption{The final loss of WSD against its cooldown fraction $c$, in the valley of \cref{fig:schedules:river} (left, river loss $\propto A^{-1/2}$) and in the same valley with initial river losses $\propto j^{-3}$ (right, $\propto A^{-1}$), at the peak $0.04$, for $T=1000$, $4000$, $16000$ (top to bottom). Curves: exact expected loss (\cref{thm:schedules:river}); points: mean of $40$ runs of stochastic gradient descent (seed $31$), within sampling error of the curves (at most $4\%$ apart, about two standard errors); diamonds: minimizer of the continuum approximation, with the river loss evaluated at the area of the schedule and the hills by \eqref{eq:schedules:cooldownlaw}, plotted at the exact loss. Script: \texttt{12-schedules-cooldown.py}.}
\label{fig:schedules:cooldown}
\end{figure}

\section{What the theory misses}\label{sec:schedules:misses}

\paragraph{What was idealized away.} The convex bound holds for every convex Lipschitz loss and so cannot see the scales of a particular one (\cref{fig:schedules:bound}). The river valley keeps two time scales but fixes everything else: quadratic hills of constant curvature, decoupled from the river; additive noise with a fixed covariance, whereas in a network the gradient covariance follows the Hessian and changes along the run; and plain stochastic gradient descent, without momentum, Adam's normalization, clipping or weight decay. The river's exponent $\alpha$ and the spread of the noise over the hills are inputs, not outputs, and they set the constants of every prediction. Neither model needs warmup, which \cref{ch:warmup} explains by the curvature that the early steps meet.

\paragraph{Where it breaks.} (1) In the river valley only stability limits the peak of a long decay, so with tuned peaks the model ranks cosine well ahead of WSD, whereas WSD matches cosine in language models (\cref{obs:schedules:cooldown}); what limits the peak in practice, instabilities and loss spikes, is outside the model (\cref{ch:eos,ch:spikes}). (2) The landscape moves with the schedule. In full-batch training the sharpness rises while it is below $2/\eta$ and then stays near $2/\eta$ \parencite{CohenEtAl2021}, so a decaying learning rate lets the sharpest directions sharpen again, and part of the oscillation that a cooldown removes is the oscillation at the edge of stability (\cref{ch:eos}), whose size is set by the dynamics rather than by noise. The two mechanisms differ in a testable way: the noise floor of \cref{cor:schedules:floor} scales like $1/B$, the oscillation at the edge does not. (3) The empirical laws relax in steps, or in area with a fitted correction, not exactly in area as a fixed quadratic hill does (\cref{sec:schedules:river}). (4) With AdamW the norms of normalized weights equilibrate on the time scale $1/(\eta\lambda)$, which a cooldown lengthens (\cref{ch:decay}), so part of what a cooldown does acts through weight norms that the river valley does not have. (5) The convex bound prefers a decay over the whole run, as do the experiments of \textcite{DefazioEtAl2023} and \textcite{BergsmaEtAl2025}, but it underestimates the penalty of short cooldowns on the cone by an order of magnitude, so its agreement with practice is one of direction. And the river valley is itself a conjecture about language models, supported by a toy derivation and by the laws it implies, not derived from a network \parencite{WenEtAl2025}.

\paragraph{Repairs and open problems.} The empirical laws are repairs of the kind the model suggests: \textcite{TissueEtAl2025} and \textcite{LuoEtAl2025} keep the area term and the annealing structure of \eqref{eq:schedules:law} and fit the saturation, and \textcite{LiEtAl2025} derive a functional law from kernel regression with power-law spectra, a noisy quadratic model whose river and hills are the two ends of one spectrum. \Textcite{DefazioEtAl2023} refine the schedule with observed gradient norms and obtain a warmup and a fast final decay; schedule-free methods \parencite{DefazioEtAl2024} average instead of decaying, and \textcite{HageleEtAl2024} compare weight averaging with cooldowns. For least squares the last iterate under step decay \parencite{GeEtAl2019} and under decaying steps in the overparameterized regime \parencite{WuEtAl2022} is understood. Open: a river valley derived from a network (Exercise~\ref{ex:schedules:derive}); a solved model whose hills sit at the edge of stability; a mechanism for the learning-rate dependence of the relaxation (Exercise~\ref{ex:schedules:relax}); and a derivation of the cooldown fraction near $20\%$ that practitioners use.

\section{Back in the wild}\label{sec:schedules:back}

\fieldnotes{Warmup, then cosine to a tenth of the peak, linear decay to zero, or a constant peak with a cooldown over the last $10$--$20\%$ (WSD).}
{Convex Lipschitz SGD; the river valley, a slow river between stiff noisy hills.}
{A last-iterate bound for every schedule; linear decay attains the averaging rate $DG/\sqrt T$ exactly; exact second moments and the annealing decomposition in the river valley.}
{Loss $=$ area term $+$ recent-schedule term; a noise floor $\approx\eta\tr\Sigma_{\mathrm{hill}}/(4B)$ that a cooldown removes, leaving $\propto(\eta hK)^{-1/2}$; best cooldown fraction $\propto T^{(2\alpha-1)/3}$; in the bound, a floor $\rho$ costs $\rho\ln T$.}
{A landscape that moves with the schedule; position-dependent noise; Adam, momentum, weight decay; tuned peaks; why $20\%$.}
{Decay to (nearly) zero (proved: linear decay is worst-case optimal, and in the river valley a final step $\rho\eta$ keeps about a fraction $\rho$ of the floor); expect the cooldown drop to scale with the peak and $1/B$, and cool down for longer than the slowest relaxation time (derived); the exact $20\%$ is folklore backed by experiments.}

The two idealizations agree on three points. The learning rate should end at or near zero: the bound charges a final floor $\rho$ a term $\rho\ln T$, and the river valley leaves a fraction $\rho$ of the noise floor in place. The loss at the end is an area term minus a recent-schedule term, which is why very different shapes end close together and why the empirical laws work. And the benefit of a longer cooldown saturates. They disagree on its length: the bound, with a tuned peak, wants a decay over the whole run; the river valley wants a cooldown set by the relaxation times and the river's slope, a fixed fraction of the run only if the river falls like the inverse square root of the area. A short cooldown branched from a stable checkpoint measures the noise floor, the size of the prize, before the shape is tuned. That $20\%$ is enough is supported by experiments \parencite{HageleEtAl2024}, not by these theories.

The chapter closes a thread from \cref{ch:prologue}: sign descent never converges at a fixed step size, and only a schedule removes its two-cycle (Exercise~\ref{ex:schedules:sign}). Momentum and schedule-free methods remove such floors by averaging instead (\cref{ch:momentum}); the river valley is the noisy quadratic model of \cref{ch:noise} at the time scale of a run; and with decoupled weight decay the schedule is the step size of a Frank--Wolfe method, a cooldown its decaying step (\cref{ch:decay}).

\begin{exercises}
\item \consolidation (Two exact bounds.) Prove \cref{cor:schedules:shapes}, and show that at $T=10^4$ the constant step $D/(G\sqrt T)$ gives a bound $1.79$ times that of the best constant step.\label{ex:schedules:exact}

\item \consolidation (Shapes of the cooldown.) For $\varphi=1$ on $[0,1-c]$ and $\varphi(s)=((1-s)/c)^p$ after, show that $\Phi(0)=1-cp/(p+1)$ and $J[\varphi]=1+1/p+\ln((p+1-cp)/c)$. For $c=0.2$, find the $p$ that minimizes $\sqrt{J/\Phi(0)}$ (about $1.35$) and compare with the $1-\sqrt\cdot$ cooldown of \cref{tab:schedules:bound}. \hint{with $u=(1-s)/c$, the integrand of \eqref{eq:schedules:J} on the cooldown is $(p+1)^2u^{p-1}/(c(2p+1))$.}\label{ex:schedules:shape}

\item \consolidation (A floor costs a logarithm.) Let $\eta_t=\eta\varphi(t/(T+1))$ with $\varphi$ positive and Lipschitz on $[0,1]$ and $\varphi(1)=\rho$. Show that the last-iterate term of \eqref{eq:schedules:lastiterate} is $\frac12G^2\eta(\rho\ln T+O(1))$, and check the entry $1.775+0.1\ln T$ of \cref{tab:schedules:bound}. \hint{near the end, $R_k\approx\rho\eta(T-k+1)$ and $V_k\approx\rho^2\eta^2(T-k+1)$.}\label{ex:schedules:floor}

\item \consolidation (Step decay.) A hill starts at $v_i(\eta)$ and the step size drops to $\eta/10$. (a) Show that the excess over the new floor is below $1\%$ of the old floor after about $22.5/(\eta h)$ steps. (b) Show that an increase of the step size enters \eqref{eq:schedules:annealing} with a positive sign; describe the hill during a linear warmup.\label{ex:schedules:step}

\item \consolidation (The cooldown law.) (a) Derive \eqref{eq:schedules:cooldownlaw} from $dv/dA=-2hv+\eta(A)\sigma^2$. (b) Show that $\chi(\kappa)=1-\kappa/3+O(\kappa^2)$ and $\chi(\kappa)\sim\frac{\sqrt\pi}2\kappa^{-1/2}$. (c) Derive \eqref{eq:schedules:Kstar}, and show that at $K^\star$ the residual noise is twice the area cost $\omega\eta K^\star/2$.\label{ex:schedules:law}

\item \variation (Weighted conversions; \textcite{DefazioEtAl2023}.) Let $w_t>0$, $W=\sum_{t\le T}w_t$, $W_t=\sum_{s=t}^Tw_s$, $z_{t+1}=z_t-\gamma w_tg_t$ and $\eta_t=\gamma w_tW_{t+1}/W$. (a) Show that $W\theta_T=\sum_tw_tz_t$, and generalize the proof of \cref{prop:schedules:linear} to $\E\Loss(\theta_T)-\Loss^\star\le W^{-1}\bigl(D^2/(2\gamma)+\frac\gamma2G^2\sum_tw_t^2\bigr)$. (b) Show by the Cauchy--Schwarz inequality that uniform weights, that is linear decay, are best. (c) If $\E\norm{g_t}_2^2\le G_t^2$, show that the best weights are $w_t\propto G_t^{-2}$; what does $\eta_t\propto G_t^{-2}\sum_{s>t}G_s^{-2}$ look like when gradient norms fall?\label{ex:schedules:weights}

\item \variation (Normalized steps need a schedule.) Sign descent on $\frac12hx^2$ is $x_{t+1}=x_t-\eta_t\sign(x_t)$ (\cref{ch:prologue}). (a) With a constant step, show that the iterates end on a two-cycle with $\abs x\le\eta$, or at the fixed point $0$ when $x_1/\eta$ is an integer. (b) Let $\eta_t\ge\eta_{t-1}/2$ for all $t$; show that $\abs{x_t}\le\eta_{t-1}$ implies $\abs{x_{t+1}}\le\eta_t$, so the loss then stays below $h\eta_t^2/2$. (c) Compare this floor, quadratic in the step and present without noise, with the noise floor of \cref{cor:schedules:floor}.\label{ex:schedules:sign}

\item \variation (The noise-optimal cooldown.) In one hill, choose $\eta_t\ge0$ to minimize $v_K$ in $v_{t+1}=(1-\eta_th)^2v_t+\eta_t^2\sigma^2$. (a) Show that the step minimizing $v_{t+1}$ given $v_t$ is $hv_t/(\sigma^2+h^2v_t)$ and gives $1/v_{t+1}=1/v_t+h^2/\sigma^2$. (b) Show that this greedy choice is optimal. \hint{$v\mapsto\min_\eta\{(1-\eta h)^2v+\eta^2\sigma^2\}$ is increasing.} (c) Deduce $v_K=\sigma^2/(h^2(K+n_0))$ with $n_0=\sigma^2/(h^2v_0)$, reached by $\eta_t=1/(h(t+n_0+1))$, and compare with \eqref{eq:schedules:cooldownlaw}.\label{ex:schedules:greedy}

\item \variation (Averaging instead of cooling.) In one hill at equilibrium at the step $\eta$, let $\bar y$ average the last $K$ iterates; put $b=1-\eta h$ and $v=v_i(\eta)$. (a) Show that $\E\bar y^2=\frac vK\bigl(\frac{1+b}{1-b}-\frac{2b(1-b^K)}{K(1-b)^2}\bigr)$, so that the hill loss at $\bar y$ is about $2/(\eta hK)$ times the floor. (b) Show that with a linear river the average gives up the same area as a linear cooldown of $K$ steps, up to one step. (c) Conclude that tail averaging beats a linear cooldown when $\eta hK$ is large; compare with \cref{ch:momentum} and with \textcite{HageleEtAl2024}, who compare averaging with cooldowns.\label{ex:schedules:averaging}

\item \simulation (Shapes and peaks.) In \texttt{12-schedules-river.py}: (a) replace the linear cooldown by a $1-\sqrt\cdot$ cooldown of the same length and check that the final loss falls from $0.131$ to $0.129$; (b) raise the peak to $0.2$, check that cosine ($0.080$) now beats WSD ($0.096$), and explain why with the river and hill parts; (c) add a linear warmup over $100$ steps and check that, at the end of the run, only the river parts change.\label{ex:schedules:sim1}

\item \simulation (The cone.) In \texttt{12-schedules-bound.py}: (a) check that constant steps settle at $\eta G^2/2$, and that after the step is halved the excess decays by a factor $e$ in about $130$ steps, at $\eta=0.0032$ and at $\eta=0.01$; (b) set $d=400$ and explain, with the first term of \eqref{eq:schedules:lastiterate}, why cosine and linear decay now end above constant steps.\label{ex:schedules:sim2}

\item \research (Derive the river.) Find a network model on structured data in which a river valley, with its exponent $\alpha$ and hill spectrum, can be derived from the data. \Textcite{WenEtAl2025} give a bigram example; a model with learned features is open.\label{ex:schedules:derive}

\item \research (Steps or area?) The river valley predicts that the loss relaxes after a decrease of the learning rate in area, the cone that it relaxes in steps. Which mechanism (a nonsmooth minimum, curvature adapting at the edge of stability, Adam's normalization) produces the learning-rate dependence that the empirical laws fit, and with what exponent? This is open, as far as we know.\label{ex:schedules:relax}
\end{exercises}

\notesandrefs
The step-size conditions of \textcite{RobbinsMonro1951} and the averaging of \textcite{PolyakJuditsky1992} are the classical theory of schedules, and \textcite{NemirovskiEtAl2009} showed that constant steps with averaging are optimal for convex Lipschitz problems. For the last iterate, \textcite{ShamirZhang2013} proved the $\log T/\sqrt T$ rate of the steps $1/\sqrt t$, \textcite{HarveyEtAl2019} showed it tight, \textcite{JainNagarajNetrapalli2019} designed a horizon-dependent schedule with the optimal rate, \textcite{Orabona2020} removed the bounded domain, and \textcite{LiuZhou2024} extended the analysis further. \Textcite{ZamaniGlineur2025} found linear decay optimal for the subgradient method, and \textcite{DefazioEtAl2023} for stochastic gradient descent. Our proof of \cref{thm:schedules:lastiterate} follows the comparison argument; our proof of \cref{prop:schedules:linear} is a conversion of the kind used by \textcite{DefazioEtAl2023}, written out for uniform weights.

WSD was introduced for MiniCPM \parencite{HuEtAl2024} and studied by \textcite{HageleEtAl2024}, and \textcite{SchaippEtAl2025} found its agreement with the convex bound; the river valley is due to \textcite{WenEtAl2025}, the laws of loss curves across schedules to \textcite{TissueEtAl2025} and \textcite{LuoEtAl2025}. We have not found the cooldown law \eqref{eq:schedules:cooldownlaw}, the optimal cooldowns \eqref{eq:schedules:Kstar} and \eqref{eq:schedules:cstar}, or the gamma-type saturation in the literature we could check; all are derived here, not proved.

\paragraph{Translating the literature.} \Textcite{DefazioEtAl2023} and \textcite{SchaippEtAl2025} write $x_t$ for our $\theta_t$ and $f$ for $\Loss$, and start at $x_1$ as we do. \Textcite{SchaippEtAl2025} write the schedule as a base learning rate $\gamma$ times a shape $\eta_t$, our peak $\eta$ times $\varphi(t/T)$, call the bound $\Omega_t$, allow a bound $G_t^2\ge\E\norm{g_t}^2$ that varies with $t$, and write $D=\norm{x_1-x_\star}$; their Theorem~3.1, which they present as a version of Theorem~10 of \textcite{DefazioEtAl2023}, is built from the same sums as \eqref{eq:schedules:lastiterate}, a term in $D^2$ and the squared steps and a term with the steps still to come, with each $\eta_t^2$ weighted by its own $G_t^2$. \Textcite{DefazioEtAl2023} work with weights $w_t$; their schedule $\eta_t\propto w_t\sum_{s>t}w_s$ is that of Exercise~\ref{ex:schedules:weights}. In \textcite{TissueEtAl2025} and \textcite{LuoEtAl2025} the forward area $S_1(t)=\sum_{s\le t}\eta_s$ is our $A_{t+1}$, the area $\sum_{s=k}^{t}\eta_s$ since a decrease at step $k$ is our $A_{t+1}-A_k$, the loss drop $\mathrm{LD}(t)$ is our annealing sum in \eqref{eq:schedules:law}, their $L_0$ and $A$ are our $L_\infty$ and $C$, and their $\alpha$ is our river exponent. The decay factor $\lambda$ (typically $0.999$) in the annealing area of \textcite{TissueEtAl2025} is our $e^{-2\bar\eta h}$, and the exponent of the power-law saturation of \textcite{LuoEtAl2025} plays the role of our gamma shape $\nu$.

\printchapterbib
\end{refsection}

% Chapter 13 of Optimizers in the Wild: weight decay.
\newcommand{\decayball}{\mathcal{B}_\lambda}
\newcommand{\decayeff}{\eta_{\mathrm{eff}}}
\newcommand{\decayfont}{\csname footnotesize\endcsname}

\casechapter{Weight Decay}{Frank--Wolfe, rotational equilibrium and the effective learning rate}{AdamW's decoupled weight decay}{ch:decay}
\begin{refsection}[chapters/13-decay.bib]

\begin{chaptersummary}
Pretraining recipes multiply the weights by $1-\eta\lambda$ before every AdamW step, with $\lambda=0.1$, although a weight matrix followed by a normalization layer computes the same function at every scale. When the step has a fixed length in a norm, decoupled decay is exactly the Frank--Wolfe method on the ball of radius $1/\lambda$: we prove that the ball attracts every initialization at rate $(1-\eta\lambda)^t$, that the learning-rate schedule acts as a Frank--Wolfe step-size schedule, and where in the ball consistent and noisy coordinates settle. When the loss is scale invariant, we prove that the norm converges from every initialization to an equilibrium at which the weights rotate by $\arccos(1-\eta\lambda)\approx\sqrt{2\eta\lambda}$ per step whatever the size of the gradients, so that only the product $\eta\lambda$ matters, and that the rotation lags any schedule. Momentum, preconditioning, gradients that respond to the effective learning rate, and the effects of decay on generalization are what these idealizations miss.
\end{chaptersummary}

\section{In the wild: a regularizer that does not regularize}\label{sec:decay:wild}

For stochastic gradient descent, weight decay and an $L_2$ penalty are the same thing: adding $\frac\lambda2\norm{\theta}_2^2$ to the loss adds $\lambda\theta$ to the gradient, and $\theta\leftarrow\theta-\eta(g+\lambda\theta)$ is $\theta\leftarrow(1-\eta\lambda)\theta-\eta g$. For Adam they differ, because the gradient of the penalty is divided by $\sqrt{\hat v}$ with the rest of the gradient. \Textcite{LoshchilovHutter2019} applied the shrinkage outside the adaptive step and found that this \term{decoupled weight decay}\index{weight decay!decoupled} makes the best decay less dependent on the learning rate and improves Adam's generalization; the result is AdamW. Every optimizer of this chapter has the form
\begin{equation}\label{eq:decay:update}
  \theta_{t+1}=(1-\eta_t\lambda)\,\theta_t-\eta_t\,u_t ,
\end{equation}
where $u_t$ is the optimizer's own step direction: the gradient for SGD, $\hat m_t/(\sqrt{\hat v_t}+\epsilon)$ for AdamW, the orthogonalized momentum for Muon. We say learning rate for $\eta_t$ in this section and step size in the theory; the two are the same.

\begin{practice}[GPT-3, \textcite{BrownEtAl2020}; LLaMA, \textcite{TouvronEtAl2023}; Muon at scale, \textcite{LiuEtAl2025Muon}]
Train with AdamW and decoupled weight decay $\lambda=0.1$. The peak learning rates of the GPT-3 models range from $6\times10^{-5}$ (the largest) to $6\times10^{-4}$ (the smallest), so each step shrinks the weights by a factor between $1-6\times10^{-6}$ and $1-6\times10^{-5}$; the factor follows the schedule, while $\lambda$ stays fixed. Decay the weight matrices; by a convention of open training code (folklore), biases and the gains of normalization layers are not decayed. When Muon replaces AdamW on the hidden matrices, keep a decoupled decay of the same form. PyTorch's AdamW multiplies the weights by $1-\eta_t\lambda$, with the default $\lambda=10^{-2}$.
\end{practice}

Two features of this practice do not fit the picture of a penalty. AdamW at $\lambda=0.1$ does not minimize the loss plus $\frac\lambda2\norm{\theta}_2^2$, which for a model with billions of weights would dwarf the loss; we shall see that it keeps the weights in the ball $\norm{\theta}_\infty\le1/\lambda=10$, far larger than the weights of a trained network. And almost every weight matrix of a transformer feeds a normalization layer, which erases the scale of its input.

\begin{observation}[\textcite{VanLaarhoven2017}; \textcite{ZhangEtAl2019}]\label{obs:decay:elr}
If the output of a weight matrix $W$ is normalized (by batch, layer or RMS normalization), the network computes the same function with $W$ and with $cW$ for every $c>0$, so an $L_2$ penalty on $W$ does not constrain the function. It changes the scale of $W$, and with it the size of each step relative to $W$: an effective learning rate. \Textcite{VanLaarhoven2017} made this argument for SGD and found that Adam removes the dependence on the scale only partly; \textcite{ZhangEtAl2019} identified the effective learning rate as one of three mechanisms by which weight decay improves the test error of vision networks.
\end{observation}

\begin{observation}[\textcite{WanEtAl2021}; \textcite{KossonMessmerJaggi2024}]\label{obs:decay:rotation}
With normalization layers and weight decay, the norm of each scale-invariant weight and its angle of rotation per step settle to a steady state. \Textcite{WanEtAl2021} derived the steady-state angle $\sqrt{2\eta\lambda}$ for SGD and confirmed their predictions on ImageNet and MS COCO. \Textcite{KossonMessmerJaggi2024} called the steady state \term{rotational equilibrium}\index{rotational equilibrium}, analysed and measured it for single neurons under AdamW, Lion and SGD with momentum, and argued that its balance across layers and neurons is part of why AdamW works better than Adam with an $L_2$ penalty.
\end{observation}

\begin{observation}[\textcite{LoshchilovHutter2019}; \textcite{WangAitchison2024}]\label{obs:decay:horizon}
The more batch passes, the smaller the optimal $\lambda$. \Textcite{LoshchilovHutter2019} proposed $\lambda=\lambda_{\mathrm{norm}}\sqrt{b/(BT)}$ for batch size $b$, $B$ training points and $T$ epochs, ``merely one possibility informed by few experiments''. \Textcite{WangAitchison2024} read AdamW's weights as an exponential moving average of its updates with a timescale of $1/(\eta\lambda)$ iterations, and report that the optimal timescale measured in epochs is roughly the same across model and dataset sizes.
\end{observation}

Muon taught the same lesson: \textcite{LiuEtAl2025Muon} name the addition of weight decay as one of the two changes that made Muon work at scale. A counter-observation: \textcite{LobachevaEtAl2021} found that with batch normalization and weight decay, training can be periodic rather than stationary; the norm shrinks, the effective learning rate grows, training destabilizes without diverging, and a new period begins.\sidenote{\Textcite{DAngeloEtAl2024} argue that for language models trained for about one epoch, weight decay acts on optimization rather than on generalization: it lowers the training loss and stabilizes training in low precision. The mechanisms of this chapter are optimization mechanisms.}

What, then, does $\lambda$ set in a training run? This chapter gives two answers, one for each of two idealizations: the radius of a ball that the weights cannot leave, and the speed at which scale-invariant weights rotate.

\section{The idealization: a shrinking step and a loss that ignores scale}\label{sec:decay:idealization}

The first idealization keeps the geometry of the step and nothing else.

\begin{assumption}[Steepest descent with decoupled decay]\label{ass:decay:nsd}
Fix a norm $\norm{\cdot}$ on $\R^d$ with dual norm $\norm{\cdot}_*$, a decay $\lambda>0$, and step sizes $\eta_t>0$ with $\eta_t\lambda\le1$. The iteration is
\[
  \theta_{t+1}=(1-\eta_t\lambda)\,\theta_t+\eta_t\lmo(g_t),
\]
that is, \eqref{eq:decay:update} with $u_t=-\lmo(g_t)$, where $g_t$ is whatever vector the optimizer feeds to its linear minimization oracle: a gradient, a stochastic gradient or a moving average of gradients.
\end{assumption}

AdamW without memory ($\beta_1=\beta_2=0$, $\epsilon=0$) is the $\ell_\infty$ case, $\lmo(g)=-\sign(g)$ (\cref{ch:sign}); Lion with decay is the $\ell_\infty$ case applied to an interpolation of momentum and gradient; Muon with decay is the spectral-norm case, $\lmo(G)=-\polar(G)$, applied to the momentum (\cref{ch:spectral}). Norms, oracles and dual norms are those of \cref{ch:steepest}. The second idealization keeps only the symmetry that normalization layers create.

\begin{definition}[Scale invariance]\label[definition]{def:decay:scaleinv}
A function $\Loss\colon\R^d\setminus\{0\}\to\R$ is \term{scale invariant}\index{scale invariance} if $\Loss(c\theta)=\Loss(\theta)$ for every $\theta\ne0$ and $c>0$.
\end{definition}

A layer normalization without gain satisfies $\mathrm{LN}(cWx)=\mathrm{LN}(Wx)$ for $c>0$, up to the small constant inside its square root, so a loss that depends on $W$ only through $\mathrm{LN}(Wx)$ is scale invariant in $W$; batch and RMS normalization, weight normalization and the normalization of queries and keys give the same structure. Throughout, $\norm{\theta}_2$ is the Euclidean norm of all entries of the weight in question (the Frobenius norm of a matrix) and $\hat\theta=\theta/\norm{\theta}_2$ its direction.

\begin{assumption}[Rotation on a sphere]\label{ass:decay:sphere}
The loss $\Loss$ and each per-example loss $\ell(\cdot;\xi)$ are scale invariant and differentiable; the step direction $u_t$ in \eqref{eq:decay:update} is orthogonal to $\theta_t$; and, in the main theorem, the stochastic gradient has a constant length on the unit sphere, $\norm{g_t}_2\norm{\theta_t}_2=\Gamma$.
\end{assumption}

\begin{dictionary}
the ball $\decayball=\{\theta:\norm{\theta}\le1/\lambda\}$ & the set AdamW's weights are drawn into: $\abs{\theta_i}\le10$ at $\lambda=0.1$\\
Frank--Wolfe step size $\gamma_t=\eta_t\lambda$ & the learning-rate schedule times the decay\\
the vertex $\lmo(g_t)/\lambda$ & where a step points: $-\sign(g_t)/\lambda$ for AdamW without memory\\
scale invariance, $\Loss(c\theta)=\Loss(\theta)$ & a weight matrix followed by a normalization layer\\
the direction $\hat\theta$; the norm $\norm{\theta}_2$ & what the layer computes; a free variable\\
$\Gamma=\norm{g}_2\norm{\theta}_2$ & the gradient norm of the same layer rescaled to unit norm\\
the angle $\omega_t$ between $\theta_t$ and $\theta_{t+1}$ & the rotation of the weight matrix in one step\\
$\arccos(1-\eta\lambda)\approx\sqrt{2\eta\lambda}$ & equilibrium rotation: $0.0077$ rad per step at $\eta=3\times10^{-4}$, $\lambda=0.1$\\
relaxation time $1/(2\eta\lambda)$ & steps to forget the initial norm, about $17{,}000$ in that run\\
the timescale $1/(\eta\lambda)$ & the number of recent updates the weights average over
\end{dictionary}

\paragraph{What was thrown away.} For the first idealization, everything except that each step has a fixed length in a norm; convergence will also need a convex loss. For the second, everything about the layer except that the loss ignores its scale; the main theorem also freezes the size of the gradient on the sphere, which in a network changes during training, and assumes steps orthogonal to the weights, which is exact for SGD and false for Adam's preconditioned and momentum steps. Each simplification isolates one mechanism: the ball holds for every sequence of bounded steps, whatever produced them, and the equilibrium needs only orthogonality and a balance of two forces.

\paragraph{What one had to notice.} In \eqref{eq:decay:update} the factor $1-\eta\lambda$ is a convex-combination weight: $(1-\eta\lambda)\theta+\eta\lambda\cdot(-u/\lambda)$ interpolates between the current point and the point $-u/\lambda$, and when the step has a fixed length this is a Frank--Wolfe step. For a scale-invariant layer the gradient step is tangent to the sphere and can only enlarge the norm, while the decay shrinks it; their balance fixes the norm, and the norm fixes how far each step turns the weights.

\begin{fieldguide}{AdamW}{\textcite{LoshchilovHutter2019}}
\fgrow{update}{$m_t=\beta_1m_{t-1}+(1-\beta_1)g_t$, \ $v_t=\beta_2v_{t-1}+(1-\beta_2)g_t^2$, \ $\theta_{t+1}=(1-\eta_t\lambda)\theta_t-\eta_t\,\hat m_t/(\sqrt{\hat v_t}+\epsilon)$}
\fgrow{norm}{$\ell_\infty$ when $\beta_1=\beta_2=0$ (\cref{ch:sign}); with the decay, Frank--Wolfe on $\norm{\theta}_\infty\le1/\lambda$ (\cref{thm:decay:frankwolfe})}
\fgrow{preconditioner}{$\diag(\sqrt{\hat v_t})+\epsilon$ (\cref{ch:adaptive})}
\fgrow{averages}{EMAs of $g$ and $g^2$; through the decay, $\theta_t$ is an EMA of $-u_s/\lambda$ over $1/(\eta\lambda)$ steps (\cref{prop:decay:ema})}
\fgrow{state}{two numbers per parameter}
\fgrow{in the wild}{$\lambda=0.1$, peak $\eta$ from $6\times10^{-5}$ to $6\times10^{-4}$ for language models; matrices decayed, gains and biases usually not}
\fgrow{caution}{codebases differ in what multiplies $\lambda$ (see the notes); Adam with an $L_2$ penalty is a different method (\cref{ex:decay:l2})}
\end{fieldguide}

\begin{fieldguide}{MuonW}{\textcite{JordanEtAl2024Muon}; \textcite{LiuEtAl2025Muon}}
\fgrow{update}{$M_t=\beta M_{t-1}+G_t$, \ $O_t\approx\polar(M_t)$ by Newton--Schulz, \ $W_{t+1}=(1-\eta_t\lambda)W_t-\eta_t\,s\,O_t$}
\fgrow{norm}{spectral (\cref{ch:spectral}); with the decay, Frank--Wolfe on $\norm{W}_{\op}\le s/\lambda$}
\fgrow{averages}{momentum on $G$; the decay averages updates over $1/(\eta\lambda)$ steps}
\fgrow{in the wild}{hidden matrices only, AdamW elsewhere; scale $s$ set to match AdamW's update size; decoupled decay as in AdamW \parencite{LiuEtAl2025Muon}}
\end{fieldguide}

\begin{fieldguide}{Hyperball}{\textcite{WenEtAl2026Hyperball}}
\fgrow{update}{$W_{t+1}=R\,(\hat W_t-\alpha_t\hat U_t)/\norm{\hat W_t-\alpha_t\hat U_t}_F$, with $\hat W=W/R$, $\norm{W_0}_F=R$, and $\hat U_t$ the update of a base optimizer (Adam, Muon) scaled to unit norm}
\fgrow{norm}{the base optimizer's, on the sphere $\norm{W}_F=R$; no decay}
\fgrow{in the wild}{unrefereed 2026 preprint, with reported speedups over weight decay up to $1.2$B parameters (\cref{sec:decay:misses})}
\end{fieldguide}

\section{The theory: decoupled weight decay is Frank--Wolfe}\label{sec:decay:frankwolfe}

To minimize a smooth convex $f$ over a compact convex set $\cK$, the \term{Frank--Wolfe method}\index{Frank--Wolfe method} \parencite{FrankWolfe1956}, or conditional gradient method, minimizes the linearization of $f$ over $\cK$ and moves a fraction $\gamma_t$ of the way to the minimizer:
\[
  v_t\in\argmin_{v\in\cK}\ip{\nabla f(x_t),v},\qquad x_{t+1}=(1-\gamma_t)\,x_t+\gamma_t\,v_t .
\]
It needs only the linear minimization oracle of $\cK$, and its iterates stay in $\cK$ as convex combinations of points of $\cK$. For the ball of a norm the oracle is $\lmo$ rescaled (\cref{fig:decay:pictures}, left). Write $\decayball=\{\theta:\norm{\theta}\le1/\lambda\}$, $\gamma_t=\eta_t\lambda$ and $\Pi_t=\prod_{s<t}(1-\gamma_s)$, with $\Pi_0=1$.

\begin{figure}[t]
\centering
\begin{tikzpicture}[scale=0.92,>=stealth]
  \begin{scope}
    \fill[seamist!70] (-1.6,-1.6) rectangle (1.6,1.6);
    \draw[seadeep,line width=0.7pt] (-1.6,-1.6) rectangle (1.6,1.6);
    \node[font=\decayfont,text=seadeep,anchor=south west] at (-1.6,1.62) {$\norm{\theta}_\infty\le1/\lambda$};
    \coordinate (th) at (0.56,0.08);
    \coordinate (v) at (-1.6,-1.6);
    \coordinate (thn) at ($(th)!0.25!(v)$);
    \draw[muted,densely dashed,line width=0.5pt] (th) -- (v);
    \draw[->,sandline!55!black,line width=0.8pt] (th) -- ++(0.55,0.75) node[anchor=west,inner sep=2pt,font=\decayfont] {$g_t$};
    \draw[->,sea,line width=1.0pt] (th) -- (thn);
    \fill[ink] (th) circle (1.4pt) node[anchor=west,font=\decayfont,xshift=2pt,yshift=-2pt] {$\theta_t$};
    \fill[sea] (thn) circle (1.4pt) node[anchor=north east,font=\decayfont,inner sep=1pt] {$\theta_{t+1}$};
    \fill[seadeep] (v) circle (1.6pt) node[anchor=north,font=\decayfont,yshift=-1pt] {$v_t=-\sign(g_t)/\lambda$};
    \coordinate (a0) at (2.55,1.95);
    \draw[muted,line width=0.5pt,densely dotted] (a0) -- (2.15,1.62) -- (1.82,1.12) -- (1.47,0.95);
    \fill[muted] (a0) circle (1.2pt) node[anchor=west,font=\decayfont] {$\theta_0$};
  \end{scope}
  \begin{scope}[xshift=5.4cm,yshift=-1.4cm]
    \def\R{2.6}
    \draw[muted!70,line width=0.4pt] (0:\R) arc[start angle=0,end angle=62,radius=\R];
    \coordinate (O) at (0,0);
    \coordinate (T) at (\R,0);
    \coordinate (S) at ({0.75*\R},0);
    \coordinate (N) at ({0.75*\R},{sqrt(1-0.75*0.75)*\R});
    \draw[->,ink,line width=0.9pt] (O) -- (T) node[anchor=north,font=\decayfont] {$\theta_t$};
    \draw[->,sandline!55!black,line width=0.9pt] (O) -- (S);
    \node[anchor=north,font=\decayfont,text=sandline!55!black] at ($(S)+(-0.35,-0.04)$) {$(1-\eta\lambda)\theta_t$};
    \draw[->,sea,line width=0.9pt] (S) -- (N) node[midway,anchor=west,font=\decayfont] {$-\eta g_t$};
    \draw[->,seadeep,line width=0.9pt] (O) -- (N) node[anchor=south,font=\decayfont] {$\theta_{t+1}$};
    \draw[muted,line width=0.4pt] ($(S)+(0,0.18)$) -- ++(-0.18,0) -- ++(0,-0.18);
    \draw[muted,line width=0.5pt] (0.75,0) arc[start angle=0,end angle=41.4,radius=0.75];
    \node[font=\decayfont] at (0.98,0.32) {$\omega$};
  \end{scope}
\end{tikzpicture}
\caption{Two pictures of decoupled weight decay, drawn with $\eta\lambda=\frac14$. Left: a step of AdamW without memory moves $\theta_t$ a fraction $\eta\lambda$ of the way to the vertex $-\sign(g_t)/\lambda$ of the ball $\norm{\theta}_\infty\le1/\lambda$, a Frank--Wolfe step; an initialization outside the ball is drawn in (\cref{thm:decay:frankwolfe}). Right: for a scale-invariant layer the gradient is orthogonal to the weights; in equilibrium $\norm{\theta_{t+1}}_2=\norm{\theta_t}_2$, and the right triangle gives $\cos\omega=1-\eta\lambda$ (\cref{thm:decay:equilibrium}).}
\label{fig:decay:pictures}
\end{figure}

\begin{theorem}[Decoupled weight decay is Frank--Wolfe]\label{thm:decay:frankwolfe}
Let the iteration satisfy \cref{ass:decay:nsd}.
\begin{enumerate}
\item With $v_t=\lmo(g_t)/\lambda$ we have $v_t\in\argmin_{v\in\decayball}\ip{g_t,v}$ and $\theta_{t+1}=(1-\gamma_t)\theta_t+\gamma_tv_t$: every step is a Frank--Wolfe step on $\decayball$ with step size $\gamma_t$ for the linear function $\ip{g_t,\cdot}$.
\item For every $t$, $\theta_t-\Pi_t\theta_0\in(1-\Pi_t)\,\decayball$. Consequently the distance from $\theta_t$ to $\decayball$ satisfies
\[
  \bigl(\norm{\theta_t}-1/\lambda\bigr)_+\le\Pi_t\,\bigl(\norm{\theta_0}-1/\lambda\bigr)_+ ,
\]
with equality for every $t$ when $\lmo(g_s)=\theta_0/\norm{\theta_0}$ for all $s$. So $\decayball$ is invariant and, at a constant step size, attracts every initialization at the rate $(1-\eta\lambda)^t$.
\item Let $\Loss$ be convex and $L$-smooth with respect to $\norm{\cdot}$, let $g_t=\nabla\Loss(\theta_t)$, and let $h_t=\Loss(\theta_t)-\Loss^\star_\lambda$ with $\Loss^\star_\lambda=\min_{\decayball}\Loss$. Then for every $t$
\begin{equation}\label{eq:decay:fwrec}
  h_{t+1}\le(1-\gamma_t)\,h_t+\tfrac L2\,\gamma_t^2\,\norm{v_t-\theta_t}^2 .
\end{equation}
For $\gamma_t=2/(t+2)$, that is $\eta_t=2/(\lambda(t+2))$, and $D=1/\lambda+\max\{1/\lambda,\norm{\theta_0}\}$, we have $h_t\le2LD^2/(t+2)$ for $t\ge1$, which is $8L/(\lambda^2(t+2))$ if $\theta_0\in\decayball$. For a constant $\gamma_t=\gamma$, $\limsup_th_t\le2L\gamma/\lambda^2=2L\eta/\lambda$, and $h_t\le(1-\gamma)^th_0+2L\gamma/\lambda^2$ for every $t$ if $\theta_0\in\decayball$.
\item If instead $g_t=\nabla\Loss(\theta_t)+\zeta_t$ with an error $\zeta_t$, then \eqref{eq:decay:fwrec} holds with the additional term $2\gamma_t\norm{\zeta_t}_*/\lambda$ on its right side.
\end{enumerate}
\end{theorem}

\begin{proof}
The plan is to read (i) off the definition of $\lmo$, to unroll the recursion for (ii), and to repeat the classical Frank--Wolfe argument for (iii) and (iv), keeping track of iterates that start outside the ball.

(i) The map $u\mapsto u/\lambda$ is a bijection from the unit ball onto $\decayball$ and multiplies $\ip{g_t,u}$ by $1/\lambda>0$, so it maps the minimizers of $\ip{g_t,\cdot}$ over the unit ball onto those over $\decayball$. The update is the identity $\eta_t\lmo(g_t)=\eta_t\lambda\cdot\lmo(g_t)/\lambda=\gamma_tv_t$.

(ii) Write $\Pi_{s,t}=\prod_{s\le k<t}(1-\gamma_k)$, so that $\Pi_{0,t}=\Pi_t$ and $\Pi_{t,t}=1$. By induction on $t$, (i) gives
\[
  \theta_t=\Pi_t\theta_0+\sum_{s<t}w_s\,v_s,\qquad w_s=\gamma_s\Pi_{s+1,t} .
\]
The weights are nonnegative because $0<\gamma_k\le1$, and they telescope, $w_s=\Pi_{s+1,t}-\Pi_{s,t}$, to $1-\Pi_t$. If $\Pi_t<1$, then $\sum_sw_sv_s/(1-\Pi_t)$ is a convex combination of points of $\decayball$, so $\sum_sw_sv_s\in(1-\Pi_t)\decayball$; if $\Pi_t=1$ the sum is zero. The triangle inequality gives $\norm{\theta_t}-1/\lambda\le\Pi_t(\norm{\theta_0}-1/\lambda)$, and we take positive parts, as $\Pi_t\ge0$. If $\lmo(g_s)=\theta_0/\norm{\theta_0}$ for all $s$, every $v_s$ is a positive multiple of $\theta_0$ and the triangle inequality is an equality. Finally, in every norm the distance from $\theta$ to $\decayball$ is $(\norm{\theta}-1/\lambda)_+$: if $\norm{\theta}>1/\lambda$, the point $\theta/(\lambda\norm{\theta})$ of the ball is at distance $\norm{\theta}-1/\lambda$, and every $x\in\decayball$ has $\norm{\theta-x}\ge\norm{\theta}-\norm{x}\ge\norm{\theta}-1/\lambda$.

(iii) The minimum $\Loss^\star_\lambda$ is attained at some $\theta^\star\in\decayball$, since $\decayball$ is compact. The descent lemma (\cref{lem:toolkit:descent}), applied to the step $\gamma_t(v_t-\theta_t)$, gives
\[
  \Loss(\theta_{t+1})\le\Loss(\theta_t)+\gamma_t\ip{\nabla\Loss(\theta_t),v_t-\theta_t}+\tfrac L2\gamma_t^2\norm{v_t-\theta_t}^2 .
\]
Because $v_t$ minimizes $\ip{\nabla\Loss(\theta_t),\cdot}$ over $\decayball\ni\theta^\star$, and $\Loss$ is convex,
\[
  \ip{\nabla\Loss(\theta_t),v_t-\theta_t}\le\ip{\nabla\Loss(\theta_t),\theta^\star-\theta_t}\le\Loss(\theta^\star)-\Loss(\theta_t)=-h_t .
\]
Substituting gives \eqref{eq:decay:fwrec}, which does not use the sign of $h_t$ (negative values are possible while $\theta_t$ is outside the ball). By (ii), $\norm{v_t-\theta_t}\le1/\lambda+\norm{\theta_t}\le2/\lambda+\Pi_te_0$ with $e_0=(\norm{\theta_0}-1/\lambda)_+$.

For $\gamma_t=2/(t+2)$ we have $\gamma_0=1$, so the first step lands on the vertex $v_0$ and $\Pi_t=0$ for $t\ge1$; hence $\norm{v_t-\theta_t}\le2/\lambda\le D$ for $t\ge1$, and $\norm{v_0-\theta_0}\le1/\lambda+\norm{\theta_0}\le D$. We show $h_t\le2LD^2/(t+2)$ by induction. At $t=1$, \eqref{eq:decay:fwrec} with $\gamma_0=1$ gives $h_1\le\frac L2D^2\le\frac23LD^2$. If the bound holds at $t\ge1$, then
\[
  h_{t+1}\le\frac t{t+2}\cdot\frac{2LD^2}{t+2}+\frac L2\cdot\frac4{(t+2)^2}\,D^2=\frac{2LD^2(t+1)}{(t+2)^2}\le\frac{2LD^2}{t+3},
\]
because $(t+1)(t+3)\le(t+2)^2$. If $\theta_0\in\decayball$, then $D=2/\lambda$.

For a constant $\gamma$, unrolling \eqref{eq:decay:fwrec} gives
\[
  h_t\le(1-\gamma)^th_0+\frac L2\gamma^2\sum_{s<t}(1-\gamma)^{t-1-s}\Bigl(\frac2\lambda+(1-\gamma)^se_0\Bigr)^2 .
\]
Expanding the square, the sum is at most $4/(\lambda^2\gamma)+(4e_0/\lambda)\,t(1-\gamma)^{t-1}+e_0^2(1-\gamma)^{t-1}/\gamma$, whose last two terms tend to zero; if $\theta_0\in\decayball$, they vanish.

(iv) Now $v_t$ minimizes $\ip{g_t,\cdot}$, so $\ip{g_t,v_t}\le\ip{g_t,\theta^\star}$, and with $\nabla\Loss(\theta_t)=g_t-\zeta_t$,
\[
  \ip{\nabla\Loss(\theta_t),v_t-\theta_t}\le\ip{\nabla\Loss(\theta_t),\theta^\star-\theta_t}+\ip{\zeta_t,\theta^\star-v_t} .
\]
The first term is at most $-h_t$ as in (iii); the second is at most $\norm{\zeta_t}_*\norm{\theta^\star-v_t}\le2\norm{\zeta_t}_*/\lambda$ by Hölder's inequality (\cref{lem:toolkit:dualnorm}). The rest of (iii) is unchanged.
\end{proof}

So the schedule $\eta_t$ is a Frank--Wolfe step-size schedule, $\gamma_t=\eta_t\lambda$, on a ball whose radius $1/\lambda$ does not depend on the schedule.\sidenote{The radius depends on what multiplies $\lambda$. In PyTorch the decay per step is $\eta_t\lambda$ and the radius $1/\lambda$; in the AdamW of \textcite{LoshchilovHutter2019} the decay is the schedule multiplier times $\lambda$ while the step also carries the base learning rate $\alpha$, and the radius is $\alpha/\lambda$.} The classical $\gamma_t=2/(t+2)$ is a learning rate that starts at $1/\lambda$ and decays like $1/t$; every schedule with $\gamma_t\to0$ and $\sum_t\gamma_t=\infty$ drives $h_t$ to zero for exact gradients (\cref{ex:decay:schedules}), and a constant one leaves a neighborhood of size at most $2L\eta/\lambda$, which a decaying schedule removes, as in \cref{ch:schedules}. These are worst cases: on the quadratics of \cref{fig:decay:frankwolfe} the gap falls like $1/t^2$ under $2/(t+2)$ and is of order $\gamma^2$ at a constant step (\cref{ex:decay:cycle}). Part (iv) adds a floor that no schedule removes: if $\E\norm{\zeta_t}_*\le\sigma_*$, the bound on $\E h_t$ does not fall below $2\sigma_*/\lambda$. The error can be that of a stochastic gradient or of a momentum average, which \cref{ch:momentum} shows how to make small; this is why Lion and Muon feed a momentum to their oracle.

Three optimizers of this book are instances. AdamW without memory is Frank--Wolfe on the $\ell_\infty$ ball, and for AdamW with memory \textcite{XieLi2024} proved that if the iterates converge under a nonincreasing schedule with divergent sum, the limit is a Karush--Kuhn--Tucker point of $\min\Loss$ subject to $\norm{\theta}_\infty\le1/\lambda$. \Textcite{ChenEtAl2024Lion} proved the corresponding statement for Lion with a Lyapunov function, Muon with decay solves a problem constrained in the spectral norm \parencite{ChenLiLiu2025}, and \textcite{PethickEtAl2025} built optimizers directly as conditional-gradient methods with the oracles of \cref{ch:steepest}. Bounded steps that are not normalized are drawn in the same way.

\begin{proposition}[Bounded steps]\label[proposition]{prop:decay:bounded}
If $\theta_{t+1}=(1-\eta_t\lambda)\theta_t-\eta_tu_t$ with $0<\eta_t\lambda\le1$ and $\norm{u_t}\le C$ for every $t$, then $\theta_t-\Pi_t\theta_0\in(1-\Pi_t)\,C\decayball$ and $(\norm{\theta_t}-C/\lambda)_+\le\Pi_t(\norm{\theta_0}-C/\lambda)_+$.
\end{proposition}

\begin{proof}
Put $v_t=-u_t/\lambda\in C\decayball$ and repeat the proof of \cref{thm:decay:frankwolfe}(ii) with $C\decayball$ in place of $\decayball$.
\end{proof}

Adam's step is bounded coordinatewise by $(1-\beta_1)/\sqrt{(1-\beta_2)(1-\beta_1^2/\beta_2)}$ when $\beta_1^2<\beta_2$ (\cref{ch:sign}), which is $1.17$ for $(\beta_1,\beta_2)=(0.9,0.95)$ and $7.27$ for $(0.9,0.999)$, so AdamW's weights are drawn into the $\ell_\infty$ ball of radius $1.17/\lambda$ or $7.27/\lambda$.\sidenote{The bound is proved in \cref{ch:sign} without bias correction. With bias correction the Cauchy--Schwarz bound after $t$ steps increases with $t$ to the same limit, as we checked numerically for these pairs.} Where in the ball do the weights sit? One coordinate of sign descent answers exactly.

\begin{proposition}[Consistent and noisy coordinates]\label[proposition]{prop:decay:regimes}
Let $0<\gamma=\eta\lambda<1$ and $x_{t+1}=(1-\gamma)x_t-\eta s_t$, where $\abs{s_t}\le1$ and $\E[s_t\mid\cF_t]=\mu$ for every $t$, with $\cF_t$ generated by $x_0,s_0,\dots,s_{t-1}$.
\begin{enumerate}
\item $\E x_t=-\mu/\lambda+(1-\gamma)^t(x_0+\mu/\lambda)$; if $s_t=\mu\in\{-1,1\}$ is deterministic, $x_t$ follows this curve to the vertex $-\mu/\lambda$.
\item If $\mu=0$ and $\E[s_t^2\mid\cF_t]=1$, as for random signs, then $\E x_t^2=(1-\gamma)^{2t}x_0^2+\eta^2\bigl(1-(1-\gamma)^{2t}\bigr)/\bigl(1-(1-\gamma)^2\bigr)$, which tends to $\eta/(\lambda(2-\eta\lambda))$: the coordinate settles at root mean square $\lambda^{-1}\sqrt{\gamma/(2-\gamma)}\approx\lambda^{-1}\sqrt{\eta\lambda/2}$.
\end{enumerate}
\end{proposition}

\begin{proof}
Taking expectations, $\E x_{t+1}=(1-\gamma)\E x_t-\eta\mu$, whose solution is (i). For (ii), $\E[x_ts_t]=\E[x_t\,\E[s_t\mid\cF_t]]=0$ because $x_t$ is determined by $\cF_t$, so $\E x_{t+1}^2=(1-\gamma)^2\E x_t^2+\eta^2$.
\end{proof}

For sign descent on a coordinate whose gradient has mean $\bar g_i$ and Gaussian noise of standard deviation $\sigma_i$, the step has mean $\mu=\erf(\bar g_i/(\sigma_i\sqrt2))$, so the coordinate settles at $\erf(\bar g_i/(\sigma_i\sqrt2))/\lambda$ (with the sign that opposes the gradient): at the vertex when the gradient is consistent, deep inside the ball when it is noise. For the GPT-3 models $\eta\lambda$ is between $6\times10^{-6}$ and $6\times10^{-5}$, and a coordinate whose gradient is noise sits at $0.2\%$ to $0.5\%$ of the radius. The ball is a constraint only for coordinates whose updates keep one sign for about $1/(\eta\lambda)$ steps; for the others, the decay sets a size by balancing noise against shrinkage. For scale-invariant weights this balance is the whole story.

\section{The theory: scale invariance and the rotational equilibrium}\label{sec:decay:equilibrium}

Gradients in this section are Euclidean. Two facts make scale-invariant weights special.

\begin{lemma}[Scale invariance]\label[lemma]{lem:decay:scaleinv}
Let $\Loss$ be scale invariant and twice differentiable on $\R^d\setminus\{0\}$. For every $\theta\ne0$ and $c>0$:
(i) $\ip{\nabla\Loss(\theta),\theta}=0$;
(ii) $\nabla\Loss(c\theta)=c^{-1}\nabla\Loss(\theta)$, so $\nabla\Loss(\theta)=\nabla\Loss(\hat\theta)/\norm{\theta}_2$;
(iii) $\nabla^2\Loss(c\theta)=c^{-2}\nabla^2\Loss(\theta)$, so the sharpness is $S(\theta)=S(\hat\theta)/\norm{\theta}_2^2$.
\end{lemma}

\begin{proof}
Differentiate $\Loss(c\theta)=\Loss(\theta)$ in $c$ at $c=1$ for (i), and in $\theta$, once and twice, for $c\nabla\Loss(c\theta)=\nabla\Loss(\theta)$ and $c^2\nabla^2\Loss(c\theta)=\nabla^2\Loss(\theta)$; then take $c=\norm{\theta}_2$ and replace $\theta$ by $\hat\theta$.
\end{proof}

A gradient step is therefore tangent to the sphere, and gradient and curvature fall as the norm grows. By (iii), the stability threshold $\eta S<2$ of \cref{ch:eos} reads $(\eta/\norm{\theta}_2^2)\,S(\hat\theta)<2$: the step size that matters is $\eta/\norm{\theta}_2^2$. If each per-example loss is scale invariant, (i) holds for every stochastic gradient, and one step of \eqref{eq:decay:update} is a right triangle (\cref{fig:decay:pictures}, right).

\begin{lemma}[One step]\label[lemma]{lem:decay:step}
Let $\theta\ne0$, $0<\eta\lambda<1$, $\theta'=(1-\eta\lambda)\theta-\eta u$ with $\ip{u,\theta}=0$, and let $\omega$ be the angle between $\theta$ and $\theta'$. Then
\begin{enumerate}
\item $\norm{\theta'}_2^2=(1-\eta\lambda)^2\norm{\theta}_2^2+\eta^2\norm{u}_2^2$;
\item $\cos\omega=(1-\eta\lambda)\norm{\theta}_2/\norm{\theta'}_2$ and $\tan\omega=\eta\norm{u}_2/((1-\eta\lambda)\norm{\theta}_2)$;
\item if $u\ne0$, then $\hat\theta'=(\hat\theta-\alpha\hat u)/\norm{\hat\theta-\alpha\hat u}_2$ with $\hat u=u/\norm{u}_2$ and $\alpha=\tan\omega$;
\item $\norm{\theta'-\theta}_2^2/\norm{\theta}_2^2=\eta^2\lambda^2+\eta^2\norm{u}_2^2/\norm{\theta}_2^2$.
\end{enumerate}
\end{lemma}

\begin{proof}
Pythagoras' theorem for the orthogonal vectors $\theta$ and $u$, with $\ip{\theta,\theta'}=(1-\eta\lambda)\norm{\theta}_2^2$ for (ii); for (iii), divide $\theta'$ by $(1-\eta\lambda)\norm{\theta}_2>0$; for (iv), $\theta'-\theta=-\eta\lambda\theta-\eta u$.
\end{proof}

Write $a=(1-\eta\lambda)^2$; then $\cos\omega=\sqrt a$ exactly when the norm does not change. The weakest form of the equilibrium needs nothing more.

\begin{proposition}[The equilibrium angle is a time average]\label[proposition]{prop:decay:average}
Let $\theta_{t+1}=(1-\eta\lambda)\theta_t-\eta u_t$ with $\ip{u_t,\theta_t}=0$, $\theta_t\ne0$, constant $\eta$ and $\lambda$, and $q_t=\norm{\theta_t}_2^2$. For all $T_1<T_2$,
\[
  \sum_{t=T_1}^{T_2-1}q_t\tan^2\omega_t=\frac{1-a}a\sum_{t=T_1}^{T_2-1}q_t+\frac{q_{T_2}-q_{T_1}}a .
\]
If $q_{T_2}=q_{T_1}$, the $q$-weighted mean of $\tan^2\omega_t$ over the window is exactly $(1-a)/a=\tan^2\bigl(\arccos(1-\eta\lambda)\bigr)$, whatever the sizes of the steps.
\end{proposition}

\begin{proof}
By \cref{lem:decay:step}(i) and (ii), $aq_t\tan^2\omega_t=\eta^2\norm{u_t}_2^2=q_{t+1}-aq_t=(q_{t+1}-q_t)+(1-a)q_t$. Sum over the window and divide by $a$.
\end{proof}

This is why measured angles are robust: gradients may be large in some steps and small in others, but once decay and growth cancel on average, the mean rotation is fixed by $\eta\lambda$ alone. With a constant gradient size on the sphere, the norm is driven to that balance from every start.

\begin{theorem}[Rotational equilibrium of SGD]\label{thm:decay:equilibrium}
Let $\theta_0\ne0$, $0<\eta\lambda<1$, and $\theta_{t+1}=(1-\eta\lambda)\theta_t-\eta g_t$ with $\ip{g_t,\theta_t}=0$ and $\norm{g_t}_2\norm{\theta_t}_2=\Gamma>0$ for every $t$. Let $q_t=\norm{\theta_t}_2^2$ and
\begin{equation}\label{eq:decay:qinf}
  q_\infty=\frac{\eta\Gamma}{\sqrt{1-a}}=\Gamma\sqrt{\frac{\eta}{\lambda(2-\eta\lambda)}}\approx\Gamma\sqrt{\frac{\eta}{2\lambda}} .
\end{equation}
\begin{enumerate}
\item $q_{t+1}=aq_t+\eta^2\Gamma^2/q_t$, and $q_t\to q_\infty$ for every $\theta_0\ne0$; the distance $\abs{\log(q_t/q_\infty)}$ never increases, and decreases strictly unless $q_t=q_\infty$.
\item Suppose $\eta\lambda\le1-1/\sqrt2$, so that $a\ge\frac12$. If $q_t\ge q_\infty$, then $0\le q_{t+1}-q_\infty\le a\,(q_t-q_\infty)$. If $q_t<q_\infty$, then $q_{t+1}>q_t$, and either $q_{t+1}\ge q_\infty$ or $0<q_\infty-q_{t+1}\le(2a-1)(q_\infty-q_t)$. Near $q_\infty$ the error is multiplied by $2a-1=1-4\eta\lambda+2\eta^2\lambda^2$ per step.
\item As $t\to\infty$: the angle $\omega_t$ between $\theta_t$ and $\theta_{t+1}$ tends to $\omega_\infty=\arccos(1-\eta\lambda)=\sqrt{2\eta\lambda}\,(1+O(\eta\lambda))$; the relative step $\norm{\theta_{t+1}-\theta_t}_2/\norm{\theta_t}_2$ tends to $\sqrt{2\eta\lambda}$ exactly; and the direction takes gradient steps on the unit sphere, $\hat\theta_{t+1}=(\hat\theta_t-\eta_{\mathrm{eff},t}\hat g_t)/\norm{\hat\theta_t-\eta_{\mathrm{eff},t}\hat g_t}_2$ with $\hat g_t=\norm{\theta_t}_2\,g_t$ the gradient at $\hat\theta_t$, whose \term{effective step size}\index{effective learning rate} $\eta_{\mathrm{eff},t}=\eta/((1-\eta\lambda)q_t)$ tends to $\tan(\omega_\infty)/\Gamma$.
\item If instead $\Gamma_t=\norm{g_t}_2\norm{\theta_t}_2\in[\Gamma_1,\Gamma_2]$ for every $t$, with $a\ge\frac12$ and $\Gamma_1/\Gamma_2>(1-a)/a$, then $q_\infty(\Gamma_1)\le\liminf_tq_t\le\limsup_tq_t\le q_\infty(\Gamma_2)$, where $q_\infty(\Gamma)$ is given by \eqref{eq:decay:qinf}.
\end{enumerate}
\end{theorem}

\begin{proof}
The plan is to reduce the dynamics to one scalar map, to show that the map moves every point toward its fixed point on a logarithmic scale, and to read the angles off \cref{lem:decay:step}.

(i) \Cref{lem:decay:step}(i) with $\norm{g_t}_2^2=\Gamma^2/q_t$ gives the recursion, and $q_{t+1}\ge aq_t>0$. Put $x_t=q_t/q_\infty$. Since $\eta^2\Gamma^2=(1-a)q_\infty^2$, dividing by $q_\infty$ gives $x_{t+1}=\varphi(x_t)$ with
\[
  \varphi(x)=a\,x+(1-a)\,\frac1x ,
\]
a convex combination, with both weights in $(0,1)$, of $x$ and $1/x$. If $x\ne1$, these two numbers lie on opposite sides of $1$ and $\varphi(x)$ lies strictly between them, so $\abs{\log\varphi(x)}<\abs{\log x}$. The numbers $\ell_t=\abs{\log x_t}$ therefore decrease to a limit $\ell\ge0$. Suppose $\ell>0$. The continuous function $y\mapsto\abs y-\abs{\log\varphi(e^y)}$ is positive on the compact set $\{y:\ell\le\abs y\le\ell_0\}$, which contains every $\log x_t$, so it has a positive minimum $\delta$ there, and $\ell_t\le\ell_0-t\delta$, which contradicts $\ell_t\ge\ell$ for $t>\ell_0/\delta$. Hence $\ell=0$.

(ii) Factor the distances:
\begin{equation}\label{eq:decay:factor}
  \varphi(x)-1=(x-1)\Bigl(a-\frac{1-a}x\Bigr),\qquad \varphi(x)-x=(1-a)\,\frac{1-x^2}x .
\end{equation}
If $x\ge1$, the second factor of the first identity lies in $[2a-1,a)\subseteq[0,a)$ because $a\ge\frac12$, so $0\le\varphi(x)-1\le a(x-1)$. If $x<1$, the second identity gives $\varphi(x)>x$, and the second factor of the first is less than $2a-1$: either it is negative and $\varphi(x)>1$, or $0\le1-\varphi(x)<(2a-1)(1-x)$, where $\varphi(x)=1$ is impossible because $\varphi(x)$ lies strictly between $x$ and $1/x$. Multiply by $q_\infty$. The local factor is $\varphi'(1)=2a-1$.

(iii) Since $q_{t+1}/q_t\to1$, \cref{lem:decay:step}(ii) gives $\cos\omega_t=(1-\eta\lambda)\sqrt{q_t/q_{t+1}}\to1-\eta\lambda$, and \cref{lem:decay:step}(iv) gives $\norm{\theta_{t+1}-\theta_t}_2^2/q_t=\eta^2\lambda^2+\eta^2\Gamma^2/q_t^2\to\eta^2\lambda^2+(1-a)=2\eta\lambda$; the expansion $\arccos(1-z)=\sqrt{2z}\,(1+z/12+O(z^2))$ gives the size of $\omega_\infty$. For the direction, \cref{lem:decay:scaleinv}(ii) for the per-example loss gives $g_t=\hat g_t/\norm{\theta_t}_2$, and \cref{lem:decay:step}(iii) gives the update with $\alpha_t\hat u_t=\eta g_t/((1-\eta\lambda)\norm{\theta_t}_2)=\eta_{\mathrm{eff},t}\hat g_t$. The limit of $\eta_{\mathrm{eff},t}$ is $\eta/((1-\eta\lambda)q_\infty)=\sqrt{1-a}/(\sqrt a\,\Gamma)$, and $\sqrt{1-a}/\sqrt a=\tan\omega_\infty$ because $\cos\omega_\infty=\sqrt a$.

(iv) Let $b_t=\eta^2\Gamma_t^2$, $b_i=\eta^2\Gamma_i^2$ and $q_i=q_\infty(\Gamma_i)=\sqrt{b_i/(1-a)}$, so that $q_{t+1}=aq_t+b_t/q_t$. \emph{Lower bound.} If $q_t<q_1$, then $q_{t+1}-q_t\ge b_1/q_t-(1-a)q_t=(1-a)(q_1^2-q_t^2)/q_t\ge(1-a)(q_1-q_t)$, so $q_1-q_{t+1}\le a(q_1-q_t)$. If $q_t\ge q_1$, then $q_{t+1}\ge aq_t+b_1/q_t\ge q_1$ by (ii) with $\Gamma=\Gamma_1$. Hence $(q_1-q_t)_+\le a^t(q_1-q_0)_+$. \emph{Upper bound.} If $q_t\ge q_2$, then $q_{t+1}\le aq_t+b_2/q_2=aq_t+(1-a)q_2$, so $q_{t+1}-q_2\le a(q_t-q_2)$. If $q_t<q_2$, then $q_{t+1}\le aq_t+b_2/q_t$, which by \eqref{eq:decay:factor} with $x=q_t/q_2$ is at most $q_2$ once $q_t\ge(1-a)q_2/a$. Since $q_1/q_2=\Gamma_1/\Gamma_2>(1-a)/a$, the lower bound gives $q_t\ge(1-a)q_2/a$ for all large $t$, and from then on $(q_{t+1}-q_2)_+\le a(q_t-q_2)_+$.
\end{proof}

The theorem corrects a natural guess about time scales. Far above equilibrium the squared norm decays like $a^t\approx e^{-2\eta\lambda t}$, as if the gradient were absent; near equilibrium it relaxes at the rate $4\eta\lambda$, twice as fast, because a smaller norm means a larger gradient (\cref{lem:decay:scaleinv}(ii)). From far below, it rises like $\sqrt{2(1-a)t}\,q_\infty$ and arrives within about $1/(4\eta\lambda)$ steps, whatever its starting value. Optimizers whose step has a fixed size, such as AdamW, Lion and Muon, have no such restoring force, and their recursion is linear.

\begin{proposition}[Steps of fixed size: AdamW, Lion, Muon]\label[proposition]{prop:decay:adamw}
Let $\theta_{t+1}=(1-\eta_t\lambda)\theta_t-\eta_tu_t$ with $0<\eta_t\lambda<1$ and $q_t=\norm{\theta_t}_2^2$.
\begin{enumerate}
\item If $\ip{u_t,\theta_t}=0$ and $\norm{u_t}_2^2=\rho^2d$ for every $t$, then $q_{t+1}=(1-\eta_t\lambda)^2q_t+\eta_t^2\rho^2d$ exactly. At a constant step size, $q_t-q_\infty=a^t(q_0-q_\infty)$ with $q_\infty=\eta^2\rho^2d/(1-a)=\rho^2\eta d/(\lambda(2-\eta\lambda))$, and $\cos\omega_t\to1-\eta\lambda$.
\item If instead $\E[\ip{u_t,\theta_t}\mid\cF_t]=0$ and $\E[\norm{u_t}_2^2\mid\cF_t]=\rho^2d$, with $\cF_t$ the past up to $\theta_t$, the same recursion holds for $\E q_t$.
\end{enumerate}
\end{proposition}

\begin{proof}
Part (i) is \cref{lem:decay:step}(i) and (ii) and the solution of a linear recursion. For (ii), expand $q_{t+1}=(1-\eta_t\lambda)^2q_t-2\eta_t(1-\eta_t\lambda)\ip{u_t,\theta_t}+\eta_t^2\norm{u_t}_2^2$ and take conditional expectations.
\end{proof}

With steps of root-mean-square size $\rho$ per entry, the equilibrium entries have root mean square $\rho\sqrt{\eta/(\lambda(2-\eta\lambda))}\approx\rho\sqrt{\eta/(2\lambda)}$, the noisy-coordinate value of \cref{prop:decay:regimes}, and the angle is that of SGD. Sign steps have $\rho=1$; an orthogonalized $d_{\mathrm{out}}\times d_{\mathrm{in}}$ step scaled by $s$ has $\rho^2d=s^2\min\{d_{\mathrm{out}},d_{\mathrm{in}}\}$. Orthogonality holds exactly for normalized gradient steps $u_t=\rho\sqrt d\,g_t/\norm{g_t}_2$; for sign steps it holds on average when the gradient on the sphere is symmetric noise, which is the hypothesis of (ii); for preconditioned steps and for momentum it fails (\cref{sec:decay:reading}). For momentum we have a derivation rather than a proof.

\begin{heuristic}
\emph{Momentum and the angle per step.} Model AdamW's step, when the gradient is noise, by $u_t=m_t=\beta_1m_{t-1}+(1-\beta_1)\xi_t$ with independent $\xi_t$ of mean zero and $\E\norm{\xi_t}_2^2=d$, the role of $g/\sqrt{\hat v}$. Then $\E\norm{u_t}_2^2\to d(1-\beta_1)/(1+\beta_1)$, while a sum of $k\gg1/(1-\beta_1)$ consecutive steps has second moment about $kd$, because the weights $(1-\beta_1)\beta_1^j$ sum to one: momentum smooths the walk without slowing its diffusion. If $1/(\eta\lambda)\gg1/(1-\beta_1)$, the weights average about $1/(\eta\lambda)$ steps, so their squared norm keeps the value $\eta d/(2\lambda)$ of \cref{prop:decay:adamw} (\cref{ex:decay:momentum} computes the correction), and the angle per step is
\begin{equation}\label{eq:decay:momentum}
  \omega\approx\sqrt{\frac{\eta^2\E\norm{u_t}_2^2}{\eta d/(2\lambda)}}=\sqrt{2\eta\lambda\,\frac{1-\beta_1}{1+\beta_1}},
\end{equation}
a factor $0.23$ below $\sqrt{2\eta\lambda}$ for $\beta_1=0.9$, while over $k$ steps the direction still turns by about $\sqrt{2\eta\lambda k}$. In the balance of \cref{prop:decay:average}, the rest of the decay is carried by the component of $m_t$ along $\theta_t$: the averaged past gradients were tangent to past directions and point partly outward from the present one.
\end{heuristic}

The equilibrium moves when the step size does, and the norm follows with a lag.

\begin{proposition}[The rotation lags the schedule]\label[proposition]{prop:decay:schedule}
In \cref{prop:decay:adamw}(i), let $q_\infty(\eta)=\rho^2\eta d/(\lambda(2-\eta\lambda))$ and $\omega_\infty(\eta)=\arccos(1-\eta\lambda)$. If $\eta_t$ is nonincreasing and $q_0\ge q_\infty(\eta_0)$, then $q_t\ge q_\infty(\eta_t)$ and $\omega_t\le\omega_\infty(\eta_t)$ for every $t$. If $\eta_t$ is nondecreasing and $q_0\le q_\infty(\eta_0)$, both inequalities reverse.
\end{proposition}

\begin{proof}
The function $q_\infty$ increases on $(0,1/\lambda)$, and $q_\infty(\eta_t)$ is the fixed point of the increasing affine map $q\mapsto(1-\eta_t\lambda)^2q+\eta_t^2\rho^2d$. So $q_t\ge q_\infty(\eta_t)$ implies $q_{t+1}\ge q_\infty(\eta_t)\ge q_\infty(\eta_{t+1})$, and induction proves the first claim. By \cref{lem:decay:step}(ii), $\tan\omega_t=\eta_t\rho\sqrt d/((1-\eta_t\lambda)\sqrt{q_t})$, which is at most the same expression with $q_\infty(\eta_t)$ in place of $q_t$, namely $\tan\omega_\infty(\eta_t)$. The nondecreasing case reverses every inequality.
\end{proof}

A decaying schedule therefore cools a scale-invariant layer faster than the quasi-static value $\sqrt{2\eta_t\lambda}$ suggests: the relaxation time $1/(2\eta_t\lambda)$ grows as $\eta_t$ falls, the norm cannot shrink in time, and near the end of a decay to zero the rotation falls like $\eta_t$ rather than $\sqrt{\eta_t}$ (\cref{ex:decay:lag}). A warmup that starts from a small norm rotates the weights faster than its equilibrium value. Because only the direction matters, all of this can be said with a learning rate on the sphere.

\begin{proposition}[Weight decay is a Hyperball schedule]\label[proposition]{prop:decay:hyperball}
Let $\Loss$ and each $\ell(\cdot;\xi)$ be scale invariant, and let the base step be $u_t=U(\hat\theta_t,\xi_t)$, a function of the direction and of the randomness $\xi_t$ only, with $\ip{u_t,\theta_t}=0$ and $\norm{u_t}_2^2=\rho^2d$. Let $q_t$ solve the recursion of \cref{prop:decay:adamw}(i) from $q_0=\norm{\theta_0}_2^2$, and run Hyperball on the unit sphere, $\hat w_{t+1}=(\hat w_t-\alpha_t\hat U_t)/\norm{\hat w_t-\alpha_t\hat U_t}_2$ with $\hat U_t=U(\hat w_t,\xi_t)/\norm{U(\hat w_t,\xi_t)}_2$, from $\hat w_0=\hat\theta_0$, with the same $\xi_t$ and
\[
  \alpha_t=\frac{\eta_t\rho\sqrt d}{(1-\eta_t\lambda)\sqrt{q_t}} .
\]
Then $\hat w_t=\hat\theta_t$ for every $t$, and the two runs have the same losses. At a constant step size, $\alpha_t\to\tan\omega_\infty\approx\sqrt{2\eta\lambda}$.
\end{proposition}

\begin{proof}
By \cref{prop:decay:adamw}(i), $\norm{\theta_t}_2^2=q_t$. If $\hat w_t=\hat\theta_t$, then $\hat U_t=\hat u_t$, and \cref{lem:decay:step}(iii) gives $\hat\theta_{t+1}=\hat w_{t+1}$. Induct; scale invariance gives equal losses; the limit follows from $q_t\to q_\infty$.
\end{proof}

The hypotheses hold exactly for normalized gradient steps, whose direction depends on $\theta$ only through $\hat\theta$ by \cref{lem:decay:scaleinv}(ii); for sign, orthogonalized and momentum steps they hold approximately, up to a radial component that averages out and a norm that changes slowly. The last result looks at the weights themselves.

\begin{proposition}[AdamW is an exponential moving average of its updates]\label[proposition]{prop:decay:ema}
For \eqref{eq:decay:update} with $0<\eta_t\lambda\le1$ and $\Pi_{s,t}=\prod_{s\le k<t}(1-\eta_k\lambda)$,
\[
  \theta_t=\Pi_{0,t}\,\theta_0-\sum_{s<t}\eta_s\,\Pi_{s+1,t}\,u_s ,\qquad \Pi_{s,t}\le\exp\Bigl(-\lambda\sum_{s\le k<t}\eta_k\Bigr).
\]
At a constant step size, $\theta_t=(1-\eta\lambda)^t\theta_0+\sum_{s<t}\eta\lambda(1-\eta\lambda)^{t-1-s}(-u_s/\lambda)$, an exponential moving average of the points $-u_s/\lambda$ with timescale $\tau=1/(\eta\lambda)$ steps, in which the updates older than $k$ steps carry total weight at most $(1-\eta\lambda)^k$.
\end{proposition}

\begin{proof}
Induction on $t$; $1-x\le e^{-x}$; and $\sum_{s<t-k}\eta\lambda(1-\eta\lambda)^{t-1-s}\le(1-\eta\lambda)^k$.
\end{proof}

The proposition is \cref{thm:decay:frankwolfe}(ii) read as an average, and with one hypothesis it gives scaling rules. A run of $T$ steps covers $\lambda\sum_t\eta_t$ timescales, $\eta\lambda T$ at a constant step size. If the right number of timescales is a fixed fraction of the run, as the evidence of \textcite{WangAitchison2024} suggests (\cref{obs:decay:horizon}), then $\eta\lambda T$ should be held fixed: at fixed learning rate and batch size, $\lambda$ should fall like the inverse of the number of training tokens, not like its inverse square root as in the rule of \textcite{LoshchilovHutter2019}. When $\mu$P scales AdamW's hidden-layer learning rate like $1/d_{\mathrm{in}}$ (\cref{ch:modular}), holding $\eta\lambda$ fixed requires $\lambda\propto d_{\mathrm{in}}$; with a fixed $\lambda$, wider models rotate more slowly and reach equilibrium later. Both rules follow from the identity and the hypothesis; neither is proved optimal.

\section{Reading the theory}\label{sec:decay:reading}

\paragraph{Only the product.} In equilibrium the rotation of a scale-invariant weight depends on $\eta$ and $\lambda$ only through $\eta\lambda$ (\cref{thm:decay:equilibrium,prop:decay:adamw}); the quotient $\eta/\lambda$ sets the norm, which the function ignores. So $(\eta,\lambda)\mapsto(c\eta,\lambda/c)$ leaves the equilibrium dynamics of every scale-invariant layer unchanged, and sweeping $\lambda$ at a fixed learning rate sweeps the speed of rotation; the two settings with equal products in \cref{fig:decay:equilibrium} (middle) give the same angles. Embeddings, output layers, gains and biases see $\eta$ and $\lambda$ separately.

\paragraph{Sizes and times.} At $\eta=3\times10^{-4}$ and $\lambda=0.1$ the equilibrium angle is $\arccos(1-3\times10^{-5})=7.75\times10^{-3}$ radians per step; with $\beta_1=0.9$ it is about $1.8\times10^{-3}$ per step by \eqref{eq:decay:momentum}, while over $1000$ steps the direction turns by about $\sqrt{2\eta\lambda\cdot1000}\approx0.24$ radians either way. The norm relaxes by a factor $e$ every $1/(2\eta\lambda)\approx1.7\times10^4$ steps for AdamW, and twice as fast near equilibrium for SGD. The number of timescales $\lambda\sum_t\eta_t$ in a run varies widely: with the hyperparameters of \textcite{BrownEtAl2020} ($300$B tokens, cosine decay to $10\%$), the smallest GPT-3 model ($6\times10^{-4}$, batches of $0.5$M tokens) covers about $18$, and the largest ($6\times10^{-5}$, batches of $3.2$M tokens) about $0.3$. If its layers behaved like the idealization, the largest model never reached equilibrium.

\begin{figure}[t]
\begin{fullwidth}
\centering
\begin{tikzpicture}
\begin{groupplot}[group style={group size=3 by 1,horizontal sep=15mm},wplot,width=59mm,height=50mm]
\nextgroupplot[xlabel={$\eta\lambda t$},ylabel={$\norm{\theta_t}_2^2/q_\infty$},ymode=log,ymin=0.045,ymax=22,xmin=0,xmax=4,
  ytick={0.0625,0.25,1,4,16},yticklabels={$\frac1{16}$,$\frac14$,$1$,$4$,$16$},
  legend style={at={(0.97,0.03)},anchor=south east,font=\scriptsize}]
\addplot[theory] table[x=s,y=sgd_hi] {figures/data/13-decay-relaxation-theory.dat};
\addplot[theory2] table[x=s,y=sgn_hi] {figures/data/13-decay-relaxation-theory.dat};
\addplot[theory,forget plot] table[x=s,y=sgd_lo] {figures/data/13-decay-relaxation-theory.dat};
\addplot[theory2,forget plot] table[x=s,y=sgn_lo] {figures/data/13-decay-relaxation-theory.dat};
\addplot[simulation,forget plot] table[x=s,y=sgd_hi] {figures/data/13-decay-relaxation-sim.dat};
\addplot[simulation2,forget plot] table[x=s,y=sgn_hi] {figures/data/13-decay-relaxation-sim.dat};
\addplot[simulation,forget plot] table[x=s,y=sgd_lo] {figures/data/13-decay-relaxation-sim.dat};
\addplot[simulation2,forget plot] table[x=s,y=sgn_lo] {figures/data/13-decay-relaxation-sim.dat};
\legend{SGD,sign descent}
\nextgroupplot[xlabel={$\eta\lambda$},ylabel={angle$\,/\arccos(1-\eta\lambda)$},xmode=log,xmin=3.5e-4,xmax=2.8e-2,ymin=0,ymax=1.17,
  ytick={0,0.23,0.5,0.75,1},yticklabels={$0$,$0.23$,$0.5$,$0.75$,$1$},
  legend style={at={(0.5,0.37)},anchor=south,font=\scriptsize},legend columns=2]
\addplot[theory,domain=3.5e-4:2.8e-2,samples=2,forget plot] {1};
\addplot[theory2,domain=3.5e-4:2.8e-2,samples=2,forget plot] {0.229416};
\addplot[simulation] table[x=el,y expr=\thisrow{sgd1a}/\thisrow{exact}] {figures/data/13-decay-sweep.dat};
\addplot[simulation2,mark size=1.9pt] table[x=el,y expr=\thisrow{adamz1a}/\thisrow{exact}] {figures/data/13-decay-sweep.dat};
\addplot[only marks,mark=square,mark size=1.9pt,color=seadeep] table[x=el,y expr=\thisrow{adamz2a}/\thisrow{exact}] {figures/data/13-decay-sweep.dat};
\addplot[only marks,mark=triangle*,mark size=2.3pt,color=sea,mark options={fill=sea,draw=white,line width=0.2pt}] table[x=el,y expr=\thisrow{adamm1a}/\thisrow{exact}] {figures/data/13-decay-sweep.dat};
\addplot[only marks,mark=triangle,mark size=2.3pt,color=sea] table[x=el,y expr=\thisrow{adamm2a}/\thisrow{exact}] {figures/data/13-decay-sweep.dat};
\addplot[simulation,forget plot] table[x=el,y expr=\thisrow{sgd2a}/\thisrow{exact}] {figures/data/13-decay-sweep.dat};
\addplot[simulation,forget plot] table[x=el,y expr=\thisrow{sgd1b}/\thisrow{exact}] {figures/data/13-decay-sweep.dat};
\addplot[simulation,forget plot] table[x=el,y expr=\thisrow{sgd2b}/\thisrow{exact}] {figures/data/13-decay-sweep.dat};
\addplot[simulation2,mark size=1.9pt,forget plot] table[x=el,y expr=\thisrow{adamz1b}/\thisrow{exact}] {figures/data/13-decay-sweep.dat};
\addplot[only marks,mark=square,mark size=1.9pt,color=seadeep,forget plot] table[x=el,y expr=\thisrow{adamz2b}/\thisrow{exact}] {figures/data/13-decay-sweep.dat};
\addplot[only marks,mark=triangle*,mark size=2.3pt,color=sea,mark options={fill=sea,draw=white,line width=0.2pt},forget plot] table[x=el,y expr=\thisrow{adamm1b}/\thisrow{exact}] {figures/data/13-decay-sweep.dat};
\addplot[only marks,mark=triangle,mark size=2.3pt,color=sea,forget plot] table[x=el,y expr=\thisrow{adamm2b}/\thisrow{exact}] {figures/data/13-decay-sweep.dat};
\legend{SGD,{$\beta_1{=}0$, $W_1$},{$\beta_1{=}0$, $W_2$},{$\beta_1{=}0.9$, $W_1$},{$\beta_1{=}0.9$, $W_2$}}
\nextgroupplot[xlabel={$t/T$},ylabel={angle per step (rad)},ymode=log,ymin=1e-5,ymax=0.1,xmin=0,xmax=1,
  legend style={at={(0.03,0.03)},anchor=south west,font=\scriptsize}]
\addplot[theory] table[x=tT,y=cos_rec] {figures/data/13-decay-schedule.dat};
\addplot[theory3] table[x=tT,y=wsd_rec] {figures/data/13-decay-schedule.dat};
\addplot[guide] table[x=tT,y=cos_quasi] {figures/data/13-decay-schedule.dat};
\addplot[guide,forget plot] table[x=tT,y=wsd_quasi] {figures/data/13-decay-schedule.dat};
\addplot[simulation2,each nth point=4,forget plot] table[x=tT,y=cos_sim] {figures/data/13-decay-schedule.dat};
\addplot[simulation,each nth point=4,forget plot] table[x=tT,y=wsd_sim] {figures/data/13-decay-schedule.dat};
\legend{cosine,WSD,quasi-static}
\end{groupplot}
\end{tikzpicture}
\end{fullwidth}
\caption{The rotational equilibrium: exact predictions (curves) and runs (marks). Left: squared norm of the linear model $x\mapsto\ip{x,\hat\theta}$ in $\R^{1024}$, trained by SGD (circles) and by sign descent, which is AdamW without memory (squares), with decoupled decay at $\eta\lambda=10^{-3}$ and batch $8$ on labels independent of the inputs, so that the gradient on the sphere is noise with $\Gamma^2=1023\cdot2/8$; one run per start (seeds $11$, $12$); curves from \cref{thm:decay:equilibrium}(i) and \cref{prop:decay:adamw}. Middle: root-mean-square angle per step, over the last $4000$ of $10^4$ steps, divided by $\arccos(1-\eta\lambda)$, for the two scale-invariant matrices of a network with two hidden layers of width $32$ followed by layer normalization, trained on teacher labels with batch $16$; SGD (circles; both layers), AdamW with $\beta_1=0$ (squares) and with $\beta_1=0.9$ (triangles; filled: $W_1$, open: $W_2$), $\beta_2=0.95$; two $(\eta,\lambda)$ pairs per product, which overlap; lines at $1$ and at the heuristic value $\sqrt{(1-\beta_1)/(1+\beta_1)}$ of \eqref{eq:decay:momentum}. Right: angle per step of the model on the left ($d=2048$, sign steps) under cosine (squares) and warmup-free WSD (circles) schedules with $\eta_{\max}\lambda=10^{-3}$ over $T=2\times10^4$ steps: recursion of \cref{prop:decay:adamw}, one run (every fourth window of $200$ steps), and the quasi-static $\arccos(1-\eta_t\lambda)$ (dashed). Scripts: \texttt{13-decay-relaxation.py}, \texttt{13-decay-sweep.py}, \texttt{13-decay-schedule.py}.}
\label{fig:decay:equilibrium}
\end{figure}

\paragraph{The figure.} In \cref{fig:decay:equilibrium} (left) the norm relaxes as the recursions predict, within the run-to-run fluctuation of about $4\%$ expected from a random walk in $1024$ dimensions. From $q_\infty/16$, SGD arrives within $\eta\lambda t\approx0.25$ while sign descent needs the full linear relaxation; from $16q_\infty$ both first decay like $a^t$. In the middle, SGD's angle in a network with layer normalization lies within $1.5\%$ of $\arccos(1-\eta\lambda)$ for every product and within $0.1\%$ for $\eta\lambda\ge2\times10^{-3}$; the small excess at the smallest products is the drift term of \cref{prop:decay:average}, since the norm is still growing slowly there. AdamW without momentum matches on the first layer and falls $10$ to $15\%$ short on the second, where its preconditioned step has an outward component along $W_2$ that we measured to carry about a fifth of the decay budget. With $\beta_1=0.9$ the ratio is $0.23$ on the first layer, as \eqref{eq:decay:momentum} predicts, and about $0.2$ on the second. On the right the recursion follows the run within $2\%$ under both schedules, and the angle falls well below its quasi-static value, to $0.36$ of it at $90\%$ of the cosine schedule and $0.06$ at $99\%$, as \cref{prop:decay:schedule} requires.

\begin{figure}[t]
\begin{fullwidth}
\centering
\begin{tikzpicture}
\begin{groupplot}[group style={group size=2 by 1,horizontal sep=16mm},wplot,width=80mm,height=52mm]
\nextgroupplot[xlabel={$\eta\lambda t$},ylabel={$\lambda\abs{\theta_i}$},ymode=log,xmin=0,xmax=10,ymin=0.012,ymax=4,
  ytick={0.02,0.1,0.5,1,3},yticklabels={$0.02$,$0.1$,$0.5$,$1$,$3$},
  legend style={at={(0.98,0.45)},anchor=east,font=\scriptsize},legend columns=1]
\addplot[guide,domain=0:10,forget plot] {1};
\addplot[theory] table[x=s,y=envC] {figures/data/13-decay-fwball.dat};
\addplot[theory3] table[x=s,y=meanS] {figures/data/13-decay-fwball.dat};
\addplot[theory2] table[x=s,y=rmsN] {figures/data/13-decay-fwball.dat};
\addplot[simulation,each nth point=5,forget plot] table[x=s,y=sgnC] {figures/data/13-decay-fwball.dat};
\addplot[simulation,each nth point=5,forget plot] table[x=s,y=sgnS] {figures/data/13-decay-fwball.dat};
\addplot[simulation,each nth point=5,forget plot] table[x=s,y=sgnN] {figures/data/13-decay-fwball.dat};
\addplot[only marks,mark=square,mark size=2.0pt,color=seadeep,each nth point=5,forget plot] table[x=s,y=adamC] {figures/data/13-decay-fwball.dat};
\addplot[only marks,mark=square,mark size=2.0pt,color=seadeep,each nth point=5,forget plot] table[x=s,y=adamS] {figures/data/13-decay-fwball.dat};
\addplot[only marks,mark=square,mark size=2.0pt,color=seadeep,each nth point=5,forget plot] table[x=s,y=adamN] {figures/data/13-decay-fwball.dat};
\legend{consistent,{signal-to-noise $1$},noise}
\nextgroupplot[xlabel={step $t$},ylabel={$\Loss(\theta_t)-\Loss^\star_\lambda$},xmode=log,ymode=log,xmin=1,xmax=1e5,ymin=1e-10,ymax=30,
  legend style={at={(0.03,0.03)},anchor=south west,font=\scriptsize},legend columns=1]
\addplot[theory3] table[x=t,y=boundfw] {figures/data/13-decay-fwgap.dat};
\addplot[theory2] table[x=t,y=boundc2] {figures/data/13-decay-fwgap.dat};
\addplot[theory] table[x=t,y=boundc3] {figures/data/13-decay-fwgap.dat};
\addplot[simulation,forget plot] table[x=t,y=fw] {figures/data/13-decay-fwgap.dat};
\addplot[simulation2,forget plot] table[x=t,y=c2] {figures/data/13-decay-fwgap.dat};
\addplot[only marks,mark=triangle*,mark size=2.0pt,color=sea,mark options={fill=sea,draw=white,line width=0.2pt},forget plot] table[x=t,y=c3] {figures/data/13-decay-fwgap.dat};
\legend{$\gamma_t=2/(t+2)$,$\gamma=10^{-2}$,$\gamma=10^{-3}$}
\end{groupplot}
\end{tikzpicture}
\end{fullwidth}
\caption{AdamW and the Frank--Wolfe ball $\norm{\theta}_\infty\le1/\lambda$. Left: sign descent (circles) and AdamW with $(\beta_1,\beta_2)=(0.9,0.95)$ (open squares), $\eta=0.01$, $\lambda=0.1$, on a linear loss with Gaussian gradient noise, every coordinate started at $3/\lambda$ (seed $31$): $20$ coordinates with a consistent gradient (largest $\lambda\abs{\theta_i}$; curve: the equality case of \cref{thm:decay:frankwolfe}(ii)), $200$ with signal-to-noise ratio $1$ (mean; curve: \cref{prop:decay:regimes}(i), limit $\erf(1/\sqrt2)=0.683$) and $2000$ with pure noise (root mean square; curve: \cref{prop:decay:regimes}(ii), limit $0.0224$). Right: the constrained gap of sign descent with decay on a convex quadratic in $\R^{10}$ with $\lambda=1$ and five active constraints at the optimum (seed $7$; $L=6.35$ by enumeration of sign vectors); marks: runs with the step rule of the curve of the same color; curves: the guarantees of \cref{thm:decay:frankwolfe}(iii). Script: \texttt{13-decay-frankwolfe.py}.}
\label{fig:decay:frankwolfe}
\end{figure}

\paragraph{Where the weights sit.} In \cref{fig:decay:frankwolfe} (left), coordinates with a consistent gradient are pulled to the face of the ball exactly along $1+(1-\eta\lambda)^t(\lambda\theta_0-1)$, so the attraction rate of \cref{thm:decay:frankwolfe}(ii) is attained; coordinates with signal-to-noise ratio one settle at $0.683/\lambda$ under sign descent and $0.705/\lambda$ under AdamW, whose averaged step is slightly larger; noise coordinates settle at $0.0224/\lambda$ under both. On the right the guarantees hold with room to spare: under $2/(t+2)$ the gap falls like a few times $1/t^2$, and at a constant step it fluctuates around $0.4\gamma^2$, far below $2L\gamma/\lambda^2$.

\paragraph{Where the agreement must fail.} The equilibrium predictions are exact for orthogonal steps, so they fail where steps are not orthogonal (momentum and preconditioning, the middle panel), where the gradient on the sphere changes faster than the norm can follow (early training, and the feedback discussed next), and where training is shorter than $1/(\eta\lambda)$. The Frank--Wolfe rates are worst cases over convex $L$-smooth losses and say nothing about a network; what transfers is the ball.

\section{What the theory misses}\label{sec:decay:misses}

\paragraph{What was idealized away.} \emph{Exact scale invariance:} embeddings, output layers, gains and biases are not scale invariant, and a matrix whose output enters the residual stream before a normalization is so only approximately; for these parameters $\eta$ and $\lambda$ act separately, and the decay does regularize. \emph{A fixed gradient size on the sphere:} in a network $\Gamma_t$ falls with the loss and depends on the effective step itself, since a larger effective step keeps the direction farther from the minimizers on the sphere. In a teacher--student version of \cref{fig:decay:equilibrium} (left), the student's wandering raised $\Gamma^2$ by $17\%$ above its noise value and the norm settled $8\%$ above the prediction made with the noise value; \cref{thm:decay:equilibrium}(iv) brackets such effects but does not model the feedback. \emph{Orthogonal steps:} preconditioners and momentum add to \cref{prop:decay:average} the radial term $-2\eta(1-\eta\lambda)\sum_t\ip{u_t,\theta_t}$, which in our network carried a fifth of the decay budget without momentum and most of it with momentum. \emph{Convexity:} the rates of \cref{thm:decay:frankwolfe} need a convex loss with a useful smoothness constant in the chosen norm; the ball statements need neither.

\paragraph{Where it breaks.} (1) The equilibrium need not be stable. \Textcite{LobachevaEtAl2021} found periodic cycles with batch normalization and weight decay, and \textcite{AminChangKhanna2026} (unrefereed, 2026) argue that a constant learning rate with weight decay cannot hold a stable interior equilibrium in such networks. By \cref{lem:decay:scaleinv}(iii), a shrinking norm raises the curvature that the step meets, so the instability of \cref{ch:eos} can return in this form; with $\Gamma$ fixed, our theorem cannot see it. (2) In vision networks weight decay does more than set a rotation speed: \textcite{ZhangEtAl2019} attribute part of its effect on test error to other mechanisms, and nothing here predicts test error. (3) The constrained solutions of the Frank--Wolfe view are not where pretraining ends: at $\lambda=0.1$ the ball is far larger than the weights, and the limit theorems of \textcite{XieLi2024} and \textcite{ChenEtAl2024Lion} assume convergence, which a finite schedule does not deliver; the constraint matters for consistent coordinates and for the implicit bias of long training (\cref{ch:bias}). (4) At scale, the effective-learning-rate view is supported but not settled. \Textcite{LiuEtAl2026ELR} (unrefereed, 2026) report that loss curves of language models with different learning rates and norms collapse, to within a few thousandths and below the variation between seeds, when indexed by the effective learning rate, and that weight decay and Hyperball act mainly through the effective-learning-rate schedules they induce. \Textcite{WenEtAl2026Hyperball} (unrefereed, 2026) report that fixing the norms with Hyperball speeds up Muon by $20$ to $30\%$ in token-equivalent terms over weight-decay baselines up to $1.2$B parameters and reduces the drift of the optimal learning rate across depths, which suggests that the equilibrium is not the best way to set the effective learning rate; \textcite{XiaoEtAl2026} (unrefereed, 2026) trace the difference between Muon with Hyperball and Muon with decay to the time course of the effective step size rather than to a better update direction, and \textcite{YuanEtAl2026} (unrefereed, 2026) show that for scale-invariant networks the two are dynamically equivalent, as in \cref{prop:decay:hyperball}. (5) Equilibrium takes time: for the largest runs the transient, not the equilibrium, may be the regime of interest (\cref{sec:decay:reading}).

\paragraph{Repairs and open problems.} Fixing the norm removes the transient and makes the angle a hyperparameter: Hyperball, and the rotational variants of AdamW, SGD and Lion of \textcite{KossonMessmerJaggi2024}, which control the rotation explicitly and reduce the need for warmup (\cref{ch:warmup}). \Cref{prop:decay:hyperball} says what such a method must do to imitate weight decay, and \cref{prop:decay:schedule} what it gains: an angle that follows the schedule without lag. Scaling rules for $\lambda$ \parencite{LoshchilovHutter2019,WangAitchison2024} rest on hypotheses about the right number of timescales that no theory derives. Open: a model in which $\Gamma$ depends on the effective step and the stability of the equilibrium, or the period of the cycles, can be computed (\cref{ex:decay:feedback}); a theory of decay for parameters that are not scale invariant; and a joint theory of decay and schedule in the river valley of \cref{ch:schedules} (\cref{ex:decay:river}).

\section{Back in the wild}\label{sec:decay:back}

\fieldnotes{AdamW with decoupled decay $\lambda=0.1$ at $\eta\sim10^{-4}$; decay for Muon at scale; rescaling $\lambda$ with training length and width.}
{(a) Steps of fixed length in a norm, decoupled decay, convex losses for the rates. (b) A scale-invariant loss, steps orthogonal to the weights, constant gradient size on the unit sphere.}
{Decay is Frank--Wolfe on the ball of radius $1/\lambda$: invariance, attraction at rate $(1-\eta\lambda)^t$, rate $O(1/t)$, an $O(\eta/\lambda)$ neighborhood. The norm of a scale-invariant weight converges from every start, and the angle per step tends to $\arccos(1-\eta\lambda)$; it has that mean whenever the norm is stationary. The rotation lags the schedule; decay is a Hyperball schedule; AdamW is an EMA of its updates.}
{Only $\eta\lambda$ matters in equilibrium: $7.7\times10^{-3}$ rad per step at $\eta\lambda=3\times10^{-5}$, after about $1/(2\eta\lambda)$ steps; momentum divides the angle per step by $\sqrt{(1+\beta_1)/(1-\beta_1)}$ (derived); a decaying schedule cools faster than $\sqrt{\eta_t\lambda}$.}
{Layers that are not scale invariant; gradients that respond to the effective step; radial components of preconditioned and momentum steps; generalization; periodic instabilities.}
{Treat $\eta\lambda$ as the hyperparameter of scale-invariant layers and keep it when $\eta$ is rescaled (proved in the idealization); keep $\eta\lambda T$ fixed when the run is lengthened (suggested, with evidence); expect the effective learning rate to lag the schedule (proved in the model).}

The theory tells a practitioner what $\lambda$ does, not what it should be. For the layers that feed a normalization, the decay is a speed: once the norm has settled, the weights turn by about $\sqrt{2\eta\lambda}$ radians per step whatever the gradients, so a sweep over $\lambda$ is a sweep over that speed, and a change of learning rate that keeps $\eta\lambda$ fixed changes nothing after the transient. The ball of radius $1/\lambda$ is real but, at $\lambda=0.1$, far from the weights of a language model; it shapes where long training ends (\cref{ch:bias}) more than how pretraining proceeds. Which product $\eta\lambda$ is best, how many timescales a run should cover, and whether an equilibrium beats a fixed norm are empirical questions, and the 2026 evidence on them is unrefereed.

The chapter closes several threads. The Frank--Wolfe identity is the normalized steepest descent of \cref{ch:steepest}, in the norms of \cref{ch:sign,ch:spectral}, made convergent by its schedule, the thread of \cref{ch:momentum,ch:schedules}. The equilibrium angle is an effective learning rate: the per-layer rates of $\mu$P (\cref{ch:modular}) and the angular updates of warmup (\cref{ch:warmup}) seen from the sphere. And by \cref{lem:decay:scaleinv}(iii), the stability threshold of \cref{ch:eos} for a scale-invariant layer is a threshold on $\eta/\norm{\theta}_2^2$, which the decay controls.

\begin{exercises}
\item \consolidation (Frank--Wolfe on an interval.) Let $f(\theta)=\frac h2\theta^2$ and $\theta_{t+1}=(1-\gamma_t)\theta_t-\gamma_t\lambda^{-1}\sign(\theta_t)$. (a) For a constant $\gamma\in(0,1)$, show that the two-cycle $\pm\gamma/(\lambda(2-\gamma))$ attracts every start, so that $f$ settles at $\frac h2\gamma^2/(\lambda^2(2-\gamma)^2)$, about $\gamma/16$ times the guarantee $2h\gamma/\lambda^2$. (b) For $\gamma_t=2/(t+2)$ and $\theta_0\ne0$, show that $\abs{\theta_t}=1/(\lambda(2\lfloor t/2\rfloor+1))$ for $t\ge1$.\label{ex:decay:cycle}

\item \consolidation (The Frank--Wolfe gap.) For convex differentiable $f$ and $\theta\in\decayball$, show that $\mathcal G(\theta)=\ip{\nabla f(\theta),\theta}+\norm{\nabla f(\theta)}_*/\lambda$ is at least $f(\theta)-\min_{\decayball}f$, and vanishes exactly at the minimizers over $\decayball$. For the $\ell_\infty$ ball, deduce that at a minimizer every coordinate has zero partial derivative or sits at $-\sign(\partial_if)/\lambda$.\label{ex:decay:gap}

\item \consolidation (Schedules for Frank--Wolfe.) Let $h_{t+1}\le(1-\gamma_t)h_t+C\gamma_t^2$ with $\gamma_t\in[0,1]$, $\gamma_t\to0$ and $\sum_t\gamma_t=\infty$. Show that $\limsup_th_t\le0$. Compute the bound on $h_T$ for the linear decay $\gamma_t=\gamma_0(1-t/T)$ and compare it with the constant step $\gamma_0$.\label{ex:decay:schedules}

\item \consolidation (The scalar map.) For $\varphi(x)=ax+(1-a)/x$: (a) verify \eqref{eq:decay:factor}; (b) show that for $a<\frac12$ the iterates approach $1$ from alternating sides; (c) show that for $a=0$ every point lies on a two-cycle; (d) for $a=(1-\eta\lambda)^2$ and $x_0=e^{10}$, show that about $10/(2\eta\lambda)$ steps are needed to reach $x_t\le e$.\label{ex:decay:map}

\item \consolidation (Angles.) Show that $\sin\omega_\infty=\sqrt{\eta\lambda(2-\eta\lambda)}$ and $2\sin(\omega_\infty/2)=\sqrt{2\eta\lambda}$. At $\eta=0.05$ and $\lambda=0.01$, evaluate $\omega_\infty$, $\sin\omega_\infty$ and $\sqrt{2\eta\lambda}$ to six digits, and explain why a simulation of the angle gives $0.031624$.\label{ex:decay:angles}

\item \variation (Adam with an $L_2$ penalty.) In Adam with $\beta_1=0$, replace decoupled decay by a penalty: $u_t=(g_t+\lambda\theta_t)/\sqrt{\hat v_t}$, $\theta_{t+1}=\theta_t-\eta u_t$, on a scale-invariant layer whose gradient on the unit sphere is noise of variance $\sigma^2$ per coordinate. Assuming $\hat v_t\approx\sigma^2/\norm{\theta_t}_2^2$, show that the decay per step is $\eta\lambda\norm{\theta_t}_2/\sigma$ and that the norm settles at $\norm{\theta}_2^3\approx\eta\sigma d/(2\lambda)$, so that the equilibrium angle depends on the size of the gradient, unlike AdamW's \parencite{KossonMessmerJaggi2024}. (A simulation with $\beta_2=0.999$ agrees to $0.3\%$.)\label{ex:decay:l2}

\item \variation (Momentum, exactly.) Let $\theta_{t+1}=c\theta_t-\eta m_t$ and $m_t=\beta m_{t-1}+(1-\beta)\xi_t$ with $c=1-\eta\lambda$ and independent $\xi_t$ of mean zero and covariance $I_d$. (a) Using \cref{lem:toolkit:lyapunov} for the pair $(\theta_{t,i},m_{t-1,i})$, show that the stationary $\E\norm{\theta}_2^2$ is $\eta^2d\,(1-\beta)(1+c\beta)/((1+\beta)(1-c\beta)(1-c^2))$, which is $\eta^2d/(1-c^2)$ times $1+O(\eta\lambda/(1-\beta))$. (b) Show that $\E\norm{m}_2^2=d(1-\beta)/(1+\beta)$ and deduce \eqref{eq:decay:momentum}. (c) For PyTorch's SGD with momentum, $m_t=\beta m_{t-1}+g_t+\lambda\theta_t$ and $\theta_{t+1}=\theta_t-\eta m_t$, show in the same model that the angle per step is about $\sqrt{2\eta\lambda/(1+\beta)}$.\label{ex:decay:momentum}

\item \variation (MuonW.) A scale-invariant $d_{\mathrm{out}}\times d_{\mathrm{in}}$ matrix is trained by $W_{t+1}=(1-\eta\lambda)W_t-\eta s\polar(G_t)$ with full-rank gradients $G_t$. (a) Show with a $2\times2$ example that $\ip{\polar(G_t),W_t}$ need not vanish although $\ip{G_t,W_t}=0$. (b) Assuming it averages to zero, compute the equilibrium $\norm{W}_F$, and compare the spectral norm of a Gaussian matrix with that Frobenius norm with the radius $s/\lambda$ of the Frank--Wolfe ball, for $s=0.2\sqrt{\max\{d_{\mathrm{out}},d_{\mathrm{in}}\}}$, $\eta=10^{-3}$, $\lambda=0.1$ and $d_{\mathrm{in}}=d_{\mathrm{out}}=1024$.\label{ex:decay:muon}

\item \variation (The lag under linear decay.) In \cref{prop:decay:adamw}(i), let $\eta_t=\eta_0(1-t/T)$ with $\eta_0\lambda T\gg1$, started at $q_\infty(\eta_0)$. Treating $t$ as continuous, show that the norm tracks $q_\infty(\eta_t)$ while $\eta_t\lambda(T-t)\gg1$ and freezes once this is of order one, after which the angle falls like $\eta_t$. Show that the transition happens when the remaining fraction of the run is about $(\eta_0\lambda T)^{-1/2}$.\label{ex:decay:lag}

\item \simulation (Universality.) Run \texttt{13-decay-sweep.py} with $\beta_1=0.5$, with $\beta_2=0.999$, and with width $64$. Check that SGD's angle stays at $\arccos(1-\eta\lambda)$ and that AdamW's with $\beta_1=0.5$ is close to $1/\sqrt3$ of it on the first layer, as \eqref{eq:decay:momentum} predicts; measure the radial component of AdamW's step on each layer and explain why it can differ between layers.\label{ex:decay:sim1}

\item \simulation (Where in the ball.) In \texttt{13-decay-frankwolfe.py}, vary the signal-to-noise ratio of the middle group over $\{0.25,0.5,1,2,4\}$; check that sign descent settles at $\erf(\mathrm{snr}/\sqrt2)/\lambda$, measure where AdamW settles, and explain why the two differ most at intermediate ratios. With $\eta=3\times10^{-4}$ and $\lambda=0.1$, check that the noise coordinates settle at $0.39\%$ of the radius.\label{ex:decay:sim2}

\item \research (Feedback and stability.) Let the gradient size on the sphere depend on the effective step, $\Gamma^2=\Gamma_0^2+\kappa\,\decayeff$, as for a student that wanders around a teacher. Find the equilibrium of \cref{thm:decay:equilibrium} and its stability. Is there a natural model, for instance one in which $\decayeff S(\hat\theta)$ is capped at $2$ as in \cref{ch:eos}, where the equilibrium loses stability and the cycles of \textcite{LobachevaEtAl2021} appear? As far as we know this is open.\label{ex:decay:feedback}

\item \research (Weight decay in the river valley.) Combine the river-valley model of \cref{ch:schedules} with a scale-invariant layer whose direction carries the river and the hills. For a given schedule, which $\lambda$ minimizes the final loss, and how does the answer depend on the length of training? Compare with the timescale rule of \cref{prop:decay:ema}.\label{ex:decay:river}
\end{exercises}

\notesandrefs
Decoupled weight decay is due to \textcite{LoshchilovHutter2019}, who showed that it differs from an $L_2$ penalty for Adam \parencite{KingmaBa2015}. That normalization turns a penalty on the weights into a control of the effective learning rate was observed by \textcite{VanLaarhoven2017}, and \textcite{ZhangEtAl2019} placed it among three mechanisms of weight decay. \Textcite{WanEtAl2021} derived the equilibrium angle for SGD with momentum and proved, under their assumptions, convergence to the equilibrium at a linear rate; \textcite{KossonMessmerJaggi2024} measured it per neuron for several optimizers, and \textcite{LobachevaEtAl2021} found the periodic alternative. \Cref{thm:decay:equilibrium} is the deterministic core of these analyses, with convergence from every initialization and a comparison bound; we know no earlier statement of the identity of \cref{prop:decay:average} in this form.

The Frank--Wolfe method was introduced by \textcite{FrankWolfe1956} for quadratic programs; \textcite{Jaggi2013} gave the analysis with the open-loop step $2/(t+2)$ that made it popular in machine learning, and \textcite{Hazan2016} treats it among projection-free methods. The identification of decoupled decay with Frank--Wolfe on a norm ball runs through \textcite{XieLi2024} for AdamW, \textcite{ChenEtAl2024Lion} for Lion, \textcite{ChenLiLiu2025} for Muon and \textcite{PethickEtAl2025} for general oracles; \textcite{BernsteinNewhouse2024} cast these optimizers as steepest descent under norms, and \textcite{Newhouse2025} connects duality, weight decay and norm constraints in a thesis. The moving-average reading of AdamW and its scaling rules are those of \textcite{WangAitchison2024}. The 2026 literature on fixed norms and effective learning rates is unrefereed \parencite{WenEtAl2026Hyperball,LiuEtAl2026ELR,XiaoEtAl2026,YuanEtAl2026,AminChangKhanna2026}. The figures were computed for this chapter; the network of \cref{fig:decay:equilibrium} is small, and its agreement with the theory checks the identities, not the behavior of large models.

\paragraph{Translating the literature.} \Textcite{LoshchilovHutter2019} write AdamW as $\theta_t=\theta_{t-1}-\eta_t\bigl(\alpha\hat m_t/(\sqrt{\hat v_t}+\epsilon)+\lambda\theta_{t-1}\bigr)$, with $\alpha$ the learning rate and $\eta_t$ a schedule multiplier: their $\alpha\eta_t$ is our $\eta_t$ and their $\lambda$ is our $\alpha\lambda$. PyTorch's AdamW decays by \texttt{lr}$\cdot$\texttt{weight\_decay} per step and is our $(\eta,\lambda)$ \parencite{Schaipp2024}. \Textcite{WangAitchison2024} write $\tau_{\mathrm{iter}}=1/(\eta\lambda)$ for our $\tau$, and $\tau_{\mathrm{epoch}}=\tau_{\mathrm{iter}}/M$ with $M$ iterations per epoch. \Textcite{KossonMessmerJaggi2024} follow the weight vector of one neuron, a row of our $W$; their angular update, like that of \textcite{WanEtAl2021}, is our $\omega_t$. Hyperball's weights have norm $R$, and its step of norm $\eta_tR$ before renormalization makes its $\eta_t$ our $\alpha_t$ of \cref{prop:decay:hyperball}. \Textcite{LiuEtAl2026ELR} define the effective learning rate as the learning rate divided by the parameter norm, proportional to our $\tan\omega_t$ for steps of fixed size; for SGD our effective step size divides by the squared norm. In the Frank--Wolfe literature $\gamma_t$ is the step size and $\cK$ the constraint set, here $\eta_t\lambda$ and $\decayball$.

\printchapterbib
\end{refsection}


\casepart{Choosing an Optimizer}{The two lines of the book meet in two questions: why Adam trains transformers so much faster than gradient descent does, and where each optimizer ends up when it runs for a long time. Both answers depend on the geometry of the problem as much as on the optimizer.}{90}
% Chapter 14 of Optimizers in the Wild: why Adam beats SGD on transformers.

\casechapter{Why Adam Beats SGD on Transformers}{Heterogeneous curvature, heavy-tailed classes and the batch-size regime}{Adam for transformers, SGD for convnets}{ch:transformers}
\begin{refsection}[chapters/14-transformers.bib]

\begin{chaptersummary}
Transformers are trained with Adam because stochastic gradient descent trains them far more slowly, while on convolutional networks the two optimizers are close; the gap survives full batches and widens with the batch size. We study three solvable idealizations: a table whose rows are used with Zipfian frequencies, a quadratic whose parameter blocks sit at different curvature scales, and the noisy quadratic model with a diagonal preconditioner. One scale invariance connects them: gradient descent turns the weight of a row or a block into a clock, while sign descent and Adam do not see the weight at all. We prove that under Zipf's law gradient descent removes only a fraction $\ln10/Z_K\approx\ln10/\ln K$ of the loss per decade of steps, that blockwise step sizes replace the global condition number by the worst block's, and that a preconditioner's advantage grows with the batch size from a noise-weighted constant to its condition-number factor. The models leave out parameter sharing and the interactions of the mechanisms, and with averaged iterates the small-batch advantage disappears, which is the regime in which plain SGD has been reported to match Adam.
\end{chaptersummary}

\section{In the wild: Adam for language, SGD for vision}\label{sec:transformers:wild}

The transformer was trained with Adam from the start \parencite{VaswaniEtAl2017,KingmaBa2015}, and nearly all large language models since have been trained with Adam, AdamW or a close relative. Convolutional networks for image classification are still often trained with stochastic gradient descent with momentum, and there the choice is not a matter of speed: on image classification, SGD reaches test accuracy as good as that of adaptive methods or better \parencite{WilsonEtAl2017}. In this chapter \emph{learning rate} (the practitioners' word) and \emph{step size} (the theorist's) mean the same thing.

\begin{practice}[GPT-3, \textcite{BrownEtAl2020}; LLaMA, \textcite{TouvronEtAl2023}; ResNets on ImageNet, \textcite{GoyalEtAl2017}]
Train a transformer language model with Adam or AdamW, $\beta_1=0.9$, $\beta_2=0.95$, weight decay $0.1$, a warmup and a cosine decay (GPT-3 also sets $\epsilon=10^{-8}$); SGD is not used. Train a ResNet on ImageNet with SGD with momentum $0.9$, learning rate $0.1\cdot B/256$ at batch size $B$, weight decay $10^{-4}$ and a step decay; Adam is not needed.
\end{practice}

The rule is folklore in the precise sense of this book: it is what everyone does, and the evidence for it is a long list of training runs rather than a theorem. What has been measured about the gap, and by whom, is the following.

\begin{observation}[The gap; \textcite{ZhangEtAl2020,KunstnerEtAl2023,ZhangEtAl2024}]\label[observation]{obs:transformers:gap}
On transformers, and most of all on language models, tuned SGD with momentum reduces the training loss much more slowly than Adam; on convolutional networks for image classification the two are close. The gap is a gap in optimization: it is visible in the training loss.
\end{observation}

The observation is robust across the papers cited, but its size depends on how carefully each optimizer is tuned: at the scale of a billion parameters, the measured speedups of the best optimizers over a well-tuned AdamW are much smaller than many earlier reports \parencite{WenEtAl2026}. The first explanation offered was the noise. \Textcite{ZhangEtAl2020} found the stochastic gradients of attention models heavy-tailed, and argued that adaptive methods and clipping cope with heavy-tailed noise where SGD does not. A test of that explanation came three years later.

\begin{observation}[Not the noise; \textcite{KunstnerEtAl2023}]\label[observation]{obs:transformers:noise}
As the batch size grows, up to the full training set, the gap between SGD and Adam on transformers does not close: SGD gains less from larger batches than Adam does. With full batches, Adam behaves much like sign descent, and sign descent with momentum closes most of the gap.
\end{observation}

Sign descent is steepest descent in the $\ell_\infty$ norm (\cref{ch:sign}), so the observation shifts the question from the noise to the geometry: what is it about a transformer's loss that makes a step of fixed size in every coordinate so much better than a step along the gradient? Two answers have been proposed, and both have been tested.

\begin{observation}[Rare classes are slow; \textcite{KunstnerEtAl2024}]\label[observation]{obs:transformers:rare}
Group the tokens of a language model into buckets of similar frequency and follow the training loss on each bucket. Under gradient descent the loss on the frequent tokens falls first, and the loss on the infrequent tokens, which together account for most of the examples, falls much more slowly; under Adam and sign descent all buckets progress at comparable speeds. The same picture appears in vision models and in linear models when the classes are heavy-tailed.
\end{observation}

\begin{observation}[Block heterogeneity; \textcite{ZhangEtAl2024}]\label[observation]{obs:transformers:blocks}
Compute the spectrum of the Hessian restricted to each parameter block (embeddings, attention, feed-forward layers, normalization gains, output layer). In transformers the spectra differ widely from block to block; in convolutional networks they are similar. SGD performs on par with Adam on problems without this \term{block heterogeneity}\index{block heterogeneity}, including quadratic problems built for the test, and worse than Adam when it is present.
\end{observation}

\Textcite{KunstnerEtAl2024} trace the slowness of rare classes to the gradients and the Hessian of a linear model, which become imbalanced across classes and correlated with each other; \textcite{ZhangEtAl2024} trace the block effect to the single step size that SGD uses for all blocks.\sidenote{The class-imbalance explanation first appeared as a workshop paper, \textcite{YadavEtAl2023}. Frederik Kunstner's PhD thesis (University of British Columbia, 2024) collects much of this line of work.} Two further findings point the same way. \Textcite{ZhaoEtAl2025} report that most of Adam's advantage is kept when adaptivity is restricted to the last layer and the normalization parameters and SGD trains the rest, and \textcite{XieEtAl2025} report that rotating the parameters of a GPT-2 model slows Adam down but leaves SGD unaffected, which says that Adam's advantage lives in the coordinate system. The last observation brings the batch size back.

\begin{observation}[Small batches; \textcite{MarekEtAl2025}]\label[observation]{obs:transformers:small}
Language models train stably with batch sizes down to one when Adam's $\beta_2$ is rescaled to keep its half-life fixed in tokens, and small batches are more robust to hyperparameters and as efficient per FLOP as large ones or more. At small batch sizes even vanilla SGD, without momentum, trains them stably and comes close to Adam; the gap grows with the batch size.
\end{observation}

\Textcite{MarekEtAl2025} is refereed; an unrefereed study reaches a similar conclusion for SGD with momentum \parencite{SreckovicEtAl2025}. So the explanations point in five directions: heavy-tailed noise, heavy-tailed classes, heterogeneous curvature, the coordinate system and the batch size. This chapter puts three of them into models that can be solved exactly, and asks of each model: \emph{which structure of the loss makes a per-coordinate normalization worth an order of magnitude in steps, why is the benefit specific to transformers, and why does it depend on the batch size?}

\section{The idealization: frequencies, blocks and batches}\label{sec:transformers:ideal}

Each of the three models isolates one mechanism, and one property of optimizers connects them.

\begin{definition}[Table models]\label[definition]{def:transformers:table}
Let $q_1\ge q_2\ge\dots\ge q_K>0$, $\sum_jq_j=1$, be the frequencies of $K$ items (tokens or classes). A \term{table model}\index{table model} has one row $w_j\in\R^r$ per item and the loss
\begin{equation}\label{eq:transformers:table}
  \Loss(W)=\sum_{j=1}^Kq_j\,\ell_j(w_j),\qquad W=(w_1,\dots,w_K),
\end{equation}
with differentiable row losses $\ell_j$. With $\ell_j(w)=\frac12\norm{w-y_j}_2^2$ it is the \emph{square-loss table}; with $\ell_j(w)=-\sum_k\nu_{jk}\log\softmax(w)_k$ for probability vectors $\nu_j\in\R^r$ it is the \emph{softmax table}. The frequencies follow \term{Zipf's law}\index{Zipf's law} if $q_j=1/(jZ_K)$, where $Z_K=1+\frac12+\dots+\frac1K$ is the $K$-th harmonic number.
\end{definition}

The softmax table is a bigram language model: row $j$ holds the logits of the token that follows token $j$, $\nu_j$ is the conditional distribution of that token, and $q_j$ is the frequency of token $j$ as an input. An embedding table has the same structure, since its row $j$ receives a gradient only from the examples that contain token $j$. The square-loss table with $y_j=\nu_j$ is the bigram model trained with the square loss on one-hot targets, because $\E[\frac12\norm{w_X-e_Y}_2^2\mid X=j]=\frac12\norm{w_j-\nu_j}_2^2+\frac12(1-\norm{\nu_j}_2^2)$; this is the model of \textcite{KunstnerBach2025}. For the output layer of a classifier the frequencies are those of the labels, and we shall see that near the optimum they enter as curvatures in the same way (\cref{prop:transformers:softmax}). The model of \cref{sec:transformers:freq} is the square-loss table with Zipfian frequencies and rows that all start at distance one from their targets.

The second model is a quadratic $\Loss(\theta)=\frac12(\theta-\theta^\star)^\top H(\theta-\theta^\star)$ with a block-diagonal Hessian $H=\diag(H_1,\dots,H_N)$, $H_b\succ0$, the blocks standing for the parameter groups of \cref{obs:transformers:blocks}. A table model is the special case in which the blocks are rows. The third model is the noisy quadratic model of \cref{ch:noise}.

\begin{assumption}[Noisy quadratic model; \textcite{ZhangEtAl2019}]\label[assumption]{ass:transformers:nqm}
The loss is $\Loss(\theta)=\frac12\sum_{i=1}^dh_i\theta_i^2$ with $h_1\ge\dots\ge h_d>0$. At step $t$ the optimizer receives $g_t=H\theta_t+\zeta_t$, where the $\zeta_t$ are independent of each other and of the past, with mean zero and covariance $\diag(c_1,\dots,c_d)/B$ for a batch size $B$. Preconditioned SGD with multipliers $a_1,\dots,a_d>0$ updates $\theta_{t+1,i}=\theta_{t,i}-\eta a_ig_{t,i}$; SGD is $a_i=1$. We write $\lambda_i=a_ih_i$ for the curvatures that the preconditioned method sees.
\end{assumption}

Near a minimum of a well-specified model the gradient covariance is close to the Hessian, which suggests $c_i\propto h_i$; our main example takes $h_i=1/i$ and $c_i=h_i$, and \cref{sec:transformers:reading} varies the noise. The model of Adam is $a_i=h_i^{-1/2}$. The reason is that Adam's second moment estimates $\E g_{t,i}^2=h_i^2\theta_{t,i}^2+c_i/B$; at small batch sizes the noise term dominates, the preconditioner is $\diag(c)^{1/2}/\sqrt B$, and with $c\propto h$ this is a constant times $H^{1/2}$ (Exercise~\ref{ex:transformers:adamnqm}). At large batch sizes the gradient term dominates, and Adam is closer to sign descent.\sidenote{The noisy quadratic model of \textcite{ZhangEtAl2019} uses the same device, preconditioning by a power of the curvature, to stand for adaptive and second-order optimizers. That Adam's preconditioner is not the Hessian, and not the Fisher information either, is the point of \textcite{KunstnerHennigBalles2019}; see \cref{ch:adaptive}.}

The property that connects the three models is a symmetry of the optimizer.

\begin{definition}[Scale invariance]\label[definition]{def:transformers:invariant}
An optimizer computes its update $u_t$ from the gradients $g_0,\dots,g_t$ seen so far and from its hyperparameters, and sets $\theta_{t+1}=\theta_t+u_t$. Partition the coordinates into blocks. The optimizer \emph{acts blockwise} if the update of each block depends only on the gradients of that block, and it is \term{scale invariant}\index{scale invariance} if multiplying the gradients of every block, at all steps, by a positive constant of that block leaves all updates unchanged. With blocks of one coordinate we say \emph{coordinatewise}.
\end{definition}

Sign descent $u_t=-\eta\sign(g_t)$ is coordinatewise scale invariant, and so are signum and Lion, which take signs of moving averages (\cref{ch:sign,ch:momentum}). So are Adam, RMSProp and AdaGrad with $\epsilon=0$: if every $g_{t,i}$ is multiplied by $s_i>0$, then $m_{t,i}$ becomes $s_im_{t,i}$ and $v_{t,i}$ becomes $s_i^2v_{t,i}$, and $m_{t,i}/\sqrt{v_{t,i}}$ does not change (we read $0/0$ as $0$). Variants of Adam that keep one second moment per block, such as Adalayer \parencite{ZhaoEtAl2025} and Adam-mini \parencite{ZhangEtAl2025}, are blockwise scale invariant, and so is blockwise normalized descent $u_{t,b}=-\eta g_{t,b}/\norm{g_{t,b}}_2$. Gradient descent is not; normalized gradient descent $u_t=-\eta g_t/\norm{g_t}_2$ is invariant only when all blocks are scaled by the same constant; and Adam with $\epsilon>0$ is invariant only for gradients much larger than $\epsilon$, below which it is SGD with step size $\eta/\epsilon$ (\cref{ch:adaptive}).

\begin{dictionary}
the frequency $q_j$ of row $j$ & how often token $j$ occurs, as an input (embedding and bigram rows) or as a label (rows of the output layer)\\
gradient descent's clock $\eta q_jt$ on row $j$ & the learning rate times the number of times token $j$ has been seen\\
Zipf's law $q_j\propto1/j$; the harmonic number $Z_K\approx\ln K+0.58$ & the token frequencies of text; $Z_K\approx11.4$ for a vocabulary of $50{,}000$ tokens\\
a scale-invariant optimizer & Adam with a small $\epsilon$, sign descent, signum, Lion; Adam-mini and Adalayer blockwise\\
blocks $H_b$ at scales $s_b$ & parameter groups whose Hessian blocks have different sizes (\cref{obs:transformers:blocks})\\
the noisy quadratic model, $c\propto h$ & training near a minimum, where the gradient covariance is close to the Hessian\\
the small-batch limit, $B\to0$ at fixed $\eta/B$ & small batches with the learning rate scaled linearly in $B$\\
the large-batch limit & training at the largest stable learning rate, limited by curvature\\
the multipliers $a_i=h_i^{-1/2}$ & Adam at small batch sizes, whose second moment is dominated by noise
\end{dictionary}

\paragraph{What was thrown away.} The models keep one mechanism each and drop everything else. The rows of a table share no parameters, while in a transformer the rows of the embedding feed shared attention and feed-forward layers. The table and block models are deterministic and have fixed curvature; the noisy quadratic model has noise but fixed curvature and Gaussian noise. There is no nonlinearity and no feature learning, and Adam's moving averages appear only as limits. Each simplification is the right first one because each model can be solved exactly, so that the signature of its mechanism, a loss that decays logarithmically, a rate set by the blockwise spectra, a gap that depends on the batch size, can be compared with the observations one at a time.

\paragraph{What one had to notice.} That the frequency of a token is a curvature. In a table model the Hessian is block diagonal with blocks $q_j\nabla^2\ell_j$, so the heavy tail of the token distribution becomes a heavy tail of the Hessian's spectrum, and Zipf's law becomes the spectrum $1/j$, the hardest one for gradient descent \parencite{KunstnerBach2025}. And that per-coordinate normalization is exactly the operation that removes such a weight, so that what looks like a statement about curvature is a statement about a symmetry.

\section{The theory: frequencies are clocks}\label{sec:transformers:freq}

The first result is the symmetry in its simplest form.

\begin{proposition}[Weights are clocks]\label[proposition]{prop:transformers:clock}
Let $\Loss(\theta)=\sum_{b=1}^Ns_b\,\ell_b(\theta_b)$ with weights $s_b>0$ and functions $\ell_b$ with locally Lipschitz gradients, where $\theta=(\theta_1,\dots,\theta_N)$ is split into blocks.
\begin{enumerate}
\item Gradient descent with step size $\eta$ on $\Loss$ moves each block $\theta_b$ as gradient descent with step size $\eta s_b$ on $\ell_b$ alone. Gradient flow on $\Loss$ is $\theta_b(t)=\varphi_b(s_bt)$, where $\varphi_b$ is gradient flow on $\ell_b$ from $\theta_b(0)$.
\item An optimizer that acts blockwise and is scale invariant moves each block exactly as it moves when run on $\ell_b$ alone, from the same point and with the same hyperparameters. In particular, its iterates do not depend on $s_1,\dots,s_N$.
\end{enumerate}
\end{proposition}

\begin{proof}
The gradient of $\Loss$ with respect to $\theta_b$ is $s_b\nabla\ell_b(\theta_b)$, which gives (i) for gradient descent. For the flow, $\psi(t)=\varphi_b(s_bt)$ satisfies $\psi'(t)=s_b\varphi_b'(s_bt)=-s_b\nabla\ell_b(\psi(t))$, the equation of block $b$ under gradient flow on $\Loss$, with the same initial point; solutions are unique because the gradients are locally Lipschitz.

For (ii), let $\tilde\theta_{b,t}$ be the iterates of the optimizer run on $\ell_b$ alone, and argue by induction that $\theta_{b,t}=\tilde\theta_{b,t}$. This holds at $t=0$. If it holds up to step $t$, then the gradients of block $b$ seen so far are $s_b\nabla\ell_b(\tilde\theta_{b,\tau})$ for $\tau\le t$, which are $s_b$ times those of the solo run. Because the optimizer acts blockwise, the update of block $b$ depends only on these gradients, and because it is scale invariant, it equals the update of the solo run. Hence $\theta_{b,t+1}=\tilde\theta_{b,t+1}$.
\end{proof}

For gradient descent the step size must be stable for every block, so on a table whose row losses have curvature at most $\Lambda$ it must satisfy $\eta q_1\Lambda<2$: the most frequent row caps the step size, and the rarest row then needs about $1/(\eta q_K)>q_1/(2q_K)$ steps. The ratio $q_1/q_K$ is a condition number made of frequencies alone, and under Zipf's law it is $K$. What this costs is the content of the next theorem. Its proof rests on a lemma about harmonic sums with exponential weights; we write $\gamma=0.5772\ldots$ for Euler's constant and $E_1(z)=\int_z^\infty u^{-1}e^{-u}\dd u$ for the exponential integral.

\begin{lemma}[Exponentially weighted harmonic sums]\label[lemma]{lem:transformers:harmonic}
For $b>0$ and an integer $K\ge1$, the sum $E_K(b)=\sum_{j=1}^Kj^{-1}e^{-b/j}$ satisfies
\[
  -\frac1{eb}-\frac{e^{-b}}b\le E_K(b)-\ln\frac Kb+\gamma\le\frac1{eb}+\frac bK .
\]
Moreover, if a function $f\ge0$ on $[1,K]$ is nondecreasing up to some point and nonincreasing after it, then $\abs{\sum_{j=1}^Kf(j)-\int_1^Kf}\le\max f$.
\end{lemma}

\begin{proof}
We prove the second statement first. Let $f$ increase on $[1,m]$ and decrease on $[m,K]$, and let $p=\floor m$. On each $[j,j+1]$ with $j<p$ the integral of $f$ lies between $f(j)$ and $f(j+1)$, with $f(j)$ the smaller; for $j>p$ the same holds with $f(j+1)$ the smaller; and on $[p,p+1]$ the integral lies between $\min(f(p),f(p+1))$ and $\max f$. Adding the lower bounds gives $\int_1^Kf\ge\sum_jf(j)-\max(f(p),f(p+1))$, and adding the upper bounds gives $\int_1^Kf\le\sum_jf(j)-f(1)-f(K)+\max f$.

For the first statement take $f(x)=x^{-1}e^{-b/x}$, which increases up to $x=b$ and decreases afterwards, with maximum $f(b)=1/(eb)$. The substitution $u=b/x$ gives $\int_1^Kf=\int_{b/K}^bu^{-1}e^{-u}\dd u=E_1(b/K)-E_1(b)$. The function $E_1(z)+\ln z$ has derivative $(1-e^{-z})/z$, so $E_1(z)=-\gamma-\ln z+\operatorname{Ein}(z)$ with $\operatorname{Ein}(z)=\int_0^z(1-e^{-u})u^{-1}\dd u$, where the constant is the classical limit $\lim_{z\to0}(E_1(z)+\ln z)=-\gamma$. Since $0\le(1-e^{-u})/u\le1$, we have $0\le\operatorname{Ein}(b/K)\le b/K$, and $0\le E_1(b)\le b^{-1}\int_b^\infty e^{-u}\dd u=e^{-b}/b$. Hence $\int_1^Kf=\ln(K/b)-\gamma+\operatorname{Ein}(b/K)-E_1(b)$, and adding the error $\abs{E_K(b)-\int_1^Kf}\le1/(eb)$ of the second statement gives both bounds.
\end{proof}

\begin{theorem}[Zipf frequencies]\label{thm:transformers:zipf}
Consider the square-loss table with Zipfian frequencies $q_j=1/(jZ_K)$, $K\ge2$, started from rows with $\norm{w_j(0)-y_j}_2=1$, and run gradient descent with a step size $\eta\in(0,2/q_1)$; write $\Loss_t$ for the loss after $t$ steps, so that $\Loss_0=\frac12$.
\begin{enumerate}
\item For every $\eta$, $\Loss_t=\frac12\sum_{j=1}^Kq_j(1-\eta q_j)^{2t}$.
\item For $\eta=1/q_1$ and every $t\ge1$,
\[
  \frac{\Loss_t}{\Loss_0}=1-\frac{\cA(t)-R_K(t)}{Z_K},\qquad \cA(t)=\sum_{j\ge1}\frac1j\Bigl[1-\Bigl(1-\frac1j\Bigr)^{2t}\Bigr],
\]
with $0\le R_K(t)\le2t/K$ and $\abs{\cA(t)-\ln(2t)-2\gamma}\le3/(2t)$. Equivalently, $\Loss_t/\Loss_0=\bigl(\ln\frac K{2t}-\gamma+r\bigr)/Z_K$ with $\abs r\le\frac3{2t}+\frac{2t}K$.
\item For every $\eta\in(0,2/q_1)$ and every $t\ge1$, $\Loss_t/\Loss_0\ge\bigl(\ln\frac K{4t}-\gamma-\frac3{2t}\bigr)/Z_K$.
\item Consequently, for every $\eps\in(0,1)$ and every stable step size, the number $T_\eps$ of steps needed to reach $\Loss_t\le\eps\Loss_0$ satisfies $T_\eps\ge\frac K4e^{-\gamma-\eps Z_K-3/(2T_\eps)}$; as $K\to\infty$ with $\eps$ fixed, $T_\eps\ge(1-o(1))\frac{e^{-\gamma}}4Ke^{-\eps Z_K}$, which is of order $K^{1-\eps}$.
\end{enumerate}
\end{theorem}

The amount $\cA(t)$ removed from the numerator does not depend on $K$, and it grows by $\ln10$ per decade of $t$: \emph{between steps $t$ and $10t$, for $1\ll t\ll K$, gradient descent removes the fraction $\ln10/Z_K$ of the initial loss}, that is $\ln10/(2Z_K)$ in absolute terms, the same amount in every decade, until $t$ approaches $K$ and the rarest rows are reached.

\begin{proof}
The plan is to reduce the loss to a harmonic sum with geometric weights, compare those weights with exponential ones, and apply \cref{lem:transformers:harmonic}. By \cref{prop:transformers:clock}(i), row $j$ follows gradient descent with step size $\eta q_j$ on $\frac12\norm{w-y_j}_2^2$, which multiplies its error by $1-\eta q_j$ at every step; this gives (i).

For (ii), $\eta=1/q_1=Z_K$ gives $\eta q_j=1/j$, so $\Loss_t/\Loss_0=F_K(t)/Z_K$ with $F_K(t)=\sum_{j=1}^Kj^{-1}(1-1/j)^{2t}$. Writing $Z_K=\sum_{j\le K}1/j$ and splitting off the tail,
\[
  F_K(t)=Z_K-\cA(t)+R_K(t),\qquad R_K(t)=\sum_{j>K}\frac1j\Bigl[1-\Bigl(1-\frac1j\Bigr)^{2t}\Bigr].
\]
Bernoulli's inequality gives $0\le1-(1-1/j)^{2t}\le2t/j$, so the series for $\cA(t)$ converges and $0\le R_K(t)\le\sum_{j>K}2t/j^2\le2t/K$. To find the size of $F_K(t)$, compare it with $E_K(2t)$. For $j\ge2$ put $x=1/j\le\frac12$; the inequalities $-x-x^2\le\ln(1-x)\le-x$, valid for $x\in[0,\frac12]$, give
\[
  e^{-2t/j}\Bigl(1-\frac{2t}{j^2}\Bigr)\le e^{-2t/j-2t/j^2}\le\Bigl(1-\frac1j\Bigr)^{2t}\le e^{-2t/j},
\]
and the term $j=1$ of $F_K$ is zero. Hence $0\le E_K(2t)-F_K(t)\le e^{-2t}+\sum_{j\ge2}g(j)$ with $g(x)=2tx^{-3}e^{-2t/x}$. The function $g$ increases up to $x=2t/3$ and decreases afterwards, its maximum is $\frac{27}{4t^2}e^{-3}<0.34/t^2$, and $\int_0^\infty g=1/(2t)$; by the second part of the lemma, $\sum_{j\ge2}g(j)\le\frac1{2t}+\frac{0.34}{t^2}$. Combining this with the first part of the lemma at $b=2t$,
\[
  -\frac1{2et}-\frac{e^{-2t}}{2t}-e^{-2t}-\frac1{2t}-\frac{0.34}{t^2}\le F_K(t)-\ln\frac K{2t}+\gamma\le\frac1{2et}+\frac{2t}K .
\]
For $t\ge1$ the left side is at least $-1.25/t$, because $te^{-2t}\le e^{-2}$ and $e^{-2t}\le e^{-2}$ there. This is the second form of (ii), with $\abs r\le\frac{1.25}t+\frac{2t}K$. For the first form, let $K\to\infty$ in $\cA(t)=Z_K-F_K(t)+R_K(t)$: since $Z_K-\ln K\to\gamma$ and $R_K(t)\to0$, the bounds on $F_K(t)-\ln K$ pass to the limit and give $\abs{\cA(t)-\ln(2t)-2\gamma}\le1.25/t$.

For (iii), let $j\ge2$. Then $0<\eta q_j<2q_j/q_1=2/j\le1$, so $(1-\eta q_j)^{2t}\ge(1-2/j)^{2t}\ge0$, and dropping the terms $j\le2$,
\[
  \frac{\Loss_t}{\Loss_0}\ge\frac1{Z_K}\sum_{j=3}^K\frac1j\Bigl(1-\frac2j\Bigr)^{2t}\ge\frac1{Z_K}\Bigl[E_K(4t)-\sum_{j\le3}\frac{e^{-4t/j}}j-\sum_{j\ge4}\frac{8t}{j^3}e^{-4t/j}\Bigr],
\]
where the second inequality uses $\ln(1-x)\ge-x-x^2$ with $x=2/j\le\frac12$ for $j\ge4$, and $(1-2/3)^{2t}\ge0$ for $j=3$. The last sum is at most $\frac1{2t}+\frac{0.17}{t^2}$, by the second part of the lemma applied to $8tx^{-3}e^{-4t/x}$; the middle sum is at most $\frac{11}6e^{-4t/3}\le0.49/t$ for $t\ge1$; and the first part of the lemma at $b=4t$ gives $E_K(4t)\ge\ln\frac K{4t}-\gamma-\frac1{4et}-\frac{e^{-4t}}{4t}$. The errors add up to less than $1.3/t\le3/(2t)$.

For (iv), if $\Loss_t\le\eps\Loss_0$, then (iii) gives $\ln\frac K{4t}-\gamma-\frac3{2t}\le\eps Z_K$, which rearranges to $t\ge\frac K4e^{-\gamma-\eps Z_K-3/(2t)}$. Applied at $t=T_\eps$ this is the first claim, and since $T_\eps\ge1$ it gives $T_\eps\ge\frac K4e^{-\gamma-\eps Z_K-3/2}\to\infty$; hence $e^{-3/(2T_\eps)}\to1$. Finally $Z_K=\ln K+\gamma+O(1/K)$, so $Ke^{-\eps Z_K}=e^{-\eps\gamma}K^{1-\eps}(1+o(1))$.
\end{proof}

The constant in (ii) is not sharp: numerically $\cA(t)-\ln(2t)-2\gamma$ is close to $1/(4t)$. Part (iii) says that the logarithmic decay is not an artifact of the step size: the best stable step size gains at most a factor two in time, a shift of $\log_{10}2\approx0.3$ decades, and part (iv) says that the number of steps grows almost linearly with $K$. The asymptotic form of (iv) converges slowly, with relative errors of order $K^{-\eps}$; for $K=10^4$ and $\eps=0.1$ the exact formula (i) gives $1361$ steps at $\eta=1/q_1$ and $683$ at the best constant step size, against the lower bound $526$ of (iv). A scale-invariant optimizer sees none of this.

\begin{corollary}[Sign descent and Adam on a table]\label[corollary]{cor:transformers:sign}
Consider the square-loss table with arbitrary frequencies, started with the same error $e=w_j(0)-y_j\in\R^r$ in every row. For every coordinatewise scale-invariant optimizer, $\Loss_t/\Loss_0$ is the same for all frequency vectors and all $K$: it is the normalized loss of the optimizer on the single row $\frac12\norm{w-y}_2^2$ from the error $e$. For sign descent with step size $\eta$, the $i$-th coordinate of the error has absolute value $\abs{e_i}-t\eta$ for $t\le\floor{\abs{e_i}/\eta}$, and afterwards it alternates between two values whose sum is $\eta$ (or stays at $0$, if $\abs{e_i}/\eta$ is an integer). Hence $\Loss_t/\Loss_0\le r\eta^2/\norm e_2^2$ for all $t>\max_i\abs{e_i}/\eta$, whatever the frequencies.
\end{corollary}

\begin{proof}
A coordinatewise scale-invariant optimizer acts blockwise for the blocks given by the rows and is invariant under rescalings that are constant on rows, so \cref{prop:transformers:clock}(ii) applies: every row follows the solo run on $\frac12\norm{w-y_j}_2^2$, whose error trajectory is the same for all rows, so that $\Loss_t=\sum_jq_j\ell(t)$ for the solo loss $\ell(t)$, and $\Loss_t/\Loss_0=\ell(t)/\ell(0)$. For sign descent the gradient of the solo row is its error, so each coordinate of the error is updated by $e_i\leftarrow e_i-\eta\sign(e_i)$: its size decreases by $\eta$ while it is at least $\eta$; once it is below $\eta$ the next step flips its sign and maps the size $\rho$ to $\eta-\rho$, and so on, unless it hits $0$ exactly, where it stays. In either case it is at most $\eta$ after $\floor{\abs{e_i}/\eta}+1$ steps.
\end{proof}

So on the Zipf table with $K=10^4$, gradient descent needs over a thousand steps to reach a tenth of the initial loss, and sign descent with a step size of $0.005$ needs $55$, for every $K$ (\cref{fig:transformers:zipf}). The corollary removes the frequencies, not the spread of the targets within a row, and the latter matters for sign descent with a constant step: when the targets are themselves Zipfian, as the conditional distributions of a bigram model are, the step must be small enough for the smallest targets and the run long enough for the largest. \Textcite{KunstnerBach2025} find that sign descent then needs a number of steps that grows only like the square root of the dimension, far fewer than gradient descent, whose count grows almost linearly. Exercise~\ref{ex:transformers:signzipf} derives the square root for one row, and shows that it concerns the loss at consecutive steps: sign descent oscillates, and its loss at even steps alone behaves differently.

The square loss is a surrogate. For the softmax table the frequency enters in the same way, exactly; for the output layer, where the frequencies are those of the labels, it enters as the curvature at the optimum.

\begin{proposition}[Frequencies in the softmax]\label[proposition]{prop:transformers:softmax}
\begin{enumerate}
\item For the softmax table, with $s_j=\softmax(w_j)$,
\[
  \nabla_{w_j}\Loss=q_j\,(s_j-\nu_j),\qquad \nabla^2_{w_j}\Loss=q_j\bigl(\diag(s_j)-s_js_j^\top\bigr),
\]
the Hessian is block diagonal over the rows, and $0\preceq\diag(s)-ss^\top\preceq\frac12I$ for every probability vector $s$. Hence \cref{prop:transformers:clock} applies with weights $q_j$: gradient descent runs row $j$ at step size $\eta q_j$, and coordinatewise scale-invariant optimizers do not see the frequencies.
\item For the output logits $z\in\R^K$ of classes with frequencies $q_1>q_2>\dots>q_K>0$, the loss $\ell(z)=-\sum_kq_k\log\softmax(z)_k$ is minimized exactly at $z=\log q+c\mathbf1$, $c\in\R$, where its Hessian is $\diag(q)-qq^\top$. The eigenvalues of this matrix are $0$, with eigenvector $\mathbf1$, and $\mu_1\ge\dots\ge\mu_{K-1}$ with $q_{k+1}\le\mu_k\le q_k$. For Zipfian $q$ the condition number $\mu_1/\mu_{K-1}$ on the subspace orthogonal to $\mathbf1$ lies between $(K-1)/2$ and $K$.
\end{enumerate}
\end{proposition}

\begin{proof}
(i) Since $\sum_k\nu_{jk}=1$, $\ell_j(w)=-\sum_k\nu_{jk}w_k+\log\sum_ke^{w_k}$, whose gradient is $s_j-\nu_j$; differentiating $s_k=e^{w_k}/\sum_le^{w_l}$ gives $\partial s_k/\partial w_l=s_k(\delta_{kl}-s_l)$, the Hessian. Row $j$ of $\Loss$ depends on $w_j$ alone, so the Hessian is block diagonal. For $u\in\R^r$, $u^\top(\diag(s)-ss^\top)u=\sum_ks_ku_k^2-(\sum_ks_ku_k)^2$ is the variance of $u_k$ when $k$ is drawn from $s$, which is nonnegative. Its values lie between $\min_ku_k$ and $\max_ku_k$, so by Popoviciu's inequality the variance is at most $\frac14(\max_ku_k-\min_ku_k)^2\le\frac12(\max_ku_k^2+\min_ku_k^2)\le\frac12\norm u_2^2$.

(ii) The gradient $\softmax(z)-q$ vanishes exactly when $\softmax(z)=q$, that is, when $z-\log q$ is a multiple of $\mathbf1$, and $\ell$ is convex by (i), so these points are the minimizers; the Hessian there is $\diag(q)-qq^\top$ by (i). A vector $u$ with $(\diag(q)-qq^\top)u=0$ has $q_ku_k=q_k\ip{q,u}$ for all $k$, so $u\propto\mathbf1$. For the interlacing use the Courant--Fischer characterization $\mu_k=\max_{\dim U=k}\min_{u\in U,\norm u_2=1}u^\top(\diag(q)-qq^\top)u$. Since $-qq^\top\preceq0$, $\mu_k\le\lambda_k(\diag(q))=q_k$. For the other side, the subspace of vectors in $\operatorname{span}(e_1,\dots,e_{k+1})$ orthogonal to $q$ has dimension at least $k$, and on it the quadratic form equals $\sum_{i\le k+1}q_iu_i^2\ge q_{k+1}\norm u_2^2$, so $\mu_k\ge q_{k+1}$. For Zipfian frequencies, $\mu_1/\mu_{K-1}$ lies between $q_2/q_{K-1}=(K-1)/2$ and $q_1/q_K=K$.
\end{proof}

Near its optimum the output layer is therefore a quadratic whose curvatures interlace the frequencies, the square-loss table of \cref{thm:transformers:zipf} up to a rank-one correction; at initialization, with uniform predictions, the gradient $1/K-q_k$ of a frequent class carries its frequency. This is the structure that \textcite{KunstnerEtAl2024} find in linear models, imbalanced gradients and Hessians correlated with each other, and it produces their frequency-bucket picture. Group the rows into buckets of comparable frequency $q$. By \cref{prop:transformers:clock}, gradient descent fits a bucket at the time $1/(\eta q)$, one decade of time per decade of frequency, while a scale-invariant optimizer fits all buckets together (\cref{fig:transformers:zipf}, right).\sidenote{\Textcite{YadavEtAl2025}, a workshop paper, prove in a minimal next-token-prediction model with heavy-tailed classes that sign descent converges faster than normalized gradient descent.}

\section{The theory: blocks and batches}\label{sec:transformers:blocks}

Rows are blocks of a special kind: they do not interact, and each has its own weight. The block quadratic keeps the first property and replaces the weights by arbitrary curvature scales. For a block $H_b$ write $\mu_b=\lmin(H_b)$, $L_b=\lmax(H_b)$ and $\kappa_b=L_b/\mu_b$; for the whole Hessian write $L=\max_bL_b$, $\mu=\min_b\mu_b$ and $\kappa=L/\mu$; and call $\omega=\max_bL_b/\min_bL_b$ the \term{scale spread}\index{scale spread} of the blocks.

\begin{theorem}[Block heterogeneity]\label{thm:transformers:blocks}
Let $\Loss(\theta)=\frac12(\theta-\theta^\star)^\top H(\theta-\theta^\star)$ with $H=\diag(H_1,\dots,H_N)\succ0$.
\begin{enumerate}
\item Gradient descent with step size $\eta$ satisfies $\max_{\theta_0\ne\theta^\star}\norm{\theta_t-\theta^\star}_2/\norm{\theta_0-\theta^\star}_2=\rho(\eta)^t$ with $\rho(\eta)=\max(\abs{1-\eta L},\abs{1-\eta\mu})$, and $\min_\eta\rho(\eta)=(\kappa-1)/(\kappa+1)$, attained at $\eta=2/(L+\mu)$.
\item Gradient descent with a step size $\eta_b$ for each block has worst-case rate $\max_b\max(\abs{1-\eta_bL_b},\abs{1-\eta_b\mu_b})$, and the best choice of the $\eta_b$ gives $\max_b(\kappa_b-1)/(\kappa_b+1)$.
\item $\max(\omega,\max_b\kappa_b)\le\kappa\le\omega\max_b\kappa_b$.
\item If $H_b=s_b\bar H_b$ with $s_b>0$, then every optimizer that acts blockwise and is scale invariant produces the same iterates for all choices of the scales $s_b$.
\end{enumerate}
\end{theorem}

\begin{proof}
(i) Expand $\theta_0-\theta^\star$ in an orthonormal eigenbasis of $H$; the component along an eigenvector with eigenvalue $\lambda_k$ is multiplied by $1-\eta\lambda_k$ at each step. The norm ratio is therefore at most $\max_k\abs{1-\eta\lambda_k}^t$, with equality for an eigenvector attaining the maximum. The function $\lambda\mapsto\abs{1-\eta\lambda}$ is convex, so its maximum over the spectrum is attained at $\mu$ or at $L$. Minimizing over $\eta$, the decreasing branch $1-\eta\mu$ and the increasing branch $\eta L-1$ cross at $\eta=2/(L+\mu)$, where both equal $(L-\mu)/(L+\mu)=(\kappa-1)/(\kappa+1)$.

(ii) The blocks evolve independently, block $b$ by the iteration of (i) with the spectrum of $H_b$ and the step size $\eta_b$; the worst case over $\theta_0$ is the worst block, and each block's rate is minimized separately by (i).

(iii) Since $\mu_b\le L_b$ for every $b$, $\kappa=\max_bL_b/\min_b\mu_b\ge\max_bL_b/\min_bL_b=\omega$, and $\kappa\ge L_b/\mu_b$ for every $b$. For the upper bound let $b^*$ attain $\min_b\mu_b$; then $\kappa=(\max_bL_b/L_{b^*})(L_{b^*}/\mu_{b^*})\le\omega\kappa_{b^*}$.

(iv) This is \cref{prop:transformers:clock}(ii) with $\ell_b(\theta_b)=\frac12(\theta_b-\theta_b^\star)^\top\bar H_b(\theta_b-\theta_b^\star)$.
\end{proof}

To reduce the error by a factor $\eps$, gradient descent needs about $\frac\kappa2\ln\frac1\eps$ steps and blockwise step sizes about $\frac12\max_b\kappa_b\ln\frac1\eps$; by (iii) the gain $\kappa/\max_b\kappa_b$ is at most the scale spread $\omega$, and it is $\omega$ when the blocks are well conditioned inside. Part (iii) also says what block heterogeneity is. Take the same spectrum twice: once distributed so that each block holds a narrow band at its own scale, and once so that every block holds the whole range. Gradient descent cannot tell the two apart, since its rate depends on the spectrum alone, but blockwise step sizes gain the factor $\omega$ in the first arrangement and nothing in the second, where $\omega=1$. In the language of \cref{obs:transformers:blocks}, transformers are the first arrangement and convolutional networks the second. Part (iv) says that a scale-invariant optimizer, Adam among them, captures the gain without being told the scales; it does not say how fast such an optimizer is within a block, which depends on the method.

Both invariances are tied to the coordinate system. Gradient descent is rotation equivariant: for an orthogonal $Q$, gradient descent on $\theta\mapsto\Loss(Q\theta)$ produces $Q^\top$ times the iterates on $\Loss$, so its rate is unchanged. Scale invariance is not preserved: after a rotation that mixes rows or blocks, the gradient of a coordinate no longer carries a single weight, and for sign descent the smoothness constant in the $\ell_\infty$ norm of a quadratic can grow from $\tr H$ to about $\frac2\pi d\lmax(H)$ (\cref{ch:sign}). This is the effect that \textcite{XieEtAl2025} measure in GPT-2, and Exercise~\ref{ex:transformers:rotation} shows it on the Zipf table.

The last model adds noise. The advantage of a preconditioner has two limits, one for large and one for small batches, and they differ. In the small-batch limit the relevant quantity is a number of samples, defined through the two functions
\begin{equation}\label{eq:transformers:BV}
  \cB_a(\tau)=\frac12\sum_{i=1}^dh_i\theta_{0,i}^2e^{-2\lambda_i\tau},\qquad \cV_a(\tau)=\frac14\sum_{i=1}^da_ic_i\bigl(1-e^{-2\lambda_i\tau}\bigr),
\end{equation}
the bias and the noise, per unit of $\eta/B$, of the gradient flow in the time $\tau=\eta t$; we write $\Phi_a(\tau,\alpha)=\cB_a(\tau)+\alpha\cV_a(\tau)$, and $\Delta_{0,i}=\frac12h_i\theta_{0,i}^2$, so that $\sum_i\Delta_{0,i}=\Loss(\theta_0)$.

\begin{lemma}[Small steps]\label[lemma]{lem:transformers:discrete}
In the noisy quadratic model, let $\lmax=\max_i\lambda_i$ and $0<\eta\le1/(2\lmax)$, and put $\alpha=\eta/B$. Then for every $t$,
\[
  \E\Loss(\theta_t)=\sum_{i=1}^d\Bigl[\Delta_{0,i}u_i+\frac{\alpha a_ic_i}4\,\frac{2(1-u_i)}{2-\eta\lambda_i}\Bigr],\qquad u_i=(1-\eta\lambda_i)^{2t},
\]
and $\abs{\E\Loss(\theta_t)-\Phi_a(\eta t,\alpha)}\le\eta\lmax\bigl(\Loss(\theta_0)+\alpha\sum_ia_ic_i\bigr)$. The formula holds for every $\eta\in(0,2/\lmax)$.
\end{lemma}

\begin{proof}
Squaring $\theta_{t+1,i}=(1-\eta\lambda_i)\theta_{t,i}-\eta a_i\zeta_{t,i}$ and taking expectations, with $\zeta_t$ independent of $\theta_t$ and of mean zero, gives $\E\theta_{t+1,i}^2=(1-\eta\lambda_i)^2\E\theta_{t,i}^2+\eta^2a_i^2c_i/B$, so $\E\theta_{t,i}^2=u_i\theta_{0,i}^2+(\eta^2a_i^2c_i/B)(1-u_i)/(1-(1-\eta\lambda_i)^2)$ (\cref{ch:noise}). Multiplying by $\frac12h_i$ and using $1-(1-x)^2=x(2-x)$ with $x=\eta\lambda_i=\eta a_ih_i$ gives the formula. Now fix $i$, put $x=\eta\lambda_i\le\frac12$ and $\bar u=e^{-2xt}$. The inequalities $-x-x^2\le\ln(1-x)\le-x$ give $\bar ue^{-2tx^2}\le u_i\le\bar u$, hence
\[
  0\le\bar u-u_i\le\bar u\,2tx^2=2x\,(xt)e^{-2xt}\le x/e,
\]
because $se^{-2s}\le1/(2e)$ for $s\ge0$. The term of $\Phi_a$ for coordinate $i$ is $\Delta_{0,i}\bar u+\frac14\alpha a_ic_i(1-\bar u)$, and $\frac{2(1-u_i)}{2-x}-(1-\bar u)=\frac x{2-x}(1-u_i)+(\bar u-u_i)$ lies in $[0,(\frac23+\frac1e)x]\subset[0,2x]$. So the two terms differ by at most $x\Delta_{0,i}/e+x\alpha a_ic_i/2\le x(\Delta_{0,i}+\alpha a_ic_i)$, and summing over $i$ with $x\le\eta\lmax$ gives the bound.
\end{proof}

\begin{theorem}[The gap and the batch size]\label{thm:transformers:batch}
In the noisy quadratic model let $\theta_{0,i}\ne0$ and $c_i>0$ for all $i$, let $\lmin$ and $\lmax$ be the extreme effective curvatures, and let $0<\eps<\Loss(\theta_0)$. Let $T_a(B)$ be the least $t$ such that some $\eta\in(0,2/\lmax)$ gives $\E\Loss(\theta_t)\le\eps$, and $T_a(\infty)$ the same count for exact gradients ($c=0$).
\begin{enumerate}
\item (Large batches.) $T_a(B)$ is nonincreasing in $B$ and at least $T_a(\infty)$. If $\eps$ is not one of the countably many numbers $\inf_\eta\Loss(\theta_t)$, $t\in\N$, for exact gradients, then $T_a(B)=T_a(\infty)$ for all large $B$.
\item (Small batches.) As $B\to0$, $B\,T_a(B)$ converges to the number of samples
\[
  n_a(\eps)=\min\Bigl\{\frac{\tau\,\cV_a(\tau)}{\eps-\cB_a(\tau)}\colon\tau>0,\ \cB_a(\tau)<\eps\Bigr\}.
\]
\item (High accuracy.) If $\lmin<\lmax$, then as $\eps\to0$,
\[
  T_a(\infty)\sim\frac{\ln(1/\eps)}{4\operatorname{artanh}(\lmin/\lmax)},\qquad n_a(\eps)\sim\frac{\sum_ia_ic_i}{8\lmin}\cdot\frac{\ln(1/\eps)}\eps .
\]
\end{enumerate}
\end{theorem}

\begin{proof}
The plan is to settle the large-batch limit with the exact formula of \cref{lem:transformers:discrete}, the small-batch limit by matching the discrete dynamics to its flow from above and from below, and the high-accuracy limits by direct estimates.

(i) In the formula of the lemma the noise term is nonnegative and proportional to $1/B$, so $\E\Loss(\theta_t)$ is nonincreasing in $B$ and at least its value for exact gradients; this gives the first two claims. Let $t^*=T_a(\infty)$ and suppose that $\inf_\eta\Loss(\theta_{t^*})<\eps$ for exact gradients, which holds unless $\eps$ equals that infimum. Then some $\eta^*$ gives a deterministic loss below $\eps$ at step $t^*$, the noise term at $(\eta^*,t^*)$ tends to $0$ as $B\to\infty$, and $T_a(B)\le t^*$ for large $B$.

(ii) The infimum defining $n_a(\eps)$ is attained: the quotient is continuous on the open set of admissible $\tau$, which is an interval $(\tau_\eps,\infty)$ because $\cB_a$ decreases continuously from $\Loss(\theta_0)$ to $0$; it tends to $\infty$ as $\tau\downarrow\tau_\eps$, where the denominator vanishes, and as $\tau\to\infty$, where $\cV_a(\tau)$ tends to $\frac14\sum_ia_ic_i>0$.

\emph{Upper bound.} Let $\tau$ be admissible and $0<\alpha<(\eps-\cB_a(\tau))/\cV_a(\tau)$, so that $\Phi_a(\tau,\alpha)<\eps$. Put $\eta=\alpha B$ and $t=\ceil{\tau/\eta}$, so that $\eta t\in[\tau,\tau+\eta]$. As $B\to0$, $\eta\to0$, $\Phi_a(\eta t,\alpha)\to\Phi_a(\tau,\alpha)$ by continuity, and the error bound of the lemma tends to $0$. So for small $B$, $\E\Loss(\theta_t)\le\eps$ and $B\,T_a(B)\le Bt\le\tau/\alpha+B$. Letting $B\to0$, then $\alpha\uparrow(\eps-\cB_a(\tau))/\cV_a(\tau)$, and minimizing over $\tau$ gives $\limsup_{B\to0}B\,T_a(B)\le n_a(\eps)$.

\emph{Lower bound.} For each $B$ let $\eta_B$ attain the count $t_B=T_a(B)$, and put $\alpha_B=\eta_B/B$. First, $\alpha_B$ stays bounded. In the formula of the lemma, $2/(2-\eta\lambda_i)\ge1$ for every stable $\eta$, so the term of coordinate $i$ is a convex combination of $\Delta_{0,i}$ and a number at least $\frac14\alpha a_ic_i$, and is at least $\min(\Delta_{0,i},\frac14\alpha a_ic_i)$. If $\alpha_B$ were unbounded, these lower bounds would sum to nearly $\Loss(\theta_0)>\eps$ along a subsequence; so $\alpha_B\le\bar\alpha$ for some $\bar\alpha$, and $\eta_B\le\bar\alpha B$. For small $B$ the lemma then applies, and with $\tau_B=\eta_Bt_B$,
\[
  \eps\ge\E\Loss(\theta_{t_B})\ge\Phi_a(\tau_B,\alpha_B)-\delta_B,\qquad\delta_B=\bar\alpha B\lmax\bigl(\Loss(\theta_0)+\bar\alpha\textstyle\sum_ia_ic_i\bigr)\to0 .
\]
Hence $\cB_a(\tau_B)+\alpha_B\cV_a(\tau_B)\le\eps+\delta_B$, and $B\,t_B=\tau_B/\alpha_B\ge\tau_B\cV_a(\tau_B)/(\eps+\delta_B-\cB_a(\tau_B))\ge n_a(\eps+\delta_B)$. It remains to show $\liminf_{\delta\downarrow0}n_a(\eps+\delta)\ge n_a(\eps)$. Let $\delta_k\downarrow0$ and let $\tau_k$ attain $n_a(\eps+\delta_k)\le n_a(\eps)$. If $\tau_k\to\infty$ or $\tau_k\to\tau_\eps$ along a subsequence, the quotient at $\tau_k$ tends to $\infty$, as above; so the $\tau_k$ have a limit point in $(\tau_\eps,\infty)$, where the quotient is continuous in $(\tau,\eps)$, and the limit of $n_a(\eps+\delta_k)$ along it is at least $n_a(\eps)$.

(iii) Put $\rho=(\lmax-\lmin)/(\lmax+\lmin)$, so that $\ln(1/\rho)=2\operatorname{artanh}(\lmin/\lmax)$. By the argument of \cref{thm:transformers:blocks}(i), every $\eta$ has $\abs{1-\eta\lambda_i}\ge\rho$ for some $i$ with $\lambda_i\in\{\lmin,\lmax\}$, so the deterministic loss at step $t$ is at least $m\rho^{2t}$, with $m$ the smallest $\Delta_{0,i}$ among those $i$; and $\eta=2/(\lmax+\lmin)$ gives a loss at most $\Loss(\theta_0)\rho^{2t}$. Hence $\ln(m/\eps)/(2\ln(1/\rho))\le T_a(\infty)\le\ln(\Loss(\theta_0)/\eps)/(2\ln(1/\rho))+1$, and both bounds are asymptotic to the claim. For $n_a$, write $V_a=\sum_ia_ic_i$ and let $m_1>0$ be the sum of $\Delta_{0,i}$ over the $i$ with $\lambda_i=\lmin$. Every admissible $\tau$ has $m_1e^{-2\lmin\tau}\le\cB_a(\tau)<\eps$, so $\tau>\tau_1=\ln(m_1/\eps)/(2\lmin)$, and since $\tau\cV_a(\tau)$ increases in $\tau$, $n_a(\eps)\ge\tau_1\cV_a(\tau_1)/\eps$, where $\cV_a(\tau_1)\to V_a/4$. Conversely, at $\tau=(\ln(\Loss(\theta_0)/\eps)+s)/(2\lmin)$ with $s=\ln\ln(1/\eps)$ we have $\cB_a(\tau)\le\eps e^{-s}$ and $\cV_a\le V_a/4$, so $n_a(\eps)\le\frac{V_a}{8\lmin\eps}\cdot\frac{\ln(\Loss(\theta_0)/\eps)+s}{1-e^{-s}}$. Both bounds are $\frac{V_a}{8\lmin}\frac{\ln(1/\eps)}\eps(1+o(1))$.
\end{proof}

The small-batch count has a reading that the formula hides. By \eqref{eq:transformers:BV}, coordinate $i$ of the preconditioned method at time $\tau$ and noise level $\alpha$ is coordinate $i$ of SGD at time $a_i\tau$ and noise level $a_i\alpha$, and the number of samples $\tau/\alpha$ is the same for both. \emph{In the small-batch limit a diagonal preconditioner does not create samples; it reallocates time among the coordinates at a fixed number of samples.} The reallocation helps when SGD's single time is wrong for some coordinates, and it helps less than at large batch, where the preconditioner shortens the time itself.

\begin{corollary}[The two ends of the gap]\label[corollary]{cor:transformers:gap}
Compare SGD with the preconditioner $a_i=h_i^{-p}$, $0\le p<1$, write $\kappa=h_1/h_d>1$, and assume the hypotheses of \cref{thm:transformers:batch}. As $\eps\to0$,
\[
  \frac{T_1(\infty)}{T_a(\infty)}\to\frac{\operatorname{artanh}(\kappa^{p-1})}{\operatorname{artanh}(\kappa^{-1})}\ge\kappa^p,\qquad \frac{n_1(\eps)}{n_a(\eps)}\to\frac{\sum_ic_i}{\sum_i(h_d/h_i)^pc_i}\in[1,\kappa^p].
\]
For $h_i=1/i$, $c_i=h_i$ and $0<p<1$ the small-batch limit is $pZ_d(1+o(1))$ as $d\to\infty$, while the large-batch limit is at least $d^p$.
\end{corollary}

\begin{proof}
Both limits follow from \cref{thm:transformers:batch}(iii): for SGD $\lmin/\lmax=1/\kappa$, $\lmin=h_d$ and $\sum_ia_ic_i=\sum_ic_i$; for the preconditioner $\lmin/\lmax=\kappa^{p-1}$, $\lmin=h_d^{1-p}$ and $\sum_ia_ic_i=\sum_ih_i^{-p}c_i$. The first limit is at least $\kappa^p$ because $\operatorname{artanh}(x)/x$ increases on $(0,1)$. In the second, $(h_d/h_i)^p\in[\kappa^{-p},1]$ gives the bounds. For $h_i=1/i$ and $c_i=h_i$ the numerator is $Z_d$ and the denominator is $d^{-p}\sum_{i\le d}i^{p-1}$, which tends to $1/p$ by comparison with $\int_0^dx^{p-1}\dd x=d^p/p$.
\end{proof}

So the gap opens with the batch size: in the high-accuracy limit it is $pZ_d\approx p\ln d$ at small batch and $d^p$ at large batch, and between the two it moves the knee of the steps-to-target curve, the \term{critical batch size}\index{critical batch size} $n_a(\eps)/T_a(\infty)$ at which the two asymptotes cross, to the right by the ratio of the two ends (\cref{ch:noise}).\sidenote{\Textcite{McCandlishEtAl2018} fit steps-to-target curves of the form $S_{\min}(1+B_{\mathrm{crit}}/B)$; in our notation $S_{\min}=T_a(\infty)$ and $S_{\min}B_{\mathrm{crit}}=n_a(\eps)$.} The small-batch end is $1$ only if the noise sits in the flattest directions, which is not the natural case; with a constant step size the model does not produce parity at batch size one. One standard ingredient of training changes this.

\begin{heuristic}[Heuristic derivation: averaging removes the small-batch advantage]
In the small-batch limit, coordinate $i$ of preconditioned SGD in the time $\tau=\eta t$ follows, to leading order, the Ornstein--Uhlenbeck process $\dd\theta_i=-\lambda_i\theta_i\dd\tau+a_i\sqrt{\alpha c_i}\dd W_\tau$, with $\alpha=\eta/B$. Its stationary autocovariance is $(\alpha a_i^2c_i/(2\lambda_i))e^{-\lambda_i\abs s}$, whose integral over $s\in\R$ is $\alpha a_i^2c_i/\lambda_i^2$. Average the iterates over a window of length $\tau_w\gg1/\lambda_i$, which uses $n_w=\tau_w/\alpha$ samples: the average has variance about $\alpha a_i^2c_i/(\lambda_i^2\tau_w)=c_i/(h_i^2n_w)$ and contributes $c_i/(2h_in_w)$ to the loss. Neither the multiplier $a_i$ nor the step size appears; this is the variance of the best estimate of $\theta^\star_i=0$ from $n_w$ samples with noise of variance $c_i$ and curvature $h_i$. Once every direction has relaxed, averaged SGD and averaged preconditioned SGD therefore reach the same loss with the same number of samples, and the preconditioner helps only through the burn-in and the relaxation of the flattest directions. Their cost grows like $\ln(1/\eps)$ as the target falls, while the cost of the variance grows like $1/\eps$, so the small-batch gap tends to one.
\end{heuristic}

\section{Reading the theory}\label{sec:transformers:reading}

\paragraph{Logarithmic loss curves.} \Cref{thm:transformers:zipf} predicts that gradient descent on Zipfian data makes the same progress, $\ln10/Z_K$ of the initial loss, in every decade of steps: on a plot against $\log t$ its loss is a straight line, and the line is flatter for larger vocabularies, with slope $0.44$, $0.31$ and $0.24$ per decade for $K=10^2$, $10^3$ and $10^4$ (\cref{fig:transformers:zipf}, left). The asymptotic line $(\ln\frac K{2t}-\gamma)/Z_K$ is within $0.003$ of the exact loss at $t=10$ and $t=100$ for $K=10^4$, and it falls below the exact curve as $t$ approaches $K$, where the term $R_K(t)$, at most $2t/K$, matters and the curve turns into the exponential tail $(1-1/K)^{2t}$ of the rarest rows. Sign descent and Adam fall on one curve for every $K$, as \cref{cor:transformers:sign} requires: the runs on the Zipf table with $K=10^4$ and $K=10^2$ agree with the single-row predictions to rounding error. The shapes differ: sign descent decreases the error of each coordinate linearly and then stops at its floor, while Adam with $\beta_2=0.999$, after its first hundred steps, behaves like gradient descent with a frozen diagonal preconditioner and converges linearly. For a vocabulary of $50{,}000$ tokens the theorem predicts that gradient descent removes $\ln10/Z_K\approx20\%$ of the initial loss per decade.

\begin{figure}[t]
\begin{fullwidth}
\centering
\begin{tikzpicture}
\begin{groupplot}[group style={group size=2 by 1,horizontal sep=17mm},wplot,width=82mm,height=56mm,xmode=log,xmin=1,xmax=1e5,ymin=0,ymax=1.02,log basis x=10]
\nextgroupplot[xlabel={step $t$},ylabel={normalized loss $\Loss_t/\Loss_0$},legend style={at={(0.5,1.03)},anchor=south,font=\scriptsize},legend columns=3]
\addplot[theory3] table[x=t,y=gd2] {figures/data/14-transformers-zipf-curves.dat};
\addplot[theory2] table[x=t,y=gd3] {figures/data/14-transformers-zipf-curves.dat};
\addplot[theory] table[x=t,y=gd4] {figures/data/14-transformers-zipf-curves.dat};
\addplot[muted,line width=0.9pt,densely dotted,no marks] table[x=t,y=sign] {figures/data/14-transformers-zipf-curves.dat};
\addplot[ink,line width=0.7pt,no marks] table[x=t,y=adam] {figures/data/14-transformers-zipf-curves.dat};
\addplot[guide,forget plot] table[x=t,y=asy2] {figures/data/14-transformers-zipf-curves.dat};
\addplot[guide,forget plot] table[x=t,y=asy3] {figures/data/14-transformers-zipf-curves.dat};
\addplot[guide] table[x=t,y=asy4] {figures/data/14-transformers-zipf-curves.dat};
\addplot[simulation,forget plot] table[x=t,y=gd2] {figures/data/14-transformers-zipf-points.dat};
\addplot[simulation,forget plot] table[x=t,y=gd3] {figures/data/14-transformers-zipf-points.dat};
\addplot[simulation] table[x=t,y=gd4] {figures/data/14-transformers-zipf-points.dat};
\addplot[simulation2] table[x=t,y=sign4] {figures/data/14-transformers-zipf-points.dat};
\addplot[only marks,mark=triangle*,mark size=1.9pt,color=ink,mark options={fill=ink!70,draw=white,line width=0.2pt}] table[x=t,y=adam4] {figures/data/14-transformers-zipf-points.dat};
\legend{{GD, $K=10^2$},{GD, $K=10^3$},{GD, $K=10^4$},sign descent,Adam,asymptote,{GD runs},{sign runs},{Adam runs}}
\node[font=\scriptsize,text=sandline!55!black,anchor=west] at (axis cs:12,0.13) {$10^2$};
\node[font=\scriptsize,text=seadeep,anchor=west] at (axis cs:160,0.13) {$10^3$};
\node[font=\scriptsize,text=sea,anchor=west] at (axis cs:2600,0.13) {$10^4$};
\nextgroupplot[xlabel={step $t$},ylabel={loss of the bucket / its initial loss},legend style={at={(0.5,1.03)},anchor=south,font=\scriptsize},legend columns=3]
\addplot[theory] table[x=t,y=b1] {figures/data/14-transformers-zipf-buckets.dat};
\addplot[theory,forget plot] table[x=t,y=b2] {figures/data/14-transformers-zipf-buckets.dat};
\addplot[theory,forget plot] table[x=t,y=b3] {figures/data/14-transformers-zipf-buckets.dat};
\addplot[theory,forget plot] table[x=t,y=b4] {figures/data/14-transformers-zipf-buckets.dat};
\addplot[ink,line width=0.7pt,no marks] table[x=t,y=adam] {figures/data/14-transformers-zipf-buckets.dat};
\addplot[simulation] table[x=t,y=gdb1] {figures/data/14-transformers-zipf-bucketpts.dat};
\addplot[simulation,forget plot] table[x=t,y=gdb2] {figures/data/14-transformers-zipf-bucketpts.dat};
\addplot[simulation,forget plot] table[x=t,y=gdb3] {figures/data/14-transformers-zipf-bucketpts.dat};
\addplot[simulation,forget plot] table[x=t,y=gdb4] {figures/data/14-transformers-zipf-bucketpts.dat};
\addplot[only marks,mark=triangle*,mark size=1.9pt,color=ink,mark options={fill=ink!70,draw=white,line width=0.2pt}] table[x=t,y=adamb4] {figures/data/14-transformers-zipf-bucketpts.dat};
\legend{{GD, exact},Adam,{GD runs},{Adam runs}}
\node[font=\scriptsize,text=sea,anchor=east] at (axis cs:3.2,0.5) {$j<10$};
\node[font=\scriptsize,text=sea,anchor=west] at (axis cs:30,0.62) {$10^1$--$10^2$};
\node[font=\scriptsize,text=sea,anchor=west] at (axis cs:300,0.62) {$10^2$--$10^3$};
\node[font=\scriptsize,text=sea,anchor=west] at (axis cs:3000,0.62) {$10^3$--$10^4$};
\end{groupplot}
\end{tikzpicture}
\end{fullwidth}
\caption{Gradient descent, sign descent and Adam on the square-loss table with Zipfian frequencies, every row starting at distance one from its target (rows of $16$ coordinates, the Gaussian quantiles rescaled to unit norm). Left: the normalized loss. Curves: the exact formula of \cref{thm:transformers:zipf}(i) for gradient descent with $\eta=1/q_1$ ($K=10^2,10^3,10^4$), its asymptote $(\ln\frac K{2t}-\gamma)/Z_K$ (dashed), and the frequency-free predictions of \cref{cor:transformers:sign}: the exact sign-descent loss of one row ($\eta=0.005$) and Adam run on one row ($\eta=0.005$, $\beta_1=0.9$, $\beta_2=0.999$, $\epsilon=0$). Points: the optimizers run on the full table ($K=10^4$ for sign descent and Adam). Right: for $K=10^4$, the loss of the rows in each frequency bucket $j\in[10^k,10^{k+1})$, each relative to its own initial loss; gradient descent fits one bucket per decade, Adam fits all four together (its four buckets coincide; the points of the rarest are shown). Deterministic; script \texttt{14-transformers-zipf.py}.}
\label{fig:transformers:zipf}
\end{figure}

\paragraph{The staircase.} The right panel of \cref{fig:transformers:zipf} is the frequency-bucket picture of \cref{obs:transformers:rare} in a model where it is exact: the bucket of rows with frequency near $q$ is fitted around $t=1/(\eta q)$, each bucket a decade after the previous one, while Adam fits all buckets together. Most of the examples come from the rare buckets: under Zipf's law each bucket $[10^k,10^{k+1})$ with $k\ge1$ carries about the fraction $\ln10/Z_K$ of the data, the same for every decade, so that the average loss is dominated by buckets that gradient descent has not reached yet. The picture predicts that gradient descent's disadvantage concentrates in the rows that serve rare tokens, the embedding and output rows. It does not predict the role of the normalization gains in the findings of \textcite{ZhaoEtAl2025}, which is better described as a block effect, the subject of \cref{thm:transformers:blocks}.

\paragraph{The gap and the batch size.} \Cref{fig:transformers:batch} shows the batch-size theory on an example in the spirit of \textcite{ZhangEtAl2019}: $d=1000$, $h_i=1/i$, $c_i=h_i$, $\theta_0=\mathbf1$, and a target of $5\%$ of the initial loss. SGD needs $10065$, $2577$, $723$, $308$, $221$ and $184$ steps at $B=1,4,16,64,256$ and $4096$, and the preconditioner $a_i=h_i^{-1/2}$ needs $3393$, $849$, $214$, $55$, $16$ and $10$. At small batches both curves follow their asymptotes $n_a(\eps)/B$ of \cref{thm:transformers:batch}(ii), with $n_1(\eps)=9987$ and $n_a(\eps)=3390$ samples; at large batches they level off near the deterministic counts $177$ and $10$. The advantage therefore grows from $2.95$ to $17.7$, and the critical batch size $n_a(\eps)/T_a(\infty)$ grows from $56$ to $339$. In the high-accuracy limit of \cref{cor:transformers:gap} the two ends would be $3.83$ and $31.6$. The Monte Carlo runs agree with the exact counts to within $2\%$, or one step, at every batch size. So do runs of RMSProp, that is, Adam with $\beta_1=0$ and $\beta_2=0.999$: up to $B=64$ they need $3404$, $865$, $210$ and $54$ steps, the counts of the $H^{1/2}$ preconditioner. Beyond that the gradient dominates Adam's second moment for the stiff coordinates, Adam becomes sign-like, and the model $a_i=h_i^{-1/2}$ no longer describes it.

\begin{figure}[t]
\begin{fullwidth}
\centering
\begin{tikzpicture}
\begin{groupplot}[group style={group size=2 by 1,horizontal sep=17mm},wplot,width=82mm,height=58mm,xmode=log,log basis x=2,xmin=1,xmax=8192,xtick={1,4,16,64,256,1024,4096},xticklabels={$1$,$4$,$16$,$64$,$256$,$1024$,$4096$}]
\nextgroupplot[xlabel={batch size $B$},ylabel={steps to target $T(B)$},ymode=log,ymin=5,ymax=2e4,legend style={at={(0.5,1.03)},anchor=south,font=\scriptsize},legend columns=3]
\addplot[theory] table[x=B,y=sgd] {figures/data/14-transformers-batch.dat};
\addplot[theory2] table[x=B,y=pre] {figures/data/14-transformers-batch.dat};
\addplot[guide] table[x=B,y=sgdsmall] {figures/data/14-transformers-batch.dat};
\addplot[guide,forget plot] table[x=B,y=presmall] {figures/data/14-transformers-batch.dat};
\addplot[guide,forget plot] table[x=B,y=sgdbig] {figures/data/14-transformers-batch.dat};
\addplot[guide,forget plot] table[x=B,y=prebig] {figures/data/14-transformers-batch.dat};
\addplot[simulation] table[x=B,y=mcsgd] {figures/data/14-transformers-batch-mc.dat};
\addplot[simulation2] table[x=B,y=mcpre] {figures/data/14-transformers-batch-mc.dat};
\addplot[only marks,mark=triangle*,mark size=2.1pt,color=ink,mark options={fill=ink!70,draw=white,line width=0.2pt}] table[x=B,y=mcrms] {figures/data/14-transformers-batch-mc.dat};
\legend{SGD,{$a_i=h_i^{-1/2}$},asymptotes,{SGD runs},{preconditioned runs},{RMSProp runs}}
\nextgroupplot[xlabel={batch size $B$},ylabel={advantage $T_{\mathrm{SGD}}(B)/T_a(B)$},ymode=log,ymin=0.7,ymax=40,ytick={1,2,5,10,20},yticklabels={$1$,$2$,$5$,$10$,$20$},legend style={at={(0.5,1.03)},anchor=south,font=\scriptsize},legend columns=3]
\addplot[theory] table[x=B,y=rconst] {figures/data/14-transformers-batch.dat};
\addplot[theory3] table[x=B,y=rlin] {figures/data/14-transformers-batch.dat};
\addplot[theory2] table[x=B,y=ravg] {figures/data/14-transformers-batch.dat};
\addplot[muted,line width=0.9pt,dashdotted,no marks] table[x=B,y=riso] {figures/data/14-transformers-batch.dat};
\addplot[simulation] table[x=B,y expr=\thisrow{mcsgd}/\thisrow{mcpre}] {figures/data/14-transformers-batch-mc.dat};
\addplot[guide,domain=1:8192,forget plot] {2.946};
\addplot[guide,domain=1:8192,forget plot] {17.7};
\legend{constant step,linear decay,tail average,{isotropic noise},runs}
\end{groupplot}
\end{tikzpicture}
\end{fullwidth}
\caption{Steps to target in the noisy quadratic model ($d=1000$, $h_i=1/i$, $c_i=h_i$, $\theta_0=\mathbf1$, target $5\%$ of the initial loss), SGD against the preconditioner $a_i=h_i^{-1/2}$, each at its best constant step size. Left: the exact counts (curves; step sizes optimized over $400$ values), the small-batch asymptotes $n_a(\eps)/B$ and the large-batch counts $T_a(\infty)$ of \cref{thm:transformers:batch} (dashed), and Monte Carlo runs at the optimal step sizes (points; $32$ runs, the first step at which the run average, smoothed over the last $2\%$ of the steps, is below the target; without the smoothing, fluctuations of the stiff coordinates make SGD's first passage up to $10\%$ early at small $B$). Triangles: RMSProp ($\beta_2=0.999$, $16$ runs, learning rate tuned over four values), shown where its second moment is dominated by noise. Right: the advantage of the preconditioner, exact, for a constant step (with the two limits of the left panel, dashed), for a linear decay of the step size to zero, for the average of the last iterates (fraction optimized), and for isotropic noise $c_i=\frac1d\sum_jh_j$ with a constant step. Script \texttt{14-transformers-batch.py}, seeds fixed.}
\label{fig:transformers:batch}
\end{figure}

\paragraph{What closes the gap at small batch.} The right panel of \cref{fig:transformers:batch} answers the question whether the small-batch advantage tends to one. With a constant step size it does not: it tends to $2.95$ here, and to $3.83$, close to $pZ_d=3.74$, as the target falls (\cref{cor:transformers:gap}). It is close to one in two cases. If the noise is isotropic, $c_i$ constant, the advantage at small batch is $1.11$, and it tends to $1+p=1.5$ as the target falls (Exercise~\ref{ex:transformers:noise}). If the iterates are averaged, as the heuristic derivation of \cref{sec:transformers:blocks} predicts, the advantage at $B=1$ is $0.99$: both methods need about $2400$ steps, against $10065$ and $3393$ for their last iterates. At lower targets the averaged advantage stays at one ($0.998$ at $1\%$, $1.000$ at $0.1\%$), and at $0.1\%$ both averaged methods need $133{,}300$ steps, close to the $\sum_ic_i/(2h_i\eps)=d/(2\eps)\approx133{,}600$ samples of the efficient estimator in the heuristic derivation. A linear decay of the step size to zero, which acts much like averaging (\cref{ch:schedules}), gives $1.31$ at $B=1$ and $B=4$. In every case the advantage grows with the batch size and reaches the curvature factor, $18$ to $22$, at the largest batches. The noisy quadratic model therefore reproduces the trend of \cref{obs:transformers:small}, a gap that opens with the batch size, and it reproduces parity at batch size one only with a decaying step or averaging, which practical recipes do use. This partly reconciles the batch-size observation with the curvature explanations: at small batch, progress is limited by samples, and no preconditioner creates samples; at large batch it is limited by curvature, where the curvature explanations apply in full.

\paragraph{Where the agreement must fail.} The table and block models are exact for their own losses and say nothing about rows that share parameters. The batch-size theory compares methods at their best constant step sizes and in expectation; its small-batch limit holds once $\eta\lambda_{\max}$ is small, which here happens already at $B=1$, but would fail for problems with much stiffer directions at the same batch size. And Adam is $H^{1/2}$ preconditioning only while noise dominates its second moment; with momentum $\beta_1>0$ it is, at small batch, the same method with a rescaled step size (\cref{ch:momentum}), but at large batch it is closer to sign descent.

\section{What the theory misses}\label{sec:transformers:misses}

\paragraph{What was idealized away.} \emph{Parameter sharing.} In a transformer the embedding rows of rare tokens feed attention and feed-forward layers that every token shares, so the frequency of a token is not the curvature of a separate block; the table model is exact for a bigram model and for the embedding and output rows, and an approximation for everything above them. \emph{Separability of the mechanisms.} The three models act one at a time, but heavy-tailed classes produce block heterogeneity, since the rows of the embedding and output layers are blocks whose scales are the frequencies; the frequency model is a special case of the block model, and the batch-size model can contain either. \emph{Fixed curvature and Gaussian noise.} The noisy quadratic model fixes the Hessian and the noise; in training both change, and the noise of language models is heavy-tailed \parencite{ZhangEtAl2020}. \emph{Adam's state.} With $\epsilon>0$ a row whose gradients are below $\epsilon$ falls back to SGD with step size $\eta/\epsilon$, which matters for the rarest rows; $\beta_2$ sets a memory, and the time Adam's preconditioner takes to adapt is ignored. \emph{Constant step sizes.} The theorems compare constant step sizes; schedules enter only through the heuristic and the numerics of \cref{fig:transformers:batch}.

\paragraph{Where it breaks.} (1) No single mechanism accounts for the gap. \Textcite{KunstnerEtAl2024} show that class imbalance reproduces the gap in linear models and vision models, and \textcite{ZhangEtAl2024} that block heterogeneity reproduces it in quadratics and accompanies it in networks; neither excludes the other, and the models of this chapter cannot weigh them against each other. (2) The size of the gap at scale depends on tuning: in the comparisons of \textcite{WenEtAl2026}, the speedups of other optimizers over a well-tuned AdamW shrink as the models grow, to about $1.1\times$ at a billion parameters. (3) The batch-size picture rests on a recent refereed study \parencite{MarekEtAl2025} and an unrefereed one \parencite{SreckovicEtAl2025}. With constant steps the noisy quadratic model predicts a residual small-batch advantage of about $p\ln d$, larger than what is reported; it needs averaging or a decaying step to give parity. (4) Adam is not $H^{1/2}$ preconditioning outside the noise-dominated regime, and at full batch it is closer to sign descent, whose rate in the table model is set by the spread of the targets within rows rather than by the frequencies \parencite{KunstnerBach2025}. (5) The block theorem compares best constant step sizes on a fixed quadratic. It says that blockwise steps can gain the scale spread, not that Adam's per-coordinate steps realize the gain; a per-coordinate method can also lose within a block, for instance after a rotation (\cref{ch:sign}), and \textcite{XieEtAl2025} find that rotating GPT-2's parameters slows Adam. (6) The explanation for convolutional networks is negative. The theory says why SGD does not lose there, if their blockwise spectra are homogeneous, as \textcite{ZhangEtAl2024} report; it does not derive the homogeneity.

\paragraph{Repairs and open problems.} Blockwise optimizers keep the invariance that the theory says matters, at a fraction of Adam's memory: Adalayer, which keeps one second moment per layer \parencite{ZhaoEtAl2025}, and Adam-mini, which keeps one learning rate per block \parencite{ZhangEtAl2025}. Rotation-equivariant preconditioners, such as Shampoo, SOAP and Muon, remove the dependence on the coordinate system (\cref{ch:kronecker,ch:spectral}). The scaling laws of \textcite{KunstnerBach2025} extend \cref{thm:transformers:zipf} to Zipf exponents other than one and to targets that are themselves Zipfian. Open are a model in which the three mechanisms act together, with shared parameters and minibatch sampling of tokens, and a theory of Adam at small batch that includes the estimation noise of its second moment and the half-life rule of \textcite{MarekEtAl2025} (Exercise~\ref{ex:transformers:unified}).

\section{Back in the wild}\label{sec:transformers:back}

\fieldnotes{Adam(W) for transformers, SGD with momentum for convnets; the gap survives full batch, grows with the batch size, and nearly vanishes at batch size one.}
{A table with Zipfian row frequencies; a block-diagonal quadratic with blocks at different scales; the noisy quadratic model with the preconditioner $H^{1/2}$ as Adam.}
{Weights are clocks for gradient descent and invisible to scale-invariant methods; on Zipf tables gradient descent removes $\ln10/Z_K$ of the loss per decade and needs order $K^{1-\eps}$ steps; blockwise steps replace $\kappa$ by $\max_b\kappa_b$; the gap tends to $\kappa^p$ at large batch and to $1/\E_c[(h_d/h)^p]$ at small batch.}
{Logarithmic loss curves and a bucket staircase for GD on heavy-tailed data; a gap independent of vocabulary size for Adam; a gap that opens with the batch size and a critical batch size raised by preconditioning.}
{Parameter sharing, interaction of the mechanisms, Adam's $\epsilon$ and $\beta_2$ memory, schedules (except numerically), the measured size of the gap at scale.}
{Expect Adam's advantage to be largest with heavy-tailed vocabularies and large batches, and smallest at batch size one with a decaying schedule (proved in the idealizations, suggested for networks); blockwise learning rates should keep most of it (suggested).}

Two practical readings follow, each with its status. First, the gap is not a mystery of attention: it is what a scale invariance buys on a loss whose blocks and rows live at different scales, and Zipfian vocabularies and heterogeneous parameter blocks make such losses. An optimizer that keeps one statistic per block and normalizes by it should keep most of the advantage, which is what Adalayer and Adam-mini report; that it keeps all of it is not proved. Second, the batch size changes the answer. At large batch, training is limited by curvature and the preconditioner pays its full condition-number factor; at small batch with a decaying schedule, training is limited by samples, and SGD is close to Adam in the noisy quadratic model, in agreement with \textcite{MarekEtAl2025}. The rule ``SGD for convnets'' is consistent with the block theorem if convnets are block-homogeneous; that they are is an observation, not a derivation.

The chapter collects threads of the whole book. The scale invariance is the coordinatewise geometry of sign descent (\cref{ch:sign}) and of Adam (\cref{ch:adaptive}), and its loss under rotation is the reason for Kronecker and orthogonalized optimizers (\cref{ch:kronecker,ch:spectral}). The batch-size theorem is \cref{ch:noise}'s noisy quadratic model read as a comparison of two optimizers, and the role of averaging connects it to schedules (\cref{ch:schedules}) and to momentum as a rescaled step size (\cref{ch:momentum}). Where the bias of each optimizer ends up is the subject of \cref{ch:bias}.

\begin{exercises}
\item \consolidation (Other Zipf exponents.) Let $q_j\propto j^{-\alpha}$, rows of unit initial error and $\eta=1/q_1$. (a) Show that $\Loss_t/\Loss_0=\sum_jj^{-\alpha}(1-j^{-\alpha})^{2t}/\sum_jj^{-\alpha}$. (b) For $\alpha>1$ and $K=\infty$, show that $\Loss_t/\Loss_0\sim\Gamma(1-\frac1\alpha)(2t)^{-(\alpha-1)/\alpha}/(\alpha\zeta(\alpha))$ as $t\to\infty$. (c) For $\alpha<1$, show that $\Loss_t/\Loss_0$ stays above $\frac12$ until $t$ is of order $K^\alpha$. Conclude that $\alpha=1$ is the only exponent for which the loss falls by a fixed amount per decade over a range of times that grows with $K$. \hint{in (b) replace the sum by $\int_1^\infty x^{-\alpha}e^{-2tx^{-\alpha}}\dd x$ and substitute $u=2tx^{-\alpha}$.}\label{ex:transformers:alpha}

\item \consolidation (Symmetric errors.) In \cref{cor:transformers:sign}, let $\abs{e_i}=a$ for every coordinate. Show that sign descent with step size $a(1-\sqrt\eps)$ reaches $\Loss\le\eps\Loss_0$ in one step for every $\eps\in(0,1)$ and all frequencies, and explain why this says more about the initial error than about sign descent.\label{ex:transformers:symmetric}

\item \consolidation\hard (Zipfian targets.) Let a single row have targets $y_i=1/(iZ_r)$, $i=1,\dots,r$, start at $w=0$, and run sign descent with a constant step size $\eta$. (a) Show that every coordinate with $y_i<\eta/2$ has error $y_i$ at even steps and $\eta-y_i$ at odd steps, so that the loss at every odd step is at least $\frac18\eta^2\cdot\#\{i:y_i<\eta/2\}$. (b) Deduce that if the loss is at most $\eps$ times its initial value at two consecutive steps, then, for $r$ large compared with $1/\eps$, the step size is at most of order $\sqrt{\eps/r}/Z_r$ and the number of steps at least of order $\sqrt{r/\eps}$: the square-root law of \textcite{KunstnerBach2025}. (c) Show that the loss at even steps alone falls below $\eps$ times its initial value after a number of steps of order $1/\eps$, whatever $r$, and say what this means for measuring the progress of a method that oscillates. \hint{the initial loss is $\frac12\norm y_2^2\approx\pi^2/(12Z_r^2)$; at even steps the coordinates with $y_i\ge\eta$ have errors at most $\eta$, and there are at most $1/(\eta Z_r)$ of them.}\label{ex:transformers:signzipf}

\item \consolidation (The softmax Hessian.) (a) Show that $\lmax(\diag(s)-ss^\top)=\frac12$ for a probability vector $s$ if and only if $s$ has two coordinates equal to $\frac12$. (b) Deduce from the descent lemma (\cref{lem:toolkit:descent}) that gradient descent with every step size below $4/q_1$ decreases the loss of the softmax table at every step. (c) For $K=1000$ Zipfian frequencies, compute the eigenvalues of $\diag(q)-qq^\top$ numerically and check the interlacing of \cref{prop:transformers:softmax}(ii); the condition number is $885$.\label{ex:transformers:softmax}

\item \consolidation (The bounds of \cref{thm:transformers:blocks}(iii).) Give block quadratics with two blocks for which (a) $\kappa=\max_b\kappa_b$, (b) $\kappa=\omega$, and (c) $\kappa=\omega\max_b\kappa_b$. In each case compute the gain of blockwise step sizes.\label{ex:transformers:bounds}

\item \variation (Momentum on the Zipf table.) Run heavy ball $\theta_{t+1}=\theta_t-\eta\nabla\Loss(\theta_t)+\beta(\theta_t-\theta_{t-1})$ on the Zipf table with $\eta=1/q_1$. (a) For rows with $\eta q_j\ll(1-\beta)^2$, show that the error contracts by about $1-\eta q_j/(1-\beta)$ per step. (b) Deduce that for $t\gg1/(1-\beta)$ the loss of heavy ball at step $t$ is about that of gradient descent at step $t/(1-\beta)$: momentum shifts the curve by $\log_{10}\frac1{1-\beta}$ decades without changing its slope. Check numerically for $K=10^4$ and $\beta=0.9$, where heavy ball at steps $100$ and $1000$ matches gradient descent at steps $1000$ and $10^4$ to within $3\%$. (c) What do the rows with $\eta q_j\gtrsim(1-\beta)^2$ do?\label{ex:transformers:momentum}

\item \variation (Noise that is not proportional to curvature.) In \cref{cor:transformers:gap} let $h_i=1/i$ and $c_i=i^{-\vartheta}$. Show that the small-batch limit tends, as $d\to\infty$, to $1+p/(1-\vartheta)$ for $0\le\vartheta<1$, and grows like $p\ln d$ for $\vartheta=1$. Check the value $1.5$ for isotropic noise against \cref{fig:transformers:batch}, where the target is not small.\label{ex:transformers:noise}

\item \variation (Adam's preconditioner in the noisy quadratic model.) (a) Show that $\E g_{t,i}^2=h_i^2\theta_{t,i}^2+c_i/B$ given $\theta_t$. (b) Deduce that RMSProp with $\beta_2$ close to one behaves near $\theta_t$ like preconditioned SGD with $a_i=\sqrt B/\sqrt{c_i+Bh_i^2\theta_{t,i}^2}$. (c) For $c_i=c_0h_i$ and small $B$, show that RMSProp with learning rate $\eta$ is the preconditioner $h_i^{-1/2}$ with step size $\eta\sqrt{B/c_0}$, and that keeping the small-batch dynamics fixed as $B$ changes requires $\eta\propto\sqrt B$, the square-root rule of \textcite{MalladiEtAl2022}. (d) What does RMSProp become at large $B$?\label{ex:transformers:adamnqm}

\item \variation (Blockwise normalized steps.) Let $H=\diag(s_1\bar H,s_2\bar H)$ with $\bar H=\diag(1,1/\kappa_0)$. Compare the best rate of gradient descent, as a function of $s_1/s_2$, with blockwise normalized descent $\theta_b\leftarrow\theta_b-\eta_tg_b/\norm{g_b}_2$ under a decaying schedule, whose iterates do not depend on $s_1,s_2$ by \cref{thm:transformers:blocks}(iv).\label{ex:transformers:blocknorm}

\item \simulation (The softmax table.) Train a softmax bigram table with $K=1000$ Zipfian input frequencies and Zipfian conditional distributions, from $W=0$, by gradient descent with $\eta=1/q_1$ and by Adam. Plot the loss on each input-frequency bucket. Check that under gradient descent the loss of row $j$ at step $t$ equals the loss of the single-row problem after $t$ steps of step size $q_j/q_1$ (\cref{prop:transformers:softmax}(i)), and that under Adam all buckets coincide.\label{ex:transformers:simsoftmax}

\item \simulation (Rotating the table.) Rotate the square-loss Zipf table with $K=1000$ by a random orthogonal matrix $Q\in\R^{K\times K}$ acting on the row index, $W\mapsto QW$. Run gradient descent, sign descent and Adam. Check that the loss curve of gradient descent is unchanged, as rotation equivariance predicts, and that sign descent and Adam lose most of their advantage; explain the difference with \cref{prop:transformers:clock}.\label{ex:transformers:rotation}

\item \simulation (Averaging and decay.) Using \texttt{14-transformers-batch.py}, reproduce the advantages $2.97$ (constant step), $1.31$ (linear decay) and $0.99$ (tail averaging) at $B=1$. Then lower the target to $1\%$ and $0.1\%$ of the initial loss and check that the averaged advantage stays near one while the constant-step advantage rises toward $3.83$ ($3.24$ and $3.43$), as \cref{cor:transformers:gap} and the heuristic derivation predict.\label{ex:transformers:simavg}

\item \research (One model for three mechanisms.) Build a model in which heavy-tailed classes, block heterogeneity and the batch-size regime act together, for instance a softmax bigram model with a shared embedding of dimension $k\ll K$ trained on minibatches of sampled tokens, and derive the advantage of a blockwise scale-invariant optimizer over SGD as a function of $K$, $k$ and $B$. No such model is solved as far as we know.\label{ex:transformers:unified}

\item \research (Adam at batch size one.) Extend \cref{thm:transformers:batch} to RMSProp with $\beta_2$ chosen by the half-life rule of \textcite{MarekEtAl2025}, $1-\beta_2\propto B$, including the estimation noise of its second moment, and decide whether Adam keeps an advantage over SGD with a decaying schedule as $B\to0$.\label{ex:transformers:adambs}
\end{exercises}

\notesandrefs
The gap between Adam and SGD on transformers was a practitioner's observation long before it was a research question; the transformer itself was trained with Adam \parencite{VaswaniEtAl2017}. \Textcite{ZhangEtAl2020} proposed heavy-tailed gradient noise as the cause and adaptive clipping as the remedy. \Textcite{KunstnerEtAl2023} tested the noise explanation by growing the batch size to the full data set, found that the gap persists, and pointed to sign descent. Two structural explanations followed in 2024: heavy-tailed class imbalance \parencite{YadavEtAl2023,KunstnerEtAl2024} and block heterogeneity of the Hessian \parencite{ZhangEtAl2024}. \Textcite{ZhaoEtAl2025} compared optimizers systematically and located Adam's advantage in a few parameter groups, \textcite{XieEtAl2025} in the coordinate system, and \textcite{MarekEtAl2025} showed that it depends on the batch size.

The Zipf table is the linear bigram model of \textcite{KunstnerBach2025}, who derive scaling laws for gradient descent and sign descent under power-law frequencies $q_k\propto1/k^\alpha$, find that $\alpha=1$ is the worst case for gradient descent, and find that sign descent needs a number of steps that grows only like the square root of the dimension. \Cref{thm:transformers:zipf} isolates the effect of the input frequencies by giving every row the same initial error, and then the constants come out exactly: the slope $\ln10/Z_K$ per decade and the offset $\ln2+\gamma$. \Textcite{YadavEtAl2025} prove an advantage of sign descent in a minimal next-token model. \Cref{prop:transformers:clock}, the corollary for sign descent and the interlacing in \cref{prop:transformers:softmax} are elementary, and we know no single reference for them.

The noisy quadratic model is due to \textcite{ZhangEtAl2019}, who found that preconditioning and momentum extend perfect batch-size scaling to larger batches; the critical batch size is due to \textcite{McCandlishEtAl2018}; scaling rules for adaptive methods from stochastic differential equations are in \textcite{MalladiEtAl2022}. The small-batch formula of \cref{thm:transformers:batch}, the two ends of \cref{cor:transformers:gap}, and the observation that averaging removes the small-batch advantage are, as far as we know, not in the literature in this form. Adam as sign descent scaled by the signal-to-noise ratio is \textcite{BallesHennig2018}; that its preconditioner is not the Fisher information is \textcite{KunstnerHennigBalles2019}. \Textcite{WenEtAl2026} measure how optimizer speedups change with scale and tuning.

\paragraph{Translating the literature.} \Textcite{KunstnerBach2025} write $\pi_k$ for the token frequencies, our $q_k$, and $\alpha$ for the exponent of the power law, as we do in Exercise~\ref{ex:transformers:alpha}. In the noisy quadratic model of \textcite{ZhangEtAl2019}, as in \cref{ch:noise}, $h_i$ and $c_i$ are the curvatures and the noise variances along the eigendirections, and preconditioning by $H^{-p}$ corresponds to our multipliers $a_i=h_i^{-p}$. \Textcite{McCandlishEtAl2018} write $S_{\min}$ and $E_{\min}$ for the minimal numbers of steps and of examples, our $T_a(\infty)$ and $n_a(\eps)$, and $B_{\mathrm{crit}}=E_{\min}/S_{\min}$ for the critical batch size. What \textcite{ZhangEtAl2024} call Adam's coordinate-wise learning rates is, for deterministic gradients, our scale invariance: the learning rate of each coordinate is set by that coordinate's own gradients.

\printchapterbib
\end{refsection}

% Chapter 15 of Optimizers in the Wild: implicit bias and the margin of a norm.
\newcommand{\biasP}{\mathcal{P}_n}
\newcommand{\biasqmin}{q_{\min}}
\newcommand{\biasbar}{\bar\gamma}
\newcommand{\biasU}{\mathcal{U}^\star}
\newcommand{\biasartanh}{\operatorname{artanh}}

\casechapter{Where Steepest Descent Ends}{Implicit bias and the margin of a norm}{Adaptive methods and generalization}{ch:bias}
\begin{refsection}[chapters/15-bias.bib]

\begin{chaptersummary}
When a model can fit its training data exactly, training does not stop at zero training error: the loss keeps falling, the weights keep growing, and what the optimizer still decides is the direction in which they grow; adaptive methods have been reported to choose directions that generalize worse than those of stochastic gradient descent. For a linear classifier on separable data with the exponential loss we prove that normalized steepest descent in a norm maximizes the margin measured in that norm, with an explicit bound for every sequence of step sizes: gradient descent finds the $\ell_2$ maximum-margin classifier, sign descent and Adam without memory the $\ell_\infty$ one, coordinate descent (boosting) the $\ell_1$ one, and orthogonalized updates the spectral one, while decoupled weight decay stops the march at radius $1/\lambda$, within $\lambda\log n$ of the maximum margin. On a dataset in the spirit of Wilson et al.\ the $\ell_\infty$ bias misclassifies every negative test example and the $\ell_2$ bias none; on another the roles reverse. The idealization misses depth, noise, Adam's $\epsilon$, finite training times, and the evidence that well-tuned adaptive methods need not generalize worse.
\end{chaptersummary}

\section{In the wild: the solutions optimizers choose}\label{sec:bias:wild}

Two optimizers that reach the same training loss need not reach the same network. Most of this book compares optimizers by how fast they reduce the loss; this chapter asks where they end, and whether the end point explains a difference in test error. The question became pressing in 2017, when \textcite{WilsonEtAl2017} compared AdaGrad, RMSProp and Adam with stochastic gradient descent on image classification, character-level language modeling and parsing, and concluded that practitioners should reconsider the use of adaptive methods for training neural networks. Practice in computer vision had settled on stochastic gradient descent with momentum already; their paper gave the habit a reason.\sidenote{In this chapter ``learning rate'' (in the practice sections) and ``step size'' (in the theory) are synonyms, as everywhere in the book.}

\begin{practice}[ResNets, \textcite{HeEtAl2016}; \textcite{WilsonEtAl2017}; \textcite{KeskarSocher2017}]
Train convolutional networks for image classification with SGD with momentum $0.9$ and weight decay; for ResNets on ImageNet, a learning rate of $0.1$ for a batch of $256$ examples, divided by $10$ whenever the error plateaus, and weight decay $10^{-4}$. When test accuracy is what counts, prefer SGD to Adam, or start with Adam and switch to SGD later in training. Train transformers with Adam or AdamW.
\end{practice}

The first sentence is a documented recipe. The second is partly folklore: it rests on comparisons whose outcome, as we shall see, depends on how the optimizers were tuned. The third is the subject of \cref{ch:transformers}. The evidence comes in three kinds.

\begin{observation}[\textcite{WilsonEtAl2017}]\label{obs:bias:wilson}
On image classification (CIFAR-10), character-level language modeling and two parsing benchmarks, the solutions found by AdaGrad, RMSProp and Adam had higher test error than those found by well-tuned SGD, in some cases even when the adaptive methods reached a lower training loss. On a constructed linear classification problem trained with the square loss, the adaptive methods provably converge to a classifier that labels every new example positive, while gradient descent converges to the minimum-norm interpolant, which labels every new example correctly.
\end{observation}

The constructed problem is the clean part of the evidence, and \cref{ex:bias:wilson} rebuilds it for the margin. The benchmark comparisons are the contested part. The second kind of evidence is that training past zero training error still changes the solution. \Textcite{SoudryEtAl2018} trained a convolutional network on CIFAR-10 with SGD, momentum and the cross-entropy loss beyond the point of zero training error: the training loss kept decaying roughly like $1/t$, the norm of the last layer's weights grew roughly like $\log t$, and while the validation loss rose, the validation error kept improving slowly.\sidenote{\Textcite{HofferHubaraSoudry2017} made a related observation for large batches: much of the generalization gap of large-batch training closes if one trains for more steps.}

\begin{observation}[\textcite{ChoiEtAl2019}; \textcite{LoshchilovHutter2019}]\label{obs:bias:tuning}
When every hyperparameter of each optimizer is tuned, including Adam's $\epsilon$ and its momentum, the adaptive methods matched or beat momentum SGD on every workload of \textcite{ChoiEtAl2019} (an unrefereed preprint). They explain the result by inclusion: Adam with a large $\epsilon$ behaves like momentum SGD with a rescaled learning rate, so a search over Adam's hyperparameters contains SGD. With weight decay decoupled from the gradient step, Adam's test error on image classification comes close to that of SGD with momentum \parencite{LoshchilovHutter2019}.
\end{observation}

The third kind of evidence is geometric. \Cref{ch:decay} shows that AdamW with weight decay $\lambda$ keeps its weights in an $\ell_\infty$ ball of radius $1/\lambda$, and that without memory it is the Frank--Wolfe method on that ball \parencite{XieLi2024}. So the optimizers do not merely move at different speeds: they move in different geometries, and the geometry leaves a trace in the weights.

The status, then: that adaptive methods generalize worse is observed in some settings and not in others, and it is not a law. What is solid is that the solution reached at a given training loss depends on the optimizer. The question for the theory is: \emph{when the training data can be fit exactly, where does steepest descent in a given norm end up, and can the answer explain why two optimizers with the same training loss generalize differently?}

\section{The idealization: separable data and the margin of a norm}\label{sec:bias:ideal}

\begin{assumption}[Separable linear classification]\label{ass:bias:separable}
The examples $x_1,\dots,x_n$ lie in a finite-dimensional real inner-product space $V$: either $\R^d$ with $\ip{\theta,x}=\theta^\top x$, or $\R^{K\times d}$ with $\ip{W,X}=\tr(W^\top X)$. The labels are $y_i\in\{-1,+1\}$. The space carries a norm $\norm{\cdot}$ with dual norm $\norm{\cdot}_*$, and $R=\max_i\norm{x_i}_*$. The data are \term{linearly separable}\index{separable data} through the origin: some $\theta$ has $y_i\ip{\theta,x_i}>0$ for every $i$. The classifier is $x\mapsto\sign\ip{\theta,x}$, and the training loss is the exponential loss
\begin{equation}\label{eq:bias:loss}
  \Loss(\theta)=\frac1n\sum_{i=1}^n\exp\bigl(-y_i\ip{\theta,x_i}\bigr).
\end{equation}
\end{assumption}

Only the products $z_i=y_ix_i$ matter, and we call them the \emph{signed examples}. We write $q_i(\theta)=\ip{\theta,z_i}$ for the margin of example $i$, which is positive when the example is classified correctly, and $\biasqmin(\theta)=\min_iq_i(\theta)$.

\begin{definition}[The margin of a norm]\label[definition]{def:bias:margin}
For $\theta\ne0$ the \term{normalized margin}\index{margin!normalized} is $\biasbar(\theta)=\biasqmin(\theta)/\norm\theta$. The \term{maximum margin}\index{margin!of a norm} of the data in the norm $\norm{\cdot}$ is
\[
  \gamma=\max_{\norm u\le1}\ \min_{i\le n}\ \ip{u,z_i},
\]
and the \term{max-margin set}\index{maximum-margin classifier} is $\biasU=\{u:\norm u\le1,\ \biasqmin(u)=\gamma\}$; its elements are the maximum-margin classifiers. When several norms are in play we write $\gamma_2,\gamma_\infty,\gamma_1,\gamma_{\op}$.
\end{definition}

The maximum exists because $\biasqmin$ is continuous and the unit ball is compact; $\gamma\ge0$ because $u=0$ is allowed, and $\gamma>0$ exactly when the data are separable. Every $u\in\biasU$ has $\norm u=1$ when $\gamma>0$, since $u/\norm u$ would otherwise have a larger margin. The normalized margin is scale-free, and $\biasbar(\theta)\le\gamma$ for every $\theta\ne0$. Its geometric meaning: the distance from a point $x$ to the hyperplane $\{\ip{\theta,\cdot}=0\}$, measured in the dual norm, is $\abs{\ip{\theta,x}}/\norm\theta$ (Exercise~\ref{ex:bias:distance}). So the $\ell_2$ maximum-margin classifier is the hard-margin support vector machine through the origin \parencite{CortesVapnik1995}, the $\ell_\infty$ one maximizes the $\ell_1$ distance from the hyperplane to the nearest example, and the $\ell_1$ one maximizes the $\ell_\infty$ distance.\sidenote{Separating planes that maximize the distance in an arbitrary norm were studied by \textcite{Mangasarian1999}, who showed that the distance is measured by the dual norm of the normal vector.}

The optimizers are the normalized steepest descent methods of \cref{ch:steepest}:
\begin{equation}\label{eq:bias:nsd}
  \theta_{t+1}=\theta_t+\eta_t\lmo\bigl(\nabla\Loss(\theta_t)\bigr),\qquad \lmo(g)\in\argmin_{\norm u\le1}\ip{g,u},
\end{equation}
with step sizes $\eta_t>0$ that may depend on the iterates so far, and with every choice of $\lmo$ when the minimizer is not unique. Each step has length at most $\eta_t$ in the norm, and $\ip{g,\lmo(g)}=-\norm g_*$. In the $\ell_2$ norm this is normalized gradient descent; in the $\ell_\infty$ norm it is sign descent, which is Adam with $\beta_1=\beta_2=0$ and $\epsilon=0$ (\cref{ch:sign}); in the $\ell_1$ norm it is greedy coordinate descent; and in the spectral norm on matrices it is the orthogonalized update of Muon without momentum (\cref{ch:spectral}).

\paragraph{What was thrown away, and why.} \emph{Features.} The model is linear in fixed features: the last layer of a network whose lower layers are frozen. It is the smallest model in which the loss has no minimizer and the optimizer must choose a direction; \cref{sec:bias:beyond} returns to homogeneous networks. \emph{Non-separable data.} Overparameterized classifiers interpolate their training data, and the interesting regime is the one after the training error has reached zero. \emph{The loss.} The exponential loss has the tail of the logistic and cross-entropy losses and makes the algebra exact: its gradient is a softmax-weighted sum of the signed examples. The logistic loss has the same limits \parencite{SoudryEtAl2018,JiTelgarsky2021}. \emph{Noise, memory and preconditioning.} Full-batch gradients, no momentum, no second-moment estimate; their effect on the limit is collected in \cref{sec:bias:beyond}. \emph{Finite time.} The theorem describes the end of training, with explicit rates, and the rates turn out to be the crux.

\begin{dictionary}
signed examples $z_i=y_ix_i$, margins $q_i(\theta)$ & training examples and the logit of the correct class minus the largest other logit\\
separability, $\Loss\to0$ without a minimizer & training past zero training error with the cross-entropy loss\\
the norm of steepest descent & the optimizer's geometry: $\ell_2$ for SGD, $\ell_\infty$ for sign descent and Adam, spectral for Muon, $\ell_1$ for boosting\\
normalized margin $\biasbar(\theta)$ & the worst example's logit gap per unit of weight norm, in the optimizer's norm\\
maximum margin $\gamma$ & the best worst-case confidence the geometry allows\\
softmax weights $p_i(\theta)$ & per-example weights in the gradient: hard examples dominate\\
smoothed margin $-\log(n\Loss)$ & the training loss on a logarithmic scale\\
distance travelled $\tau_T=\sum_{t<T}\eta_t$ & the area under the learning-rate schedule\\
the ball $\norm\theta\le1/\lambda$ & decoupled weight decay $\lambda$ (\cref{ch:decay})
\end{dictionary}

The chapter derives one optimizer that the book has not met: boosting is $\ell_1$ steepest descent on the exponential loss, with the features given by weak classifiers $h_j$, $x_{ij}=h_j(x_i)\in\{-1,1\}$.\sidenote{That AdaBoost is coordinate descent on the exponential loss was pointed out by Breiman and developed by \textcite{FriedmanHastieTibshirani2000}; \textcite{Friedman2001} turned it into gradient boosting with general losses and shrinkage.}

\begin{fieldguide}{AdaBoost}{\textcite{FreundSchapire1997}}
\fgrow{update}{example weights $p_i\propto e^{-y_i\ip{\theta,x_i}}$; pick $j\in\argmax_j\abs{\sum_ip_iy_ih_j(x_i)}$; $\theta_j\leftarrow\theta_j+\alpha\,\sign(\sum_ip_iy_ih_j(x_i))$}
\fgrow{norm}{$\ell_1$ on the weights of the weak classifiers: normalized steepest descent \eqref{eq:bias:nsd} with exact line search, $\alpha=\biasartanh(r)$, where $r$ is the edge \eqref{eq:bias:edge}}
\fgrow{preconditioner}{none}
\fgrow{averages}{none; the example weights $p$ are the dual variables of \cref{lem:bias:duality}}
\fgrow{state}{one weight per example and per weak classifier used}
\fgrow{in the wild}{decision stumps or small trees as weak classifiers; in gradient boosting, shrinkage of each step by a small factor \parencite{Friedman2001}}
\fgrow{caution}{with its full line-search step AdaBoost can stop short of the maximum $\ell_1$ margin \parencite{RudinDaubechiesSchapire2004}; small steps reach it (\cref{cor:bias:steps})}
\end{fieldguide}

\paragraph{What one had to notice.} On separable data the training loss has infimum zero and no minimizer, so every descent method sends $\norm{\theta_t}$ to infinity, and the only thing that can converge is the direction $\theta_t/\norm{\theta_t}$. The exponential tail turns the gradient into a softmax over the examples, concentrated on those with the smallest margins; late in training the optimizer sees only these support vectors, and it pushes the weights in the direction that raises the smallest margins fastest \emph{per unit of its own norm}. \Textcite{SoudryEtAl2018} made this precise for gradient descent, and \textcite{GunasekarEtAl2018} saw that the norm of the step decides which margin is maximized.

\section{The theory: the margin of a norm}\label{sec:bias:theory}

Throughout this section \cref{ass:bias:separable} holds, so that $\gamma>0$. We write $\biasP=\{p\in\R^n:p_i\ge0,\ \sum_ip_i=1\}$ for the probability simplex.

\subsection{Duality, and a gradient that never vanishes}

The maximum margin has a dual description, which identifies the examples that matter.

\begin{lemma}[Duality of margins]\label[lemma]{lem:bias:duality}
Let $C=\conv\{z_1,\dots,z_n\}$. Then
\[
  \gamma=\min_{p\in\biasP}\Bigl\lVert\sum_{i=1}^np_iz_i\Bigr\rVert_*=\min_{c\in C}\norm c_*,
\]
so the maximum margin is the dual-norm distance from the origin to the convex hull of the signed examples, and the data are separable if and only if $0\notin C$. If $u^\star\in\biasU$ and $p^\star$ attains the minimum, with $c^\star=\sum_ip^\star_iz_i$, then:
\begin{enumerate}
\item $\ip{u^\star,c^\star}=\norm{c^\star}_*=\gamma$, so $u^\star$ maximizes $\ip{u,c^\star}$ over the unit ball: it is a value of $\lmo(-c^\star)$;
\item $p^\star_i=0$ for every $i$ with $q_i(u^\star)>\gamma$.
\end{enumerate}
\end{lemma}

\begin{proof}
The plan is to prove the two inequalities between $\gamma$ and $\delta=\min_{c\in C}\norm c_*$ separately: one is Hölder's inequality, the other the separating hyperplane theorem. The minimum $\delta$ is attained because $C$ is compact and the dual norm is continuous.

For every $u$ with $\norm u\le1$ and every $p\in\biasP$, a minimum is at most an average, and Hölder's inequality (\cref{lem:toolkit:dualnorm}) bounds the average:
\[
  \biasqmin(u)\le\sum_ip_iq_i(u)=\Bigl\langle u,\sum_ip_iz_i\Bigr\rangle\le\Bigl\lVert\sum_ip_iz_i\Bigr\rVert_* .
\]
Maximizing over $u$ and minimizing over $p$ gives $\gamma\le\delta$.

For the reverse inequality, if $\delta=0$ there is nothing to prove, because $\gamma\ge0$. If $\delta>0$, the open set $B=\{v:\norm v_*<\delta\}$ and $C$ are disjoint convex sets, and the separating hyperplane theorem gives $\theta\ne0$ and $b\in\R$ with $\ip{\theta,v}\le b\le\ip{\theta,c}$ for all $v\in B$ and $c\in C$. The supremum of $\ip{\theta,v}$ over $B$ is $\delta\norm\theta$, by the bidual formula $\norm\theta=\max_{\norm v_*\le1}\ip{\theta,v}$ (\cref{lem:toolkit:dualnorm}). Hence $u=\theta/\norm\theta$ satisfies $q_i(u)\ge b/\norm\theta\ge\delta$ for every $i$, and $\gamma\ge\delta$. The data are separable if and only if some $\theta$ has $\biasqmin(\theta)>0$, that is, if and only if $\gamma>0$, which we have just shown to be equivalent to $\delta>0$, that is, to $0\notin C$.

Finally, for $u^\star$ and $p^\star$ as in the statement, the same chain reads
\[
  \gamma=\biasqmin(u^\star)\le\sum_ip^\star_iq_i(u^\star)=\ip{u^\star,c^\star}\le\norm{u^\star}\norm{c^\star}_*\le\norm{c^\star}_*=\gamma ,
\]
so every inequality is an equality. The middle one gives (i), since $\max_{\norm u\le1}\ip{u,c^\star}=\norm{c^\star}_*$, and the maximizers of $\ip{u,c^\star}$ are the minimizers of $\ip{-c^\star,u}$. The first one says that $\sum_ip^\star_i(q_i(u^\star)-\gamma)=0$, a sum of nonnegative terms, which gives (ii).
\end{proof}

The examples with $p^\star_i>0$ are \term{support vectors}\index{support vector}: they attain the minimum margin. Part (i) says something about steepest descent, which the proofs below make quantitative: if the optimizer's weights on the examples settle at $p^\star$, its step settles at a maximum-margin direction.

\begin{lemma}[The gradient is at least the margin times the loss]\label[lemma]{lem:bias:gradient}
For every $\theta\in V$, with the softmax weights $p_i(\theta)=e^{-q_i(\theta)}/\sum_je^{-q_j(\theta)}$,
\[
  \nabla\Loss(\theta)=-\Loss(\theta)\sum_{i=1}^np_i(\theta)\,z_i,\qquad \gamma\,\Loss(\theta)\le\norm{\nabla\Loss(\theta)}_*\le R\,\Loss(\theta).
\]
In particular $\nabla\Loss$ never vanishes, $\Loss$ has no minimizer, and $\inf\Loss=0$.
\end{lemma}

\begin{proof}
Differentiating \eqref{eq:bias:loss} gives $\nabla\Loss(\theta)=-\frac1n\sum_ie^{-q_i(\theta)}z_i$, which is the formula because $\frac1n\sum_je^{-q_j(\theta)}=\Loss(\theta)$. The weights $p(\theta)$ lie in $\biasP$, so \cref{lem:bias:duality} gives $\norm{\sum_ip_iz_i}_*\ge\gamma$, and the triangle inequality gives $\norm{\sum_ip_iz_i}_*\le\sum_ip_i\norm{z_i}_*\le R$. Since $\gamma>0$ and $\Loss>0$, the gradient never vanishes, and a convex differentiable function without stationary points has no minimizer. Finally $\Loss(su^\star)\le e^{-s\gamma}$ for $u^\star\in\biasU$ and $s>0$, which tends to $0$.
\end{proof}

We call the ratio
\begin{equation}\label{eq:bias:edge}
  r(\theta)=\frac{\norm{\nabla\Loss(\theta)}_*}{\Loss(\theta)}\in[\gamma,R]
\end{equation}
the \term{edge}\index{edge} of the step at $\theta$, the name it has in boosting, where $r$ is the weighted correlation of the best weak classifier with the labels. The lemma is the key inequality of the chapter: the gradient never becomes small faster than the loss.

\begin{lemma}[One normalized step]\label[lemma]{lem:bias:step}
For every $\theta\in V$ and $\eta>0$, with $u=\lmo(\nabla\Loss(\theta))$ and $r=r(\theta)$,
\[
  \Loss(\theta+\eta u)\le\Bigl(\cosh(\eta R)-\frac rR\sinh(\eta R)\Bigr)\Loss(\theta)\le\exp\Bigl(-\eta r+\frac{\eta^2R^2}2\Bigr)\Loss(\theta).
\]
\end{lemma}

\begin{proof}
The plan is to write the ratio of the two losses as an average of exponentials of bounded numbers whose mean is $\eta r$, and to bound that average first by the chord of the exponential and then by a Gaussian bound. Let $s_i=\ip{u,z_i}$. Since $\norm u\le1$, Hölder's inequality gives $\abs{s_i}\le R$, and by \cref{lem:bias:gradient} and $\ip{g,\lmo(g)}=-\norm g_*$,
\[
  \sum_ip_i(\theta)\,s_i=\Bigl\langle u,\sum_ip_i(\theta)z_i\Bigr\rangle=-\frac{\ip{u,\nabla\Loss(\theta)}}{\Loss(\theta)}=r .
\]
Because $q_i(\theta+\eta u)=q_i(\theta)+\eta s_i$, the ratio of losses is $\Loss(\theta+\eta u)/\Loss(\theta)=\sum_ip_i(\theta)e^{-\eta s_i}$. The function $s\mapsto e^{-\eta s}$ is convex, so on $[-R,R]$ it lies below its chord:
\[
  e^{-\eta s}\le\frac{R-s}{2R}\,e^{\eta R}+\frac{R+s}{2R}\,e^{-\eta R}=\cosh(\eta R)-\frac sR\sinh(\eta R).
\]
Averaging with the weights $p_i(\theta)$ gives the first inequality. For the second, fix $\rho\in[-1,1]$ and put $D(x)=\cosh x-\rho\sinh x$ and $N(x)=\sinh x-\rho\cosh x$ for $x\ge0$, so that $D\ge e^{-x}>0$, $D'=N$ and $N'=D$. The function $f(x)=\log D(x)+\rho x$ has $f(0)=0$, $f'(x)=N(x)/D(x)+\rho$, so $f'(0)=0$, and
\[
  f''(x)=\frac{D(x)^2-N(x)^2}{D(x)^2}=1-\Bigl(\frac{N(x)}{D(x)}\Bigr)^2\le1 .
\]
Taylor's theorem with integral remainder gives $f(x)\le x^2/2$, that is, $D(x)\le\exp(-\rho x+x^2/2)$. With $x=\eta R$ and $\rho=r/R\in[0,1]$ this is the second inequality.
\end{proof}

The second bound is Hoeffding's lemma for a variable with values in $[-R,R]$; the first is attained when every $s_i$ equals $\pm R$, as in AdaBoost with weak classifiers valued in $\{-1,1\}$.

\subsection{The theorem}

The proof follows the \term{smoothed margin}\index{smoothed margin}
\[
  \psi(\theta)=-\log\bigl(n\Loss(\theta)\bigr)=-\log\sum_{i=1}^ne^{-q_i(\theta)},
\]
a soft minimum of the margins, which satisfies
\begin{equation}\label{eq:bias:psi}
  \biasqmin(\theta)-\log n\le\psi(\theta)\le\biasqmin(\theta),
\end{equation}
because $e^{-\biasqmin}\le\sum_ie^{-q_i}\le ne^{-\biasqmin}$. We write
\[
  \tau_T=\sum_{t<T}\eta_t
\]
for the distance travelled in $T$ steps, an upper bound on $\norm{\theta_T}$.

\begin{theorem}[The margin of a norm; after \textcite{GunasekarEtAl2018}]\label{thm:bias:margin}
Assume \cref{ass:bias:separable}. Let $\theta_0=0$ and let $\theta_t$ follow \eqref{eq:bias:nsd} with step sizes $\eta_t>0$, which may depend on $\theta_0,\dots,\theta_t$, and with every choice of $\lmo$. Then for every $T\ge1$:
\begin{enumerate}
\item $\Loss(\theta_T)\le\exp\bigl(-\gamma\tau_T+\frac{R^2}2\sum_{t<T}\eta_t^2\bigr)$ and $\norm{\theta_T}\le\tau_T$;
\item if $\gamma\tau_T>\log n+\frac{R^2}2\sum_{t<T}\eta_t^2$, then $\theta_T$ classifies every training example correctly, and
\begin{equation}\label{eq:bias:bound}
  \gamma-\frac{\log n+\frac{R^2}2\sum_{t<T}\eta_t^2}{\tau_T}\ \le\ \biasbar(\theta_T)\ \le\ \gamma .
\end{equation}
\end{enumerate}
\end{theorem}

\begin{proof}
The plan is to show that the smoothed margin grows by at least $\gamma\eta_t-R^2\eta_t^2/2$ per step while the norm grows by at most $\eta_t$, and to divide. By \cref{lem:bias:step} at $\theta=\theta_t$, and by $r(\theta_t)\ge\gamma$ (\cref{lem:bias:gradient}),
\[
  \psi(\theta_{t+1})-\psi(\theta_t)=-\log\frac{\Loss(\theta_{t+1})}{\Loss(\theta_t)}\ge\eta_tr(\theta_t)-\frac{R^2\eta_t^2}2\ge\gamma\eta_t-\frac{R^2\eta_t^2}2 .
\]
Summing over $t<T$, with $\psi(\theta_0)=\psi(0)=-\log n$,
\begin{equation}\label{eq:bias:psigrowth}
  \psi(\theta_T)\ge\gamma\tau_T-\frac{R^2}2\sum_{t<T}\eta_t^2-\log n .
\end{equation}
Since $n\Loss(\theta_T)=e^{-\psi(\theta_T)}$, this is the bound on the loss in (i). The bound on the norm follows from the triangle inequality, because each step has norm $\eta_t\norm{\lmo(\cdot)}\le\eta_t$.

For (ii), let $A_T$ denote the right side of \eqref{eq:bias:psigrowth}, which is positive by hypothesis. By \eqref{eq:bias:psi}, $\biasqmin(\theta_T)\ge\psi(\theta_T)\ge A_T>0$, so every example is classified correctly and $\theta_T\ne0$. Dividing by $\norm{\theta_T}\le\tau_T$ gives $\biasbar(\theta_T)\ge A_T/\norm{\theta_T}\ge A_T/\tau_T$, which is the lower bound in \eqref{eq:bias:bound}. The upper bound holds because $\theta_T/\norm{\theta_T}$ lies in the unit ball, so $\biasbar(\theta_T)=\biasqmin(\theta_T/\norm{\theta_T})\le\gamma$.
\end{proof}

The bound has two error terms. The term $\log n/\tau_T$ is the price of starting from uniform weights on the examples; it is all that remains when the steps shrink at a fixed distance $\tau_T$, the limit of a flow. The term $R^2\sum\eta_t^2/(2\tau_T)$ is the price of taking finite steps. What the step sizes must do follows at once.

\begin{corollary}[Step sizes]\label[corollary]{cor:bias:steps}
In the setting of \cref{thm:bias:margin}:
\begin{enumerate}
\item If $\eta_t\to0$ and $\sum_t\eta_t=\infty$, then $\biasbar(\theta_t)\to\gamma$, every limit point of $\theta_t/\norm{\theta_t}$ lies in $\biasU$, and if $\biasU=\{u^\star\}$, then $\theta_t/\norm{\theta_t}\to u^\star$.
\item For $\eta_t=c/(R\sqrt{t+1})$ with $c>0$, whenever the right side is below $\gamma$,
\[
  \gamma-\biasbar(\theta_T)\le R\,\frac{\log n+\frac{c^2}2(1+\log T)}{2c\,(\sqrt{T+1}-1)} .
\]
\item For a constant step $\eta_t=\eta$, whenever the right side is below $\gamma$, $\gamma-\biasbar(\theta_T)\le\log n/(\eta T)+R^2\eta/2$. Hence $\liminf_t\biasbar(\theta_t)\ge\gamma-R^2\eta/2$, and for a known horizon $T$ the step $\eta=\sqrt{2\log n/T}/R$ gives $\gamma-\biasbar(\theta_T)\le R\sqrt{2\log n/T}$.
\end{enumerate}
\end{corollary}

\begin{proof}
The plan is to show that the error term of \eqref{eq:bias:bound} tends to zero and to use compactness for the directions. (i) Given $\eps>0$, choose $t_0$ with $\eta_t\le\eps$ for $t\ge t_0$. Then $\sum_{t<T}\eta_t^2\le\sum_{t<t_0}\eta_t^2+\eps\tau_T$, and since $\tau_T\to\infty$ the error term of \eqref{eq:bias:bound} has limit superior at most $R^2\eps/2$; as $\eps$ is arbitrary, $\biasbar(\theta_t)\to\gamma$. For large $t$ the iterate is nonzero, and $u_t=\theta_t/\norm{\theta_t}$ lies on the unit sphere, which is compact. If $u_{t_k}\to\bar u$, then $\norm{\bar u}=1$ and, by continuity of $\biasqmin$, $\biasqmin(\bar u)=\lim_k\biasbar(\theta_{t_k})=\gamma$, so $\bar u\in\biasU$. If $\biasU=\{u^\star\}$, every subsequence of $(u_t)$ has a further subsequence converging to $u^\star$, so $u_t\to u^\star$. (ii) Compare sums with integrals: $\sum_{k=1}^Tk^{-1/2}\ge2(\sqrt{T+1}-1)$ and $\sum_{k=1}^Tk^{-1}\le1+\log T$, and insert both in \eqref{eq:bias:bound}. (iii) Here $\tau_T=\eta T$ and $\sum_{t<T}\eta_t^2=\eta^2T$; the bound $\log n/(\eta T)+R^2\eta/2$ is minimized at the stated $\eta$, where its two terms are equal.
\end{proof}

So step sizes that decay to zero but not too fast reach the maximum margin, at the rate $\tilde O(1/\sqrt T)$ for $\eta_t\propto1/\sqrt t$. A constant step is guaranteed only to come within $R^2\eta/2$; whether it reaches the maximum is a question the theorem leaves open, and \cref{sec:bias:reading} shows that in our experiments it does.

\begin{remark}[Hedge against a best response]\label[remark]{rem:bias:hedge}
The proof is a regret bound in disguise. The weights $p_t=p(\theta_t)$ evolve by multiplicative weights, $p_{t+1,i}\propto p_{t,i}\exp(-\eta_t\ip{u_t,z_i})$ with $u_t=\lmo(\nabla\Loss(\theta_t))$, and by \cref{lem:bias:gradient} the step $u_t$ maximizes $\ip{u,\sum_ip_{t,i}z_i}$ over the unit ball: it is a best response to $p_t$ in the zero-sum game $\max_u\min_p\ip{u,\sum_ip_iz_i}$, whose value is $\gamma$ by \cref{lem:bias:duality}. Normalized steepest descent is the Hedge algorithm over the examples played against a best-responding classifier, and \eqref{eq:bias:psigrowth} is the regret bound of Hedge for the losses $\eta_t\ip{u_t,z_i}$. \Textcite{FreundSchapire1999} explained boosting through this game.
\end{remark}

The theorem of \textcite{GunasekarEtAl2018} concerns the unnormalized method $\theta_{t+1}=\theta_t-\eta_t\nabla\Loss(\theta_t)^\sharp$, whose step has length $\eta_t\norm{\nabla\Loss(\theta_t)}_*$ (\cref{ch:steepest}); they prove, for every norm and every initialization, that the normalized margin converges to $\gamma$ when the step sizes are bounded and small enough relative to $1/(R^2\Loss)$. The case of a constant step from $\theta_0=0$ follows from \cref{thm:bias:margin}, with a rate.

\begin{corollary}[Unnormalized steepest descent]\label[corollary]{cor:bias:unnormalized}
Assume \cref{ass:bias:separable}, let $0<\eta\le1/R^2$, $\theta_0=0$, and $\theta_{t+1}=\theta_t-\eta\nabla\Loss(\theta_t)^\sharp$. Then $\Loss(\theta_T)\le1/(1+\eta\gamma^2T/2)$, and whenever the right side is below $\gamma$,
\[
  \gamma-\biasbar(\theta_T)\le R\,\frac{\log n+\frac12+R^2/\gamma^2}{\log(1+\eta\gamma^2T/2)} .
\]
In particular $\biasbar(\theta_T)\to\gamma$, and the limit points of $\theta_T/\norm{\theta_T}$ lie in $\biasU$.
\end{corollary}

\begin{proof}
The plan is to read the method as normalized steepest descent with the adaptive step $\eta_t=\eta\norm{\nabla\Loss(\theta_t)}_*$, which is allowed in \cref{thm:bias:margin}, and to control $\tau_T$ and $\sum\eta_t^2$ through the loss. Since $g^\sharp=-\norm g_*\lmo(g)$, the iteration is \eqref{eq:bias:nsd} with this $\eta_t$, and $\eta_t=\eta r_t\Loss_t$ with $r_t=r(\theta_t)$ and $\Loss_t=\Loss(\theta_t)$. By \cref{lem:bias:step},
\[
  \Loss_{t+1}\le\Loss_t\exp\Bigl(-\eta r_t^2\Loss_t\Bigl(1-\frac{\eta R^2\Loss_t}2\Bigr)\Bigr).
\]
We show by induction that $\Loss_t\le1$, which holds at $t=0$. If $\Loss_t\le1$, then $\eta R^2\Loss_t\le1$, the inner factor is at least $\frac12$, and with $r_t\ge\gamma$ we get $\Loss_{t+1}\le\Loss_t\exp(-\eta\gamma^2\Loss_t/2)\le\Loss_t\le1$. The inequality $e^{-x}\le1/(1+x)$ then gives $1/\Loss_{t+1}\ge1/\Loss_t+\eta\gamma^2/2$, hence $\Loss_T\le1/(1+\eta\gamma^2T/2)$.

Since $q_i(\theta)\le R\norm\theta$, we have $\Loss(\theta)\ge e^{-R\norm\theta}$, so $\tau_T\ge\norm{\theta_T}\ge\log(1/\Loss_T)/R\ge\log(1+\eta\gamma^2T/2)/R$. And since $\eta_t\le\eta R\Loss_t$,
\[
  \frac{R^2}2\sum_{t\ge0}\eta_t^2\le\frac{\eta^2R^4}2\sum_{t\ge0}\frac1{(1+\eta\gamma^2t/2)^2}\le\frac{\eta^2R^4}2\Bigl(1+\frac2{\eta\gamma^2}\Bigr)\le\frac12+\frac{R^2}{\gamma^2},
\]
using $\eta R^2\le1$. Inserting the two bounds in \eqref{eq:bias:bound} gives the rate, and the limit points are handled as in \cref{cor:bias:steps}.
\end{proof}

The rate is logarithmic: to halve the gap one must square the number of steps. For gradient descent, \textcite{SoudryEtAl2018} proved that the direction converges to the support vector machine at the rate $1/\log t$. The reason is visible in the proof: an unnormalized step shrinks with the gradient, the gradient shrinks with the loss, and the distance travelled grows only like $\log t$. A normalized method keeps a step of fixed length, travels a distance proportional to $t$ for a constant step, and reaches the regime of its bias polynomially fast. This is the margin side of the book's recurring theme (\cref{ch:steepest}): normalized updates.

\subsection{Four geometries}

\begin{corollary}[What each optimizer maximizes]\label[corollary]{cor:bias:instances}
Let $\eta_t\to0$ with $\sum_t\eta_t=\infty$, and $\theta_0=0$.
\begin{enumerate}
\item \emph{Normalized gradient descent} ($\ell_2$; $\lmo(g)=-g/\norm g_2$; $R=\max_i\norm{x_i}_2$): $\theta_t/\norm{\theta_t}_2$ converges to the hard-margin support vector machine direction, the unique element of $\biasU$.
\item \emph{Sign descent} ($\ell_\infty$; $\lmo(g)=-\sign(g)$, with every value in $[-1,1]$ allowed in coordinates where $g$ vanishes; $R=\max_i\norm{x_i}_1$): the normalized $\ell_\infty$ margin converges to $\gamma_\infty$, the value of the linear program $\max\{s:\ip{u,z_i}\ge s\ \text{for all}\ i,\ \norm u_\infty\le1\}$.
\item \emph{Greedy coordinate descent} ($\ell_1$; $\lmo(g)=-\sign(g_j)e_j$ for some $j\in\argmax_j\abs{g_j}$; $R=\max_i\norm{x_i}_\infty$): the normalized $\ell_1$ margin converges to $\gamma_1$, the value of a linear program; with weak classifiers as features this is boosting with small steps.
\item \emph{Orthogonalized updates} (spectral norm on $\R^{K\times d}$; $\lmo(G)=-\polar(G)$; $R=\max_i\norm{X_i}_{\nuc}$): the normalized spectral margin converges to $\gamma_{\op}=\max\{\min_i\ip{W,Z_i}:\norm W_{\op}\le1\}$, the value of a semidefinite program.
\end{enumerate}
In (ii)--(iv), every limit point of the normalized iterates is a maximum-margin classifier, and the iterates converge in direction when $\biasU$ is a single point.
\end{corollary}

\begin{proof}
The linear minimization oracles and dual norms are those of \cref{ch:steepest,ch:sign,ch:spectral} (the spectral one through von Neumann's trace inequality, \cref{thm:toolkit:vonneumann}), so everything but the uniqueness in (i) is \cref{cor:bias:steps}(i). If $u\ne u'$ both lay in $\biasU$ for the $\ell_2$ norm, then $v=(u+u')/2$ would have $\biasqmin(v)\ge\gamma$ and $\norm v_2<1$, because the Euclidean ball is strictly convex, and $v/\norm v_2$ would have a margin larger than $\gamma$.
\end{proof}

A multiclass linear classifier fits case (iv). For $K$ classes and $W\in\R^{K\times d}$ with rows $w_c$, the margin of example $i$ is $\min_{c\ne y_i}(w_{y_i}-w_c)^\top x_i=\min_{c\ne y_i}\ip{W,Z_{ic}}$ with the rank-one matrices $Z_{ic}=(e_{y_i}-e_c)x_i^\top$, whose nuclear norm is $\sqrt2\norm{x_i}_2$. The multiclass exponential loss $\frac1{n(K-1)}\sum_i\sum_{c\ne y_i}e^{-\ip{W,Z_{ic}}}$, a relative of the cross-entropy $\log(1+\sum_{c\ne y_i}e^{-\ip{W,Z_{ic}}})$ with the same tail, is the loss \eqref{eq:bias:loss} for the $n(K-1)$ signed examples $Z_{ic}$. So orthogonalized updates on the last layer maximize its spectral-norm margin; in practice Muon is used for hidden layers and Adam for the output layer (\cref{ch:spectral}), which is one reason the statement should be read as a model, not as a description.

The $\ell_1$ and $\ell_\infty$ problems are linear programs, and their solutions have a shape that the Euclidean one lacks.

\begin{proposition}[Sparse and dense maximum margins]\label[proposition]{prop:bias:vertices}
Let $V=\R^d$. (i) Some $\ell_1$-maximum-margin classifier has at most $n$ nonzero coordinates. (ii) Some $\ell_\infty$-maximum-margin classifier has at least $d+1-n$ coordinates equal to $\pm1$. (iii) If $\biasU$ is a single point, that point has the property of (i), respectively (ii).
\end{proposition}

\begin{proof}
The plan is to write each problem as a linear program and use that a linear program over a polyhedron without lines, with a finite optimum, has an optimal vertex. For $\ell_\infty$ the program is $\max s$ over $(u,s)\in\R^{d+1}$ subject to $\ip{u,z_i}\ge s$ ($i\le n$) and $-1\le u_j\le1$ ($j\le d$). Its feasible set contains no line, because the box bounds $u$ and then $s$ is bounded above, and the optimum is finite. At a vertex, $d+1$ linearly independent constraints are active, of which at most $n$ are margin constraints, so at least $d+1-n$ are box constraints, and no coordinate can have both of its box constraints active: at least $d+1-n$ coordinates equal $\pm1$. For $\ell_1$, write $u=u^+-u^-$ with $u^\pm\ge0$; the program is $\max s$ subject to $\ip{u^+-u^-,z_i}\ge s$, $\sum_j(u^+_j+u^-_j)\le1$ and $u^\pm\ge0$, in $2d+1$ variables. Its feasible set contains no line, and at a vertex $2d+1$ linearly independent constraints are active, at most $n+1$ of them margin or norm constraints; so at least $2d-n$ of the $2d$ nonnegativity constraints are active, at most $n$ of the numbers $u^\pm_j$ are nonzero, and $u$ has at most $n$ nonzero coordinates. Every optimal $(u^+,u^-,s)$ gives a classifier in $\biasU$, because $\norm{u^+-u^-}_1\le\sum_j(u^+_j+u^-_j)$.

For (iii), suppose $\biasU=\{u^\star\}$. In the $\ell_\infty$ program, every optimum has $u=u^\star$ and $s=\gamma$, so the optimum is unique and is the optimal vertex. In the $\ell_1$ program, every optimum has $u^+-u^-=u^\star$ with $\norm{u^\star}_1=1$, and $\sum_j(u_j^++u_j^-)\le1=\norm{u^+-u^-}_1$ forces $u^+$ and $u^-$ to be the positive and negative parts of $u^\star$: again the optimum is unique and is a vertex.
\end{proof}

So when $d\gg n$, the end point of sign descent saturates all but at most $n-1$ of its weights at $\pm1$: the $\ell_\infty$ geometry charges nothing for using a feature at full weight, and every feature that helps the worst-case margin, however little, gets full weight. The end point of coordinate descent uses at most $n$ features. The end point of gradient descent lies in the span of the support vectors (\cref{lem:bias:duality}(i): in the $\ell_2$ case $u^\star=c^\star/\norm{c^\star}_2$). These are three different priors on the interpolating solution, and the next example shows that the choice can decide the test error.

\begin{example}[A dataset on which the $\ell_\infty$ bias fails]\label[example]{ex:bias:wilson}
Take $n_+\ge1$ positive and $n_-\ge1$ negative examples, in the spirit of the construction of \textcite{WilsonEtAl2017}. Each example has three shared features and some features of its own. Let $e_{(i)}$ be the indicator vector of the features owned by example $i$: one feature for a positive example and five for a negative one, disjoint across examples. The examples are
\[
  x_i=\bigl(y_i,\ 1,\ 1,\ e_{(i)}\bigr)\in\R^{3+n_++5n_-}.
\]
The first feature is the label, the second and third are constant, and the owned features let every example be memorized. A new example with label $y$ has the same three shared features, $(y,1,1)$, and features of its own that no training example has. Every classifier trained on the sample puts weight zero on those (the maximum-margin classifiers live on the training features, and a steepest descent method started at zero never moves a coordinate in which every gradient vanishes), so it scores a new example by $y\,\theta_1+\theta_2+\theta_3$. New labels are $+1$ with probability $p\in(0,1)$.
\end{example}

\begin{proposition}[The two biases on the example]\label[proposition]{prop:bias:wilson}
On the dataset of \cref{ex:bias:wilson}, for every $n_+,n_-\ge1$:
\begin{enumerate}
\item The $\ell_\infty$ max-margin set is the single point $u_\infty$ with $u_1=u_2=u_3=1$, $u=1$ on the features owned by positive examples and $u=-1$ on those owned by negative examples, and $\gamma_\infty=4$. It scores every new example by $y+2>0$, so its test error is $1-p$.
\item Put $P=3+1/n_+$, $Q=3+5/n_-$, $A_+=(Q+1)/(PQ-1)$ and $A_-=(P+1)/(PQ-1)$. The $\ell_2$ max-margin direction is that of $\theta_2=\sum_i\alpha_iz_i$ with $\alpha_i=A_+/n_+$ for positive and $\alpha_i=A_-/n_-$ for negative examples. Every training example is a support vector, and $\theta_2$ is also the minimum-norm solution of $\ip{\theta,x_i}=y_i$ for all $i$. It scores a new positive example by $3A_+-A_->0$ and a new negative one by $A_+-3A_-<0$, so its test error is $0$.
\item The $\ell_1$ max-margin set is $\{e_1\}$, with $\gamma_1=1$; its test error is $0$.
\end{enumerate}
Consequently, from $\theta_0=0$ with $\eta_t\to0$ and $\sum\eta_t=\infty$, the test error of sign descent tends to $1-p$ and those of normalized gradient descent and of coordinate descent tend to $0$. If $n_+>n_-$, sign descent and Adam with $\epsilon=0$, for all $\beta_1,\beta_2$ and all step sizes, move along the ray $\{su_\infty:s>0\}$ from the first step on, so their test error is $1-p$ at every step.
\end{proposition}

\begin{proof}
The plan is to identify each maximizer by a direct bound for $\ell_\infty$ and $\ell_1$ and by the KKT conditions for $\ell_2$, and then to follow sign descent and Adam along the ray.

(i) For $\norm u_\infty\le1$ and a positive example $i$, $q_i(u)=u_1+u_2+u_3+u_{(i)}\le4$, with equality only if all four coordinates equal $1$. For a negative example $j$ and such a $u$, $q_j(u)=u_1-u_2-u_3-\sum_{k\in(j)}u_k=-1-\sum_{k\in(j)}u_k$, which is at least $4$ only if $u_k=-1$ on the five features of $j$. Since $u_\infty$ gives every example the margin $4$, it is the only point of $\biasU$, and $\gamma_\infty=4$. It scores a new example by $y+1+1\ge1$.

(ii) The $\ell_2$ problem is equivalent to minimizing $\frac12\norm\theta_2^2$ subject to $q_i(\theta)\ge1$, whose solution $\theta_2$ has the direction of the $\ell_2$ max-margin classifier. Its KKT conditions, $\theta=\sum_i\alpha_iz_i$ with $\alpha\ge0$, $q_i(\theta)\ge1$ and $\alpha_i(q_i(\theta)-1)=0$, are sufficient for this convex problem. With the stated $\alpha$, the coordinates of $\theta=\sum_i\alpha_iz_i$ are $a=A_++A_-$ in the first feature, $c=A_+-A_-$ in the second and third, $A_+/n_+$ in each feature owned by a positive example and $-A_-/n_-$ in each feature owned by a negative one. The margins of the positive and of the negative examples are therefore
\[
  a+2c+\frac{A_+}{n_+}=PA_+-A_-\qquad\text{and}\qquad a-2c+\frac{5A_-}{n_-}=QA_--A_+ ,
\]
and the stated $A_\pm$ solve the system that sets both to $1$. They are positive because $P,Q>3$, so the KKT conditions hold with every constraint active; the solution is unique because the objective is strictly convex. Since all constraints are active, $\theta_2$ also minimizes the norm subject to $q_i(\theta)=1$, that is, to $\ip{\theta,x_i}=y_i$. The scores of new examples are $a+2c=3A_+-A_-=(3Q+2-P)/(PQ-1)$, positive because $3Q>9>P$, and $-a+2c=A_+-3A_-=(Q-3P-2)/(PQ-1)$, negative because $Q\le8<3P+2$.

(iii) With uniform weights, $\bar z=\frac1n\sum_iz_i$ has first coordinate $1$ and all other coordinates of absolute value less than $1$, so $\gamma_1\le\norm{\bar z}_\infty=1$ by \cref{lem:bias:duality}, and $e_1$ gives every example the margin $1$. If $u\in\biasU$, take a positive $i$ and a negative $j$; then $q_i(u)+q_j(u)=2u_1+u_{(i)}-\sum_{k\in(j)}u_k\ge2$, while the left side is at most $2\abs{u_1}+\abs{u_{(i)}}+\sum_{k\in(j)}\abs{u_k}\le2\norm u_1\le2$, with equality only if $u=e_1$. It scores a new example by $y$.

For the consequences, \cref{cor:bias:instances} gives convergence of the directions to $u_\infty$, $\theta_2/\norm{\theta_2}_2$ and $e_1$, and the score of a new example is a continuous function of the direction that does not vanish at these limits. Finally, let $n_+>n_-$ and $\theta=su_\infty$ with $s\ge0$. All margins equal $4s$, the weights $p_i(\theta)$ are uniform, and $\nabla\Loss(\theta)=-\Loss(\theta)\,\bar z$, where $\sign(\bar z)=u_\infty$: the second and third coordinates of $\bar z$ are $(n_+-n_-)/n>0$ and the owned features carry the signs of their owners. Sign descent therefore moves to $(s+\eta_t)u_\infty$. For Adam with $\epsilon=0$, if all gradients so far are positive multiples of $-\bar z$, then in coordinate $j$ the moments are $m_{t,j}=-\bar z_jM_t$ and $v_{t,j}=\bar z_j^2V_t$ for scalars $M_t,V_t>0$ common to all coordinates, so the update $-\eta_t\hat m_t/\sqrt{\hat v_t}$ is a positive multiple of $\sign(\bar z)=u_\infty$; by induction the iterates stay on the ray.
\end{proof}

The $\ell_\infty$ geometry gives the two constant features full weight because they help the positive examples, which own one feature each, to catch up with the negative ones, which own five; on a new example the constant features then outvote the label. The $\ell_2$ geometry gives each feature a weight proportional to its usefulness, and the constant features get the small weight $c$. Nothing in this is special to $\ell_\infty$ being bad: with a dense ground truth the roles reverse (Exercise~\ref{ex:bias:dense}). In the original construction the adaptive methods are trained on the square loss, and the solution they reach is proportional to $\sign(X^\top y)$, which for $n_+>n_-$ is our $u_\infty$; gradient descent reaches the minimum-norm interpolant, which is our $\theta_2$. On this dataset the square loss and the exponential loss end at the same classifiers, and our version needs no condition on $n_+$ and $n_-$ for the limits.

\section{The theory: weight decay, adaptive methods and networks}\label{sec:bias:beyond}

\subsection{Weight decay stops the march}

With decoupled weight decay the iterates no longer run to infinity. \Cref{ch:decay} shows that normalized steepest descent with decoupled weight decay $\lambda$,
\[
  \theta_{t+1}=(1-\eta_t\lambda)\theta_t+\eta_t\lmo\bigl(\nabla\Loss(\theta_t)\bigr),
\]
is the Frank--Wolfe method on the ball $\{\norm\theta\le1/\lambda\}$, which for a convex objective that is smooth on the ball, and the Frank--Wolfe schedule $\eta_t\lambda=2/(t+2)$, converges to the minimum of the objective over the ball. The exponential loss is convex and smooth on bounded sets, so AdamW without memory ends at a minimizer of $\Loss$ over the $\ell_\infty$ ball of radius $1/\lambda$. Where is that?

\begin{proposition}[The end point at radius $B$]\label[proposition]{prop:bias:decay}
Assume \cref{ass:bias:separable}, let $B>0$, and let $\theta_B$ minimize $\Loss$ over $\{\norm\theta\le B\}$. Then $\norm{\theta_B}=B$ and $\Loss(\theta_B)\le e^{-\gamma B}$, and if $B>\log(n)/\gamma$,
\[
  \gamma-\frac{\log n}B\le\biasbar(\theta_B)\le\gamma .
\]
\end{proposition}

\begin{proof}
If $\norm{\theta_B}<B$, then $\theta_B$ would be a local and hence, by convexity, a global minimizer of $\Loss$ on $V$, which does not exist (\cref{lem:bias:gradient}). Comparing with $Bu^\star$ for some $u^\star\in\biasU$ gives $\Loss(\theta_B)\le\Loss(Bu^\star)\le e^{-\gamma B}$. Then $e^{-\biasqmin(\theta_B)}\le n\Loss(\theta_B)\le ne^{-\gamma B}$, so $\biasqmin(\theta_B)\ge\gamma B-\log n$, and dividing by $\norm{\theta_B}=B$ gives the lower bound.
\end{proof}

So decoupled weight decay replaces the endless march along a direction by a point at radius $1/\lambda$, whose margin is within $\lambda\log n$ of the maximum: AdamW without memory ends within $\lambda\log n$ of the $\ell_\infty$ maximum margin, and its end point approaches the $\ell_\infty$ maximum-margin direction as $\lambda\to0$. That the norm-constrained or regularized minimizers of the exponential and logistic losses approach the maximum-margin direction as the constraint is relaxed is the \term{regularization path}\index{regularization path} picture of \textcite{RossetZhuHastie2004}; the proposition gives it a rate. The same proposition covers the ordinary $L_2$ penalty (Exercise~\ref{ex:bias:penalty}).

\subsection{Adam, AdamW and Muon}

The theorem covers the memoryless versions of the optimizers. For the versions with memory, the following has been proved; we state the results and do not prove them.
\begin{itemize}
\item \emph{Momentum.} Heavy-ball and Nesterov momentum on top of gradient descent do not change the limit: it is still the $\ell_2$ maximum-margin direction \parencite{WangEtAl2022}. Averaging changes the speed, not the norm.
\item \emph{Adam.} On separable data with the logistic loss and decaying step sizes, and with the $\epsilon$ in the denominator neglected, Adam converges in direction to an $\ell_\infty$ maximum-margin classifier, in polynomial time for a general class of decaying step sizes \parencite{ZhangZouCao2024}. With a fixed $\epsilon>0$ the limit is different: the gradients on separable data tend to zero, $\sqrt{\hat v_t}$ eventually falls below $\epsilon$, and Adam becomes gradient descent with momentum and a rescaled step; \textcite{WangEtAl2021} proved that RMSProp and Adam then maximize the $\ell_2$ margin, for homogeneous networks and in particular for linear models, while AdaGrad, whose preconditioner keeps the whole history, does not.
\item \emph{AdamW.} If AdamW converges with non-increasing step sizes whose sum diverges, its limit is a KKT point of minimizing the loss subject to $\norm\theta_\infty\le1/\lambda$ \parencite{XieLi2024}; for a linear model on separable data this is the end point of \cref{prop:bias:decay}.
\item \emph{Spectral descent and Muon.} For multiclass separable data and the cross-entropy loss, normalized steepest descent in the spectral norm and Muon converge to the spectral-norm maximum-margin classifier \parencite{FanSchmidtThrampoulidis2025}.
\item \emph{Mirror descent} with the potential $\norm\theta_p^p$ maximizes the $\ell_p$ margin, a way to reach a chosen geometry with a smooth potential \parencite{SunEtAl2022}.
\end{itemize}

\subsection{Homogeneous networks}

The linear analysis extends to networks whose output is a homogeneous function of the parameters: a ReLU network without biases, of depth $a$, satisfies $f(x;c\theta)=c^af(x;\theta)$ for $c>0$. Write $q_i(\theta)=y_if(x_i;\theta)$ and define $\Loss$ and $\psi$ by \eqref{eq:bias:loss} and \eqref{eq:bias:psi} with these margins. The witness $u^\star$ of the linear case is no longer available, but Euler's identity for homogeneous functions takes its place.

\begin{proposition}[The smoothed margin never decreases]\label[proposition]{prop:bias:homogeneous}
Let $f(x;\theta)$ be continuously differentiable in $\theta\in V$ and homogeneous of degree $a>0$. Let $\theta(t)$, $t\ge t_0$, be absolutely continuous with $\norm{\dot\theta(t)}\le1$ and $\ip{\nabla\Loss(\theta(t)),\dot\theta(t)}=-\norm{\nabla\Loss(\theta(t))}_*$ for almost every $t$, a normalized steepest flow in the norm $\norm{\cdot}$, and suppose that $n\Loss(\theta(t_0))<1$. Then the smoothed normalized margin $\tilde\gamma(t)=\psi(\theta(t))/\norm{\theta(t)}^a$ is nondecreasing on $[t_0,\infty)$, and $\biasqmin(\theta(t))/\norm{\theta(t)}^a\ge\tilde\gamma(t_0)>0$ for every $t\ge t_0$.
\end{proposition}

\begin{proof}
The plan is to show that $\log\psi$ grows at least as fast as $a\log\norm\theta$. Differentiating $f(x;c\theta)=c^af(x;\theta)$ at $c=1$ gives Euler's identity $\ip{\nabla q_i(\theta),\theta}=a\,q_i(\theta)$. Hence, using $\biasqmin\ge\psi$ from \eqref{eq:bias:psi},
\[
  \ip{-\nabla\Loss(\theta),\theta}=\frac an\sum_ie^{-q_i(\theta)}q_i(\theta)\ge a\Loss(\theta)\biasqmin(\theta)\ge a\Loss(\theta)\psi(\theta),
\]
and the definition of the dual norm gives $\norm{\nabla\Loss(\theta)}_*\ge\ip{-\nabla\Loss(\theta),\theta}/\norm\theta\ge a\Loss(\theta)\psi(\theta)/\norm\theta$ whenever $\theta\ne0$. Along the flow, for almost every $t$,
\[
  \frac{d}{dt}\psi(\theta(t))=-\frac{\ip{\nabla\Loss(\theta),\dot\theta}}{\Loss(\theta)}=\frac{\norm{\nabla\Loss(\theta)}_*}{\Loss(\theta)}\ge\frac{a\,\psi(\theta)}{\norm\theta}.
\]
So $\psi$ is nondecreasing while it is positive; it is positive at $t_0$, so it stays at least $\psi(\theta(t_0))>0$, and then $\theta(t)$ stays away from $0$, because $\abs{q_i(\theta)}\le M\norm\theta^a$ with $M$ the maximum of $\abs{f(x_i;\cdot)}$ over the unit sphere. The function $t\mapsto\norm{\theta(t)}$ is absolutely continuous with derivative at most $\norm{\dot\theta(t)}\le1$. Therefore, for almost every $t$,
\[
  \frac{d}{dt}\log\psi(\theta(t))\ge\frac a{\norm{\theta(t)}}\ge\frac a{\norm{\theta(t)}}\,\frac{d}{dt}\norm{\theta(t)}=\frac{d}{dt}\Bigl(a\log\norm{\theta(t)}\Bigr),
\]
so $\log\tilde\gamma(t)=\log\psi(\theta(t))-a\log\norm{\theta(t)}$ is an absolutely continuous function with a nonnegative derivative almost everywhere, hence nondecreasing. The last claim follows from $\biasqmin\ge\psi$.
\end{proof}

For linear models ($a=1$) the proposition is the continuous-time shadow of \cref{thm:bias:margin}. For networks it is the first half of the known picture.

\begin{theorem*}[\textcite{LyuLi2020}; \textcite{JiTelgarsky2020}; \textcite{TsilivisVardiKempe2025}]
For homogeneous networks trained by gradient flow on the exponential or logistic loss from a point where $n\Loss<1$, so that every training example is classified correctly, the smoothed normalized margin increases, the loss tends to zero, and every limit point of $\theta/\norm\theta_2$ is the direction of a KKT point of $\min\frac12\norm\theta_2^2$ subject to $q_i(\theta)\ge1$ for all $i$ (Lyu and Li, who also treat ReLU networks through Clarke subdifferentials, and gradient descent with small steps); under a definability condition the direction converges (Ji and Telgarsky). For steepest flows in other norms, the margin measured in that norm increases and the limit points are KKT points of the corresponding max-margin problem (Tsilivis, Vardi and Kempe).
\end{theorem*}

\begin{proofsketch}
\Cref{prop:bias:homogeneous} is the monotonicity for differentiable networks. At a KKT point $\bar\theta$, rescaled so that $\biasqmin(\bar\theta)=1$, the parameter vector is a nonnegative combination of the gradients $\nabla q_i(\bar\theta)$ of the examples that attain the minimum margin. Along the flow, $-\nabla\Loss=\frac1n\sum_ie^{-q_i}\nabla q_i$ is such a combination, with weights concentrating on the examples of smallest margin as $\norm\theta\to\infty$. It therefore suffices to show that $\theta$ and $-\nabla\Loss$ align along a sequence of times. The inequality in the proof of \cref{prop:bias:homogeneous} is an equality only when they are aligned, and since $\tilde\gamma$ is bounded by the maximum normalized margin, the integrated deficit is finite, which forces alignment along a sequence. Omitted: the quantitative version of both steps (approximate KKT points and their limits), which is the core of \textcite{LyuLi2020}, and its extension to dual norms, which is in \textcite{TsilivisVardiKempe2025}.
\end{proofsketch}

A KKT point is a stationary point of the margin problem, not necessarily a maximizer. \Textcite{VardiShamirSrebro2022} show that for ReLU networks the limit need not be even a local maximizer of the margin, while for some linear networks it is.

\section{Reading the theory}\label{sec:bias:reading}

\paragraph{Three norms, three classifiers.} \Cref{fig:bias:geometry} shows the duality of \cref{lem:bias:duality} on six examples in the plane. The three maximum-margin classifiers are different, and each is determined by a different support set: the $\ell_1$ classifier uses one feature, the $\ell_\infty$ classifier gives both features full weight, and the $\ell_2$ classifier lies between them, determined by two support vectors. In the plane of signed examples, each classifier is the normal of the line that separates the convex hull from the largest dual ball around the origin.

\begin{figure}[t]
\begin{fullwidth}
\centering
\begin{tikzpicture}
\begin{groupplot}[group style={group size=2 by 1,horizontal sep=18mm},wplot,width=66mm,height=66mm,axis equal image]
\nextgroupplot[xlabel={feature $1$},ylabel={feature $2$},xmin=-1.05,xmax=1.05,ymin=-1.05,ymax=1.05,xtick={-1,0,1},ytick={-1,0,1},title={examples and classifiers}]
\addplot[theory3] table[x=sepl1x,y=sepl1y] {figures/data/15-bias-geometry-lines.dat};
\addplot[theory] table[x=sepl2x,y=sepl2y] {figures/data/15-bias-geometry-lines.dat};
\addplot[theory2] table[x=seplinfx,y=seplinfy] {figures/data/15-bias-geometry-lines.dat};
\addplot[only marks,mark=*,mark size=2.2pt,color=sea,mark options={fill=sea,draw=white}] table[x=x1,y=x2,restrict expr to domain={\thisrow{y}}{0.5:1.5}] {figures/data/15-bias-geometry-points.dat};
\addplot[only marks,mark=square*,mark size=2.0pt,color=seadeep,mark options={fill=seadeep,draw=white}] table[x=x1,y=x2,restrict expr to domain={\thisrow{y}}{-1.5:-0.5}] {figures/data/15-bias-geometry-points.dat};
\legend{$\ell_1$,$\ell_2$,$\ell_\infty$,positive,negative}
\nextgroupplot[xlabel={$z_1$},ylabel={$z_2$},xmin=-0.15,xmax=1.05,ymin=-0.6,ymax=0.6,xtick={0,0.5,1},ytick={-0.5,0,0.5},title={signed examples and dual balls}]
\addplot[fill=seafoam,draw=sea!60,line width=0.5pt] table[x=z1,y=z2] {figures/data/15-bias-geometry-hull.dat};
\addplot[sandline!55!black,line width=0.8pt] table[x=sx,y=sy] {figures/data/15-bias-geometry-balls.dat};
\addplot[sea,line width=0.8pt] table[x=cx,y=cy] {figures/data/15-bias-geometry-balls.dat};
\addplot[seadeep,line width=0.8pt,dashed] table[x=dx,y=dy] {figures/data/15-bias-geometry-balls.dat};
\addplot[guide] table[x=supl1x,y=supl1y] {figures/data/15-bias-geometry-lines.dat};
\addplot[guide] table[x=supl2x,y=supl2y] {figures/data/15-bias-geometry-lines.dat};
\addplot[guide] table[x=suplinfx,y=suplinfy] {figures/data/15-bias-geometry-lines.dat};
\addplot[-stealth,sandline!55!black,line width=0.9pt] table[x=dirl1x,y=dirl1y] {figures/data/15-bias-geometry-lines.dat};
\addplot[-stealth,sea,line width=0.9pt] table[x=dirl2x,y=dirl2y] {figures/data/15-bias-geometry-lines.dat};
\addplot[-stealth,seadeep,line width=0.9pt] table[x=dirlinfx,y=dirlinfy] {figures/data/15-bias-geometry-lines.dat};
\addplot[only marks,mark=*,mark size=1.6pt,color=ink] table[x=z1,y=z2] {figures/data/15-bias-geometry-points.dat};
\end{groupplot}
\end{tikzpicture}
\end{fullwidth}
\caption{Three geometries on one dataset (script \texttt{15-bias-geometry.py}). Left: six examples (circles positive, squares negative) and the three maximum-margin classifiers through the origin: $\ell_1$ (solid, the vertical line: one feature), $\ell_2$ (solid, between), $\ell_\infty$ (dashed: both features with full weight). Right: the signed examples $z_i=y_ix_i$, their convex hull, and the largest dual balls around the origin that it does not enter: the $\ell_\infty$ square for the $\ell_1$ margin ($\gamma_1=0.42$), the circle for the $\ell_2$ margin ($\gamma_2=0.506$) and the $\ell_1$ diamond for the $\ell_\infty$ margin ($\gamma_\infty=0.57$). Dashed: the separating lines $\ip{u^\star,c}=\gamma$; arrows: the maximum-margin directions $u^\star$, normal to them. The values agree with the linear and quadratic programs to $10^{-6}$.}
\label{fig:bias:geometry}
\end{figure}

\paragraph{Each optimizer maximizes its own margin.} \Cref{fig:bias:margins} runs the four optimizers of \cref{cor:bias:instances} from zero. In each panel the points rise to the maximum margin of the optimizer's own norm, computed by a linear or quadratic program or, for the spectral norm, by cutting planes on the dual problem of \cref{lem:bias:duality}. The decaying-step bound of \cref{cor:bias:steps}(ii) lies below its points and reaches the line slowly; the constant-step bound stalls at $\gamma-R^2\eta/2$, a gap of $0.025R$ here, while the constant-step points do not stall. \Cref{tab:bias:cross} measures the end points in all norms: each optimizer comes within $2\cdot10^{-3}$ of the maximum margin of its own norm and stays well below the maximum margins of the others; coordinate descent, for instance, ends with an $\ell_\infty$ margin $11\%$ below $\gamma_\infty$.

\begin{figure}[t]
\begin{fullwidth}
\centering
\begin{tikzpicture}
\begin{groupplot}[group style={group size=2 by 2,horizontal sep=17mm,vertical sep=19mm},wplot,width=80mm,height=50mm,xmode=log,xmin=1,xmax=1e5,xlabel={step $t$},ylabel={normalized margin}]
\nextgroupplot[title={$\ell_2$: normalized gradient descent},ymin=0,ymax=0.34,legend style={at={(0.98,0.04)},anchor=south east,font=\scriptsize},legend cell align=left]
\addplot[theory] table[x=t,y=l2] {figures/data/15-bias-margins-gamma.dat};
\addplot[theory3,restrict y to domain=-0.5:1] table[x=t,y=l2decb] {figures/data/15-bias-margins.dat};
\addplot[theory2,restrict y to domain=-0.5:1] table[x=t,y=l2conb] {figures/data/15-bias-margins.dat};
\addplot[simulation] table[x=t,y=l2dec] {figures/data/15-bias-margins.dat};
\addplot[simulation2] table[x=t,y=l2con] {figures/data/15-bias-margins.dat};
\legend{maximum margin $\gamma$,bound: decaying,bound: constant,run: decaying,run: constant}
\nextgroupplot[title={$\ell_\infty$: sign descent},ymin=0,ymax=0.41]
\addplot[theory] table[x=t,y=linf] {figures/data/15-bias-margins-gamma.dat};
\addplot[theory3,restrict y to domain=-0.5:1] table[x=t,y=linfdecb] {figures/data/15-bias-margins.dat};
\addplot[theory2,restrict y to domain=-0.5:1] table[x=t,y=linfconb] {figures/data/15-bias-margins.dat};
\addplot[simulation] table[x=t,y=linfdec] {figures/data/15-bias-margins.dat};
\addplot[simulation2] table[x=t,y=linfcon] {figures/data/15-bias-margins.dat};
\nextgroupplot[title={$\ell_1$: coordinate descent},ymin=0,ymax=0.19]
\addplot[theory] table[x=t,y=l1] {figures/data/15-bias-margins-gamma.dat};
\addplot[theory3,restrict y to domain=-0.5:1] table[x=t,y=l1decb] {figures/data/15-bias-margins.dat};
\addplot[theory2,restrict y to domain=-0.5:1] table[x=t,y=l1conb] {figures/data/15-bias-margins.dat};
\addplot[simulation] table[x=t,y=l1dec] {figures/data/15-bias-margins.dat};
\addplot[simulation2] table[x=t,y=l1con] {figures/data/15-bias-margins.dat};
\nextgroupplot[title={spectral: orthogonalized updates},ymin=0,ymax=0.25]
\addplot[theory] table[x=t,y=op] {figures/data/15-bias-margins-gamma.dat};
\addplot[theory3,restrict y to domain=-0.5:1] table[x=t,y=opdecb] {figures/data/15-bias-margins.dat};
\addplot[theory2,restrict y to domain=-0.5:1] table[x=t,y=opconb] {figures/data/15-bias-margins.dat};
\addplot[simulation] table[x=t,y=opdec] {figures/data/15-bias-margins.dat};
\addplot[simulation2] table[x=t,y=opcon] {figures/data/15-bias-margins.dat};
\end{groupplot}
\end{tikzpicture}
\end{fullwidth}
\caption{Normalized margins of normalized steepest descent in four norms, from $\theta_0=0$, with decaying steps $\eta_t=1/(R\sqrt{t+1})$ (circles) and a constant step $\eta=0.05/R$ (squares), $R=\max_i\norm{z_i}_*$. Horizontal lines: the exact maximum margins. Curves: the lower bound \eqref{eq:bias:bound} of \cref{thm:bias:margin} for each schedule. Vector panels: $n=16$ separable examples in $\R^5$ on the unit sphere; spectral panel: a linear classifier $W\in\R^{3\times4}$ on $12$ examples in three classes, as $24$ signed matrices of nuclear norm one. One run per schedule, fixed seeds (script \texttt{15-bias-margins.py}).}
\label{fig:bias:margins}
\end{figure}

\begin{table}[t]
\centering\small
\begin{tabular}{@{}lccc@{}}
\toprule
end point of & \multicolumn{3}{c}{normalized margin measured in} \\
\cmidrule(l){2-4}
 & $\ell_2$ & $\ell_\infty$ & $\ell_1$ \\
\midrule
normalized gradient descent & $\mathbf{0.3108}$ & $0.3666$ & $0.1673$\\
sign descent & $0.3020$ & $\mathbf{0.3750}$ & $0.1598$\\
coordinate descent & $0.2921$ & $0.3350$ & $\mathbf{0.1709}$\\
maximum margin $\gamma$ & $0.3115$ & $0.3766$ & $0.1717$\\
\midrule
 & Frobenius & max-entry & spectral\\
\midrule
normalized gradient descent & $\mathbf{0.2085}$ & $0.2891$ & $0.2231$\\
sign descent & $0.1855$ & $\mathbf{0.3451}$ & $0.2099$\\
orthogonalized updates & $0.2067$ & $0.2992$ & $\mathbf{0.2276}$\\
maximum margin $\gamma$ & $0.2094$ & $0.3467$ & $0.2286$\\
\bottomrule
\end{tabular}
\caption{The end points of \cref{fig:bias:margins} (decaying steps, $T=10^5$) measured in every norm, for the vector data (top) and the matrix data (bottom; ``max-entry'' is the $\ell_\infty$ norm of the entries). Each optimizer approaches the maximum margin of its own norm (bold) and not those of the others.}
\label{tab:bias:cross}
\end{table}

\paragraph{The distance travelled is the clock.} The gaps $\gamma-\biasbar(\theta_t)$ of the runs in \cref{fig:bias:margins} obey a simpler law than the theorem: at $t=10^5$, the gap times the distance travelled, $(\gamma-\biasbar(\theta_t))\,\tau_t$, is $0.46$ for $\ell_2$, $0.49$--$0.51$ for $\ell_\infty$, $0.57$ for $\ell_1$ and $0.68$ for the spectral norm, and it is the same for the two schedules within $4\%$, although their distances differ by a factor of eight. So the margin is a function of the distance travelled, as for the flow, and it approaches the maximum like $C/\tau_t$ with $C$ well below the flow bound $\log n=2.77$ (for the spectral data $\log24=3.18$). The finite-step term $R^2\sum\eta_t^2/(2\tau_t)$ of \eqref{eq:bias:bound} is a price that these runs do not pay. The reason, as far as we can see, is that near the maximum-margin direction a step changes the small margins by nearly the same amount, so the variance that Hoeffding's lemma charges for is small, and what is left cancels over consecutive steps. Whether this holds for every dataset and every norm is open as far as we know (Exercise~\ref{ex:bias:constant}); for the $\ell_2$ norm, where the dual problem is smooth, \textcite{JiTelgarsky2021} prove an $O(1/t)$ margin rate for gradient descent with aggressive step sizes.

The law turns into time scales. A normalized method with a constant step $\eta$ comes within $\eps$ of its maximum margin after about $C/(\eps\eta)$ steps; with $\eta_t\propto1/\sqrt t$ after about $(C/\eps)^2$; unnormalized gradient descent with a constant step travels only a distance of order $\log t$ (\cref{cor:bias:unnormalized}) and needs a number of steps exponential in $1/\eps$. On separable data, Adam, sign descent and Muon are in the regime of their implicit bias after polynomially many steps; plain gradient descent creeps toward its own.

\paragraph{Where AdamW ends.} \Cref{fig:bias:decay} follows \cref{prop:bias:decay} on the vector data. The exact end point at radius $B=1/\lambda$, computed by constrained minimization, and AdamW without memory, run with the Frank--Wolfe schedule, agree to $2\cdot10^{-4}$. At $\lambda=0.1$, the default of language-model training (\cref{ch:decay}), the end point has $\ell_\infty$ margin $0.340$, about $10\%$ below $\gamma_\infty=0.377$; the guaranteed $\gamma_\infty-\lambda\log n$ is $0.099$ there. Over the range shown the exact gap is close to $0.4/B$, so the bound's $\log n/B$ has the right order in $B$ and a constant about seven times too large.

\begin{figure}[t]
\centering
\begin{tikzpicture}
\begin{axis}[wplot,xmode=log,xmin=1,xmax=300,ymin=0,ymax=0.42,xlabel={radius $B=1/\lambda$ of the $\ell_\infty$ ball},ylabel={$\ell_\infty$ margin of the end point},legend style={at={(0.98,0.04)},anchor=south east,font=\scriptsize}]
\addplot[guide] table[x=B,y=gamma] {figures/data/15-bias-decay.dat};
\addplot[theory] table[x=B,y=exact] {figures/data/15-bias-decay.dat};
\addplot[theory2,restrict y to domain=-0.5:1] table[x=B,y=bound] {figures/data/15-bias-decay.dat};
\addplot[simulation] table[x=B,y=fw] {figures/data/15-bias-decay-fw.dat};
\addplot[muted,line width=0.4pt,densely dotted] coordinates {(10,0) (10,0.42)};
\node[font=\scriptsize,text=muted,anchor=south west] at (axis cs:10.5,0.005) {$\lambda=0.1$};
\legend{$\gamma_\infty$,exact end point,bound $\gamma_\infty-\lambda\log n$,AdamW without memory}
\end{axis}
\end{tikzpicture}
\caption{Where AdamW without memory ends on the data of \cref{fig:bias:margins}: the $\ell_\infty$ margin of the minimizer of the exponential loss over $\{\norm\theta_\infty\le B\}$ (curve, constrained minimization), the lower bound of \cref{prop:bias:decay} (dashed), and AdamW with $\beta_1=\beta_2=0$, $\epsilon=0$ and $\eta_t\lambda=2/(t+2)$ after $2\cdot10^4$ steps (points). Script \texttt{15-bias-decay.py}.}
\label{fig:bias:decay}
\end{figure}

\paragraph{The example, in numbers.} On the dataset of \cref{ex:bias:wilson} the predictions are exact: test error $1-p$ for sign descent and for Adam with $\epsilon=0$, and $0$ for gradient descent and coordinate descent, for every sample with both classes. For $n_+=6$ and $n_-=4$, scaled so that every training margin is at least $1$, the $\ell_2$ classifier scores a new positive example by $0.93$ and a new negative one by $-0.58$, while the $\ell_\infty$ classifier scores them by $0.75$ and $0.25$: both classify the training set with room to spare, and only one of them reads the label.

\paragraph{Where agreement must fail.} The predictions concern directions, so they fail when the direction is not determined. If $\biasU$ has more than one point, which happens for the $\ell_1$ and $\ell_\infty$ norms, the theorem fixes the margin but not the limit inside $\biasU$, which then depends on the step sizes and even on the floating-point arithmetic (Exercise~\ref{ex:bias:tie}). And the theorem says nothing about the early steps, before $\gamma\tau_t$ exceeds $\log n$, when the iterates are not yet in the regime of their bias.

\section{What the theory misses}\label{sec:bias:misses}

\paragraph{What was idealized away.} \emph{Fixed features.} In a network the features move, and the bias depends on the parametrization as well as on the optimizer: for diagonal linear networks, gradient descent from a small initialization behaves like an $\ell_1$ method and from a large one like an $\ell_2$ method \parencite{WoodworthEtAl2020,MoroshkoEtAl2020}. For homogeneous networks the limit is a KKT point of a margin problem, not a maximizer. \emph{Separability and infinite time.} The theory describes the end of a training run that fits its data exactly and runs long enough for $\gamma\tau_t\gg\log n$. Language-model pretraining makes about one pass over its data and never interpolates it, so the theory does not apply there at all; in vision, training stops after a fixed budget, often not long after the training error reaches zero. \emph{Full batch, no noise.} Minibatch noise changes the trajectory and the time scales. \emph{Memory and $\epsilon$.} Momentum leaves the limit unchanged, but Adam's $\epsilon$ does not (\cref{sec:bias:beyond}). \emph{The example.} The dataset of \cref{ex:bias:wilson} is built to make a point; it shows that an implicit bias can decide the test error, not that it does so on real data.

\paragraph{Where it breaks.} (1) \emph{The generalization gap is not robust.} With every hyperparameter tuned, adaptive methods matched or beat momentum SGD in the study of \textcite{ChoiEtAl2019}, and decoupled weight decay closed most of the gap on image classification \parencite{LoshchilovHutter2019}. On transformers the adaptive methods are preferred, for reasons of optimization speed \parencite{ZhangEtAl2020,KunstnerEtAl2023} (\cref{ch:transformers}). The max-margin mechanism predicts a gap only when the data are interpolated, training runs into the regime of the bias, and the data favour one geometry over another; the comparisons that report gaps do not check these conditions, and the ones that find none do not either. (2) \emph{Adam's $\epsilon$ decides the limit.} With a fixed $\epsilon>0$, Adam's long-run bias is the $\ell_2$ margin, not the $\ell_\infty$ one \parencite{WangEtAl2021}: the $\ell_\infty$ bias is a transient that lasts while the gradients are large compared with $\epsilon$. On separable data the gradients decay like the loss, so the transient ends; whether it ends within a training run depends on $\epsilon$ and on how fast the loss falls (Exercise~\ref{ex:bias:eps}). (3) \emph{KKT is not maximality.} For ReLU networks the limit direction can fail to be even a local maximizer of the margin \parencite{VardiShamirSrebro2022}. (4) \emph{Margins do not reliably rank deep networks.} Spectrally normalized margins separate networks trained on true labels from networks trained on random ones \parencite{BartlettFosterTelgarsky2017}, but in a large-scale comparison of complexity measures, norm- and margin-based measures predicted generalization less well than sharpness-based ones \parencite{JiangEtAl2020}, and \textcite{NagarajanKolter2019} built settings in which every uniform-convergence bound, margin-based or not, is vacuous although the classifier generalizes. (5) \emph{The step-size conditions are too cautious, and the loss matters.} Gradient descent on the logistic loss with an arbitrary constant step size, even one that makes the loss oscillate, still converges in direction to the $\ell_2$ maximum margin, while on the exponential loss large steps can diverge \parencite{WuBravermanLee2023}; and constant normalized steps converge in our experiments where \cref{cor:bias:steps} guarantees only $\gamma-R^2\eta/2$.

\paragraph{Repairs and open problems.} Faster margin rates: \textcite{NacsonEtAl2019} improved the $1/\log t$ rate of gradient descent with aggressive step sizes, \textcite{JiTelgarsky2021} obtained $O(1/t)$ through a primal--dual analysis in which gradient descent induces mirror descent on the example weights, and \textcite{JiSrebroTelgarsky2021} accelerated the dual to obtain $\tilde O(1/t^2)$. Other geometries: mirror descent with $\ell_p$ potentials \parencite{SunEtAl2022}; steepest flows on homogeneous networks \parencite{TsilivisVardiKempe2025}; Adam, AdamW, spectral descent and Muon (\cref{sec:bias:beyond}). Non-separable data: gradient descent converges in a well-defined sense even when the data cannot be separated \parencite{JiTelgarsky2019}. Open: (a) whether constant normalized steps reach the maximum margin in the $\ell_1$, $\ell_\infty$ and spectral norms, and at the rate $1/\tau_t$ observed in \cref{fig:bias:margins}; (b) the implicit bias of Muon, with momentum and with a different norm for each layer, in a deep network; (c) a theory that predicts, from properties of the data, which geometry generalizes better, and so the sign of the difference between Adam and SGD.

\section{Back in the wild}\label{sec:bias:back}

\fieldnotes{Prefer SGD to adaptive methods when test accuracy counts, or switch from Adam to SGD late; partly folklore, with mixed evidence.}
{A linear classifier on separable data, the exponential loss, normalized steepest descent in a norm; weight decay as a norm ball.}
{The normalized margin is at least $\gamma-(\log n+\frac{R^2}2\sum\eta_t^2)/\sum\eta_t$ for every step-size sequence, so steepest descent maximizes the margin of its own norm; the end point with weight decay is within $\lambda\log n$; on a constructed dataset the $\ell_\infty$ bias errs on all negative examples and the $\ell_2$ bias on none.}
{SGD $\to$ $\ell_2$ SVM, sign descent and Adam with negligible $\epsilon$ $\to$ $\ell_\infty$, boosting $\to$ $\ell_1$, orthogonalized updates $\to$ spectral margin; gaps of order $1/\tau_t$ for normalized steps and $1/\log t$ for plain gradient descent.}
{Depth and moving features, noise, Adam's $\epsilon$, finite training, non-interpolating regimes such as language-model pretraining, and the empirical fragility of the generalization gap.}
{An optimizer's norm is a prior on the interpolating solution (proved in the idealization); which prior generalizes depends on the data, in both directions (shown by examples); ``SGD generalizes better'' is folklore whose evidence depends on tuning.}

What should a practitioner take from this? First, when a model can interpolate its training data and is trained far past zero training error, the choice of optimizer is also a choice of solution, and the chapter says which one: the classifier that maximizes the worst-case margin in the optimizer's own norm. That is proved in the idealization, and it is the cleanest sense in which the optimizer matters beyond speed. Second, nothing in the theory says that the $\ell_2$ prior is the better one. The dataset of \cref{ex:bias:wilson} penalizes the $\ell_\infty$ prior, and a dataset with a dense ground truth penalizes the $\ell_2$ prior (Exercise~\ref{ex:bias:dense}); the rule that SGD generalizes better is folklore, supported in some comparisons and removed by tuning in others. Third, the knobs that decide the end point are known: the norm of the update, Adam's $\epsilon$ (which turns the $\ell_\infty$ bias into an $\ell_2$ one once gradients become small), the weight decay (which stops the march at radius $1/\lambda$), and the length of training measured as distance travelled.

The chapter closes the unification line of the book on the side of solutions rather than speed. The norms of \cref{ch:steepest,ch:sign,ch:spectral} reappear as the norms in which margins are maximized; the moving averages of \cref{ch:momentum} leave the limit unchanged; the preconditioners of \cref{ch:adaptive} change it, through the norm they define (Exercise~\ref{ex:bias:precond}); and the norm ball of decoupled weight decay (\cref{ch:decay}) turns the limit direction into a point. The other half of Part V, \cref{ch:transformers}, is about speed, and there the evidence favours Adam.

\begin{exercises}
\item \consolidation (Three margins of two examples.) Let the signed examples be $z_1=(5,1)$ and $z_2=(4,3)$. (a) Show that $\gamma_1=4$ with $\biasU=\{e_1\}$, $\gamma_2=11/\sqrt5$ with $\biasU=\{(2,1)/\sqrt5\}$, and $\gamma_\infty=6$ with $\biasU=\{(1,1)\}$. (b) Verify \cref{lem:bias:duality} by finding the point of the segment $[z_1,z_2]$ closest to the origin in the $\ell_\infty$, $\ell_2$ and $\ell_1$ norms. (c) Which example is a support vector in each geometry?\label{ex:bias:twopoints}

\item \consolidation (Distance in the dual norm.) For $\theta\ne0$ and $x\in V$, show that $\min\{\norm{x-h}_*:\ip{\theta,h}=0\}=\abs{\ip{\theta,x}}/\norm\theta$. Deduce that $\biasbar(\theta)$ is the smallest dual-norm distance from an example to the hyperplane $\{\ip{\theta,\cdot}=0\}$, signed by correctness.\label{ex:bias:distance}

\item \consolidation (Every initialization.) Prove \cref{thm:bias:margin} from an arbitrary $\theta_0$: whenever the right side is below $\gamma$,
\[
  \gamma-\biasbar(\theta_T)\le\frac{(R+\gamma)\norm{\theta_0}+\log n+\frac{R^2}2\sum_{t<T}\eta_t^2}{\norm{\theta_0}+\tau_T}.
\]
\hint{$\psi(\theta_0)\ge-R\norm{\theta_0}-\log n$ and $\norm{\theta_T}\le\norm{\theta_0}+\tau_T$.}\label{ex:bias:init}

\item \consolidation\hard (AdaBoost's margin.) Let every $\abs{z_{ij}}$ equal $1$, as for weak classifiers valued in $\{-1,1\}$, so that $R=1$ in the $\ell_1$ geometry. (a) Show that the step $\eta=\biasartanh(r)$ minimizes the first bound of \cref{lem:bias:step} and that, with it, $\Loss(\theta_{t+1})=\sqrt{1-r_t^2}\,\Loss(\theta_t)$ exactly; this is AdaBoost. (b) Deduce that AdaBoost's normalized $\ell_1$ margin satisfies $\liminf_t\biasbar(\theta_t)\ge-\frac12\log(1-\gamma^2)/\biasartanh(\gamma)$, which is about $\gamma/2$ for small $\gamma$. \hint{the function $r\mapsto-\frac12\log(1-r^2)/\biasartanh(r)$ is increasing on $(0,1)$, and $r_t\ge\gamma$.}\label{ex:bias:adaboost}

\item \consolidation (A tie.) Let $z_1=(2,0)$ and $z_2=(1,2)$. (a) Show that $\gamma_\infty=2$ and that the $\ell_\infty$ max-margin set is $\{(1,s):\frac12\le s\le1\}$. (b) Show that sign descent from $0$ moves exactly along $(1,1)$, so that $\biasbar(\theta_t)=2$ from the first step on. (c) In double precision, $e^{-x}$ underflows to $0$ for $x$ above about $745$. Explain why a run of sign descent with a constant step, which computes the softmax weights in floating point, leaves the ray $(1,1)$ once $\tau_t$ exceeds about $745$, and why its margin stays $2$.\label{ex:bias:tie}

\item \variation (Preconditioning changes the norm.) Let $P\succ0$ be fixed and run normalized steepest descent in the norm $\norm\theta_P=\sqrt{\theta^\top P\theta}$. (a) Show that the step is $-\eta_tP^{-1}g/\norm g_{P^{-1}}$: preconditioned normalized gradient descent. (b) Show that its limit direction is the $\ell_2$ maximum-margin direction of the transformed examples $P^{-1/2}x_i$, mapped back by $P^{-1/2}$. (c) For a diagonal $D\succ0$ and the norm $\norm{D\theta}_\infty$, write the step and the max-margin problem. Conclude that a frozen preconditioner changes the implicit bias, and a moving average does not change it in the sense of \textcite{WangEtAl2022}.\label{ex:bias:precond}

\item \variation (A dense ground truth.) Let the labels of $x\sim\cN(0,I_2)$ be $\sign(x_1+x_2)$ and let the training set be the single example $x=(1,\frac12)$ with label $+1$. (a) Show that $u_\infty=(1,1)$ and $u_2=(2,1)/\sqrt5$. (b) For a classifier $u$, show that the test error is the angle between $u$ and $(1,1)$ divided by $\pi$. (c) Deduce test errors $0$ for the $\ell_\infty$ classifier and $\arccos(3/\sqrt{10})/\pi\approx0.102$ for the $\ell_2$ one.\label{ex:bias:dense}

\item \variation (Multiclass margins.) For $K$ classes and $W\in\R^{K\times d}$, verify that the multiclass margin is $\min_{c\ne y_i}\ip{W,Z_{ic}}$ with $Z_{ic}=(e_{y_i}-e_c)x_i^\top$, that $\norm{Z_{ic}}_{\nuc}=\sqrt2\norm{x_i}_2$ and $\norm{Z_{ic}}_F=\sqrt2\norm{x_i}_2$, and write down the four instances of \cref{cor:bias:instances} for $W$ (Frobenius, entrywise $\ell_\infty$, entrywise $\ell_1$, spectral). Which of them has a limit that is unique for generic data?\label{ex:bias:multiclass}

\item \variation (Coupled weight decay.) Let $\theta_\lambda$ minimize $\Loss(\theta)+\frac\lambda2\norm\theta_2^2$. (a) Show that $\theta_\lambda$ minimizes $\Loss$ over the ball of radius $\norm{\theta_\lambda}_2$, so that \cref{prop:bias:decay} applies with $B=\norm{\theta_\lambda}_2$. (b) By comparing with $Bu^\star$ for $B=\log(1/\lambda)/\gamma_2$, show that $\Loss(\theta_\lambda)\to0$ as $\lambda\to0$, hence $\norm{\theta_\lambda}_2\to\infty$ and $\biasbar(\theta_\lambda)\to\gamma_2$.\label{ex:bias:penalty}

\item \simulation (Adam's $\epsilon$ on the example.) On the dataset of \cref{ex:bias:wilson} with $n_+=6$ and $n_-=4$, run Adam with $\beta_1=0.9$, $\beta_2=0.999$, bias correction, steps $0.05/\sqrt{t+1}$ and $\epsilon\in\{0,10^{-8},10^{-3}\}$ for $10^5$ steps from zero, and record the score $(\theta_2+\theta_3-\theta_1)/\norm\theta_\infty$ of a new negative example. Observe that it stays at $1$ (misclassified) for $\epsilon=0$, as \cref{prop:bias:wilson} proves, and for $\epsilon=10^{-8}$ over this horizon, but that for $\epsilon=10^{-3}$ it starts to fall once the loss is small enough that $\sqrt{\hat v_t}$ is comparable to $\epsilon$. Explain why it can become negative, and estimate when.\label{ex:bias:eps}

\item \simulation (The distance is the clock.) Rerun \texttt{15-bias-margins.py} with constant steps $0.2/R$ and $0.01/R$ and with decaying steps $0.3/(R\sqrt{t+1})$. Check that $(\gamma-\biasbar(\theta_t))\,\tau_t$ approaches the same constant for every schedule in each norm, and compare that constant with $\log n$ and with the asymptote $R^2\eta/2$ of the constant-step bound.\label{ex:bias:clock}

\item \simulation (AdaBoost and shrinkage.) Draw random matrices $Z\in\{-1,1\}^{n\times d}$ with $4\le n,d\le12$, keep the separable ones, and run AdaBoost (step $\biasartanh(r_t)$) and coordinate descent with the constant step $0.05$ for $3000$ steps. Compare their final $\ell_1$ margins with $\gamma_1$ from the linear program and with the bound of Exercise~\ref{ex:bias:adaboost}. Observe that AdaBoost sometimes ends visibly below $\gamma_1$, while the small steps do not; this is the phenomenon studied by \textcite{RudinDaubechiesSchapire2004} and \textcite{RatschWarmuth2005}.\label{ex:bias:boost}

\item \research (Constant steps.) Prove that normalized steepest descent with a constant step $\eta$ in the $\ell_\infty$ or $\ell_1$ norm satisfies $\biasbar(\theta_t)\to\gamma$, ideally at the rate $O(1/\tau_t)$ of \cref{fig:bias:margins}, or find a dataset on which the margin stays below $\gamma$. For the $\ell_2$ norm the smoothness of the dual problem gives such rates \parencite{JiTelgarsky2021}; the difficulty for $\ell_1$ and $\ell_\infty$ is that the step jumps between vertices of the ball and never settles.\label{ex:bias:constant}

\item \research (Muon in depth.) For a two-layer linear network $x\mapsto W_2W_1x$ trained with orthogonalized updates applied to each layer (with a shared step length in a max norm over the layers, \cref{ch:modular}), \cref{prop:bias:homogeneous} applies with $a=2$. Characterize the KKT points of the corresponding max-margin problem in terms of the end-to-end matrix $W_2W_1$, compare with gradient descent (whose limits relate to the nuclear norm of $W_2W_1$ in related settings), and determine the effect of momentum. As far as we know the answer is open.\label{ex:bias:muon}
\end{exercises}

\notesandrefs
Maximum margins entered machine learning with the support vector machine \parencite{CortesVapnik1995}, and margins in other norms with boosting. \Textcite{FreundSchapire1997} introduced AdaBoost, \textcite{SchapireEtAl1998} explained its resistance to overfitting by the growth of the $\ell_1$ margins of the training examples, and \textcite{FreundSchapire1999} cast boosting as a repeated zero-sum game played by multiplicative weights, the view of \cref{rem:bias:hedge}. \Textcite{FriedmanHastieTibshirani2000} and \textcite{Friedman2001} read boosting as coordinate descent on the exponential loss; \textcite{RossetZhuHastie2004} showed that boosting with small steps follows a regularization path that ends at the $\ell_1$ maximum margin; \textcite{RudinDaubechiesSchapire2004} found that AdaBoost itself can cycle and stop short of it; \textcite{RatschWarmuth2005} and \textcite{Telgarsky2013} gave step sizes and shrinkage schemes with margin guarantees. Long before gradient descent, boosting was understood as an optimizer whose end point maximizes a margin, the $\ell_1$ one.

\Textcite{SoudryEtAl2018} proved that gradient descent on separable data converges in direction to the $\ell_2$ support vector machine at the rate $1/\log t$, and \textcite{JiTelgarsky2019} extended the picture to non-separable data. \Textcite{GunasekarEtAl2018} asked which geometry an optimizer carries and showed that steepest descent in a norm converges to the maximum margin in that norm; \cref{thm:bias:margin} proves a version of their result for normalized steps, through the Hoeffding step of \cref{lem:bias:step} and the regret argument of \cref{rem:bias:hedge}, which gives an explicit bound for every step-size sequence. Rates were improved by aggressive and normalized steps \parencite{NacsonEtAl2019,JiTelgarsky2021} and by acceleration of the dual \parencite{JiSrebroTelgarsky2021}. The geometric reading of margins in a general norm as distances in the dual norm is due to \textcite{Mangasarian1999}.

For networks, \textcite{LyuLi2020} proved that gradient flow and gradient descent on homogeneous networks increase a smoothed margin and converge to KKT points of the $\ell_2$ max-margin problem, \textcite{JiTelgarsky2020} proved directional convergence, \textcite{VardiShamirSrebro2022} showed that the KKT points need not be local maximizers, and \textcite{TsilivisVardiKempe2025} extended the picture to steepest flows in general norms. For adaptive methods, \textcite{WangEtAl2021} showed that Adam and RMSProp with a fixed $\epsilon$ share the $\ell_2$ bias of gradient descent, \textcite{WangEtAl2022} that momentum does too, \textcite{ZhangZouCao2024} that Adam with negligible $\epsilon$ converges to the $\ell_\infty$ maximum margin, \textcite{XieLi2024} that AdamW solves an $\ell_\infty$-constrained problem, and \textcite{FanSchmidtThrampoulidis2025} that spectral descent and Muon maximize the spectral margin on multiclass data. The view of optimizers as steepest descent in norms that organizes this book is that of \textcite{BernsteinNewhouse2024}. Volume 2 of this series, \emph{The Hydrogen Atoms of Deep Learning}, Chapter 4, treats the Euclidean case for gradient flow.

The empirical debate began with \textcite{WilsonEtAl2017}, was continued by \textcite{KeskarSocher2017}, who switched from Adam to SGD during training, and by \textcite{LoshchilovHutter2019}, whose decoupled weight decay narrowed the gap, and was reframed by \textcite{ChoiEtAl2019}, who found that comparisons follow the inclusion relations among optimizers once all hyperparameters are tuned.

\paragraph{Translating the literature.} \Textcite{GunasekarEtAl2018} write $w$ for our $\theta$ and $\{(x_n,y_n)\}_{n=1}^N$ for the data (our $n$ is their $N$); their steepest descent step $\Delta w_t\in\argmin_v\ip{\nabla\Loss(w_t),v}+\frac12\norm v^2$ is our $-\nabla\Loss^\sharp$, so their method is the unnormalized one of \cref{cor:bias:unnormalized}, and their step size $\eta_t$ corresponds to our normalized step $\eta_t\norm{\nabla\Loss(w_t)}_*$. \Textcite{SoudryEtAl2018} write $\hat w$ for the $\ell_2$ max-margin solution scaled to have minimum margin $1$ (our $u_2/\gamma_2$). \Textcite{LyuLi2020} write $L$ for the order of homogeneity (our $a$; in this book $L$ is a smoothness constant), $\bar\gamma$ for the normalized margin $\biasqmin/\norm\theta_2^a$ and $\tilde\gamma$ for the smoothed one, and use the sum of the losses, so that their condition $\Loss<1$ is our $n\Loss<1$. In the boosting literature the example weights $D_t$ are our $p(\theta_t)$, the edge $1-2\eps_t$ of the chosen weak classifier with weighted error $\eps_t$ is our $r(\theta_t)$, and AdaBoost's step $\alpha_t=\frac12\log((1-\eps_t)/\eps_t)$ is our $\biasartanh(r(\theta_t))$.

\printchapterbib
\end{refsection}


\bookmarksetup{startatroot}
% Epilogue of Optimizers in the Wild.
\specialchapter{Epilogue}{What theory owes practice}{30}\label{ch:epilogue}
\begin{refsection}[chapters/16-epilogue.bib]
% Tables outside the numbered chapters are numbered in one plain sequence: the prologue has
% Tables 1 and 2, so the ledger is Table 3. The redefinition is local to this environment.
\renewcommand{\thetable}{\arabic{table}}\setcounter{table}{2}

The prologue listed the decisions that a practitioner makes before a training run and asked what theory can say about each. This epilogue keeps the account. It records, practice by practice, what the chapters proved and in which idealization, what they predicted, and what they left unexplained; it says what the idealizations miss when they are taken together; and it lists the questions that the book raises and cannot answer.

\addsec{A ledger}

\Cref{tab:epilogue:ledger} has one row for each practice. The verdict is one of three. A practice is \emph{explained} when a theorem in a faithful idealization accounts for it and the prediction survives a comparison with simulation or with the evidence from training runs. It is \emph{partly explained} when a mechanism has been identified and proved in an idealization, but either the idealization is far from training or the prediction is only qualitative. It is \emph{unexplained} when the book has nothing to say beyond a description, or when the practice is folklore. The last column says what the verdict rests on, which is where a reader who doubts it should look.

\begin{table}[p]
\begin{fullwidth}
\centering\footnotesize
\renewcommand{\arraystretch}{1.14}
\begin{tabular}{@{}>{\raggedright\arraybackslash}p{31mm}r>{\raggedright\arraybackslash}p{17mm}>{\raggedright\arraybackslash}p{104mm}@{}}
\toprule
\textsc{practice} & \textsc{ch.} & \textsc{verdict} & \textsc{what the verdict rests on}\\
\midrule
Normalized updates & \ref{ch:steepest} & explained & Every update in use is a steepest-descent step in some norm; rates depend on the norm only through a smoothness constant and a dual norm\\
Sign-like steps (Adam, Lion) & \ref{ch:sign} & explained & Memoryless Adam and Lion take sign steps; with memory Adam's step is at most $C(\beta_1,\beta_2)$ learning rates, $7.27$ at $(0.9,0.999)$ and $1.16$ at $(0.9,0.95)$\\
When $\ell_\infty$ steps help & \ref{ch:sign} & partly & A gradient denser than the effective rank of a curvature that is nearly diagonal in the parameters' coordinates; the comparison is between guarantees\\
Muon's orthogonalization & \ref{ch:spectral} & explained & The steepest step in the spectral norm is the polar factor; Muon's quintic does not converge but holds all but the smallest singular values in $[0.681,1.203]$\\
When Muon helps & \ref{ch:spectral} & partly & Nuclear rank of the gradient against stable rank of the inputs, for one-step guarantees; many steps and the shrinking gains at scale are not explained\\
Width rules: $\mu$P, Muon's shape factor & \ref{ch:modular} & partly & Width exponents follow from an aligned step in the RMS$\to$RMS norm and match measured slopes, but alignment fades; base rates are not derived\\
Momentum & \ref{ch:momentum} & partly & Stability regions and stationary variances are exact on quadratics; normalized steps need memory when the noise is skewed; why $\beta=0.9$ is not explained\\
Iterate and weight averages & \ref{ch:momentum} & partly & For convex losses a uniform average at a constant step matches a decayed step; on quadratics a weight average cools only the stiff directions (derived)\\
Adam's $\beta_2$ & \ref{ch:adaptive} & partly & Rare gradients need a long memory $1/(1-\beta_2)$, which costs an excess move of about $2\ln2/(1-\beta_2)$ learning rates after a quiet spell; the values are not derived\\
Shampoo, K-FAC, SOAP & \ref{ch:kronecker} & partly & The Kronecker bound and Shampoo's regret are proved, but the regret never beats an adaptive gradient step; the best power depends on what the factors estimate; the wins over AdamW are not derived\\
Training at $2/\eta$ & \ref{ch:eos} & partly & Thresholds for gradient descent, momentum, and Adam with a frozen preconditioner are exact on quadratics; self-stabilization is proved in a cubic model, in continuous time; the rise of sharpness is an input\\
Linear warmup & \ref{ch:warmup} & partly & If curvature is bounded by the loss, then a warmup pays the initial curvature once and a constant step at every scale, and the best stable rise is hyperbolic, not linear; the bound is measured, not derived, and the length is not explained\\
Gradient clipping & \ref{ch:spikes} & partly & Under $(L_0,L_1)$-smoothness, measured but not derived, clipping avoids paying for the largest gradient, which no fixed step can; the threshold $1.0$ is folklore\\
z-loss and QK-norm & \ref{ch:spikes} & partly & Each removes a mechanism of instability (a flat direction of the softmax; attention logits quadratic in the weights), though a z-loss of $10^{-4}$ is too weak to pin a unigram model's logits; real spikes have interacting causes\\
Batch-size rules & \ref{ch:noise} & partly & In the noisy quadratic model the best step grows linearly with the batch up to a critical batch size of order $\tr\Sigma/(S\eps)$, which preconditioning raises; Adam's square-root rule is derived, not proved; curvature and noise are frozen\\
Cosine and WSD & \ref{ch:schedules} & partly & Convex last-iterate bounds for every schedule, and an exact river valley with a noise floor and a cooldown law; with tuned peaks the valley ranks cosine well ahead of WSD, while in language models the two match\\
Decoupled weight decay & \ref{ch:decay} & explained & Steepest descent with decoupled decay is Frank--Wolfe on a norm ball; under SGD scale-invariant layers settle at a rotation of $\arccos(1-\eta\lambda)$ per step, as vision runs confirm; momentum shrinks the angle\\
The value $\lambda=0.1$ & \ref{ch:decay} & unexplained & Only the product $\eta\lambda$ and the timescale $1/(\eta\lambda)$ matter; neither value is derived\\
Adam over SGD for transformers & \ref{ch:transformers} & partly & Three solvable mechanisms: heavy-tailed classes slow gradient descent to a logarithmic decay; blocks of curvature at different scales; and a gap that grows with the batch size; their interaction is open\\
Adaptive methods and generalization & \ref{ch:bias} & partly & Steepest descent ends at the maximum margin of its own norm on separable data, a prior that can favor either optimizer; the evidence for SGD is mixed\\
\bottomrule
\end{tabular}
\caption{The practices of the book, with the chapter that treats each, its verdict, and what the verdict rests on.}
\label{tab:epilogue:ledger}
\end{fullwidth}
\end{table}

Three things stand out. First, the update rules themselves are well explained: the first line of the book, that optimizers are steepest descent in norms with preconditioning and moving averages, accounts for the form of every rule in the field guide, with decoupled weight decay read as a Frank--Wolfe step, and it predicts how learning rates should scale with width while updates stay aligned with the features. Second, almost every practice of the second line is partly explained, and in the same way: a mechanism is proved in an idealization, and the idealization contains one or two parameters, a curvature bound, a noise level, a rate of sharpening, that are measured rather than derived. Third, the numbers that practitioners set in the recipes this book quotes, $\beta_1=0.9$, $\beta_2=0.95$, weight decay $0.1$, a clipping threshold of $1.0$, a warmup of two thousand steps, a cooldown of a fifth of training, are unexplained. The theory says which combinations matter and how they should scale; it does not say why these values.

\addsec{What the idealizations miss together}

Each chapter listed what its own idealization misses. Read together, the lists repeat five omissions.

\paragraph{The landscape is fixed.} Almost every idealization in the book runs an optimizer on a landscape given in advance: a quadratic, a convex function, a cubic model of the edge of stability, a valley with a river. Along the trajectory of a network the local landscape changes: the sharpness rises, features form, norms grow and settle, and the gradient noise changes with them. Where the book needed such a change, it put it in by hand, as the rate of progressive sharpening in \cref{ch:eos} or the loss-dependent curvature in \cref{ch:warmup}.

\paragraph{One mechanism at a time.} Warmup is explained in \cref{ch:warmup} by curvature that falls with the loss, by the edge of stability, and by the statistics of Adam's first steps, and in \cref{ch:decay} by the size of early angular updates; the advantage of Adam on transformers by three mechanisms in \cref{ch:transformers}. In a training run they all operate at once, and nothing in the book says how they combine.

\paragraph{Guarantees, not performance.} Several comparisons in the book are comparisons of worst-case guarantees: which norm gives the better descent bound (\cref{ch:steepest,ch:sign,ch:spectral}), which schedule gives the better last-iterate bound (\cref{ch:schedules}). Where the chapters measured how far a guarantee is from what the optimizer does, the answer ranged from nothing, for sign descent on a diagonal quadratic, to more than a factor of ten, for AdaGrad's regret and for normalized momentum (\cref{ch:steepest,ch:adaptive,ch:momentum}). A guarantee for one step can also misrank two methods over a run: on generic targets the first spectral step beats the first gradient step, yet spectral descent can need more than twice as many steps to reach the same loss (\cref{ch:spectral}).

\paragraph{Simple noise.} The noise in most chapters has bounded variance, often Gaussian, often proportional to the curvature. Heavy tails get a theorem only in \cref{ch:spikes}, where clipping restores a rate, and an exercise in \cref{ch:momentum}. The structure of the data enters through rare classes (\cref{ch:sign,ch:spikes}), sparse features (\cref{ch:adaptive}), and the Zipf law of token frequencies (\cref{ch:transformers}), and each time it changed the answer: a sign method without memory ignores a rare class, and gradient descent slows to a logarithmic decay on Zipfian data.

\paragraph{No scale.} A careful benchmark found that the speedups of newer optimizers over a well-tuned AdamW shrink as models grow, for the fastest of them from about $1.4\times$ at $0.1$ billion parameters to about $1.1\times$ at $1.2$ billion \parencite{WenEtAl2026Fantastic}. No idealization in this book has a size parameter that reproduces this. The nearest candidate is the noisy quadratic model of \cref{ch:noise,ch:transformers}: there the advantage of a preconditioner falls from its condition-number factor at large batch sizes to a noise-weighted constant at small ones, and the critical batch size grows as the target loss falls, so that a run with a fixed batch size moves toward the small-batch end as it aims lower. The book does not test whether this is why the gains shrink.

\addsec{What is open}

The chapters end with research exercises. Here are the questions that recur across them, in roughly the order of the book.

\begin{enumerate}
\item \emph{The right norm.} The norm of a network is assembled from norms on its layers, but the norms for embeddings, for attention, for normalization gains, and for the output layer are chosen, not derived (\cref{ch:spectral,ch:modular}). Is there a principle, perhaps from the statistics of activations and gradients, that selects them?
\item \emph{The averaging constants.} Why do $\beta_1=0.9$ and $\beta_2\in[0.95,0.999]$ work so widely (\cref{ch:momentum,ch:adaptive})? The theory explains which timescales these constants set, not why those timescales are right.
\item \emph{Progressive sharpening.} Full-batch gradient descent on a network raises the sharpness until the step is barely stable (\cref{ch:eos}). Why does the sharpness rise in the first place? A model in which progressive sharpening is an output rather than an input would bear on most of Parts III and IV.
\item \emph{Long warmups.} The curvature model of \cref{ch:warmup} sets the length of a warmup by the initial excess curvature, about $K_1\Delta_0/\mu$ steps, which is unknown for a network, and Adam's statistics suggest $2/(1-\beta_2)$ steps, forty at $\beta_2=0.95$, while LLaMA, with that $\beta_2$, warms up for $2000$. What sets the length?
\item \emph{Cooldowns.} The river valley of \cref{ch:schedules} predicts the noise floor that a cooldown removes and how the best cooldown fraction scales with the length of the run, but through an exponent and a spread of the noise that must be measured; with tuned peaks it ranks cosine well ahead of warmup--stable--decay, while in language models the two match. What limits the peak, and what sets a cooldown near a fifth of the run?
\item \emph{Shrinking gains.} Why do the gains of Muon, SOAP, and other matrix methods over AdamW shrink with model size (\cref{ch:spectral,ch:kronecker})? An unrefereed preprint reports that Hyperball, which fixes the norms of weights and updates, gives Muon token-equivalent speedups of $20$ to $30\%$ over weight-decay baselines up to $1.2$ billion parameters \parencite{WenEtAl2026Hyperball}; whether such gains survive at larger scale, and whether the effective learning rate that weight decay sets explains the difference (\cref{ch:decay}), are open.
\item \emph{Hyperparameter scaling laws.} How should the learning rate, the batch size, and the weight decay scale with model size and data size together? The book gives the exponent of width (\cref{ch:modular}), the critical batch size of a quadratic (\cref{ch:noise}), and the timescale of weight decay (\cref{ch:decay}); a joint theory is open.
\item \emph{Spikes.} A solvable model of a loss spike in a transformer, with the optimizer state, the data, and the architecture all playing their parts, is open (\cref{ch:spikes}).
\item \emph{Beyond pretraining.} The optimizers of this book are tuned for pretraining. Fine-tuning and reinforcement learning are not treated, mixtures of experts appear only through experts that a router stops choosing (\cref{ch:spikes}), and low precision only through rounding in bfloat16 (\cref{ch:spectral,ch:spikes}); each raises questions that the book does not answer.
\end{enumerate}

\addsec{How to turn a practice into a theorem}

The chapters followed one method, and it may be useful to state it. Begin with the practice as practitioners state it, with its numbers and its sources, and with the evidence for and against it; the first section of every chapter is an attempt at this, and the attempt often showed that the practice was folklore. Find the smallest idealization in which the practice has a reason: one coordinate for sign descent, one matrix for Muon, a cubic for the edge of stability, a valley with a river for schedules, a sphere for weight decay. Prove what the idealization implies, completely, so that its consequences are not in doubt. Extract a number from the proof: a threshold, an exponent, a rate, a timescale, an equilibrium. Compare the number with a simulation, and then with training runs, and report the disagreement. Finally, list what the idealization threw away, and look in that list for the next idealization.

The method has a limit that the ledger makes plain: it explains mechanisms one at a time, in models that are much smaller than the networks they describe. Whether the mechanisms add up to an account of large-scale training is the question this book leaves open.

\notesandrefs

The benchmark of \textcite{WenEtAl2026Fantastic} compares ten optimizers on language models of $0.1$ to $1.2$ billion parameters, trained on one to eight times the Chinchilla-optimal amount of data, with careful tuning of each, and finds that the speedups over AdamW shrink with scale; its follow-up \parencite{WenEtAl2026Hyperball}, an unrefereed preprint, proposes fixing the norms of weights and updates directly. Several practices of the ledger, among them warmup, clipping, the batch size, and the schedule, are discussed from the practitioner's side in the \emph{Deep Learning Tuning Playbook} \parencite{GodboleEtAl2023}. The individual chapters give the sources for each entry of the ledger.

\printchapterbib
\end{refsection}


\appendix
% Appendix A of Optimizers in the Wild: the shared mathematical toolkit.
\casechapter{A Toolkit}{Duality, matrices, recurrences and online learning}{}{ch:toolkit}
\begin{refsection}[chapters/17-toolkit.bib]

This appendix collects the mathematics that the chapters share, in the form in which they use it, with a proof wherever one fits in a page; it is a reference, not a chapter to be read in order. Norms and duality (\cref{sec:toolkit:duality}) are the language of \cref{ch:steepest,ch:sign,ch:spectral,ch:modular} and return in \cref{ch:momentum,ch:decay,ch:bias}. Matrix analysis (\cref{sec:toolkit:matrices}) supplies the polar factor of \cref{ch:spectral} and the L\"owner order, Kronecker products, and geometric mean behind the preconditioners of \cref{ch:adaptive,ch:kronecker}. Linear recurrences (\cref{sec:toolkit:recurrences}) give the stability thresholds of \cref{ch:momentum,ch:eos}, the noise floors of \cref{ch:momentum,ch:noise,ch:schedules}, and, through Gr\"onwall's inequality, the local smoothness bounds of \cref{ch:spikes}. The descent lemma and the Polyak--\L{}ojasiewicz inequality (\cref{sec:toolkit:smooth}) start the convergence proofs of the prologue and of \cref{ch:steepest,ch:sign,ch:warmup}; online convex optimization (\cref{sec:toolkit:online}) is the setting of \cref{ch:adaptive,ch:kronecker,ch:schedules}; and martingale and tail bounds (\cref{sec:toolkit:prob}) control the noise in \cref{ch:momentum} and the sign errors in \cref{ch:sign}.

\paragraph{Conventions.} Vectors live in $\R^d$. A linear functional, such as a gradient, is identified with a vector $z$ through the pairing $\ip{z,x}=\sum_iz_ix_i$,\sidenote{The identification uses the pairing, not the norm. Changing the norm changes which step is steepest; it does not change the gradient.} and matrices carry the pairing $\ip{A,B}=\tr(A^\top B)$. An unadorned $\norm{\cdot}$ is an arbitrary norm. For symmetric matrices, $A\preceq B$ means that $B-A$ is positive semidefinite, written $B-A\succeq0$, and $A\prec B$ that it is positive definite.

\section{Norms and duality}\label{sec:toolkit:duality}

A norm measures steps. A gradient turns a step into a rate of change of the loss, and the dual norm measures it by the largest rate it offers per unit step.

\begin{definition}[Dual norm]\label{def:toolkit:dual}
The \term{dual norm}\index{dual norm} of a norm $\norm{\cdot}$ on $\R^d$ is $\norm{z}_*=\max_{\norm{x}\le1}\ip{z,x}$. The maximum exists because the unit ball is closed and bounded, hence compact.
\end{definition}

One fact from convex geometry is used twice, here and as the engine of online gradient descent (\cref{thm:toolkit:ogd}). For $P\succ0$ write $\norm{x}_P=\sqrt{x^\top Px}$.

\begin{lemma}[Projection onto a convex set]\label{lem:toolkit:projection}
Let $P\succ0$ and let $\cK\subseteq\R^d$ be nonempty, closed, and convex. Every $y\in\R^d$ has a unique nearest point $\Pi(y)=\argmin_{x\in\cK}\norm{x-y}_P$, and for every $u\in\cK$,
\[
  \ip{P(y-\Pi(y)),u-\Pi(y)}\le0,\qquad\norm{\Pi(y)-u}_P^2\le\norm{y-u}_P^2-\norm{y-\Pi(y)}_P^2 .
\]
\end{lemma}

\begin{proof}
A minimizer $p$ exists because $x\mapsto\norm{x-y}_P$ is continuous and, for a fixed $x_0\in\cK$, the points of $\cK$ with $\norm{x-y}_P\le\norm{x_0-y}_P$ form a compact set. For $u\in\cK$ and $s\in(0,1]$ the point $p+s(u-p)$ lies in $\cK$, so $0\le\norm{y-p-s(u-p)}_P^2-\norm{y-p}_P^2=-2s\ip{P(y-p),u-p}+s^2\norm{u-p}_P^2$; dividing by $s$ and letting $s\downarrow0$ gives the first inequality. Expanding $\norm{y-u}_P^2=\norm{(y-p)+(p-u)}_P^2$ and using the first inequality gives the second, and the second with $u=p'$ shows that every minimizer $p'$ equals $p$.
\end{proof}

\begin{lemma}[Dual norms]\label{lem:toolkit:dualnorm}
Let $\norm{\cdot}$ be a norm on $\R^d$ and $\norm{z}_*=\max_{\norm{x}\le1}\ip{z,x}$ its dual (\cref{def:toolkit:dual}).
\begin{enumerate}
\item The dual norm is a norm, and \term{H\"older's inequality}\index{H\"older's inequality} $\abs{\ip{z,x}}\le\norm{z}_*\norm{x}$ holds for all $z$ and $x$.
\item (Bidual.) For every $x$, $\norm{x}=\max_{\norm{z}_*\le1}\ip{z,x}$.
\item Let $1\le p\le\infty$ and $1/p+1/q=1$, with $q=\infty$ for $p=1$ and $q=1$ for $p=\infty$. The dual of $\norm{\cdot}_p$ is $\norm{\cdot}_q$. For $z\ne0$ the maximum of $\ip{z,x}$ over $\norm{x}_p\le1$ is attained at $x=\sign(z)$ if $q=1$, at $x=\sign(z_i)e_i$ with $\abs{z_i}=\norm{z}_\infty$ if $q=\infty$, and at $x_i=\sign(z_i)\abs{z_i}^{q-1}/\norm{z}_q^{q-1}$ if $1<q<\infty$.
\end{enumerate}
\end{lemma}

\begin{proof}
Parts (i) and (iii) are computations; for (ii) the plan is to separate a point of the unit sphere from the unit ball by a hyperplane, obtained as a limit of projections. (i) By the equivalence of norms on $\R^d$, $c\norm{x}_2\le\norm{x}\le C\norm{x}_2$ for some $0<c\le C$, so $\norm{z}_2/C\le\norm{z}_*\le\norm{z}_2/c$. The dual norm is positively homogeneous by definition, even because the unit ball is symmetric, and subadditive because a maximum of a sum is at most the sum of the maxima. Applying the definition to $\pm x/\norm{x}$ gives H\"older's inequality.

(ii) H\"older's inequality gives $\max_{\norm{z}_*\le1}\ip{z,x}\le\norm{x}$; for the reverse we may assume $\norm{x}=1$. Let $B$ be the unit ball, let $p_k$ be the Euclidean projection of $y_k=(1+1/k)x\notin B$ onto $B$ (\cref{lem:toolkit:projection} with $P=I$), and let $z_k=(y_k-p_k)/\norm{y_k-p_k}_2$. By the lemma $\ip{z_k,u-p_k}\le0$ for $u\in B$, while $\ip{z_k,y_k-p_k}>0$, so $\ip{z_k,u}<\ip{z_k,y_k}$ for every $u\in B$. A subsequence of the unit vectors $z_k$ converges to a unit vector $z$ with $\ip{z,u}\le\ip{z,x}$ for all $u\in B$, that is, $0<\norm{z}_*\le\ip{z,x}$; so $z/\norm{z}_*$ has dual norm $1$ and pairs with $x$ to at least $1=\norm{x}$.

(iii) For $p\in\{1,\infty\}$, $\abs{\ip{z,x}}\le\sum_i\abs{z_i}\abs{x_i}$, which is at most $\norm{z}_\infty\norm{x}_1$ and at most $\norm{z}_1\norm{x}_\infty$, and the stated vectors attain these bounds. For $1<p<\infty$, the concavity of the logarithm gives Young's inequality $ab\le a^p/p+b^q/q$ for $a,b\ge0$; summed over the coordinates of $x$ and $z$ with $\norm{x}_p=\norm{z}_q=1$ it gives $\abs{\ip{z,x}}\le1$, hence $\abs{\ip{z,x}}\le\norm{z}_q\norm{x}_p$ by homogeneity. Since $(q-1)p=q$, the stated $x$ has $\norm{x}_p^p=\sum_i\abs{z_i}^q/\norm{z}_q^q=1$ and $\ip{z,x}=\norm{z}_q$.
\end{proof}

For a convex function $f$ on $\R^d$, a vector $z$ is a subgradient of $f$ at $x$ if $f(y)\ge f(x)+\ip{z,y-x}$ for all $y$; the set $\partial f(x)$ of subgradients is the \term{subdifferential}\index{subdifferential}. A convex function is differentiable at $x$ if and only if $\partial f(x)$ has a single element, which is then $\nabla f(x)$ \parencite[Section~25]{Rockafellar1970}.

\begin{lemma}[The subdifferential of a norm and the linear minimization oracle]\label{lem:toolkit:subdiff}
For every $x\in\R^d$, $\partial\norm{\cdot}(x)=\{z:\norm{z}_*\le1,\ \ip{z,x}=\norm{x}\}$. In particular $\partial\norm{\cdot}(0)$ is the dual unit ball, and every subgradient at $x\ne0$ has dual norm $1$. Applied to the dual norm, this says that for every $g\in\R^d$
\begin{equation}\label{eq:toolkit:lmo}
  \argmin_{\norm{u}\le1}\ip{g,u}=-\partial\norm{\cdot}_*(g)=\{u:\norm{u}\le1,\ \ip{g,u}=-\norm{g}_*\}.
\end{equation}
Hence every \term{linear minimization oracle}\index{linear minimization oracle} $\lmo(g)$ satisfies $\ip{g,\lmo(g)}=-\norm{g}_*$, and $\norm{\lmo(g)}=1$ when $g\ne0$, while every point of the unit ball is a minimizer when $g=0$. The oracle is unique exactly when $\norm{\cdot}_*$ is differentiable at $g$, and then $\lmo(g)=-\nabla\norm{\cdot}_*(g)$.
\end{lemma}

\begin{proof}
If $\norm{z}_*\le1$ and $\ip{z,x}=\norm{x}$, then $\norm{x}+\ip{z,y-x}=\ip{z,y}\le\norm{y}$ for every $y$, so $z$ is a subgradient. Conversely, for a subgradient $z$ the subgradient inequality at $y=0$ and at $y=2x$ gives $\ip{z,x}=\norm{x}$, after which it reads $\ip{z,y}\le\norm{y}$ for all $y$, that is, $\norm{z}_*\le1$; and at $x\ne0$, $\norm{x}=\ip{z,x}\le\norm{z}_*\norm{x}$ forces $\norm{z}_*=1$. By \cref{lem:toolkit:dualnorm}(ii) the norm $\norm{\cdot}$ is the dual of $\norm{\cdot}_*$, so the first part applied to $\norm{\cdot}_*$ at $g$ identifies $\partial\norm{\cdot}_*(g)$ with the maximizers of $\ip{g,u}$ over the unit ball; since the ball is symmetric, the minimizers are their negatives. The last sentence is the criterion for differentiability quoted above.
\end{proof}

The \term{convex conjugate}\index{convex conjugate} of a function $f$ on $\R^d$ is $f^*(z)=\sup_x\{\ip{z,x}-f(x)\}$: at a gradient, the best first-order gain of a step minus its cost.

\begin{lemma}[The conjugate of half a squared norm, and the sharp operator]\label{lem:toolkit:conjugate}
The conjugate of $f(x)=\frac12\norm{x}^2$ is $f^*(z)=\frac12\norm{z}_*^2$, and the supremum defining $f^*(z)$ is attained exactly at the points $x=-\norm{z}_*u$ with $u\in\argmin_{\norm{u}\le1}\ip{z,u}$. Consequently the \term{sharp operator}\index{sharp operator}
\[
  g^\sharp\in\argmax_{u\in\R^d}\bigl\{\ip{g,u}-\tfrac12\norm{u}^2\bigr\}
\]
is $g^\sharp=-\norm{g}_*\lmo(g)$, one value for each choice of the oracle; every value satisfies $\norm{g^\sharp}=\norm{g}_*$ and $\ip{g,g^\sharp}=\norm{g}_*^2$; and for every $L>0$,
\begin{equation}\label{eq:toolkit:quadmin}
  \min_{u\in\R^d}\Bigl\{\ip{g,u}+\frac L2\norm{u}^2\Bigr\}=-\frac{\norm{g}_*^2}{2L},\qquad\text{attained at }u=-\frac{g^\sharp}{L}.
\end{equation}
\end{lemma}

\begin{proof}
For every $x$, H\"older's inequality and $ab\le\frac12a^2+\frac12b^2$ give $\ip{z,x}-\frac12\norm{x}^2\le\norm{z}_*\norm{x}-\frac12\norm{x}^2\le\frac12\norm{z}_*^2$, with equality in both steps if and only if $\norm{x}=\norm{z}_*$ and $\ip{z,x}=\norm{z}_*\norm{x}$. If $z=0$, then $x=0$. If $z\ne0$, write $x=-\norm{z}_*u$; the conditions become $\norm{u}=1$ and $\ip{z,u}=-\norm{z}_*$, which by \cref{lem:toolkit:subdiff} describe the minimizers of $\ip{z,u}$ over the unit ball, the set \eqref{eq:toolkit:lmo}. Such minimizers exist, so the supremum is attained and equals $\frac12\norm{z}_*^2$; the claims about $g^\sharp$ are the case $z=g$. For \eqref{eq:toolkit:quadmin}, the substitution $u=-v/L$ gives $\ip{g,u}+\frac L2\norm{u}^2=-\frac1L\bigl(\ip{g,v}-\frac12\norm{v}^2\bigr)$, whose minimum is $-\frac1{2L}\norm{g}_*^2$, attained at $v=g^\sharp$.
\end{proof}

When the oracle is not unique, $g^\sharp$ stands for each maximizer; all of them have the same norm and the same pairing with $g$, which is all that the convergence proofs use. In the Euclidean norm $g^\sharp=g$, and preconditioning is the same construction for a quadratic norm; \cref{tab:toolkit:duality} lists the duality maps of the norms used in this book.

\begin{example}[Quadratic norms and preconditioning]\label{exa:toolkit:quadnorm}
Let $P\succ0$. The dual of $\norm{x}_P$ is $\norm{z}_{P^{-1}}$, $g^\sharp=P^{-1}g$, and for $g\ne0$ the oracle is unique, $\lmo(g)=-P^{-1}g/\norm{g}_{P^{-1}}$. Indeed $\ip{z,x}=\ip{P^{-1/2}z,P^{1/2}x}\le\norm{z}_{P^{-1}}\norm{x}_P$ by the Cauchy--Schwarz inequality, with equality exactly when $x$ is a nonnegative multiple of $P^{-1}z$. So preconditioned gradient descent $\theta\leftarrow\theta-\eta P^{-1}g$ is unnormalized steepest descent in $\norm{\cdot}_P$: choosing a preconditioner is choosing a quadratic norm (\cref{ch:adaptive}).
\end{example}

\begin{table}[t]
\centering\small
\renewcommand{\arraystretch}{1.18}
\begin{tabular}{@{}llll@{}}
\toprule
norm of a step & dual norm & an oracle $\lmo(g)$ & sharp operator $g^\sharp$\\
\midrule
$\norm{u}_2$ & $\norm{g}_2$ & $-g/\norm{g}_2$ & $g$\\
$\norm{u}_\infty$ & $\norm{g}_1$ & $-\sign(g)$ & $\norm{g}_1\sign(g)$\\
$\norm{u}_1$ & $\norm{g}_\infty$ & $-\sign(g_i)\,e_i$ & $g_i\,e_i$\\
$\norm{u}_P$ & $\norm{g}_{P^{-1}}$ & $-P^{-1}g/\norm{g}_{P^{-1}}$ & $P^{-1}g$\\
$\norm{U}_{\op}$ & $\norm{G}_{\nuc}$ & $-\polar(G)$ & $\norm{G}_{\nuc}\polar(G)$\\
$\norm{U}_F$ & $\norm{G}_F$ & $-G/\norm{G}_F$ & $G$\\
$\norm{U}_{\nuc}$ & $\norm{G}_{\op}$ & $-u_1v_1^\top$ & $\sigma_1(G)\,u_1v_1^\top$\\
\bottomrule
\end{tabular}
\caption{Duality maps, for $g\ne0$ and $G\ne0$. In the $\ell_1$ row, $\abs{g_i}=\norm{g}_\infty$; in the nuclear row, $(u_1,v_1)$ is a top singular pair of $G$. The oracle is unique except in four cases, where the table lists one choice: $\ell_\infty$ when some $g_i=0$, $\ell_1$ when the largest $\abs{g_i}$ occurs at two or more coordinates, the spectral norm when $G$ does not have full rank, and the nuclear norm when $\sigma_1(G)=\sigma_2(G)$. Proofs: \cref{lem:toolkit:dualnorm}(iii) and \cref{ch:steepest} for vectors, \cref{exa:toolkit:quadnorm} for $\norm{\cdot}_P$, and \cref{thm:toolkit:vonneumann} with \cref{ch:spectral} for matrices.}
\label{tab:toolkit:duality}
\end{table}

\section{Matrix analysis}\label{sec:toolkit:matrices}

\paragraph{Singular values.} Let $G\in\R^{m\times n}$ and $k=\min\{m,n\}$. A singular value decomposition $G=\sum_{i=1}^k\sigma_iu_iv_i^\top$ has $\sigma_1\ge\dots\ge\sigma_k\ge0$ and orthonormal families $(u_i)$ in $\R^m$ and $(v_i)$ in $\R^n$; it exists by the spectral theorem for $G^\top G$ \parencite{HornJohnson2013}, and $\sigma(G)=(\sigma_1,\dots,\sigma_k)$ is the vector of singular values. A function acts on a symmetric $A=\sum_i\lambda_ia_ia_i^\top$, a spectral decomposition, by $f(A)=\sum_if(\lambda_i)a_ia_i^\top$, a polynomial in $A$ that does not depend on the decomposition; this defines $A^s$ for $A\succeq0$ and $s\ge0$, with $A^0=I$, and for $A\succ0$ and all real $s$, and $A^{1/2}$ is the only positive semidefinite square root of $A$ \parencite{HornJohnson2013}. Keeping the $r=\rank G$ terms with $\sigma_i>0$ gives the compact decomposition $G=U\diag(\sigma)V^\top$ with $U^\top U=V^\top V=I_r$, and $\polar(G)=UV^\top$, which does not depend on the decomposition: $\polar(G)=Gf(G^\top G)$ for $f(x)=x^{-1/2}$ ($x>0$) and $f(0)=0$. The \term{Schatten norm}\index{Schatten norm} $\norm{G}_{S_p}=\norm{\sigma(G)}_p$ is the spectral norm for $p=\infty$, the Frobenius norm for $p=2$, and the nuclear norm for $p=1$. A norm is unitarily invariant if $\norm{Q_1GQ_2}=\norm{G}$ for all orthogonal $Q_1,Q_2$; Schatten norms are, and that they are norms is part of the next theorem.

\begin{theorem}[von Neumann's trace inequality and Schatten duality]\label{thm:toolkit:vonneumann}\index{von Neumann's trace inequality}
For all $A,B\in\R^{m\times n}$,
\begin{equation}\label{eq:toolkit:vn}
  \abs{\tr(A^\top B)}\le\sum_{i=1}^k\sigma_i(A)\,\sigma_i(B),
\end{equation}
with equality when $A=\sum_ia_iu_iv_i^\top$ and $B=\sum_ib_iu_iv_i^\top$ are singular value decompositions with the same vectors. Consequently, for $1\le p\le\infty$ and $1/p+1/q=1$, $\max_{\norm{B}_{S_p}\le1}\tr(A^\top B)=\norm{A}_{S_q}$: the Schatten functions are norms, and the dual of $\norm{\cdot}_{S_p}$ is $\norm{\cdot}_{S_q}$. In particular the spectral and nuclear norms are dual to each other, the Frobenius norm is its own dual, and $\tr(A^\top\polar(A))=\norm{A}_{\nuc}$ with $\norm{\polar(A)}_{\op}\le1$.
\end{theorem}

\begin{proof}
The plan is to reduce \eqref{eq:toolkit:vn} to an inequality between the two vectors of singular values, weighted by a matrix whose rows and columns sum to at most one, and to prove that inequality by summation by parts. Let $a=\sigma(A)$, $b=\sigma(B)$, and take full decompositions $A=U_1D_aV_1^\top$, $B=U_2D_bV_2^\top$ with orthogonal $U_1,U_2,V_1,V_2$, where $D_a\in\R^{m\times n}$ has $a_i$ in position $(i,i)$ and zeros elsewhere, and likewise $D_b$. With the orthogonal matrices $X=U_1^\top U_2$ and $Y=V_2^\top V_1$, the cyclic property of the trace gives
\[
  \tr(A^\top B)=\tr(D_a^\top XD_bY)=\sum_{i,j\le k}a_iX_{ij}b_jY_{ji},
\]
so $\abs{\tr(A^\top B)}\le\sum_{i,j\le k}a_iw_{ij}b_j$ with $w_{ij}=\frac12(X_{ij}^2+Y_{ji}^2)\ge\abs{X_{ij}Y_{ji}}$. Rows and columns of orthogonal matrices are unit vectors, so every row and every column of $(w_{ij})_{i,j\le k}$ sums to at most $1$. Put $a_{k+1}=b_{k+1}=0$, $\alpha_l=a_l-a_{l+1}\ge0$ and $\gamma_l=b_l-b_{l+1}\ge0$, so that $a_i=\sum_{l\ge i}\alpha_l$ and $b_j=\sum_{l\ge j}\gamma_l$. Exchanging the order of summation,
\[
  \sum_{i,j\le k}a_iw_{ij}b_j=\sum_{l,l'\le k}\alpha_l\gamma_{l'}\sum_{i\le l,\,j\le l'}w_{ij}\le\sum_{l,l'\le k}\alpha_l\gamma_{l'}\min\{l,l'\}=\sum_{i\le k}a_ib_i ,
\]
because the inner sum is at most $l$ by the row sums and at most $l'$ by the column sums, and because the last identity is the first one with $w$ replaced by the identity. With common singular vectors, $\tr(A^\top B)=\sum_{i,j}a_ib_j(u_i^\top u_j)(v_j^\top v_i)=\sum_ia_ib_i$.

For the duality, \eqref{eq:toolkit:vn} and H\"older's inequality for vectors bound $\tr(A^\top B)$ by $\norm{\sigma(A)}_q\norm{\sigma(B)}_p$. Let $A=\sum_ia_iu_iv_i^\top\ne0$, and let $c$ maximize $\ip{\sigma(A),c}$ over the unit ball of $\ell_p$, as in \cref{lem:toolkit:dualnorm}(iii), with $c=e_1$ if $p=1$. Its entries are nonnegative and nonincreasing, so $B=\sum_ic_iu_iv_i^\top$ has $\sigma(B)=c$ and $\tr(A^\top B)=\norm{\sigma(A)}_q$. Hence $\norm{A}_{S_q}$ is the maximum of the linear functions $A\mapsto\tr(A^\top B)$ over $\{\norm{\sigma(B)}_p\le1\}$, so it is subadditive, hence a norm, and it is the dual of $\norm{\cdot}_{S_p}$. For $p=\infty$, $c=\sign(\sigma(A))$ and $B=\polar(A)$.
\end{proof}

Applied to the spectral norm, the theorem says that $-\polar(G)$ is a steepest unit step for a matrix gradient $G$, which is the update of Muon (\cref{ch:spectral}).\sidenote{\textcite{vonNeumann1937} proved the inequality for square matrices, in the paper that identified the unitarily invariant norms with the symmetric gauge functions of the singular values.} When $G$ has full rank, the polar factor is also the semi-orthogonal matrix closest to $G$.

\begin{theorem}[The polar factor]\label{thm:toolkit:polar}\index{polar decomposition}
Let $G\in\R^{m\times n}$ have rank $r$ and compact decomposition $G=U\diag(\sigma)V^\top$.
\begin{enumerate}
\item $G=\polar(G)(G^\top G)^{1/2}=(GG^\top)^{1/2}\polar(G)$. The matrix $\polar(G)$ has $r$ singular values equal to $1$ and the others $0$; it is semi-orthogonal ($\polar(G)^\top\polar(G)=I_n$ if $m\ge n$, $\polar(G)\polar(G)^\top=I_m$ if $m\le n$) if and only if $r=\min\{m,n\}$, and then $\polar(G)=G(G^\top G)^{-1/2}$ for $m\ge n$ and $\polar(G)=(GG^\top)^{-1/2}G$ for $m\le n$.
\item Let $m\ge n$. Every $Q\in\R^{m\times n}$ with $Q^\top Q=I_n$ satisfies
\[
  \norm{G-Q}_F^2\ge\sum_{i=1}^r(\sigma_i-1)^2+(n-r),
\]
with equality if and only if $QV=U$. If $r=n$, then $\polar(G)$ is the unique nearest semi-orthogonal matrix in the Frobenius norm; if $r<n$, then the nearest ones are the semi-orthogonal $Q$ with $QV=U$. For $m<n$ the same holds with $QQ^\top=I_m$, $U^\top Q=V^\top$, and $m-r$ for $n-r$.
\item The semi-orthogonal matrices nearest to $G$ in the Frobenius norm are also nearest in every unitarily invariant norm.
\end{enumerate}
\end{theorem}

\begin{proof}[Proof of (i) and (ii)]
The plan is to read (i) off the compact decomposition and to reduce (ii) to a trace that the Cauchy--Schwarz inequality bounds term by term. (i) From $G^\top G=V\diag(\sigma^2)V^\top$ we get $(G^\top G)^{1/2}=V\diag(\sigma)V^\top$ and $\polar(G)(G^\top G)^{1/2}=UV^\top V\diag(\sigma)V^\top=G$; the other identity is the same computation with $GG^\top$. The expression $UV^\top$ is a compact decomposition with $r$ singular values $1$, and $\polar(G)^\top\polar(G)=VV^\top$, the projection onto the column span of $V$, equals $I_n$ exactly when $r=n$; then $G(G^\top G)^{-1/2}=U\diag(\sigma)V^\top V\diag(\sigma)^{-1}V^\top=UV^\top$. The case $m\le n$ is the transpose.

(ii) Since $\tr(Q^\top Q)=n$, $\norm{G-Q}_F^2=\norm{G}_F^2-2\tr(G^\top Q)+n$ with $\tr(G^\top Q)=\sum_{i\le r}\sigma_iu_i^\top Qv_i$. Because $\norm{Qv}_2^2=v^\top Q^\top Qv=\norm{v}_2^2$, the Cauchy--Schwarz inequality gives $u_i^\top Qv_i\le1$, with equality if and only if $Qv_i=u_i$. Hence $\norm{G-Q}_F^2\ge\sum_i\sigma_i^2-2\sum_i\sigma_i+n$, which is the bound, with equality if and only if $QV=U$. Such $Q$ exist: complete $(v_i)_{i\le r}$ to an orthonormal basis $(v_i)_{i\le n}$ of $\R^n$ and $(u_i)_{i\le r}$ to an orthonormal family $(u_i)_{i\le n}$ in $\R^m$, possible as $m\ge n$, and put $Q=\sum_{i\le n}u_iv_i^\top$. If $r=n$, then $V$ is orthogonal and $QV=U$ gives $Q=UV^\top$. For $m<n$ apply this to $G^\top$.
\end{proof}

Part (iii) is due to \textcite{FanHoffman1955} for square matrices; \textcite[Chapter~8]{Higham2008} treats the rectangular case. It is not used in this book.

\paragraph{The L\"owner order.} The \term{L\"owner order}\index{L\"owner order} is preserved by congruence: $A\preceq B$ implies $CAC^\top\preceq CBC^\top$ for every $C$ of compatible size, with the converse when $C$ is invertible. A symmetric $N$ satisfies $N\preceq I$ if and only if $\lmax(N)\le1$, and for every matrix $X$, $\norm{X}_{\op}\le1$ if and only if $X^\top X\preceq I$, if and only if $XX^\top\preceq I$. A function $f$ on $[0,\infty)$ is \term{operator monotone} if $0\preceq A\preceq B$ implies $f(A)\preceq f(B)$ for symmetric matrices of every size.

\begin{theorem}[L\"owner--Heinz inequality]\label{thm:toolkit:loewner}\index{L\"owner--Heinz inequality}
Let $0\preceq A\preceq B$ and $0\le r\le1$. Then $A^r\preceq B^r$, and if $A\succ0$, then also $B^{-r}\preceq A^{-r}$. For every $r>1$ the first statement fails for some pair; for $r=2$, $A=\bigl(\begin{smallmatrix}1&1\\1&1\end{smallmatrix}\bigr)\preceq B=\bigl(\begin{smallmatrix}2&1\\1&1\end{smallmatrix}\bigr)$, but $B^2-A^2=\bigl(\begin{smallmatrix}3&1\\1&0\end{smallmatrix}\bigr)$ has determinant $-1$.
\end{theorem}

\begin{proof}
The plan is to turn each inequality $A^s\preceq B^s$ into an operator-norm bound, to derive the case $r=\frac12$ from the case $r=1$, and to obtain every $r\in[0,1]$ by repeating that step. Assume first that $B\succ0$. For $s\ge0$, congruence with $B^{-s/2}$ shows that $A^s\preceq B^s$ is equivalent to $B^{-s/2}A^sB^{-s/2}\preceq I$, that is, to $\norm{A^{s/2}B^{-s/2}}_{\op}\le1$. Let $\cS$ be the set of $s\in[0,1]$ for which $A^s\preceq B^s$ for all such pairs, of every size; it contains $0$ and $1$. If $r,s\in\cS$, then $X=A^{r/2}B^{-r/2}$ and $Y=A^{s/2}B^{-s/2}$ have operator norm at most $1$, and so does $X^\top Y=B^{-r/2}A^{(r+s)/2}B^{-s/2}$. This matrix is similar to $N=B^{(r-s)/4}(X^\top Y)B^{(s-r)/4}=B^{-(r+s)/4}A^{(r+s)/2}B^{-(r+s)/4}$, so the eigenvalues of the positive semidefinite $N$ are at most $1$, and $N\preceq I$, which by congruence is $A^{(r+s)/2}\preceq B^{(r+s)/2}$. With $r=0$ and $s=1$ this is the case $\frac12$: $\norm{A^{1/2}B^{-1/2}}_{\op}\le1$ gives $B^{-1/4}A^{1/2}B^{-1/4}\preceq I$. By induction $\cS$ contains the dyadic rationals of $[0,1]$. For $r\in(0,1]$ and positive dyadic $r_j\to r$, $x^{r_j}\to x^r$ for all $x\ge0$, so $A^{r_j}\to A^r$ and $B^{r_j}\to B^r$, and $r\in\cS$ as the positive semidefinite cone is closed. If $B$ is singular, apply the result to $A\preceq B+\eps I$ and let $\eps\downarrow0$.

For the inverses, let $0\prec A\preceq B$, so that $X=A^{1/2}B^{-1/2}$ has $\norm{X}_{\op}\le1$. Then $XX^\top=A^{1/2}B^{-1}A^{1/2}\preceq I$, and congruence with $A^{-1/2}$ gives $B^{-1}\preceq A^{-1}$; applied to $A^r\preceq B^r$, this gives $B^{-r}\preceq A^{-r}$. In the example, $B-A=\diag(1,0)$. The case $r=2$ rules out every $r>1$: if $A^r\preceq B^r$ held for all pairs with $B\succ0$, then so would $A^{r^j}\preceq B^{r^j}$ for all $j$ (apply it repeatedly), and halving, which never used $s\le1$, and continuity would give every exponent in $[0,r^j]$, which contains $2$ for large $j$.
\end{proof}

So $t\mapsto t^r$ is operator monotone for $0\le r\le1$ and $t\mapsto t^{-r}$ reverses the order; this keeps the square roots and fourth roots of AdaGrad and Shampoo monotone as gradients accumulate (\cref{ch:adaptive,ch:kronecker}).\sidenote{The inequality is due to \textcite{Lowner1934}, who characterized all operator monotone functions, and to \textcite{Heinz1951}. The proof by repeated halving is one of several; \textcite[Chapter~V]{Bhatia1997} gives the general theory.}

\paragraph{Kronecker products.} For $A\in\R^{m\times n}$ and $B\in\R^{p\times q}$, the \term{Kronecker product}\index{Kronecker product} $A\otimes B\in\R^{mp\times nq}$ is the block matrix with blocks $A_{ij}B$, and $\vect\colon\R^{m\times n}\to\R^{mn}$ stacks the columns, $\vect(X)=(X_{11},\dots,X_{m1},X_{12},\dots,X_{mn})^\top$.

\begin{lemma}[Kronecker products and vectorization]\label{lem:toolkit:kron}
For matrices of compatible sizes:
\begin{enumerate}
\item $(A\otimes B)(C\otimes D)=(AC)\otimes(BD)$, $(A\otimes B)^\top=A^\top\otimes B^\top$, and $\tr(A\otimes B)=\tr(A)\tr(B)$; in particular $A\otimes I$ and $I\otimes B$ commute, with product $A\otimes B$.
\item With column stacking,
\[
  \vect(AXB)=(B^\top\otimes A)\vect(X),\qquad\vect(uw^\top)=w\otimes u,
\]
and $\ip{\vect(X),\vect(Y)}=\tr(X^\top Y)$. For symmetric $A\in\R^{m\times m}$ and $B\in\R^{n\times n}$, the matrix $B\otimes A$ acts on $\vect(X)$, $X\in\R^{m\times n}$, as $X\mapsto AXB$, and $\vect(X)^\top(B\otimes A)\vect(X)=\tr(X^\top AXB)$.
\item If $A=\sum_i\lambda_i(A)a_ia_i^\top$ and $B=\sum_j\lambda_j(B)b_jb_j^\top$ are spectral decompositions, then $A\otimes B=\sum_{i,j}\lambda_i(A)\lambda_j(B)(a_i\otimes b_j)(a_i\otimes b_j)^\top$ is one; so the eigenvalues of $A\otimes B$ are the products $\lambda_i(A)\lambda_j(B)$, and $A\otimes B\succeq0$ when $A,B\succeq0$.
\item For $A,B\succeq0$ and $r>0$, $(A\otimes B)^r=A^r\otimes B^r$; for $A,B\succ0$ this holds for every real $r$.
\item If $0\preceq A\preceq A'$ and $0\preceq B\preceq B'$, then $A\otimes B\preceq A'\otimes B'$.
\end{enumerate}
\end{lemma}

\begin{proof}
The idea is to check each identity on blocks or on rank-one matrices and to extend it by linearity. (i) Block $(i,j)$ of $(A\otimes B)(C\otimes D)$ is $\sum_lA_{il}C_{lj}BD$; the transpose is taken blockwise; the diagonal blocks of $A\otimes B$ are $A_{ii}B$. (ii) Column $l$ of $uw^\top$ is $w_lu$, so $\vect(uw^\top)=w\otimes u$. With $X=\sum_jx_je_j^\top$ we have $\vect(X)=\sum_je_j\otimes x_j$ and $AXB=\sum_j(Ax_j)(B^\top e_j)^\top$, so by the mixed product $\vect(AXB)=\sum_j(B^\top e_j)\otimes(Ax_j)=(B^\top\otimes A)\vect(X)$. The last claims follow from $(B\otimes A)\vect(X)=\vect(AXB^\top)$ and $B=B^\top$. (iii) By the mixed product the vectors $a_i\otimes b_j$ are orthonormal and $(a_ia_i^\top)\otimes(b_jb_j^\top)=(a_i\otimes b_j)(a_i\otimes b_j)^\top$; expand $A\otimes B$ bilinearly. (iv) Apply $x\mapsto x^r$ to (iii), using $(xy)^r=x^ry^r$ for $x,y\ge0$. (v) $A'\otimes B'-A\otimes B=(A'-A)\otimes B'+A\otimes(B'-B)$, a sum of products of positive semidefinite matrices.
\end{proof}

Under column stacking, the matrix that multiplies $X$ from the right comes first: Shampoo's step $L^{-1/4}GR^{-1/4}$, with $L\in\R^{m\times m}$ accumulating $GG^\top$ and $R\in\R^{n\times n}$ accumulating $G^\top G$, is $(R^{-1/4}\otimes L^{-1/4})\vect(G)$.\sidenote{Some papers stack rows. With $\mathrm{rvec}(X)=\vect(X^\top)$ the identity reads $\mathrm{rvec}(AXB)=(A\otimes B^\top)\,\mathrm{rvec}(X)$, and the factors trade places; check the convention before comparing formulas from different sources.}

\paragraph{The matrix geometric mean.} The geometric mean of two positive numbers is the solution $x>0$ of $x\,a^{-1}x=b$. For matrices that do not commute, $AB$ is not even symmetric, and the right generalization keeps the equation. It rests on Schur complements.

\begin{lemma}[Schur complements]\label{lem:toolkit:schur}
If $A\succ0$ and $C$ is symmetric, then $\bigl(\begin{smallmatrix}A&X\\X^\top&C\end{smallmatrix}\bigr)\succeq0$ if and only if $C\succeq X^\top A^{-1}X$.
\end{lemma}

\begin{proof}
For the invertible matrix $E=\bigl(\begin{smallmatrix}I&-A^{-1}X\\0&I\end{smallmatrix}\bigr)$, a direct multiplication gives $E^\top\bigl(\begin{smallmatrix}A&X\\X^\top&C\end{smallmatrix}\bigr)E=\diag\bigl(A,\,C-X^\top A^{-1}X\bigr)$, and congruence by an invertible matrix preserves positive semidefiniteness in both directions.
\end{proof}

\begin{theorem}[The matrix geometric mean]\label{thm:toolkit:geomean}\index{matrix geometric mean}
For positive definite $A,B\in\R^{n\times n}$ let $A\#B=A^{1/2}\bigl(A^{-1/2}BA^{-1/2}\bigr)^{1/2}A^{1/2}$.
\begin{enumerate}
\item $A\#B$ is the unique positive definite solution $X$ of $XA^{-1}X=B$. Hence $A\#B=B\#A$ and $(sA)\#(sB)=s(A\#B)$ for $s>0$.
\item (Variational characterization.) $A\#B$ is the largest symmetric $X$, in the L\"owner order, such that
\[
  \begin{pmatrix}A&X\\X&B\end{pmatrix}\succeq0 :
\]
it satisfies the condition, and every symmetric $X$ that satisfies it has $X\preceq A\#B$. Consequently $A\#B$ is monotone in each argument, and $A\#B\preceq\frac12(A+B)$.
\item (Joint concavity.) For positive definite $A_1,A_2,B_1,B_2$ and $s\in[0,1]$, $A_1\#B_1+A_2\#B_2\preceq(A_1+A_2)\#(B_1+B_2)$ and
\[
  s(A_1\#B_1)+(1-s)(A_2\#B_2)\preceq\bigl(sA_1+(1-s)A_2\bigr)\#\bigl(sB_1+(1-s)B_2\bigr).
\]
\item (Commuting case.) If $AB=BA$, then $A\#B=(AB)^{1/2}=A^{1/2}B^{1/2}$.
\end{enumerate}
\end{theorem}

\begin{proof}
The plan is to derive everything from the equation in (i) and from the block test of \cref{lem:toolkit:schur}. (i) Put $C=A^{-1/2}BA^{-1/2}\succ0$. The matrix $X=A^{1/2}C^{1/2}A^{1/2}$ is positive definite with $XA^{-1}X=A^{1/2}CA^{1/2}=B$. Conversely, if $X\succ0$ solves the equation, then $Y=A^{-1/2}XA^{-1/2}\succ0$ satisfies $Y^2=C$, so $Y=C^{1/2}$ and $X=A\#B$. If $Z=B\#A$, then $ZB^{-1}Z=A$; inverting gives $Z^{-1}BZ^{-1}=A^{-1}$, that is, $ZA^{-1}Z=B$, so $Z=A\#B$. And $sX$ solves the equation for $(sA,sB)$.

(ii) By \cref{lem:toolkit:schur} the condition is $XA^{-1}X\preceq B$, which after congruence with $A^{-1/2}$ reads $Y^2\preceq C$ for the symmetric $Y=A^{-1/2}XA^{-1/2}$. Then $(Y^2)^{1/2}\preceq C^{1/2}$ by \cref{thm:toolkit:loewner} with $r=\frac12$, and $Y\preceq(Y^2)^{1/2}$ because each eigenvalue $y$ of $Y$ satisfies $y\le\abs{y}$; so $Y\preceq C^{1/2}$, and congruence with $A^{1/2}$ gives $X\preceq A\#B$. The choice $X=A\#B$, that is $Y=C^{1/2}$, satisfies the condition. If $A\preceq A'$ and $B\preceq B'$, then adding $\diag(A'-A,B'-B)\succeq0$ shows that $A\#B$ satisfies the condition for $(A',B')$, so $A\#B\preceq A'\#B'$. For every admissible $X$ and every $x$, the vector $(x,-x)$ gives $x^\top Ax-2x^\top Xx+x^\top Bx\ge0$, that is, $X\preceq\frac12(A+B)$.

(iii) The block matrices of (ii) for $(A_1,B_1,A_1\#B_1)$ and $(A_2,B_2,A_2\#B_2)$ are positive semidefinite, hence so is their sum, the block matrix for $(A_1+A_2,B_1+B_2,A_1\#B_1+A_2\#B_2)$; now apply (ii). For $0<s<1$ the second inequality follows from the first, applied to $sA_1,sB_1,(1-s)A_2,(1-s)B_2$, and from (i); for $s\in\{0,1\}$ it is an equality.

(iv) The matrices $A^{1/4}$ and $A^{1/2}$ are polynomials in $A$, and $B^{1/2}$ is a polynomial in $B$, so they commute. Hence $X=A^{1/2}B^{1/2}=A^{1/4}B^{1/2}A^{1/4}$ is positive definite and $XA^{-1}X=A^{1/2}B^{1/2}A^{-1/2}B^{1/2}=B$, so $X=A\#B$ by (i); and $X^2=AB$.
\end{proof}

Part (ii) turns an inequality about square roots, which are hard to handle, into one about a block matrix, which is linear in its entries. The next corollary is the form in which \cref{ch:kronecker} uses it.

\begin{corollary}[Kronecker square roots are jointly concave]\label{cor:toolkit:kronmean}
For positive semidefinite $A_1,\dots,A_T\in\R^{m\times m}$ and $B_1,\dots,B_T\in\R^{n\times n}$,
\[
  \sum_{t=1}^TA_t^{1/2}\otimes B_t^{1/2}\preceq\Bigl(\sum_{t=1}^TA_t\Bigr)^{1/2}\otimes\Bigl(\sum_{t=1}^TB_t\Bigr)^{1/2}.
\]
\end{corollary}

\begin{proof}
Suppose first that all $A_t$ and $B_t$ are positive definite. The matrices $A_t\otimes I$ and $I\otimes B_t$ commute, with product $A_t\otimes B_t$, so by \cref{thm:toolkit:geomean}(iv) and \cref{lem:toolkit:kron}(iv) their geometric mean is
\[
  (A_t\otimes I)\#(I\otimes B_t)=(A_t\otimes B_t)^{1/2}=A_t^{1/2}\otimes B_t^{1/2}.
\]
Summing over $t$ and using the first inequality of \cref{thm:toolkit:geomean}(iii) $T-1$ times gives $\sum_tA_t^{1/2}\otimes B_t^{1/2}\preceq\bigl((\sum_tA_t)\otimes I\bigr)\#\bigl(I\otimes\sum_tB_t\bigr)$, which is the right side by the same computation. In general, replace $A_t$ and $B_t$ by $A_t+\eps I$ and $B_t+\eps I$ and let $\eps\downarrow0$.
\end{proof}

For $T=1$ the corollary is an equality; its content is that $(A,B)\mapsto A^{1/2}\otimes B^{1/2}$ is superadditive, a case of a theorem of \textcite{Ando1979}.\sidenote{\textcite{Ando1979} proved the variational characterization of $A\#B$, its joint concavity, and the joint concavity of $(A,B)\mapsto A^{1-s}\otimes B^s$ for $0\le s\le1$. \textcite[Chapter~4]{Bhatia2007} treats matrix means.} It underlies the comparison, in the analysis of Shampoo, between the Kronecker-factored preconditioner and the full-matrix AdaGrad preconditioner \parencite{GuptaKorenSinger2018}, which \cref{ch:kronecker} proves.

\section{Linear recurrences, Lyapunov equations, and Gr\"onwall's inequality}\label{sec:toolkit:recurrences}

On a quadratic, gradient descent, the heavy-ball method, and Adam with a frozen preconditioner are linear recurrences, one for each eigendirection of the (preconditioned) Hessian, of order one or two and possibly driven by noise. Stability thresholds are statements about the roots of their characteristic polynomials, and noise floors are their stationary variances. The \term{spectral radius}\index{spectral radius} $\rho(A)$ of a square matrix is the largest absolute value of its complex eigenvalues.

\begin{lemma}[Stability of linear recurrences]\label{lem:toolkit:recurrence}
Let $A\in\R^{d\times d}$, and let $a_0$, $a_1$, $a$, and $\beta$ be real numbers.
\begin{enumerate}
\item For every $\eps>0$ there is $C_\eps$ with $\norm{A^t}_{\op}\le C_\eps(\rho(A)+\eps)^t$ for all $t\ge0$. Hence $x_{t+1}=Ax_t$ converges to $0$ from every $x_0$ if and only if $\rho(A)<1$, and then geometrically fast; if $\rho(A)\ge1$, then some real $x_0$ gives iterates that do not converge to $0$, and if $\rho(A)>1$, then some real $x_0$ gives unbounded iterates.
\item (\term{Schur--Cohn--Jury test}\index{Schur--Cohn--Jury test}.) Both roots of $z^2+a_1z+a_0$ lie in the open unit disk if and only if $\abs{a_0}<1$ and $\abs{a_1}<1+a_0$.
\item (Momentum form.) Both roots of $z^2-(1+\beta-a)z+\beta$ lie in the open unit disk if and only if $\abs{\beta}<1$ and $0<a<2+2\beta$. Equivalently, $y_{t+1}=(1+\beta-a)y_t-\beta y_{t-1}$ converges to $0$ from every initial pair $(y_0,y_{-1})$ if and only if these conditions hold.
\end{enumerate}
\end{lemma}

\begin{proof}
(i) The plan is to make $A$ almost diagonal by a change of basis. By Schur's triangularization $A=QRQ^*$ with $Q$ unitary and $R$ upper triangular with the eigenvalues on its diagonal \parencite{HornJohnson2013}. For $D=\diag(1,\delta,\dots,\delta^{d-1})$ with $\delta\in(0,1]$, the matrix $D^{-1}RD$ has entries $\delta^{j-i}R_{ij}$: the diagonal of $R$, and above it entries of absolute value at most $\delta\abs{R_{ij}}$. So $\norm{D^{-1}RD}_{\op}\le\rho(A)+\delta\norm{R}_F\le\rho(A)+\eps$ for $\delta=\min\{1,\eps/\norm{R}_F\}$ (every $\delta$ works if $R=0$). With $S=QD$, $A^t=S(D^{-1}RD)^tS^{-1}$ and $\norm{A^t}_{\op}\le\norm{S}_{\op}\norm{S^{-1}}_{\op}(\rho(A)+\eps)^t\le\delta^{1-d}(\rho(A)+\eps)^t$, the norms being over complex vectors. If $\rho(A)<1$, take $\rho(A)+\eps<1$. Conversely, let $\lambda_1$ be an eigenvalue with $\abs{\lambda_1}=\rho(A)\ge1$ and eigenvector $v=a+ib\ne0$, $a,b$ real. Then $\norm{A^ta}_2+\norm{A^tb}_2\ge\abs{\lambda_1}^t\norm{v}_2$, so one of the real sequences $A^ta$, $A^tb$ does not tend to $0$, and if $\abs{\lambda_1}>1$, then one of them is unbounded.

(ii) Let $p(z)=z^2+a_1z+a_0=(z-z_1)(z-z_2)$. If $\abs{z_1},\abs{z_2}<1$, then $\abs{a_0}=\abs{z_1z_2}<1$, and $p(1)=(1-z_1)(1-z_2)$ and $p(-1)=(1+z_1)(1+z_2)$ are positive, as products of positive real factors or as $\abs{1\mp z_1}^2$ for a conjugate pair; so $1\pm a_1+a_0>0$, which is $\abs{a_1}<1+a_0$. Conversely, let $\abs{a_0}<1$ and $\abs{a_1}<1+a_0$, so that $p(\pm1)>0$. A conjugate pair has $\abs{z_1}^2=z_1z_2=a_0<1$. Real roots $z_1\le z_2$ lie on either side of the vertex $-a_1/2\in(-1,1)$ (as $\abs{a_1}<2$); since $p$ increases on $[-a_1/2,\infty)$ and $p(1)>0=p(z_2)$, $z_2<1$, and symmetrically $z_1>-1$.

(iii) With $a_0=\beta$ and $a_1=-(1+\beta-a)$ the conditions of (ii) are $\abs{\beta}<1$ and $\abs{1+\beta-a}<1+\beta$, that is, $0<a<2+2\beta$. The recurrence is $x_{t+1}=Ax_t$ for $x_t=(y_t,y_{t-1})$ with the companion matrix $A=\bigl(\begin{smallmatrix}1+\beta-a&-\beta\\1&0\end{smallmatrix}\bigr)$, whose characteristic polynomial is $z^2-(1+\beta-a)z+\beta$, so (i) applies.
\end{proof}

Part (iii) is the form used in \cref{ch:momentum,ch:eos}. On a mode of curvature $h$, the heavy-ball method $\theta_{t+1}=\theta_t-\eta h\theta_t+\beta(\theta_t-\theta_{t-1})$ is (iii) with $a=\eta h$, so it is stable if and only if $\abs{\beta}<1$ and $0<\eta h<2+2\beta$. The moving-average form, $m_t=\beta m_{t-1}+(1-\beta)h\theta_t$ and $\theta_{t+1}=\theta_t-\eta m_t$, satisfies
\[
  \theta_{t+1}-\theta_t=\beta(\theta_t-\theta_{t-1})-\eta(1-\beta)h\theta_t ,
\]
which is (iii) with $a=\eta(1-\beta)h$; for $0\le\beta<1$ it is stable if and only if $0<\eta h<(2+2\beta)/(1-\beta)$, a threshold equal to $38$ at $\beta=0.9$.\sidenote{Test (ii) is the quadratic case of the criterion of Schur and Cohn for the roots of a polynomial to lie in the unit disk, which Jury put in the tabular form used in control engineering; his book \parencite{Jury1964} presents it.}

\begin{lemma}[Linear recursions driven by noise]\label{lem:toolkit:lyapunov}\index{Lyapunov equation}
Let $x_{t+1}=A_tx_t+\xi_t$, with $\E\norm{x_0}_2^2<\infty$, deterministic $A_t$, and noise with $\E[\xi_t\mid\cF_t]=0$ and $\E[\xi_t\xi_t^\top]=Q_t$, where $\cF_t$ is generated by $x_0,\xi_0,\dots,\xi_{t-1}$. Let $C_t=\E[x_tx_t^\top]$.
\begin{enumerate}
\item $\E x_{t+1}=A_t\E x_t$ and $C_{t+1}=A_tC_tA_t^\top+Q_t$, and the covariance $C_t-(\E x_t)(\E x_t)^\top$ obeys the same recursion as $C_t$.
\item If $A_t=A$ with $\rho(A)<1$ and $Q_t=Q$, then the discrete Lyapunov equation $\Sigma=A\Sigma A^\top+Q$ has the unique solution $\Sigma_\infty=\sum_{j\ge0}A^jQ(A^\top)^j\succeq0$, and $C_t-\Sigma_\infty=A^t(C_0-\Sigma_\infty)(A^\top)^t\to0$ geometrically fast.
\item (Order one.) If $x_{t+1}=a_tx_t+\xi_t$ is scalar and $q_t=\E\xi_t^2$, then
\[
  \E x_T^2=\Bigl(\prod_{t<T}a_t^2\Bigr)\E x_0^2+\sum_{s<T}q_s\prod_{s<u<T}a_u^2 ;
\]
for constant $\abs{a}<1$ and $q$ this is $a^{2T}\E x_0^2+q(1-a^{2T})/(1-a^2)\to q/(1-a^2)$.
\item (Order two.) Let $y_{t+1}=c_1y_t+c_2y_{t-1}+e_t$ with square-integrable initial values, $\E[e_t\mid\cF_t]=0$ and $\E e_t^2=s^2$ for a filtration with $y_t,y_{t-1}$ measurable for $\cF_t$, and stable coefficients, $\abs{c_2}<1$ and $\abs{c_1}<1-c_2$. Then $\E y_t^2\to\gamma_0$ and $\E[y_ty_{t-1}]\to\gamma_1$, with
\begin{equation}\label{eq:toolkit:ar2}
  \gamma_0=\frac{(1-c_2)\,s^2}{(1+c_2)\bigl((1-c_2)^2-c_1^2\bigr)},\qquad\gamma_1=\frac{c_1\gamma_0}{1-c_2}.
\end{equation}
\end{enumerate}
\end{lemma}

\begin{proof}
The plan is to expand the second moment, in which the noise enters only through its covariance, and to solve the resulting linear recursion. (i) In the expansion of $x_{t+1}x_{t+1}^\top$ the cross terms vanish in expectation, because $x_t$ is $\cF_t$-measurable and so $\E[x_t\xi_t^\top]=\E\bigl[x_t\E[\xi_t\mid\cF_t]^\top\bigr]=0$; the mean and the covariance follow in the same way. (ii) By \cref{lem:toolkit:recurrence}(i), with $\rho(A)+\eps<1$, the terms of the series are bounded in norm by $C_\eps^2(\rho(A)+\eps)^{2j}\norm{Q}_{\op}$, so it converges; it is positive semidefinite and solves the equation. Subtracting the equation from (i) gives $C_{t+1}-\Sigma_\infty=A(C_t-\Sigma_\infty)A^\top$, hence the formula; and two solutions differ by some $Z=AZA^\top=A^tZ(A^\top)^t\to0$. (iii) Unroll (i). (iv) The vector $(y_t,y_{t-1})$ obeys (i) with the companion matrix $A=\bigl(\begin{smallmatrix}c_1&c_2\\1&0\end{smallmatrix}\bigr)$ and $Q=\diag(s^2,0)$; the polynomial $z^2-c_1z-c_2$ has its roots in the open unit disk by \cref{lem:toolkit:recurrence}(ii) with $a_1=-c_1$, $a_0=-c_2$, so (ii) applies. The second row of $A$ is $(1,0)$, so the $(2,2)$ entry of the Lyapunov equation makes both diagonal entries of $\Sigma_\infty$ equal, $\Sigma_\infty=\bigl(\begin{smallmatrix}\gamma_0&\gamma_1\\\gamma_1&\gamma_0\end{smallmatrix}\bigr)$, and its $(1,2)$ and $(1,1)$ entries read $\gamma_1=c_1\gamma_0+c_2\gamma_1$ and $\gamma_0=(c_1^2+c_2^2)\gamma_0+2c_1c_2\gamma_1+s^2$. Substituting the first into the second and using $(1-c_2)(1-c_1^2-c_2^2)-2c_1^2c_2=(1+c_2)\bigl((1-c_2)^2-c_1^2\bigr)$ gives \eqref{eq:toolkit:ar2}; the denominators are positive under the stability conditions.
\end{proof}

\begin{example}[Heavy ball in noise]\label{exa:toolkit:heavyball}
Let $\theta_{t+1}=\theta_t-\eta(h\theta_t+\zeta_t)+\beta(\theta_t-\theta_{t-1})$ with $h>0$, $\E[\zeta_t\mid\cF_t]=0$, and $\E\zeta_t^2=\sigma^2$, for a filtration as in \cref{lem:toolkit:lyapunov}(iv). Part (iv) applies with $c_1=1+\beta-\eta h$, $c_2=-\beta$, and $s^2=\eta^2\sigma^2$, and $(1-c_2)^2-c_1^2=\eta h(2+2\beta-\eta h)$; so for $\abs{\beta}<1$ and $0<\eta h<2+2\beta$,
\[
  \lim_{t\to\infty}\E\theta_t^2=\frac{\eta\sigma^2(1+\beta)}{h(1-\beta)(2+2\beta-\eta h)} .
\]
At $\beta=0$ this is the stationary variance $\eta\sigma^2/(h(2-\eta h))$ of stochastic gradient descent, also given by (iii); to first order in $\eta$, heavy ball has the variance of stochastic gradient descent with step $\eta/(1-\beta)$ (\cref{ch:momentum,ch:noise}).
\end{example}

\paragraph{Differential and difference inequalities.} The last tool, named after \textcite{Gronwall1919}, compares a function that grows at most linearly in itself with the exponential solving the corresponding equation.

\begin{theorem}[Gr\"onwall's inequality]\label{thm:toolkit:gronwall}\index{Gr\"onwall's inequality}
Let $T>0$.
\begin{enumerate}
\item (Differential form.) If $u$ is absolutely continuous on $[0,T]$ and $u'\le bu+a$ almost everywhere, for integrable $a$ and $b$, then with $B(t)=\int_0^tb$,
\[
  u(t)\le u(0)e^{B(t)}+\int_0^ta(s)e^{B(t)-B(s)}\dd s ;
\]
for constants $b>0$ and $a$, $u(t)\le\bigl(u(0)+a/b\bigr)e^{bt}-a/b$.
\item (Integral form.) If $u$, $a$, and $b$ are continuous on $[0,T]$, $b\ge0$, and $u(t)\le a(t)+\int_0^tb(s)u(s)\dd s$ for all $t$, then $u(t)\le a(t)+\int_0^ta(s)b(s)e^{\int_s^tb}\dd s$; if moreover $a$ is nondecreasing, then $u(t)\le a(t)\exp\bigl(\int_0^tb\bigr)$.
\item (Discrete form.) If $b_t\ge0$ and $u_{t+1}\le(1+b_t)u_t+a_t$ for all $t\ge0$, then $u_t\le u_0\prod_{s<t}(1+b_s)+\sum_{s<t}a_s\prod_{s<r<t}(1+b_r)$, which is at most $\bigl(u_0+\sum_{s<t}a_s\bigr)\exp\bigl(\sum_{s<t}b_s\bigr)$ when $u_0$ and all $a_s$ are nonnegative.
\end{enumerate}
\end{theorem}

\begin{proof}
(i) The function $v=ue^{-B}$ is absolutely continuous with $v'=(u'-bu)e^{-B}\le ae^{-B}$ almost everywhere; integrate and multiply by $e^{B(t)}$. For constants, $\int_0^tae^{b(t-s)}\dd s=(a/b)(e^{bt}-1)$. (ii) For $R(t)=\int_0^tbu$, $R'=bu\le ba+bR$ because $b\ge0$, and (i) with $R(0)=0$ bounds $R$; then use $u\le a+R$. If $a$ is nondecreasing, then the integral is at most $a(t)\int_0^tb(s)e^{\int_s^tb}\dd s=a(t)\bigl(e^{\int_0^tb}-1\bigr)$. (iii) Induction on $t$, and $1+b\le e^b$.
\end{proof}

\Cref{ch:spikes} applies the differential form to a function of $s$ built from $\nabla\Loss(\theta+su)$, which is Lipschitz, hence absolutely continuous and differentiable almost everywhere, under a Hessian bound growing with the gradient.

\section{Smoothness and the Polyak--\L{}ojasiewicz inequality}\label{sec:toolkit:smooth}

A differentiable $\Loss\colon\R^d\to\R$ is $L$-smooth with respect to $\norm{\cdot}$, for a constant $L>0$, if $\norm{\nabla\Loss(\theta)-\nabla\Loss(\theta')}_*\le L\norm{\theta-\theta'}$ for all $\theta,\theta'$; the gradient is a functional, so it is measured in the dual norm.

\begin{lemma}[Descent lemma in a general norm]\label{lem:toolkit:descent}\index{descent lemma}
Let $\Loss$ be $L$-smooth with respect to $\norm{\cdot}$.
\begin{enumerate}
\item For all $\theta$ and $u$, $\abs{\Loss(\theta+u)-\Loss(\theta)-\ip{\nabla\Loss(\theta),u}}\le\frac L2\norm{u}^2$; in particular
\begin{equation}\label{eq:toolkit:descent}
  \Loss(\theta+u)\le\Loss(\theta)+\ip{\nabla\Loss(\theta),u}+\frac L2\norm{u}^2 .
\end{equation}
\item If $\Loss$ is bounded below by $\Loss^\star$, then $\norm{\nabla\Loss(\theta)}_*^2\le2L\bigl(\Loss(\theta)-\Loss^\star\bigr)$ for every $\theta$.
\item (Converse under convexity.) If $\Loss$ is convex and differentiable and \eqref{eq:toolkit:descent} holds for all $\theta,u$, then $\Loss$ is $L$-smooth with respect to $\norm{\cdot}$, and for all $\theta,\theta'$,
\[
  \Loss(\theta')\ge\Loss(\theta)+\ip{\nabla\Loss(\theta),\theta'-\theta}+\frac1{2L}\norm{\nabla\Loss(\theta')-\nabla\Loss(\theta)}_*^2,
\]
\[
  \ip{\nabla\Loss(\theta')-\nabla\Loss(\theta),\theta'-\theta}\ge\frac1L\norm{\nabla\Loss(\theta')-\nabla\Loss(\theta)}_*^2 .
\]
In the Euclidean norm the second inequality is called \term{co-coercivity}.
\item (Hessian test.) A twice continuously differentiable $\Loss$ is $L$-smooth with respect to $\norm{\cdot}$ if and only if $u^\top\nabla^2\Loss(\theta)v\le L$ for all $\theta$ and all $\norm{u},\norm{v}\le1$; where $\nabla^2\Loss(\theta)\succeq0$ it suffices to take $u=v$. For a quadratic with Hessian $H\succeq0$ the best constant is $\max_{\norm{u}\le1}u^\top Hu$: $\lmax(H)$ in the Euclidean norm and $\max_{s\in\{\pm1\}^d}s^\top Hs$ in $\norm{\cdot}_\infty$.
\end{enumerate}
\end{lemma}

\begin{proof}
The plan is to integrate along segments for (i) and (iv) and to minimize the right side of \eqref{eq:toolkit:descent} for (ii) and (iii). (i) Since the gradient is continuous,
\[
  \Loss(\theta+u)-\Loss(\theta)-\ip{\nabla\Loss(\theta),u}=\int_0^1\ip{\nabla\Loss(\theta+su)-\nabla\Loss(\theta),u}\dd s ,
\]
and by H\"older's inequality and smoothness the integrand is at most $Ls\norm{u}^2$ in absolute value. (ii) By \eqref{eq:toolkit:quadmin}, the right side of \eqref{eq:toolkit:descent} at $u=-g^\sharp/L$, $g=\nabla\Loss(\theta)$, is $\Loss(\theta)-\norm{g}_*^2/(2L)$, and it is at least $\Loss^\star$.

(iii) We apply the argument of (ii) to $\Loss$ minus its linearization. Fix $\theta$ and let $\phi(y)=\Loss(y)-\ip{\nabla\Loss(\theta),y}$, a convex function with $\nabla\phi(\theta)=0$, hence minimized at $\theta$, which satisfies \eqref{eq:toolkit:descent} with the same $L$. For every $y$, by \eqref{eq:toolkit:quadmin},
\[
  \phi(\theta)\le\min_u\phi(y+u)\le\phi(y)-\frac1{2L}\norm{\nabla\phi(y)}_*^2 ,
\]
which at $y=\theta'$ is the first inequality. Adding it to the same inequality with $\theta$ and $\theta'$ exchanged gives the second, and H\"older's inequality bounds its left side by $\norm{\nabla\Loss(\theta')-\nabla\Loss(\theta)}_*\norm{\theta'-\theta}$, which gives $L$-smoothness.

(iv) For a symmetric $M$, $\norm{Mv}_*=\max_{\norm{u}\le1}u^\top Mv$. If the condition holds, then $\nabla\Loss(\theta')-\nabla\Loss(\theta)=\int_0^1\nabla^2\Loss(\theta+s(\theta'-\theta))(\theta'-\theta)\dd s$ has dual norm at most $L\norm{\theta'-\theta}$; conversely, $\nabla^2\Loss(\theta)v$ is the limit of $(\nabla\Loss(\theta+sv)-\nabla\Loss(\theta))/s$. For $M\succeq0$, the Cauchy--Schwarz inequality gives $u^\top Mv\le\max\{u^\top Mu,v^\top Mv\}$, and the convex function $u\mapsto u^\top Hu$ attains its maximum over the cube at a vertex.
\end{proof}

Without convexity the converse fails: $-\frac12\norm{\theta}_2^2$ satisfies \eqref{eq:toolkit:descent} for every $L>0$, however small. Two consequences of \eqref{eq:toolkit:descent}, with $g=\nabla\Loss(\theta)$, start the convergence theory of \cref{ch:steepest}: the unnormalized step $-g^\sharp/L$, which minimizes its right side, and the normalized step $\eta\lmo(g)$ of length at most $\eta$ (\cref{lem:toolkit:subdiff}) guarantee
\[
\begin{aligned}
  \Loss\bigl(\theta-g^\sharp/L\bigr)&\le\Loss(\theta)-\frac{\norm{g}_*^2}{2L},\\
  \Loss\bigl(\theta+\eta\lmo(g)\bigr)&\le\Loss(\theta)-\eta\norm{g}_*+\frac{L\eta^2}2 .
\end{aligned}
\]
The only traces of the norm in these bounds are the constant $L$ and the dual norm of the gradient.

\begin{lemma}[The Polyak--\L{}ojasiewicz inequality]\label{lem:toolkit:pl}\index{Polyak--\L{}ojasiewicz inequality}
Let $\Loss$ be differentiable with $\Loss^\star=\inf\Loss>-\infty$, let $\Delta(\theta)=\Loss(\theta)-\Loss^\star$ and $\mu>0$. The function $\Loss$ satisfies the PL inequality with constant $\mu$ with respect to $\norm{\cdot}$ if $\frac12\norm{\nabla\Loss(\theta)}_*^2\ge\mu\,\Delta(\theta)$ for all $\theta$.
\begin{enumerate}
\item If a sequence satisfies $\Loss(\theta_{t+1})\le\Loss(\theta_t)-c_t\norm{\nabla\Loss(\theta_t)}_*^2$ with $c_t\ge0$, then $\Delta_{t+1}\le(1-2\mu c_t)\Delta_t$ and $\Delta_T\le\exp\bigl(-2\mu\sum_{t<T}c_t\bigr)\Delta_0$, where $\Delta_t=\Delta(\theta_t)$.
\item If $\Loss$ is moreover $L$-smooth with respect to $\norm{\cdot}$ and not constant, then $\mu\le L$, and $\theta_{t+1}=\theta_t-\eta\nabla\Loss(\theta_t)^\sharp$ with $0<\eta<2/L$ satisfies $\Delta_t\le\bigl(1-2\mu\eta(1-L\eta/2)\bigr)^t\Delta_0$, which is $(1-\mu/L)^t\Delta_0$ at $\eta=1/L$. In the Euclidean norm this is gradient descent.
\item A function that is $\mu$-strongly convex with respect to $\norm{\cdot}$, that is, $\Loss(\theta+u)\ge\Loss(\theta)+\ip{\nabla\Loss(\theta),u}+\frac\mu2\norm{u}^2$ for all $\theta,u$, satisfies the PL inequality with constant $\mu$.
\end{enumerate}
\end{lemma}

\begin{proof}
(i) The PL inequality gives $\norm{\nabla\Loss(\theta_t)}_*^2\ge2\mu\Delta_t$, hence the first claim. Since $\Delta_{t+1}\ge0$, either $\Delta_t=0$ or $1-2\mu c_t\ge0$, so $\Delta_{t+1}\le\max\{0,1-2\mu c_t\}\Delta_t\le e^{-2\mu c_t}\Delta_t$; multiply. (ii) At a point with $\Delta>0$, the PL inequality and \cref{lem:toolkit:descent}(ii) give $2\mu\Delta\le\norm{\nabla\Loss}_*^2\le2L\Delta$. By \eqref{eq:toolkit:descent} with $u=-\eta g^\sharp$ and \cref{lem:toolkit:conjugate}, $\Loss(\theta_{t+1})\le\Loss(\theta_t)-\eta(1-L\eta/2)\norm{g}_*^2$, so (i) applies with $c_t=\eta(1-L\eta/2)\le1/(2L)$, and $1-2\mu c_t\ge1-\mu/L\ge0$. (iii) By strong convexity and \eqref{eq:toolkit:quadmin} with $L=\mu$, $\Loss^\star\ge\inf_u\{\Loss(\theta)+\ip{g,u}+\frac\mu2\norm{u}^2\}=\Loss(\theta)-\norm{g}_*^2/(2\mu)$.
\end{proof}

The PL inequality asks for neither convexity nor a unique minimizer, only that the gradient be large wherever the loss is far from its infimum; least squares satisfies it even when the minimizer is not unique (\cref{ex:toolkit:lsq}). Part (i) applies to every step-size rule that guarantees a decrease proportional to the squared gradient, which is how \cref{ch:warmup} combines it with a curvature bound that grows with the loss.\sidenote{\textcite{Polyak1963} introduced the inequality and proved linear convergence of gradient descent under it; \textcite{Lojasiewicz1963} proved an inequality of this type for real-analytic functions. \textcite{KarimiNutiniSchmidt2016} compared it with other conditions for linear convergence and popularized the name.}

\section{Online convex optimization}\label{sec:toolkit:online}

In online convex optimization a learner picks $x_t$ in a closed convex set $\cK\subseteq\R^d$ at each round $t=1,\dots,T$, then observes a convex function $f_t$ and a subgradient $g_t\in\partial f_t(x_t)$. Its \term{regret} against $u\in\cK$ is $\sum_t(f_t(x_t)-f_t(u))$, which by convexity is at most the linearized regret $\sum_t\ip{g_t,x_t-u}$. The bounds below hold for every sequence of functions, in particular for the stochastic gradients of a training run. Adaptive methods normalize each gradient by the accumulated size of the gradients so far, and the following lemma controls the resulting sums.

\begin{lemma}[Sums of normalized increments]\label{lem:toolkit:adagrad}\index{AdaGrad!sum lemma}
Let $a_1,\dots,a_T\ge0$, $S_0=0$, and $S_t=\sum_{s\le t}a_s$.
\begin{enumerate}
\item $\sum_{t=1}^Ta_t/\sqrt{S_t}\le2\sqrt{S_T}$, where terms with $S_t=0$ are read as $0$.
\item For $\eps>0$, $\sum_{t=1}^Ta_t/(\eps+S_t)\le\log(1+S_T/\eps)$.
\item (Matrix form.) If $G_0\succeq0$ and $G_1,\dots,G_T\succ0$, then $\sum_{t=1}^T\tr\bigl((G_t-G_{t-1})G_t^{-1/2}\bigr)\le2\tr(G_T^{1/2})-2\tr(G_0^{1/2})$. In particular, if $G_t=G_{t-1}+g_tg_t^\top$ and $G_0\succ0$, then $\sum_tg_t^\top G_t^{-1/2}g_t\le2\tr(G_T^{1/2})$.
\end{enumerate}
\end{lemma}

\begin{proof}
(i) If $a_t>0$, then $a_t/\sqrt{S_t}=(\sqrt{S_t}-\sqrt{S_{t-1}})(\sqrt{S_t}+\sqrt{S_{t-1}})/\sqrt{S_t}\le2(\sqrt{S_t}-\sqrt{S_{t-1}})$, and if $a_t=0$ both sides vanish; the sum telescopes. (ii) Since $1/(\eps+x)\ge1/(\eps+S_t)$ on $[S_{t-1},S_t]$, $a_t/(\eps+S_t)\le\log\frac{\eps+S_t}{\eps+S_{t-1}}$, and the sum telescopes. (iii) For $A\succ0$ and $B\succeq0$, the Cauchy--Schwarz inequality for $\ip{X,Y}=\tr(X^\top Y)$ with $X=A^{-1/4}B^{1/2}$ and $Y=A^{1/4}$, followed by $\sqrt{xy}\le\frac12(x+y)$, gives
\[
  \tr(B^{1/2})=\tr(X^\top Y)\le\tr(BA^{-1/2})^{1/2}\tr(A^{1/2})^{1/2}\le\tfrac12\tr(BA^{-1/2})+\tfrac12\tr(A^{1/2}),
\]
that is, $\tr\bigl((A-B)A^{-1/2}\bigr)\le2\tr(A^{1/2})-2\tr(B^{1/2})$. Take $A=G_t$ and $B=G_{t-1}$ and sum.
\end{proof}

The constant $2$ is best possible, as $a_t=1$ gives $\sum_{t\le T}t^{-1/2}\ge2\sqrt{T+1}-2$, and the current term must be included: with $S_{t-1}$ for $S_t$, $a_1=1$ and $a_2=M$ give at least $M$, which no multiple of $\sqrt{S_T}=\sqrt{1+M}$ bounds.\sidenote{The scalar lemma appears in the analysis of self-confident learning rates by \textcite{AuerCesaBianchiGentile2002}. Adaptive per-coordinate steps of this kind were introduced independently by \textcite{McMahanStreeter2010} and by \textcite{DuchiHazanSinger2011}, who named the method AdaGrad and also analyzed its full-matrix version.}

\begin{theorem}[Online gradient descent]\label{thm:toolkit:ogd}\index{online gradient descent}
Let $\cK\subseteq\R^d$ be nonempty, closed, and convex, $f_1,\dots,f_T$ convex with subgradients $g_t\in\partial f_t(x_t)$, $P_1,\dots,P_T\succ0$, $x_1\in\cK$, and $x_{t+1}=\argmin_{x\in\cK}\norm{x-(x_t-P_t^{-1}g_t)}_{P_t}$.
\begin{enumerate}
\item (General preconditioners.) For every $u\in\cK$, with $P_0=0$,
\[
\begin{aligned}
  \sum_{t=1}^T\bigl(f_t(x_t)-f_t(u)\bigr)&\le\sum_{t=1}^T\ip{g_t,x_t-u}\\
  &\le\frac12\sum_{t=1}^T(x_t-u)^\top(P_t-P_{t-1})(x_t-u)\\
  &\quad+\frac12\sum_{t=1}^Tg_t^\top P_t^{-1}g_t .
\end{aligned}
\]
\item (Step sizes.) If $P_t=I/\eta_t$ with $\eta_1\ge\dots\ge\eta_T>0$ and $\cK$ has Euclidean diameter $D$, then for every $u\in\cK$
\[
  \sum_{t=1}^T\bigl(f_t(x_t)-f_t(u)\bigr)\le\frac{D^2}{2\eta_T}+\frac12\sum_{t=1}^T\eta_t\norm{g_t}_2^2 .
\]
If $\norm{g_t}_2\le G$, the constant step $\eta=D/(G\sqrt T)$ gives $DG\sqrt T$ and $\eta_t=D/(G\sqrt t)$ gives $\frac32DG\sqrt T$. With $S_t=\sum_{s\le t}\norm{g_s}_2^2$, the adaptive step $\eta_t=D/\sqrt{2S_t}$ gives $\sqrt2\,D\sqrt{S_T}$, and the constant step $D/\sqrt{S_T}$, which needs $S_T$ in advance, gives $D\sqrt{S_T}$.
\item (Weighted form.) If $P_t=I/\eta_t$ with arbitrary $\eta_t>0$, then for every $u\in\cK$, with $\cK$ possibly unbounded,
\[
  \sum_{t=1}^T\eta_t\bigl(f_t(x_t)-f_t(u)\bigr)\le\frac12\norm{x_1-u}_2^2-\frac12\norm{x_{T+1}-u}_2^2+\frac12\sum_{t=1}^T\eta_t^2\norm{g_t}_2^2,
\]
and the corresponding bound on $\sum_t\eta_t\ip{g_t,x_t-u}$ is an equality when $\cK=\R^d$.
\end{enumerate}
\end{theorem}

\begin{proof}
The plan is to bound the progress toward $u$ in one round and to sum. By \cref{lem:toolkit:projection} with $P=P_t$ and $y=x_t-P_t^{-1}g_t$, $\norm{x_{t+1}-u}_{P_t}^2\le\norm{y-u}_{P_t}^2=\norm{x_t-u}_{P_t}^2-2\ip{g_t,x_t-u}+g_t^\top P_t^{-1}g_t$, that is,
\begin{equation}\label{eq:toolkit:oneround}
  \ip{g_t,x_t-u}\le\tfrac12\bigl(\norm{x_t-u}_{P_t}^2-\norm{x_{t+1}-u}_{P_t}^2\bigr)+\tfrac12g_t^\top P_t^{-1}g_t .
\end{equation}
(i) Summing \eqref{eq:toolkit:oneround} and regrouping the squared distances by the point $x_t$ instead of the matrix, $\sum_t(\norm{x_t-u}_{P_t}^2-\norm{x_{t+1}-u}_{P_t}^2)=\sum_t(x_t-u)^\top(P_t-P_{t-1})(x_t-u)-\norm{x_{T+1}-u}_{P_T}^2$; drop the last term, and use $f_t(x_t)-f_t(u)\le\ip{g_t,x_t-u}$.

(ii) Here $(x_t-u)^\top(P_t-P_{t-1})(x_t-u)=(1/\eta_t-1/\eta_{t-1})\norm{x_t-u}_2^2\le(1/\eta_t-1/\eta_{t-1})D^2$, with $1/\eta_0=0$, and the sum telescopes to $D^2/\eta_T$. For the constant step both terms are $\frac12DG\sqrt T$. For $\eta_t=D/(G\sqrt t)$ the terms are $\frac12DG\sqrt T$ and $\frac12DG\sum_tt^{-1/2}\le DG\sqrt T$, by \cref{lem:toolkit:adagrad}(i) with $a_t=1$. For the adaptive step both terms are at most $D\sqrt{S_T/2}$, by \cref{lem:toolkit:adagrad}(i) with $a_t=\norm{g_t}_2^2$; rounds with $S_t=0$ have $g_t=0$, leave the iterate in place and contribute nothing, so we read $1/\eta_t=0$ for them. The last choice makes both terms $\frac12D\sqrt{S_T}$.

(iii) With $P_t=I/\eta_t$, multiply \eqref{eq:toolkit:oneround} by $\eta_t$ and sum; the distances telescope whatever the steps. When $\cK=\R^d$ there is no projection, and the first inequality of the proof is an equality.
\end{proof}

The preconditioners pay for their growth through the first sum in (i) and are rewarded through the second. The pieces of this appendix combine as follows for full-matrix AdaGrad, and in the same way for Shampoo in \cref{ch:kronecker}.

\begin{remark}[Full-matrix AdaGrad]\label{rem:toolkit:fullmatrix}
Let $\eps>0$, $G_t=\eps I+\sum_{s\le t}g_sg_s^\top$, and $P_t=G_t^{1/2}/\eta$ on a set $\cK$ of Euclidean diameter $D$. By \cref{thm:toolkit:loewner}, $P_t-P_{t-1}\succeq0$, so $(x_t-u)^\top(P_t-P_{t-1})(x_t-u)\le D^2\tr(P_t-P_{t-1})$, which sums to $D^2\tr(P_T)$; and $\sum_tg_t^\top P_t^{-1}g_t\le2\eta\tr(G_T^{1/2})$ by \cref{lem:toolkit:adagrad}(iii). So \cref{thm:toolkit:ogd}(i) bounds the regret by $\bigl(\frac{D^2}{2\eta}+\eta\bigr)\tr(G_T^{1/2})$, which is $\sqrt2\,D\tr(G_T^{1/2})$ at $\eta=D/\sqrt2$.
\end{remark}

Averaging turns regret bounds into bounds for stochastic optimization \parencite{CesaBianchiConconiGentile2004}; part (iii) allows an arbitrary schedule.

\begin{corollary}[Stochastic gradient descent on a convex function]\label{cor:toolkit:sgd}
Let $\Loss$ be convex on a closed convex $\cK$, $x_1\in\cK$, and $x_{t+1}=\argmin_{x\in\cK}\norm{x-(x_t-\eta_tg_t)}_2$ with deterministic $\eta_t>0$, where $\E[g_t\mid g_1,\dots,g_{t-1}]\in\partial\Loss(x_t)$ and $\E\norm{g_t}_2^2\le G^2$. Then for every $u\in\cK$ and $\bar x_T=\sum_t\eta_tx_t/\sum_t\eta_t$,
\[
  \E\Loss(\bar x_T)-\Loss(u)\le\frac{\norm{x_1-u}_2^2+G^2\sum_{t\le T}\eta_t^2}{2\sum_{t\le T}\eta_t}.
\]
\end{corollary}

\begin{proof}
Apply \cref{thm:toolkit:ogd}(iii) to the linear functions $f_t(x)=\ip{g_t,x}$. Since $x_t$ is a function of $g_1,\dots,g_{t-1}$, conditioning on them and the subgradient inequality give $\E\ip{g_t,x_t-u}\ge\E[\Loss(x_t)-\Loss(u)]$. Take expectations and use Jensen's inequality $\Loss(\bar x_T)\le\sum_t\eta_t\Loss(x_t)/\sum_t\eta_t$.
\end{proof}

\section{Martingales and tail bounds}\label{sec:toolkit:prob}

Gradient noise is a sequence of martingale differences: given the past, a stochastic gradient is unbiased. The theorem below controls sums of such differences; the lemma after it bounds the probability that a noisy estimate has the wrong sign, which is what sign descent pays for noise.

\begin{theorem}[Doob's and the Azuma--Hoeffding inequalities]\label{thm:toolkit:azuma}
Let $(M_t)_{t=0}^T$ be a martingale with respect to a filtration $(\cF_t)_{t=0}^T$.
\begin{enumerate}
\item (\term{Doob's maximal inequality}\index{Doob's maximal inequality}.) If $\E M_T^2<\infty$, then for every $a>0$,
\[
  \Pr\Bigl(\max_{t\le T}\abs{M_t}\ge a\Bigr)\le\frac{\E M_T^2}{a^2},\qquad\E\Bigl[\max_{t\le T}M_t^2\Bigr]\le4\,\E M_T^2 .
\]
The same holds in $\R^d$, with $\abs{\cdot}$ replaced by $\norm{\cdot}_2$, and then $\E\norm{M_T-M_0}_2^2=\sum_{t=1}^T\E\norm{M_t-M_{t-1}}_2^2$.
\item (\term{Azuma--Hoeffding inequality}\index{Azuma--Hoeffding inequality}.) If $\abs{M_t-M_{t-1}}\le c_t$ almost surely for constants $c_t>0$, and $V=\sum_{t=1}^Tc_t^2$, then for every $a>0$,
\[
  \Pr\Bigl(\max_{t\le T}(M_t-M_0)\ge a\Bigr)\le e^{-a^2/(2V)}
\]
and $\Pr(\abs{M_T-M_0}\ge a)\le2e^{-a^2/(2V)}$.
\end{enumerate}
\end{theorem}

\begin{proof}
The plan is to prove a maximal inequality for nonnegative submartingales, processes with $\E[X_{t+1}\mid\cF_t]\ge X_t$, by stopping at the first crossing, and to apply it to $\norm{M_t}_2^2$ and to $e^{\kappa(M_t-M_0)}$. Let $X_t\ge0$ be an integrable submartingale, $X^*=\max_{t\le T}X_t$ and $b>0$. The events $E_t=\{X_t\ge b,\ X_s<b\text{ for }s<t\}$ are disjoint, lie in $\cF_t$ and have union $\{X^*\ge b\}$, and $b\Pr(E_t)\le\E[X_t;E_t]\le\E[X_T;E_t]$. Summing over $t$,
\begin{equation}\label{eq:toolkit:doob}
  b\,\Pr(X^*\ge b)\le\E\bigl[X_T;\,X^*\ge b\bigr].
\end{equation}
(i) By the conditional Jensen inequality, $\norm{M_t}_2^2\le\E[\norm{M_T}_2^2\mid\cF_t]$ is integrable, and $\norm{M_t}_2$ and $\norm{M_t}_2^2$ are nonnegative submartingales. Then \eqref{eq:toolkit:doob} for $\norm{M_t}_2^2$ with $b=a^2$ gives the first bound. For $X_t=\norm{M_t}_2$, \eqref{eq:toolkit:doob} and the Cauchy--Schwarz inequality give $\E[(X^*)^2]=\int_0^\infty2b\Pr(X^*\ge b)\dd b\le2\E[X_TX^*]\le2(\E X_T^2)^{1/2}(\E(X^*)^2)^{1/2}$, and $\E(X^*)^2\le\sum_t\E X_t^2<\infty$, which gives the second. The increments are orthogonal, $\E\ip{M_t-M_{t-1},M_s-M_{s-1}}=0$ for $s<t$, by conditioning on $\cF_{t-1}$.

(ii) For $\kappa>0$ and $\abs{x}\le c$, convexity gives $e^{\kappa x}\le\cosh(\kappa c)+(x/c)\sinh(\kappa c)$, the chord through $x=\pm c$. As $\E[M_t-M_{t-1}\mid\cF_{t-1}]=0$, this yields $\E[e^{\kappa(M_t-M_{t-1})}\mid\cF_{t-1}]\le\cosh(\kappa c_t)\le e^{\kappa^2c_t^2/2}$, comparing the series of $\cosh y$ and $e^{y^2/2}$ term by term, and conditioning successively gives $\E e^{\kappa(M_T-M_0)}\le e^{\kappa^2V/2}$. The nonnegative submartingale $e^{\kappa(M_t-M_0)}$ and \eqref{eq:toolkit:doob} with $b=e^{\kappa a}$ give $\Pr(\max_t(M_t-M_0)\ge a)\le e^{-\kappa a+\kappa^2V/2}$, and $\kappa=a/V$ gives the first bound; the second applies it to $M$ and $-M$.
\end{proof}

For independent summands in intervals centered at their means, (ii) is Hoeffding's inequality \parencite{Hoeffding1963}, which \textcite{Azuma1967} extended to martingale differences. Doob's inequalities are from \textcite{Doob1953}; \textcite{Williams1991} gives textbook proofs, and \textcite{BoucheronLugosiMassart2013} develop the exponential method much further.

\begin{lemma}[Sign errors: Markov, Chebyshev, and Gauss]\label{lem:toolkit:gauss}
Let $Y$ have mean $g\ne0$ and variance $\sigma^2<\infty$, for instance a stochastic estimate of one coordinate $g$ of a gradient.
\begin{enumerate}
\item $\Pr(\sign Y\ne\sign g)\le\Pr(\abs{Y-g}\ge\abs{g})\le\min\{\sigma/\abs{g},\sigma^2/g^2\}$, so that $\abs{g}\Pr(\sign Y\ne\sign g)\le\sigma$. If the law of $Y$ is symmetric about $g$, then the first probability is half the second.
\item (\term{Gauss's inequality}\index{Gauss's inequality}.) Suppose that the law of $Y$ is unimodal about $\nu$: apart from a possible atom at $\nu$, it has a density that is nondecreasing on $(-\infty,\nu)$ and nonincreasing on $(\nu,\infty)$. Let $\tau^2=\E(Y-\nu)^2$. Then for every $k>0$,
\[
  \Pr\bigl(\abs{Y-\nu}\ge k\bigr)\le
  \begin{cases}
    4\tau^2/(9k^2) & \text{if }k\ge2\tau/\sqrt3,\\
    1-k/(\sqrt3\,\tau) & \text{if }k\le2\tau/\sqrt3.
  \end{cases}
\]
\item If $Y$ has a density that is symmetric about $g$ and nonincreasing in $\abs{y-g}$, then $\Pr(\sign Y\ne\sign g)$ is at most $2\sigma^2/(9g^2)$ if $\abs{g}\ge2\sigma/\sqrt3$ and at most $\frac12-\abs{g}/(2\sqrt3\,\sigma)$ otherwise. Consequently $\abs{g}\Pr(\sign Y\ne\sign g)\le\frac{\sqrt3}8\sigma$, with equality for the uniform law at $\abs{g}=\frac{\sqrt3}2\sigma$.
\end{enumerate}
\end{lemma}

\begin{proof}
Part (i) is Markov's and Chebyshev's inequalities, and (iii) combines (i) and (ii); for (ii) the plan is to slice a unimodal density horizontally into uniform densities on intervals that end at the mode, and to check the inequality for each. (i) Let $g>0$, replacing $Y$ by $-Y$ if necessary. A sign error means $Y-g\le-g$, which implies $\abs{Y-g}\ge g$; apply Markov's inequality with $\E\abs{Y-g}\le\sigma$, and Chebyshev's. Under symmetry the disjoint events $\{Y-g\le-g\}$ and $\{Y-g\ge g\}$ have equal probabilities and make up $\{\abs{Y-g}\ge g\}$.

(ii) Let $\nu=0$, and let $p$ be the density away from $0$. For $s>0$ the set $\{y>0:p(y)>s\}$ is an interval with left end $0$, as $p$ is nonincreasing; let $z(s)$ be its length, finite because $s\,z(s)\le1$. Writing $p(y)=\int_0^\infty\ind\{p(y)>s\}\dd s$, Tonelli's theorem gives $\int_0^\infty\phi\,p=\int_0^\infty\int_0^{z(s)}\phi(y)\dd y\dd s$ for measurable $\phi\ge0$, so that
\[
  \Pr(Y\ge k)=\int_0^\infty(z(s)-k)_+\dd s,\qquad\Pr(Y>0)=\int_0^\infty z(s)\dd s,
\]
and $\E[Y^2;Y>0]=\int_0^\infty\frac13z(s)^3\dd s$.
If $\alpha\ge0$ and $\gamma>0$ satisfy $(1-k/z)_+\le\alpha+\gamma z^2$ for all $z>0$, then $(z-k)_+\le\alpha z+\gamma z^3$, and integrating gives $\Pr(Y\ge k)\le\alpha\Pr(Y>0)+3\gamma\E[Y^2;Y>0]$; with the same bound for $-Y$, and since the atom at $0$ contributes to neither side, $\Pr(\abs{Y}\ge k)\le\alpha+3\gamma\tau^2$. Take $z_0=\max\{3k/2,\sqrt3\,\tau\}$, $\gamma=k/(2z_0^3)$ and $\alpha=1-3k/(2z_0)\ge0$. Then $q(z)=\alpha+\gamma z^2$ and $f(z)=1-k/z$ have the same value and slope at $z_0$; as $f$ is concave and $q$ convex, $f\le q$ on $(0,\infty)$, and $q\ge0$, so $(1-k/z)_+\le q(z)$. The bound $\alpha+3\gamma\tau^2=1-3k/(2z_0)+3k\tau^2/(2z_0^3)$ equals $4\tau^2/(9k^2)$ for $z_0=3k/2$, that is, $k\ge2\tau/\sqrt3$, and $1-k/(\sqrt3\,\tau)$ for $z_0=\sqrt3\,\tau$.

(iii) Such a law is unimodal about $g$ with $\tau=\sigma$, and by (i) and (ii) with $k=\abs{g}$ the probability of a sign error is at most half the bound of (ii). With $x=\abs{g}/\sigma$, the function $x(\frac12-x/(2\sqrt3))$ is at most $\frac{\sqrt3}8$, its value at $x=\frac{\sqrt3}2$, and $x\cdot2/(9x^2)\le\frac{\sqrt3}9$ for $x\ge2/\sqrt3$. The uniform law on $[g-\sqrt3\sigma,g+\sqrt3\sigma]$ errs with probability $\frac12-g/(2\sqrt3\sigma)$ for $0<g\le\sqrt3\sigma$, which is $\frac14$ at $g=\frac{\sqrt3}2\sigma$.
\end{proof}

For Gaussian noise, $\max_g\abs{g}\Pr(\sign Y\ne\sign g)$ is about $0.17\sigma$, against the unimodal bound $\frac{\sqrt3}8\sigma\approx0.22\sigma$ and Markov's $\sigma$.\sidenote{\textcite{Gauss1823} proved the inequality in his work on least squares; \textcite{Pukelsheim1994} reviews it and gives an elementary proof of the Vysochanski\u{\i}--Petunin inequality, its analogue for deviations from the mean, and \textcite{BernsteinEtAl2018} use it for signSGD under unimodal symmetric noise.} Since $\E[g\sign Y]\ge\abs{g}\bigl(1-2\Pr(\sign Y\ne\sign g)\bigr)$, part (i) gives $\E\ip{\nabla\Loss,\sign\tilde g}\ge\norm{\nabla\Loss}_1-2\sum_i\sigma_i$ for a stochastic gradient $\tilde g$ with coordinatewise noise levels $\sigma_i$, the bound used in \cref{ch:sign}; under unimodal symmetric noise, (iii) replaces $2\sum_i\sigma_i$ by $\frac{\sqrt3}4\sum_i\sigma_i$.

\begin{exercises}
\item \consolidation Let $\theta=(\theta_1,\dots,\theta_N)$ be split into blocks with norms $\norm{\cdot}_{(l)}$, and $\norm{\theta}=\max_l\norm{\theta_l}_{(l)}$. Show that $\norm{z}_*=\sum_l\norm{z_l}_{(l),*}$ and that an oracle can be computed block by block, $\lmo(g)_l=\lmo_{(l)}(g_l)$. What are the dual norm and an oracle of $\sum_l\norm{\theta_l}_{(l)}$?\label{ex:toolkit:product}
\item \consolidation\hard Show that $g\mapsto g^\sharp$ is single-valued and linear if and only if $\norm{\cdot}=\norm{\cdot}_P$ for some $P\succ0$. \hint{If $g^\sharp=Mg$, then $\norm{g}_*^2=\ip{g,Mg}$ by \cref{lem:toolkit:conjugate}; use \cref{exa:toolkit:quadnorm} and \cref{lem:toolkit:dualnorm}(ii).}
\item \consolidation For $A,B\succ0$, show that $(A\#B)^{-1}=A^{-1}\#B^{-1}$, and deduce $2(A^{-1}+B^{-1})^{-1}\preceq A\#B\preceq\frac12(A+B)$.\label{ex:toolkit:means}
\item \consolidation Show that $\Loss(\theta)=\frac12\norm{X\theta-y}_2^2$ satisfies the PL inequality in the Euclidean norm with $\mu$ the smallest nonzero eigenvalue of $X^\top X$, even when $X^\top X$ is singular. \hint{$\Delta(\theta)=\frac12\norm{\Pi(X\theta-y)}_2^2$ for the orthogonal projection $\Pi$ onto the range of $X$.}\label{ex:toolkit:lsq}
\item \simulation For Gaussian, Laplace, uniform, and Student-$t_3$ noise of variance $1$, compute $\max_{g>0}g\Pr(\sign Y\ne\sign g)$ by quadrature and compare it with $1$, with $\frac12$, and with $\frac{\sqrt3}8$. You should find about $0.170$, $0.130$, $0.217$, and $0.113$; explain why the heavy-tailed law does best at fixed variance.
\end{exercises}

\notesandrefs

\textcite{Rockafellar1970} is the reference for dual norms, subdifferentials, and conjugates. \textcite[Section~2.1]{Nesterov2018} proves the characterizations of smooth convex functions in the Euclidean norm, and \textcite{Beck2017} works with a general norm and its dual, as this book does; the notation $g^\sharp$ follows \textcite{KelnerEtAl2014}. For matrix analysis, \textcite{Bhatia1997} covers unitarily invariant norms, von Neumann's inequality, and operator monotone functions, \textcite{Bhatia2007} matrix means, and \textcite{HornJohnson2013} everything else; \textcite[Chapter~4]{HornJohnson1991} treats Kronecker products and vectorization.

Online gradient descent is due to \textcite{Zinkevich2003}, whose regret bound with $\eta_t\propto t^{-1/2}$ is a case of \cref{thm:toolkit:ogd}(ii). \textcite{Hazan2016} and \textcite{Orabona2019} are modern introductions to online learning; Orabona's notes treat adaptive and parameter-free methods in the language of this appendix. For martingales and concentration see \textcite{Williams1991} and \textcite{BoucheronLugosiMassart2013}.

\printchapterbib
\end{refsection}


\backmatter
\printindex
\specialchapter{How This Book Was Made}{Colophon}{34}

\addsec{Roles}

\emph{Optimizers in the Wild} was written by Claude, an AI model made by Anthropic, in October 2026. Xiaoyu Li conceived the series \emph{The Missing Textbooks} and proposed this volume: its subject, its gap (optimization textbooks stop at convex analysis and stationary points, while practice is decided by Adam, warmup, weight decay, learning-rate schedules, Muon and the edge of stability), its two lines (the unification of optimizers as steepest descent under norms with preconditioning and moving averages, and the dynamics of training), and its position relative to existing work. Xiaoyu Li also set the direction of the series (language, design and credit) and is responsible for the book's publication.

\addsec{Process}

A written plan fixed the thesis, the fifteen chapters, the global notation, the labels of the shared appendix, and a specification for each chapter, including the editor's plan for each main theorem; several of these plans were checked numerically before writing began, and two were corrected as a result. A style guide fixed the anatomy of a chapter, the voice, a proof style that follows Halmos, the distinction between proofs, sketches, heuristic derivations, observations and folklore, and the rule that every chapter compares an exact prediction with a simulation. Separate instances of the same model then drafted the chapters and the appendix in parallel, each re-deriving its results, checking them numerically, verifying its references by web search, and writing the scripts behind its figures. A coordinating instance wrote the prologue, preface, notation, epilogue and this colophon, designed the volume's additions to the series typography, and integrated the whole.

Each chapter then received an independent referee pass from a separate instance, which re-derived the proofs, re-ran the figure scripts, checked constants and exercises numerically, and checked historical claims and references against sources wherever these could be reached. The referees corrected the errors they found directly.

\addsec{Limits}

No human has yet checked every proof. Several results are presented as heuristic derivations, marked with $\Diamond$, because no proof is known to the author or because a proof would not fit; they are checked against simulations but not proved. In the environment in which the book was written, references could be checked only through web search, not by reading the papers themselves, so bibliographic details were included only where a search confirmed them. Results from 2025 and 2026 are cited from preprints where no refereed version could be found, and are marked as such where a claim rests on them. Practices change quickly, and the recipes quoted in the book are dated. Corrections are welcome and will be recorded in the repository of \emph{The Missing Textbooks}.

\medskip
{\small\noindent Typeset with pdf\LaTeX{} in \texttt{missingbook}, the series class built on \textsc{koma}-Script, with the volume style \texttt{wild}; Palatino (newpx) for text and mathematics, biblatex and biber for references, Ti\emph{k}Z and pgfplots for figures, whose data were computed with Python, NumPy, SciPy and PyTorch. The emblem on the title page shows one gradient and its three steepest-descent directions, the points of the $\ell_1$, $\ell_2$ and $\ell_\infty$ unit balls that make the largest inner product with it; the wave at the foot of each opening page rises from chapter to chapter.\par}

\end{document}
